\documentclass[11pt,a4paper]{article}
\usepackage[T1]{fontenc}
\usepackage[utf8]{inputenc}
\usepackage{lmodern}
\usepackage{amsmath,amssymb,amsthm,mathtools}
\usepackage[margin=28mm]{geometry}
\usepackage{microtype}
\usepackage{enumitem}
\usepackage{xcolor}
\usepackage[colorlinks=true,linkcolor=blue!45!black,citecolor=blue!45!black,urlcolor=blue!45!black]{hyperref}
\usepackage{bookmark}
\numberwithin{equation}{section}
\newtheorem{theorem}{Theorem}[section]
\newtheorem{proposition}[theorem]{Proposition}
\newtheorem{lemma}[theorem]{Lemma}
\newtheorem{corollary}[theorem]{Corollary}
\theoremstyle{definition}
\newtheorem{definition}[theorem]{Definition}
\theoremstyle{remark}
\newtheorem{remark}[theorem]{Remark}
\newtheorem{example}[theorem]{Example}
\newcommand{\D}{\mathbb D}
\newcommand{\E}{\mathbb E}
\newcommand{\Pbb}{\mathbb P}
\newcommand{\N}{\mathbb N}
\newcommand{\C}{\mathbb C}
\newcommand{\R}{\mathbb R}
\newcommand{\dd}{\,\mathrm d}
\newcommand{\ind}{\mathbf 1}
\newcommand{\Mp}{\mathcal M_{\mathrm p}}
\newcommand{\supp}{\operatorname{supp}}
\newcommand{\Law}{\mathcal L}
\newcommand{\Cov}{\operatorname{Cov}}
\DeclareMathOperator{\Var}{Var}
\DeclareMathOperator{\Exp}{Exp}
\DeclareMathOperator{\Beta}{Beta}
\DeclareMathOperator{\dist}{dist}
\setlist{itemsep=3pt,topsep=5pt}
\setlength{\emergencystretch}{2em}
\hypersetup{pdftitle={The planar log gas in the disk: limits, conditional Gibbs laws, and hyperbolic geometry},pdfauthor={}}
\title{The planar log gas in the disk\\[3pt]
\large Infinite limits below $2$, empty limits at and above $4$,\\
conditional Gibbs laws, and hyperbolic geometry}
\author{}
\date{Revised September 8, 2026}

\begin{document}
\maketitle
\vspace{-1.5em}
\begin{abstract}
We study the unweighted logarithmic gas in the open unit disk with
density proportional to $\prod_{i<j}|z_i-z_j|^\beta$ at fixed $\beta>0$.
The unscaled point processes are tight for every positive $\beta$.
For every $0<\beta<2$, every subsequential limit is almost surely
infinite and satisfies $\sum_{z\in\Xi}(1-|z|)=\infty$. We prove this
by an exterior-only interpolation identity,
weighted H\"older inequalities, and radial bounds uniform under
holomorphic product weights. Every limiting intensity is bounded
below by a positive multiple of the Bergman intensity, and we obtain
quantitative lower-count tails and almost sure logarithmic growth of
the accumulated radial depth.
At $\beta=2$ the full limit is the Bergman process. For $\beta\ge4$
the full sequence converges to the empty configuration, locally in total
variation, with an explicit rate above $4$.
For every subsequential limit we construct analytic canonical
conditional densities, a meromorphic electric field, and consistent
augmented DLR kernels. The corrected field is independent of the
Gaussian harmonic modes. Exact Gaussian elimination gives a canonical
hyperbolic Green specification with an explicit geometric one-body
factor and a positive-kernel factorization of the unmarked density.
At $\beta=2$ we identify the complete count-resolved compensated DLR law.
An independent perturbation argument identifies the first inverse-temperature
response of every correlation order at $2$, proves its hyperbolic
covariance, and bounds the second temperature derivatives uniformly
in particle number. We now evaluate the second intensity response at every
interior point: its locally uniform limit is $1/[\pi(1-|z|^2)^2]$,
agreeing with the proposed density
$(4/\beta-1)/[\pi(1-|z|^2)^2]$. For $0<\beta<2$, signed central-charge
moments determine the radial-depth coefficient $2/\beta-1/2$ in every
positive moment along a density-one set of particle numbers. The same
coefficient holds at every fixed hyperbolic center. These results do not
prove the proposed local intensity: an explicit order-one mean-potential
defect remains undetermined. Euclidean area coordinates give an exact
screening equation, and a neutral hyperbolic plasma representation identifies
the conjecture with local neutrality. We prove an energy obstruction at the
moving finite-volume boundary and an exact connected-force correction to the
screening identity. For bounded square-integrable one-body tilts of the
Bergman process, we prove a strict nonlinear screening sign and uniqueness
of the associated self-consistency equation; this does not identify a
general-temperature gas limit. A Blaschke interpolation argument reduces
nonemptiness above $2$ to a stated, unproved boundary moment estimate.
We construct infinitely many distinct infinite, rotationally invariant
solutions of the same canonical Green equations for every $0<\beta<4$.
For $\beta\le2$ these states are mutually singular with every actual
disk-gas limit. Canonical Gibbs uniqueness therefore fails, while
uniqueness and full hyperbolic
invariance for $0<\beta<4$, $\beta\ne2$, remain unproved here;
nonemptiness remains unresolved only for $2<\beta<4$.
\end{abstract}

\tableofcontents

\section{Model, conventions, and scope}
Let $\D=\{z\in\C:|z|<1\}$, let $B_r=\{z:|z|<r\}$, and let $dA$ be
ordinary, unnormalized planar Lebesgue measure. We also write $\D_r=B_r$
when using centered polydisks. For $\beta>0$ and $N\ge1$,
put
\begin{equation}\label{eq:model}
 \Pbb_{N,\beta}(\dd z_1\cdots\dd z_N)
 =\frac{1}{Z_{N,\beta}}|\Delta_N(z)|^\beta
   \prod_{i=1}^N\ind_{\D}(z_i)\,dA(z_i),
 \qquad \Delta_N(z)=\prod_{1\le i<j\le N}(z_i-z_j).
\end{equation}
Empty products equal $1$. The partition function is finite because every
pair distance is at most $2$, and is positive because the integrand is
positive off the collision diagonals. There is no external potential in
the interior of the disk. This is the constant-potential disk model in
\cite[Section 3]{Samaj}; the boundary background in that formulation adds
only a configuration-independent energy.

The associated counting measure and disk counts are
\[
 \Xi_N=\sum_{i=1}^N\delta_{z_i},\qquad M_N(r)=\Xi_N(B_r).
\]
We work in $\Mp(\D)$, the space of locally finite, nonnegative,
integer-valued Radon measures, with the vague topology. Multiplicities are
allowed in this state space until simplicity is proved. Vague convergence
means convergence of $\xi(f)=\int f\,d\xi$ for every $f\in C_c(\D)$.
It allows particles to disappear through $\partial\D$. It does not allow
infinitely many particles to accumulate in a compact subset of $\D$.
The closed disk would be an inappropriate state space for the unnormalized
measures, whose total mass is $N$.

We use the standard Polish-space, Prokhorov, and monotone-coupling facts for
counting measures and stochastic order; see \cite{Kallenberg,Strassen}.
All estimates specific to \eqref{eq:model}, including the product association
inequality needed below, are proved here. No claim of bibliographic novelty
is made for the resulting estimates.

\begin{theorem}[Subsequential existence]\label{thm:core}
Fix $\beta>0$. The laws of $(\Xi_N)_{N\ge1}$ are tight in $\Mp(\D)$.
In particular, every sequence of particle numbers tending to infinity has
a subsequence along which $\Xi_N$ converges in distribution to a locally
finite counting measure $\Xi$. Every such limit is simple.
For $0<r<1$ and integers $m\ge2$,
\begin{equation}\label{eq:gaussian-tail}
 \sup_{N\ge1}\Pbb_{N,\beta}\{M_N(r)\ge m\}
 \le \min\left\{1,
  \exp\left(m-\frac{\beta}{4}m(m-1)\log\frac1r\right)\right\}.
\end{equation}
The same bound holds with $M_N(r)$ replaced by $M(r)=\Xi(B_r)$.
All positive moments and all real exponential moments of compact counts
are finite, uniformly before taking the limit.
\end{theorem}

The proof of tightness is in Section~\ref{sec:core-proof}; simplicity and
the passage of count bounds to limits are justified without assuming
simplicity in Section~\ref{sec:hierarchy}. The later sections establish
stronger properties. Throughout, ``a limit'' means a subsequential limit
in distribution unless full-sequence convergence is stated.

\begin{theorem}[Full-sequence uniqueness for $\beta\ge4$]\label{thm:front-uniqueness}
For every fixed $\beta\ge4$, the laws of $\Xi_N$ converge to the law
concentrated on the empty configuration. For $\beta>4$ and each $R<1$,
\[
 \sup_{|z|\le R}\rho_{1,N}(z)
 +\mathbb P\{\Xi_N(\overline{B_R})>0\}
 \le C_{\beta,R}N^{2-\beta/2}(\log(eN))^\beta.
\]
Thus convergence on every relatively compact window holds in total
variation. At $\beta=2$, the full sequence instead converges to the
Bergman determinantal process with kernel
$K(z,w)=1/[\pi(1-z\overline w)^2]$ relative to ordinary area.
\end{theorem}

The assertion for $\beta>4$ is proved in
Theorem~\ref{thm:empty-limit-beta-above-four}. The critical case $4$
is proved in Theorem~\ref{thm:critical-empty}, with an independent
conformal proof in Theorem~\ref{thm:critical-mobius-empty}; no rate
at $4$ is obtained. The assertion at $2$ is proved in
Section~\ref{sec:beta-two}. The next theorem establishes almost sure
infiniteness throughout $0<\beta<2$. Uniqueness remains unresolved
for $0<\beta<4$, $\beta\ne2$, and nonemptiness for $2<\beta<4$.
In particular, we do not identify $4$ as the exact transition parameter.

\begin{theorem}[Infinite subsequential limits below $2$]
\label{thm:front-infinite-below-two}
For every $0<\beta<2$, every subsequential limiting process satisfies
\[
 \mathbb P\{\Xi(\D)=\infty\}=1,\qquad
 \rho_1(z)\ge\frac{c_\beta}{\pi(1-|z|^2)^2},\quad z\in\D,
\]
where $c_\beta>0$ depends only on $\beta$. One may take $c_\beta=1$
for $\beta\le1$. Every factorial correlation has a strictly positive
lower bound after division by its Vandermonde factor. For each fixed
$m\ge0$,
\[
 \mathbb P\{\Xi(B_R)\le m\}
 \le C_{\beta,\tau,m}(1-R^2)^{2\tau/(m\beta+2)},
\]
where $\tau=1$ for $\beta\le1$, and any $0<\tau<2-\beta$ is
allowed for $1<\beta<2$. These conclusions concern every subsequence;
they do not assume uniqueness or hyperbolic invariance.
Moreover, the stronger boundary-mass conclusion is
\[
 \sum_{z\in\Xi}(1-|z|)=\infty\quad\text{almost surely}.
\]
Writing $H(R)=\sum_{|z|<R}-\log|z|$ and
$\kappa_\beta=\min(1,2-\beta)$, one has
\[
 \mathbb E e^{-\beta H(R)}\le C_{\beta,\tau}(1-R^2)^\tau,
 \qquad
 \liminf_{R\uparrow1}\frac{H(R)}{\log(1/(1-R^2))}
 \ge\frac{\kappa_\beta}{\beta}\quad\text{almost surely}.
\]
\end{theorem}
\begin{proof}
The intensity and factorial lower bounds are
Theorems~\ref{thm:nonemptiness-below-two}
and~\ref{thm:nonemptiness-all-orders-below-two}. Almost sure
infiniteness and the lower-count estimate are
Theorem~\ref{thm:almost-sure-infinite-below-two}
and Corollary~\ref{cor:lower-count-tail-below-two}.
The stronger non-Blaschke and depth-growth conclusions are
Theorem~\ref{thm:non-blaschke-limits}.
\end{proof}

\begin{remark}[New results concerning the intensity conjecture]
The conjecture under investigation is
\[
 \rho_\beta(z)=\frac{4/\beta-1}{\pi(1-|z|^2)^2},\qquad 0<\beta<4.
\]
It remains unproved except at $\beta=2$. The new exact perturbative
calculation, Theorem~\ref{thm:beta2-exact-second-all-anchors}, establishes
\[
 \lim_N\left.\partial_\beta^2\log\rho_{N,\beta}(z)\right|_2=0,
 \qquad
 \lim_N\left.\partial_\beta^2\rho_{N,\beta}(z)\right|_2
       =\frac1{\pi(1-|z|^2)^2}.
\]
The convergence is locally uniform in $z$. These limits concern
derivatives of the finite gases at $2$, and do not
interchange the particle-number and temperature limits.
Theorem~\ref{thm:signed-charge-radial-exponent} proves, for $0<\beta<2$,
\[
 \frac{\sum_i-\log|z_i|}{\log N}\longrightarrow\frac2\beta-\frac12
\]
in every positive absolute moment along a set of particle numbers of
natural density one. It also proves the signed-charge free energy for
every charge $t>-2$. Corollary~\ref{cor:all-anchor-depth-coefficient}
extends the depth coefficient to every fixed hyperbolic center.
The exact mean-potential defect in
Proposition~\ref{prop:intensity-mean-potential-defect} vanishes if and
only if the intensity conjecture holds. Theorem~\ref{thm:affine-meromorphic-field-intensity}
gives a sufficient affine-field symmetry criterion, whose symmetry
assumption has not been proved for the selected gas limit. Numerical
checks in Section~\ref{sec:intensity-finite-n-numerics} retain substantial
finite-particle corrections and are not used as proof.
\end{remark}

\begin{remark}[What the screening investigation proves]
Sections~\ref{sec:screening-euclidean-coordinates}--\ref{sec:screening-berezin-response}
give a rigorous test of the proposed energy strategy. The map
$T(z)=z/\sqrt{1-|z|^2}$ sends hyperbolic area $dA/(1-|z|^2)^2$
to ordinary planar area, making the conjectured intensity constant.
Proposition~\ref{prop:exact-neutral-ocp} identifies the original gas,
up to a dilation tending to the identity, with an exactly neutral
hyperbolic plasma of background density $(4/\beta-1)/\pi$.
Global neutrality alone does not prove local neutrality.
Even at $\beta=2$, positive mean electrostatic energy of order $N$
persists near the moving boundary
(Propositions~\ref{prop:screening-finite-particle-energy-escape}
and~\ref{prop:neutral-mean-boundary-energy}).

Theorem~\ref{thm:screening-entropy-energy} identifies the precise
connected-force work and cutoff boundary estimate that would force
the conjecture. Both remain unproved for a general-temperature limit.
In an explicit reference family, however, the sign is established:
Theorem~\ref{thm:screening-tilt-monotonicity} gives strict nonlinear
response monotonicity for every bounded $L^2$ one-body tilt of the
Bergman process. Corollary~\ref{cor:screening-one-body-fixedpoint}
then proves uniqueness of its self-consistent screening potential,
and Proposition~\ref{prop:screening-weighted-projection-gap} gives
a uniform quantitative response bound. These are theorems about the
stated reference family, not uniqueness or the intensity formula for
the original gas at $\beta\ne2$.
\end{remark}

The quantitative proof for $\beta>4$ has two inputs.
Corollary~\ref{cor:analytic-insertion-without-clt} bounds compact insertion
expectations uniformly in $N$, using transport and tightness under
exponential tilts; it requires no central limit theorem. Harmonic boundary
transport proves an adjacent-partition-ratio lower bound that diverges
when $\beta>4$. Their combination forces every compact intensity to
zero. Corollaries~\ref{cor:empty-correlation-rates}
and~\ref{cor:empty-bernoulli-approximation} quantify all fixed correlations,
count moments, and the probability of multiple interior points.

At every positive $\beta$, the arguments also establish analytic Gaussian
fluctuations independent of the interior, a leading partition-function
asymptotic (Theorem~\ref{thm:partition-leading}), and an exact criterion
sufficient for uniqueness (Proposition~\ref{prop:transfer-uniqueness-criterion}).
The covariance estimate in that criterion remains a hypothesis; the
large-$\beta$ uniqueness proof does not use it.

The critical proof uses a different mechanism. The independent
Gaussian harmonic modes force a component of the electric-force
variance. A hypothetical nonempty limit at $4$ has a bounded
renormalized intensity, forcing its remaining score fluctuations
to vanish under boundary translation. The translated limiting
Ward identity then contradicts retained two-point repulsion.

The other extensions supply further information about the model.
Theorem~\ref{thm:partition-ratio-density-one} identifies adjacent-ratio
exponents on a set of full natural density.
Theorem~\ref{thm:boundary-weight-exponential} controls boundary-vanishing
statistics and permits explicit limiting changes of potential.
Theorem~\ref{thm:inserted-charge-screening} proves the exact Gaussian
mean shift produced by finitely many fixed charges, and
Corollary~\ref{cor:palm-screening-relative-invariance} gives relative
correlation invariance under fixed analytic tilts. The local Ward
identity, Proposition~\ref{prop:mobius-local-ward}, separates the
proved $O(N^{-1})$ response from the stronger estimate needed for
conformal invariance. These statements concern the original fixed
temperature; no interchange of limits in temperature and particle
number is made.



\paragraph{Conditional laws and hyperbolic geometry.}
For the DLR investigation, the main unconditional statements are in
Sections~\ref{sec:unconditional-canonical-quasigibbs}
and~\ref{sec:conditional-gaussian-curvature}. If $B\Subset\D$ is a
disk and $m\ge1$ is an allowed conditional count, the ordered
conditional density has the form
\[
 p_{\eta,m}(\mathbf z)=|\Delta_m(\mathbf z)|^\beta
                                      g_{\eta,m}(\mathbf z),
 \qquad
 \left[\partial_{z_i}\partial_{\overline z_j}\log g_{\eta,m}\right]_{ij}
 \succeq\frac\beta2\left[(1-z_i\overline z_j)^{-2}\right]_{ij},
\]
with $g_{\eta,m}>0$ real analytic. This is a canonical conditional-law
theorem; it does not identify the relative probabilities of different
counts. Section~\ref{sec:consistent-meromorphic-augmentation} constructs
a simultaneous meromorphic force field and proves nested-set
consistency of the canonical resampling kernels on that augmented space.
The stronger separation theorem in
Section~\ref{sec:gaussian-separation-meromorphic-field} permits exact
elimination of the independent Gaussian modes. The resulting canonical
Green kernels, their explicit geometric factor, and their affine
isometry covariance are proved in
Section~\ref{sec:gaussian-elimination-green-specification}.
Section~\ref{sec:bergman-complete-dlr} identifies those
probabilities and the full compensated DLR equation at $\beta=2$.
The results of Section~\ref{sec:hyperbolic-score-covariance} distinguish
covariance of the score and Gaussian field from invariance of the
point-process law. Section~\ref{sec:invariant-dlr-obstruction} proves
why an invariant Green specification needs compensation, then states
precisely the additional hypotheses needed for general-temperature
DLR and intensity conclusions.


\begin{theorem}[Canonical Gibbs nonuniqueness with infinitely many particles]
\label{thm:front-infinite-canonical-nonuniqueness}
For each $0<\beta<4$, infinitely many distinct infinite point-process
laws admit augmentations satisfying the same canonical Green DLR
specification as the actual disk-gas limits. They may all be chosen
rotation and reflection invariant. None of these constructed states
is fully hyperbolically invariant. For $0<\beta\le2$, every one is
mutually singular with every actual disk-gas limit: the constructed
states have finite Blaschke mass, whereas actual limits have infinite
Blaschke mass almost surely.
\end{theorem}
\begin{proof}
The construction and all the distinctions in the statement are proved
in Theorem~\ref{thm:infinite-green-canonical-nonuniqueness}.
Proposition~\ref{prop:explicit-szego-activity-states} gives an explicit
determinantal family at $\beta=2$.
\end{proof}

This is a nonuniqueness theorem for the canonical equations. It does
not assert that the original finite-$N$ disk gas has two subsequential
limits. Section~\ref{sec:boundary-insertion-nonuniqueness} gives an
exact oscillation criterion for that different question; no test
satisfying its nonzero-oscillation condition has been exhibited.

\section{Two elementary inequalities}
\begin{lemma}[Subharmonic powers and annular means]\label{lem:subharmonic}
If $P$ is holomorphic on a planar domain and $\beta>0$, then $|P|^\beta$
is nonnegative and subharmonic. If $u\ge0$ is subharmonic on a neighborhood
of $\overline{B(z,b)}$ and $0\le a<b$, then
\[
 u(z)\le\frac{1}{\pi(b^2-a^2)}
 \int_{a<|w-z|<b}u(w)\,dA(w).
\]
The circular mean of $u$ about $z$ is nondecreasing in its radius.
\end{lemma}
\begin{proof}
Write $\alpha=\beta/2$. For $\varepsilon>0$, direct differentiation gives
\[
 \Delta (|P|^2+\varepsilon)^\alpha
 =4\alpha(|P|^2+\varepsilon)^{\alpha-2}
       (\varepsilon+\alpha|P|^2)|P'|^2\ge0.
\]
Letting $\varepsilon\downarrow0$ locally uniformly proves the first claim.
For a smooth subharmonic function the derivative of its circular mean
at radius $s$ is $(2\pi s)^{-1}\int_{B(z,s)}\Delta u\,dA\ge0$.
Regularization gives the same conclusion for a general subharmonic function.
The submean inequality on circles, integrated against $2\pi s\,ds$ for
$a<s<b$, gives the annular statement. The same argument applies to
separately subharmonic functions one variable at a time.
\end{proof}

\begin{lemma}[Association and increasing tilts]\label{lem:association}
Let $\mu$ be a product probability measure on a finite product of intervals.
If $f,g$ are bounded coordinatewise nondecreasing real functions, then
$\Cov_\mu(f,g)\ge0$. Consequently, if $F\ge0$ is bounded, nondecreasing,
and $\mu(F)>0$, and $d\nu=F\,d\mu/\mu(F)$, then
\begin{equation}\label{eq:tilt}
 \nu(A)\le\mu(A)
\end{equation}
for every coordinatewise decreasing Borel set $A$. More generally,
$\nu(g)\le\mu(g)$ for every bounded nondecreasing function of a
coordinatewise decreasing vector-valued function of the coordinates.
\end{lemma}
\begin{proof}
For one coordinate, with independent identically distributed $X,X'$, twice
the covariance is
\[
 \E[(f(X)-f(X'))(g(X)-g(X'))]\ge0.
\]
Induct on the number of coordinates. Condition on the last coordinate and
apply the induction hypothesis to the conditional covariance. The two
conditional expectations are nondecreasing functions of the last
coordinate, so the covariance of these expectations is nonnegative by the
one-dimensional case. The covariance decomposition proves association.
Apply it to $F$ and $-\ind_A$ for \eqref{eq:tilt}, or to $F$ and the
negative of the indicated decreasing function for the final statement.
Truncation extends the expectation inequality to nonnegative unbounded
functions. This is the product association property in \cite{EPW}.
\end{proof}

\section{Radial comparison and tightness}\label{sec:core-proof}
Configurations with a zero radius or repeated radii have probability zero,
by absolute continuity. Sort radii and relabel the angles so that
$0<r_1<\cdots<r_N<1$. Symmetry makes the density on this ordered region
$N!$ times the original density. This constant will be absorbed in
normalization. Set
\begin{equation}\label{eq:ratios}
 t_k=\frac{r_k}{r_{k+1}},\quad r_{N+1}=1,
 \qquad r_j=\prod_{k=j}^N t_k,
 \qquad a_k=2k+\frac{\beta}{2}k(k-1).
\end{equation}
Thus $t\in(0,1)^N$. Factoring each pair distance and computing the
triangular Jacobian gives
\[
 \prod_{j=1}^N r_j\,dr_j=\prod_{k=1}^N t_k^{2k-1}\,dt_k,
 \qquad
 \prod_{j=1}^N r_j^{\beta(j-1)}
 =\prod_{k=1}^N t_k^{\beta k(k-1)/2}.
\]
For example, the Jacobian $|\det(\partial r/\partial t)|$ is
$\prod_{k=1}^N t_k^{k-1}$, whereas $\prod_jr_j=\prod_kt_k^k$.
After angular integration the ratio density is therefore proportional to
\begin{equation}\label{eq:ratio-density}
 F_N(t)\prod_{k=1}^N t_k^{a_k-1},
 \qquad
 F_N(t)=\int_{[0,2\pi)^N}
 \prod_{i<j}\left|1-
  \left(\prod_{\ell=i}^{j-1}t_\ell\right)e^{i(\theta_i-\theta_j)}
 \right|^\beta\prod_{j=1}^N\frac{d\theta_j}{2\pi}.
\end{equation}
This is a bounded continuous nonnegative function on $[0,1]^N$, and its
integral against the product beta density is positive.

\begin{lemma}[Monotonicity of the angular factor]\label{lem:angular}
The function $F_N$ is nondecreasing in each coordinate $t_k$. It is
independent of $t_N$.
\end{lemma}
\begin{proof}
Fix $k<N$ and all other ratios. Only factors with $i\le k<j$ depend on
$t_k$; write their radial ratios as $t_kq_{ij}$ with $q_{ij}$ independent
of $t_k$. Average additionally over the change
$\theta_i\mapsto\theta_i+\phi$ for all $i\le k$, keeping the outer
angles fixed. Haar invariance makes this extra average leave $F_N$
unchanged. The noncrossing factors are unchanged under this common
rotation, and the crossing product is
\[
 \left|P_\theta(t_ke^{i\phi})\right|^\beta,
 \qquad
 P_\theta(w)=\prod_{i\le k<j}
 (1-q_{ij}e^{i(\theta_i-\theta_j)}w).
\]
The circular mean in $\phi$ is nondecreasing in $t_k$ by
Lemma~\ref{lem:subharmonic}. All remaining factors are nonnegative and
independent of $t_k$, so their integration preserves the inequality.
No ratio in \eqref{eq:ratio-density} involves $t_N$.
\end{proof}

Let $U_k$ be independent with distribution $\Beta(a_k,1)$, namely density
$a_ku^{a_k-1}\ind_{(0,1)}(u)$. Under their product law, the ratio law is
the increasing tilt by $F_N$. Thus the actual ratios stochastically
dominate $(U_1,\ldots,U_N)$. Equivalently, all bounded coordinatewise
increasing functions of the ordered depths
\[
 X_{N,j}=-\log r_j=\sum_{k=j}^N(-\log t_k)
\]
have expectation at most the corresponding function of
\begin{equation}\label{eq:finite-depth-comparison}
 T_j^{(N)}=\sum_{k=j}^N E_k,
 \qquad E_k=-\log U_k\sim\Exp(a_k)\quad\hbox{independently}.
\end{equation}
The exponential convention is $\Pbb(E_k>x)=e^{-a_kx}$ for $x\ge0$.
Indeed, each depth is decreasing in every ratio, so this is exactly
Lemma~\ref{lem:association}. In particular, for $N\ge m$ and $s=\log(1/r)$,
\begin{equation}\label{eq:count-comparison}
 \Pbb\{M_N(r)\ge m\}
 \le \Pbb\{T_m^{(N)}>s\}
 \le \Pbb\{T_m\ge s\},
 \qquad T_m=\sum_{k=m}^{\infty}E_k.
\end{equation}
The infinite sums are almost surely finite, since $\sum_k a_k^{-1}<\infty$.

For $m\ge2$, take $\lambda=\beta m(m-1)/4$. Since
$a_k\ge\beta k(k-1)/2$, for $k\ge m$ we have $\lambda/a_k\le1/2$ and
\begin{align*}
 \log\E e^{\lambda T_m^{(N)}}
 &=\sum_{k=m}^N-\log(1-\lambda/a_k)
 \le2\lambda\sum_{k=m}^N a_k^{-1}\\
 &\le\frac{4\lambda}{\beta}\sum_{k=m}^\infty\frac1{k(k-1)}=m.
\end{align*}
Exponential Markov inequality proves \eqref{eq:gaussian-tail} at finite
$N$. The case $N<m$ is automatic. It also follows that for every
$r<1$, $q>0$, and $t\ge0$,
\begin{equation}\label{eq:uniform-moments}
 \sup_N\E M_N(r)^q<\infty,
 \qquad \sup_N\E e^{tM_N(r)}<\infty.
\end{equation}
For instance, for $m\ge2$ the summands in
$\sum_m e^{tm}\Pbb(M_N(r)=m)$ are bounded by
$\exp((t+1)m-c_rm(m-1))$, with $c_r=(\beta/4)\log(1/r)>0$.
The same calculation shows $\sup_N\E e^{tM_N(r)}\le C_r e^{C_r(t+1)^2}$,
where constants may also depend on $\beta$.

\begin{proof}[Proof of tightness in Theorem~\ref{thm:core}]
Choose an increasing sequence $r_\ell\uparrow1$. Given $\varepsilon>0$,
choose integers $b_\ell$ such that
$\sup_N\Pbb(M_N(r_\ell)>b_\ell)<\varepsilon2^{-\ell}$.
The set
\[
 \mathcal K=\bigcap_{\ell\ge1}
       \{\xi\in\Mp(\D):\xi(B_{r_\ell})\le b_\ell\}
\]
is closed, since counting an open relatively compact set is lower
semicontinuous under vague convergence. Its elements have uniformly
bounded mass on every compact set. Such a set is relatively compact in
the vague topology, and its vague limits are again counting measures.
These standard facts can also be seen by diagonal extraction on compact
exhaustions, where boundedly many atom positions have convergent
subsequences; atoms escaping an exhaustion are treated at the next level.
Thus $\mathcal K$ is compact. A union bound gives
$\inf_N\Pbb(\Xi_N\in\mathcal K)\ge1-\varepsilon$.
The space $\Mp(\D)$ is Polish, so Prokhorov's theorem proves tightness
and subsequential existence.
\end{proof}

\section{Correlation densities and local regularity}
\label{hier:section}
\label{sec:hierarchy}

We continue to fix $\beta>0$.  Let $\Xi_{N_j}$ be any subsequence
converging in law, for the vague topology on $\D$, to a locally finite
integer-valued random measure $\Xi$.  At this stage multiplicities are
allowed in the definition of $\Xi$.  The compact count estimates proved
above imply, for every $R<1$ and every $t>0$,
\begin{equation}
 \sup_N \E e^{tM_N(R)}<\infty,
 \qquad \E e^{tM(R)}<\infty.
 \label{hier:exponential-moments}
\end{equation}
For the limiting bound, fix $R<S<1$ and choose $\chi\in C_c(\D)$
with $0\leq\chi\leq1$, $\chi=1$ on $\overline{B_R}$, and support
in $B_S$. The continuous mapping theorem and the lower-semicontinuous
form of Portmanteau give
\[
 \E e^{tM(R)}\leq\E e^{t\Xi(\chi)}
 \leq\liminf_j\E e^{t\Xi_{N_j}(\chi)}
 \leq\sup_N\E e^{tM_N(S)}<\infty.
\]
We use unnormalized planar area measure throughout this section, and write
\[
 (m)_p=m(m-1)\cdots(m-p+1),\qquad
 \Delta_p(z)=\prod_{1\le i<j\le p}(z_i-z_j).
\]
Empty products have value one.

\begin{theorem}[Regular correlation hierarchy]
\label{thm:hierarchy}
The limit $\Xi$ is simple.  For each integer $p\ge1$, it has a continuous
factorial correlation density of the form
\begin{equation}
 \rho_p(z)=|\Delta_p(z)|^\beta h_p(z),
 \label{hier:limiting-factorization}
\end{equation}
where $h_p$ is nonnegative, plurisubharmonic, and locally H\"older
continuous with exponent $\min\{1,\beta\}$.  Along the given subsequence,
\begin{equation}
 h_{p,N_j}\longrightarrow h_p,
 \qquad \rho_{p,N_j}\longrightarrow\rho_p
 \quad\hbox{locally uniformly on }\D^p.
 \label{hier:uniform-convergence}
\end{equation}
For every $r<1$ there is a finite constant $C_{p,r,\beta}$, independent
of $N$, such that
\begin{equation}
 \rho_{p,N}(z)\le C_{p,r,\beta}|\Delta_p(z)|^\beta,
 \qquad
 \rho_p(z)\le C_{p,r,\beta}|\Delta_p(z)|^\beta,
 \qquad z\in\overline{\D_r}^{\,p}.
 \label{hier:vandermonde-bound}
\end{equation}
Here $\rho_{p,N}=h_{p,N}=0$ when $p>N$.
\end{theorem}

We give all details, including the compactness argument needed to
identify the limiting densities.
Here plurisubharmonic means upper semicontinuous with subharmonic
restriction to every complex affine line on each component of its
intersection with the domain.

For $N\ge p$, direct integration of the joint density gives
\begin{align}
 \rho_{p,N}(z)&=|\Delta_p(z)|^\beta h_{p,N}(z),
 \label{hier:finite-factorization}\\
 h_{p,N}(z)
 &=\frac{(N)_p}{Z_{N,\beta}}
   \int_{\D^{N-p}}|\Delta_{N-p}(y)|^\beta
   \prod_{i=1}^p\prod_{\ell=1}^{N-p}|z_i-y_\ell|^\beta
   \,\dd A^{N-p}(y).
 \label{hier:h-definition}
\end{align}
In particular $h_{p,N}$ is continuous.  The integrand, as a function of
$z$, is the $\beta$th power of the modulus of a holomorphic polynomial,
and is therefore plurisubharmonic.  Integration preserves this property.
Only separate subharmonicity is needed for the following estimate.

\begin{lemma}[Separated annular averages]
\label{hier:annular-lemma}
Fix $p\ge1$, $\beta>0$, and $0<r<R<1$.  There is a constant
$C=C(p,\beta,r,R)<\infty$ such that every nonnegative separately
subharmonic continuous function $u$ on $\D_R^p$ satisfies
\begin{equation}
 \sup_{\overline{\D_r}^{\,p}}u
 \le C\int_{\D_R^p}|\Delta_p(w)|^\beta u(w)\,\dd A^p(w).
 \label{hier:annular-estimate}
\end{equation}
\end{lemma}

\begin{proof}
Fix a tuple $z^*\in\overline{\D_r}^{\,p}$.  Partition its coordinates
into groups with common values $c_1,\ldots,c_q$.  Choose $\eta>0$ so
small that
\[
 100p\eta\le R-r,
 \qquad
 100p\eta\le\min_{a\ne b}|c_a-c_b|,
\]
omitting the second condition when $q=1$.  In each group assign its
coordinates distinct labels $\ell=1,\ldots,m_a$.

Suppose that $\max_i|z_i-z_i^*|<\eta/4$.  For the coordinate with
label $\ell_i$, average over the annulus
\[
 A_i(z_i)=\{w:(3\ell_i-2)\eta\le|w-z_i|
                         \le(3\ell_i-1)\eta\}.
\]
All these annuli lie in $\D_R$.  For two coordinates in the same group,
the radial intervals have a gap of at least $2\eta$, whereas their
centers are at distance less than $\eta/2$.  Their annuli are thus at
distance at least $3\eta/2$.  For different groups, the lower bound on
the distance between the $c_a$'s gives a still larger separation.
Consequently
\[
 |w_i-w_j|\ge\eta\quad(i\ne j),
 \qquad |A_i(z_i)|\ge3\pi\eta^2.
\]
The circular submean inequality, integrated in radius, is also an
annular submean inequality.  Applying it successively in all variables
and then using the displayed separation gives
\[
 u(z)\le
 (3\pi)^{-p}\eta^{-2p-\beta p(p-1)/2}
 \int_{\D_R^p}|\Delta_p(w)|^\beta u(w)\,\dd A^p(w).
\]
The neighborhoods of the tuples $z^*$ constructed in this way cover
$\overline{\D_r}^{\,p}$.  A finite subcover gives a single finite
constant in \eqref{hier:annular-estimate}.
\end{proof}

Since
\[
 \int_{\D_R^p}\rho_{p,N}(w)\,\dd A^p(w)
 =\E(M_N(R))_p,
\]
Lemma~\ref{hier:annular-lemma} and
\eqref{hier:exponential-moments} show that $h_{p,N}$ are uniformly
bounded on compact subsets of $\D^p$.  This already proves the
finite-$N$ part of \eqref{hier:vandermonde-bound}.

\begin{lemma}[Uniform H\"older estimates]
\label{hier:holder-lemma}
For each $p\ge1$ and $r<1$, there is a constant $C<\infty$ such that
\[
 |h_{p,N}(z)-h_{p,N}(w)|
 \le C\|z-w\|^{\min\{1,\beta\}},
 \qquad z,w\in\overline{\D_r}^{\,p},
\]
uniformly in $N$.
\end{lemma}

\begin{proof}
Write \eqref{hier:h-definition} as
\[
 h_{p,N}(z)=\int |F_y(z)|^\beta\,\dd\mu_N(y),
 \qquad F_y(z)=\prod_{i,\ell}(z_i-y_\ell),
\]
where $\mu_N$ is the finite positive measure containing the prefactor
and the factor $|\Delta_{N-p}(y)|^\beta$.
Fix $r<s<t<1$.  For every holomorphic function $F$ on $\D_t^p$,
iterated scalar submean inequalities, followed by Cauchy estimates on
an intermediate polydisk, give
\begin{equation}
 \sup_{\overline{\D_s}^{\,p}}|F|^\beta
 +\sum_{i=1}^p\sup_{\overline{\D_s}^{\,p}}|\partial_iF|^\beta
 \le C\int_{\D_t^p}|F(z)|^\beta\,\dd A^p(z).
 \label{hier:analytic-derivative-estimate}
\end{equation}
These estimates hold for every $\beta>0$: one first applies the
submean inequality to $|F|^\beta$, and only then uses the usual
supremum form of Cauchy's derivative estimate.
Integrating \eqref{hier:analytic-derivative-estimate} with respect to
$\mu_N$ and using the uniform local bound on $h_{p,N}$ yields
\[
 \int(A_y^\beta+D_y^\beta)\,\dd\mu_N(y)\le C,
\]
where
\[
 A_y=\sup_{\overline{\D_s}^{\,p}}|F_y|,
 \qquad
 D_y=\sum_i\sup_{\overline{\D_s}^{\,p}}|\partial_iF_y|.
\]
The polydisk is convex, so
$|F_y(z)-F_y(w)|\le C_pD_y\|z-w\|$ for $z,w$ in the smaller
polydisk.  If $0<\beta\le1$, the inequality
\[
 \big||a|^\beta-|b|^\beta\big|\le|a-b|^\beta
\]
proves the required estimate after integration.  If $\beta>1$, use
\[
 \big||a|^\beta-|b|^\beta\big|
 \le\beta|a-b|(|a|+|b|)^{\beta-1}
\]
and H\"older's inequality to obtain
\[
 |h_{p,N}(z)-h_{p,N}(w)|
 \le C\|z-w\|
 \left(\int D_y^\beta\,\dd\mu_N\right)^{1/\beta}
 \left(\int A_y^\beta\,\dd\mu_N\right)^{(\beta-1)/\beta}.
\]
The right side is uniformly bounded by $C\|z-w\|$.
\end{proof}

For completeness, factorial moment measures also behave continuously
in the present setting.  If $\xi$ is an integer-valued locally finite
measure, its $p$th factorial measure $\xi^{[p]}$ counts ordered
selections of distinct particle labels, even when some labels occupy
the same location.  For $f\in C_c(\D^p)$,
\begin{equation}
 \langle f,\xi^{[p]}\rangle
 =\sum_{\pi\in\mathcal P_p}
 (-1)^{p-|\pi|}\prod_{B\in\pi}(|B|-1)!
 \int f(x_{\pi(1)},\ldots,x_{\pi(p)})
        \,\dd\xi^{\otimes|\pi|}(x).
 \label{hier:partition-formula}
\end{equation}
Here $\mathcal P_p$ is the set of partitions of $\{1,\ldots,p\}$;
the blocks are ordered by their smallest elements, and $\pi(i)$ is
the index of the block containing $i$. To verify the formula, expand
each product measure as a sum over labels and group terms by their
label-equality partition. The coefficient for a block of $n$ equal
labels is the sum of the signs of all permutations of $n$ elements:
a partition into cycles contributes
$(-1)^{n-|\pi|}\prod_{B\in\pi}(|B|-1)!$. This sum is $1$ for $n=1$
and $0$ for $n\geq2$, since multiplication by a transposition pairs
permutations with opposite signs. Hence only distinct labels survive.

Every term on the right is continuous for vague convergence. Indeed,
continuous functions on the relevant compact product can be uniformly
approximated by sums of products of one-variable continuous functions,
using compactly supported cutoffs. Vaguely convergent measures have
uniformly bounded mass on those compact sets, so this approximation
reduces the assertion to the defining one-variable convergence. The
diagonal substitution in each term has compact support because every
block variable occurs among the arguments of $f$. Thus
$\xi\mapsto\langle f,\xi^{[p]}\rangle$ is continuous.  Its absolute
value is at most $\|f\|_\infty\xi(K)^p$ for a compact set $K$
containing all coordinate projections of $\operatorname{supp}f$.
The count estimates give uniform integrability.  Consequently,
\begin{equation}
 \E\langle f,\Xi_{N_j}^{[p]}\rangle
 \longrightarrow\E\langle f,\Xi^{[p]}\rangle.
 \label{hier:factorial-moment-convergence}
\end{equation}

\begin{proof}[Proof of Theorem~\ref{thm:hierarchy}]
The local bounds and Lemma~\ref{hier:holder-lemma} give local uniform
precompactness of $h_{p,N_j}$ by Arzel\`a--Ascoli.  Let $h$ be the
limit along any further subsequence.  Equations
\eqref{hier:finite-factorization} and
\eqref{hier:factorial-moment-convergence} show that the $p$th
factorial moment measure of $\Xi$ has density $|\Delta_p|^\beta h$.
This identifies $h$ uniquely away from the diagonals.  Since that set
is dense and $h$ is continuous, it identifies $h$ everywhere.
Every locally uniformly convergent further subsequence therefore has
the same limit, proving \eqref{hier:uniform-convergence} along the
original process-convergent subsequence.  Local uniform limits
preserve plurisubharmonicity and the H\"older bounds.

In particular the second factorial moment measure of $\Xi$ is
absolutely continuous with respect to planar area on $\D^2$.
Its mass on the diagonal is zero.  On every compact set $K$, this
mass equals
\[
 \E\sum_{x\in K}\Xi(\{x\})\bigl(\Xi(\{x\})-1\bigr).
\]
It follows, by a countable compact exhaustion, that $\Xi$ has no
multiplicities almost surely.  The remaining assertions follow from
the estimates already proved.
\end{proof}

\begin{corollary}[Null sets and small clusters]
\label{cor:hier-clusters}
The process $\Xi$ almost surely avoids every fixed area-null Borel
subset of $\D$.  Moreover, for each compact $K\subset\D$ and each
$p\ge1$, there are constants $C,\varepsilon_0>0$ such that
\begin{equation}
 \Pbb\{\Xi(B(z,\varepsilon))\ge p\}
 \le C\varepsilon^{2p+\beta p(p-1)/2},
 \qquad z\in K,\quad 0<\varepsilon\le\varepsilon_0.
 \label{hier:cluster-estimate}
\end{equation}
The same estimate holds uniformly for $\Xi_N$.
\end{corollary}

\begin{proof}
The first assertion follows by integrating $\rho_1$ over the null
set and using a compact exhaustion.  For the second, choose a compact
disk containing all the balls in question.  Markov's inequality for
the factorial count, followed by \eqref{hier:vandermonde-bound}, gives
\[
 \Pbb\{\Xi(B(z,\varepsilon))\ge p\}
 \le\frac1{p!}\int_{B(z,\varepsilon)^p}\rho_p(w)\,\dd A^p(w)
 \le\frac{C}{p!}(2\varepsilon)^{\beta p(p-1)/2}
                    (\pi\varepsilon^2)^p.
\]
The finite-$N$ proof is identical.
\end{proof}

\begin{corollary}[Convergence of disk counts]\label{hier:disk-counts}
For each $r<1$, along the chosen subsequence,
$M_{N_j}(r)\Rightarrow M(r)$ and all their positive moments converge.
Their exponential moments converge for every real parameter. The
Gaussian tail bound \eqref{eq:gaussian-tail} holds for $M(r)$ with
the same constants as at finite $N$.
\end{corollary}
\begin{proof}
The fixed circle $\partial B_r$ is an area-null set and therefore
contains no point of $\Xi$ almost surely. Counting $B_r$ is consequently
vaguely continuous at $\Xi$ almost surely. The continuous mapping theorem
gives convergence in distribution of the integer-valued counts. The
uniform exponential bounds give uniform integrability of every fixed
power and exponential, proving convergence of the expectations.
Since the counts are integer valued, probabilities of $M\geq m$
converge by using the continuity threshold $m-1/2$. This passes
\eqref{eq:gaussian-tail} to $\Xi$ and completes all assertions of
Theorem~\ref{thm:core}.
\end{proof}

\begin{corollary}[Leading small-cluster asymptotics]
\label{cor:hier-cluster-asymptotics}
Fix $a\in\D$ and $p\ge1$, and put
\[
 d_p=2p+\frac\beta2p(p-1),\qquad
 Z_{p,\beta}=\int_{\D^p}|\Delta_p(w)|^\beta\,\dd A^p(w).
\]
As $\varepsilon\downarrow0$,
\begin{align}
 \E\bigl(\Xi(B(a,\varepsilon))\bigr)_p
 &=\varepsilon^{d_p}
      \bigl[Z_{p,\beta}h_p(a,\ldots,a)+o(1)\bigr],
 \label{hier:cluster-factorial-asymptotics}\\
 \Pbb\{\Xi(B(a,\varepsilon))\ge p\}
 &=\varepsilon^{d_p}
      \left[\frac{Z_{p,\beta}}{p!}h_p(a,\ldots,a)+o(1)\right],
 \label{hier:cluster-probability-asymptotics}
\end{align}
and the same leading asymptotic holds for the probability of exactly
$p$ points.  The coefficient is permitted to vanish.
\end{corollary}

\begin{proof}
In the integral of $\rho_p$ over $B(a,\varepsilon)^p$, make the
change of variables $z_i=a+\varepsilon w_i$.  The area factor is
$\varepsilon^{2p}$ and the Vandermonde factor contributes
$\varepsilon^{\beta p(p-1)/2}$.  Continuity of $h_p$ gives
\eqref{hier:cluster-factorial-asymptotics}.  If
$L=\Xi(B(a,\varepsilon))$, then
\[
 0\le\E\binom Lp-\Pbb(L\ge p)
 \le(p+1)\E\binom L{p+1}
 =O(\varepsilon^{d_{p+1}}).
\]
Since $d_{p+1}>d_p$, this proves
\eqref{hier:cluster-probability-asymptotics}.  The probability of
at least $p+1$ points has this same smaller order, proving the
assertion for exactly $p$ points.
\end{proof}

\begin{remark}[Even inverse temperatures]
\label{hier:even-beta}
If $\beta=2q$ with $q\in\N$, the functions $h_p$ and $\rho_p$ are
real analytic.  Indeed, polarize the representation in the proof of
Lemma~\ref{hier:holder-lemma}:
\[
 H_N(z,\overline w)
 =\int F_y(z)^q\overline{F_y(w)^q}\,\dd\mu_N(y).
\]
This is holomorphic in $z$ and antiholomorphic in $w$, and
\[
 |H_N(z,\overline w)|
 \le h_{p,N}(z)^{1/2}h_{p,N}(w)^{1/2}.
\]
The uniform local bounds and Montel's theorem give a polarized
holomorphic limit along a further subsequence.  Its diagonal is $h_p$,
so $h_p$ is real analytic.  The factor $|\Delta_p|^{2q}$ is a
polynomial in the real and imaginary coordinates, proving the same
assertion for $\rho_p$.
\end{remark}

\section{Janossy densities and convergence in local total variation}
\label{hier:janossy-section}

Let $B\subset\D$ be any relatively compact Borel set.  No regularity
or null-boundary assumption is imposed on $B$.  Write $M_B=\Xi(B)$.
Its $m$th Janossy measure is
\[
 J_{B,m}(A)
 =\E\bigl[\ind_{\{M_B=m\}}\Xi_B^{[m]}(A)\bigr],
 \qquad A\subset B^m.
\]
For $m=0$, this means the scalar $\Pbb(M_B=0)$.  The convention used
here gives total mass $m!\Pbb(M_B=m)$ for $m\ge1$.

\begin{proposition}[Janossy expansion]
\label{hier:janossy-proposition}
For every $m\ge1$, $J_{B,m}$ has a density $j_{B,m}$ satisfying
$0\le j_{B,m}\le\rho_m$ almost everywhere on $B^m$ and
\begin{equation}
 j_{B,m}(z)
 =\sum_{\ell=0}^{\infty}\frac{(-1)^\ell}{\ell!}
       \int_{B^\ell}\rho_{m+\ell}(z,y)\,\dd A^\ell(y).
 \label{hier:janossy-expansion}
\end{equation}
The series is absolutely convergent in $L^1(B^m)$.  The same formula
holds for $m=0$, with $\rho_0=1$, and gives the hole probability.
It also holds for every finite-$N$ process.
\end{proposition}

\begin{proof}
The domination $J_{B,m}\le\E\Xi_B^{[m]}$ gives absolute continuity
and the density bound.  For a finite configuration with $n$ points
in $B$, the inclusion--exclusion identity for each ordered selection
of $m$ labels is
\[
 \sum_{\ell=0}^{n-m}\frac{(-1)^\ell}{\ell!}(n-m)_\ell
 =(1-1)^{n-m}=\ind_{\{n=m\}}.
\]
For $n<m$ there are no such selections.  Integrate against a bounded
test function on $B^m$ and take expectations.  The sum of the
$L^1$ norms of the terms in \eqref{hier:janossy-expansion} is
\begin{align*}
 \sum_{\ell\ge0}\frac1{\ell!}\E(M_B)_{m+\ell}
 &=\E\bigl[(M_B)_m2^{M_B-m}\ind_{\{M_B\ge m\}}\bigr]
 <\infty.
\end{align*}
The exponential count moments therefore justify all interchanges
and prove the claimed $L^1$ identity.  The proof for $m=0$ is the
same binomial identity without the initial selection of labels.
\end{proof}

\begin{corollary}[Absolute continuity in a finite window]
For every relatively compact Borel set $B$, the law of $\Xi|_B$ is
absolutely continuous with respect to a Poisson process of unit area
intensity on $B$. On a configuration of $m$ distinct points
$z_1,\ldots,z_m$, a version of the density is
$e^{|B|}j_{B,m}(z_1,\ldots,z_m)$, where $|B|$ is its area.
If $\Pbb(M_B=m)>0$, a uniformly randomly labeled configuration
conditional on $M_B=m$ has density
$j_{B,m}/(m!\Pbb(M_B=m))$ on $B^m$.
\end{corollary}
\begin{proof}
The Poisson Janossy density on $B^m$ is the constant $e^{-|B|}$,
and the probability law is its symmetric pushforward with weight
$1/m!$. Comparing it with $j_{B,m}$ proves the first formula.
The second follows by dividing the Janossy measure by its mass
$m!\Pbb(M_B=m)$. If $|B|=0$, both processes are empty on $B$ almost
surely, and the first assertion holds with the same convention.
\end{proof}

\begin{theorem}[Local total variation convergence]
\label{hier:local-tv}
Along every subsequence $\Xi_{N_j}\Rightarrow\Xi$ considered above,
and for every relatively compact Borel set $B\subset\D$,
\begin{equation}
 \bigl\|\mathcal L(\Xi_{N_j}|_B)-\mathcal L(\Xi|_B)
       \bigr\|_{\mathrm{TV}}\longrightarrow0.
 \label{hier:tv-convergence}
\end{equation}
Here total variation is the supremum of the differences of
probabilities over measurable events in the space of finite
configurations on $B$.
\end{theorem}

\begin{proof}
For fixed $m$ and $\ell$, local uniform convergence of the
correlations implies
\[
 \int_{B^\ell}\rho_{m+\ell,N_j}(\,\cdot\,,y)\,\dd A^\ell(y)
 \longrightarrow
 \int_{B^\ell}\rho_{m+\ell}(\,\cdot\,,y)\,\dd A^\ell(y)
 \quad\hbox{in }L^1(B^m).
\]
This conclusion uses only the fact that $\overline B$ is compact
in $\D$.  For the inclusion--exclusion terms with $\ell\ge L$,
the sum of their $L^1$ norms is at most
\[
 \E\bigl[(M_B)_m2^{M_B-m}
                  \ind_{\{M_B\ge m+L\}}\bigr].
\]
The analogous expression for $\Xi_N$ tends to zero uniformly in
$N$ as $L\to\infty$, by \eqref{hier:exponential-moments} applied
to a disk containing $\overline B$.  It follows that, for every
$m\ge0$,
\begin{equation}
 j_{B,m,N_j}\longrightarrow j_{B,m}
 \quad\hbox{in }L^1(B^m),
 \label{hier:janossy-l1-convergence}
\end{equation}
with the natural scalar interpretation when $m=0$.

The law of a finite simple configuration is the pushforward, under
$z\mapsto\sum_i\delta_{z_i}$, of its symmetric Janossy measures
with weights $1/m!$.  Hence its total variation distance is bounded
by
\[
 \frac12\sum_{m=0}^{\infty}\frac1{m!}
             \|j_{B,m,N_j}-j_{B,m}\|_{L^1(B^m)}.
\]
For $m\le L$ this tends to zero by
\eqref{hier:janossy-l1-convergence}.  The remaining terms are at
most one half the sum of the two probabilities of more than $L$
points in $B$, which tend to zero uniformly as $L\to\infty$.
This proves \eqref{hier:tv-convergence}.
\end{proof}

\section{Entire transforms and determination by correlations}
\label{hier:transforms-section}

\begin{proposition}[Generating functional]
\label{hier:generating-proposition}
For every bounded complex-valued Borel function $g$ with relatively
compact support in $\D$,
\begin{equation}
 \E\prod_{x\in\Xi}(1+g(x))
 =\sum_{p=0}^{\infty}\frac1{p!}
       \int_{\D^p}\prod_{i=1}^p g(z_i)\rho_p(z)
                        \,\dd A^p(z).
 \label{hier:generating-functional}
\end{equation}
The series is absolutely convergent.  In particular, the full
correlation hierarchy determines the law of $\Xi$.
\end{proposition}

\begin{proof}
Expand the finite product over the points in a relatively compact
Borel set $B$ containing the support of $g$.  The absolute sum of
the expected terms is at most
\[
 \sum_{p\ge0}\frac{\|g\|_\infty^p}{p!}\E(M_B)_p
 =\E(1+\|g\|_\infty)^{M_B}<\infty.
\]
This proves the identity and absolute convergence.  Taking
$g=e^{-f}-1$ for nonnegative continuous compactly supported $f$
recovers the Laplace functional, which determines the law of a
locally finite random measure.
\end{proof}

\begin{corollary}[Entire moment and probability generating functions]
\label{cor:hier-entire}
For every relatively compact Borel set $B$, the function
\[
 G_B(s)=\E s^{M_B},\qquad s\in\C,
\]
is entire, and
\begin{equation}
 G_B(s)=\sum_{p=0}^{\infty}\frac{(s-1)^p}{p!}
                        \int_{B^p}\rho_p(z)\,\dd A^p(z).
 \label{hier:pgf}
\end{equation}
For bounded compactly supported real-valued Borel functions
$f_1,\ldots,f_q$, the function
\[
 (z_1,\ldots,z_q)\longmapsto
 \E\exp\left(\sum_{a=1}^qz_a\langle f_a,\Xi\rangle\right)
\]
is entire on $\C^q$.  Its moments determine the joint distribution
of the linear statistics.  Along a process-convergent subsequence, the corresponding
finite-$N$ transforms converge locally uniformly in the complex
parameters, even for these Borel test functions.
\end{corollary}

\begin{proof}
On each compact set of complex parameters, the integrands and all
their parameter derivatives are dominated by a constant times
$e^{tM_B}$ for some finite $t$, where $B$ contains all supports.
The exponential moment bounds justify differentiation and the
power-series expansions of every order.  Formula \eqref{hier:pgf}
is also the special case $g=(s-1)\ind_B$ of
\eqref{hier:generating-functional}.

For convergence, first restrict to configurations with at most
$L$ points in $B$.  The integrands are then bounded, so local total
variation convergence applies.  The contribution of the omitted
configurations tends to zero uniformly in $N$ and uniformly on
compact parameter sets, by a slightly larger exponential moment.
This proves the final assertion.  Entire transforms are uniquely
determined by their Taylor coefficients.  Their values on imaginary
parameters give the joint characteristic function, proving moment
determination.
\end{proof}

\begin{remark}[Growth from the Gaussian count tail]
\label{hier:transform-growth}
The stronger compact count estimate obtained from the radial
comparison has the form
$\Pbb(M_B\ge m)\le C e^{-cm^2}$, after adjustment of constants.
Summing the tail and completing the square gives
\[
 \E e^{tM_B}\le C\exp(Ct^2+Ct),\qquad t\ge0.
\]
Thus the entire moment generating function of each bounded compactly
supported linear statistic has growth order at most two.  Likewise,
\[
 \log\sup_{|s|\le R}|G_B(s)|
 =O\bigl((\log(1+R))^2+1\bigr),\qquad R\to\infty.
\]
All these estimates may be chosen uniformly for the finite-$N$
processes.
\end{remark}

\section{A simultaneous comparison for the ranked radii}
\label{rad:section-comparison}

For a locally finite counting measure $\xi$ on $\D$, let $R_k(\xi)$ be
the $k$th smallest modulus of its points, counting multiplicities.  Set
$R_k(\xi)=1$ if $\xi(\D)<k$, and put
\[
 X_k(\xi)=-\log R_k(\xi),
 \qquad -\log 0=+\infty.
\]
Thus the ranked depths form a nonincreasing sequence in $[0,+\infty]$;
zeros are appended when there are only finitely many points.  Recall that
$E_k$, $k\geq1$, are independent exponentials with rates
\[
 a_k=2k+\frac{\beta}{2}k(k-1),
 \qquad T_k=\sum_{j=k}^{\infty}E_j.
\]
Since $\sum_j a_j^{-1}<\infty$, all $T_k$ are finite almost surely and
$T_k\downarrow0$ almost surely.

\begin{proposition}[Coupling of all ranked depths]
\label{rad:prop-coupling}
For each fixed $N$, there is a coupling of $\Xi_N$ with $(T_k)_{k\geq1}$
such that
\begin{equation}
 X_k(\Xi_N)\leq T_k\qquad\text{for every }k\geq1
 \quad\text{almost surely}.
 \label{rad:eq-ranked-coupling}
\end{equation}
The same assertion holds, in a separate coupling, for every subsequential
limit $\Xi$ furnished by Theorem~\ref{thm:core}.  In particular, if
\[
 L(t)=\#\{k\geq1:T_k>t\},\qquad t>0,
\]
then in these couplings
\begin{equation}
 M_N(r)\leq L(\log(1/r)),\qquad
 M(r)\leq L(\log(1/r)),\qquad 0<r<1,
 \label{rad:eq-count-coupling}
\end{equation}
simultaneously in $r$.
\end{proposition}

\begin{proof}
The radial comparison used in Theorem~\ref{thm:core} shows that the gap
vector is stochastically dominated by $(E_1,\ldots,E_N)$. Summing coordinates gives
coordinatewise stochastic domination of the ranked depths by
\[
 \left(\sum_{j=k}^N E_j\right)_{1\leq k\leq N},
\]
with zeros appended.  This latter vector is coordinatewise at most
$(T_k)_{k\geq1}$. Apply Strassen's theorem first to the finite depth
vector and $(T_k^{(N)})_{k\leq N}$. Differences of the latter recover
the independent $E_1,\ldots,E_N$. Append independent $E_j$ for $j>N$;
this gives \eqref{rad:eq-ranked-coupling} for each fixed $N$.

For completeness, ranked radii are continuous under vague convergence
of locally finite counting measures on $\D$, with values in $[0,1]$.
Indeed, to locate a ranked radius one brackets it by radii $u<v<1$
whose circles contain no points of the limiting configuration; the
counts in $B_u$ and $B_v$ then converge.  Such brackets can approach the
ranked radius from both sides.  At the endpoints $0$ and $1$, only the
nontrivial side of the bracket is needed.  Consequently the map to
$(X_k)_{k\geq1}$ is continuous into the product space
$[0,+\infty]^{\N}$, where the interval has its compact topology.

Along any convergent subsequence of $\Xi_N$, take the couplings above
and extract a weakly convergent subsequence of their joint laws of
ranked depths and $(T_k)$.  These laws are tight, and their second
marginal is fixed.  The set
\[
 \{(x,t):x_k\leq t_k\text{ for all }k\geq1\}
\]
is closed.  The limiting joint law therefore couples the ranked depths
of $\Xi$ with $(T_k)$ as claimed.  If desired, a regular conditional
distribution of $\Xi$ given its ranked depths lifts this to a coupling
of the full configuration with $(T_k)$; the same observation applies
to the finite-$N$ couplings.  Finally,
$M(r)=\#\{k:X_k(\Xi)>\log(1/r)\}$, and the same
identity holds for $M_N(r)$, proving \eqref{rad:eq-count-coupling}.
\end{proof}

\section{Explicit occupancy and hole bounds}
\label{rad:section-explicit}

\begin{proposition}[A gamma-function occupancy bound]
\label{rad:prop-gamma-tail}
Let $c=4/\beta$.  For every $k\geq1$, $0<r<1$, and every $N$,
\begin{equation}
 \Pbb\{M_N(r)\geq k\}
 \leq
 \frac{\Gamma(2k+c)}{\Gamma(k+1)\Gamma(k+c)}\,
 r^{\,2k+\beta k(k-1)/2}.
 \label{rad:eq-gamma-tail}
\end{equation}
The same bound holds with $M_N$ replaced by $M$ for every subsequential
limit.  In particular,
\begin{equation}
 \Pbb\{M(r)>0\}\leq (1+4/\beta)r^2,
 \qquad
 \Pbb\{M_N(r)>0\}\leq (1+4/\beta)r^2.
 \label{rad:eq-small-disk}
\end{equation}
\end{proposition}

\begin{proof}
Write $t=\log(1/r)$.  By Proposition~\ref{rad:prop-coupling}, it
suffices to bound $\Pbb\{T_k>t\}$.  Conditioning on $T_{k+1}$ and using
its independence of $E_k$ gives
\[
 \Pbb\{E_k+T_{k+1}>t\mid T_{k+1}\}
 \leq e^{-a_kt}e^{a_kT_{k+1}}.
\]
We compute the exponential moment on the right.  Since
\[
 a_j=\frac{\beta}{2}j(j+c-1),\qquad
 a_j-a_k=\frac{\beta}{2}(j-k)(j+k+c-1),
\]
for $J>k$ independence gives
\begin{align*}
 \E\exp\left(a_k\sum_{j=k+1}^J E_j\right)
 &=\prod_{j=k+1}^J\frac{a_j}{a_j-a_k}\\
 &=\frac{\Gamma(J+1)}{\Gamma(k+1)\Gamma(J-k+1)}
   \frac{\Gamma(J+c)}{\Gamma(k+c)}
   \frac{\Gamma(2k+c)}{\Gamma(J+k+c)}.
\end{align*}
The usual ratio asymptotic $\Gamma(x+b)/\Gamma(x+d)\sim x^{b-d}$
shows that the last expression tends to
$\Gamma(2k+c)/(\Gamma(k+1)\Gamma(k+c))$.  Monotone convergence proves
that this limit is exactly $\E e^{a_kT_{k+1}}$.  Thus
\[
 \Pbb\{T_k>t\}\leq
 e^{-a_kt}\frac{\Gamma(2k+c)}{\Gamma(k+1)\Gamma(k+c)},
\]
which proves the assertions.
\end{proof}

\begin{proposition}[Uniform positive hole probabilities]
\label{rad:prop-holes}
Fix $0<r<1$ and put
\[
 t=\log(1/r),\qquad
 J=\max\left\{1,\left\lceil\frac{8}{\beta t}\right\rceil\right\},
 \qquad
 h_\beta(r)=\frac12\prod_{j=1}^J
 \left(1-e^{-a_jt/(2J)}\right)>0.
\]
Then
\begin{equation}
 \Pbb\{M_N(r)=0\}\geq h_\beta(r),\qquad
 \Pbb\{M(r)=0\}\geq h_\beta(r).
 \label{rad:eq-hole-bound}
\end{equation}
Consequently every compact subset of $\D$ has a positive hole
probability, uniformly in $N$ and in the choice of subsequential limit.
\end{proposition}

\begin{proof}
For $j\geq2$, $a_j\geq\beta j(j-1)/2$, whence
\[
 \E T_{J+1}\leq\frac{2}{\beta J}\leq\frac t4.
\]
Markov's inequality gives
$\Pbb\{T_{J+1}\leq t/2\}\geq1/2$.  Independently, the event
$E_j<t/(2J)$ for every $1\leq j\leq J$ has probability equal to the
product in the definition of $h_\beta(r)$.  Their intersection implies
$T_1<t$.  By \eqref{rad:eq-count-coupling}, this event forces both
counts, in their respective couplings, to vanish.  Any compact subset
of $\D$ is contained in some $B_r$.
\end{proof}

\section{The pure-death comparison and boundary growth}
\label{rad:section-boundary}

\begin{proposition}[A mean occupancy bound]
\label{rad:prop-mean}
For every $0<r<1$,
\begin{equation}
 \E M_N(r)\leq\frac{1}{1-r^{\beta/2}}
 \leq1+\frac{2}{\beta\log(1/r)},\qquad
 \E M(r)\leq\frac{1}{1-r^{\beta/2}}.
 \label{rad:eq-mean}
\end{equation}
\end{proposition}

\begin{proof}
For $n\geq1$, define
\[
 L_n(t)=\#\left\{1\leq k\leq n:
               \sum_{j=k}^n E_j>t\right\}.
\]
This is a pure-death chain starting from $n$: its holding time at state
$k$ is $E_k$, and the transition from $k$ to $k-1$ has rate $a_k$.
Set $a_0=0$ and $m_n(t)=\E L_n(t)$. The elementary generator identity and Jensen's
inequality give
\[
 m_n'(t)=-\E a_{L_n(t)}
 \leq-\frac{\beta}{2}\E[L_n(t)(L_n(t)-1)]
 \leq-\frac{\beta}{2}m_n(t)(m_n(t)-1).
\]
Comparing with the scalar differential equation with initial value
$n$ yields
\[
 m_n(t)\leq
 \frac{1}{1-(1-1/n)e^{-\beta t/2}}.
\]
In the construction using the common sequence $(E_j)$,
$L_n(t)\uparrow L(t)$ for each $t>0$.  Monotone convergence therefore
gives
\[
 \E L(t)\leq\frac{1}{1-e^{-\beta t/2}}.
\]
Apply Proposition~\ref{rad:prop-coupling}.  The second inequality in
\eqref{rad:eq-mean} uses
$1/(1-e^{-x})=1+1/(e^x-1)\leq1+1/x$ for $x>0$.
\end{proof}

\begin{lemma}[The reference process near the boundary]
\label{rad:lem-reference-growth}
Almost surely,
\begin{equation}
 kT_k\longrightarrow\frac{2}{\beta}\quad(k\to\infty),
 \qquad
 tL(t)\longrightarrow\frac{2}{\beta}\quad(t\downarrow0).
 \label{rad:eq-reference-growth}
\end{equation}
\end{lemma}

\begin{proof}
Since $a_j\sim\beta j^2/2$,
\[
 b_k:=\E T_k=\sum_{j=k}^{\infty}a_j^{-1}
       \sim\frac{2}{\beta k},\qquad
 \operatorname{Var}(T_k)=\sum_{j=k}^{\infty}a_j^{-2}=O(k^{-3}).
\]
For every $\varepsilon>0$, Chebyshev's inequality gives
\[
 \sum_{\ell=1}^{\infty}
 \Pbb\{\ell^2|T_{\ell^2}-b_{\ell^2}|>\varepsilon\}<\infty.
\]
The Borel--Cantelli lemma, first for rational $\varepsilon>0$, implies
$\ell^2T_{\ell^2}\to2/\beta$ almost surely.  If
$\ell^2\leq k\leq(\ell+1)^2$, monotonicity of $T_k$ gives
\[
 \ell^2 T_{(\ell+1)^2}\leq kT_k
 \leq(\ell+1)^2T_{\ell^2}.
\]
Both outer quantities tend to $2/\beta$, proving the first assertion.

All $E_j$ are strictly positive almost surely, so $T_k$ is strictly
decreasing and $L(t)\to\infty$ as $t\downarrow0$.  If $k=L(t)$, then
$T_{k+1}\leq t<T_k$.  Hence
\[
 kT_{k+1}\leq kt<kT_k.
\]
The first assertion makes both outer quantities tend to $2/\beta$.
This proves the second assertion.
\end{proof}

\begin{corollary}[Almost sure upper boundary growth]
\label{rad:cor-boundary-growth}
Every subsequential limit satisfies
\begin{equation}
 \limsup_{r\uparrow1}(1-r)M(r)\leq\frac{2}{\beta}
 \qquad\text{almost surely}.
 \label{rad:eq-boundary-growth}
\end{equation}
This is an upper bound; no positive lower boundary growth or infinitude
of the limiting configuration is asserted.
\end{corollary}

\begin{proof}
In the coupling of Proposition~\ref{rad:prop-coupling},
$M(e^{-t})\leq L(t)$ simultaneously for all $t>0$.  Apply
Lemma~\ref{rad:lem-reference-growth} and
$(1-e^{-t})/t\to1$.
\end{proof}

\section{Weighted summability of the limiting configuration}
\label{rad:section-weighted}

\begin{proposition}[Exponential moments of depth norms]
\label{rad:prop-depth-moments}
For $p>1$, define
\[
 S_p(\xi)=\sum_{z\in\xi}(-\log|z|)^p,
\]
where points are counted with multiplicity.  For every $0\leq\lambda<2$,
\begin{equation}
 \sup_{N\geq1}\E\exp\{\lambda S_p(\Xi_N)^{1/p}\}
 \leq
 \prod_{j=1}^{\infty}
 \left(1-\frac{\lambda j^{1/p}}{a_j}\right)^{-1}
 <\infty.
 \label{rad:eq-depth-mgf}
\end{equation}
The same bound holds for $\E\exp\{\lambda S_p(\Xi)^{1/p}\}$ for any
subsequential limit.  In particular, these weighted sums have all
positive moments, and
\begin{equation}
 \sum_{z\in\Xi}(1-|z|)^p<\infty
 \qquad\text{almost surely for every }p>1.
 \label{rad:eq-boundary-summability}
\end{equation}
\end{proposition}

\begin{proof}
By the ranked coupling and Minkowski's inequality in $\ell^p$,
\begin{align*}
 S_p(\Xi)^{1/p}
 &\leq\left(\sum_{k=1}^{\infty}T_k^p\right)^{1/p}\\
 &=\left\|\sum_{j=1}^{\infty}
           E_j(\ind_{\{k\leq j\}})_{k\geq1}\right\|_{\ell^p}
 \leq\sum_{j=1}^{\infty}j^{1/p}E_j.
\end{align*}
The same inequalities hold for $\Xi_N$.  One can first apply
Minkowski's inequality to finite sums and then use monotone convergence.
Since $p>1$,
\[
 \sum_{j=1}^{\infty}\frac{j^{1/p}}{a_j}<\infty,
\]
so the sum on the right is finite almost surely.  Independence and
monotone convergence give
\[
 \E\exp\left\{\lambda\sum_{j=1}^{\infty}j^{1/p}E_j\right\}
 =\prod_{j=1}^{\infty}
 \left(1-\frac{\lambda j^{1/p}}{a_j}\right)^{-1}.
\]
Every factor is finite: $a_j\geq2j$ implies
$\lambda j^{1/p}/a_j\leq\lambda/2<1$.  Moreover,
\[
 \sum_{j=1}^{\infty}
 -\log\left(1-\frac{\lambda j^{1/p}}{a_j}\right)
 \leq\frac{\lambda}{1-\lambda/2}
       \sum_{j=1}^{\infty}\frac{j^{1/p}}{a_j}<\infty.
\]
This proves \eqref{rad:eq-depth-mgf} and all asserted moment bounds.

The inequality $1-r\leq-\log r$ proves
\eqref{rad:eq-boundary-summability} for each fixed $p>1$.  To obtain the
statement simultaneously for all such $p$, first intersect the
probability-one events for rational $p>1$, and then, for any given
$p>1$, choose a rational $q$ with $1<q<p$ and use
$(1-r)^p\leq(1-r)^q$.
\end{proof}

\begin{remark}
\label{rad:rem-scope}
The radial estimates provide upper occupancy bounds, positive local
hole probabilities, and summability restrictions.  They supply neither
a lower bound on the probability of seeing a point in a compact set
nor uniqueness of the subsequential law.  In particular, they do not by
themselves exclude the identically empty limit.
\end{remark}

\section{Symmetry and the radial intensity}
\label{sec:intensity-extensions}

Throughout this section, $\Xi$ is a subsequential limit and
$\rho_p$ denotes its continuous $p$-point correlation density provided by
Theorem~\ref{thm:hierarchy}. Along the subsequence defining $\Xi$, write
$\rho_{p,N}\to\rho_p$ for the locally uniform convergence established
there. All assertions concern this arbitrary subsequential limit.

\begin{proposition}[Rotation and reflection symmetry]
\label{prop:orthogonal-symmetry}
The law of $\Xi$ is invariant under every rotation about the origin and
under complex conjugation. Consequently, for every $p\geq1$,
$\theta\in\R$, and $z_1,\ldots,z_p\in\D$,
\[
 \rho_p(e^{i\theta}z_1,\ldots,e^{i\theta}z_p)
 =\rho_p(z_1,\ldots,z_p),
 \qquad
 \rho_p(\overline z_1,\ldots,\overline z_p)
 =\rho_p(z_1,\ldots,z_p).
\]
In particular, the one-point intensity is radial.
\end{proposition}

\begin{proof}
Rotations and conjugation preserve planar area, the disk, and every
pairwise distance, and therefore preserve each finite-$N$ Gibbs law.
If $T$ is one of these maps, the induced map on point measures is
continuous for the vague topology: for $f\in C_c(\D)$,
$(T_*\xi)(f)=\xi(f\circ T)$ and $f\circ T\in C_c(\D)$.
The continuous mapping theorem passes invariance to every subsequential
limit. The corresponding correlation measures are invariant as well.
Their continuous densities satisfy the displayed identities first
almost everywhere, and then everywhere by continuity.
\end{proof}

\begin{lemma}[Log-subharmonicity of the finite-$N$ intensity]
\label{lem:log-subharmonic-intensity}
For every $N\geq1$, the function $\rho_{1,N}$ is strictly positive on
$\D$, and $\log\rho_{1,N}$ is subharmonic. The radial function
\[
 g_N(t):=\log\rho_{1,N}(e^t),\qquad -\infty<t<0,
\]
is convex and nondecreasing. In particular, $r\mapsto\rho_{1,N}(r)$ is
nondecreasing on $[0,1)$.
\end{lemma}

\begin{proof}
Integrating out $N-1$ particles gives a representation
\[
 \rho_{1,N}(z)=\int |P_y(z)|^\beta\,\dd\nu_N(y),
 \qquad
 P_y(z)=\prod_{j=1}^{N-1}(z-y_j),
\]
where $\nu_N$ is the finite positive measure containing the normalization
constant and the Vandermonde interaction among the $y_j$.
For $N=1$, the polynomial and the integral over the empty tuple have
their usual empty-product interpretation. The integrand is strictly
positive for almost every $y$, for each fixed $z$, so
$\rho_{1,N}(z)>0$.

For $\varepsilon>0$, set
\[
 \phi_{y,\varepsilon}(z)
 =\frac\beta2\log\bigl(|P_y(z)|^2+\varepsilon\bigr),
 \qquad
 H_\varepsilon(z)
 =\log\int e^{\phi_{y,\varepsilon}(z)}\,\dd\nu_N(y).
\]
The functions $\phi_{y,\varepsilon}$ are smooth and subharmonic; indeed,
\[
 \Delta\log\bigl(|P_y|^2+\varepsilon\bigr)
 =\frac{4\varepsilon|P_y'|^2}
 {\bigl(|P_y|^2+\varepsilon\bigr)^2}\geq0,
\]
where $\Delta$ is the ordinary real Laplacian. Polynomial bounds on
compact sets justify differentiating the integral. Let
$\mathbb E_{z,\varepsilon}$ denote expectation for the probability
measure proportional to
$e^{\phi_{y,\varepsilon}(z)}\dd\nu_N(y)$. Direct differentiation gives
\[
 \Delta H_\varepsilon
 =\mathbb E_{z,\varepsilon}[\Delta\phi_{y,\varepsilon}]
  +\mathbb E_{z,\varepsilon}
       [|\nabla\phi_{y,\varepsilon}|^2]
  -\bigl|\mathbb E_{z,\varepsilon}
       [\nabla\phi_{y,\varepsilon}]\bigr|^2
 \geq0.
\]
As $\varepsilon\downarrow0$, these subharmonic functions decrease to
the finite function $\log\rho_{1,N}$. The decreasing-limit property
for subharmonic functions proves the first assertion. Equivalently,
one can pass directly to their submean inequalities, using the positive
continuous function $\rho_{1,N}$ to give a local lower bound.

Rotation symmetry makes $\log\rho_{1,N}$ radial. On an annulus
$e^a<|z|<e^b$, the harmonic function with its two constant boundary
values is the affine interpolation of $g_N(a)$ and $g_N(b)$ in the
variable $\log|z|$. The maximum principle therefore gives the convexity
inequality for $g_N$. If $g_N$ had a negative secant slope, convexity
would imply $g_N(t)\to+\infty$ as $t\to-\infty$. This contradicts
boundedness of $\log\rho_{1,N}$ near the origin. Thus $g_N$ is
nondecreasing. Continuity at the origin extends monotonicity to
$r=0$.
\end{proof}

\begin{proposition}[Intensity dichotomy and logarithmic convexity]
\label{prop:intensity-dichotomy}
Exactly one of the following alternatives holds:
\begin{enumerate}
 \item $\Xi$ is almost surely the empty configuration and
       $\rho_1\equiv0$;
 \item $\Xi$ is not almost surely empty and
       $\rho_1(r)>0$ for every $0<r<1$.
\end{enumerate}
In the second alternative,
\[
 t\longmapsto\log\rho_1(e^t),\qquad -\infty<t<0,
\]
is convex and nondecreasing. The density $\rho_1$ is continuous and
radially nondecreasing on $\D$, and $\log\rho_1$, with
$\log0=-\infty$, is subharmonic. Moreover, every fixed nonempty open
set $U\subset\D$ satisfies
\[
 \Pbb\{\Xi(U)>0\}>0.
\]
The proposition does not assert $\rho_1(0)>0$ or
$\Pbb\{\Xi(\D)>0\}=1$.
\end{proposition}

\begin{proof}
The identity $\E\Xi(A)=\int_A\rho_1\,\dd A$, first for relatively
compact sets and then by monotone convergence, shows that
$\rho_1\equiv0$ is equivalent to $\Xi$ being almost surely empty.
Local uniform convergence and Lemma~\ref{lem:log-subharmonic-intensity}
show that $\rho_1$ is continuous, radial, and nondecreasing.

Suppose $\rho_1$ is not identically zero. Its intensity measure is
nonzero, so there is $r_0\in(0,1)$ with $\rho_1(r_0)>0$.
Monotonicity already gives positivity for $r\geq r_0$.
For $0<s<r_0$, choose $r_0<R<1$ and set
\[
 \alpha=\frac{\log R-\log r_0}{\log R-\log s}\in(0,1).
\]
The finite-$N$ convexity inequality is
\[
 \rho_{1,N}(r_0)
 \leq \rho_{1,N}(s)^\alpha
       \rho_{1,N}(R)^{1-\alpha}.
\]
Passing to the locally uniform limit would contradict
$\rho_1(r_0)>0$ if $\rho_1(s)=0$. Thus $\rho_1(s)>0$ for every
$s>0$. Logs now converge uniformly on each compact radial interval in
$(0,1)$, and their limit is convex and nondecreasing in $\log r$.

A radial function that is convex in $\log r$ is subharmonic on the
punctured disk. This follows, for example, by smoothing the convex
function on compact logarithmic intervals and using
$\Delta[g(\log r)]=g''(\log r)/r^2$ for the smooth approximants.
The resulting function $\log\rho_1$ is locally bounded above near
the origin. Its subharmonic extension there takes the value
\[
 \limsup_{z\to0}\log\rho_1(z)=\log\rho_1(0),
\]
by continuity and radial monotonicity of $\rho_1$; this value may be
$-\infty$. The removable-singularity property for subharmonic
functions gives the asserted subharmonicity on $\D$.

Finally, any nonempty open $U$ contains a disk with compact closure
in $\D\setminus\{0\}$. The integral of the positive continuous
intensity over that disk is strictly positive. Its point count is a
nonnegative integrable random variable with positive expectation and
therefore has positive probability of being positive. This proves the
open-set assertion.
\end{proof}

\begin{corollary}[An explicit local intensity bound]
For $0\leq r<R<1$, both the finite-$N$ and limiting radial intensities satisfy
\[
 \rho_{1,N}(r),\ \rho_1(r)
 \leq\frac{1}{\pi(R^2-r^2)(1-R^{\beta/2})}.
\]
In particular, for $z\in\D$,
\[
 \rho_{1,N}(z),\ \rho_1(z)
 \leq\frac{8}{\pi\min\{\beta/2,1\}(1-|z|)^2}.
\]
\end{corollary}
\begin{proof}
Radial monotonicity and Proposition~\ref{rad:prop-mean} give
\[
 \pi(R^2-r^2)\rho_1(r)
 \leq\int_{B_R\setminus B_r}\rho_1(z)\,dA(z)
 \leq\E M(R)\leq\frac{1}{1-R^{\beta/2}},
\]
and identically at finite $N$. Take $R=(1+r)/2$ and use
$R^2-r^2\geq(1-r)/4$ and
$1-R^{\beta/2}\geq\min\{\beta/2,1\}(1-R)$.
\end{proof}

\begin{remark}
The distinction between ``not almost surely empty'' and ``almost surely
nonempty'' is essential. The positive-intensity alternative above does
not exclude an atom at the empty configuration in the law of $\Xi$.
That atom, and indeed every finite total count, will be excluded for
$0<\beta<2$ by Theorem~\ref{thm:almost-sure-infinite-below-two}.
Nor do the Euclidean symmetries proved here establish invariance under
general conformal automorphisms of the disk.
\end{remark}


\section{The determinantal case as a check}\label{sec:beta-two}
This section verifies that the normalization and topology recover the
Bergman process when $\beta=2$. It is not used in the preceding estimates.
The functions
\[
 p_j(z)=\sqrt{\frac{j+1}{\pi}}z^j,\qquad j\ge0,
\]
are orthonormal in $L^2(\D,dA)$, by direct polar integration. The
Vandermonde determinant and expansion of determinants give
\[
 \frac{|\Delta_N(z)|^2}{Z_{N,2}}
 =\frac1{N!}\det[K_N(z_i,z_j)]_{i,j=1}^N,
 \qquad K_N(z,w)=\sum_{j=0}^{N-1}p_j(z)\overline{p_j(w)},
 \qquad Z_{N,2}=\pi^N.
\]
Indeed, integrating the square of the determinant of the orthonormal
functions leaves precisely the $N!$ matching permutations; the product
of their squared leading coefficients is $N!/\pi^N$.
Orthogonality also gives $\int K_N(z,u)K_N(u,w)dA(u)=K_N(z,w)$ and
$\int K_N(u,u)dA(u)=N$. Expanding a determinant and using these two
identities gives
\[
 \int_{\D}\det[K_N(z_i,z_j)]_{i,j=1}^k\,dA(z_k)
 =(N-k+1)\det[K_N(z_i,z_j)]_{i,j=1}^{k-1}.
\]
To check the coefficient, permutations fixing $k$ give $N$ times the
smaller determinant, whereas inserting $k$ into each of the $k-1$
positions in the cycles of a smaller permutation gives the negative
contribution $-(k-1)$ times that determinant. Iterating this identity
proves
\[
 \rho_{p,N}(z_1,\ldots,z_p)=\det[K_N(z_i,z_j)]_{i,j=1}^p.
\]
The geometric series and its derivative imply local uniform convergence
\[
 K_N(z,w)\longrightarrow
 K(z,w)=\frac{1}{\pi(1-z\overline w)^2}.
\]
Every subsequential limit therefore has correlations $\det[K(z_i,z_j)]$.
The determination-by-correlations result proved above makes their laws
equal. Tightness then gives convergence of the entire sequence to this
law, the Bergman determinantal process; see also \cite{PV,BGZNW}.

For this process,
\begin{align*}
 \E M(r)&=\int_{B_r}\frac{dA(z)}{\pi(1-|z|^2)^2}
 =\frac{r^2}{1-r^2},\\
 \Var M(r)&=\E M(r)-\int_{B_r^2}|K(z,w)|^2dA(z)dA(w)
 \le\E M(r).
\end{align*}
The variance identity follows by integrating $\rho_2=\rho_1\rho_1-|K|^2$
and using $M^2=(M)_2+M$. Chebyshev's inequality now shows that the total
number of points is infinite almost surely: for every fixed integer $L$,
$\Pbb(M(r)\le L)\to0$ as $r\uparrow1$, and
$\{\Xi(\D)\le L\}\subset\{M(r)\le L\}$.
In fact
\[
 \lim_{r\uparrow1}(1-r)M(r)=\frac12\qquad\hbox{almost surely}.
\]
For a proof, use $r_n=1-n^{-2}$, $n\ge2$. The means of
$(1-r_n)M(r_n)$ tend to $1/2$, while their variances are $O(n^{-2})$.
Chebyshev and Borel--Cantelli give almost sure convergence along this
sequence. Monotonicity of $M(r)$ and
$(1-r_n)/(1-r_{n+1})\to1$ extend it to all $r\uparrow1$.

% only standard amsmath/amssymb notation. All labels have the prefix raduniq.
\section{A positive radial expansion and its limitations for uniqueness}

We give an exact refinement of the radial comparison. It applies to every
real $\beta>0$, without an integrality assumption. Although the resulting
unnormalized angular factors are monotone in $N$, we then show that two
natural ways to turn this observation into a convergence argument fail
already at $\beta=2$.

Order the radii as $0<r_1<\cdots<r_N<1$, which is possible almost surely,
and put
\[
 t_k=r_k/r_{k+1},\qquad r_{N+1}=1,\qquad
 a_k=2k+\frac\beta2k(k-1),\qquad p=\frac\beta2.
\]
Use normalized Haar measure on each angular circle, and define
\begin{equation}\label{eq:raduniq-angular}
 F_N(t)=\int_{[0,2\pi)^N}
 \prod_{i<j}\left|1-
 \left(\prod_{\ell=i}^{j-1}t_\ell\right)
 e^{\mathrm i(\theta_i-\theta_j)}\right|^{2p}
 \prod_{i=1}^N\frac{d\theta_i}{2\pi}.
\end{equation}
In particular, $F_N$ does not depend on $t_N$. The ordered-ratio density,
with respect to $dt_1\cdots dt_N$, is
\begin{equation}\label{eq:raduniq-ratio-density}
 \frac1{W_N}F_N(t)\prod_{k=1}^N t_k^{a_k-1},\qquad
 W_N=\frac{Z_{N,\beta}}{N!(2\pi)^N}.
\end{equation}
For completeness, the change of variables gives
$\prod_j r_j\,dr_j=\prod_k t_k^{2k-1}\,dt_k$.
Factoring $r_j^\beta$ from each pair $i<j$ contributes
$\prod_k t_k^{\beta k(k-1)/2}$. Angular integration and the $N!$
possible radial orders give \eqref{eq:raduniq-ratio-density}.

\begin{proposition}[Positive Fourier expansion]\label{prop:raduniq-positive}
For $n\in\mathbb N_0$, let
\[
 b_n=(-1)^n\binom pn,
 \qquad
 \binom p0=1,\qquad
 \binom pn=\frac{p(p-1)\cdots(p-n+1)}{n!}\quad(n\ge1).
\]
For $m=(m_1,\ldots,m_{N-1})\in\mathbb N_0^{N-1}$, define
\begin{equation}\label{eq:raduniq-coefficients}
 B_{N,m}=
 \sum_{\substack{n_{ij}\in\mathbb N_0\ (i<j)\\
       \sum_{i\le k<j}n_{ij}=m_k\ (1\le k<N)}}
       \prod_{i<j}b_{n_{ij}}.
\end{equation}
Each sum in \eqref{eq:raduniq-coefficients} is finite. For every
$t\in(0,1)^N$,
\begin{equation}\label{eq:raduniq-positive-series}
 F_N(t)=\sum_{m\in\mathbb N_0^{N-1}}
            B_{N,m}^{\,2}\prod_{k=1}^{N-1}t_k^{2m_k}.
\end{equation}
The series has nonnegative terms and is finite at every such $t$.
\end{proposition}

\begin{proof}
The branch of $(1-u)^p$ that equals $1$ at $u=0$ is holomorphic for
$|u|<1$ and has the absolutely convergent series
$(1-u)^p=\sum_{n\ge0}b_nu^n$. Since $p$ is real, its modulus is
$|1-u|^p$. Apply this expansion to the analytic product
\[
 A_N(t,\theta)=\prod_{i<j}
 \left(1-\left(\prod_{\ell=i}^{j-1}t_\ell\right)
 e^{\mathrm i(\theta_i-\theta_j)}\right)^p.
\]
For fixed $t\in(0,1)^N$, all the arguments have modulus strictly below
$1$. The product of the finitely many absolutely convergent series
therefore converges absolutely and uniformly in $\theta$.

A choice of edge multiplicities $(n_{ij})_{i<j}$ contributes the radial
monomial $\prod_{k<N}t_k^{m_k}$, where
$m_k=\sum_{i\le k<j}n_{ij}$, and the angular character
$e^{\mathrm i\sum_i\alpha_i\theta_i}$, where
\[
 \alpha_i=\sum_{j>i}n_{ij}-\sum_{j<i}n_{ji}.
\]
These quantities satisfy
\[
 m_k=\sum_{i=1}^k\alpha_i.
\]
Conversely, with the conventions $m_0=m_N=0$, the angular charge is
uniquely determined by $m$ through $\alpha_i=m_i-m_{i-1}$.
Thus different crossing vectors $m$ yield different Fourier characters.
For a fixed $m$, every contributing edge multiplicity obeys
\[
 0\le n_{ij}\le\min_{i\le k<j}m_k.
\]
There are only finitely many such edge arrays. This proves the asserted
finiteness of \eqref{eq:raduniq-coefficients} and justifies regrouping
the absolutely convergent series as
\[
 A_N(t,\theta)=
 \sum_m B_{N,m}\left(\prod_{k<N}t_k^{m_k}\right)
 \exp\left(\mathrm i\sum_{i=1}^N(m_i-m_{i-1})\theta_i\right).
\]
Orthogonality of distinct characters in normalized Haar measure, or
Parseval's identity, now gives \eqref{eq:raduniq-positive-series}, since
$|A_N(t,\theta)|^2$ is the integrand in
\eqref{eq:raduniq-angular}. The coefficients are real, so their squared
moduli are $B_{N,m}^{\,2}$.
\end{proof}

\begin{corollary}[An exact mixture of independent gaps]
\label{cor:raduniq-mixture}
In every expression indexed by $m\in\mathbb N_0^{N-1}$, set $m_N=0$.
The numbers
\begin{equation}\label{eq:raduniq-mixing-probabilities}
 \pi_N(m)=\frac{B_{N,m}^{\,2}}{W_N}
             \prod_{k=1}^N\frac1{a_k+2m_k}
\end{equation}
form a probability mass function. The ordered radial gaps
$X_k=-\log t_k$ have the following exact construction: first sample
$m$ with law $\pi_N$, and then, conditionally on $m$, sample independently
\[
 X_k\sim\operatorname{Exp}(a_k+2m_k),\qquad 1\le k\le N.
\]
Here $\operatorname{Exp}(\lambda)$ denotes the exponential distribution
of rate $\lambda$. In particular, $X_N\sim\operatorname{Exp}(a_N)$ is
independent of $(X_1,\ldots,X_{N-1})$.
\end{corollary}

\begin{proof}
Tonelli's theorem applied to the nonnegative series
\eqref{eq:raduniq-positive-series} gives
\[
 W_N=\int_{(0,1)^N}F_N(t)\prod_k t_k^{a_k-1}\,dt
     =\sum_m B_{N,m}^{\,2}\prod_{k=1}^N\frac1{a_k+2m_k}.
\]
The normalization is finite by \eqref{eq:raduniq-ratio-density} and
positive because $B_{N,0}=1$. Substituting this identity into the ratio
density expresses that density as the mixture, with weights
\eqref{eq:raduniq-mixing-probabilities}, of the product densities
\[
 \prod_{k=1}^N(a_k+2m_k)t_k^{a_k+2m_k-1}.
\]
Taking negative logarithms gives the claimed exponential distributions.
Since the conditional distribution of $X_N$ is the same for every $m$,
its conditional independence also gives its unconditional independence
from all the other gaps.
\end{proof}

\begin{proposition}[Monotone unnormalized factors do not give a martingale]
\label{prop:raduniq-not-martingale}
The angular factors satisfy
\begin{equation}\label{eq:raduniq-monotone-N}
 F_{N+1}(t_1,\ldots,t_N)
 \ge F_N(t_1,\ldots,t_{N-1}),\qquad t_k\in(0,1).
\end{equation}
Nevertheless, on the common product reference space
\[
 \mu_\infty=\bigotimes_{k\ge1}\operatorname{Beta}(a_k,1),
 \qquad
 H_N=\frac{F_N(U_1,\ldots,U_{N-1})}
               {\mathbb E_{\mu_\infty}F_N},
\]
the normalized factors need not be a martingale, a submartingale, or a
supermartingale for the filtration
$\mathcal F_N=\sigma(U_1,\ldots,U_{N-1})$.
\end{proposition}

\begin{proof}
If the last crossing coordinate for $N+1$ vertices is $m_N=0$, then
$\sum_{i\le N}n_{i,N+1}=0$. Nonnegativity forces all edges into vertex
$N+1$ to have multiplicity zero. Hence
\[
 B_{N+1,(m,0)}=B_{N,m}.
\]
The terms with last coordinate zero reproduce the whole series for
$F_N$; all remaining terms in $F_{N+1}$ are nonnegative. This proves
\eqref{eq:raduniq-monotone-N}.

To test the normalized factors, take $\beta=2$, so that $a_k=k(k+1)$.
Put $x=U_1^2$ and $y=U_2^2$. Direct expansion of the Vandermonde, or
\eqref{eq:raduniq-positive-series}, gives
\[
 F_2=1+x,\qquad
 F_3=1+x+y+xy^2+x^2y+x^2y^2.
\]
The reference laws give
$\mathbb Ey=6/(6+2)=3/4$ and
$\mathbb Ey^2=6/(6+4)=3/5$. Therefore
\[
 \mathbb E[F_3\mid\mathcal F_2]
 =\frac74+\frac85x+\frac{27}{20}x^2.
\]
Since $\mathbb Ex=1/2$ and $\mathbb Ex^2=1/3$, we have
$\mathbb EF_2=3/2$ and $\mathbb EF_3=3$. It follows that
\begin{equation}\label{eq:raduniq-sign-changing}
 \mathbb E[H_3\mid\mathcal F_2]-H_2
 =\frac13\left(-\frac14-\frac25x+\frac{27}{20}x^2\right).
\end{equation}
The polynomial in parentheses is negative at $x=0$ and positive at
$x=1$. Since $x$ is uniform on $(0,1)$, both signs occur with positive
probability. This disproves each of the three asserted types of
conditional ordering for this example.
\end{proof}

\begin{remark}[Positive coefficients do not imply log-supermodularity]
There is a second obstruction to inferring a monotone coupling from the
positive expansion. With $u=t_1^2$ and $v=t_2^2$, the same $\beta=2$
example gives
\[
 F_3(u,v)=1+u+v+uv^2+u^2v+u^2v^2.
\]
At $(u,v)=(0,0)$, its value and first derivatives equal $1$, whereas
$\partial_u\partial_v F_3(0,0)=0$. Consequently,
\[
 \partial_u\partial_v\log F_3(0,0)=-1.
\]
This mixed derivative remains negative for all sufficiently small
positive $u,v$, by continuity. Thus $F_3$ is not log-supermodular even
in the interior of its domain. Equivalently, in the original positive
coordinates the mixed derivative is multiplied by $4t_1t_2>0$ and
still has negative sign. Multiplication by the product reference density
does not change a mixed derivative of its logarithm. Therefore positive
Fourier coefficients alone do not supply the log-supermodularity
hypothesis used in the corresponding FKG lattice criterion. Failure of
that sufficient criterion is not a claim that the resulting law fails
every possible association inequality.
\end{remark}

\begin{proposition}[Failure of uniform integrability at $\beta=2$]
\label{prop:raduniq-no-ui}
At $\beta=2$, the normalized factors $(H_N)$ in
Proposition~\ref{prop:raduniq-not-martingale} are not uniformly
integrable on $(\mu_\infty,\mathcal F_\infty)$.
The reference ordered depths and the Bergman ordered depths have
mutually singular laws: their almost-sure asymptotics are respectively
\begin{equation}\label{eq:raduniq-incompatible-depths}
 k\sum_{j\ge k}(-\log U_j)\longrightarrow1,
 \qquad kD_k\longrightarrow\frac12.
\end{equation}
The only determinantal input in the proof is the $\beta=2$ radial law,
which we recall and derive below; no general-$\beta$ convergence or
uniqueness statement is used.
\end{proposition}

\begin{proof}
Under $\mu_\infty$, the variables $X_j=-\log U_j$ are independent
exponentials of rates $j(j+1)$. Their tail sums $S_k=\sum_{j\ge k}X_j$
are finite almost surely, and
\[
 \mathbb ES_k=\sum_{j\ge k}\frac1{j(j+1)}=\frac1k,
 \qquad
 \operatorname{Var}S_k
 =\sum_{j\ge k}\frac1{j^2(j+1)^2}\le\frac C{k^3}.
\]
For any $\varepsilon>0$, Chebyshev's inequality gives
$\mathbb P(|kS_k-1|>\varepsilon)\le C/(\varepsilon^2k)$.
Apply Borel--Cantelli at $k=n^2$, first for every positive rational
$\varepsilon$, and then interpolate between successive squares using
the monotonicity of $S_k$. This proves the first almost-sure limit in
\eqref{eq:raduniq-incompatible-depths}.

At $\beta=2$, the unordered radii of the finite-$N$ gas have the law
of independent variables $R_j$ with densities
$2j r^{2j-1}\,dr$, $1\le j\le N$, after forgetting their labels.
Indeed, expand
$\Delta_N(z)=\det[z_i^{j-1}]_{i,j=1}^N$. In its squared modulus,
angular integration kills all terms with distinct permutations. The
remaining sum of radial monomials, multiplied by $\prod_i r_i\,dr_i$,
is exactly the symmetrization of these independent radial densities
after normalization. Thus the corresponding unordered depths have the
law of independent variables $Y_j=-\log R_j\sim\operatorname{Exp}(2j)$.
We may use one infinite independent sequence $(Y_j)$ to couple these
radial laws for every $N$.

For $x>0$, let
\[
 M(x)=\sum_{j\ge1}\mathbf1_{\{Y_j\ge x\}}.
\]
This count is almost surely finite, since
\[
 \mathbb EM(x)=\frac1{e^{2x}-1},\qquad
 \operatorname{Var}M(x)\le\mathbb EM(x).
\]
Chebyshev and Borel--Cantelli at $x=n^{-2}$ show that
$xM(x)\to1/2$ almost surely along this sequence. Monotonicity of $M$
and the fact that successive values of $n^{-2}$ have ratio tending to
$1$ extend the limit to all $x\downarrow0$.
The sequence $(Y_j)$ has no ties almost surely and contains infinitely
many positive numbers, with only finitely many above each positive
threshold. It therefore has a decreasing ordering
$D_1>D_2>\cdots>0$, with $D_k\downarrow0$ and $M(D_k)=k$.
The full limit $xM(x)\to1/2$, evaluated at $x=D_k$, gives
$kD_k\to1/2$ almost surely. This is the Bergman radial law and proves
the second limit in \eqref{eq:raduniq-incompatible-depths}. The two
almost-sure limits define disjoint measurable sets of depth sequences,
which proves the stated mutual singularity.

We give a direct argument showing why this singularity actually
precludes uniform integrability of $(H_N)$; uniform integrability alone
would not imply $L^1$ convergence of an arbitrary sequence. Define
$Q_N$ on the common reference space by $dQ_N=H_N\,d\mu_\infty$.
By \eqref{eq:raduniq-ratio-density}, the first $N$ coordinates under
$Q_N$ are exactly the finite-$N$ gap variables. In particular,
\[
 D_{k,N}:=\sum_{j=k}^N X_j
\]
under $Q_N$ has the law of the $k$th largest of $Y_1,\ldots,Y_N$.
For every fixed $k$, the latter increases almost surely to $D_k$ as
$N\to\infty$.

For $N\ge k$, set
\[
 A_{k,N}=\left\{
 \left|k\sum_{j=k}^N X_j-\frac12\right|<\frac18\right\}.
\]
For fixed $k$, as $N\to\infty$,
\begin{align*}
 \mu_\infty(A_{k,N})
 &\longrightarrow
 \mu_\infty\left(\left|kS_k-\frac12\right|<\frac18\right),\\
 Q_N(A_{k,N})
 &\longrightarrow
 \mathbb P\left(\left|kD_k-\frac12\right|<\frac18\right).
\end{align*}
There are no boundary atoms: $S_k$ is the sum of an exponential
variable and an independent variable, and $D_k$ is one of a countable
family of continuously distributed variables. The two right-hand
sides tend respectively to $0$ and $1$ as $k\to\infty$, by
\eqref{eq:raduniq-incompatible-depths}. Choose an increasing sequence
$N_k\ge k$ large enough that both preceding finite-$N$ probabilities
are within $1/k$ of their respective limits. Then
\[
 \mu_\infty(A_{k,N_k})\longrightarrow0,
 \qquad
 \int_{A_{k,N_k}}H_{N_k}\,d\mu_\infty
 =Q_{N_k}(A_{k,N_k})\longrightarrow1.
\]
This contradicts uniform absolute continuity of the integrals of a
uniformly integrable family. Hence $(H_N)$ is not uniformly integrable.
\end{proof}

\begin{remark}
The positive expansion supplies an exact finite-$N$ mixture and a
stronger form of radial domination. It does not establish convergence
of the normalized mixing probabilities $\pi_N$, nor does it identify
the angular law of a possible limit. The preceding obstructions concern
specific monotonicity, martingale, and uniform-integrability arguments;
they neither prove nor disprove uniqueness for general $\beta$.
\end{remark}



\section{Partition functions and a vanishing central-density subsequence}

Throughout this section, $dA$ is ordinary planar Lebesgue measure, so
$A(\mathbb D)=\pi$.  Set
\[
 Z_{N,\beta}=\int_{\mathbb D^N}|\Delta_N(z)|^\beta\,dA^{\otimes N}(z),
 \qquad Q_{N,\beta}=\frac{Z_{N,\beta}}{N!},
 \qquad Q_{0,\beta}=Z_{0,\beta}=1.
\]
Here $\rho_{N,\beta}=\rho_{1,N}$ denotes
the one-point correlation density relative to $dA$.

\begin{theorem}[Leading partition-function asymptotics]
\label{thm:partition-leading}
For every fixed $\beta>0$,
\begin{equation}
 \log Q_{N,\beta}
 =\left(\frac\beta2-2\right)N\log N+O_\beta(N),
 \qquad N\longrightarrow\infty.
 \label{eq:partition-leading}
\end{equation}
More explicitly, for $\beta\geq2$, $N\geq2$, and $s_N=1-1/N$,
\begin{equation}
 \pi^N(1-s_N^2)^N s_N^{\beta N(N-1)/2}
       (N!)^{\beta/2-1}
 \ \leq\ Q_{N,\beta}
 \ \leq\ \frac{\pi^N N^{(\beta/2-1)N}}{N!}.
 \label{eq:partition-bounds}
\end{equation}
For $0<\beta\leq2$ and $N\geq1$,
\begin{equation}
 \frac{\pi^N N^{(\beta/2-1)N}}{N!}
 \ \leq\ Q_{N,\beta}
 \ \leq\ \frac{(2\pi)^N(N!)^{\beta/2}}
 {\displaystyle\prod_{j=1}^N\left(2j+\frac\beta2j(j-1)\right)}.
 \label{eq:partition-bounds-low-beta}
\end{equation}
\end{theorem}

\begin{proof}
First, Hadamard's determinant inequality, applied to the Vandermonde matrix
whose $i$th row is $(1,z_i,\ldots,z_i^{N-1})$, gives
\begin{equation}
 |\Delta_N(z)|\leq N^{N/2},\qquad z\in\overline{\mathbb D}^{\,N}.
 \label{eq:hadamard-vandermonde}
\end{equation}
The exact identity $Z_{N,2}=\pi^N$ follows by expanding the two determinants
in $|\Delta_N|^2$ and using
\[
 \int_{\mathbb D}z^j\overline z^{\,k}\,dA(z)
 =\boldsymbol 1_{\{j=k\}}\frac\pi{j+1}.
\]
Indeed, the only surviving pairs of permutations coincide, and their total
contribution is
\[
 N!\prod_{j=0}^{N-1}\frac\pi{j+1}=\pi^N.
\]
Suppose first that $\beta\geq2$.  Then
\eqref{eq:hadamard-vandermonde} consequently yields
\[
 Z_{N,\beta}
 \leq N^{(\beta-2)N/2} Z_{N,2}
 =\pi^N N^{(\beta/2-1)N}.
\]
This proves the upper bound in \eqref{eq:partition-bounds}.

For the lower bound, let $dm(\theta)=d\theta/(2\pi)$ and define
\[
 A_{N,\beta}(r_1,\ldots,r_N)
 =\int_{[0,2\pi)^N}
 |\Delta_N(r_1e^{i\theta_1},\ldots,r_Ne^{i\theta_N})|^\beta
 \,dm^{\otimes N}(\theta).
\]
For fixed values of all coordinates except the $i$th, the integrand is
the $\beta$th power of the modulus of a polynomial in the $i$th complex
coordinate.  It is therefore subharmonic in that coordinate.  The circular
mean of a subharmonic function on a disk is nondecreasing in its radius.
Integrating over the other angles proves that $A_{N,\beta}$ is
nondecreasing in each $r_i$.  By homogeneity, for $r_i\in[s_N,1]$,
\[
 A_{N,\beta}(r_1,\ldots,r_N)
 \geq s_N^{\beta N(N-1)/2}A_{N,\beta}(1,\ldots,1).
\]
On the torus, orthogonality of the monomials similarly gives
\[
 A_{N,2}(1,\ldots,1)=N!.
\]
Jensen's inequality for the convex function $x\mapsto x^{\beta/2}$,
applied to normalized angular measure, now implies
\[
 A_{N,\beta}(1,\ldots,1)\geq(N!)^{\beta/2}.
\]
Integrate over the product of the annuli $s_N\leq|z_i|<1$.  Since
$dA=2\pi r\,dr\,dm(\theta)$ and each annulus has area
$\pi(1-s_N^2)$, we obtain
\[
 Z_{N,\beta}
 \geq \pi^N(1-s_N^2)^N
        s_N^{\beta N(N-1)/2}(N!)^{\beta/2}.
\]
Dividing by $N!$ proves the lower bound.

For these values of $\beta$,
\[
 \log(1-s_N^2)=-\log N+O(1),\qquad
 \frac{\beta N(N-1)}2\log s_N=O_\beta(N),
\]
and $\log(N!)=N\log N-N+O(\log N)$.  Substitution into the two bounds
proves \eqref{eq:partition-leading} when $\beta\geq2$.

Now suppose that $0<\beta\leq2$.  Since $\beta-2\leq0$, the direction
of the pointwise estimate from Hadamard's inequality is reversed:
\[
 |\Delta_N(z)|^\beta
 \geq N^{(\beta/2-1)N}|\Delta_N(z)|^2.
\]
This inequality first follows where $\Delta_N(z)\ne0$ and extends to its
zero set because both sides vanish there.  Integration and $Z_{N,2}=\pi^N$
give the lower bound in \eqref{eq:partition-bounds-low-beta}.

To prove its upper bound, use the radial-ratio representation.  With
ordered radii $0<r_1<\cdots<r_N<1$, put $r_{N+1}=1$ and
$t_j=r_j/r_{j+1}$.  In the notation of the radial comparison lemma,
\begin{equation}
 Q_{N,\beta}
 =(2\pi)^N\int_{(0,1)^N}
 F_{N,\beta}(t)\prod_{j=1}^N t_j^{a_j-1}\,dt,
 \qquad a_j=2j+\frac\beta2j(j-1).
 \label{eq:partition-radial-representation}
\end{equation}
The factor $N!$ from ordering the radii has been canceled by the definition
of $Q_{N,\beta}$; the factor $(2\pi)^N$ converts normalized angular
measure to ordinary polar area measure.  Recall that
\[
 F_{N,\beta}(t)
 =\int_{[0,2\pi)^N}
 \prod_{i<j}\left|
 1-\left(\prod_{h=i}^{j-1}t_h\right)e^{i(\theta_i-\theta_j)}
 \right|^\beta\,dm^{\otimes N}(\theta).
\]
The radial comparison lemma proves that this factor is nondecreasing in
each coordinate.  It extends continuously to $[0,1]^N$, as follows
directly from the displayed integral.  Therefore
\[
 F_{N,\beta}(t)\leq F_{N,\beta}(1,\ldots,1)
 =A_{N,\beta}(1,\ldots,1)
 \leq (N!)^{\beta/2}.
\]
The last inequality is Jensen's inequality for the concave function
$x\mapsto x^{\beta/2}$ on normalized angular measure, together with
$A_{N,2}(1,\ldots,1)=N!$.  Inserting this bound into
\eqref{eq:partition-radial-representation} and integrating each power
$t_j^{a_j-1}$ proves \eqref{eq:partition-bounds-low-beta}.

For clarity, the product in this upper bound can also be evaluated exactly.
Writing $b=4/\beta$,
\[
 \prod_{j=1}^N a_j
 =\left(\frac\beta2\right)^N N!
       \frac{\Gamma(N+b)}{\Gamma(b)}.
\]
Consequently the upper bound equals
\[
 \left(\frac{4\pi}{\beta}\right)^N
       (N!)^{\beta/2-1}\frac{\Gamma(b)}{\Gamma(N+b)}.
\]
Stirling's formula, with fixed $b>0$, shows that its logarithm is
$(\beta/2-2)N\log N+O_\beta(N)$.  The logarithm of the lower bound
in \eqref{eq:partition-bounds-low-beta} has the same leading term and
an $O_\beta(N)$ remainder.  This completes the proof for all $\beta>0$.
\end{proof}

\begin{proposition}[Exact central insertion and a telescoping lower bound]
\label{prop:central-insertion}
Fix $\beta>0$, and write $b=4/\beta$.  For every $N\geq1$,
\begin{equation}
 \rho_{N,\beta}(0)
 =\frac{Q_{N-1,\beta}}{Q_{N,\beta}}
   \mathbb E_{N-1,\beta}
     \left[\prod_{j=1}^{N-1}|z_j|^\beta\right].
 \label{eq:central-insertion-exact}
\end{equation}
In particular,
\begin{equation}
 \frac{b(b+1)}{(N-1+b)(N+b)}
       \frac{Q_{N-1,\beta}}{Q_{N,\beta}}
 \ \leq\ \rho_{N,\beta}(0)
 \ \leq\ \frac{Q_{N-1,\beta}}{Q_{N,\beta}}.
 \label{eq:central-insertion-bounds}
\end{equation}
The empty product in \eqref{eq:central-insertion-exact} equals one.
\end{proposition}

\begin{proof}
Fixing one particle at the origin in the definition of the one-point
correlation density gives
\[
 \rho_{N,\beta}(0)=\frac{N}{Z_{N,\beta}}
 \int_{\mathbb D^{N-1}}|\Delta_{N-1}(w)|^\beta
       \prod_{j=1}^{N-1}|w_j|^\beta\,dA^{\otimes(N-1)}(w).
\]
The identity $NZ_{N-1,\beta}/Z_{N,\beta}
=Q_{N-1,\beta}/Q_{N,\beta}$ proves
\eqref{eq:central-insertion-exact}; its upper bound follows from
$|w_j|\leq1$.

For completeness, apply the increasing-tilt comparison for the radial
ratios to an $n$-particle configuration.  Write its ordered radii as
$0<r_1<\cdots<r_n<1$, put $r_{n+1}=1$ and $t_k=r_k/r_{k+1}$, and recall
\[
 a_k=2k+\frac\beta2k(k-1).
\]
That comparison says that the actual ratio law has the form of the
product of the $\operatorname{Beta}(a_k,1)$ laws tilted by a
coordinatewise nondecreasing function.  Harris's inequality for product
measures consequently gives
$\mathbb E[G(t)]\geq\mathbb E[G(T)]$ for each bounded nondecreasing
$G$, where the $T_k$ are independent with those beta laws.
Use the increasing function
\[
 G(t)=\prod_{j=1}^n r_j^\beta=\prod_{k=1}^n t_k^{\beta k}.
\]
The beta moments and a telescoping product give
\begin{align*}
 \mathbb E_{n,\beta}\left[\prod_{j=1}^n|z_j|^\beta\right]
 &\geq\prod_{k=1}^n\frac{a_k}{a_k+\beta k}
 =\prod_{k=1}^n\frac{k-1+b}{k+1+b}\\
 &=\frac{b(b+1)}{(n+b)(n+b+1)}.
\end{align*}
For $n=0$ the same formula equals one.  Taking $n=N-1$ in
\eqref{eq:central-insertion-exact} proves the lower bound.
\end{proof}

\begin{theorem}[A central-density subsequence for $\beta>4$]
\label{thm:central-subsequence}
Fix $\beta>4$ and set $c_\beta=\beta/2-2>0$.  For every
$\varepsilon\in(0,c_\beta)$ there are infinitely many integers $N$ such
that
\begin{equation}
 \rho_{N,\beta}(0)\leq N^{-c_\beta+\varepsilon}.
 \label{eq:central-subsequence-rate}
\end{equation}
In particular, $\liminf_{N\to\infty}\rho_{N,\beta}(0)=0$.
\end{theorem}

\begin{proof}
The lower bound from Theorem~\ref{thm:partition-leading} gives a constant
$C_\beta$ such that
\[
 \log Q_{N,\beta}\geq c_\beta N\log N-C_\beta N.
\]
Suppose, to obtain a contradiction, that only finitely many $N$ satisfy
\[
 \frac{Q_{N,\beta}}{Q_{N-1,\beta}}
 \geq N^{c_\beta-\varepsilon}.
\]
Then this inequality is reversed strictly for all $N\geq N_0$ for some
$N_0$.  Summing logarithms of the ratios gives
\begin{align*}
 \log Q_{N,\beta}
 &\leq\log Q_{N_0-1,\beta}
   +(c_\beta-\varepsilon)\sum_{n=N_0}^N\log n\\
 &=(c_\beta-\varepsilon)N\log N+O_{\beta,\varepsilon,N_0}(N).
\end{align*}
This contradicts the preceding lower bound for large $N$.  Thus the
displayed ratio bound holds for infinitely many $N$, and the upper bound
in \eqref{eq:central-insertion-bounds} proves
\eqref{eq:central-subsequence-rate}.
\end{proof}

\paragraph{What this does and does not imply.}
If one uses the previously established compactness of the correlation
densities in the topology of local uniform convergence, the indices in
Theorem~\ref{thm:central-subsequence} admit a further subsequence whose
limiting continuous one-point density satisfies $\rho_\infty(0)=0$.
This statement by itself does not imply that the limiting point process
is empty: a nonnegative radial subharmonic function may vanish at the
origin without vanishing elsewhere.  The theorem also does not assert
full-sequence convergence of the central densities.

In particular, the $O_\beta(N)$ remainder in
\eqref{eq:partition-leading} cannot be differenced as though it were a
smooth asymptotic expansion.  For example, for fixed real $c$ and $a>0$,
the positive numerical sequence
\[
 B_N=\exp\{cN\log N+aN(-1)^N\},\qquad N\geq1,
\]
satisfies $\log B_N=cN\log N+O(N)$, whereas
\[
 \log\frac{B_N}{B_{N-1}}
 =c\bigl(N\log N-(N-1)\log(N-1)\bigr)
   +a(2N-1)(-1)^N
\]
for $N\geq2$.  Its adjacent ratios therefore have arbitrarily large
oscillations.  This example concerns only what an $O(N)$ asymptotic
allows; it is not a proposed behavior of the gas partition functions.
The later harmonic boundary transport supplies the additional
adjacent-ratio lower bound, while the uniform analytic exponential-moment
bound controls compact insertions. Together they strengthen the central
subsequence statement to a full empty-limit theorem for $\beta>4$.
The elementary partition asymptotic alone does not give this conclusion.


\section{A global Gaussian mode independent of every interior limit}

Write $S_{k,N}=\sum_{i=1}^N z_i^k$ and
\[
 c_N=2+\frac\beta2(N-1).
\]
The following statement identifies a global fluctuation without assuming
uniqueness of the interior point process.

\begin{theorem}[Gaussian center of charge and independence]
For every $\beta>0$,
\[
 \mathbb E S_{1,N}=0,\qquad
 \mathbb E|S_{1,N}|^2=\frac{N}{c_N}.
\]
Moreover, $S_{1,N}$ converges in law to a centered circular complex Gaussian
$G_\beta$ with density
\[
 \frac{\beta}{2\pi}\exp\left(-\frac\beta2|g|^2\right)dA(g).
\]
If $\Xi_{N_j}\Rightarrow\Xi$ vaguely on the open disk, then
\[
 (\Xi_{N_j},S_{1,N_j})\Rightarrow(\Xi,G_\beta),
\]
where the two variables on the right are independent.
For every fixed $t\in\mathbb C$ and every bounded continuous functional $F$
on the space of locally finite counting measures with its vague topology,
\[
 \lim_{j\to\infty}
 \mathbb E\left[F(\Xi_{N_j})e^{-2\operatorname{Re}(\bar tS_{1,N_j})}\right]
 =e^{2|t|^2/\beta}\mathbb EF(\Xi).
\]
\end{theorem}

\begin{proof}
We first record an elementary macroscopic estimate. For every fixed integer
$k\ge1$ and every $\varepsilon>0$,
\begin{equation}\label{eq:global-mode-concentration}
 \mathbb P(|S_{k,N}|\ge\varepsilon N)
 \le \exp\left[
 (1+\beta/2)N\log N+2\beta N
 -\frac{\beta4^{-k}}{2k}\varepsilon^2N^2\right],\qquad N\ge2.
\end{equation}
To prove it, put $q=1-1/N$. The annulus $q<|z|<1$ has area at least
$\pi/N$. Under normalized area measure on that annulus, independent
$Z,W$ satisfy
\[
 \mathbb E\log|Z-W|
 =\mathbb E\log\max(|Z|,|W|)\ge\log q\ge-2/N.
\]
The equality follows by angular integration of the logarithmic kernel.
Jensen's inequality applied to the partition function restricted to this
annulus gives
\begin{equation}\label{eq:annular-partition-lower}
 Z_{N,\beta}\ge(\pi/N)^N e^{-\beta(N-1)}.
\end{equation}
For arbitrary points in the disk the identity
\[
 |1-q^2z_i\bar z_j|^2-q^2|z_i-z_j|^2
 =(1-q^2|z_i|^2)(1-q^2|z_j|^2)\ge0
\]
implies
\begin{align*}
 \log|\Delta_N(z)|
 &\le-\frac{N(N-1)}2\log q
   -\frac12\sum_{\ell\ge1}\frac{q^{2\ell}}\ell
       \left(|S_{\ell,N}|^2-\sum_i|z_i|^{2\ell}\right)\\
 &\le N+\frac N2\log N-\frac{4^{-k}}{2k}|S_{k,N}|^2.
\end{align*}
The series is absolutely convergent since $q<1$. In the second line we
used $q\ge1/2$, $-\log q\le2/N$, and
$1-q^2\ge1/N$. Integrating this upper bound over a subset of the disk
product and using \eqref{eq:annular-partition-lower} proves
\eqref{eq:global-mode-concentration}. In particular,
$S_{2,N}/N\to0$ in probability.

For $a\in\mathbb D$, let
\[
 T_a(z)=\frac{z+a}{1+\bar az},\qquad
 R_{N,a}(z)=
 (1-|a|^2)^{c_NN}\prod_{i=1}^N|1+\bar az_i|^{-2c_N}.
\]
The distance and Jacobian formulas for this disk automorphism are
\[
 |T_a(z)-T_a(w)|
 =\frac{(1-|a|^2)|z-w|}{|1+\bar az|\,|1+\bar aw|},\qquad
 |T_a'(z)|^2=\frac{(1-|a|^2)^2}{|1+\bar az|^4}.
\]
A change of variables therefore gives, for every bounded measurable $F$,
\begin{equation}\label{eq:mobius-exact-test}
 \mathbb E\left[F((T_a)_*\Xi_N)R_{N,a}\right]
 =\mathbb EF(\Xi_N).
\end{equation}
In particular, $\mathbb ER_{N,a}=1$. Expansion at $a=0$ is legitimate
under the integral for each fixed $N$, since all derivatives are uniformly
bounded on a sufficiently small closed disk in the $a$ variable. The
coefficient of $|a|^2$ in $R_{N,a}$ is
$c_N^2|S_{1,N}|^2-c_NN$. Its expectation vanishes, proving the variance
identity. The mean vanishes by rotation symmetry.

Fix $t\in\mathbb C$ and set $a_N=t/c_N$ for all sufficiently large $N$.
Taylor expansion, uniformly over all particle configurations, yields
\begin{equation}\label{eq:mobius-small-tilt}
 \log R_{N,a_N}
 =-2\operatorname{Re}(\bar tS_{1,N})
   -\frac{N}{c_N}|t|^2
   +\frac1{c_N}\operatorname{Re}(\bar t^{\,2}S_{2,N})
   +O_t(N/c_N^2).
\end{equation}
Here $c_N\asymp_\beta N$, and the error bound holds deterministically.
Thus, writing
\[
 Q_{N,t}=\exp\left[-2\operatorname{Re}(\bar tS_{1,N})
                        -\frac{N}{c_N}|t|^2\right],
\]
we have $\log(R_{N,a_N}/Q_{N,t})\to0$ in probability. This logarithm
is also bounded in absolute value by a constant depending only on $t$
and $\beta$, uniformly for all sufficiently large $N$.

For clarity, uniform integrability does not need to be assumed here.
The deterministic bounded-error estimate and $\mathbb ER_{N,a_N}=1$
give
\[
 \sup_{N\ge N_t}
 \mathbb E e^{-2\operatorname{Re}(\bar tS_{1,N})}<\infty.
\]
The same argument with $pt$ in place of $t$ gives the bound for every
fixed real $p>1$. Consequently, both $(R_{N,a_N})_N$ and $(Q_{N,t})_N$
have uniformly bounded $p$th moments. Their bounded logarithmic ratio
converges to zero in probability, so
\begin{equation}\label{eq:tilt-l1}
 \mathbb E|R_{N,a_N}-Q_{N,t}|\longrightarrow0.
\end{equation}
For example, this follows from H\"older's inequality applied to
$Q_{N,t}|\exp(\log(R_{N,a_N}/Q_{N,t}))-1|$.

The tiny automorphisms do not change the vague limit. Indeed, for
$f\in C_c(\mathbb D)$, choose a compact $K'\subset\mathbb D$ containing
both $\operatorname{supp}f$ and $T_{a_N}^{-1}(\operatorname{supp}f)$ for
all sufficiently large $N$. Uniform continuity of the zero extension of
$f$ and uniform
convergence $T_{a_N}\to\mathrm{id}$ on $K'$ give
\[
 \left|\int f\,d((T_{a_N})_*\Xi_N)-\int f\,d\Xi_N\right|
 \le\omega_f(C|a_N|)\,\Xi_N(K')\longrightarrow0
\]
in probability, by the already proved tightness of compact-disk counts.
Applying this to a countable family defining a metric for the vague
topology shows that the distance between these two configurations tends
to zero in probability. Since $(\Xi_N)_N$ is tight, every $F\in C_b$
therefore satisfies
\[
 F((T_{a_N})_*\Xi_N)-F(\Xi_N)\longrightarrow0
\]
in probability. Boundedness of $F$ and the uniform $p$th moment bound
for $R_{N,a_N}$ make the corresponding weighted difference tend to zero
in $L^1$.

Combining this observation with \eqref{eq:mobius-exact-test} and
\eqref{eq:tilt-l1} gives
\[
 \mathbb E\left[F(\Xi_N)e^{-2\operatorname{Re}(\bar tS_{1,N})}\right]
 -\exp\left(\frac{N}{c_N}|t|^2\right)\mathbb EF(\Xi_N)
 \longrightarrow0.
\]
Along $N_j$, this proves the asserted mixed Laplace-transform limit.
The variance identity makes $S_{1,N}$ tight. Every joint subsequential
limit of $(\Xi_{N_j},S_{1,N_j})$ satisfies the same transform identity,
because the exponential moment bounds justify passage to the limit.
Taking $F=1$, uniqueness of the moment-generating function in
$\mathbb R^2$ identifies its second coordinate as $G_\beta$. For each
nonnegative bounded continuous $F$ with positive expectation, the
probability measure on that second coordinate obtained by weighting by
$F$ has the same Gaussian moment-generating function. It is therefore
the same Gaussian law. Since bounded continuous functionals determine
probability measures on the configuration space, this proves
independence. Every joint subsequential limit has consequently been
identified, which proves the joint convergence along $N_j$.
\end{proof}

\begin{remark}
The theorem does not identify the law of $\Xi$. It shows, in particular,
that identifying the first global harmonic mode cannot by itself establish
uniqueness of the interior limit: that Gaussian variable is asymptotically
independent of the entire interior configuration.
\end{remark}


% This file is intentionally independent of the probability/induction proof.
\section{A disk transport for each harmonic mode}
\label{sec:higher-mode-transport}

Fix an integer $k\geq 1$. Throughout this section $D=\{z:|z|<1\}$,
$S_j=\sum_{i=1}^N z_i^j$ for every integer $j$ for which this expression
is defined, and $S_0=N$. Negative powers are used only on configurations
without a point at zero. All constants in the deterministic estimates may
depend on $k$ and on a fixed smooth cutoff, but not on $N,r,a$, or the
configuration.

\begin{lemma}[A smooth global extension of a branched M\"obius map]
\label{lem:patched-lift}
There are constants $c_k,C_k>0$ with the following property. If
$0<r<1/4$ and $|a|r^{-k}<c_k$, there is a smooth orientation-preserving
diffeomorphism $T_{a,r}:\overline D\to\overline D$ which is the identity
on $\{|z|\leq r\}$ and agrees on $\{|z|\geq 2r\}$ with
\begin{equation}
 f_a(z)=z\,(1+a z^{-k})^{1/k}(1+\bar a z^k)^{-1/k}.
 \label{eq:lift-map}
\end{equation}
Here both fractional powers are defined by the logarithmic power series
at $1$. Moreover,
\begin{equation}
 \|T_{a,r}-\mathrm{id}\|_\infty\leq C_k|a|r^{1-k},\qquad
 \|D T_{a,r}-I\|_\infty\leq C_k|a|r^{-k}.
 \label{eq:patched-c1}
\end{equation}
If $\eta=C_k|a|r^{-k}<1$, then
\begin{equation}
 (1-\eta)(1-|z|)\leq 1-|T_{a,r}(z)|
 \leq(1+\eta)(1-|z|),\qquad z\in\overline D.
 \label{eq:distance-boundary}
\end{equation}
After decreasing $c_k$ if necessary,
$T_{a,r}(\{|z|<2r\})\subset\{|z|<3r\}$.
\end{lemma}

\begin{proof}
Choose a smooth function $\chi:[0,\infty)\to[0,1]$ which is zero on
$[0,1]$ and one on $[2,\infty)$. On $|z|>r$ define
\[
 T_{a,r}(z)=z+\chi(|z|/r)(f_a(z)-z),
\]
and set $T_{a,r}(z)=z$ on $|z|\leq r$. The flatness of $\chi$ at the
end of its constant interval makes this a smooth map across $|z|=r$.
On $r\leq|z|\leq1$, the power series in \eqref{eq:lift-map} and their
first derivatives give
\[
 |f_a(z)-z|\leq C_k|a||z|^{1-k},\qquad
 |f'_a(z)-1|\leq C_k|a||z|^{-k}.
\]
The product rule, including the derivative $r^{-1}\chi'(|z|/r)$,
proves \eqref{eq:patched-c1}. The estimates remain true on the closed
disk. Since the disk is convex, the derivative bound implies
\[
 (1-\eta)|z-w|\leq |T_{a,r}(z)-T_{a,r}(w)|
 \leq(1+\eta)|z-w|.
\]
In particular the map is injective, has nonsingular derivative, and is
orientation preserving when $\eta<1$.

The identity
\begin{equation}
 f_a(z)^k=\frac{z^k+a}{1+\bar a z^k}
 \label{eq:power-mobius}
\end{equation}
shows that $|f_a(z)|=1$ on $|z|=1$ and $|f_a(z)|<1$ where $|z|<1$
and the formula is used. Since the patched map is a convex combination
of $z$ and $f_a(z)$, it maps $D$ into $D$ and the boundary into the
boundary. Its interior image is nonempty and open by the inverse
function theorem. It is also relatively closed in $D$: if
$T_{a,r}(z_n)\to w\in D$, choose a subsequence with
$z_n\to z\in\overline D$. The point $z$ cannot be on the boundary,
since its image would then be on the boundary. Thus
$w=T_{a,r}(z)\in T_{a,r}(D)$. Connectedness of $D$ proves that the
interior image is $D$. Every boundary point is also in the image by
approximating it from inside $D$ and taking a convergent subsequence
of preimages. The inverse function theorem, including an extension of
the formula to a neighborhood of the boundary, proves the asserted
diffeomorphism property.
Taking the infimum of the two-sided Lipschitz estimate over
$w\in\partial D$, whose images fill $\partial D$, proves
\eqref{eq:distance-boundary}. Finally, on $|z|<2r$ the displacement is
at most $C_k|a|r^{1-k}=C_k(|a|r^{-k})r$; decreasing $c_k$ proves the
last assertion.
\end{proof}

\begin{lemma}[Uniform divided-difference estimates]
\label{lem:divided-difference-expansion}
For distinct $z,w$ with $r\leq|z|,|w|\leq1$, put
\[
 D_m(z,w)=\frac{z^m-w^m}{z-w},\qquad m\in\mathbb Z.
\]
If $|a|r^{-k}$ is sufficiently small, then the logarithm below is the
branch near zero and
\begin{align}
 \log\frac{f_a(z)-f_a(w)}{z-w}
 ={}&\frac{a}{k}D_{1-k}-\frac{\bar a}{k}D_{1+k}
 \nonumber\\
 &+\frac{a^2}{2k^2}\{(1-k)D_{1-2k}-D_{1-k}^2\}
 \nonumber\\
 &+\frac{\bar a^2}{2k^2}\{(k+1)D_{1+2k}-D_{1+k}^2\}
 \nonumber\\
 &+\frac{|a|^2}{k^2}\{-1+D_{1-k}D_{1+k}\}
 +O_k(|a|^3r^{-3k}).
 \label{eq:pair-expansion}
\end{align}
The error is uniform even when $z$ and $w$ approach one another.
Furthermore,
\begin{align}
 \log f'_a(z)
 ={}&-\frac{k-1}{k}a z^{-k}-\frac{k+1}{k}\bar a z^k
       -|a|^2
 \nonumber\\
 &+\frac{k-1}{2k}a^2z^{-2k}
       +\frac{k+1}{2k}\bar a^2z^{2k}
       +O_k(|a|^3r^{-3k}).
 \label{eq:jacobian-expansion}
\end{align}
\end{lemma}

\begin{proof}
For $m\geq1$, the polynomial identity for a divided difference gives
$|D_m(z,w)|\leq m$. For $m=-q\leq-1$, one has the exact identity
\[
 D_{-q}(z,w)=-\sum_{\ell=1}^q z^{-\ell}w^{-(q+1-\ell)},
\]
and therefore $|D_{-q}(z,w)|\leq q r^{-(q+1)}$. Also $D_0=0$.
Consequently, for $u,v\geq0$ with $u+v=n$,
\begin{equation}
 |D_{1+k(v-u)}(z,w)|\leq(1+kn)r^{-kn}.
 \label{eq:general-dd-bound}
\end{equation}

Expand the two binomial factors in \eqref{eq:lift-map}. This gives an
absolutely convergent series whose $(u,v)$ term is
$b_u c_v a^u\bar a^v z^{1+k(v-u)}$, where
$b_u=\binom{1/k}{u}$ and $c_v=\binom{-1/k}{v}$.
There is a constant $C_k$ such that $|b_u|+|c_u|\leq C_k2^u$ for all
$u\geq0$; for instance Cauchy's coefficient bound on the circle of
radius $1/2$ gives this estimate. After taking divided differences,
\eqref{eq:general-dd-bound} implies that the sum of terms of total
degree at least three is $O_k((|a|r^{-k})^3)$, provided $|a|r^{-k}$ is
smaller than a fixed constant. The terms of degree at most two give
\begin{align*}
 \frac{f_a(z)-f_a(w)}{z-w}
 ={}&1+\frac{a}{k}D_{1-k}-\frac{\bar a}{k}D_{1+k}
 +\frac{1-k}{2k^2}a^2D_{1-2k}\\
 &+\frac{k+1}{2k^2}\bar a^2D_{1+2k}
 -\frac{|a|^2}{k^2}+O_k(|a|^3r^{-3k}).
\end{align*}
The perturbation of $1$ is $O_k(|a|r^{-k})$. Applying
$\log(1+u)=u-u^2/2+O(|u|^3)$ proves \eqref{eq:pair-expansion}.
This argument uses endpoint monomial identities and does not assume that
the segment from $z$ to $w$ avoids the origin.

Differentiating \eqref{eq:power-mobius} gives the exact formula
\[
 f'_a(z)=(1-|a|^2)(f_a(z)/z)^{1-k}(1+\bar a z^k)^{-2}.
\]
Its logarithm, with the branches equal to zero when $a=0$, is
\[
 \log f'_a(z)=\log(1-|a|^2)
 -\frac{k-1}{k}\log(1+a z^{-k})
 -\frac{k+1}{k}\log(1+\bar a z^k).
\]
Taylor expansion proves \eqref{eq:jacobian-expansion}.
\end{proof}

\begin{proposition}[All aggregate quadratic coefficients]
\label{prop:aggregate-coefficients}
Suppose the points $z_1,\ldots,z_N$ are distinct and
$r\leq|z_i|\leq1$. With logarithms chosen near zero, set
\[
 L_N=\sum_{i<j}\log\frac{f_a(z_i)-f_a(z_j)}{z_i-z_j},
 \qquad J_N=\sum_i\log f'_a(z_i).
\]
Then
\begin{align}
 L_N={}&\frac{a}{2k}
 \left\{(k-1)S_{-k}-\sum_{p=1}^{k-1}S_{-p}S_{-(k-p)}\right\}
 \nonumber\\
 &-\frac{\bar a}{2k}
 \left\{(2N-k-1)S_k+\sum_{p=1}^{k-1}S_pS_{k-p}\right\}
 \nonumber\\
 &+\frac{a^2}{4k^2}\sum_{p=1}^{2k-1}|k-p|
       \{S_{-p}S_{-(2k-p)}-S_{-2k}\}
 \nonumber\\
 &+\frac{\bar a^2}{4k^2}\sum_{p=0}^{2k}|k-p|
       \{S_pS_{2k-p}-S_{2k}\}
 \nonumber\\
 &+\frac{|a|^2}{k^2}
 \left\{-\frac{kN^2}{2}+\frac{k^2N}{2}
 -\sum_{p=1}^{k-1}(k-p)S_pS_{-p}\right\}
 +O_k(N^2|a|^3r^{-3k}),
 \label{eq:aggregate-L}
\end{align}
and
\begin{align}
 J_N={}&-\frac{k-1}{k}aS_{-k}-\frac{k+1}{k}\bar aS_k
       -N|a|^2
 \nonumber\\
 &+\frac{k-1}{2k}a^2S_{-2k}
       +\frac{k+1}{2k}\bar a^2S_{2k}
       +O_k(N|a|^3r^{-3k}).
 \label{eq:aggregate-J}
\end{align}
Empty sums are zero, so these formulas also cover $k=1$.
\end{proposition}

\begin{proof}
The linear terms follow from
\begin{align*}
 \sum_{i<j}D_{1+k}(z_i,z_j)
 &=\frac12\left\{(2N-k-1)S_k+
                         \sum_{p=1}^{k-1}S_pS_{k-p}\right\},\\
 \sum_{i<j}D_{1-k}(z_i,z_j)
 &=-\frac12\left\{\sum_{p=1}^{k-1}S_{-p}S_{-(k-p)}
                                     -(k-1)S_{-k}\right\}.
\end{align*}
To verify them, expand each divided difference into monomials and use
\[
 \sum_{i\ne j}z_i^u z_j^v=S_uS_v-S_{u+v}.
\]
The coefficient arrays are symmetric in $u,v$, so an unordered pair
sum is one half of the corresponding ordered pair sum.

For clarity the quadratic identities at a single pair are
\begin{align*}
 (1-k)D_{1-2k}-D_{1-k}^2
 &=\sum_{p=1}^{2k-1}|k-p|z^{-p}w^{-(2k-p)},\\
 (k+1)D_{1+2k}-D_{1+k}^2
 &=\sum_{p=0}^{2k}|k-p|z^p w^{2k-p},\\
 D_{1-k}D_{1+k}
 &=-(k-1)-\sum_{p=1}^{k-1}(k-p)
                         (z^pw^{-p}+z^{-p}w^p).
\end{align*}
The first two identities follow by counting the coefficients in the
squares of finite geometric sums: the triangular convolution
coefficients, subtracted from $k-1$ or $k+1$, are exactly $|k-p|$.
For the last identity multiply the two geometric sums directly.
It follows that
\begin{align*}
 \sum_{i<j}\{-1+D_{1-k}(z_i,z_j)D_{1+k}(z_i,z_j)\}
 &=-\frac{kN(N-1)}2
       -\sum_{p=1}^{k-1}(k-p)(S_pS_{-p}-N)\\
 &=-\frac{kN^2}2+\frac{k^2N}2
       -\sum_{p=1}^{k-1}(k-p)S_pS_{-p}.
\end{align*}
Substitution in \eqref{eq:pair-expansion} proves
\eqref{eq:aggregate-L}. Summing \eqref{eq:jacobian-expansion} proves
\eqref{eq:aggregate-J}.
\end{proof}

\begin{lemma}[Negative powers controlled by radial deficits]
\label{lem:negative-powers}
Fix $0<\rho<1$ and an integer $j\geq1$. If all $|z_i|\geq r$ with
$0<r<\rho$, and if
$\mathcal D_N=\sum_i(1-|z_i|)$ and $M_N(\rho)=\#\{i:|z_i|<\rho\}$, then
\begin{equation}
 |S_{-j}-\overline{S_j}|
 \leq C_{j,\rho}\mathcal D_N+2r^{-j}M_N(\rho).
 \label{eq:negative-power-bound}
\end{equation}
\end{lemma}

\begin{proof}
For $z=s e^{i\theta}$,
$|z^{-j}-\bar z^j|=s^{-j}-s^j$. On $\rho\leq s\leq1$, the
mean-value theorem bounds this by $C_{j,\rho}(1-s)$. On
$r\leq s<\rho$, it is at most $2r^{-j}$. Sum these two bounds.
\end{proof}

\begin{proposition}[A probabilistic criterion for the transport expansion]
\label{prop:transport-asymptotic-criterion}
Let $k\geq1$, fix $t\in\mathbb C$, and choose
\[
 0<\alpha<\frac1{3k},\qquad r_N=N^{-\alpha},\qquad
 a_N=\frac{2kt}{\beta N},\qquad T_N=T_{a_N,r_N}.
\]
For finitely many small $N$ the map may be defined to be the identity.
Let $\mathcal L_N$ be any sequence of probability laws on configurations
of $N$ distinct points in $D$. Suppose, under these laws, that
\begin{enumerate}
\item $\mathcal L_N\{\min_i|z_i|<2r_N\}\to0$;
\item $M_N(\rho)=O_{\mathcal L_N}(1)$ for one fixed $0<\rho<1$;
\item $\mathcal D_N=O_{\mathcal L_N}(\log(N+1))$;
\item $S_j/N\to0$ in probability for every $1\leq j\leq2k$;
\item $S_j=O_{\mathcal L_N}(1)$ for $1\leq j<k$.
\end{enumerate}
Define the deterministic likelihood function
\begin{equation}
 R_N(z)=
 \prod_{i<j}\left|\frac{T_N(z_i)-T_N(z_j)}{z_i-z_j}\right|^\beta
 \prod_i\det DT_N(z_i).
 \label{eq:transport-likelihood}
\end{equation}
Then
\begin{equation}
 \log R_N=-2\operatorname{Re}(\bar tS_k)
                       -\frac{2k}{\beta}|t|^2
                       +o_{\mathcal L_N}(1).
 \label{eq:higher-mode-likelihood-expansion}
\end{equation}
No claim about exponential moments of the error is part of this statement.
\end{proposition}

\begin{proof}
Work on the event $G_N=\{\min_i|z_i|\geq2r_N\}$, whose probability
tends to one. The preceding negative-power bound implies
\[
 S_{-j}=\overline{S_j}
       +O_{\mathcal L_N}(\log(N+1)+N^{\alpha j}),\qquad 1\leq j\leq2k.
\]
In particular $S_{-j}=o_{\mathcal L_N}(N)$ for $j\leq2k$, since
$2k\alpha<1$. For $j<k$, the positive modes are tight, so
\[
 S_{-j}=O_{\mathcal L_N}(\log(N+1)+N^{\alpha j}).
\]
For $1\leq p<k$ this gives
\[
 N^{-1}S_{-p}S_{-(k-p)}=o_{\mathcal L_N}(1),
 \qquad N^{-1}S_pS_{k-p}=o_{\mathcal L_N}(1),
\]
because $\alpha k<1$. Consequently the first line of
\eqref{eq:aggregate-L} is $o_{\mathcal L_N}(1)$, and its second line
is $-\bar a_N N S_k/k+o_{\mathcal L_N}(1)$.

Every factor $S_j$ or $S_{-j}$ of nonzero degree appearing in the
quadratic terms is $o_{\mathcal L_N}(N)$. The terms involving $S_0=N$
are paired with $S_{2k}=o_{\mathcal L_N}(N)$. Thus the pure $a_N^2$
and $\bar a_N^2$ contributions are $o_{\mathcal L_N}(1)$.
The mixed contribution is
$-|a_N|^2N^2/(2k)+o_{\mathcal L_N}(1)$.
All terms in \eqref{eq:aggregate-J} are $o_{\mathcal L_N}(1)$.
Finally the largest deterministic remainder satisfies
\[
 N^2|a_N|^3r_N^{-3k}=O_{k,t,\beta}(N^{-1+3k\alpha})\longrightarrow0.
\]
On $G_N$ the map is holomorphic at every point, so its real Jacobian
is $|f'_{a_N}(z_i)|^2$ and
$\log R_N=\beta\operatorname{Re}L_N+2\operatorname{Re}J_N$.
Substitution of $a_N=2kt/(\beta N)$ gives
\eqref{eq:higher-mode-likelihood-expansion}. Discarding $G_N^c$ is
legitimate for convergence in probability.
\end{proof}

\begin{lemma}[Effect on power sums]
\label{lem:transport-mode-shifts}
For every fixed integer $j\geq1$, on $G_N$,
\begin{equation}
 S_j(T_Nz)-S_j(z)
 =\frac{j}{k}\{a_NS_{j-k}-\bar a_NS_{j+k}\}
       +O_{j,k}(N|a_N|^2r_N^{-2k}).
 \label{eq:mode-shift-expansion}
\end{equation}
In particular, under any laws satisfying the small-hole, radial-deficit,
and local-count assumptions above and satisfying $S_\ell/N\to0$ for
the positive index $\ell=j+k$ and, when $j\ne k$, also
$\ell=|j-k|$,
\[
 S_j(T_Nz)-S_j(z)\longrightarrow
 \begin{cases}2kt/\beta,&j=k,\\0,&j\ne k,\end{cases}
\]
in probability. Regardless of macroscopic-mode convergence, on $G_N$
there is the deterministic bound
\begin{equation}
 |S_k(T_Nz)-S_k(z)|\leq C_{k,t,\beta}.
 \label{eq:k-shift-uniform-bound}
\end{equation}
\end{lemma}

\begin{proof}
Raise \eqref{eq:lift-map} to the $j$th power and expand both binomial
factors to first order. The pointwise remainder is at most
$C_{j,k}|a_N|^2r_N^{-2k}$, because $|z|^j\leq1$.
This proves \eqref{eq:mode-shift-expansion}. Its total remainder is
$O(N^{-1+2k\alpha})=o(1)$. The negative modes with indices between
$1$ and $k-1$, when they occur, are $o(N)$ by
\eqref{eq:negative-power-bound}; the other nonconstant modes are
$o(N)$ by hypothesis. For $j=k$ the constant mode contributes
$a_NN=2kt/\beta$.
Finally, \eqref{eq:power-mobius} gives directly
\[
 |f_{a_N}(z)^k-z^k|
 =\left|\frac{a_N-\bar a_Nz^{2k}}{1+\bar a_Nz^k}\right|
 \leq\frac{2|a_N|}{1-|a_N|}.
\]
Summing proves \eqref{eq:k-shift-uniform-bound}.
\end{proof}

\begin{lemma}[Stability under the transported law]
\label{lem:transported-controls}
Let $P_N$ be the disk-gas law and let $Q_N$ have density
$R_N$ from \eqref{eq:transport-likelihood} with respect to $P_N$.
The change of variables gives $\mathbb E_{P_N}R_N=1$ and makes
$Y=T_N Z$ under $Q_N$ have law $P_N$. Assume under $P_N$ that
\[
 P_N\{M_N(s)>0\}\leq C s^2\quad(0<s<s_0),\qquad
 \mathcal D_N=O_{P_N}(\log(N+1)),
\]
that compact-disk counts are tight, that $S_j/N\to0$ for all fixed
positive $j$, and that $S_j$ is tight for $1\leq j<k$.
Then all five hypotheses of
Proposition~\ref{prop:transport-asymptotic-criterion} hold under $Q_N$,
as well as under $P_N$.
\end{lemma}

\begin{proof}
Under $Q_N$, $Y=T_N Z$ has the original law. If some $|z_i|<2r_N$,
then $|Y_i|<3r_N$. Consequently
$Q_N(G_N^c)\leq P_N\{M_N(3r_N)>0\}=O(r_N^2)$.
The uniform displacement in \eqref{eq:patched-c1} tends to zero.
Thus, for any fixed $\rho<\rho'<1$, and all sufficiently large $N$,
$M_Z(\rho)\leq M_Y(\rho')$. This proves tightness of local counts
under $Q_N$. Equation~\eqref{eq:distance-boundary} gives
$\mathcal D_N(Z)\leq(1-\eta_N)^{-1}\mathcal D_N(Y)$ with $\eta_N\to0$, which proves
the radial-deficit bound under $Q_N$. For each fixed positive $j$,
\[
 N^{-1}|S_j(Y)-S_j(Z)|
 \leq j\|T_N-\mathrm{id}\|_\infty\longrightarrow0,
\]
so macroscopic-mode convergence transfers too.

It remains to check tightness of the lower modes without assuming it
under $Q_N$. On $G_N$, the first-order bound on the binomial factors
gives, for $1\leq j<k$,
\[
 |f_{a_N}(z)^j-z^j|
 \leq C_{j,k}|a_N|(|z|^{j-k}+|z|^{j+k}).
\]
Splitting the sum at any fixed $0<\rho<1$ yields
\[
 |S_j(Y)-S_j(Z)|
 \leq C_{j,k,\rho}|a_N|
                   \{N+r_N^{j-k}M_Z(\rho)\}.
\]
The first term is bounded and the second is $o_{Q_N}(1)$ because
$\alpha(k-j)<1$ and the local count is tight. Since $S_j(Y)$ has the
original tight law and $Q_N(G_N^c)\to0$, this proves tightness of
$S_j(Z)$ under $Q_N$, completing the proof.
\end{proof}


% Probability half of the higher harmonic-mode theorem.
% To follow the deterministic transport lemma in higher_modes_transport_notes.tex.
\section{All harmonic modes are Gaussian and independent of the interior}

For integers $j$, set $S_{j,N}=\sum_{i=1}^N z_i^j$ whenever all $z_i\ne0$;
this holds almost surely. Set $S_{0,N}=N$. Let
$\mathcal D_N=\sum_i(1-|z_i|)$. Three estimates already proved for the disk gas will
be used:
\begin{enumerate}
\item For every fixed $0<r_0<1$, the counts $\Xi_N(B_{r_0})$ are tight,
and $\sup_N\sup_{|z|\le r_0}\rho_{1,N}(z)<\infty$.
\item $\mathcal D_N=O_{\mathbb P}(\log(N+1))$.
\item For every fixed positive integer $j$, $S_{j,N}/N\to0$ in probability.
\end{enumerate}
Here the second estimate follows directly from radial stochastic domination:
if $a_j=2j+\beta j(j-1)/2$ are the comparison exponential rates, then
\[
 \mathbb E\mathcal D_N\le\mathbb E\sum_i(-\log|z_i|)
 \le\sum_{j=1}^N\frac{j}{a_j}=O_\beta(\log(N+1)).
\]
The last estimate follows from the elementary exponential concentration
bound established above using the annular lower bound on $Z_{N,\beta}$.

\begin{theorem}[Joint Gaussian harmonic modes]\label{thm:all-harmonic-modes}
For every $\beta>0$ and every fixed integer $m\ge1$,
\[
 (S_{1,N},\ldots,S_{m,N})\ \Longrightarrow\ (G_1,\ldots,G_m),
\]
where the $G_j$ are independent centered circular complex Gaussians with
$\mathbb E|G_j|^2=2j/\beta$. If
$\Xi_{N_\ell}\Rightarrow\Xi$ in the vague topology of the open disk, then
\[
 (\Xi_{N_\ell},S_{1,N_\ell},\ldots,S_{m,N_\ell})
 \ \Longrightarrow\ (\Xi,G_1,\ldots,G_m),
\]
where the Gaussian vector is independent of $\Xi$.
\end{theorem}

\begin{proof}
The proof is by induction on the mode $k$. The already established
M\"obius argument supplies $k=1$. We give the induction step in detail.
Fix $k\ge2$, suppose that all conclusions have been proved through $k-1$,
and put
\[
 r_N=N^{-\alpha},\qquad \alpha=\frac1{100k},\qquad
 a_N=\frac{2kt}{\beta N},\quad t\in\mathbb C.
\]
Use the smooth disk diffeomorphism $T_{N,t}$ from the preceding transport
lemma. It agrees for $|z|\ge2r_N$ with
\[
 f_{a_N}(z)
 =z(1+a_Nz^{-k})^{1/k}(1+\bar a_Nz^k)^{-1/k},
\]
where the powers are defined by the analytic logarithms near $1$.
Its exact likelihood ratio under change of variables is
\[
 R_{N,t}(z)=
 \prod_i J_{T_{N,t}}(z_i)
 \prod_{i<j}\left|\frac{T_{N,t}(z_i)-T_{N,t}(z_j)}{z_i-z_j}\right|^\beta.
\]
Define $Q_{N,t}$ by $dQ_{N,t}=R_{N,t}\,dP_{N,\beta}$. Then
\begin{equation}\label{eq:allmodes-pullback}
 \mathbb E_PR_{N,t}=1,\qquad
 (T_{N,t})_*Q_{N,t}=P_{N,\beta},
\end{equation}
where the second formula denotes coordinatewise transport of configurations.

Let $A_N$ be the event that every particle has modulus at least $2r_N$.
Uniform local intensity bounds give $P(A_N^c)=O(r_N^2)$. The construction
also gives $T_{N,t}(B_{2r_N})\subset B_{3r_N}$ for all sufficiently large
$N$. Hence the exact pullback identity gives
\begin{equation}\label{eq:allmodes-hole-both-laws}
 Q_{N,t}(A_N^c)\le
 P\{\Xi_N(B_{3r_N})\ge1\}=O(r_N^2).
\end{equation}
In particular, expansions on $A_N$ need not be integrated over its
complement; that complement has vanishing probability under \emph{both}
measures.

We first verify all input bounds for the expansions under both $P$ and
$Q=Q_{N,t}$. For $p\ge1$, on any event where $\min_i|z_i|\ge r_N$,
\begin{equation}\label{eq:allmodes-negative-moment}
 |S_{-p,N}-\overline{S_{p,N}}|
 \le C_{p,r_0}\mathcal D_N+C_p r_N^{-p}\Xi_N(B_{r_0}),
 \qquad 0<r_0<1\text{ fixed}.
\end{equation}
Indeed, for $|z|\ge r_0$ one has
$|z^{-p}-\bar z^p|=|z|^{-p}-|z|^p\le C_{p,r_0}(1-|z|)$;
for the remaining points use $|z|^{-p}+|z|^p\le2r_N^{-p}$.
Under $P$, this proves
\[
 S_{-p,N}/N\to0\quad(1\le p\le2k),
\]
on $A_N$, since $\alpha p<1$. For $p<k$, induction further gives
\begin{equation}\label{eq:allmodes-negative-small}
 S_{-p,N}=O_P(\log(N+1)+r_N^{-p})\quad\hbox{on }A_N.
\end{equation}

Under $Q$, write $y_i=T_{N,t}(z_i)$; the configuration $y$ has law $P$.
Uniform convergence $T_{N,t}\to\mathrm{id}$ implies tightness of every
fixed compact-disk count for $z$ and also
$S_{j,N}(z)/N\to0$ for every fixed positive $j$. The maps and their
inverses are uniformly bilipschitz and map the unit circle onto itself.
Consequently
\[
 1-|z|\le C(1-|T_{N,t}(z)|),
\]
with a constant independent of large $N$. Therefore $\mathcal D_N(z)=O_Q(\log(N+1))$.
Formula \eqref{eq:allmodes-negative-moment} now proves
$S_{-p,N}(z)/N\to0$ on $A_N$ under $Q$ for $p\le2k$.

The lower positive modes are tight under $Q$ as well. To see the stronger
fact needed later, for every fixed positive $j$ use the uniform Taylor
expansion from Lemma~\ref{lem:transport-mode-shifts}, on $A_N$:
\begin{equation}\label{eq:allmodes-mode-increment}
 S_{j,N}(y)-S_{j,N}(z)
 =\frac jk\left(a_NS_{j-k,N}(z)-\bar a_NS_{j+k,N}(z)\right)
   +O_{j,k,t}(N|a_N|^2r_N^{-2k}).
\end{equation}
For $j<k$, the negative index has absolute value at most $k-1$;
the convergence $S_{-(k-j),N}(z)/N\to0$ under $Q$ was proved in the
preceding paragraph without assuming tightness of any lower positive
mode under $Q$. Both linear terms in
\eqref{eq:allmodes-mode-increment} are therefore $o_Q(1)$.
The deterministic error is
$O(N^{-1+2\alpha k})=o(1)$. It follows that
\begin{equation}\label{eq:allmodes-lower-invariant}
 S_{j,N}(y)-S_{j,N}(z)\to0\quad\hbox{in }Q\hbox{-probability},\qquad j<k.
\end{equation}
Thus the lower positive modes are tight under $Q$. Applying
\eqref{eq:allmodes-negative-moment} once more proves
\eqref{eq:allmodes-negative-small} under $Q$ too.

The deterministic likelihood expansion from the transport lemma now applies
under both measures. The linear terms other than
$-(\beta N/k)\operatorname{Re}(\bar a_NS_{k,N})$ contain only single
$S_{\pm k,N}$ or products $S_{p,N}S_{k-p,N}$ and
$S_{-p,N}S_{-(k-p),N}$ with $1\le p<k$. After multiplication by
$a_N=O(N^{-1})$, all are $o(1)$ in probability: for the negative products
use \eqref{eq:allmodes-negative-small} and $\alpha k<1$.
All pure quadratic terms vanish by macroscopic positive and negative moment
bounds. The mixed quadratic term has limit
$-\beta N^2|a_N|^2/(2k)=-2k|t|^2/\beta$.
The total deterministic remainder is bounded by
\[
 C_{k,t,\beta}\left(N^2|a_N|^3r_N^{-3k}
                         +N|a_N|^2r_N^{-2k}\right)=o(1).
\]
Together with \eqref{eq:allmodes-hole-both-laws}, this proves
\begin{equation}\label{eq:allmodes-loglikelihood}
 \log R_{N,t}
 =-2\operatorname{Re}(\bar tS_{k,N})-\frac{2k}{\beta}|t|^2+o(1)
\end{equation}
in probability under \emph{each} of $P$ and $Q_{N,t}$.
No assertion of integrability of the error is used.

This identity first gives tightness of $S_{k,N}$ under $P$. For example,
use $t=-1$, let $c=2k/\beta$, and denote the probabilistic error by $E_N$.
Markov's inequality and $\mathbb E_PR_{N,-1}=1$ imply
\[
 P\{\operatorname{Re}S_{k,N}>L\}
 \le P\{|E_N|>1\}+e^{-2L+c+1}.
\]
The choices $t=1,i,-i$ control the other three tails. Hence $S_{k,N}$
is tight. It is tight under $Q_{N,t}$ too: on $A_N$ the exact identity
$f_a(z)^k=(z^k+a)/(1+\bar az^k)$ gives
\[
 S_{k,N}(y)-S_{k,N}(z)
 =a_NN-\bar a_NS_{2k,N}(z)+O_t(N|a_N|^2)
 =\frac{2kt}{\beta}+o_Q(1).
\]
The $y$ mode has the just-proved baseline tightness, proving the claim.

Now \eqref{eq:allmodes-loglikelihood} under $Q$ shows that $\log R_{N,t}$
is tight under $Q$. This proves uniform integrability of the exact
likelihood ratios, since
\begin{equation}\label{eq:allmodes-ui}
 \mathbb E_P[R_{N,t}\mathbf1_{\{R_{N,t}>K\}}]
 =Q_{N,t}\{\log R_{N,t}>\log K\}\longrightarrow0
\end{equation}
as $K\to\infty$, uniformly in all sufficiently large $N$. The finitely
many excluded $N$ do not affect uniform integrability. This argument
is why the small patch near the origin causes no missing tail estimate.

Let $Y_N=(\Xi_N,S_{1,N},\ldots,S_{k-1,N})$, and write $Y_N^T$ for the same
observables after transport. Under $Q$, the vague distance between the two
configurations tends to zero: the maps tend uniformly to the identity,
and compact-disk counts are tight under $Q$. Equation
\eqref{eq:allmodes-lower-invariant} deals with the remaining coordinates.
Thus $Y_N^T-Y_N\to0$ in the product-space metric in $Q$-probability.
Tightness under $Q$ implies, for every $F\in C_b$ of this product space,
\[
 \mathbb E_Q[F(Y_N^T)-F(Y_N)]\longrightarrow0.
\]
The exact pullback identity consequently gives
\begin{equation}\label{eq:allmodes-test-identity}
 \mathbb E_PF(Y_N)-\mathbb E_P[R_{N,t}F(Y_N)]\longrightarrow0.
\end{equation}

Take any jointly convergent subsequence of $(Y_N,S_{k,N})$, with limit
$(Y,G_k)$. Such subsequences exist by tightness. By
\eqref{eq:allmodes-loglikelihood} under $P$ and the continuous mapping
theorem,
\[
 (Y_N,S_{k,N},R_{N,t})\Rightarrow
 \left(Y,G_k,\exp\left[-2\operatorname{Re}(\bar tG_k)
                                  -\frac{2k}{\beta}|t|^2\right]\right).
\]
Uniform integrability \eqref{eq:allmodes-ui} permits passage to the
expectation in \eqref{eq:allmodes-test-identity}. We obtain
\begin{equation}\label{eq:allmodes-limit-transform}
 \mathbb E\left[F(Y)e^{-2\operatorname{Re}(\bar tG_k)}\right]
 =e^{2k|t|^2/\beta}\mathbb EF(Y),\qquad t\in\mathbb C.
\end{equation}
In particular, the limiting exponential moments are finite. We have not
assumed, and do not need, uniform integrability of the corresponding
\emph{unmodified} prelimit exponential statistic on the small-hole event.
The later Theorem~\ref{thm:harmonic-exponential-moments} proves the
stronger prelimit exponential-moment bounds by a separate fixed-radius
argument.

Taking $F=1$, uniqueness of the moment-generating function in $\mathbb R^2$
identifies $G_k$ as the centered circular complex Gaussian of variance
$2k/\beta$. For nonnegative $F$ with positive expectation, divide
\eqref{eq:allmodes-limit-transform} by that expectation and apply the same
uniqueness theorem to the $F$-weighted probability law of $G_k$.
It equals the unweighted Gaussian law. Bounded continuous functions
therefore show that $G_k$ is independent of $Y$.
The induction hypothesis identifies the lower coordinates of $Y$ as
independent Gaussians independent of the limiting interior process.
This completes the induction and proves the stated convergences.
\end{proof}

\begin{corollary}[Harmonic polynomials and the countable mode vector]
If $h(z)=\operatorname{Re}\sum_{j=1}^m b_j z^j$, then
\[
 \sum_{i=1}^N h(z_i)
 \Longrightarrow\mathcal N\left(0,\frac1\beta\sum_{j=1}^m j|b_j|^2\right).
\]
Its limiting variance equals
$(\pi\beta)^{-1}\int_{\D}|\nabla h|^2\,dA$.
The entire countable vector $(S_{j,N})_{j\ge1}$ converges in the product
topology of $\C^{\N}$ to $(G_j)_{j\ge1}$. Along every subsequence on
which $\Xi_N$ converges, this convergence is joint with $\Xi_N$, and the
limiting countable Gaussian vector is independent of the interior limit.
\end{corollary}
\begin{proof}
The first statement is a real linear image of the Gaussian vector in
Theorem~\ref{thm:all-harmonic-modes}. Circular symmetry gives
$\Var\operatorname{Re}(b_jG_j)=j|b_j|^2/\beta$, and independence makes
the variances add. Polar integration and orthogonality give
$\int_{\D}|\nabla h|^2dA=\pi\sum_j j|b_j|^2$.
For countable product convergence, each coordinate family is tight.
Given $\varepsilon>0$, choose finite bounds $L_j$ such that
$\sup_N\Pbb(|S_{j,N}|>L_j)\le\varepsilon2^{-j}$.
The compact product $\prod_j\{|z|\le L_j\}$ then has probability at
least $1-\varepsilon$ for every $N$. This proves tightness in the
countable product, also jointly with any tight family of interior
configurations. Finite-dimensional distributions, already identified
by the theorem, determine every possible subsequential limit in that
product. They also prove independence from the interior configuration.
\end{proof}


% Section fragment for the disk-log-gas manuscript.
% Uses its existing theorem environments and macros.
\section{Uniform exponential moments of every fixed harmonic mode}
\label{sec:harmonic-exponential-moments}

Write $\mathcal D_N=\sum_{i=1}^N(1-|z_i|)$ and
$S_{j,N}=\sum_{i=1}^N z_i^j$ for $j\geq1$.

\begin{lemma}[A quadratic exponential bound for the radial deficit]
\label{lem:deficit-quadratic-exponential}
For every $c>0$ and $\beta>0$,
\[
 \sup_{N\geq1}\E_N\exp\{c\mathcal D_N^2/N\}<\infty.
\]
\end{lemma}
\begin{proof}
Use the established simultaneous radial domination with independent
exponentials $E_j$ of rates $a_j=2j+\beta j(j-1)/2$. It gives a coupling
under which the ordered depths satisfy
$-\log R_k\leq\sum_{j=k}^N E_j$. Fix an integer $J\geq1$. Since
$1-s\leq-\log s$ for $0<s\leq1$, this coupling gives
\[
 \mathcal D_N\leq J+Y_N,
 \qquad Y_N=\sum_{j=J+1}^N(j-J)E_j,
\]
where an empty sum is zero. Put $\mu_N=\E Y_N$. The elementary estimate
$a_j\geq\beta j(j-1)/2$ shows
$\mu_N\leq C_\beta\log(N+1)$.

For each fixed $\lambda>0$, choose $J$ so large that
$\lambda(j-J)/a_j\leq1/2$ for every $j>J$; this is possible because
$a_j\geq\beta j(j-1)/2$ and
$(j-J)/(j(j-1))\leq1/(j-1)\leq1/J$. Independence and
$-\log(1-u)-u\leq u^2$ for $0\leq u\leq1/2$ give
\[
 \E e^{\lambda(Y_N-\mu_N)}
 =\prod_{j=J+1}^N
    \frac{\exp[-\lambda(j-J)/a_j]}
         {1-\lambda(j-J)/a_j}
 \leq\exp\!\left(\lambda^2
        \sum_{j=J+1}^\infty\frac{(j-J)^2}{a_j^2}\right)<\infty,
\]
uniformly in $N$. The displayed infinite series converges by comparison
with $\sum j^{-2}$.

Take $\lambda=2c$ and set $b_N=J+\mu_N$ and
$X_N=(\mathcal D_N-b_N)_+$. On the coupling,
$X_N\leq(Y_N-\mu_N)_+$; also $X_N\leq\mathcal D_N\leq N$.
Consequently
\[
 \frac{\mathcal D_N^2}{N}
 \leq\frac{2b_N^2}{N}+\frac{2X_N^2}{N}
 \leq\frac{2b_N^2}{N}+2X_N.
\]
The first term is bounded uniformly in $N$, while
$\E e^{2cX_N}\leq1+\E e^{2c(Y_N-\mu_N)}$ is uniformly bounded by the
preceding product estimate. This proves the claim.
\end{proof}

\begin{lemma}[A global bound for a fixed-radius transport]
\label{lem:fixed-radius-loglikelihood-bound}
Fix $k\geq1$, $t\in\C$, and $0<r<1/4$. For sufficiently large $N$,
put $a=2kt/(\beta N)$ and let $T=T_{a,r}$ be the map of
Lemma~\ref{lem:patched-lift}. Let $R$ denote its exact change-of-variables
likelihood under $\Pbb_{N,\beta}$. Put
$M=\#\{i:|z_i|<2r\}$. There is a constant $C$, independent of $N$
and the configuration of distinct points, such that
\begin{equation}\label{eq:fixed-radius-global-error}
 \left|\log R+2\operatorname{Re}(\bar t S_{k,N})\right|
 \leq C\left[1+M+
   \frac{\mathcal D_N^2+M^2+
                     \sum_{p=1}^{k-1}|S_{p,N}|^2}{N}\right].
\end{equation}
\end{lemma}
\begin{proof}
All constants below may depend on $k,t,r,\beta$. The derivative estimate
for the patched map is $\|DT-I\|_\infty\leq C/N$. Convexity of the disk
therefore implies
\[
 \left|\log\left|\frac{T(z)-T(w)}{z-w}\right|\right|\leq C/N,
 \qquad |\log\det DT(z)|\leq C/N
\]
for distinct $z,w$ and sufficiently large $N$. There are at most $NM$
pairs meeting $\{|z|<2r\}$. Their total contribution to $\log R$ is
bounded in absolute value by $CM$, and the entire Jacobian contribution
is bounded by $C$.

Let $I=\{i:|z_i|\geq2r\}$, $n=|I|=N-M$, and
$U_j=\sum_{i\in I}z_i^j$ for integers $j$. For pairs with both indices
in $I$, the map agrees with $f_a$. The first-order part of
Proposition~\ref{prop:aggregate-coefficients}, with $n$ in place of $N$,
gives
\begin{align*}
 \sum_{\substack{i<j\\i,j\in I}}
   \log\frac{f_a(z_i)-f_a(z_j)}{z_i-z_j}
 ={}&\frac{a}{2k}
       \left[(k-1)U_{-k}-\sum_{p=1}^{k-1}U_{-p}U_{-(k-p)}\right]\\
 &-\frac{\bar a}{2k}
       \left[(2n-k-1)U_k+\sum_{p=1}^{k-1}U_pU_{k-p}\right]+O(1).
\end{align*}
The remainder is uniformly $O(n^2|a|^2)=O(1)$ because $r$ is fixed;
this also follows by summing the second-order error in the individual
divided-difference expansion. The identity is valid when $n=0$ or $1$
as well, with the corresponding empty pair sum.

For every fixed $1\leq p\leq k$, the mean-value theorem on $[2r,1]$
gives
\[
 |U_{-p}-\overline{U_p}|
 \leq C\sum_{i\in I}(1-|z_i|)\leq C\mathcal D_N,
 \qquad |U_p-S_{p,N}|\leq M.
\]
Thus, for $p<k$, each positive or negative product in the preceding
display is bounded by
$C[\mathcal D_N^2+M^2+\sum_{q=1}^{k-1}|S_{q,N}|^2]$.
The isolated $U_{-k}$ term, multiplied by $a$, is bounded because
$|U_{-k}|\leq N(2r)^{-k}$. Finally,
\[
 \left|\frac{2n-k-1}{N}U_k-2S_{k,N}\right|
 \leq C(1+M),
\]
using $|S_{k,N}|\leq N$, $|U_k-S_{k,N}|\leq M$, and $M\leq N$.
Taking real parts, multiplying the pair expression by $\beta$, and
using $a=2kt/(\beta N)$ proves
\eqref{eq:fixed-radius-global-error}.
\end{proof}

\begin{theorem}[Genuine exponential-moment bounds and convergence]
\label{thm:harmonic-exponential-moments}
For every $\beta>0$, integer $m\geq1$, and $c>0$,
\begin{equation}\label{eq:harmonic-exponential-bound}
 \sup_{N\geq1}\E_N\exp\!\left(c\sum_{j=1}^m|S_{j,N}|\right)<\infty.
\end{equation}
Consequently all mixed moments of every fixed finite vector of harmonic
modes converge to those of the independent Gaussian vector identified in
Theorem~\ref{thm:all-harmonic-modes}. Along every interior-convergent
subsequence, the Gaussian independence identity extends from bounded
tests in the modes to every continuous mode test of at most exponential
growth, multiplied by a bounded continuous function of the interior.
\end{theorem}
\begin{proof}
We prove the moment bound by induction on $k$. The base $k=1$ follows
from the already proved M\"obius exponential-moment estimate; alternatively
the following argument works for $k=1$ with no lower-mode terms.
Fix $t\in\C$, $0<r<1/4$, and use the preceding map and likelihood for
all sufficiently large $N$. Define $Q$ by $dQ=R\,dP$, where
$P=\Pbb_{N,\beta}$. Under $Q$, the variables $y_i=T(z_i)$ have law $P$.
The deterministic transport bounds give, for each fixed $p\geq1$,
\[
 |S_{p,N}(z)-S_{p,N}(y)|\leq C_p,
 \qquad \mathcal D_N(z)\leq2\mathcal D_N(y),
 \qquad M(z)\leq\#\{i:|y_i|<3r\}.
\]
The first bound uses $\|T-\mathrm{id}\|_\infty=O(N^{-1})$ and
$|z^p-y^p|\leq p|z-y|$ on the closed disk; the others follow from
Lemma~\ref{lem:patched-lift}.

Every constant multiple of the right-hand side of
\eqref{eq:fixed-radius-global-error} has a uniformly bounded exponential
moment under $Q$. Here are the individual verifications. The compact
count has all exponential moments uniformly in $N$ by the established
Gaussian count tail. Since $M\leq N$, $M^2/N\leq M$. The deficit term
is controlled by Lemma~\ref{lem:deficit-quadratic-exponential} and the
transport bound. For $p<k$, $|S_{p,N}(y)|\leq N$ yields
$|S_{p,N}(y)|^2/N\leq|S_{p,N}(y)|$, so the induction hypothesis and
the bounded transport shift control each lower-mode term. H\"older's
inequality combines these finitely many estimates, whose coefficients
may be arbitrarily large fixed constants.

The exact change of variables now gives
\begin{align*}
 \E_P e^{-2\operatorname{Re}(\bar t S_{k,N})}
 &=\E_Q\exp\{-\log R-2\operatorname{Re}(\bar tS_{k,N})\}\\
 &\leq \E_Q\exp\!\left\{C\left[1+M+
   \frac{\mathcal D_N^2+M^2+
                     \sum_{p=1}^{k-1}|S_{p,N}|^2}{N}\right]\right\}
 \leq C_{k,t,r,\beta}.
\end{align*}
The finitely many small values of $N$ are harmless because $|S_{k,N}|\leq N$.
Using $t$ in the four real and imaginary directions, and
$|w|\leq|\operatorname{Re}w|+|\operatorname{Im}w|$, proves uniform
exponential moments of $|S_{k,N}|$. A further application of
H\"older's inequality proves \eqref{eq:harmonic-exponential-bound}.

The established Gaussian convergence and
\eqref{eq:harmonic-exponential-bound} imply uniform integrability of
every continuous test of at most exponential growth: use the same bound
with a larger coefficient. Truncation and weak convergence prove the
moment and independence assertions.
\end{proof}

\section{Analytic harmonic statistics: tightness and Laplace transforms}
\label{sec:analytic-harmonic-statistics}

This section upgrades the finite-mode central limit theorem to analytic
harmonic statistics and to convergence of their exponential moments.
The argument uses a smooth rational vector field. Its singularity is
removed on a fixed compact disk; the resulting error is controlled by
that disk's particle count.

Let $f$ be holomorphic on a neighborhood of $\overline{\D}$, with
$f(0)=0$. Define
\[
 Y_N(f)=2\operatorname{Re}\sum_{i=1}^N f(z_i),\qquad
 Z_N(f)=\mathbb E_{N,\beta}e^{Y_N(f)},\qquad
 dP_N^f=Z_N(f)^{-1}e^{Y_N(f)}dP_{N,\beta}.
\]
Here $Z_N(f)$ is a normalized moment-generating function, not the
unweighted particle partition function. Also put
\[
 S_{k,N}=\sum_i z_i^k,\qquad
 X_N(u)=\sum_i\frac{u z_i}{1-u z_i}
       =\sum_{k\geq1}u^kS_{k,N},\qquad |u|<1,
\]
and $D_N=\sum_i(1-|z_i|)$. Throughout this section,
$B_R=\{z:|z|<R\}$.

\begin{lemma}[Uniform radial estimates under holomorphic tilts]
\label{lem:analytic-tilt-radial}
The radial stochastic comparison for the unweighted disk gas holds,
with the same comparison rates
\[
 a_j=2j+\frac\beta2j(j-1),
\]
under every law $P_N^f$. Consequently, for every $R<1$ there are
constants $C_R,c_R>0$, depending on $\beta,R$ but not on $N$ or $f$,
such that
\begin{equation}\label{eq:analytic-tilt-count}
 P_N^f\{\Xi_N(B_R)\geq m\}\leq C_R e^{-c_Rm^2},\qquad m\geq0.
\end{equation}
Moreover, $\mathbb E_N^fD_N=O_\beta(\log(N+1))$. For each fixed $f$
and $s<1$,
\begin{equation}\label{eq:analytic-tilt-macro}
 \frac1N\sup_{|u|\leq s}|X_N(u)|\longrightarrow0
 \quad\hbox{in }P_N^f\hbox{-probability}.
\end{equation}
\end{lemma}

\begin{proof}
We explain why the previously proved radial comparison applies without
change. In the ordered-radius proof, fix all angular variables except a
common rotation of an inner block, and scale that block by a complex
parameter $q$. The factors involving that block have the form
\[
 |P(q)|^\beta
 \exp\left(2\operatorname{Re}\sum_{i\text{ in the block}}f(qw_i)\right)
 =\left|P(q)\exp\left(\frac2\beta
                    \sum_{i\text{ in the block}}f(qw_i)\right)\right|^\beta,
\]
where $P$ is the same holomorphic cross-interaction polynomial as in the
unweighted proof. The expression is subharmonic in $q$. Its circular
mean is therefore nondecreasing in $|q|$. The radial Jacobian and the
homogeneous Vandermonde exponents are unchanged. Thus the same product
association argument gives exactly the same stochastic domination of
radial gaps by independent exponentials of rates $a_j$. This proves
\eqref{eq:analytic-tilt-count} by the established count-tail estimate.
It also gives
\[
 \mathbb E_N^fD_N\leq
 \mathbb E_N^f\sum_i(-\log|z_i|)
 \leq\sum_{j=1}^N\frac{j}{a_j}=O_\beta(\log(N+1)).
\]

Rotation invariance of the unweighted gas and $f(0)=0$ give
$\mathbb E_{N,\beta}Y_N(f)=0$. Jensen's inequality therefore gives
$Z_N(f)\geq1$. Hence, for any event $A$,
\[
 P_N^f(A)\leq e^{2N\|f\|_{\infty,\overline\D}}P_{N,\beta}(A).
\]
The already proved macroscopic Fourier concentration bound, whose
logarithm is $-c_{k,\varepsilon}N^2+O_\beta(N\log N)$ for the event
$\{|S_{k,N}|>\varepsilon N\}$, implies $S_{k,N}/N\to0$ under $P_N^f$
for each fixed $k$. Finally,
\[
 \frac1N\sup_{|u|\leq s}|X_N(u)|
 \leq\sum_{k=1}^K s^k\frac{|S_{k,N}|}{N}
          +\frac{s^{K+1}}{1-s}.
\]
First choose $K$ and then let $N\to\infty$. This proves
\eqref{eq:analytic-tilt-macro}.
\end{proof}

\begin{lemma}[A rational transport with a compact-count error]
\label{lem:rational-resolvent-transport}
Fix $0<s<R_0<R_1<1$ and a fixed holomorphic tilt $f$ as above.
For every $u$ with $|u|\leq s$ and every $t\in\{1,-1,i,-i\}$ there is a
smooth diffeomorphism $T_{N,u,t}:\overline\D\to\overline\D$.
Its likelihood ratio $R_{N,u,t}^f$ under change of variables in $P_N^f$
satisfies
\[
 \mathbb E_N^f R_{N,u,t}^f=1
\]
and the deterministic bound
\begin{equation}\label{eq:rational-likelihood-error}
 \left|\log R_{N,u,t}^f
             +\beta\operatorname{Re}(\bar t X_N(u))\right|
 \leq C\left(1+M_N+
        \frac{|X_N(u)|^2+D_N^2+M_N^2}{N}\right),
\end{equation}
where $M_N=\Xi_N(B_{R_1})$ and $C$ does not depend on $N,u,t$ or the
particle configuration.
The assertion is needed only off the collision diagonals, which have
probability zero.
\end{lemma}

\begin{proof}
For $u$ fixed write
\[
 x_u(z)=\frac{uz}{1-uz},\qquad
 y_u(z)=\frac{\bar u}{z-\bar u},\qquad
 A_u(z)=z y_u(z),\qquad B_u(z)=z x_u(z).
\]
Choose a smooth real radial cutoff $\chi$, equal to zero on $B_{R_0}$
and equal to one on $\overline\D\setminus B_{R_1}$. Define
\[
 v_{u,t}(z)=\chi(z)\{tA_u(z)-\bar t B_u(z)\},
\]
with the zero definition throughout $B_{R_0}$. This removes the pole of
$A_u$, since $|u|\leq s<R_0$. The vector fields and their derivatives
of any fixed finite order are uniformly bounded in $u,t$. On $|z|=1$
one has $y_u(z)=\overline{x_u(z)}$ and therefore
\[
 \operatorname{Re}(\bar z v_{u,t}(z))
 =\operatorname{Re}\{t\overline{x_u(z)}-\bar t x_u(z)\}=0.
\]
Their flows preserve the disk and its boundary. Let $T_{N,u,t}$ be the
flow at time $1/N$.

We record explicitly the uniform expansion used for this flow. If $v$
is one of these vector fields and $T_h$ is its flow, then
\[
 \log\frac{|T_h(z)-T_h(w)|}{|z-w|}
 =h\operatorname{Re}\frac{v(z)-v(w)}{z-w}+O(h^2),
 \qquad z\ne w,
\]
uniformly on $\overline\D^2$. To see this, the flow equation gives
$T_h=\mathrm{id}+hv+O(h^2)$ in the $C^1$ norm. The $O(h^2)$ remainder
is Lipschitz with constant $O(h^2)$, so its divided difference has the
same bound; expanding $\log|1+q|$ proves the formula. Also
$\log J_{T_h}(z)=O(h)$ uniformly. Since $f$ has bounded first derivative,
$|Y_N(f)(T_hz)-Y_N(f)(z)|=O(Nh)$. The exact density and Jacobian change
of variables consequently give, at $h=1/N$,
\begin{equation}\label{eq:rational-flow-first-order}
 \log R_{N,u,t}^f
 =\frac\beta N\operatorname{Re}
        \sum_{i<j}\frac{v_{u,t}(z_i)-v_{u,t}(z_j)}{z_i-z_j}+O(1).
\end{equation}
The total pairwise remainder is $O(N^2h^2)=O(1)$, so this estimate is
valid uniformly even for very close particles. Exact change of variables
gives $\mathbb E_N^f R_{N,u,t}^f=1$.

Let $J=\{i:|z_i|\geq R_1\}$ and $n=|J|=N-M_N$. Since $v_{u,t}$ is
uniformly Lipschitz, deleting all pairs with at least one endpoint outside
$J$ changes the sum in \eqref{eq:rational-flow-first-order}, after division
by $N$, by at most $CM_N$. For points indexed by $J$, the exact identities
\[
 \frac{A_u(z)-A_u(w)}{z-w}=-y_u(z)y_u(w),\qquad
 \frac{B_u(z)-B_u(w)}{z-w}
       =x_u(z)+x_u(w)+x_u(z)x_u(w)
\]
apply. With $X_J=\sum_{i\in J}x_u(z_i)$ and
$Y_J=\sum_{i\in J}y_u(z_i)$, the remaining sum is
\begin{align*}
 &-\frac t2\left(Y_J^2-\sum_{i\in J}y_u(z_i)^2\right)\\
 &\qquad-\bar t\left((n-1)X_J+
             \frac12\left(X_J^2-\sum_{i\in J}x_u(z_i)^2\right)\right).
\end{align*}
All single-particle $x_u$ values, and all $y_u$ values on $J$, are bounded
by fixed constants. Thus the diagonal terms contribute $O(1)$ after
division by $N$. Replacing $(n-1)X_J/N$ by $X_N(u)$ costs at most
$C(M_N+1)$. Finally, on $|z|\geq R_1$,
\[
 y_u(z)-\overline{x_u(z)}
 =\frac{\bar u(1-|z|^2)}{(z-\bar u)(1-\bar u\bar z)},
 \qquad
 |y_u(z)-\overline{x_u(z)}|\leq C(1-|z|).
\]
It follows that $|Y_J-\overline{X_J}|\leq CD_N$, and that
\[
 |Y_J|^2+|X_J|^2
 \leq C\{|X_N(u)|^2+D_N^2+M_N^2\}.
\]
Substitution proves \eqref{eq:rational-likelihood-error}.
\end{proof}

\begin{proposition}[Tightness of the analytic resolvent under all fixed tilts]
\label{prop:analytic-resolvent-tightness}
For every fixed $f$ holomorphic near $\overline\D$, the laws of $X_N$
under $P_N^f$ are tight in the space $\mathcal H(\D)$ of holomorphic
functions, with the topology of uniform convergence on compact subsets.
\end{proposition}

\begin{proof}
Fix $s<R_0<R_1<1$ as in the preceding lemma. Let $C$ be its error
constant. Choose positive $\delta,\varepsilon$ sufficiently small that
\[
 C\delta+C\varepsilon+C\delta^2\varepsilon
       <\frac{\beta}{2\sqrt2}.
\]
Define the single event, common to every parameter $|u|\leq s$,
\[
 G_N=\left\{D_N\leq N^{1/3},\quad
           \sup_{|u|\leq s}|X_N(u)|\leq\varepsilon N\right\}.
\]
Lemma~\ref{lem:analytic-tilt-radial} and Markov's inequality imply
$P_N^f(G_N)\to1$.

For every complex number $x$ one of $t\in\{1,-1,i,-i\}$ satisfies
$-\operatorname{Re}(\bar t x)\geq|x|/\sqrt2$. On
$G_N\cap\{M_N\leq\delta|X_N(u)|\}$, the error in
\eqref{eq:rational-likelihood-error} is at most
\[
 C_0+(C\delta+C\varepsilon+C\delta^2\varepsilon)|X_N(u)|.
\]
Here $D_N^2/N\leq1$, $|X_N(u)|^2/N\leq\varepsilon|X_N(u)|$, and
$M_N^2/N\leq\delta^2\varepsilon|X_N(u)|$. For the indicated choice of
$t$ we consequently have
\[
 \log R_{N,u,t}^f\geq c|X_N(u)|-C_0,
 \qquad c=\frac{\beta}{2\sqrt2}>0.
\]
By Markov's inequality, $\mathbb E_N^fR_{N,u,t}^f=1$, a union bound over
the four choices of $t$, and \eqref{eq:analytic-tilt-count},
\begin{equation}\label{eq:analytic-resolvent-good-tail}
 P_N^f\bigl(G_N\cap\{|X_N(u)|>L\}\bigr)
 \leq 4e^{C_0-cL}+C_1e^{-c_1L^2},\qquad L\geq0,
\end{equation}
with constants uniform in $N$ and $|u|\leq s$. In the second term we
used that $M_N>\delta|X_N(u)|$ and $|X_N(u)|>L$ imply
$M_N>\delta L$. Integrating the tail estimate gives
\begin{equation}\label{eq:analytic-resolvent-truncated-mean}
 \sup_{N,|u|\leq s}\mathbb E_N^f[\mathbf1_{G_N}|X_N(u)|]<\infty.
\end{equation}

For $0<r<s$, Cauchy's integral formula gives, configuration by
configuration,
\[
 \sup_{|u|\leq r}|X_N(u)|
 \leq\frac{s}{s-r}\frac1{2\pi}
        \int_0^{2\pi}|X_N(se^{i\theta})|\,d\theta.
\]
Multiplying by the indicator of the \emph{same} event $G_N$, taking
expectations, and applying Fubini's theorem and
\eqref{eq:analytic-resolvent-truncated-mean} bounds the resulting
expectation uniformly in $N$. Since $P_N^f(G_N^c)\to0$, the local suprema
are tight. This holds for every $r<1$. On an increasing sequence of such
disks, choose bounds with summable exceptional probabilities. Montel's
theorem then supplies compact subsets of $\mathcal H(\D)$ carrying
arbitrarily high probability, proving tightness.
\end{proof}

\begin{lemma}[Tight exponential families give uniform integrability]
\label{lem:tight-exponential-family}
Let $W_N$ be real random variables, with $\mathbb Ee^{tW_N}<\infty$ for
all real $t$. Suppose that for every fixed real $t$, the variables $W_N$
are tight under the tilted laws
\[
 dQ_N^t=\frac{e^{tW_N}}{\mathbb Ee^{tW_N}}\,dP_N.
\]
Then $\sup_N\mathbb Ee^{tW_N}<\infty$ for every $t$, and
$(e^{tW_N})_N$ is uniformly integrable for every $t$. If also
$W_N\Rightarrow W$, then
$\mathbb Ee^{tW_N}\to\mathbb Ee^{tW}$ for every real $t$.
\end{lemma}

\begin{proof}
For $t$ fixed, tightness under $Q_N^t$ provides $A<\infty$ with
$Q_N^t(|W_N|\leq A)\geq1/2$ for all $N$. Consequently
\[
 \frac12\mathbb Ee^{tW_N}
 \leq\mathbb E[e^{tW_N}\mathbf1_{\{|W_N|\leq A\}}]
 \leq e^{|t|A}.
\]
Applying the same bound at $pt$ for any $p>1$ proves a uniform $L^p$
bound for $e^{tW_N}$, hence uniform integrability. Convergence in law and
uniform integrability give the last assertion, for instance by bounded
continuous truncation. These bounds also show that all the asserted
limiting exponential moments are finite.
\end{proof}

\begin{corollary}[Uniform analytic insertion bounds without a central limit theorem]
\label{cor:analytic-insertion-without-clt}
For every $f$ holomorphic near $\overline\D$ with $f(0)=0$, and every
real $t$, one has
\[
 \sup_N\mathbb E_{N,\beta}e^{tY_N(f)}<\infty.
\]
In particular, for every real $q$ and $R<1$,
\begin{equation}\label{eq:analytic-insertion-without-clt}
 \sup_N\sup_{|a|\le R}
 \mathbb E_{N,\beta}\left|\prod_{i=1}^N(1-a z_i)\right|^q<\infty,
 \qquad
 \sup_N\sup_{|z|\le R}
 \mathbb E_{N,\beta}\prod_{i=1}^N|z-z_i|^\beta<\infty.
\end{equation}
These bounds use no central limit theorem or identification of any
limiting distribution.
\end{corollary}

\begin{proof}
Write $f(z)=\sum_{k\ge1}f_kz^k$ and choose $a>1$ within its disk of
holomorphy. Cauchy's estimate gives $|f_k|\le C a^{-k}$, after decreasing
$a$ if necessary. Choose $a^{-1}<r<1$. For $h(u)=\sum h_ku^k$ in
$\mathcal H(\D)$, another Cauchy estimate gives
\[
 \left|\sum_{k\ge1}f_kh_k\right|
 \le C\sup_{|u|\le r}|h(u)|\sum_{k\ge1}(ar)^{-k}.
\]
Thus $h\mapsto2\operatorname{Re}\sum f_kh_k$ is a continuous real
linear functional on $\mathcal H(\D)$, and it takes $X_N$ to $Y_N(f)$.
Proposition~\ref{prop:analytic-resolvent-tightness} applied to every
tilt $tf$ proves that $Y_N(f)$ is tight under its own exponential
tilt by $e^{tY_N(f)}$. Lemma~\ref{lem:tight-exponential-family} proves
the first assertion.

Apply it to $f(z)=(q/2)\log(1-a z)$ for each fixed $|a|<1$, with the
holomorphic logarithm vanishing at zero. This proves the harmonic-product
bound for each fixed parameter. As a function of the parameter $a$, the
integrand is the squared modulus of a holomorphic function, for every
real $q$. Its circular means are therefore nondecreasing. Rotation
invariance of the gas makes its expectation radial, so the supremum on
$|a|\le R$ is its value at $a=R$, already bounded uniformly in $N$.
Finally, $|z-w|\le|1-\bar z w|$ on the disk gives the insertion bound.
\end{proof}

\begin{theorem}[Analytic Gaussian fluctuations and harmonic Laplace limits]
\label{thm:analytic-harmonic-laplace}
Let $G_1,G_2,\ldots$ be independent centered circular complex Gaussians
with $\mathbb E|G_k|^2=2k/\beta$. Then
\[
 X_N\Rightarrow X,\qquad X(u)=\sum_{k\geq1}G_ku^k,
\]
in $\mathcal H(\D)$. Along every subsequence for which
$\Xi_{N_j}\Rightarrow\Xi$ vaguely in the open disk,
\[
 (\Xi_{N_j},X_{N_j})\Rightarrow(\Xi,X),
\]
where $X$ and $\Xi$ are independent.

For $f(z)=\sum_{k\geq1}f_kz^k$ holomorphic near $\overline\D$,
\begin{equation}\label{eq:analytic-strong-szego}
 \lim_{N\to\infty}\mathbb E_{N,\beta}
           \exp\left(2\operatorname{Re}\sum_i f(z_i)\right)
 =\exp\left(\frac2\beta\sum_{k\geq1}k|f_k|^2\right).
\end{equation}
Moreover, for every bounded continuous functional $F$ on the space of
locally finite counting measures and every subsequence as above,
\begin{equation}\label{eq:analytic-mixed-laplace}
 \lim_{j\to\infty}\mathbb E_{N_j,\beta}
       [F(\Xi_{N_j})e^{Y_{N_j}(f)}]
 =\mathbb EF(\Xi)\,
       \exp\left(\frac2\beta\sum_{k\geq1}k|f_k|^2\right).
\end{equation}
\end{theorem}

\begin{proof}
Proposition~\ref{prop:analytic-resolvent-tightness}, with $f=0$, proves
tightness in $\mathcal H(\D)$. Taylor coefficient evaluation is continuous
in that space. The already proved unweighted finite-mode theorem
therefore identifies every subsequential limit with the indicated
Gaussian series. That series converges locally uniformly almost surely:
for each $r<1$,
\[
 \mathbb E\sum_{k\geq1}|G_k|r^k
 \leq\sum_{k\geq1}(2k/\beta)^{1/2}r^k<\infty.
\]
The finite-mode theorem also identifies joint coefficients with the
interior limit and gives independence for every finite coefficient
vector. Taylor coefficients generate the Borel sigma field of
$\mathcal H(\D)$ (they determine each function, and their partial sums
recover it locally). Hence the full Gaussian analytic function is
independent of $\Xi$, and the joint convergence follows.

Choose $R>1$ such that $f$ is holomorphic on $\{|z|<R\}$, and choose
$R^{-1}<r<1$. If $h(u)=\sum_{k\geq0}h_ku^k$ is holomorphic in $\D$,
then
\[
 L_f(h)=\sum_{k\geq1}f_kh_k
\]
is a continuous linear functional on $\mathcal H(\D)$. Indeed, choose
$R'$ with $1/r<R'<R$. Cauchy's estimates give
$|f_k|\leq C(R')^{-k}$ and
$|h_k|\leq r^{-k}\sup_{|u|\leq r}|h(u)|$, so the series is bounded by a
convergent geometric series times that local supremum. In particular,
$Y_N(f)=2\operatorname{Re}L_f(X_N)$. The analytic convergence proves
\[
 Y_N(f)\Rightarrow Y(f)=2\operatorname{Re}\sum_{k\geq1}f_kG_k,
 \qquad
 \operatorname{Var}Y(f)=\frac4\beta\sum_{k\geq1}k|f_k|^2,
\]
with joint independence from the subsequential interior limit.

For each real $t$, the exponential tilt by $e^{tY_N(f)}$ is exactly
$P_N^{tf}$. Proposition~\ref{prop:analytic-resolvent-tightness} and
continuity of $L_f$ show that $Y_N(f)$ is tight under this tilted law.
Lemma~\ref{lem:tight-exponential-family} consequently proves uniform
integrability of $e^{Y_N(f)}$, as well as of every fixed positive power
of it. Passing to the limit gives \eqref{eq:analytic-strong-szego}, using
the elementary real Gaussian moment-generating function. Joint
convergence, boundedness of $F$, and the same uniform integrability give
\eqref{eq:analytic-mixed-laplace}.
\end{proof}

\begin{corollary}[Tilted finite modes and invariance of interior limits]
\label{cor:analytic-tilted-limits}
Under $P_N^f$, every finite mode vector converges to independent circular
complex Gaussians with means
\[
 m_k=\frac{2k}{\beta}\overline{f_k}
\]
and variances $2k/\beta$. Along any unweighted interior-convergent
subsequence, the interior process under $P_N^f$ converges to that same
interior limit; it is independent of the shifted Gaussian analytic
function in the joint limit.
\end{corollary}

\begin{proof}
The preceding proof gives joint convergence of the interior, the analytic
resolvent, and $Y_N(f)$, as well as uniform integrability of $e^{Y_N(f)}$.
Exponential change of measure is therefore legitimate. Completing the
square in each Gaussian coordinate yields the stated means and unchanged
variances. The tilt depends only on the Gaussian analytic function,
which was independent of the interior; the marginal interior law is
therefore unchanged. Infinite-dimensional statements follow first on
finite coefficient vectors and then by the already established analytic
tightness under $P_N^f$.
\end{proof}

\begin{corollary}[Uniform bounds for harmonic products and insertions]
\label{cor:harmonic-product-insertion}
For $|u|<1$ set
\[
 H_N(u)=\sum_{i=1}^N\log|1-u z_i|.
\]
For every real $q$,
\begin{equation}\label{eq:harmonic-product-limit}
 \lim_{N\to\infty}\mathbb E_{N,\beta}e^{qH_N(u)}
       =(1-|u|^2)^{-q^2/(2\beta)}.
\end{equation}
For every $r<1$,
\begin{equation}\label{eq:harmonic-product-uniform}
 \sup_N\sup_{|u|\leq r}\mathbb E_{N,\beta}e^{qH_N(u)}<\infty.
\end{equation}
Consequently the insertion functions
\[
 I_N(z)=\mathbb E_{N,\beta}\prod_{i=1}^N|z-z_i|^\beta
\]
satisfy $\sup_N\sup_{|z|\leq r}I_N(z)<\infty$ for every $r<1$, and
\[
 \limsup_{N\to\infty}I_N(z)\leq(1-|z|^2)^{-\beta/2}.
\]
\end{corollary}

\begin{proof}
Apply \eqref{eq:analytic-strong-szego} to
$f(z)=(q/2)\log(1-uz)$, using the holomorphic logarithm which vanishes
at $z=0$. Its coefficients are $f_k=-q u^k/(2k)$, which give
\eqref{eq:harmonic-product-limit}.

For completeness, pointwise convergence also gives the uniform bound
\eqref{eq:harmonic-product-uniform}. For each fixed configuration,
$\exp(qH_N(u))$ is the squared modulus of the holomorphic function
\[
 \exp\left(\frac q2\sum_i\log(1-u z_i)\right).
\]
Its angular mean in $u$ is nondecreasing with $|u|$, by subharmonicity.
After expectation, rotation invariance of the unweighted gas makes
$u\mapsto\mathbb E e^{qH_N(u)}$ radial, so this function itself is
nondecreasing with $|u|$. Its supremum on $|u|\leq r$ is thus its value
at the positive real number $r$, and those values are bounded by
\eqref{eq:harmonic-product-limit}.

Finally,
\[
 |1-\bar z w|^2-|z-w|^2=(1-|z|^2)(1-|w|^2)\geq0,
 \qquad z,w\in\D.
\]
Therefore $I_N(z)\leq\mathbb E e^{\beta H_N(\bar z)}$.
The claimed bounds follow from
\eqref{eq:harmonic-product-limit}--\eqref{eq:harmonic-product-uniform}.
\end{proof}

\begin{corollary}[Exponential bounds for analytic local suprema]
\label{cor:analytic-supremum-exponential}
For every $r<1$ and $\lambda\geq0$,
\[
 \sup_N\mathbb E_{N,\beta}
       \exp\left(\lambda\sup_{|u|\leq r}|X_N(u)|\right)<\infty.
\]
In particular, all fixed moments of these local suprema are uniformly
integrable, and every finite collection of analytic linear statistics
has convergence of all joint polynomial moments.
\end{corollary}

\begin{proof}
Choose $r<s<1$. The analytic Laplace theorem, applied to constant complex
multiples of $z\mapsto sz/(1-sz)$, bounds the moment-generating function
of every fixed real linear projection of $X_N(s)$ uniformly in $N$.
Since, for $a\geq0$,
\[
 e^{a|x|}\leq e^{a(|\operatorname{Re}x|+|\operatorname{Im}x|)}
 \leq\sum_{\varepsilon,\eta\in\{-1,1\}}
       e^{a(\varepsilon\operatorname{Re}x+\eta\operatorname{Im}x)},
\]
it follows that $\sup_N\mathbb E e^{a|X_N(s)|}<\infty$ for every $a$.
Rotation invariance makes the distribution of $X_N(se^{i\theta})$
independent of $\theta$. Cauchy's estimate from the preceding proof and
Jensen's inequality now give
\begin{align*}
 \mathbb E\exp\left(\lambda\sup_{|u|\leq r}|X_N(u)|\right)
 &\leq\frac1{2\pi}\int_0^{2\pi}
       \mathbb E\exp\left(\frac{\lambda s}{s-r}
                           |X_N(se^{i\theta})|\right)d\theta\\
 &=\mathbb E\exp\left(\frac{\lambda s}{s-r}|X_N(s)|\right).
\end{align*}
This is uniformly bounded. Higher exponential bounds imply uniform
integrability of every power. Finally, the continuous coefficient
pairings $L_f$ used above are bounded by a constant times one such local
supremum, so the same conclusion holds for their finite products.
\end{proof}

\section{Harmonic boundary transport and uniqueness for \texorpdfstring{$\beta>4$}{beta greater than 4}}
\label{sec:harmonic-boundary-breakthrough}

Write
\[
 P_N(w)=\prod_{i=1}^N(1-wz_i),\qquad
 M_{N,\beta}(r)=\mathbb E_{N,\beta}|P_N(r)|^\beta,
 \qquad Q_{N,\beta}=Z_{N,\beta}/N!.
\]
The following result concerns the original gas with flat area measure in
the disk. In particular, no change of model or of the inverse temperature
is involved.

\begin{theorem}[Boundary insertion and the empty limit above $4$]
\label{thm:empty-limit-beta-above-four}
For every fixed $\beta>0$, there are positive constants $c_\beta,C_\beta$
such that, for all sufficiently large $N$,
\begin{align}
 M_{N,\beta}(1)
 &\geq c_\beta\frac{N^{\beta/2}}{(\log N)^\beta},
 \label{eq:harmonic-boundary-moment-lower}\\
 \frac{Q_{N+1,\beta}}{Q_{N,\beta}}
 &\geq c_\beta\frac{N^{\beta/2-2}}{(\log N)^\beta}.
 \label{eq:harmonic-boundary-ratio-lower}
\end{align}
Consequently, for every $\beta>4$, the full sequence of point processes
converges in probability, for the vague topology on the open disk, to
the empty configuration. More quantitatively, for every $R<1$,
\begin{equation}
 \sup_{|z|\leq R}\rho_{N,\beta}(z)
 +\mathbb P_{N,\beta}\{\Xi_N(\{|z|\leq R\})>0\}
 \leq C_{\beta,R}N^{2-\beta/2}(\log N)^\beta.
 \label{eq:empty-limit-beta-above-four-rate}
\end{equation}
The restriction to each such disk converges in total variation to the
point mass at the empty configuration, with the same upper bound up to
a constant.
\end{theorem}

We prove the theorem by a deterministic boundary perturbation and a
relative-entropy estimate. The use of a harmonic extension, instead of
rotating each particle by a function of its angle alone, is essential:
the resulting quadratic cost involves the analytic modes
$S_{j,N}=\sum_i z_i^j$ directly.

\begin{lemma}[Two estimates used in the transport argument]
\label{lem:harmonic-transport-inputs}
For every fixed $\beta>0$, with $B_N=\sum_i-\log|z_i|$, one has
\begin{equation}
 \mathbb E B_N\leq C_\beta\log(eN),\qquad
 \sum_{j=1}^N\frac{\mathbb E|S_{j,N}|^2}{j}
 \leq C_\beta N\log(eN).
 \label{eq:harmonic-transport-inputs}
\end{equation}
\end{lemma}

\begin{proof}
The radial-gap domination bounds $B_N$ in a coupling by
$\sum_{k=1}^N kE_k$, where the $E_k$ are independent exponentials with
rates $a_k=2k+\frac\beta2k(k-1)$. Its mean is
$\sum_{k\leq N}k/a_k=O_\beta(\log(eN))$.

For the second assertion let $q=1-1/N$ and $L_N=\log|\Delta_N|$.
The identity
\[
 |1-q^2z\overline w|^2-q^2|z-w|^2
 =(1-q^2|z|^2)(1-q^2|w|^2)\geq0
\]
and the logarithm power series yield
\[
 \sum_{j\geq1}\frac{q^{2j}}j|S_{j,N}|^2
 \leq-2L_N-N(N-1)\log q-N\log(1-q^2).
\]
Relative entropy with respect to normalized product area gives
\[
 \beta\mathbb E L_N\geq
 \log(Z_{N,\beta}/\pi^N)
 =\left(\frac\beta2-1\right)N\log N+O_\beta(N),
\]
where the last equality is the established partition estimate.
The expected right-hand side of the preceding inequality is therefore
$O_\beta(N\log(eN))$. For $1\leq j\leq N$,
$q^{2j}\geq q^{2N}\geq1/16$ when $N\geq2$.
This proves the claim; the case $N=1$ is absorbed into the constant.
\end{proof}

\begin{lemma}[Harmonic extension of a boundary perturbation]
\label{lem:harmonic-boundary-map}
Let $r\in(1/2,1)$, $m=(1-r)^{-1}$, and
\[
 \psi_r(\theta)=\sum_{k\geq1}\frac{r^k}{k}\sin(k\theta),
 \qquad \ell_r=-\log(1-r^2).
\]
Let $\varepsilon>0$, and let $T$ be the harmonic Poisson extension to
$\mathbb D$ of
\[
 e^{i\theta}\longmapsto e^{i\theta}e^{i\varepsilon\psi_r(\theta)}.
\]
Write $T(z)=z+V(z)$ and write the Fourier coefficients of its boundary
displacement as
\[
 e^{i\theta}(e^{i\varepsilon\psi_r(\theta)}-1)
 =\sum_{n\in\mathbb Z}c_ne^{in\theta}.
\]
There are absolute constants $c,C>0$ such that, if
$\varepsilon m\leq c$, then $T$ is a smooth diffeomorphism of the disk,
extending to a homeomorphism of its closure, and
\begin{align}
 \|V\|_\infty&\leq C\varepsilon,
 &\|DV\|_\infty&\leq C\varepsilon m,
 \label{eq:harmonic-map-uniform-bounds}\\
 \sum_{n\in\mathbb Z}|n||c_n|^2
 &\leq\varepsilon^2(\ell_r/2+C),
 &|c_1|&\leq C\varepsilon^2,
 \label{eq:harmonic-map-half-derivative}\\
 \sum_{n\in\mathbb Z}n^2|c_n|^2e^{|n|/(2m)}
 &\leq C\varepsilon^2m.
 \label{eq:harmonic-map-strip-bound}
\end{align}
Here $DV$ is the real derivative matrix. The coefficient $c_1$ is real
and nonpositive.
\end{lemma}

\begin{proof}
Since $\psi_r=-\operatorname{Arg}(1-re^{i\theta})$ on the circle,
\[
 \|\psi_r\|_\infty\leq\pi/2,
 \quad \|\psi_r'\|_\infty\leq m,
 \quad \|\psi_r\|_2^2=\tfrac12\sum_{k\geq1}r^{2k}/k^2\leq C.
\]
The boundary displacement has modulus at most $\varepsilon\pi/2$;
the Poisson representation therefore gives the first bound in
\eqref{eq:harmonic-map-uniform-bounds}. Moreover,
\[
 c_1=\int_0^{2\pi}(e^{i\varepsilon\psi_r(\theta)}-1)
                         \frac{d\theta}{2\pi}.
\]
Oddness of $\psi_r$ makes this real, and the inequalities
$0\leq1-\cos x\leq x^2/2$ give the stated bounds on $c_1$.

Set $u=e^{i\varepsilon\psi_r}-1$ and denote its Fourier coefficients
by $u_k$. The Fourier identity
\[
 \sum_{k\in\mathbb Z}|k||u_k|^2
 =\int\!\!\int
   \frac{|u(\theta)-u(\phi)|^2}{4\sin^2((\theta-\phi)/2)}
        \frac{d\theta}{2\pi}\frac{d\phi}{2\pi}
\]
follows by Fourier orthogonality and the Fej\'er identity
$\int(1-\cos(k\theta))/(1-\cos\theta)\,d\theta/(2\pi)=|k|$.
It is first proved for trigonometric polynomials and then by smooth
approximation; here all functions are analytic on a strip, so convergence
in every Sobolev norm is available. Since
$|e^{i\varepsilon a}-e^{i\varepsilon b}|\leq\varepsilon|a-b|$
for real $a,b$, the displayed identity gives
\[
 \sum_k|k||u_k|^2
 \leq\varepsilon^2\sum_k|k||\widehat\psi_r(k)|^2
 =\varepsilon^2\ell_r/2.
\]
As $c_n=u_{n-1}$ and $|n|\leq|n-1|+1$, adding
$\|u\|_2^2\leq\varepsilon^2\|\psi_r\|_2^2$ proves
\eqref{eq:harmonic-map-half-derivative}.

For completeness, we prove the exponentially weighted estimate rather
than using a derivative bound for a singular boundary kernel. Set
$h=1/(4m)$. The series for $\psi_r$ extends holomorphically to
$|\operatorname{Im}\theta|\leq h$, because
$re^h\leq e^{-3/(4m)}<1$. On either line $\theta=x\pm ih$, Fourier
orthogonality gives
\[
 \|\psi_r(\cdot\pm ih)\|_2\leq C,
 \qquad
 \|\psi_r'(\cdot\pm ih)\|_2\leq C\sqrt m,
 \qquad
 \|\psi_r(\cdot\pm ih)\|_\infty\leq C\log(em).
\]
The condition $\varepsilon m\leq c$ bounds
$e^{\varepsilon|\operatorname{Im}\psi_r|}$ by an absolute constant.
Differentiating
$b(\theta)=e^{i\theta}(e^{i\varepsilon\psi_r(\theta)}-1)$ and using
$|e^w-1|\leq|w|e^{|w|}$ therefore shows
\[
 \|b'(\cdot\pm ih)\|_2^2\leq C\varepsilon^2m.
\]
Parseval's identity on the two lines, followed by addition, yields
\eqref{eq:harmonic-map-strip-bound}.
The harmonic extension is explicitly
\[
 V(z)=c_0+\sum_{n\geq1}c_nz^n+\sum_{n\leq-1}c_n\overline z^{\,|n|}.
\]
Cauchy--Schwarz and \eqref{eq:harmonic-map-strip-bound} give
\[
 \|DV\|_\infty
 \leq\sum_n|n||c_n|
 \leq
 \left(\sum_n n^2|c_n|^2e^{|n|/(2m)}\right)^{1/2}
 \left(\sum_{n\ne0}e^{-|n|/(2m)}\right)^{1/2}
 \leq C\varepsilon m.
\]
Shrinking $c$ if necessary makes this smaller than $1/2$.

The boundary map is a circle diffeomorphism, since its angular derivative
is $1+\varepsilon\psi_r'>0$. The Poisson extension maps the disk strictly
inside the disk: it is an average of unit vectors, and equality in the
triangle inequality would force the nonconstant boundary map to be
constant. Also, on the convex closed disk,
\[
 |T(z)-T(w)|\geq(1-\|DV\|_\infty)|z-w|.
\]
Thus $T$ is injective and has everywhere invertible real derivative.
Its image of the open disk is open by the inverse function theorem.
It is also closed relative to the open disk: if $T(z_n)\to y\in\mathbb D$,
a convergent subsequence in the closed disk gives $T(z)=y$, and $z$ cannot
be on the boundary because $T$ maps that boundary to the unit circle.
Connectedness of the disk proves surjectivity. The inverse function
theorem and compactness now give all claimed regularity properties.
\end{proof}

\begin{lemma}[A pairwise bound involving only analytic modes]
\label{lem:harmonic-pair-balayage}
For $p\geq1$ put $D_p(z,w)=(z^p-w^p)/(z-w)$, extended polynomially to
the diagonal. For any $z_1,\ldots,z_N\in\mathbb D$,
\begin{equation}
 \sum_{i<j}|D_p(z_i,z_j)|^2
 \leq\frac{pN^2}{2}
       +\sum_{j=1}^{p-1}(p-j)|S_{j,N}|^2.
 \label{eq:harmonic-pair-balayage}
\end{equation}
Consequently, for the harmonic displacement in
Lemma~\ref{lem:harmonic-boundary-map}, rotation invariance gives
\begin{equation}
 \mathbb E\sum_{i<j}
       \left|\frac{V(z_i)-V(z_j)}{z_i-z_j}\right|^2
 \leq\frac{N^2}{2}\sum_n|n||c_n|^2
       +\sum_{j\geq1}K_j\mathbb E|S_{j,N}|^2,
 \label{eq:harmonic-pair-quadratic}
\end{equation}
where $K_j=\sum_{|n|>j}(|n|-j)|c_n|^2$.
\end{lemma}

\begin{proof}
For each $i$, let $\zeta_i$ be an independent point on the unit circle
with Poisson distribution based at $z_i$. Its density is
$(1-|z_i|^2)/|e^{i\theta}-z_i|^2$ relative to $d\theta/(2\pi)$,
and its moments are $\mathbb E\zeta_i^j=z_i^j$ for $j\geq0$.
These statements follow directly by expanding the Poisson kernel.
Since $D_p$ is holomorphic in each variable,
\[
 D_p(z_i,z_l)=\mathbb E[D_p(\zeta_i,\zeta_l)],\qquad i\ne l,
\]
and Jensen's inequality bounds its squared modulus by the corresponding
second moment. On the circle, the Fej\'er identity gives
\[
 \sum_{i<l}|D_p(\zeta_i,\zeta_l)|^2
 =\frac{pN(N-1)}2+
   \sum_{j=1}^{p-1}(p-j)\left(\left|\sum_i\zeta_i^j\right|^2-N\right).
\]
Independence and the Poisson moments imply
\[
 \mathbb E\left|\sum_i\zeta_i^j\right|^2
 =|S_{j,N}|^2+\sum_i(1-|z_i|^{2j}).
\]
Substitution and deletion of nonpositive terms prove
\eqref{eq:harmonic-pair-balayage}.

For the second assertion, apply a common rotation $z_i\mapsto e^{i\alpha}z_i$.
In the divided difference of $V$, the positive Fourier mode $n\geq1$
has rotation frequency $n-1$, while the negative mode $n=-p\leq-1$
has frequency $-p-1$. These frequencies are all distinct. Orthogonality
therefore makes the rotational average of its squared modulus equal to
\[
 \sum_{n\ne0}|c_n|^2|D_{|n|}(z_i,z_l)|^2.
\]
The harmless unimodular factor in an antiholomorphic divided difference
disappears in its squared modulus. Use rotation invariance, sum
\eqref{eq:harmonic-pair-balayage} with these nonnegative weights, and
interchange nonnegative sums to obtain \eqref{eq:harmonic-pair-quadratic}.
\end{proof}

\begin{proposition}[Mesoscopic moment lower bound by harmonic transport]
\label{prop:harmonic-mesoscopic-moment}
Suppose $m_N\to\infty$ and $m_N\log(eN)/N\to0$, and put
$r_N=1-1/m_N$. Then, for fixed $\beta>0$,
\begin{equation}
 \log M_{N,\beta}(r_N)
 \geq\frac\beta2\log\frac1{1-r_N^2}-C_\beta-o(1).
 \label{eq:harmonic-mesoscopic-moment}
\end{equation}
\end{proposition}

\begin{proof}
Suppress subscripts on $r,m$ and put $\varepsilon=2/N$.
Use the map $T=\mathrm{id}+V$ of Lemma~\ref{lem:harmonic-boundary-map},
and let $\mathsf Q$ be the pushforward of the gas law $\mathsf P$ under
$T$ in all coordinates. Set $\eta=\|DV\|_\infty=o(1)$ and
\[
 W(z)=\beta\sum_i\log|1-rz_i|.
\]
The Gibbs variational inequality, which follows from nonnegativity of
relative entropy with respect to $e^W\mathsf P/\mathbb E_{\mathsf P}e^W$,
says
\begin{equation}
 \log M_{N,\beta}(r)
 \geq\mathbb E_{\mathsf Q}W-
                  \operatorname{Ent}(\mathsf Q\mid\mathsf P).
 \label{eq:harmonic-gibbs-variational}
\end{equation}
All changes of variables are legitimate: $T$ is a diffeomorphism,
its Jacobian is bounded above and below, and
$|Tz-Tw|/|z-w|$ lies between $1-\eta$ and $1+\eta$.

\emph{Entropy cost.}
For any complex $q$ with $|q|\leq\eta<1$,
Taylor's formula gives
\[
 \log|1+q|\geq\operatorname{Re}q-
                      \frac{|q|^2}{2(1-\eta)^2}.
\]
Likewise, the real matrix logarithm, or Taylor's formula for the
determinant in dimension two, gives
\[
 \log\det(I+DV)\geq\operatorname{div}V-C|DV|^2
\]
when $\eta\leq1/2$. The exact density change hence implies
\begin{align}
 \operatorname{Ent}(\mathsf Q\mid\mathsf P)
 \leq{}&-\mathbb E\sum_i\operatorname{div}V(z_i)
        +C\mathbb E\sum_i|DV(z_i)|^2\notag\\
 &-\beta\mathbb E\operatorname{Re}\sum_{i<j}
                  \frac{V(z_i)-V(z_j)}{z_i-z_j}
 +\frac\beta{2(1-\eta)^2}
       \mathbb E\sum_{i<j}
          \left|\frac{V(z_i)-V(z_j)}{z_i-z_j}\right|^2.
 \label{eq:harmonic-entropy-decomposition}
\end{align}
Global rotational averaging gives
\[
 \mathbb E\sum_i\operatorname{div}V(z_i)=2N\operatorname{Re}c_1,
 \qquad
 \mathbb E\sum_{i<j}\frac{V(z_i)-V(z_j)}{z_i-z_j}
 =\frac{N(N-1)}2c_1.
\]
Both linear costs are therefore $O_\beta(N^2\varepsilon^2)=O_\beta(1)$.
The explicit harmonic series also gives, by angular averaging at each
radius and then expectation,
\[
 \mathbb E\sum_i|DV(z_i)|^2
 \leq C N\sum_n n^2|c_n|^2
 \leq C N\varepsilon^2m.
\]
This estimate holds for the Frobenius matrix norm and hence for any
fixed choice of real matrix norm, after changing the absolute constant.

To control the final term in \eqref{eq:harmonic-entropy-decomposition},
notice from \eqref{eq:harmonic-map-strip-bound} that
\[
 \sup_{j\geq1}jK_j\leq\tfrac14\sum_n n^2|c_n|^2
 \leq C\varepsilon^2m,
\]
because $j(p-j)\leq p^2/4$. Also,
\[
 \sum_{j>N}K_j\leq\tfrac12\sum_{|n|>N}n^2|c_n|^2
 \leq C\varepsilon^2m e^{-N/(2m)}.
\]
Lemma~\ref{lem:harmonic-transport-inputs} for $j\leq N$, and the
deterministic bound $|S_{j,N}|\leq N$ for $j>N$, consequently give
\begin{equation}
 \sum_{j\geq1}K_j\mathbb E|S_{j,N}|^2
 \leq C_\beta\varepsilon^2mN\log(eN)
       +CN^2\varepsilon^2m e^{-N/(2m)}.
 \label{eq:harmonic-quadratic-error}
\end{equation}
Combine \eqref{eq:harmonic-pair-quadratic},
\eqref{eq:harmonic-map-half-derivative}, and
\eqref{eq:harmonic-quadratic-error}. Since
$\eta\ell_r=O(m\log(eN)/N)=o(1)$, this proves
\begin{equation}
 \operatorname{Ent}(\mathsf Q\mid\mathsf P)
 \leq\frac{\beta\varepsilon^2N^2}{8}\ell_r+C_\beta+o(1)
 =\frac\beta2\ell_r+C_\beta+o(1).
 \label{eq:harmonic-entropy-final}
\end{equation}
Here the tail error tends to zero: the hypothesis implies
$N/m\gg\log N$, and therefore $m e^{-N/(2m)}\to0$.

\emph{Gain in the inserted potential.}
Let $v$ be the harmonic extension of $i e^{i\theta}\psi_r(\theta)$.
It is given explicitly by
\[
 v(z)=\frac12\sum_{k\geq1}\frac{r^k}{k}
                   (z^{k+1}-\overline z^{\,k-1}).
\]
The boundary inequality
$|e^{i\varepsilon x}-1-i\varepsilon x|\leq\varepsilon^2x^2/2$
and the Poisson representation show
\begin{equation}
 \|V-\varepsilon v\|_\infty\leq C\varepsilon^2,
 \qquad \|v\|_\infty\leq\pi/2.
 \label{eq:harmonic-map-linearization}
\end{equation}
Use the holomorphic logarithm $F(z)=\beta\log(1-rz)$, with $F(0)=0$,
so that $W=\operatorname{Re}\sum_iF(z_i)$.
Rotation invariance implies $\mathbb E_{\mathsf P}W=0$.
Expanding $F'(z)=-\beta\sum_{k\geq1}r^kz^{k-1}$ and averaging a
common rotation gives the exact identity
\[
 \varepsilon\mathbb E\operatorname{Re}\sum_iF'(z_i)v(z_i)
 =\frac{\varepsilon\beta}{2}
       \sum_{k\geq1}\frac{r^{2k}}k
                     \mathbb E\sum_i|z_i|^{2k-2}.
\]
For $k\geq1$, $1-u^{2k-2}\leq2k(-\log u)$ on $0<u\leq1$.
Lemma~\ref{lem:harmonic-transport-inputs} therefore bounds the preceding
quantity below by
\[
 \frac{\varepsilon\beta N}{2}\ell_r
       -C_\beta\varepsilon m\log(eN)
 =\beta\ell_r-o(1).
\]
The linear error from $V-\varepsilon v$ is at most
$CN\|F'\|_\infty\varepsilon^2\leq C_\beta N\varepsilon^2m=o(1)$.

The quadratic Taylor remainder also has this latter bound; the following
averaged estimate avoids the incorrect larger bound of order
$N\varepsilon^2m^2$. For $0\leq s\leq1$, the uniform displacement
estimate and $\varepsilon m=o(1)$ imply
\[
 |1-r(z+sV(z))|\geq(1-C\varepsilon m)|1-rz|
                          \geq\tfrac12|1-rz|.
\]
Consequently
$|F''(z+sV(z))|\leq4\beta r^2/|1-rz|^2$.
At any fixed radius $u\leq1$,
\[
 \int_0^{2\pi}\frac{d\theta/(2\pi)}{|1-ru e^{i\theta}|^2}
 =\frac1{1-r^2u^2}\leq\frac1{1-r^2}\leq m.
\]
Using $\|V\|_\infty\leq C\varepsilon$, Taylor's formula and rotation
invariance thus bound the expected total absolute remainder by
$C_\beta N\varepsilon^2m=o(1)$. We have proved
\begin{equation}
 \mathbb E_{\mathsf Q}W\geq\beta\ell_r-o(1).
 \label{eq:harmonic-gain-final}
\end{equation}
Equations \eqref{eq:harmonic-gibbs-variational},
\eqref{eq:harmonic-entropy-final}, and
\eqref{eq:harmonic-gain-final} prove the proposition.
\end{proof}

\begin{proof}[Proof of Theorem~\ref{thm:empty-limit-beta-above-four}]
Take $m_N=N/(\log N)^2$ and $r_N=1-1/m_N$ for all sufficiently large
$N$. Proposition~\ref{prop:harmonic-mesoscopic-moment} applies and gives
\[
 M_{N,\beta}(r_N)
 \geq c_\beta(1-r_N^2)^{-\beta/2}
 \geq c_\beta m_N^{\beta/2}
 =c_\beta\frac{N^{\beta/2}}{(\log N)^\beta}.
\]
For each configuration, $w\mapsto|P_N(w)|^\beta$ is subharmonic, so
its circular means are nondecreasing. Rotation invariance identifies
the expectation of those means with $M_{N,\beta}(r)$. Thus
$M_{N,\beta}(1)\geq M_{N,\beta}(r_N)$, proving
\eqref{eq:harmonic-boundary-moment-lower}.

We include the exact boundary insertion recursion. Choose the particle
of largest radius in an $(N+1)$-particle configuration, write it as
$te^{i\theta}$, and express the remaining particles as
$te^{i\theta}w_i$, $w_i\in\mathbb D$. Ties have zero probability.
The polar Jacobian and homogeneity of the Vandermonde give
\[
 Z_{N+1,\beta}
 =\frac{2\pi(N+1)}{a_{N+1}}Z_{N,\beta}M_{N,\beta}(1),
 \qquad a_{N+1}=2(N+1)+\frac\beta2N(N+1).
\]
Indeed, the radial integral is
$\int_0^1t^{a_{N+1}-1}\,dt=a_{N+1}^{-1}$, and the common angle
contributes $2\pi$. Dividing by factorials yields
\[
 \frac{Q_{N+1,\beta}}{Q_{N,\beta}}
 =\frac{2\pi}{a_{N+1}}M_{N,\beta}(1),
\]
and proves \eqref{eq:harmonic-boundary-ratio-lower}.

For the conclusion about the process,
Corollary~\ref{cor:analytic-insertion-without-clt} has already proved,
without using a central limit theorem,
\[
 \sup_n\sup_{|a|\leq R}
 \mathbb E_{n,\beta}\left|\prod_{i=1}^n(1-a z_i)\right|^\beta
 <\infty\qquad(R<1).
\]
The disk identity
\[
 |1-\overline z w|^2-|z-w|^2
 =(1-|z|^2)(1-|w|^2)\geq0
\]
and the one-point insertion formula imply
\begin{align*}
 \rho_{N,\beta}(z)
 &=\frac{Q_{N-1,\beta}}{Q_{N,\beta}}
       \mathbb E_{N-1,\beta}\prod_i|z-w_i|^\beta\\
 &\leq\frac{Q_{N-1,\beta}}{Q_{N,\beta}}
       \mathbb E_{N-1,\beta}
          \left|\prod_i(1-\overline z w_i)\right|^\beta.
\end{align*}
The uniform harmonic-moment bound and
\eqref{eq:harmonic-boundary-ratio-lower} give the intensity assertion in
\eqref{eq:empty-limit-beta-above-four-rate}. Integrating over the disk
of radius $R$ and applying Markov's inequality gives the probability
assertion. This tends to zero when $\beta>4$.

For a probability law $\nu$ on configurations, its total variation
distance to the point mass at the empty configuration is
$1-\nu(\{\varnothing\})$, using the convention
$\|\nu-\mu\|_{\mathrm{TV}}=\sup_A|\nu(A)-\mu(A)|$.
This proves the local total variation assertion. Finally, every compact
subset of the open disk lies in a disk of radius $R<1$. For each compactly
supported continuous test function, its integral against $\Xi_N$ is
therefore zero with probability tending to one. Equivalently, on any
countable family of test functions defining a metric for vague
convergence, finite unions and then the summable tail of the metric show
convergence in probability to the empty measure. This proves the
full-sequence assertion and hence uniqueness of the limiting law for
$\beta>4$.
\end{proof}

\section{Quantitative consequences of the empty limit}
\label{sec:empty-limit-correlations}

Fix $\beta>4$ and put
\[
 b_N=N^{2-\beta/2}(\log(eN))^\beta.
\]
Thus $b_N\to0$. In this section we use the adjacent-ratio lower bound
\eqref{eq:harmonic-boundary-ratio-lower}, in the equivalent form
\begin{equation}\label{eq:empty-ratio-input}
 \frac{Q_{N-1,\beta}}{Q_{N,\beta}}\leq C_\beta b_N
\end{equation}
for all sufficiently large $N$, and the uniform harmonic-product bound
of Corollary~\ref{cor:harmonic-product-insertion}.

\begin{corollary}[Decay of every correlation and compact count moment]
\label{cor:empty-correlation-rates}
For every fixed integer $p\geq1$ and $r<1$, there is a constant
$C_{\beta,p,r}<\infty$ such that, for all sufficiently large $N$,
\begin{equation}\label{eq:empty-p-correlation-rate}
 \rho_{p,N}(z_1,\ldots,z_p)
 \leq C_{\beta,p,r}|\Delta_p(z_1,\ldots,z_p)|^\beta b_N^p,
 \qquad |z_1|,\ldots,|z_p|\leq r.
\end{equation}
Here $\Delta_1=1$. Consequently, for every relatively compact Borel
set $B\subset\D$, with $M_N=\Xi_N(B)$, and every fixed integer $p\geq1$,
\begin{equation}\label{eq:empty-factorial-hit-rates}
 \mathbb E(M_N)_p=O_{\beta,p,B}(b_N^p),\qquad
 \mathbb P(M_N\geq p)=O_{\beta,p,B}(b_N^p),\qquad
 \mathbb E M_N^p=O_{\beta,p,B}(b_N).
\end{equation}
For each fixed $t\geq0$, one also has
\begin{equation}\label{eq:empty-count-exponential-rate}
 \mathbb E(e^{tM_N}-1)=O_{\beta,t,B}(b_N).
\end{equation}
In particular, the local counting generating functions converge to $1$
uniformly on compact subsets of the complex parameter plane, at rate
$O(b_N)$.
\end{corollary}

\begin{proof}
The exact $p$-point insertion formula, for $N\geq p$, is
\begin{align*}
 \rho_{p,N}(z_1,\ldots,z_p)
 ={}&\frac{Q_{N-p,\beta}}{Q_{N,\beta}}|\Delta_p(z_1,\ldots,z_p)|^\beta\\
 &\quad\times\mathbb E_{N-p,\beta}
    \prod_{i=1}^p\prod_{j=1}^{N-p}|z_i-w_j|^\beta.
\end{align*}
Indeed, the combinatorial factor $N!/(N-p)!$ in the correlation function
cancels against the factorials in the definition $Q_n=Z_n/n!$.
The disk inequality $|z-w|\leq|1-\overline z w|$ and H\"older's
inequality give
\begin{align*}
 \mathbb E_{N-p,\beta}
    \prod_{i=1}^p\prod_j|z_i-w_j|^\beta
 &\leq\mathbb E_{N-p,\beta}
        \prod_{i=1}^p e^{\beta H_{N-p}(\overline z_i)}\\
 &\leq\prod_{i=1}^p
       \left(\mathbb E_{N-p,\beta}
                   e^{p\beta H_{N-p}(\overline z_i)}\right)^{1/p}
 \leq C_{\beta,p,r}.
\end{align*}
The last bound is uniform in the tuple on the specified compact disk.
For fixed $p$, multiplying \eqref{eq:empty-ratio-input} at
$N,N-1,\ldots,N-p+1$ gives
\[
 \frac{Q_{N-p,\beta}}{Q_{N,\beta}}
 \leq C_{\beta,p}\prod_{j=0}^{p-1}b_{N-j}
 \leq C'_{\beta,p}b_N^p.
\]
The last inequality uses $b_{N-j}/b_N\to1$ for fixed $j$. This proves
\eqref{eq:empty-p-correlation-rate}.

Choose $r<1$ with $B\subset B_r$. Integrating the correlation bound
over $B^p$ proves the factorial-moment assertion; the Vandermonde is
bounded on this compact product. Since $(M_N)_p\geq p!$ on
$\{M_N\geq p\}$, Markov's inequality proves the probability bound.
The identity
\[
 M_N^p=\sum_{j=1}^p
       \left\{\begin{matrix}p\\j\end{matrix}\right\}(M_N)_j
\]
with the Stirling numbers of the second kind gives the ordinary-moment
bound, because $b_N\to0$.

Finally, the earlier Gaussian compact-count bound gives
$\sup_N\mathbb E e^{2tM_N}<\infty$. Splitting at counts zero, one, and
at least two, and applying Cauchy--Schwarz on the last event, yields
\begin{align*}
 \mathbb E(e^{tM_N}-1)
 &\leq(e^t-1)\mathbb P(M_N=1)
       +\mathbb E[e^{tM_N}\mathbf1_{\{M_N\geq2\}}]\\
 &\leq(e^t-1)\mathbb EM_N
       +(\mathbb Ee^{2tM_N})^{1/2}
             \mathbb P(M_N\geq2)^{1/2}
 =O(b_N).
\end{align*}
For complex $w$ with $|w|\leq t$, the elementary inequality
$|e^{wM_N}-1|\leq e^{tM_N}-1$ proves the locally uniform assertion.
\end{proof}

\begin{corollary}[A quantitative description of the rare interior points]
\label{cor:empty-bernoulli-approximation}
Let $B\Subset\D$ be Borel and put
$\mu_N=\int_B\rho_{1,N}(z)\,dA(z)$. For all sufficiently large $N$,
$0\leq\mu_N\leq1$. Let $\mathsf B_N$ be the probability law which is
empty with probability $1-\mu_N$ and has one point with intensity measure
$\mathbf1_B(z)\rho_{1,N}(z)\,dA(z)$. Then
\[
 d_{\mathrm{TV}}\bigl(\mathcal L(\Xi_N|_B),\mathsf B_N\bigr)
 \leq\mathbb E[M_N\mathbf1_{\{M_N\geq2\}}]
 \leq\mathbb E(M_N)_2=O_{\beta,B}(b_N^2).
\]
This approximation uses the actual finite-$N$ intensity; it does not
assert convergence of a rescaled intensity profile.
\end{corollary}

\begin{proof}
The first corollary gives $\mu_N=O(b_N)$, so $\mathsf B_N$ is well-defined
for large $N$. On one-point configurations, the actual probability
measure is dominated by the intensity measure defining $\mathsf B_N$.
Its deficit in total mass is
$\mu_N-\mathbb P(M_N=1)=\mathbb E[M_N\mathbf1_{\{M_N\geq2\}}]$.
On the empty configuration the actual law has excess mass
$\mu_N-\mathbb P(M_N\geq1)$, and on configurations with at least two
points it has excess mass $\mathbb P(M_N\geq2)$. The latter two masses
sum to the one-point deficit. The total variation distance, with the
convention $d_{\mathrm{TV}}(P,Q)=\sup_A|P(A)-Q(A)|$, therefore equals
that deficit. Since $m\leq m(m-1)$ for integers $m\geq2$, the claimed
bound follows from \eqref{eq:empty-factorial-hit-rates} at $p=2$.
\end{proof}

\section{Adjacent partition ratios on a set of full natural density}
\label{sec:partition-density-one}

Set $c_\beta=\beta/2-2$ and
\[
 R_{n,\beta}=\frac{Q_{n+1,\beta}}{Q_{n,\beta}},\qquad n\geq1.
\]
The next theorem is weaker than a full-sequence ratio asymptotic.
It makes no assertion about uniqueness of the interior point process.

\begin{theorem}[Full-density exponent for adjacent ratios]
\label{thm:partition-ratio-density-one}
For every $\beta>0$ and every $\varepsilon>0$,
\begin{equation}
 \#\left\{2\leq n\leq N:
 \left|\frac{\log R_{n,\beta}}{\log n}-c_\beta\right|>\varepsilon
 \right\}
 =O_{\beta,\varepsilon}\left(
                   \frac{N\log\log(e^eN)}{\log(eN)}\right).
 \label{eq:partition-ratio-density-exceptions}
\end{equation}
There exists a set $\mathcal D_\beta\subset\mathbb N$ of natural density
one, meaning
$\#(\mathcal D_\beta\cap\{1,\ldots,N\})/N\to1$, such that
\begin{equation}
 R_{n,\beta}=n^{c_\beta+o(1)},
 \qquad n\to\infty,\quad n\in\mathcal D_\beta.
 \label{eq:partition-ratio-density-one}
\end{equation}
\end{theorem}

\begin{proof}
The boundary-transport estimate
\eqref{eq:harmonic-boundary-ratio-lower} gives, after increasing a
constant to cover finitely many indices,
\begin{equation}
 s_n:=\log R_{n,\beta}-c_\beta\log n
 \geq-C_\beta-\beta\log\log(en),\qquad n\geq2.
 \label{eq:partition-centered-ratio-lower}
\end{equation}
On the other hand, telescoping and
$\log Q_{N,\beta}=c_\beta N\log N+O_\beta(N)$ imply
\begin{equation}
 \sum_{n=2}^N s_n=O_\beta(N).
 \label{eq:partition-centered-ratio-sum}
\end{equation}
Indeed, the sum of the log ratios is
$\log Q_{N+1,\beta}-\log Q_{2,\beta}$, whereas
$\sum_{n=2}^N\log n=\log(N!)=N\log N-N+O(\log N)$.

Let $A_{N,\varepsilon}=\{2\leq n\leq N:s_n>\varepsilon\log n\}$.
Combining \eqref{eq:partition-centered-ratio-lower} and
\eqref{eq:partition-centered-ratio-sum} gives
\[
 \varepsilon\sum_{n\in A_{N,\varepsilon}}\log n
 \leq C_\beta N\log\log(e^eN).
\]
At most $\sqrt N$ of these indices are below $\sqrt N$, and every
remaining logarithm is at least $\tfrac12\log N$. Consequently
\[
 \#A_{N,\varepsilon}
 \leq\sqrt N+
    \frac{C_\beta N\log\log(e^eN)}{\varepsilon\log N}.
\]
Equation \eqref{eq:partition-centered-ratio-lower} rules out
$s_n<-\varepsilon\log n$ for all sufficiently large $n$, so this proves
\eqref{eq:partition-ratio-density-exceptions}.

We give an explicit full-density selection. For all sufficiently large
$n$, let
\[
 \delta_n=\sqrt{\frac{\log\log n}{\log n}}.
\]
This is positive, decreases eventually, and tends to zero. Include an
index $n$ in $\mathcal D_\beta$ precisely when
$s_n\leq\delta_n\log n$, allowing an arbitrary choice at finitely many
small indices. If $n\leq N$ is a sufficiently large excluded index,
then $s_n>\delta_N\log n$. The preceding counting argument, with
$\varepsilon=\delta_N$, therefore gives
\[
 \#\bigl(\{1,\ldots,N\}\setminus\mathcal D_\beta\bigr)
 \leq O_\beta(1)+\sqrt N+
     C_\beta N\sqrt{\frac{\log\log N}{\log N}}=o(N).
\]
On $\mathcal D_\beta$, the upper bound $s_n/\log n\leq\delta_n\to0$
and the lower bound \eqref{eq:partition-centered-ratio-lower} squeeze
$s_n/\log n$ to zero. This proves
\eqref{eq:partition-ratio-density-one}.
\end{proof}

\begin{corollary}[Diverging total radial depth below $4$ on a full-density set]
\label{cor:depth-density-one-below-four}
Fix $0<\beta<4$ and put
\[
 B_n=\sum_{i=1}^n-\log|z_i|,\qquad
 d_\beta=\frac2\beta-\frac12>0.
\]
On the full-density set $\mathcal D_\beta$ supplied by the preceding
theorem, for every $0\leq a<d_\beta$,
\begin{equation}
 \mathbb P_{n,\beta}\{B_n\leq a\log n\}
 \leq n^{\beta a+c_\beta+o(1)}\longrightarrow0,
 \qquad n\to\infty,\quad n\in\mathcal D_\beta.
 \label{eq:depth-density-one-below-four}
\end{equation}
In particular, $B_n\to\infty$ in probability along this set of indices.
\end{corollary}

\begin{proof}
The central insertion identity at size $n+1$ gives
\[
 \mathbb E_{n,\beta}e^{-\beta B_n}
 =\rho_{n+1,\beta}(0)\,R_{n,\beta}.
\]
The already established compact correlation estimates imply
$\sup_n\rho_{n+1,\beta}(0)<\infty$. Hence
\eqref{eq:partition-ratio-density-one} gives
\[
 \mathbb E_{n,\beta}e^{-\beta B_n}
 \leq n^{c_\beta+o(1)}
 \qquad(n\to\infty,\ n\in\mathcal D_\beta).
\]
On the event $B_n\leq a\log n$, the integrand is at least
$n^{-\beta a}$. This proves
\eqref{eq:depth-density-one-below-four}; its exponent is negative exactly
when $a<-c_\beta/\beta=d_\beta$.
Choosing any fixed positive $a<d_\beta$ also proves divergence in
probability.
\end{proof}

\paragraph{Limitation.}
The total radial depth is not a compactly supported observable on the
open disk. Its divergence can be caused by particles approaching the
boundary. Therefore the last corollary does not prove that an interior
subsequential limit is nonempty, nor does it identify that limit. Likewise,
the full-density ratio theorem does not control the exceptional indices
and cannot be upgraded to a full-sequence asymptotic without a further
regularity argument for the actual adjacent partition ratios.

\section{Boundary-vanishing statistics and changes of potential}
\label{sec:weighted-boundary-tilts}

Write $d(z)=1-|z|$. The distinction between $d^p$ for $p>1$ and the
borderline weight $d$ is useful when analyzing the exterior field.
Here we obtain exponential, rather than only polynomial, control of
the first class of statistics.

\begin{lemma}[Uniform exponential moments of scaled disk counts]
\label{lem:scaled-count-exponential}
For every $\beta>0$ and $\lambda\ge0$,
\[
 \sup_{N\ge1}\sup_{0<t<1}
 \mathbb E_{N,\beta}\exp\{\lambda t M_N(1-t)\}<\infty.
\]
The same assertion holds for every subsequential limit, uniformly over
all such limits.
\end{lemma}

\begin{proof}
Put $M=M_N(1-t)$ and $a=\beta/16$. The compact-count bound, together
with $\log(1/(1-t))\ge t$, gives
\[
 \mathbb P(M\ge m)\le
 \exp\{m-\tfrac\beta4t m(m-1)\}.
\]
If $m\ge K_t:=\lceil 8/(\beta t)\rceil+2$, the exponent is at most
$-a t m^2$. For $q=\lambda t$, summation by parts gives
\[
 \mathbb Ee^{qM}
 =1+(e^q-1)\sum_{m\ge1}e^{q(m-1)}\mathbb P(M\ge m).
\]
This is first an identity for truncated integer-valued variables and
then follows in general by monotone convergence. The terms below
$K_t$ contribute at most $e^{qK_t}$, which is bounded uniformly in $t$.
For the rest, use $e^q-1\le\lambda e^\lambda t$ and
\[
 \lambda t m-a t m^2
 \le \frac{\lambda^2t}{2a}-\frac a2t m^2.
\]
The elementary integral bound
$\sum_{m\ge1}e^{-at m^2/2}\le C_a/\sqrt t$ shows that the remaining
contribution is bounded uniformly as well. The limiting count-tail
bound is the same, so the proof applies verbatim to every limit.
\end{proof}

\begin{theorem}[Exponential control and convergence of weighted statistics]
\label{thm:boundary-weight-exponential}
Fix $p>1$ and define
\[
 U_p(\xi)=\sum_{z\in\xi}d(z)^p,\qquad
 V_{p,\delta}(\xi)=\sum_{\substack{z\in\xi\\d(z)<\delta}}d(z)^p,
 \quad 0<\delta<1.
\]
For every $c\ge0$ there is a constant $C_{\beta,p,c}$ such that
\begin{align}
 \sup_N\mathbb E e^{cU_p(\Xi_N)}&\le C_{\beta,p,c},
 \label{eq:boundary-weight-all-exponentials}\\
 \sup_N\mathbb E
  \exp\{c\delta^{1-p}V_{p,\delta}(\Xi_N)\}
 &\le C_{\beta,p,c},\qquad 0<\delta<1.
 \label{eq:boundary-weight-tail-exponentials}
\end{align}
These bounds also hold for every subsequential limit with the same
constants. In particular, for every fixed $q\ge1$,
\[
 \sup_N\|V_{p,\delta}(\Xi_N)\|_{L^q}
 +\sup_\Xi\|V_{p,\delta}(\Xi)\|_{L^q}
 \le C_{\beta,p,q}\delta^{p-1}.
\]
If $g:\D\to\mathbb R$ is continuous and
$|g(z)|\le C d(z)^p$, then, along any subsequence
$\Xi_{N_j}\Rightarrow\Xi$,
\[
 (\Xi_{N_j},\Xi_{N_j}(g))\Rightarrow(\Xi,\Xi(g)),
 \qquad
 \mathbb E e^{t\Xi_{N_j}(g)}\longrightarrow
                  \mathbb E e^{t\Xi(g)}\quad(t\in\mathbb R).
\]
The series defining $\Xi(g)$ is absolutely convergent almost surely.
The assertions extend jointly to any finite collection of such functions.
\end{theorem}

\begin{proof}
Layer integration, justified by Tonelli's theorem, gives
\[
 U_p(\Xi_N)=\int_0^1 p t^{p-1}M_N(1-t)\,dt.
\]
The weight $w(t)=(p-1)t^{p-2}$ integrates to one. Jensen's inequality
therefore yields
\[
 e^{cU_p(\Xi_N)}
 \le\int_0^1 w(t)
    \exp\left\{\frac{cp}{p-1}tM_N(1-t)\right\}\,dt.
\]
Take expectations and apply Lemma~\ref{lem:scaled-count-exponential}.
For the boundary tail, the pointwise inequality
\[
 V_{p,\delta}(\Xi_N)
 \le\int_0^\delta p t^{p-1}M_N(1-t)\,dt
\]
and the probability weight
$w_\delta(t)=(p-1)t^{p-2}/\delta^{p-1}$ on $(0,\delta)$ give the second
bound by the same argument. The scaled-count lemma for the limit and
Tonelli give identical estimates for $\Xi$; one may first truncate the
integrals away from zero and then use monotone convergence. Exponential
moments imply the stated $L^q$ bounds, for example by
$x^q\le C_q e^x$ for $x\ge0$.

Choose continuous cutoffs $\chi_\delta$ with $0\le\chi_\delta\le1$,
equal to one on $\{|z|\le1-\delta\}$ and zero sufficiently close to
the unit circle. Then $g_\delta=\chi_\delta g$ has compact support, and
\[
 |\xi(g)-\xi(g_\delta)|\le C V_{p,\delta}(\xi).
\]
Vague convergence gives joint convergence of the process and its
$g_\delta$ statistic. The uniform $L^1$ tail bound just proved, first
at finite $N$ and then for $\Xi$, permits $\delta\downarrow0$ and proves
the claimed joint convergence. This standard approximation step can
also be verified directly on bounded Lipschitz tests: the replacement
error is bounded by the $L^1$ tail, and the middle term converges for
each fixed cutoff. Bound \eqref{eq:boundary-weight-all-exponentials},
applied with an arbitrarily larger coefficient, makes
$e^{t\Xi_N(g)}$ uniformly integrable. Truncation proves convergence
of its expectation. Finite collections are treated by the same cutoff
and the sum of their individual bounds.
\end{proof}

\begin{corollary}[Explicit limiting changes of potential]
\label{cor:boundary-vanishing-potential-change}
Let $g$ satisfy the preceding theorem, and let $f$ be holomorphic near
$\overline\D$ with $f(0)=0$. Define a perturbed law by
\[
 d\mathsf P_N^{f,g}
 =\frac{\exp\{Y_N(f)+\Xi_N(g)\}}
        {\mathbb E\exp\{Y_N(f)+\Xi_N(g)\}}\,d\mathsf P_{N,\beta},
 \qquad Y_N(f)=2\operatorname{Re}\sum_i f(z_i).
\]
Along every subsequence $\Xi_{N_j}\Rightarrow\Xi$, the normalizer tends to
\[
 \exp\left\{\frac2\beta\sum_{k\ge1}k|f_k|^2\right\}
                   \mathbb E e^{\Xi(g)},
\]
and the interior law under $\mathsf P_{N_j}^{f,g}$ tends to the law
\[
 d\mathsf P^g(\xi)=
       \frac{e^{\xi(g)}}{\mathbb E e^{\Xi(g)}}\,d\mathcal L(\Xi)(\xi).
\]
Its joint analytic Gaussian limit has means
$(2k/\beta)\overline{f_k}$, variances $2k/\beta$, and is independent
of that tilted interior process. For $\beta>4$, the full perturbed
sequence still tends to the empty configuration. In this range,
$U_p(\Xi_N)\to0$ in probability and in every fixed positive moment.
\end{corollary}

\begin{proof}
The analytic Gaussian theorem gives joint convergence of $\Xi_N$ and
$X_N$, with independence in the limit. The cutoff argument in the
preceding proof adjoins $\Xi_N(g)$ to that joint convergence. The limit
of this last statistic is the measurable function $\Xi(g)$ of the
interior process. H\"older's inequality and the separate exponential
bounds for $Y_N(f)$ and $U_p(\Xi_N)$ give uniform integrability of
every fixed positive power of $e^{Y_N(f)+\Xi_N(g)}$.
Passing to expectations against bounded continuous tests proves both
the normalizer formula and the limiting change of measure. Gaussian
completion of the square gives the stated mode means and variances.
Independence follows because the limiting weight factors into a function
of the interior and a function of the independent Gaussian field.

For $\beta>4$, the unweighted full interior limit is empty, so all
subsequences have the same perturbed empty limit by the preceding result.
Alternatively, uniform integrability transfers the vanishing compact
hitting probabilities directly. Apply the weighted-statistic theorem
with $g(z)=d(z)^p$ to see that $U_p(\Xi_N)\Rightarrow0$; the exponential
bound gives convergence of every positive moment, and convergence to
the constant zero in law is convergence in probability.
\end{proof}

The strict inequality $p>1$ is used in the integrability of
$t^{p-2}$ at zero. The theorem gives no estimate for the first-order
boundary field weighted by $1-|z|$, which remains relevant to uniqueness
below and at $4$.

\section{Inserted charges and harmonic screening}
\label{sec:palm-harmonic-screening}

Fix points $a_1,\ldots,a_p\in\D$, exponents $\gamma_1,\ldots,\gamma_p>0$,
and a function $f$ holomorphic near $\overline\D$, with $f(0)=0$.
Repeated points are permitted. Define
\[
 U(z)=2\operatorname{Re}f(z)+\sum_{j=1}^p\gamma_j\log|z-a_j|,
\]
and let $P_N^U$ be the probability law proportional to
\[
 |\Delta_N(z)|^\beta\prod_{i=1}^N e^{U(z_i)}\,dA(z_i)
\]
on $\D^N$. Values at the finitely many zeros of the weight are irrelevant
to this density. The normalizing integral is finite and positive.
This section identifies its harmonic fluctuations without asserting
uniqueness of its interior point-process limit.

\begin{theorem}[Screening of finitely many inserted charges]
\label{thm:inserted-charge-screening}
Under $P_N^U$, the analytic functions
\[
 X_N(u)=\sum_i\frac{uz_i}{1-uz_i}=\sum_{k\geq1}S_{k,N}u^k
\]
converge in law in $\mathcal H(\D)$ to
\[
 X^U(u)=\sum_{k\geq1}(G_k+m_k)u^k,\qquad
 m_k=\frac{2k}{\beta}\overline{f_k}
           -\frac1\beta\sum_{j=1}^p\gamma_j a_j^k.
\]
Here $f(z)=\sum_{k\geq1}f_kz^k$, and the $G_k$ are independent centered
circular complex Gaussians with $\mathbb E|G_k|^2=2k/\beta$.
The laws of the interior point processes under $P_N^U$ are tight.
Along every subsequence on which the interior converges, the displayed
analytic convergence is joint with the interior, and the Gaussian
analytic function is independent of that interior limit.

For every $h(z)=\sum_{k\geq1}h_kz^k$ holomorphic near $\overline\D$,
\begin{align}
 \lim_{N\to\infty}\mathbb E_N^U
       e^{2\operatorname{Re}\sum_i h(z_i)}
 =\exp\bigg\{&\frac2\beta\sum_{k\geq1}k|h_k|^2
       +\frac4\beta\operatorname{Re}\sum_{k\geq1}kh_k\overline{f_k}
       \notag\\[-2pt]
       &-\frac2\beta\sum_{j=1}^p\gamma_j
                               \operatorname{Re}h(a_j)\bigg\}.
 \label{eq:inserted-charge-laplace}
\end{align}
The convergence also holds after multiplication by a bounded continuous
interior functional along an interior-convergent subsequence, with its
limiting expectation as an additional factor.
\end{theorem}

\begin{proof}
We give the details needed to extend the preceding analytic-tilt argument
across the zeros of the present weight.

\emph{Uniform local and macroscopic bounds.}
The ordered-radius comparison remains valid with the same rates $a_k$.
Indeed, in the inner-block scaling argument the extra factors are
products of $|qw_i-a_j|^{\gamma_j}$ and a nonvanishing holomorphic
exponential. Their product with the Vandermonde cross factors is
subharmonic in the complex scaling parameter. One way to verify this
last assertion is to take its logarithm off the finitely many zeros:
it is a sum of positive multiples of logarithms of holomorphic moduli
and a harmonic function. The product is continuous at those zeros and
extends subharmonically there. Its circular means are therefore
nondecreasing, exactly as required by the original comparison proof.
Thus compact counts have the same uniform Gaussian upper tails and
\[
 \mathbb E_N^U\sum_i(-\log|z_i|)=O_\beta(\log(eN)).
\]
The one-point density is subharmonic for the same reason. The disk
submean inequality and the compact count bound therefore bound that
density uniformly on each compact subset of $\D$. This also proves
point-process tightness.

Write $\Gamma=\sum_j\gamma_j$. If $A_N$ denotes the normalizing
expectation of $e^{\sum_iU(z_i)}$ under the unweighted gas, then Jensen's
inequality, rotation invariance, and the angular mean formula give
\begin{align*}
 \log A_N
 &\geq\sum_j\gamma_j\mathbb E_{N,\beta}\sum_i\log|z_i-a_j|\\
 &=\sum_j\gamma_j\mathbb E_{N,\beta}\sum_i
                          \log\max(|z_i|,|a_j|)
 \geq-C_{\beta,\Gamma}\log(eN).
\end{align*}
The expectation of the analytic part vanishes because $f(0)=0$.
Since $|z-a_j|\leq2$ in the disk, the density of $P_N^U$ relative to the
unweighted gas is bounded above by $e^{CN}(eN)^C$ for a fixed constant.
The unweighted macroscopic Fourier concentration estimate consequently
implies $S_{k,N}/N\to0$ under $P_N^U$ for each fixed $k$. As in
Lemma~\ref{lem:analytic-tilt-radial}, geometric truncation then proves
\[
 N^{-1}\sup_{|u|\leq s}|X_N(u)|\to0\qquad(s<1).
\]
Together with the radial bound, these statements imply convergence of
the normalized empirical measure on the closed disk to uniform measure
on its boundary.

\emph{Analytic tightness before identifying the limit.}
The rational transport proof of
Proposition~\ref{prop:analytic-resolvent-tightness} applies directly.
For a parameter disk $|u|\leq s$, choose its fixed cutoff disk to contain
both that disk and all the $a_j$, with a strictly larger radius less
than one. The vector field is identically zero on a neighborhood of
all the inserted charges. The change of $U$ along its time-$1/N$ flow
is therefore uniformly bounded by $C/N$ per particle: on the region
where the vector field is nonzero, all the logarithmic derivatives of
$U$ are bounded; on its zero region the points remain fixed.
Thus the additional log likelihood is $O(1)$, deterministically.
The rational divided-difference estimate, the four-direction tail
argument on the common good event, and Fubini--Cauchy apply without
any other change. They prove tightness of $X_N$ in $\mathcal H(\D)$
under $P_N^U$. In particular all its fixed positive modes are tight.
This argument works with $f+th$ in place of $f$, for any fixed real $t$
and fixed analytic $h$.

\emph{The small transport and the logarithmic singularities.}
Fix $k\geq1$ and $t\in\C$. Put
\[
 r_N=N^{-1/(100k)},\qquad b_N=\frac{2kt}{\beta N},
\]
and use the already constructed patched branched-M\"obius map
$T_N=T_{b_N,r_N}$. Let $P=P_N^U$, and let $Q$ have density $R_N$
relative to $P$, where $R_N$ is the exact change-of-variables likelihood.
Thus $(T_N)_*Q=P$ and $\mathbb E_PR_N=1$.
The uniform displacement is $O(N^{-1}r_N^{1-k})=o(r_N)$.
The compact density bound just proved shows that the event that a point
lies within $4r_N$ of the origin or of one of the inserted charges has
probability $O(r_N^2)$ under $P$. It also has probability $O(r_N^2)$
under $Q$: its image under $T_N$ is contained in the corresponding
union of disks of radius $5r_N$, for all sufficiently large $N$.
Denote the complementary event by $A_N$.

Compact counts are tight under $Q$, since $T_N\to\mathrm{id}$ uniformly
and its pushforward has law $P$. The boundary-distance estimate for
$T_N$ also gives
$\sum_i(1-|z_i|)=O_Q(\log(eN))$.
Every fixed positive mode is tight under $Q$. To see this directly,
write $y_i=T_N(z_i)$ and split particles at a fixed radius $r_0>0$.
Outside that radius $|T_N(z)-z|\leq C_{r_0}/N$, while the number of
points inside is tight and the global displacement tends to zero.
It follows that $|S_{j,N}(y)-S_{j,N}(z)|=O_Q(1)$ for each fixed $j$;
the $y$ mode is tight under $P$.
The empirical measures under $Q$ have the same uniform boundary limit,
again by uniform convergence of $T_N$ to the identity.

On $A_N$, the negative-mode bound used in the original proof applies:
for every fixed $q\geq1$ and any fixed $r_0\in(0,1)$,
\[
 |S_{-q,N}-\overline{S_{q,N}}|
 \leq C_{q,r_0}\sum_i(1-|z_i|)
        +C_qr_N^{-q}\Xi_N(B_{r_0}).
\]
Consequently all the inputs to the deterministic branched-transport
likelihood expansion hold under both $P$ and $Q$. Specifically,
$S_{-q,N}=O(\log(eN)+r_N^{-q})$ in probability for $q<k$,
$S_{\pm q,N}/N\to0$ for $1\leq q\leq2k$, and every lower positive mode
is tight. Products of lower modes, multiplied by $N^{-1}$, are therefore
negligible. The deterministic remainders are
$O(N^{-1}r_N^{-3k}+N^{-1}r_N^{-2k})=o(1)$.
The unweighted part of the log likelihood is thus
\[
 -2\operatorname{Re}(\bar t S_{k,N})
             -\frac{2k}{\beta}|t|^2+o(1)
\]
in probability under both laws. This is exactly the deterministic
expansion already established in the finite-mode proof; it has not been
integrated over the exceptional small disks.

We now compute the extra change $\sum_i[U(T_Nz_i)-U(z_i)]$.
Choose a fixed $r_0<1$ whose disk contains all the inserted charges and
the origin. On $A_N$, for particles in that disk the logarithmic
singularities are at distance at least $4r_N$. The mean value theorem
and the displacement bound make the absolute change of $U$ per such
particle at most $C/(Nr_N^k)=o(1)$. There are only a tight number of
such particles, under either law. Their total contribution is negligible.
Outside that fixed disk, $U$ is smooth with bounded derivatives and
\[
 N(T_N(z)-z)\longrightarrow
 v_{k,t}(z)=\frac2\beta(tz^{1-k}-\bar t z^{1+k})
\]
uniformly. The boundary empirical limit therefore shows that the total
extra change converges in probability, under both laws, to the boundary
average of the directional derivative $dU(z)[v_{k,t}(z)]$.

On $|z|=1$ the vector $v_{k,t}$ is tangent to the boundary, and
\[
 U(z)=2\operatorname{Re}g(z),\qquad
 g(z)=f(z)+\frac12\sum_{j=1}^p\gamma_j\log(1-\overline{a_j}z),
\]
with all logarithms zero at $z=0$. If $g(z)=\sum_{j\geq1}g_jz^j$,
then Fourier orthogonality gives that directional average as
$(4k/\beta)\operatorname{Re}(t g_k)$. Thus
\begin{equation}\label{eq:palm-screening-likelihood}
 \log R_N=-2\operatorname{Re}(\bar tS_{k,N})
       -\frac{2k}{\beta}|t|^2
       +\frac{4k}{\beta}\operatorname{Re}(t g_k)+o(1)
\end{equation}
in probability under both $P$ and $Q$.

\emph{Identification and independence.}
For every fixed positive $j$, the existing mode-shift expansion gives,
on $A_N$,
\[
 S_{j,N}(T_Nz)-S_{j,N}(z)
 =\frac jk\{b_NS_{j-k,N}-\overline b_NS_{j+k,N}\}+o(1).
\]
The negative-mode estimates and the empirical limit imply that this
converges in $Q$-probability to zero if $j\ne k$, and to $2kt/\beta$
if $j=k$. The interior configuration also changes by a quantity tending
to zero in the vague metric, by uniform convergence of $T_N$ and tight
compact counts.

The right side of \eqref{eq:palm-screening-likelihood} is tight under
$Q$, so the exact likelihoods $R_N$ are uniformly integrable under $P$:
\[
 \mathbb E_P[R_N\mathbf1_{\{R_N>K\}}]
       =Q(\log R_N>\log K)\longrightarrow0
\]
uniformly in sufficiently large $N$, as $K\to\infty$.
Let $Y_N$ comprise the interior configuration and any finite collection
of modes other than $k$. Its transported version differs from it by
$o_Q(1)$. Exact change of variables gives, for $F\in C_b$,
\[
 \mathbb E_PF(Y_N)-\mathbb E_P[R_NF(Y_N)]\longrightarrow0.
\]
On any jointly convergent subsequence, denote the limiting $k$th mode
by $Z_k$ and the remaining variables by $Y$. Uniform integrability and
\eqref{eq:palm-screening-likelihood} yield
\[
 \mathbb E[F(Y)e^{-2\operatorname{Re}(\bar t Z_k)}]
 =\exp\left\{\frac{2k}{\beta}|t|^2
           -\frac{4k}{\beta}\operatorname{Re}(tg_k)\right\}\mathbb EF(Y).
\]
The uniqueness theorem for real Gaussian moment-generating functions
identifies $Z_k$ as a circular complex Gaussian of mean
$(2k/\beta)\overline{g_k}$ and variance $2k/\beta$, independent of $Y$.
This holds for each $k$ and every finite collection of the other modes.
It therefore identifies the full finite-dimensional joint law and its
independence from the interior. Analytic tightness, continuity of Taylor
coefficients, and their determination of a holomorphic function prove
the asserted convergence in $\mathcal H(\D)$. Finally,
\[
 \frac{2k}{\beta}\overline{g_k}
 =\frac{2k}{\beta}\overline{f_k}
             -\frac1\beta\sum_j\gamma_j a_j^k,
\]
which is the stated mean.

\emph{Laplace transforms.}
The continuous analytic pairing from
Theorem~\ref{thm:analytic-harmonic-laplace} gives convergence in law of
$2\operatorname{Re}\sum_i h(z_i)$ to the corresponding real Gaussian
pairing of $X^U$. Under its exponential tilt by any fixed real multiple,
the model has the same inserted charges and analytic part $f+th$.
The analytic tightness established above still applies. Thus
Lemma~\ref{lem:tight-exponential-family} proves uniform integrability
of every fixed exponential and allows passage to its expectation.
The Gaussian variance is $(4/\beta)\sum_k k|h_k|^2$, and its mean is
\[
 \frac4\beta\operatorname{Re}\sum_k kh_k\overline{f_k}
       -\frac2\beta\sum_j\gamma_j\operatorname{Re}h(a_j).
\]
This proves \eqref{eq:inserted-charge-laplace}. Joint convergence and
independence, with the same uniform integrability, prove the mixed
assertion.
\end{proof}

\begin{corollary}[Palm screening and relative correlation invariance]
\label{cor:palm-screening-relative-invariance}
Fix distinct $a_1,\ldots,a_p\in\D$. The reduced Palm law of the
$(N+p)$-particle gas at these locations is the $N$-particle law with
$f=0$ and $\gamma_j=\beta$ above. Its $k$th harmonic mode therefore
has limiting mean $-\sum_j a_j^k$ and variance $2k/\beta$.
For every analytic $h$ as above,
\[
 \lim_{N\to\infty}\mathbb E_N^{\mathrm{Palm}(a)}
       e^{2\operatorname{Re}\sum_i h(z_i)}
 =\exp\left\{\frac2\beta\sum_{k\geq1}k|h_k|^2
                     -2\operatorname{Re}\sum_{j=1}^p h(a_j)\right\}.
\]
If $\rho_{p,N}^h$ denotes the $p$-point density of the gas tilted by
$e^{2\operatorname{Re}\sum_i h(z_i)}$, then
\[
 \frac{\rho_{p,N}^h(a_1,\ldots,a_p)}
      {\rho_{p,N}(a_1,\ldots,a_p)}\longrightarrow1.
\]
This is a relative statement and does not require a nonzero limiting
intensity.
\end{corollary}

\begin{proof}
The Palm density follows immediately by fixing the $p$ inserted
coordinates in the Vandermonde density and canceling their mutual
Vandermonde factor in normalization. The first two conclusions are
Theorem~\ref{thm:inserted-charge-screening}. All finite-$N$ denominators
in the last display are positive for a distinct tuple. Exact insertion
and exponential tilting give
\[
 \frac{\rho_{p,N}^h(a)}{\rho_{p,N}(a)}
 =e^{2\operatorname{Re}\sum_jh(a_j)}
   \frac{\mathbb E_{N-p}^{\mathrm{Palm}(a)}
                    e^{2\operatorname{Re}\sum_i h(z_i)}}
        {\mathbb E_{N,\beta}e^{2\operatorname{Re}\sum_i h(z_i)}}.
\]
The Palm Laplace limit and
\eqref{eq:analytic-strong-szego} cancel the two charge terms and the
common Gaussian variance term, proving the limit one.
\end{proof}

\section{The exact M\"obius identity and interior symmetries}
\label{sec:mobius-ward-interior}

For $t\in\C$ let $v_t(z)=t-\overline t z^2$, viewed as a real vector
field. A smooth local cylinder functional means a function
\[
 F(\eta)=\Phi\left(\int h_1\,d\eta,\ldots,\int h_q\,d\eta\right),
\]
where $h_j\in C_c^1(\D;\R)$ and $\Phi\in C_b^1(\R^q;\R)$ has bounded
first derivatives. Define its directional derivative by
\[
 \mathcal L_tF(\eta)
 =\sum_{j=1}^q\partial_j\Phi\left(\int h_1d\eta,\ldots,\int h_qd\eta\right)
       \int dh_j(z)[v_t(z)]\,d\eta(z).
\]

\begin{proposition}[An exact local Ward identity and its first-order limit]
\label{prop:mobius-local-ward}
For every such $F$ and every $t\in\C$,
\begin{equation}\label{eq:mobius-local-ward}
 \mathbb E_{N,\beta}\mathcal L_tF(\Xi_N)
 =2c_N\operatorname{Cov}_{N,\beta}
       \left(F(\Xi_N),\operatorname{Re}(\overline t S_{1,N})\right),
 \qquad c_N=2+\frac\beta2(N-1).
\end{equation}
In particular, the covariance is $O_{\beta,F,t}(N^{-1})$.
If $\Xi_{N_j}\Rightarrow\Xi$ along a subsequence, then
\begin{equation}\label{eq:mobius-first-order-covariance}
 N_j\operatorname{Cov}_{N_j,\beta}
       \left(F(\Xi_{N_j}),\operatorname{Re}(\overline t S_{1,N_j})\right)
 \longrightarrow\frac1\beta\mathbb E\mathcal L_tF(\Xi).
\end{equation}
\end{proposition}

\begin{proof}
Use the exact disk automorphism
$T_a(z)=(z+a)/(1+\overline a z)$ and its likelihood
\[
 R_{N,a}=(1-|a|^2)^{c_NN}\prod_i|1+\overline a z_i|^{-2c_N}.
\]
Its change-of-variables identity is
$\mathbb E[F((T_a)_*\Xi_N)R_{N,a}]=\mathbb EF(\Xi_N)$.
Differentiating at $a=st$, $s=0$, gives
\[
 \left.\frac d{ds}T_{st}(z)\right|_{s=0}=v_t(z),\qquad
 \left.\frac d{ds}R_{N,st}\right|_{s=0}
       =-2c_N\operatorname{Re}(\overline t S_{1,N}).
\]
For each fixed $N$ all derivatives are bounded on a sufficiently small
parameter interval, so differentiation under the integral is legitimate.
Rotation invariance gives $\mathbb ES_{1,N}=0$, proving
\eqref{eq:mobius-local-ward}.

For some compact $K\Subset\D$ and a fixed constant $C$,
$|\mathcal L_tF(\eta)|\leq C\eta(K)$. The compact-count bounds therefore
bound its expectation uniformly in $N$, proving the covariance estimate.
The functional $\mathcal L_tF$ is continuous for the vague topology:
it is a finite sum of products of continuous cylinder functions and
integrals of continuous compactly supported functions. Uniform moments
of the compact count give uniform integrability. Thus
$\mathbb E\mathcal L_tF(\Xi_{N_j})\to\mathbb E\mathcal L_tF(\Xi)$.
Finally, $N/(2c_N)\to1/\beta$, which proves
\eqref{eq:mobius-first-order-covariance}.
\end{proof}

\begin{corollary}[A sufficient symmetry criterion and the resulting intensity]
\label{cor:mobius-invariance-criterion}
If, along an interior-convergent subsequence, the covariances in
\eqref{eq:mobius-first-order-covariance} are $o(N_j^{-1})$ for every
smooth local cylinder $F$ and every $t\in\C$, then the limiting law is
invariant under all conformal automorphisms of $\D$. Its intensity is
necessarily
\[
 \rho_1(z)=\frac{c}{(1-|z|^2)^2},\qquad c\geq0.
\]
The case $c=0$ is exactly the almost surely empty process.
\end{corollary}

\begin{proof}
The proposition gives $\mathbb E\mathcal L_tF(\Xi)=0$ for all the
specified tests. Let $\varphi_s$ be the flow of $v_t$ on the disk.
It consists of disk automorphisms. For bounded $s$-intervals the
functionals $F\circ(\varphi_s)_*$ remain smooth local cylinders whose
supports lie in a common compact set. Their derivatives are dominated
by a fixed constant times its particle count. Hence
\[
 \frac d{ds}\mathbb EF((\varphi_s)_*\Xi)
 =\mathbb E\mathcal L_t(F\circ(\varphi_s)_*)(\Xi)=0.
\]
Thus the law is invariant under each of these flows. Together with the
already established rotation invariance, these flows generate every
disk automorphism: the flow with $t=e^{i\theta}$ sends zero to
$e^{i\theta}\tanh s$, and an automorphism fixing zero is a rotation.
Smooth local cylinders determine the point-process law, for example by
approximating compactly supported continuous Laplace tests, so the
claimed invariance follows.

The limiting intensity is continuous by the earlier correlation theorem.
For $T_a(0)=a$ the derivative has modulus $|T_a'(0)|=1-|a|^2$.
Invariance of the intensity measure and change of variables therefore
imply
$\rho_1(a)(1-|a|^2)^2=\rho_1(0)$, proving the formula with
$c=\rho_1(0)$. If $c=0$, every compact count has expectation zero, and
a countable exhaustion shows that the configuration is empty almost
surely. The converse is immediate.
\end{proof}

\begin{remark}[The scale missing from an argument based only on Gaussian independence]
The proved bound in Proposition~\ref{prop:mobius-local-ward} is
$O(N^{-1})$, whereas the symmetry criterion requires $o(N^{-1})$.
The first-order limit in \eqref{eq:mobius-first-order-covariance} explains
this difference exactly: its value is the interior response to the
M\"obius flow. Leading-order Gaussian independence does not determine
that response.

A fixed nontrivial disk automorphism also cannot be inserted directly
into the fixed-analytic-tilt invariance theorem. Its exact likelihood
contains the exponent $2c_N=4+\beta(N-1)$, which grows with $N$.
The previously established Laplace theorem applies to each fixed
analytic potential. Substituting this growing exponent would require
a new uniform estimate. Neither the symmetry criterion nor that new
estimate is asserted here, and the displayed hyperbolic intensity
formula is a conditional conclusion, not a formula established for
all inverse temperatures.
\end{remark}

\section{An exact Gaussian Hankel representation at \texorpdfstring{$\beta=4$}{beta equals 4}}
\label{sec:critical-hankel}

The endpoint $\beta=4$, outside the scope of the quantitative
boundary-transport theorem, admits the following exact
representation. The endpoint is settled later by the force
argument; the representation here does not, by itself,
prove convergence of adjacent partition ratios or of the point process.

\begin{proposition}[Gaussian Hankel determinant identity]
\label{prop:critical-gaussian-hankel}
Let $a_0,a_1,\ldots$ be independent standard circular complex Gaussian
random variables, normalized by $\mathbb E|a_j|^2=1$. Define the nested
complex symmetric Hankel matrices
\[
 H_N=\left(\frac{a_{i+j}}{\sqrt{i+j+1}}\right)_{0\leq i,j<N}.
\]
Then, with ordinary area in the gas partition function,
\begin{equation}
 \frac{Q_{N,4}}{\pi^N}=\mathbb E|\det H_N|^2.
 \label{eq:critical-gaussian-hankel}
\end{equation}
In particular,
\begin{equation}
 \frac{Q_{N+1,4}}{Q_{N,4}}
 =\pi\frac{\mathbb E|\det H_{N+1}|^2}
             {\mathbb E|\det H_N|^2}.
 \label{eq:critical-hankel-adjacent-ratio}
\end{equation}
\end{proposition}

\begin{proof}
Put $d\mu=dA/\pi$ on $\mathbb D$. The matrix-entry covariances are
\[
 \mathbb E[H_{ij}\overline{H_{kl}}]
 =\frac{\mathbf1_{\{i+j=k+l\}}}{i+j+1}
 =\int_{\mathbb D}z^{i+j}\overline z^{\,k+l}\,d\mu(z),
 \qquad \mathbb E[H_{ij}H_{kl}]=0.
\]
The complex Gaussian Wick formula says that a product of $N$ unbarred
entries and $N$ barred entries has expectation equal to the sum, over
all bijections between them, of the products of their covariances.
For clarity, this formula follows by differentiating the identity
\[
 \mathbb E\exp\left(\sum_j u_ja_j+\sum_jv_j\overline a_j\right)
 =\exp\left(\sum_j u_jv_j\right),
\]
first for finitely many independent Gaussians and then for their linear
combinations. All expressions used here involve finitely many variables.

Index permutations of $\{0,\ldots,N-1\}$ by $\sigma,\tau,\pi$.
Expanding the two determinants and applying Wick gives
\begin{align*}
 \mathbb E|\det H_N|^2
 &=\sum_{\sigma,\tau,\pi\in S_N}
       \operatorname{sgn}(\sigma)\operatorname{sgn}(\tau)
       \prod_{i=0}^{N-1}
       \int_{\mathbb D}z_i^{i+\sigma(i)}
                \overline z_i^{\,\pi(i)+\tau(\pi(i))}\,d\mu(z_i).
\end{align*}
Set $\kappa=\tau\circ\pi$, so that
$\operatorname{sgn}(\tau)=\operatorname{sgn}(\kappa)
\operatorname{sgn}(\pi)$. Write
\[
 V_N(z)=\det[z_i^j]_{0\leq i,j<N},\qquad |V_N(z)|=|\Delta_N(z)|.
\]
Collecting the three permutation sums inside the product integral yields
\begin{equation}
 \mathbb E|\det H_N|^2
 =\int_{\mathbb D^N}
       \left(\prod_{i=0}^{N-1}z_i^i\right)
       V_N(z)\overline{V_N(z)}^{\,2}
                \,d\mu^{\otimes N}(z).
 \label{eq:critical-hankel-presymmetrization}
\end{equation}
Finally symmetrize this integral over all permutations of the variables.
The factor $\overline{V_N}^{\,2}$ is symmetric, while
$V_N$ changes by the sign of the permutation. Moreover,
\[
 \sum_{\lambda\in S_N}\operatorname{sgn}(\lambda)
           \prod_{i=0}^{N-1}z_{\lambda(i)}^i=V_N(z).
\]
Thus \eqref{eq:critical-hankel-presymmetrization} equals
\[
 \frac1{N!}\int_{\mathbb D^N}|\Delta_N(z)|^4\,d\mu^{\otimes N}(z)
 =\frac{Z_{N,4}}{N!\pi^N}=\frac{Q_{N,4}}{\pi^N}.
\]
This proves \eqref{eq:critical-gaussian-hankel}, and division proves
\eqref{eq:critical-hankel-adjacent-ratio}.
\end{proof}

\paragraph{Normalization check.}
For $N=2$,
\[
 \det H_2=\frac{a_0a_2}{\sqrt3}-\frac{a_1^2}{2},
 \qquad
 \mathbb E|\det H_2|^2=\frac13+\frac{\mathbb E|a_1|^4}{4}
 =\frac13+\frac12=\frac56.
\]
This agrees with the direct disk integral $Q_{2,4}=5\pi^2/6$.

\paragraph{What remains to be proved by this route.}
The matrices $H_N$ are nested, but their determinant moments are not
ordinary leading principal minors of a deterministic positive matrix.
Consequently the Hadamard--Fischer determinant inequality does not imply
log-concavity of the sequence in \eqref{eq:critical-gaussian-hankel}.
A convergence or monotonicity theorem for the adjacent determinant-moment
ratios in \eqref{eq:critical-hankel-adjacent-ratio} is still needed.
Even such a ratio theorem would have to be combined with convergence of
the local insertion functions to establish uniqueness of the interior
point process at $\beta=4$ by this determinant route. The independent
score proof below does not need either adjacent-ratio convergence
or such an insertion limit.

\begin{proposition}[Explicit exponential bounds at $\beta=4$]
\label{prop:critical-partition-exponential-bounds}
For every $N\geq1$,
\begin{equation}
 \pi^N\frac{(2N-1)!!}{N^N}\leq Q_{N,4}\leq\pi^N.
 \label{eq:critical-partition-exponential-bounds}
\end{equation}
Consequently the radius of convergence $\lambda_c$ of
$\sum_{N\geq0}Q_{N,4}\lambda^N$ satisfies
\begin{equation}
 \frac1\pi\leq\lambda_c\leq\frac e{2\pi}.
 \label{eq:critical-grandcanonical-radius-bounds}
\end{equation}
These bounds do not assert convergence of the adjacent ratios.
\end{proposition}

\begin{proof}
We first compute the circle integral needed in the argument, without
using a Selberg or Dyson integral formula. On the unit circle,
\[
 |\Delta_N(z)|^4
 =\frac{\Delta_N(z)^4}{\prod_i z_i^{2N-2}}.
\]
Its integral under normalized product angular measure is therefore the
coefficient of $\prod_i z_i^{2N-2}$ in $\Delta_N^4$.
The confluent Vandermonde identity expresses $\Delta_N^4$ as the
determinant with two rows for each $z_i$, namely
\[
 (1,z_i,\ldots,z_i^{2N-1}),\qquad
 (0,1,2z_i,\ldots,(2N-1)z_i^{2N-2}).
\]
This identity follows from the ordinary Vandermonde determinant on
$(z_1,z_1+h_1,\ldots,z_N,z_N+h_N)$ by subtracting each first row
from its following row, dividing by $\prod_i h_i$, and letting all
$h_i\to0$.

Expand the confluent determinant in its two-row blocks. Choosing columns
$p<q$, numbered from $0$ to $2N-1$, in one block contributes
$(q-p)z_i^{p+q-1}$. The required exponent forces $p+q=2N-1$.
Thus the chosen column pairs must be exactly
$\{p,2N-1-p\}$, $0\leq p<N$, assigned in $N!$ ways to the row blocks.
All their signs are positive: in increasing order of $p$, the column
shuffle has $N(N-1)$ inversions, an even number, and exchanging any two
two-column blocks preserves the sign. Their common product is
$\prod_{p=0}^{N-1}(2N-1-2p)=(2N-1)!!$. We have proved
\begin{equation}
 \int_{\mathbb T^N}|\Delta_N(z)|^4\,dm^{\otimes N}(z)
 =N!(2N-1)!!.
 \label{eq:critical-circle-elementary-integral}
\end{equation}

Let $\widetilde H_N=(a_{i+j})_{0\leq i,j<N}$. The Wick-and-symmetrization
proof of Proposition~\ref{prop:critical-gaussian-hankel} applies verbatim
with normalized circle measure in place of normalized area, because
$\mathbb E[a_{i+j}\overline a_{k+l}]=\mathbf1_{\{i+j=k+l\}}$
is now the corresponding monomial covariance. Hence
\eqref{eq:critical-circle-elementary-integral} gives
\[
 \mathbb E|\det\widetilde H_N|^2=(2N-1)!!.
\]
Group the determinant polynomial into Gaussian monomials,
\[
 \det\widetilde H_N=\sum_{\mathbf m}b_{\mathbf m}
                         \prod_{s=0}^{2N-2}a_s^{m_s}.
\]
Distinct monomials are orthogonal under complex Gaussian measure, with
squared norm $\prod_s m_s!$. Replacing $a_s$ by $a_s/\sqrt{s+1}$
therefore gives
\[
 \mathbb E|\det H_N|^2
 =\sum_{\mathbf m}|b_{\mathbf m}|^2
             \frac{\prod_s m_s!}{\prod_s(s+1)^{m_s}},
 \qquad
 \sum_{\mathbf m}|b_{\mathbf m}|^2\prod_s m_s!=(2N-1)!!.
\]
Every nonzero monomial arises from at least one determinant permutation
$\sigma$. Its denominator is
$\prod_s(s+1)^{m_s}=\prod_{i=0}^{N-1}(i+\sigma(i)+1)$.
The sum of these $N$ factors is $N^2$, so the arithmetic--geometric mean
inequality bounds their product above by $N^N$.
For its lower bound, whenever $i<j$ and $\sigma(i)=b>\sigma(j)=a$,
replacing this inversion by the ordered pair changes the two-factor
product from $(i+b+1)(j+a+1)$ to $(i+a+1)(j+b+1)$, a decrease of
$(j-i)(b-a)>0$. Removing all inversions shows that the minimum is attained
at the identity permutation and equals $\prod_i(2i+1)=(2N-1)!!$.
Thus every denominator in the nonnegative sum lies between
$(2N-1)!!$ and $N^N$. The last two displays and
\eqref{eq:critical-gaussian-hankel} prove
\eqref{eq:critical-partition-exponential-bounds}.

Finally, $(2N-1)!!=(2N)!/(2^NN!)$ and Stirling's formula give
$((2N-1)!!/N^N)^{1/N}\to2/e$.
The root test for the power series proves
\eqref{eq:critical-grandcanonical-radius-bounds}.
\end{proof}

\section{An exact outermost-particle transfer and the remaining covariance estimate}
\label{sec:uniqueness-transfer}

Throughout this section $\beta>0$ is fixed, $Z_N=Z_{N,\beta}$, and
$\E_N$ denotes expectation under the unweighted disk gas with $N$
particles. We use the previously proved tightness and the compact-count
estimate
\begin{equation}\label{transfer:count-input}
 B_{\beta,r}:=\sup_{N\ge1}\E_N[M_N(r)^2]<\infty,
 \qquad 0<r<1.
\end{equation}
For completeness, an explicit permissible value of $B_{\beta,r}$ is
\[
 \sum_{k=1}^{\infty}(2k-1)b_k,
 \qquad b_1=1,\quad
 b_k=\min\left\{1,
 \exp\left(k-\frac\beta4k(k-1)\log\frac1r\right)\right\}
 \quad(k\ge2).
\]
This follows from $\E X^2=\sum_{k\ge1}(2k-1)\Pbb(X\ge k)$ for
nonnegative integer-valued $X$ and the established count-tail bound.

The following transfer formula is exact. It is important that the
remaining particles acquire a nonconstant weight after deletion: their
law is not asserted to be the original $(N-1)$-particle law.

\begin{proposition}[Deletion and rescaling of the outermost particle]
\label{prop:outermost-transfer}
For $N\ge2$, let
\[
 R_N=\max_{1\le j\le N}|z_j|,
 \qquad
 a_N=2N+\frac\beta2N(N-1).
\]
Almost surely there is a unique maximizing index $J_N$, and $R_N>0$.
Write $z_{J_N}=R_Ne^{i\Theta_N}$, with
$\Theta_N\in[0,2\pi)$, and let $W_N=(w_1,\ldots,w_{N-1})$ be the
remaining labeled particles divided by $R_N$, retaining their relative
order of labels. Define the point configuration
\[
 \widehat\Xi_N=\sum_{j=1}^{N-1}\delta_{w_j}
\]
and, for $n\ge1$,
\begin{equation}\label{transfer:boundary-weight}
 G_n(w_1,\ldots,w_n)
 =\frac1{2\pi}\int_0^{2\pi}
   \prod_{j=1}^{n}|e^{i\theta}-w_j|^\beta\,d\theta.
\end{equation}
Then $R_N$ is independent of the pair $(W_N,\Theta_N)$ and has density
$a_N r^{a_N-1}\ind_{(0,1)}(r)$. Moreover,
\begin{equation}\label{transfer:tilted-law}
 \Law(W_N)(dw)
 =\frac{G_{N-1}(w)}{\E_{N-1}G_{N-1}}
     \Pbb_{N-1,\beta}(dw),
\end{equation}
and
\begin{equation}\label{transfer:partition-identity}
 Z_N=\frac{2\pi N}{a_N}\,Z_{N-1}\,\E_{N-1}G_{N-1}.
\end{equation}
The same tilted-law formula holds for the corresponding unlabeled point
configurations.
\end{proposition}

\begin{proof}
The original law is absolutely continuous with respect to area measure
on $\D^N$. Each set $\{|z_i|=|z_j|\}$, $i\ne j$, and the set on
which all particles are zero have area measure zero. Thus $J_N$ is
unique and $R_N>0$ almost surely.

On the event that label $N$ is the maximizing label, make the change of
variables
\[
 z_N=re^{i\theta},\qquad z_j=rw_j\quad(1\le j\le N-1),
 \qquad 0<r<1,\quad w\in\D^{N-1}.
\]
Its area Jacobian is $r^{2(N-1)+1}$. One way to see this without
ignoring the dependence on $r$ is to order the coordinates as
$(z_N,z_1,\ldots,z_{N-1})$: the derivative matrix is block triangular,
with polar Jacobian $r$ for $z_N$ and two-dimensional dilation
Jacobian $r^2$ for each remaining coordinate. Also,
\[
 |\Delta_N(z)|^\beta
 =r^{\beta N(N-1)/2}|\Delta_{N-1}(w)|^\beta
   \prod_{j=1}^{N-1}|e^{i\theta}-w_j|^\beta.
\]
Consequently, for every bounded measurable function $H$ of
$(r,w,\theta)$, exchangeability and summation over the possible
maximizing labels give
\begin{align}
 \E_N H(R_N,W_N,\Theta_N)
 &=\frac{N}{Z_N}\int_0^1 r^{a_N-1}\,dr
   \int_{\D^{N-1}}|\Delta_{N-1}(w)|^\beta
       \prod_{j=1}^{N-1}dA(w_j)\notag\\
 &\hspace{12mm}\times
   \int_0^{2\pi}\prod_{j=1}^{N-1}|e^{i\theta}-w_j|^\beta
       H(r,w,\theta)\,d\theta.
 \label{transfer:joint-integral}
\end{align}
There is no additional ordering restriction on $w$: the maximizing
particle was selected by its radius, while the remaining original
labels are still exchangeable.

Taking $H=1$ and using $\int_0^1r^{a_N-1}\,dr=a_N^{-1}$ proves
\eqref{transfer:partition-identity}. The factorization in
\eqref{transfer:joint-integral} proves the asserted independence and
radial density. Integrating over $\theta$ proves
\eqref{transfer:tilted-law}. The denominator is finite and strictly
positive: $G_n\le2^{\beta n}$, and its integrand is positive whenever
all $w_j\in\D$. Pushing the labeled law forward to its counting
measure proves the final assertion.
\end{proof}

We next quantify the difference between the rescaled remaining
configuration and the original configuration on a fixed compact set.
For a nonnegative $f\in C_c^1(\D)$, write
\[
 F_f(\xi)=\exp[-\xi(f)],\qquad L_f=\|\nabla f\|_\infty,
\]
where $f$ is extended by zero to $\C$. This extension is $C^1$ because
$f$ has compact support in $\D$.

\begin{proposition}[A summable error for compact Laplace functionals]
\label{prop:transfer-error}
Choose $0<r<r_0<1$ such that $\supp f\subset B_r$. For every $N\ge2$,
\begin{align}
 \left|\E_N F_f(\Xi_N)-
 \frac{\E_{N-1}[G_{N-1}F_f(\Xi_{N-1})]}
      {\E_{N-1}G_{N-1}}\right|
 &\le r_0^{a_N}
  +\frac{L_fr}{r_0}\sqrt{B_{\beta,r}}
       \sqrt{\frac{2}{(a_N+1)(a_N+2)}}.
 \label{transfer:error-bound}
\end{align}
In particular, the left-hand side is $O_{\beta,f,r,r_0}(N^{-2})$.
\end{proposition}

\begin{proof}
Use the coupling in Proposition~\ref{prop:outermost-transfer}.
On $\{R_N\ge r_0\}$, the deleted particle lies outside $\supp f$.
For every remaining particle, if $|z_j|\ge r$, both
$f(z_j)$ and $f(z_j/R_N)$ vanish. If $|z_j|<r$, the global Lipschitz
bound on $f$ gives
\[
 |f(z_j/R_N)-f(z_j)|
 \le L_f|z_j|\left(\frac1{R_N}-1\right)
 \le\frac{L_fr}{r_0}(1-R_N).
\]
Therefore
\[
 |\widehat\Xi_N(f)-\Xi_N(f)|
 \le\frac{L_fr}{r_0}(1-R_N)M_N(r)
 \qquad\hbox{on }\{R_N\ge r_0\}.
\]
The function $s\mapsto e^{-s}$ is $1$-Lipschitz on $[0,\infty)$.
On the complementary event the two Laplace functionals differ by at
most $1$. It follows that
\begin{align*}
 |\E_NF_f(\widehat\Xi_N)-\E_NF_f(\Xi_N)|
 &\le \Pbb_N(R_N<r_0)
    +\frac{L_fr}{r_0}\E_N[(1-R_N)M_N(r)]\\
 &\le r_0^{a_N}
    +\frac{L_fr}{r_0}\sqrt{B_{\beta,r}}
          \sqrt{\E_N(1-R_N)^2}.
\end{align*}
Here the second step uses Cauchy--Schwarz, not independence between
$R_N$ and $M_N(r)$; such independence is not being asserted.
Direct integration of the radial density gives
\[
 \E_N(1-R_N)^2
 =a_N\int_0^1(1-s)^2s^{a_N-1}\,ds
 =\frac{2}{(a_N+1)(a_N+2)}.
\]
Equation~\eqref{transfer:tilted-law} now proves
\eqref{transfer:error-bound}. Finally, for $N\ge2$,
$a_N\ge\beta N^2/4$. Thus the second term in
\eqref{transfer:error-bound} is $O(N^{-2})$, while its first term
decays at least as fast as
$\exp[-(\beta/4)\log(1/r_0)N^2]$.
\end{proof}

The preceding two propositions identify an explicit estimate that
would suffice to prove uniqueness at any fixed temperature where it holds.

\begin{proposition}[A one-sided covariance criterion for uniqueness]
\label{prop:transfer-uniqueness-criterion}
For $n\ge1$ and nonnegative $f\in C_c^1(\D)$, put
\[
 c_n(f)=
 \frac{\Cov_n(F_f(\Xi_n),G_n)}{\E_nG_n},
 \qquad x_-:=\max\{-x,0\}.
\]
Suppose that, for every nonnegative $f\in C_c^1(\D)$,
\begin{equation}\label{transfer:missing-covariance-bound}
 \sum_{n=1}^{\infty} [c_n(f)]_-<\infty.
\end{equation}
Then the entire sequence $\Xi_N$ converges in distribution in
$\Mp(\D)$ to a unique law. The same conclusion follows if
\eqref{transfer:missing-covariance-bound} is verified only on a
countable family of nonnegative smooth compactly supported functions
whose Laplace functionals determine laws on $\Mp(\D)$.

For each test function satisfying
\eqref{transfer:missing-covariance-bound}, the sequence
$\E_NF_f(\Xi_N)$ in fact has finite total variation, and
$\sum_n|c_n(f)|<\infty$.
\end{proposition}

\begin{proof}
Fix $f$ and put $x_n=\E_nF_f(\Xi_n)\in[0,1]$. By
Proposition~\ref{prop:transfer-error},
\begin{equation}\label{transfer:increment-identity}
 x_{n+1}-x_n=c_n(f)+\varepsilon_n(f),
 \qquad \sum_{n=1}^{\infty}|\varepsilon_n(f)|<\infty.
\end{equation}
Consequently the sum of the negative parts of the increments is finite:
\[
 \sum_{n\ge1}(x_n-x_{n+1})_+
 \le \sum_{n\ge1}[c_n(f)]_-
      +\sum_{n\ge1}|\varepsilon_n(f)|<\infty.
\]
Writing $p_n=(x_{n+1}-x_n)_+$ and
$d_n=(x_n-x_{n+1})_+$, telescoping gives
\[
 \sum_{n=1}^m p_n
 =x_{m+1}-x_1+\sum_{n=1}^m d_n
 \le1+\sum_{n\ge1}d_n.
\]
Thus $\sum_n|x_{n+1}-x_n|<\infty$, so $x_n$ converges. Equation
\eqref{transfer:increment-identity} also gives
$\sum_n|c_n(f)|<\infty$. The one-sided form of the hypothesis may
therefore be convenient to establish, but is not claimed here to be
strictly weaker than absolute covariance summability in the presence
of \eqref{transfer:increment-identity} and $0\le x_n\le1$.

By the already established tightness, every subsequence of particle
numbers tending to infinity has a further subsequence with a weak
limit. The function $F_f$ is bounded and vaguely continuous. Hence
all such limits have the same Laplace functional at every test for
which $x_n$ converges.

We recall why these functionals determine the limiting law. Choose a
countable collection $(g_j)_{j\ge1}\subset C_c^\infty(\D)$ of
nonnegative functions separating nonnegative Radon measures; it can
be obtained from countable uniform approximation families on a compact
exhaustion of $\D$. Include all finite sums
$\sum_j m_jg_j$ with nonnegative integer coefficients. Equality of
Laplace functionals for these sums gives equality of all mixed moments
of every finite vector
\[
 (e^{-\xi(g_1)},\ldots,e^{-\xi(g_k)})\in[0,1]^k.
\]
Polynomials are uniformly dense in $C([0,1]^k)$, so those moments
determine each finite-dimensional distribution. The coordinates
$\xi\mapsto\xi(g_j)$ form a countable measure-determining family
for the Borel sigma field of the vague topology on locally finite
Radon measures. Thus equality of their joint distributions determines
the law of $\xi$. This also explains the stated countable-family
version of the criterion.

It follows that all weak subsequential limits have the same law.
Tightness then implies convergence of the full sequence: otherwise a
subsequence outside some neighborhood of that law would have a
further weakly convergent subsequence, giving a contradiction.
\end{proof}

\begin{remark}[What remains unproved]
\label{rem:transfer-gap}
The transfer formula and the $O(N^{-2})$ comparison error above are
unconditional consequences of the finite-particle density and the
proved compact-count bounds. In contrast,
\eqref{transfer:missing-covariance-bound} has not been established
here for general $\beta$. In particular, neither independence of
$R_N$ and $\widehat\Xi_N$ nor positivity of the boundary weight
$G_n$ implies the required covariance estimate. No uniqueness theorem
for general $\beta$ is being deduced from this conditional criterion.
Theorems~\ref{thm:empty-limit-beta-above-four} and~\ref{thm:critical-empty}
prove uniqueness for $\beta\ge4$ without verifying that covariance estimate.
If the criterion is eventually verified, the previously proved
subsequential local total variation result will also upgrade to the
full sequence: a subsequence violating that conclusion would, by
tightness, have a further weakly convergent subsequence, whose limit
must be the unique limit and along which local total variation
convergence is already known.
\end{remark}


\section{All positive harmonic modes can miss an exterior field}
\label{sec:boundary-memory}

The following deterministic example isolates a difficulty in identifying
the Gibbs conditionals of a limit. It is not a construction of different
subsequential limits of the gas.

\begin{proposition}[Boundary information beyond the global harmonic modes]
\label{prop:boundary-memory}
Fix an integer $k\ge1$ and a real number $b$. There exist configurations
$\eta_n=\sum_{j=0}^{n-1}\delta_{w_{j,n}}$ of distinct points in $\D$ such that
\begin{align*}
 \eta_n&\longrightarrow0\quad\hbox{vaguely on }\D,\\
 \frac1n\eta_n&\Longrightarrow\frac{d\theta}{2\pi}
       \quad\hbox{weakly on }\overline\D,\\
 \sum_{j=0}^{n-1}w_{j,n}^{\ell}&\longrightarrow0
       \quad\hbox{for every fixed integer }\ell\ge1,\\
 \sum_{j=0}^{n-1}(-\log|w_{j,n}|)^p&\longrightarrow0
       \quad\hbox{for every real }p>1,
\end{align*}
whereas, locally uniformly for $z\in\D$,
\begin{equation}\label{eq:boundary-memory-field}
 \sum_{j=0}^{n-1}\log\left|1-\frac{z}{w_{j,n}}\right|
 \longrightarrow b\operatorname{Re}(z^k).
\end{equation}
Consequently, for every $0<R<1$, the probability densities on $B_R$
proportional to $\prod_j|z-w_{j,n}|^\beta$ converge uniformly to
\begin{equation}\label{eq:boundary-memory-density}
 \frac{\exp\{\beta b\operatorname{Re}(z^k)\}}
 {\int_{B_R}\exp\{\beta b\operatorname{Re}(u^k)\}\,dA(u)}.
\end{equation}
These limiting densities vary with $b$, although all the limits in the
first display are unchanged.
\end{proposition}

\begin{proof}
Choose a constant $c>|b|$. For all sufficiently large $n$, put
\[
 \theta_j=\frac{2\pi j}{n},\qquad
 w_{j,n}=e^{-c/n}e^{i\theta_j}
                   \exp\left(\frac b n e^{ik\theta_j}\right).
\]
Their moduli and arguments satisfy
\[
 -\log|w_{j,n}|=\frac{c-b\cos(k\theta_j)}n>0,
 \qquad
 \arg w_{j,n}=\theta_j+\frac b n\sin(k\theta_j)\pmod{2\pi}.
\]
The latter circle map is strictly increasing when $n>|b|k$, proving
distinctness. The moduli tend uniformly to one and the angular
perturbations tend uniformly to zero. This proves both the vague limit
and the weak empirical limit. It also gives
\[
 \sum_j(-\log|w_{j,n}|)^p
 \le(c+|b|)^p n^{1-p}\longrightarrow0\qquad(p>1).
\]

Fix a positive integer $\ell$. Taylor expansion with a uniform error
gives
\[
 w_{j,n}^{\ell}
 =e^{i\ell\theta_j}
   \left(1-\frac{\ell c}n
          +\frac{\ell b}n e^{ik\theta_j}+O_{\ell,b,c}(n^{-2})\right).
\]
For $n>\ell+k$, the sums of both $e^{i\ell\theta_j}$ and
$e^{i(\ell+k)\theta_j}$ vanish. Summing the display therefore gives
$\sum_jw_{j,n}^{\ell}=O_{\ell,b,c}(n^{-1})$. This proves the claim
simultaneously for each fixed positive harmonic mode, without requiring
any relation between $\ell$ and a growing particle number.

Fix $R<1$. For $|z|\le R$, write
\[
 v_j=z e^{-i\theta_j},\qquad
 \delta_j=\frac{c-b e^{ik\theta_j}}n,
 \qquad \frac z{w_{j,n}}=v_j e^{\delta_j}.
\]
Choose $R'$ strictly between $R$ and $1$. For all sufficiently large $n$,
$|v_j e^{s\delta_j}|\le R'$ for $0\le s\le1$. The branch of
$\log(1-v)$ that is zero at $v=0$ is consequently defined throughout
the expansion. Its second derivative in $\delta$ is bounded uniformly
there. Taylor's theorem in the complex parameter gives
\[
 \log(1-v_j e^{\delta_j})
 =\log(1-v_j)-\delta_j\frac{v_j}{1-v_j}
                  +O_{R,b,c}(n^{-2}).
\]
The product over roots of unity, with equality of branches fixed at $z=0$,
gives $\sum_j\log(1-v_j)=\log(1-z^n)$. Expanding
$v_j/(1-v_j)=\sum_{h\ge1}z^h e^{-ih\theta_j}$ and summing characters gives,
for $n>k$,
\[
 \frac1n\sum_j\frac{v_j}{1-v_j}=\frac{z^n}{1-z^n},\qquad
 \frac1n\sum_j e^{ik\theta_j}\frac{v_j}{1-v_j}
       =\frac{z^k}{1-z^n}.
\]
Thus, uniformly for $|z|\le R$,
\[
 \sum_j\log(1-z/w_{j,n})
 =\log(1-z^n)-c\frac{z^n}{1-z^n}
                 +b\frac{z^k}{1-z^n}+O_{R,b,c}(n^{-1}).
\]
Taking real parts and letting $n\to\infty$ proves
\eqref{eq:boundary-memory-field}.

Finally, write
\[
 \prod_j|z-w_{j,n}|^\beta
 =\left(\prod_j|w_{j,n}|^\beta\right)
   \exp\left\{\beta\sum_j\log|1-z/w_{j,n}|\right\}.
\]
The first factor cancels upon normalization. Uniform convergence on
$\overline{B_R}$ proves \eqref{eq:boundary-memory-density}.
The limiting densities for distinct $b$ differ, since
$\operatorname{Re}(z^k)$ is not constant on $B_R$.
\end{proof}

For a tagged particle conditioned to lie in $B_R$, with all remaining
particles fixed at $w_{j,n}$, the finite-gas conditional density is exactly
the density used in the proposition. These boundary configurations are
not asserted to be typical for the remaining-particle marginal of the
$(n+1)$-particle gas. The proposition therefore does not disprove
uniqueness. It proves that the displayed limiting data, even including
\emph{every fixed positive harmonic sum}, do not determine the exterior
logarithmic field. A uniqueness argument using conditional Gibbs laws
still has to control that field under the actual gas probabilities.


\section{An exact hyperbolic reformulation and a high-temperature obstruction}
\label{sec:hyperbolic-reduction}

This section records an exact reformulation.  It does not assert uniqueness
of an infinite-volume Gibbs measure.  In particular, the elementary
integrable-interaction criterion for a high-temperature cluster expansion
does not apply to the interaction obtained below.

Put
\[
 g(z,w)=\log\frac{|1-z\overline w|}{|z-w|},\qquad
 d\lambda(z)=\frac{dA(z)}{(1-|z|^2)^2}.
\]
Here $g$ is the nonnegative Green function of the disk normalized by
$-\Delta_z g(z,w)=2\pi\delta_w$, with zero boundary values and the
convention $g(z,z)=+\infty$.  Indeed $\log|1-z\overline w|$ is harmonic
on the disk, while $\Delta_z\log|z-w|=2\pi\delta_w$; the latter
distributional identity follows by integrating the normal derivative
$1/\varepsilon$ on a circle of radius $\varepsilon$ around $w$ and then
letting $\varepsilon\downarrow0$.  The identity
\[
 |1-z\overline w|^2-|z-w|^2=(1-|z|^2)(1-|w|^2)
\]
proves nonnegativity and the zero boundary values.  Both $g$ and
$\lambda$ are invariant under disk automorphisms.  Our measure $\lambda$
is one quarter of the area measure for the curvature $-1$ metric
$4|dz|^2/(1-|z|^2)^2$.

\begin{proposition}[Exact decomposition]\label{prop:hyperbolic-exact}
For any distinct $z_1,\ldots,z_N\in\mathbb D$, put
$S_k=\sum_{j=1}^N z_j^k$.  Then
\begin{align}
 |\Delta_N(z)|^\beta\prod_{j=1}^N dA(z_j)
 ={}&\exp\left\{-\beta\sum_{i<j}g(z_i,z_j)
       -\frac\beta2\sum_{k=1}^{\infty}\frac{|S_k|^2}{k}\right\}
 \prod_{j=1}^N(1-|z_j|^2)^{2-\beta/2}d\lambda(z_j).
 \label{eq:hyperbolic-exact}
\end{align}
All series in this identity converge absolutely.
\end{proposition}

\begin{proof}
Writing $r_* =\max_j|z_j|<1$ gives $|S_k|\le Nr_*^k$,
which proves the convergence assertion.  The absolutely convergent Taylor
series for $\log(1-\zeta)$ gives
\begin{align*}
 \sum_{i<j}\log|1-z_i\overline z_j|
 &=-\frac12\sum_{k\ge1}\frac1k
   \left(|S_k|^2-\sum_i|z_i|^{2k}\right)\\
 &=-\frac12\sum_{k\ge1}\frac{|S_k|^2}{k}
   -\frac12\sum_i\log(1-|z_i|^2).
\end{align*}
Use $\log|z_i-z_j|=-g(z_i,z_j)+\log|1-z_i\overline z_j|$
and $dA(z)=(1-|z|^2)^2d\lambda(z)$.
\end{proof}

\begin{proposition}[Failure of absolute Mayer integrability]
\label{prop:hyperbolic-mayer}
For every $\beta>0$ and every $z\in\mathbb D$,
\[
 \int_{\mathbb D}\left(1-e^{-\beta g(z,w)}\right)d\lambda(w)=+\infty,
 \qquad
 \int_{\mathbb D}\left(1-e^{-\beta g(z,w)}\right)^2d\lambda(w)<\infty.
\]
The second integral is independent of $z$ and equals
\[
 \pi\int_0^1\frac{(1-t^{\beta/2})^2}{(1-t)^2}\,dt.
\]
\end{proposition}

\begin{proof}
For $a\in\mathbb D$, the map
$\varphi_a(w)=(w-a)/(1-\overline a w)$ satisfies
\[
 1-|\varphi_a(w)|^2
   =\frac{(1-|a|^2)(1-|w|^2)}{|1-\overline a w|^2},\qquad
 |\varphi_a'(w)|
   =\frac{1-|a|^2}{|1-\overline a w|^2}.
\]
These identities prove invariance of $\lambda$; the corresponding
pseudohyperbolic-distance identity proves invariance of $g$.  Thus it
suffices to set $z=0$, in which case $e^{-\beta g(0,w)}=|w|^\beta$.
Polar coordinates and $t=|w|^2$ reduce the first integral to
\[
 \pi\int_0^1\frac{1-t^{\beta/2}}{(1-t)^2}\,dt.
\]
Since $(1-t^{\beta/2})/(1-t)\to\beta/2$ as $t\uparrow1$, this
integral diverges logarithmically.  Squaring the numerator instead gives
a bounded integrand near $1$, and an integrand tending to $1$ near $0$.
This proves all assertions.
\end{proof}

There is a separate issue in turning \eqref{eq:hyperbolic-exact} into a
limiting specification: the positive harmonic moments $S_k$ do not by
themselves determine the centered Green field created by particles
approaching the boundary.  The following identity locates this issue
without making any limiting interchange.

\begin{lemma}[The centered exterior Green field]
\label{lem:centered-exterior-green}
Let $0<r<R<1$.  For $|z|\le r$ and $R<|w|<1$,
\[
 g(z,w)-g(0,w)
  =\operatorname{Re}\sum_{k=1}^{\infty}
       \frac{z^k}{k}\left(w^{-k}-\overline w^{\,k}\right).
\]
The series converges uniformly in these $z,w$.  In particular, for any
finite configuration $\xi$ supported on $\{|w|>R\}$,
\[
 \sum_{w\in\xi}\bigl(g(z,w)-g(0,w)\bigr)
 =\operatorname{Re}\sum_{k=1}^{\infty}\frac{z^k}{k}
     \sum_{w\in\xi}(w^{-k}-\overline w^{\,k}).
\]
Furthermore,
\[
 |g(z,w)-g(0,w)|\le C_{r,R}(1-|w|).
\]
\end{lemma}

\begin{proof}
Use
\[
 g(z,w)-g(0,w)
 =\log|1-z\overline w|-\log|1-z/w|
\]
and the Taylor series for both logarithms.  Uniform convergence follows
from $|z/w|\le r/R<1$ and $|z\overline w|\le r<1$.
For the final bound write $w=s e^{i\theta}$.  Since
\[
 s^{-k}-s^k=\int_s^1 k(t^{-k-1}+t^{k-1})\,dt
 \le k(1-s)(R^{-k-1}+R^{-1}),
\]
the sum of the absolute values is at most
\[
 (1-s)\sum_{k\ge1}r^k(R^{-k-1}+R^{-1})
 \le C_{r,R}(1-s).
\]
\end{proof}

The last bound involves the first boundary deficit
$\sum_w(1-|w|)$, rather than a power strictly larger than one.  Uniform
control of higher powers of the boundary deficit therefore does not
justify discarding this field.  Similarly, the first proposition of this
section cannot be combined with an independent Gaussian harmonic limit
by simply factoring expectations: the Green factor is a statistic of
particles throughout the disk.  Establishing the required boundary-tail
control, or a different uniqueness principle that avoids it, remains
necessary for this approach.

\section{A precise variance criterion for the exterior inverse field}
\label{sec:inverse-field-variance-criterion}

The result in this section is conditional. It isolates a quantitative
estimate that would control the exterior field which is invisible to
vague convergence. It does not assert that the estimate holds for the gas,
and it does not by itself assert uniqueness.

For $k\geq1$ and $s<1$, define the angular count
\[
 A_{k,N}(s)=\sum_{|z_i|\leq s}e^{-ik\arg z_i}.
\]
The convention for the argument at zero is immaterial, since a particle
is at zero with probability zero. Consider the following hypothesis:
there exist $s_*\in(0,1)$, $C<\infty$, $b\geq0$, and $0\leq\gamma<2$ such that
\begin{equation}\label{eq:inverse-variance-hypothesis}
 \sup_N\mathbb E|A_{k,N}(s)|^2
 \leq C k^b(1-s)^{-\gamma},
 \qquad k\geq1,\quad s_*\leq s<1.
\end{equation}
The previously proved count bound only gives an estimate with exponent
$\gamma=2$. Thus \eqref{eq:inverse-variance-hypothesis} is a genuine
additional assumption.

\begin{lemma}[Verification at the determinantal point]
\label{lem:inverse-variance-beta-two}
For $\beta=2$, hypothesis~\eqref{eq:inverse-variance-hypothesis} holds
with $b=0$ and $\gamma=1$.
\end{lemma}
\begin{proof}
Let $f(z)=\mathbf1_{\{|z|\le s\}}e^{-ik\arg z}$. Rotation invariance
gives $\mathbb E\Xi_N(f)=0$. The determinantal one- and two-point
formulas proved in Section~\ref{sec:beta-two} give
\[
 \mathbb E|\Xi_N(f)|^2
 =\int |f(z)|^2\rho_{1,N}(z)\,dA(z)
  -\iint f(z)\overline{f(w)}|K_N(z,w)|^2\,dA(z)dA(w).
\]
The second integral is nonnegative: expanding the finite orthonormal
basis of $K_N$ writes it as a sum of quantities
$|\int f(z)\phi_i(z)\overline{\phi_j(z)}\,dA(z)|^2$.
Consequently
\[
 \mathbb E|A_{k,N}(s)|^2\le\mathbb E M_N(s)
 =\sum_{j=1}^N s^{2j}
 \le\frac{s^2}{1-s^2}\le\frac1{1-s},
\]
uniformly in $k,N,s$. This is the claimed hypothesis at $2$.
\end{proof}

\begin{proposition}[Vanishing boundary tails of the inverse correction]
\label{prop:inverse-correction-tail}
Assume \eqref{eq:inverse-variance-hypothesis}, and put
$\delta=1-\gamma/2>0$. For $R<1$ define
\[
 C_{k,N}(R)=\sum_{|z_i|>R}(z_i^{-k}-\overline{z_i}^{\,k}).
\]
For every $R_0\in(\max\{s_*,0\},1)$, every $R\geq R_0$, and every
$k\geq1$,
\begin{equation}\label{eq:inverse-tail-L2}
 \sup_N\|C_{k,N}(R)\|_{L^2}
 \leq C_{R_0} k^{1+b/2}R_0^{-k-1}(1-R)^\delta.
\end{equation}
In particular, if $q<R_0$ and
\[
 H_{N,R}(z)=\sum_{k\geq1}\frac{z^k}{k}C_{k,N}(R),
 \qquad |z|<R,
\]
then
\begin{equation}\label{eq:inverse-tail-analytic-L2}
 \sup_N\left\|\sup_{|z|\leq q}|H_{N,R}(z)|\right\|_{L^2}
 \leq C_{q,R_0}(1-R)^\delta.
\end{equation}
\end{proposition}
\begin{proof}
Writing $z_i=u_i e^{i\theta_i}$ and using
\[
 u^{-k}-u^k=\int_u^1 k(s^{-k-1}+s^{k-1})\,ds,
\]
finite summation and Fubini give the exact identity
\begin{equation}\label{eq:inverse-tail-Abel}
 C_{k,N}(R)=\int_R^1 k(s^{-k-1}+s^{k-1})
                   [A_{k,N}(s)-A_{k,N}(R)]\,ds.
\end{equation}
There is no boundary atom at a deterministic radius, and changing the
endpoint convention does not affect this integral. The coefficient in
the integrand is at most $2kR_0^{-k-1}$. Minkowski's integral inequality
and \eqref{eq:inverse-variance-hypothesis} bound the $L^2$ norm by
\[
 C k^{1+b/2}R_0^{-k-1}
   \int_R^1[(1-s)^{-\gamma/2}+(1-R)^{-\gamma/2}]\,ds,
\]
which is \eqref{eq:inverse-tail-L2}. Another application of Minkowski,
now to the power series, gives \eqref{eq:inverse-tail-analytic-L2}, since
$\sum_{k\geq1}k^{b/2}(q/R_0)^k<\infty$. All series used before this
bound converge absolutely for each finite configuration when $|z|<R$.
\end{proof}

\begin{proposition}[Identification from the limiting configuration]
\label{prop:inverse-field-identification}
Assume \eqref{eq:inverse-variance-hypothesis} and
$\Xi_{N_j}\Rightarrow\Xi$. Fix $r\in(0,1)$. There is a measurable
random holomorphic function on $B_r$,
\begin{equation}\label{eq:inverse-exterior-limit-definition}
 H_r^{\Xi}(z)=\lim_{R\uparrow1}
    \sum_{r<|w|\leq R,\ w\in\Xi}\ \sum_{k\geq1}
       \frac{z^k}{k}(w^{-k}-\bar w^k),
\end{equation}
where the limit exists in $L^2$ for the supremum norm on every smaller
closed disk. It is measurable with respect to the exterior configuration
$\Xi|_{\{|w|>r\}}$. Jointly with the point process and its positive
Gaussian harmonic modes,
\[
 \sum_{|z_i|>r}\sum_{k\geq1}\frac{z^k}{k}
                         (z_i^{-k}-\bar z_i^k)
 \ \Longrightarrow\ H_r^{\Xi}(z)
 \quad\hbox{in }\mathcal H(B_r).
\]
The convergence is in the topology of uniform convergence on compact
subsets; the statement about the positive modes uses their already proved
joint analytic convergence and independence from $\Xi$.
\end{proposition}
\begin{proof}
First truncate the sums to $r<|w|\leq R<1$. They are continuous functions
of the configuration with values in $\mathcal H(B_r)$ at every
configuration with no point on the two delimiting circles. The limiting
process has this property almost surely, by its continuous intensity.
Compact count moments are uniformly bounded. Consequently convergence
on each finite annulus holds jointly with the process and its positive
analytic modes, and also gives convergence of second moments of each
bounded-annulus sum in any smaller-disk supremum norm. For the latter
assertion the supremum is bounded by a fixed constant times the number
of particles in that annulus, so uniform fourth count moments provide
uniform integrability.

For $R<S<1$, the finite-particle annulus sum between $R$ and $S$ equals
$H_{N,R}-H_{N,S}$. Proposition~\ref{prop:inverse-correction-tail} bounds
its $L^2$ supremum norm by $C(1-R)^\delta$, uniformly in $S,N$, once
$R$ exceeds a fixed radius larger than the parameter disk. Passing to
the limit on the finite annulus gives the same bound for $\Xi$.
The truncated sums for $\Xi$ are therefore Cauchy in $L^2$ on every
smaller disk. Their limits are compatible and holomorphic; for example,
take an almost surely uniformly convergent subsequence on a countable
compact exhaustion. They are exterior-measurable because every truncated
sum is. The uniform finite-$N$ tail bound and the corresponding limiting
tail bound now justify removal of the truncation in the joint convergence,
by the usual bounded-Lipschitz test argument: fix $R$, pass to the limit,
and then let $R\uparrow1$. This proves the assertions.
\end{proof}

\begin{proposition}[Same-count conditional density ratios]
\label{prop:inverse-field-conditional-ratios}
Under the same hypothesis, fix $r\in(0,1)$ and an integer $m\geq1$
with $\mathbb P\{\Xi(B_r)=m\}>0$. Condition the finite processes on
$\Xi_{N_j}(B_r)=m$. For their exterior configurations $\eta_N$, let
$p_N(x\mid\eta_N,m)$ be the conditional density of a uniformly randomly
labeled inside tuple $x\in B_r^m$. For distinct-coordinate tuples $x,y$,
one has the exact finite-$N$ identity
\begin{equation}\label{eq:same-count-exact-ratio}
 \frac{p_N(x\mid\eta_N,m)}{p_N(y\mid\eta_N,m)}
 =\frac{|\Delta_m(x)|^\beta}{|\Delta_m(y)|^\beta}
    \exp\left\{-\beta\operatorname{Re}
       \left[\sum_{i=1}^m E_N(x_i)-\sum_{i=1}^m E_N(y_i)\right]\right\},
\end{equation}
where
\[
 E_N(z)=\sum_{k\geq1}\frac{z^k}{k}
                           \sum_{w\in\eta_N}w^{-k}.
\]
As random holomorphic functions on $B_r$, these exterior fields converge
jointly with the conditioned point process to
\begin{equation}\label{eq:augmented-exterior-field}
 E_r(z)=\sum_{k\geq1}\frac{z^k}{k}
       \left(\bar G_k-
        \sum_{w\in\Xi\cap B_r}\bar w^k\right)+H_r^{\Xi}(z).
\end{equation}
Here $\Xi$ has its law conditioned on $\Xi(B_r)=m$, and the Gaussian
vector $(G_k)$ is independent of it. In particular, the random ratios
in \eqref{eq:same-count-exact-ratio} converge in law, locally uniformly
in $(x,y)$ away from collision diagonals, to the same expression with
$E_r$ in place of $E_N$.
\end{proposition}
\begin{proof}
Conditional on the exterior and the number of inside particles, the
finite density is proportional to
\[
 |\Delta_m(x)|^\beta
            \prod_{i=1}^m\prod_{w\in\eta_N}|x_i-w|^\beta.
\]
Its normalizer is finite and positive: there are finitely many exterior
points, all strictly outside the closed radius-$r$ disk almost surely,
and the displayed product is bounded and positive off its collision
diagonals. In the ratio of two such densities the normalizer cancels.
For $|z|<r<|w|$,
\[
 \log|z-w|=\log|w|-\operatorname{Re}
                      \sum_{k\geq1}\frac{z^k}{kw^k}.
\]
The constant exterior terms cancel because both tuples have $m$ entries.
This proves \eqref{eq:same-count-exact-ratio}.

The exact decomposition
\[
 \sum_{w\in\eta_N}w^{-k}
 =\overline{S_{k,N}}-
       \sum_{w\in\Xi_N\cap B_r}\bar w^k+C_{k,N}(r)
\]
and Proposition~\ref{prop:inverse-field-identification} identify the
limit field. The Gaussian series converges locally uniformly almost
surely by its second moments and geometric weights; the inside sum is
finite. Conditioning causes no loss of convergence: the count event
has a positive limiting probability, its indicator is almost surely
continuous at the limiting configuration, and joint convergence hence
passes to the conditional laws. Independence of the Gaussian vector
from $\Xi$ is preserved by this conditioning. Finally, exponentiation
and evaluation on compact tuple sets are continuous for local uniform
convergence of analytic functions, proving the ratio convergence.
\end{proof}

\begin{remark}[The remaining two issues]
Neither \eqref{eq:inverse-variance-hypothesis} nor a uniqueness conclusion
is proved in this section. The limiting ratio formula includes the
additional Gaussian field and the actual inside power sums. Convergence
of these random conditional-density ratios must not be identified,
without a separate argument, with convergence of conditional laws given
the limiting exterior configuration alone. Moreover, fixing the number
of inside particles removes the activity factor from a Gibbs density:
it gives no rule for selecting the conditional distribution of that
number. These are distinct issues from the vanishing-tail estimate.
Sections~\ref{sec:unconditional-canonical-quasigibbs}
and~\ref{sec:conditional-gaussian-curvature} later give an unconditional
canonical mixture theorem and prove that the additional harmonic-field
randomness cannot simply be discarded. The complete count-resolved
equation is obtained at $\beta=2$ in
Section~\ref{sec:bergman-complete-dlr}.
\end{remark}

\section{A positive radial mixture at every even inverse temperature}
\label{sec:even-beta-radial-mixture}

This section uses the polynomial structure available when $\beta=2m$,
where $m$ is a positive integer. It is a different finite-dimensional
representation from the ordered-gap expansion. In particular, it improves
some of the general bounds on the total radial depth.

\begin{proposition}[Unordered radial mixture and majorization]
\label{prop:even-beta-radial-mixture}
Expand the polynomial
\[
 \Delta_N(z)^m=\sum_{\mathbf k}c_{\mathbf k}
                    z_1^{k_1}\cdots z_N^{k_N}.
\]
All indices in this sum satisfy $k_i\geq0$ and
$\sum_i k_i=mN(N-1)/2$. There is a probability distribution
\begin{equation}
 p_N(\mathbf k)=
 \frac{|c_{\mathbf k}|^2\prod_i(k_i+1)^{-1}}
      {\sum_{\boldsymbol\ell}|c_{\boldsymbol\ell}|^2
                           \prod_i(\ell_i+1)^{-1}}
 \label{eq:even-beta-occupation-mixture}
\end{equation}
such that, conditionally on $\mathbf k$, the variables
$|z_i|^2$ are independent and have densities
$(k_i+1)x^{k_i}\mathbf1_{(0,1)}(x)\,dx$.
Every $\mathbf k$ with nonzero coefficient is majorized by
$\mathbf b=(0,m,2m,\ldots,m(N-1))$.
\end{proposition}

\begin{proof}
Write $z_i=\sqrt{x_i}e^{i\theta_i}$. Orthogonality of angular monomials
gives
\[
 \int_{[0,2\pi]^N}|\Delta_N(z)|^{2m}
                         \prod_i\frac{d\theta_i}{2\pi}
 =\sum_{\mathbf k}|c_{\mathbf k}|^2\prod_i x_i^{k_i}.
\]
Since $dA(z_i)=\pi\,dx_i\,d\theta_i/(2\pi)$, integrating the $x_i$
proves the asserted mixture, including its normalization.

The ordinary determinant expansion writes $\Delta_N$ as a signed sum
of monomials whose exponent vectors are permutations of
$(0,1,\ldots,N-1)$. A monomial appearing in its $m$-th power therefore
has exponent vector $\mathbf k$ equal to a sum of $m$ such permutations.
Cancellation can delete exponent vectors, but cannot create new ones.
For a set $I$ of $q$ coordinate indices,
\[
 \sum_{i\in I}k_i
 \leq m\big((N-q)+(N-q+1)+\cdots+(N-1)\big).
\]
Taking the largest $q$ coordinates gives the majorization inequalities;
the equality of total sums was already observed.
\end{proof}

\begin{theorem}[Sharp extremal product bounds for the radial mixture]
\label{thm:even-beta-radial-products}
Put $B_N=\sum_{i=1}^N-\log|z_i|$ and $\beta=2m$ as above.
For every $\gamma>0$,
\begin{align}
 \prod_{j=0}^{N-1}\frac{mj+1}{mj+1+\gamma/2}
 &\leq \mathbb E_{N,2m}e^{-\gamma B_N}
 \leq
 \left(\frac{1+m(N-1)/2}
 {1+m(N-1)/2+\gamma/2}\right)^N.
 \label{eq:even-beta-negative-depth-mgf}
\end{align}
For $0<\theta<2$ one also has
\begin{equation}
 \mathbb E_{N,2m}e^{\theta B_N}
 \leq\prod_{j=0}^{N-1}\frac{mj+1}{mj+1-\theta/2}.
 \label{eq:even-beta-positive-depth-mgf}
\end{equation}
Consequently
\begin{equation}
 \mathbb E_{N,2m}B_N
 \leq\frac12\sum_{j=0}^{N-1}\frac1{mj+1}
 =\frac1\beta\log N+O_\beta(1),
 \label{eq:even-beta-mean-depth}
\end{equation}
and, for each $0<\theta<2$, a finite constant $C_{\beta,\theta}$ satisfies
\begin{equation}
 \mathbb P_{N,\beta}\left\{B_N\geq\frac1\beta\log N+t\right\}
 \leq C_{\beta,\theta}e^{-\theta t},\qquad t\geq0.
 \label{eq:even-beta-depth-upper-tail}
\end{equation}
In the charge-matched case there is the particularly simple bound
\begin{equation}
 \mathbb E_{N,\beta}\prod_i|z_i|^\beta
 \geq\frac1{mN+1}=\frac1{\beta N/2+1}.
 \label{eq:even-beta-charge-matched-bound}
\end{equation}
\end{theorem}

\begin{proof}
Conditioning on the mixture index gives
\[
 \mathbb E[e^{-\gamma B_N}\mid\mathbf k]
 =\prod_i\frac{k_i+1}{k_i+1+\gamma/2},\qquad
 \mathbb E[e^{\theta B_N}\mid\mathbf k]
 =\prod_i\frac{k_i+1}{k_i+1-\theta/2}.
\]
For $a>0$, the function
$g_a(x)=\log(x+1)-\log(x+1+a)$ is concave on $[0,\infty)$.
The majorization inequalities imply
$\sum_i g_a(k_i)\geq\sum_i g_a(b_i)$, while Jensen's inequality gives
$\sum_i g_a(k_i)\leq N g_a(m(N-1)/2)$.
In the present setting this use of majorization needs no general theorem:
there are permutations $\sigma_1,\ldots,\sigma_m$ such that
$\mathbf k=m^{-1}\sum_{\ell=1}^m\sigma_\ell\mathbf b$.
Coordinatewise Jensen, followed by summation in the coordinate index,
gives
$\sum_i g_a(k_i)\geq m^{-1}\sum_\ell\sum_i
g_a((\sigma_\ell\mathbf b)_i)=\sum_i g_a(b_i)$.
These bounds hold for each index, and hence after mixing.

For $a=\theta/2\in(0,1)$, the function
$h_a(x)=\log(x+1)-\log(x+1-a)$ is convex. Coordinatewise Jensen in the
same average-of-permutations representation reverses the preceding
inequality and proves \eqref{eq:even-beta-positive-depth-mgf}.
The function $x\mapsto1/(x+1)$ is also convex. The conditional expectation
of $B_N$ is $\frac12\sum_i(k_i+1)^{-1}$, so the same argument gives
\eqref{eq:even-beta-mean-depth}.

Finally,
\[
 \sum_{j=0}^{N-1}\log\frac{mj+1}{mj+1-\theta/2}
 =\frac{\theta}{2m}\log N+O_{m,\theta}(1).
\]
Indeed the $j=0$ term is finite, and the Taylor remainder after the
linear term for $j\geq1$ is bounded by a constant times $(mj+1)^{-2}$;
these remainders are summable. Exponential Markov gives
\eqref{eq:even-beta-depth-upper-tail}. When $\gamma=2m$, the lower product
in \eqref{eq:even-beta-negative-depth-mgf} telescopes to $1/(mN+1)$.
\end{proof}

\begin{remark}[Scope of the improvement]
At $\beta=2$, the only exponent vectors are permutations of
$(0,1,\ldots,N-1)$, so both product formulas before their Jensen upper
bound are exact. At $\beta=4$ the charge-matched estimate reads
$\mathbb E\prod_i|z_i|^4\geq(2N+1)^{-1}$.
These are finite-$N$ bounds; they provide no lower bound on
$B_N$ tending to infinity, and in particular do not prove that the
critical process is empty.
\end{remark}


\section{An exact Bergman intensity bound at even inverse temperature}
\label{sec:bergman-exclusion-even-beta}

Let $\mathcal B=A^2(\mathbb D,dA)$ be the Bergman Hilbert space with
reproducing kernel
\[
 K(z,w)=\frac1{\pi(1-z\overline w)^2}.
\]
For a closed subspace $\mathcal H\subset\mathcal B$, write $K_{\mathcal H}$
for its reproducing kernel. Point evaluation remains continuous on every
such subspace.

\begin{lemma}[Exclusion for a reproducing-kernel direction]
\label{lem:bergman-reproducing-exclusion}
Let $F\in\mathcal H^{\otimes N}$ be a holomorphic function of nonzero
norm, and suppose $F$ vanishes whenever any two coordinates coincide.
Give the labeled points density $|F|^2/\|F\|^2$ with respect to $dA^N$,
and let $\rho_F$ be the intensity of their counting measure. Then
\begin{equation}
 \rho_F(a)\leq K_{\mathcal H}(a,a),\qquad a\in\mathbb D.
 \label{eq:bergman-reproducing-exclusion}
\end{equation}
No permutation symmetry assumption is necessary.
\end{lemma}

\begin{proof}
If $K_{\mathcal H}(a,a)=0$, evaluation at $a$ vanishes identically on
$\mathcal H$, so every marginal density at $a$ is zero. Otherwise let
$u_a=K_{\mathcal H}(\cdot,a)/\sqrt{K_{\mathcal H}(a,a)}$, and let $P_i$
be the orthogonal rank-one projection onto $u_a$ in tensor coordinate
$i$, acting as the identity in the other coordinates. Put $\Psi=F/\|F\|$.
The reproducing identity gives
\[
 \|P_i\Psi\|^2
 =\frac1{K_{\mathcal H}(a,a)}
   \int_{\mathbb D^{N-1}}|\Psi(z_1,\ldots,z_{i-1},a,
                                  z_{i+1},\ldots,z_N)|^2\,dA^{N-1}.
\]
The projections commute, and $P_iP_j\Psi=0$ for $i\ne j$, because
the corresponding double evaluation is zero. For $Q=\sum_iP_i$,
\[
 \langle\Psi,Q^2\Psi\rangle=\langle\Psi,Q\Psi\rangle=:q.
\]
Cauchy--Schwarz implies $q^2\leq\|\Psi\|^2\|Q\Psi\|^2=q$;
as $q\geq0$, this gives $q\leq1$. Summing the preceding marginal
identity gives $q=\rho_F(a)/K_{\mathcal H}(a,a)$, proving the result.
\end{proof}

\begin{theorem}[Finite-$N$ bounds for the gas and its reduced Palm laws]
\label{thm:even-beta-bergman-bound}
Let $\beta=2m$ with $m$ a positive integer. For every $N\geq1$,
\begin{equation}
 \rho_{1,N}(z)\leq\frac1{\pi(1-|z|^2)^2},\qquad
 \mathbb E\Xi_N(B_r)\leq\frac{r^2}{1-r^2}.
 \label{eq:even-beta-bergman-intensity-count}
\end{equation}
For distinct fixed points $a_1,\ldots,a_p$ with $p<N$, define
\[
 \mathcal H_{\mathbf a,m}
 =\{f\in\mathcal B:f^{(j)}(a_i)=0,
                   \ 1\leq i\leq p,\ 0\leq j<m\}.
\]
The reduced Palm intensity, defined by the finite-$N$ correlation ratio,
satisfies
\begin{equation}
 \frac{\rho_{p+1,N}(a_1,\ldots,a_p,z)}
      {\rho_{p,N}(a_1,\ldots,a_p)}
 \leq K_{\mathcal H_{\mathbf a,m}}(z,z).
 \label{eq:even-beta-palm-bergman-bound}
\end{equation}
In particular, with $t=|(z-a)/(1-\overline a z)|^2$,
\begin{equation}
 \rho_{2,N}(a,z)
 \leq \rho_{1,N}(a)\,
       \frac{t^m(m+1-mt)}{\pi(1-|z|^2)^2}.
 \label{eq:even-beta-pair-bergman-bound}
\end{equation}
These bounds pass to every subsequential point-process limit through
the already proved locally uniform correlation convergence.
\end{theorem}

\begin{proof}
Use $F=\Delta_N^m$ and $\mathcal H=\mathcal B$ in the lemma; the density
is exactly the gas density at $\beta=2m$. Integrating the kernel diagonal
over $B_r$ gives $r^2/(1-r^2)$.

Conditioning on $p$ distinct marked points leaves the normalized squared
modulus of the polynomial
\[
 F_{\mathbf a}(w)
 =\Delta_{N-p}(w)^m
      \prod_{j=1}^{N-p}\prod_{i=1}^p(w_j-a_i)^m.
\]
Each coordinate belongs to $\mathcal H_{\mathbf a,m}$, and the polynomial
still vanishes on every collision of the remaining coordinates.
Applying the lemma proves the Palm bound. The denominator correlation
is strictly positive for distinct marked points, since its insertion
integral is positive.

For a single marked point at zero, the subspace consists of Bergman
functions whose first $m$ Taylor coefficients vanish. Its kernel diagonal is
\[
 K_{\mathcal H_{0,m}}(z,z)
 =\frac1\pi\sum_{k=m}^{\infty}(k+1)|z|^{2k}
 =\frac{|z|^{2m}(m+1-m|z|^2)}{\pi(1-|z|^2)^2}.
\]
The disk automorphism $\varphi_a(z)=(z-a)/(1-\overline a z)$ induces a
unitary Bergman operator $f\mapsto\varphi_a'\,(f\circ\varphi_a)$,
as is seen directly by changing variables in the area integral.
It maps the subspace vanishing to order $m$ at zero onto the subspace
vanishing to order $m$ at $a$. Consequently its kernel diagonal is
$|\varphi_a'(z)|^2K_{\mathcal H_{0,m}}(\varphi_a(z),\varphi_a(z))$.
The identity
$|\varphi_a'(z)|/(1-|\varphi_a(z)|^2)=1/(1-|z|^2)$ proves the displayed
pair bound. The limit assertion follows by passage to the inequalities
on compact sets; for the Palm bound it is enough to retain the
undivided form $\rho_{p+1}\leq\rho_p K_{\mathcal H}$ if the limiting
denominator vanishes.
\end{proof}

\begin{corollary}[A determinantal upper bound and uniform count moments]
\label{cor:even-beta-determinantal-upper-bound}
For every even $\beta=2m$, every $p\geq1$, and all distinct points,
\begin{equation}
 \rho_{p,N}(z_1,\ldots,z_p)
 \leq\det[K(z_i,z_j)]_{i,j=1}^p
 \leq\prod_{i=1}^p K(z_i,z_i).
 \label{eq:even-beta-determinantal-upper-bound}
\end{equation}
If $A\Subset\mathbb D$ is Borel and
$\mu_A=\int_A K(z,z)\,dA(z)$, then for every $t\geq0$,
\begin{equation}
 \mathbb E e^{t\Xi_N(A)}
 \leq \exp\{\mu_A(e^t-1)\}.
 \label{eq:even-beta-poisson-exponential-bound}
\end{equation}
All these inequalities hold also for subsequential limits and for
their images under arbitrary disk automorphisms. They are inequalities
for correlations and moments; they do not assert stochastic domination
of point-process laws.
\end{corollary}

\begin{proof}
Iterate \eqref{eq:even-beta-palm-bergman-bound}. The subspace vanishing
to order $m$ at each of the preceding marks is contained in the subspace
vanishing to order one there, so its kernel diagonal is smaller.
The order-one subspace is the orthogonal complement of the span of the
corresponding reproducing vectors. Its kernel diagonal at the next
point is the squared norm of the next vector after orthogonal projection
off that span. The product of these successive squared norms is the
Gram determinant $\det[K(z_i,z_j)]$. The vectors at distinct points
are linearly independent, since polynomial interpolation separates their
evaluations. This proves the first inequality for $p\leq N$; for $p>N$
the left side is zero. Hadamard's inequality, or the same projection
description, gives the second inequality.

Integrating it over $A^p$ bounds the $p$-th factorial moment by
$\mu_A^p$. For an integer-valued $M\geq0$,
$e^{tM}=\sum_{p\geq0}(M)_{\underline p}(e^t-1)^p/p!$.
All terms are nonnegative for $t\geq0$, so monotone convergence and
the factorial moment bounds prove the exponential estimate.
The locally uniform correlation convergence proves the limiting
inequalities; the same moment argument applies to the limiting process.
Finally the Bergman kernel and its vanishing subspaces transform
unitarily under disk automorphisms, as in the preceding proof.
Both sides of the correlation inequalities acquire the same product
Jacobian factor. This proves their covariance, and hence the assertion
for transformed laws.
\end{proof}

\begin{remark}[What the bound does not settle]
At $\beta=4$, the bound gives the exact finite-$N$ estimate
$\pi(1-|z|^2)^2\rho_{1,N}(z)\leq1$. It does not show that its left side
tends to zero by itself; the subsequent critical emptiness proof combines
it with the canonical-score estimates and a second limiting procedure.
The rank-one argument bounds occupation of each coherent
kernel direction; it does not assert domination of the entire one-body
operator by the identity for arbitrary symmetric holomorphic states.
\end{remark}


\section{Electric forces at particles and an augmented limiting Ward identity}
\label{sec:campbell-force}

The estimates in this section concern the full electric force evaluated
at an actual particle. They do not assume control of the exterior
inverse field. In particular, the marks constructed below are not
asserted to be measurable functions of the limiting configuration.

For a distinct finite tuple define the complex force
\[
 C_{i,N}=\sum_{j\ne i}\frac1{z_i-z_j}.
\]
The real directional derivative of the log density in a direction
$v\in\C$ at the $i$th coordinate is
$\beta\operatorname{Re}(v C_{i,N})$. Thus the real force vector,
in the usual identification of $\R^2$ with $\C$, is $\beta\bar C_{i,N}$.
Write $\partial=(\partial_x-i\partial_y)/2$ and
$\Delta=\partial_x^2+\partial_y^2$.

\begin{proposition}[Exact Campbell force identities]
\label{prop:campbell-force-identities}
For every fixed $N$ and $\beta>0$, $\rho_{1,N}$ is $C^2$ in $\D$.
For every $\phi\in C_c^\infty(\D)$,
\begin{align}
 \beta^2\E\sum_i\phi(z_i)|C_{i,N}|^2
 &=\int_{\D}\rho_{1,N}(z)\Delta\phi(z)\,dA(z),
 \label{eq:campbell-force-square}\\
 \rho_{1,N}(a)\E_N^{!a}C_{a,N}
 &=\frac2\beta\partial\rho_{1,N}(a).
 \label{eq:campbell-force-mean}
\end{align}
Here $\E_N^{!a}$ is the reduced Palm expectation at $a$, and
$C_{a,N}=\sum_{w\in\Xi_N}1/(a-w)$ under that reduced Palm law.
The first identity is valid also for signed $\phi$.
Moreover,
\begin{equation}\label{eq:campbell-force-log-curvature}
 \Delta\log\rho_{1,N}(a)
 =\beta^2\left(\E_N^{!a}|C_{a,N}|^2
                       -|\E_N^{!a}C_{a,N}|^2\right).
\end{equation}
\end{proposition}

\begin{proof}
Use the finite positive measure $\nu_N$ from the insertion formula,
so that
\[
 \rho_{1,N}(a)=\int |P_y(a)|^\beta\,d\nu_N(y),\qquad
 P_y(a)=\prod_{j=1}^{N-1}(a-y_j).
\]
For every polynomial $P$, in the sense of distributions,
\begin{equation}\label{eq:polynomial-force-laplacian}
 \Delta|P|^\beta=\beta^2|P|^\beta|P'/P|^2.
\end{equation}
At a zero of multiplicity $q$, the right side is locally bounded by
a constant times $|a-y|^{q\beta-2}$ and hence is integrable. To verify
that there is no additional point mass, integrate by parts outside
radius-$\varepsilon$ circles around the zeros. The two boundary terms
are $O(\varepsilon^{q\beta})$ and $O(\varepsilon^{q\beta+1})$,
respectively, and tend to zero. Away from the zeros the identity follows
by two ordinary derivatives of $\exp(\beta\log|P|)$.

For nonnegative $\phi$, integrate
\eqref{eq:polynomial-force-laplacian} against $\phi$, and then against
$\nu_N$. Tonelli applies on the nonnegative right side; the other side
is absolutely bounded by
$\|\Delta\phi\|_\infty\int_{\supp\phi}\rho_{1,N}\,dA$.
This proves the first formula and its local integrability. Signed tests
follow by domination with a compactly supported nonnegative cutoff.
One distributional derivative similarly gives
\[
 \partial\rho_{1,N}(a)
 =\frac\beta2\int |P_y(a)|^\beta
                   \sum_j(a-y_j)^{-1}\,d\nu_N(y).
\]

Here the density of $\nu_N$ with respect to area on $\D^{N-1}$ is
bounded, with a bound allowed to depend on $N$. Each integrand in a
first or second real derivative is bounded in absolute value by a
finite sum of constant multiples of
\[
 |a-y_j|^{\beta-2}\prod_{\ell\ne j}|a-y_\ell|^\beta
 \quad\hbox{and}\quad
 |a-y_j|^{\beta-1}|a-y_\ell|^{\beta-1}
                       \prod_{h\notin\{j,\ell\}}|a-y_h|^\beta.
\]
These expressions are locally integrable uniformly in $a$ on compact
sets. Their integrals over the sets where any singular denominator is
within $\varepsilon$ of zero tend to zero uniformly in $a$: for the
first expression use $\int_{|u|<\varepsilon}|u|^{\beta-2}dA(u)
=2\pi\varepsilon^\beta/\beta$, and for the second use the analogous
integral with exponent $\beta-1$. If an exponent is nonnegative, ordinary
boundedness suffices. Off those sets, dominated convergence applies.
Thus the distributional first and second derivatives are continuous,
which proves $C^2$ regularity and the pointwise Palm formulas.
Finally $\rho_{1,N}>0$, and differentiating its logarithm and using
$|\nabla\rho_{1,N}|^2=4|\partial\rho_{1,N}|^2$ gives
\eqref{eq:campbell-force-log-curvature}.
\end{proof}

\begin{corollary}[Uniform local force energy and Fisher information]
\label{cor:uniform-campbell-force}
For every compact $K\Subset\D$ there is $C_{\beta,K}<\infty$ such that
\begin{equation}\label{eq:uniform-campbell-force}
 \sup_N\E\sum_{z_i\in K}|C_{i,N}|^2\le C_{\beta,K}.
\end{equation}
More explicitly, for $0\le r<1$,
\begin{equation}\label{eq:explicit-campbell-force}
 \sup_N\E\sum_{|z_i|\le r}|C_{i,N}|^2
 \le C_\beta(1-r)^{-3}.
\end{equation}
For every nonnegative $\phi\in C_c^\infty(\D)$,
\begin{equation}\label{eq:force-fisher}
 \int\phi\,|\nabla\sqrt{\rho_{1,N}}|^2\,dA
 \le\frac14\int\rho_{1,N}\Delta\phi\,dA.
\end{equation}
Consequently every limiting intensity satisfies
$\sqrt{\rho_1}\in H^1_{\mathrm{loc}}(\D)$ and the same inequality,
with the right side interpreted distributionally.
\end{corollary}

\begin{proof}
Choose a nonnegative smooth cutoff equal to one on $K$ and apply
\eqref{eq:campbell-force-square}, using the uniform compact count bound.
For the explicit version use a cutoff equal to one on $B_r$, supported
in $B_R$, where $R=(1+r)/2$, whose second derivatives are bounded by
$C(1-r)^{-2}$. Such cutoffs also exist for $r$ near zero, by choosing
the cutoff constant on a fixed smaller centered disk. The radial count
estimate $\E M_N(R)\le (1-R^{\beta/2})^{-1}$, and
$1-R^{\beta/2}\ge c_\beta(1-R)$, give the stated power.
By Jensen in the Palm law,
\[
 |\nabla\sqrt{\rho_{1,N}}|^2
 =\frac{\beta^2}4\rho_{1,N}|\E_N^{!a}C_{a,N}|^2
 \le\frac{\beta^2}4\rho_{1,N}\E_N^{!a}|C_{a,N}|^2.
\]
Integration proves \eqref{eq:force-fisher}. Along a process-convergent
subsequence the intensities converge locally uniformly, hence their
square roots converge in local $L^2$. Cutoffs equal to one on compact
sets give a locally bounded sequence in $H^1$. Weak compactness and
lower semicontinuity, including after multiplication of gradients by
$\sqrt\phi$, prove the assertion for the limiting intensity.
\end{proof}

\begin{theorem}[A force-marked subsequential limit]
\label{thm:force-marked-limit}
Suppose $\Xi_{N_j}\Rightarrow\Xi$. There is a further subsequence along
which the particles and their force marks have a joint limit
\[
 \left(\Xi_{N_j},\sum_i\delta_{(z_i,C_{i,N_j})}\right)
 \ \Longrightarrow\
 \left(\Xi,\sum_{z\in\Xi}\delta_{(z,C_z)}\right).
\]
The marked measures have the vague topology on $\D\times\C$.
Every particle of $\Xi$ has exactly one finite complex mark $C_z$.
For nonnegative $\phi\in C_c^\infty(\D)$,
\begin{equation}\label{eq:limiting-force-energy}
 \beta^2\E\sum_{z\in\Xi}\phi(z)|C_z|^2
 \le\int\rho_1\Delta\phi\,dA.
\end{equation}

Let $v\in C_c^1(\D;\C)$ be a real vector field and let
\[
 F(\eta)=\Phi\left(\int h_1d\eta,\ldots,\int h_qd\eta\right),
 \qquad h_j\in C_c^1(\D;\R),\quad\Phi\in C_b^1(\R^q;\R),
\]
with bounded first derivatives. Define
\[
 \mathcal L_vF(\eta)
 =\sum_j\partial_j\Phi\left(\int h_1d\eta,\ldots,\int h_qd\eta\right)
                   \int dh_j(z)[v(z)]\,d\eta(z).
\]
The augmented limiting law satisfies
\begin{equation}\label{eq:augmented-limit-ward}
 \E\left[\mathcal L_vF(\Xi)
       +F(\Xi)\sum_{z\in\Xi}\operatorname{div}v(z)
       +\beta F(\Xi)\operatorname{Re}
                         \sum_{z\in\Xi}v(z)C_z\right]=0.
\end{equation}
All three expectations are absolutely convergent.
\end{theorem}

\begin{proof}
On a compact base set $K$, the number of marks is bounded by $\Xi_N(K)$,
whose moments are uniformly bounded. This gives tightness of the marked
measures on compact subsets of $\D\times\C$, jointly with $\Xi_N$.
Moreover, \eqref{eq:uniform-campbell-force} gives
\[
 \E\#\{i:z_i\in K,\ |C_{i,N}|>T\}\le C_{\beta,K}T^{-2}.
\]
Thus no spatial particle is lost through escape of its mark to infinity.
To formalize the identification of the projection, test against a
continuous function with compact support in the base, multiply it by
a continuous mark cutoff equal to one on $|c|\le T$, pass to a joint
subsequential limit, and then let $T\to\infty$. The displayed bound
makes the discrepancy tend to zero in $L^1$. The projection is therefore
$\Xi$ almost surely. Since $\Xi$ is simple, there is exactly one mark
over each of its particles. Continuous compact mark cutoffs followed
by monotone approximation of $|c|^2$ yield
\eqref{eq:limiting-force-energy} from
\eqref{eq:campbell-force-square} and local uniform intensity convergence.

At finite $N$, integration by parts in the $i$th coordinate, then
summation in $i$, gives
\begin{equation}\label{eq:finite-force-ward}
 \E\left[\mathcal L_vF(\Xi_N)
       +F(\Xi_N)\sum_i\operatorname{div}v(z_i)
       +\beta F(\Xi_N)\operatorname{Re}\sum_iv(z_i)C_{i,N}\right]=0.
\end{equation}
There is no boundary term on $\partial\D$ because $v$ is compactly
supported. For the collision diagonals this can either be justified
by deleting small circles and observing the boundary factors
$O(\varepsilon^{\beta+1})$, followed by integration in the other
coordinates, or by using the locally integrable first distributional
derivative of the density. Its absolute value is bounded by a finite
sum of products containing one factor $|z_i-z_j|^{\beta-1}$;
these are integrable for $\beta>0$.

The first two terms in \eqref{eq:finite-force-ward} pass to the limit
by vague continuity and uniform compact count moments. For the last
term use a continuous mark cutoff at level $T$. The cut off statistic
is continuous and bounded by a constant times $T\Xi_N(K)$, so its
expectation also converges. The omitted tail is bounded in absolute
expectation, uniformly in $N$, by
\[
 \|F\|_\infty\|v\|_\infty
   \E\sum_{z_i\in K}|C_{i,N}|\mathbf1_{\{|C_{i,N}|>T\}}
 \le C_{\beta,F,v,K}/T.
\]
The limiting energy bound gives the same tail bound for the limit.
Letting $T\to\infty$ proves \eqref{eq:augmented-limit-ward} and absolute
integrability. The statements about tightness and subsequences here
can also be obtained directly by exhausting both spaces by compact
sets, bounding their point counts, and applying diagonal compactness
of bounded finite point measures on each compact set.
\end{proof}

\begin{remark}[What the force marks do and do not resolve]
The theorem constructs finite, locally square-integrable forces at the
particles of every subsequential limit and supplies an unconditional
integration-by-parts identity on this augmented probability space.
It does not show that $C_z$ is determined by $\Xi$, identify it with a
particular summation of the infinite Coulomb force, or establish
uniqueness of either the unmarked or the augmented law. A boundary
harmonic contribution may still survive in the marks. In particular,
the angular-count variance hypothesis in
Section~\ref{sec:inverse-field-variance-criterion} has not been deduced
from the local force-energy bound.
\end{remark}


\begin{corollary}[A canonical configuration-measurable logarithmic derivative]
\label{cor:canonical-force-score}
For every subsequential limiting law there is a measurable complex
function $b(z,\eta)$, defined for $z\in\eta$ up to Campbell-null sets,
which satisfies
\[
 \beta^2\E\sum_{z\in\Xi}\phi(z)|b(z,\Xi)|^2
 \le \int\rho_1\Delta\phi\,dA,
 \qquad 0\le\phi\in C_c^\infty(\D),
\]
and, for the cylinder functions and vector fields of
Theorem~\ref{thm:force-marked-limit},
\[
 \E\left[\mathcal L_vF(\Xi)
 +F(\Xi)\sum_{z\in\Xi}\operatorname{div}v(z)
 +\beta F(\Xi)\operatorname{Re}
       \sum_{z\in\Xi}v(z)b(z,\Xi)\right]=0.
\]
Among locally Campbell-integrable functions, this identity determines
$b$ uniquely up to Campbell-null sets. Thus the real vector field
$\beta\overline{b(z,\eta)}$ is a canonical logarithmic derivative of the
limiting law. This assertion does not identify it with a prescribed
summation of the infinite Coulomb force.
\end{corollary}

\begin{proof}
Consider the Campbell measure of the augmented law on triples
$(\eta,z,c)$:
\[
 \widehat\Gamma(H)=\E\sum_{z\in\Xi}H(\Xi,z,C_z).
\]
Its projection to $(\eta,z)$ is the ordinary Campbell measure
$\Gamma$, and both measures are sigma-finite, with finite mass over
compact base sets. The configuration space and the other factors are
standard Borel. Disintegrate $\widehat\Gamma$ with respect to $\Gamma$
and denote the conditional probability kernel of the mark by
$\kappa_{\eta,z}(dc)$. This can, equivalently, be done after multiplication
by a strictly positive integrable weight depending only on $z$, which
turns the Campbell measures into finite measures without changing the
conditional kernel. Set
\[
 b(z,\eta)=\int c\,\kappa_{\eta,z}(dc).
\]
The local second-moment bound of
Theorem~\ref{thm:force-marked-limit} makes this integral finite
$\Gamma$-almost everywhere. Conditional Jensen gives the asserted
energy bound. Conditional expectation of the augmented Ward identity
gives the unmarked identity, because its force coefficient
$F(\eta)v(z)$ is measurable with respect to $(\eta,z)$.

For uniqueness, let $b_1,b_2$ be two locally Campbell-integrable
functions satisfying the identity. Subtract their identities. For every
admissible real cylinder function $F$ and every complex vector field
$v\in C_c^1(\D;\C)$ this gives
\[
 \operatorname{Re}\int F(\eta)v(z)
              (b_1-b_2)(z,\eta)\,d\Gamma(\eta,z)=0.
\]
Using $iv$ as well gives the same identity without the real part.
The bounded smooth cylinder functions form a unital algebra generating
the Borel sigma-algebra of the configuration space: integrals against
a countable separating family in $C_c^1(\D)$ generate the vague
sigma-algebra, and bounded smooth scalar transforms recover their
measurability. For each fixed $v$, the monotone-class theorem therefore
extends the last equality to every bounded Borel function $F$.
On each compact base region the measure
$(b_1-b_2)\,d\Gamma$ has finite total variation. Continuous compactly
supported spatial tests, approximated by $C_c^1$ tests, determine finite
measures. Hence the equality for all product tests $F(\eta)v(z)$ implies
$(b_1-b_2)\,d\Gamma=0$ there. Exhausting the disk proves uniqueness.
\end{proof}


\section{Gaussian force regression and a lower bound on intensity curvature}
\label{sec:gaussian-force-curvature}

The independent Gaussian harmonic field forces a specific component
of the electric-force fluctuations. This gives a constraint on every
limiting intensity, and a second proof of emptiness above $4$.
The proof in this section uses compactness, compact correlation bounds,
the finite-mode Gaussian theorem and its independence from the limiting
configuration, and Section~\ref{sec:campbell-force}. It does not use the
boundary transport estimate for adjacent partition functions.

\begin{theorem}[Regression of force marks on the Gaussian modes]
\label{thm:gaussian-force-regression}
Refine the subsequence of Theorem~\ref{thm:force-marked-limit} so that
the force-marked configuration converges jointly with every fixed
positive harmonic mode. Let $(G_k)_{k\ge1}$ be their limiting independent
circular Gaussian variables, independent of $\Xi$, with
$\E|G_k|^2=2k/\beta$. Let $b(z,\Xi)$ be the canonical score of
Corollary~\ref{cor:canonical-force-score}. Then, in the sense of
conditional expectations under the Campbell measure,
\begin{equation}\label{eq:gaussian-force-regression}
 \E[C_z\mid\Xi,z,G_1,\ldots,G_m]
 =b(z,\Xi)-\sum_{k=1}^m\overline{G_k}z^{k-1}.
\end{equation}
In particular, for every nonnegative $\phi\in C_c^\infty(\D)$,
\begin{equation}\label{eq:gaussian-force-energy-lower}
 \int\rho_1\Delta\phi\,dA
 \ge\beta^2\E\sum_{z\in\Xi}\phi(z)|b(z,\Xi)|^2
       +2\beta\int\frac{\phi(z)\rho_1(z)}{(1-|z|^2)^2}\,dA(z).
\end{equation}
The infinite series in \eqref{eq:gaussian-force-regression}, with
$m=\infty$, also converges in local Campbell $L^2$ and gives the
conditional expectation given all the $G_k$.
\end{theorem}

\begin{proof}
Tightness of the marked configurations and of every finite collection
of harmonic modes gives the asserted joint subsequence by diagonal
extraction. The marginal law of $(\Xi,(G_k))$ retains the proved
Gaussian independence. In particular its Campbell measure factors as
the Campbell measure of $\Xi$ times the Gaussian probability law.

Take a cylinder function $F$ and vector field $v$ as in
Theorem~\ref{thm:force-marked-limit}, and a real bounded $C^1$ function
$h$ on $\C^m\simeq\R^{2m}$ with bounded derivatives. Apply the finite
integration-by-parts calculation to
$F(\Xi_N)h(S_{1,N},\ldots,S_{m,N})$. Although the power sums are not
local, at fixed $N$ they and their derivatives are bounded, and the
vector field remains compactly supported. In addition to the three
terms of \eqref{eq:finite-force-ward}, multiplied by $h$, the derivative
of $h$ contributes
\[
 F(\Xi_N)\,Dh(S_{1,N},\ldots,S_{m,N})
       \left[\left(k\sum_i v(z_i)z_i^{k-1}\right)_{k=1}^m\right].
\]
All terms pass to the joint limit: the first and derivative terms are
bounded by constants times compact particle counts, and the force term
is handled by the mark truncation and $L^2$ tail bound in the proof of
Theorem~\ref{thm:force-marked-limit}.

For a circular complex Gaussian of variance $2k/\beta$, ordinary real
Gaussian integration by parts, or its Wirtinger form, gives
\[
 \E\,\partial_{G_k}h(G)=\frac\beta{2k}\E[h(G)\overline{G_k}].
\]
Indeed its density is proportional to
$\exp(-\beta|g|^2/(2k))$. Since
$Dh(G)[d]=2\operatorname{Re}\sum_k\partial_{G_k}h(G)d_k$,
conditioning first on $\Xi$ turns the preceding derivative term into
\[
 \beta\E\left[F(\Xi)h(G)\operatorname{Re}
          \sum_{z\in\Xi}v(z)\sum_{k=1}^m\overline{G_k}z^{k-1}\right].
\]
Write $T_m(z)=\sum_{k=1}^m\overline{G_k}z^{k-1}$. The other two
nonforce terms factor by independence of $G$ and $\Xi$.
Subtract the unmarked Ward identity for $b$, multiplied by $\E h(G)$.
The result is
\[
 \E\left[F(\Xi)h(G)\operatorname{Re}
         \sum_{z\in\Xi}v(z)
                   \bigl(C_z+T_m(z)-b(z,\Xi)\bigr)\right]=0.
\]
Replacing $v$ by $iv$ removes the real-part restriction. Product tests
$F(\eta)h(g)v(z)$ determine the locally finite Campbell measure:
bounded smooth cylinder functions and bounded smooth functions of $g$
are generating algebras, and compactly supported smooth spatial tests
determine finite measures. The monotone-class argument used in
Corollary~\ref{cor:canonical-force-score} therefore proves the
conditional-expectation identity.

Conditional Jensen and Gaussian independence now give
\[
 \E\sum_{z\in\Xi}\phi(z)|C_z|^2
 \ge\E\sum_{z\in\Xi}\phi(z)|b(z,\Xi)|^2
       +\frac2\beta\int\phi(z)\rho_1(z)
                         \sum_{k=1}^m k|z|^{2k-2}\,dA(z).
\]
There is no cross term because the Gaussian variables have zero mean
and are independent of the unmarked Campbell measure. Combine this
with \eqref{eq:limiting-force-energy}, let $m\to\infty$, and use
$\sum_{k\ge1}kx^{k-1}=(1-x)^{-2}$ to obtain
\eqref{eq:gaussian-force-energy-lower}. On a compact spatial set,
the same geometric series proves $T_m$ is Cauchy in Campbell $L^2$.
Conditional expectations converge in $L^2$ to the expectation given
the sigma-algebra generated by all the modes. One can apply this usual
probability-space assertion after restriction to any finite-mass
compact Campbell window; these restrictions exhaust the disk.
\end{proof}



\begin{theorem}[Curvature and boundary growth of a nonempty intensity]
\label{thm:limiting-intensity-curvature}
Every subsequential limiting intensity satisfies the measure inequality
\begin{equation}\label{eq:intensity-curvature-raw}
 \Delta\rho_1
 \ge\frac{|\nabla\rho_1|^2}{\rho_1}\,dA
       +\frac{2\beta\rho_1}{(1-|z|^2)^2}\,dA,
\end{equation}
where the first quotient is assigned zero on $\{\rho_1=0\}$.
On every open set where $\rho_1>0$,
\begin{equation}\label{eq:intensity-log-curvature-lower}
 \Delta\log\rho_1\ge\frac{2\beta}{(1-|z|^2)^2}\,dA.
\end{equation}
If the limiting process is not almost surely empty, then
\begin{equation}\label{eq:weighted-intensity-monotonicity}
 r\longmapsto (1-r^2)^{\beta/2}\rho_1(r)
 \quad\hbox{is nondecreasing on }(0,1).
\end{equation}
Consequently, for $0<r_0\le r<1$,
\begin{equation}\label{eq:nonempty-boundary-growth}
 \rho_1(r)\ge\rho_1(r_0)
             \left(\frac{1-r_0^2}{1-r^2}\right)^{\beta/2}.
\end{equation}
\end{theorem}

\begin{proof}
Disintegrate the unmarked Campbell measure with respect to its spatial
marginal $\rho_1(z)dA(z)$, and write $m_b(z)$ for the mean of
$b(z,\Xi)$ under this conditional law, and put $J(z)=\rho_1(z)m_b(z)$.
The Ward identity for $F=1$ gives, distributionally,
\[
 \partial_x\rho_1=\beta\operatorname{Re}J,
 \qquad \partial_y\rho_1=-\beta\operatorname{Im}J.
\]
Local boundedness of $\rho_1$ and the Campbell $L^2$ bound imply
$J\in L^2_{\rm loc}$ by conditional Jensen. Thus
$\rho_1\in H^1_{\rm loc}$, and another conditional Jensen inequality
gives
\[
 \frac{|\nabla\rho_1(z)|^2}{\rho_1(z)}
 \le\beta^2\rho_1(z)\E[|b(z,\Xi)|^2\mid z]
\]
almost everywhere, with the stated zero convention. Combining this
with \eqref{eq:gaussian-force-energy-lower} proves
\eqref{eq:intensity-curvature-raw}. The measure $\Delta\rho_1$ is
positive and locally finite, also directly from subharmonicity.

Here is a weak chain-rule justification of the logarithmic inequality.
On a relatively compact subset of $\{\rho_1>0\}$ the continuous
function $\rho_1$ is bounded above and below by positive constants.
For a nonnegative smooth test $\phi$ supported there, the function
$\psi=\phi/\rho_1$ is continuous, compactly supported and in $H^1$.
Nonnegative mollifications approximate it uniformly and in $H^1$,
with support in a common positive-intensity region. Testing the Radon
measure $\Delta\rho_1$ against these mollifications and passing to
the limit gives
\[
 \int\frac\phi{\rho_1}\,d(\Delta\rho_1)
 =-\int\nabla\rho_1\cdot\nabla(\phi/\rho_1)\,dA
 =\langle\Delta\log\rho_1,\phi\rangle
                  +\int\phi\frac{|\nabla\rho_1|^2}{\rho_1^2}\,dA.
\]
Inserting \eqref{eq:intensity-curvature-raw} cancels the gradient term
and proves \eqref{eq:intensity-log-curvature-lower}.

For a nonempty limiting law, the intensity dichotomy gives
$\rho_1(r)>0$ for every $r>0$. Since
$\Delta\log(1-|z|^2)=-4/(1-|z|^2)^2$, the radial function
\[
 h(z)=\log\rho_1(z)+\frac\beta2\log(1-|z|^2)
\]
is subharmonic on the punctured disk. It is locally bounded above near
zero, since $\rho_1$ is continuous and locally bounded. The removable
singularity property gives a subharmonic extension across zero,
possibly with value $-\infty$ there. Therefore $t\mapsto h(e^t)$ is
convex and nondecreasing, by the radial subharmonicity argument used
earlier for the intensity. Exponentiating proves the last two claims.
\end{proof}

\begin{corollary}[An independent empty-limit proof and the critical constraint]
\label{cor:force-empty-critical}
For every $\beta>4$ the full process sequence converges to the empty
configuration. At $\beta=4$, any subsequential limit that is not almost
surely empty must satisfy
\begin{equation}\label{eq:critical-boundary-intensity-coefficient}
 \lim_{r\uparrow1}(1-r^2)^2\rho_1(r)=c,
 \qquad 0<c\le1/\pi.
\end{equation}
The critical assertion is a necessary condition for a nonempty limit,
not an assertion that such a limit exists.
\end{corollary}

\begin{proof}
Radial monotonicity and the count comparison, applied with
$R=(1+r)/2$, give
\[
 \rho_1(r)\pi(R^2-r^2)
 \le\E\Xi(B_R)\le\frac1{1-R^{\beta/2}},
\]
and hence $\rho_1(r)\le C_\beta(1-r)^{-2}$.
For $\beta>4$, this contradicts
\eqref{eq:nonempty-boundary-growth} if the limit is nonempty.
Every subsequential limit is therefore empty, and tightness implies
convergence of the full sequence. This argument gives no rate in $N$.

At $\beta=4$, the weighted intensity is positive, nondecreasing, and
bounded above, and thus has a finite positive limit $c$. Moreover,
\[
 \E\Xi(B_r)=2\pi\int_0^r s\rho_1(s)\,ds,
 \qquad \E\Xi(B_r)\le(1-r^2)^{-1}.
\]
The assumed positive limit of $(1-r^2)^2\rho_1(r)$ implies
$(1-r^2)\E\Xi(B_r)\to\pi c$, by integrating upper and lower bounds
$(c\pm\varepsilon)/(1-s^2)^2$ on an outer annulus. Therefore
$\pi c\le1$.
\end{proof}

\begin{proposition}[The variance that remains at the critical value]
\label{prop:critical-score-variance}
For a nonempty limiting law, let
\[
 V(z)=\operatorname{Var}(b(z,\Xi)\mid z)
\]
be the conditional variance under disintegration of the Campbell
measure with respect to its spatial coordinate. On $\D\setminus\{0\}$,
the curvature bound strengthens to
\begin{equation}\label{eq:curvature-residual-variance}
 \Delta\log\rho_1
 \ge\left(\frac{2\beta}{(1-|z|^2)^2}+\beta^2V(z)\right)dA(z).
\end{equation}
At $\beta=4$, for every $r_0\in(0,1)$,
\begin{equation}\label{eq:critical-score-integrability}
 \int_{r_0<|z|<1}(1-|z|)V(z)\,dA(z)<\infty.
\end{equation}
Thus a lower bound $V(z)\ge c(1-|z|)^{-2}$ almost everywhere on an
entire outer annulus, for any $c>0$, would rule out a nonempty critical
limit. The proof below uses boundary translation and repulsion to obtain
the contradiction without first establishing this particular lower bound.
\end{proposition}

\begin{proof}
In the proof of Theorem~\ref{thm:limiting-intensity-curvature}, retain
the conditional variance instead of discarding it in Jensen:
\[
 \beta^2\rho_1\E[|b|^2\mid z]
 =\frac{|\nabla\rho_1|^2}{\rho_1}+\beta^2\rho_1 V.
\]
The same weak chain rule gives
\eqref{eq:curvature-residual-variance}.

At $\beta=4$ put $h(r)=\log[(1-r^2)^2\rho_1(r)]$.
Corollary~\ref{cor:force-empty-critical} shows that $h$ has a finite
limit at $1$, and $h$ is radial subharmonic on the punctured disk.
Let $g(t)=h(e^t)$, which is convex for $t<0$. Denote its distributional
second-derivative measure by $\mu_g$, which is nonnegative. For any fixed $T<0$ and
$T<t<0$, the integral representation of a convex function gives
\[
 g(t)=g(T)+(t-T)g'_+(T)
              +\int_{(T,t]}(t-s)\,d\mu_g(s).
\]
Letting $t\uparrow0$ and using the
finite boundary value proves
$\int_{(T,0)}(-s)\,d\mu_g(s)<\infty$.
The radial Laplacian relation gives
\[
 \int_{e^T<|z|<1}\log(1/|z|)\,d(\Delta h)(z)
 =2\pi\int_{(T,0)}(-s)\,d\mu_g(s)<\infty.
\]
Since $\Delta h\ge16VdA$ and
$\log(1/r)\ge1-r$, this implies
\eqref{eq:critical-score-integrability}. The stated hypothetical
lower bound would make its integral diverge logarithmically.
\end{proof}


\section{Uniform local disk bounds and boundary blowup compactness}
\label{sec:local-disk-domination}

The radial comparison is local in a useful sense: conditioning on the
exterior of any disk produces a holomorphic tilt of the same disk gas.
It therefore supplies estimates independent of both the location and
the radius of that disk.

\begin{proposition}[Domination in every interior disk]
\label{prop:local-disk-domination}
Let $B(a,R)\subset\D$ be an open Euclidean disk, and let $0<q<1$.
For every $\beta>0$, every $N$, and every integer $n\ge2$,
\begin{equation}\label{eq:local-disk-count-tail}
 \mathbb P\{\Xi_N(B(a,qR))\ge n\}
 \le \min\left\{1,
   \exp\left(n-\frac\beta4 n(n-1)\log\frac1q\right)\right\}.
\end{equation}
The same estimate holds conditionally on the exterior configuration
and the number of particles in $B(a,R)$, for almost every conditioning
configuration. It also holds with $\Xi_N$ replaced by any subsequential
limit $\Xi$. In particular, for every $t\ge0$ and $s>0$ there are finite
constants depending only on $\beta,q,t,s$ such that
\[
 \sup_{N,a,R}\E e^{t\Xi_N(B(a,qR))}<\infty,
 \qquad
 \sup_{N,a,R}\E\Xi_N(B(a,qR))^s<\infty,
\]
and the same uniform bounds hold for all subsequential limits.
\end{proposition}

\begin{proof}
Fix the disk and condition on its exterior configuration
$y_1,\ldots,y_{N-m}$ and the inside count $m$. The joint density is
absolutely continuous, so almost surely no particle lies on the
conditioning circle. Consequently $|y_j-a|>R$ for every exterior
particle. After rescaling the inside coordinates as $z_i=a+R u_i$,
their conditional density on $\D^m$ is proportional to
\[
 |\Delta_m(u)|^\beta
 \prod_{i=1}^m\prod_{j=1}^{N-m}|a+Ru_i-y_j|^\beta.
\]
Each function $u\mapsto a+Ru-y_j$ is nonzero on a neighborhood of
$\overline\D$ and has a holomorphic logarithm there, possibly after
shrinking the neighborhood. There are finitely many exterior particles,
so a common neighborhood suffices. Hence the second factor is
$\exp(2\operatorname{Re}\sum_i f(u_i))$, where
\[
 f(u)=\frac\beta2\sum_{j=1}^{N-m}\log(a+Ru-y_j)
\]
is holomorphic on that neighborhood. Subtracting $f(0)$ does not change
the normalized conditional law. The radial comparison of
Lemma~\ref{lem:analytic-tilt-radial} applies with constants independent
of this holomorphic function and of $m$. The exponential-gap calculation
leading to \eqref{eq:gaussian-tail} therefore proves
\eqref{eq:local-disk-count-tail} conditionally. The case $m<n$ is
automatic. Integrating the conditional inequality proves the
unconditional one.

For a subsequential limit, the map
$\eta\mapsto\eta(B(a,qR))$ is lower semicontinuous for the vague
topology, because the disk is open and relatively compact in $\D$.
Its integer-valued upper-level event $\{\eta(B(a,qR))\ge n\}$ is open.
The portmanteau inequality transfers the same upper bound to $\Xi$.
Summation of the displayed Gaussian tail proves all the claimed
moment bounds. This argument applies to every deterministic disk
contained in $\D$, so its constants remain uniform in its center and
radius.
\end{proof}

\begin{theorem}[Boundary blowup compactness and scaled repulsion]
\label{thm:boundary-blowup-compactness}
Let $\Xi$ be any subsequential limit at inverse temperature $\beta>0$,
and let
\[
 \mathbb H=\{w\in\mathbb C:\operatorname{Re}w<0\}.
\]
For $0<\delta<1/2$ define the trimmed boundary blowup
\begin{equation}\label{eq:trimmed-boundary-blowup}
 \widehat\Xi_\delta
 =\sum_{\substack{z\in\Xi\\ |z|<1-\delta^2}}
                  \delta_{(z-1)/\delta}.
\end{equation}
These are finite point processes on $\mathbb H$, and their laws are tight
as $\delta\downarrow0$. For every compact $K\Subset\mathbb H$ there
are $C_K,c_K>0$ such that
\begin{equation}\label{eq:boundary-blowup-count-tail}
 \sup_{0<\delta<1/2}
 \mathbb P\{\widehat\Xi_\delta(K)\ge n\}
 \le C_K e^{-c_K n^2},\qquad n\ge0.
\end{equation}
In particular all local count moments and exponential moments are
uniformly bounded.

Write $\alpha=\min\{1,\beta\}$ and let $h_p$ be the normalized
correlations of $\Xi$ from Theorem~\ref{thm:hierarchy}. For each compact
$K\Subset\mathbb H$ and each $p\ge1$, there are $C_{K,p}<\infty$ and
$\delta_{K,p}>0$ such that, for $0<\delta<\delta_{K,p}$, the functions
\[
 H_{p,\delta}(w)
 =\delta^{2p+\beta p(p-1)/2}
            h_p(1+\delta w_1,\ldots,1+\delta w_p)
\]
are bounded by $C_{K,p}$ and have $\alpha$-H\"older seminorm at most
$C_{K,p}$ on $K^p$. In particular,
\begin{equation}\label{eq:scaled-vandermonde-bound}
 \delta^{2p}\rho_p(1+\delta w_1,\ldots,1+\delta w_p)
 \le C_{K,p}|\Delta_p(w)|^\beta,
 \qquad w\in K^p.
\end{equation}
These are exactly the correlations of $\widehat\Xi_\delta$ on $K^p$
for all sufficiently small $\delta$.

Every boundary blowup subsequential limit is simple and has continuous
correlation densities $\varrho_p=|\Delta_p|^\beta H_p$, where $H_p$ is
nonnegative, plurisubharmonic and locally $\alpha$-H\"older. Along a
process-convergent blowup subsequence, both normalized and ordinary
correlations converge locally uniformly. Convergence also holds in
total variation on every relatively compact window.
\end{theorem}

\begin{proof}
First record a minor issue concerning the ambient space. The literal
rescaling of an infinite configuration in $\D$ need not be locally
finite on all of $\mathbb H$: its particles may accumulate at other
points of the rescaled circular boundary, which lie inside this
half-plane. The trimming in \eqref{eq:trimmed-boundary-blowup} avoids
this issue. It removes particles in a region that is eventually
disjoint from every fixed compact subset of $\mathbb H$ after scaling.
Indeed, for compact $K\Subset\mathbb H$ choose $d>0$ and $L<\infty$
with $\operatorname{Re}w\le-d$ and $|w|\le L$ on $K$. Then
\[
 |1+\delta w|^2\le1-2d\delta+L^2\delta^2
                   <(1-\delta^2)^2
\]
for sufficiently small $\delta$. Thus the trimmed and literal
rescalings agree on $K$ eventually. Equivalently, the theorem gives
local convergence on compact subsets of the expanding domains
$\{w:|1+\delta w|<1\}$ without choosing a trim.

Every compact subset of $\mathbb H$ is contained in a Euclidean disk
whose closure lies in $\mathbb H$. To verify this, place the center at
$-L_0$ on the real axis, choose $0<\varepsilon<d$, and take radius
$L_0-\varepsilon$. The inequality
\[
 |w+L_0|^2-(L_0-\varepsilon)^2
 =2L_0(\operatorname{Re}w+\varepsilon)+|w|^2-\varepsilon^2<0
\]
holds on the compact set when $L_0$ is large. We may therefore choose
concentric disks
\[
 K\subset B(c,r),\qquad
 0<r<s<t<R,\qquad \overline{B(c,R)}\subset\mathbb H.
\]
For sufficiently small $\delta$, the physical disk
$B(1+\delta c,\delta R)$ lies in $\D$, since its rescaled closure is a
compact subset of $\mathbb H$. Proposition~\ref{prop:local-disk-domination},
with $q=r/R$, bounds the count in the physical image of $K$ by a
Gaussian tail independent of $\delta$. Trimming only decreases counts.
For the remaining $\delta\ge\delta_0>0$, all retained particles lie
in $\{|z|<1-\delta_0^2\}$, whose count has a Gaussian tail by the
original radial estimate. Combining the two bounds proves
\eqref{eq:boundary-blowup-count-tail}. Compact exhaustion and the usual
point-measure tightness criterion now prove tightness on $\mathbb H$.

We prove the correlation bounds first at finite $N$. On the rescaled
domain define
\[
 H_{p,N,\delta}(w)
 =\delta^{2p+\beta p(p-1)/2}
       h_{p,N}(1+\delta w_1,\ldots,1+\delta w_p).
\]
These functions are nonnegative and plurisubharmonic. The insertion
formula \eqref{hier:h-definition} also writes each of them as
an integral $\int|F_y(w)|^\beta\,d\mu_{N,\delta}(y)$ with positive
finite measure and holomorphic polynomial $F_y$. Their normalization
is chosen so that
\begin{align*}
 |\Delta_p(w)|^\beta H_{p,N,\delta}(w)
 &=\delta^{2p}\rho_{p,N}(1+\delta w_1,\ldots,1+\delta w_p),\\
 \int_{B(c,t)^p}|\Delta_p(w)|^\beta H_{p,N,\delta}(w)\,dA^p(w)
 &=\E\bigl(\Xi_N(B(1+\delta c,\delta t))\bigr)_p.
\end{align*}
The right side is bounded uniformly in $N$ and small $\delta$ by
Proposition~\ref{prop:local-disk-domination}, applied inside the larger
physical disk of radius $\delta R$.

Translate and rescale the separated-annular-average inequality of
Lemma~\ref{hier:annular-lemma} from the unit disk to the fixed nested
disks $B(c,s)\subset B(c,t)$. All change-of-variable constants depend
only on these fixed disks and on $p,\beta$. It follows that
$H_{p,N,\delta}$ is uniformly bounded on $B(c,s)^p$. Applying this
argument on slightly larger intermediate disks as needed gives a
bound on a neighborhood of that compact polydisk. For completeness,
the subsequent H\"older estimate uses only this bound and the positive
holomorphic-integral representation. Submean and Cauchy estimates give
\[
 \int\left(\sup_{B(c,r)^p}|F_y|^\beta
       +\sum_{j=1}^p\sup_{B(c,r)^p}|\partial_jF_y|^\beta\right)
       d\mu_{N,\delta}(y)\le C_{K,p}.
\]
One can use an intermediate radius between $r$ and $s$ for the
suprema to control differences on the closed smaller polydisk.
For $\beta\le1$, use
$\big||u|^\beta-|v|^\beta\big|\le|u-v|^\beta$; for $\beta>1$, use
$\big||u|^\beta-|v|^\beta\big|
\le\beta|u-v|(|u|+|v|)^{\beta-1}$ and H\"older's inequality.
Integration then gives the asserted uniform exponent $\alpha$,
exactly as in Lemma~\ref{hier:holder-lemma}. This proves all the bounds
uniformly at finite $N$ and at small scales.

Now hold $\delta$ fixed and take the original process-convergent
subsequence $N_j\to\infty$. The locally uniform correlation convergence
of Theorem~\ref{thm:hierarchy}, on the fixed physical compact disks,
passes these same bounds to $H_{p,\delta}$ and proves
\eqref{eq:scaled-vandermonde-bound}. The initial observation about the
trim identifies these densities with those of $\widehat\Xi_\delta$
on every fixed compact window once $\delta$ is sufficiently small.

Finally let a sequence $\delta_j\downarrow0$ have a process limit on
$\mathbb H$. The uniform H\"older bounds give locally uniform
subsequential limits of each $H_{p,\delta_j}$. The factorial-measure
partition identity \eqref{hier:partition-formula} and the uniform
compact count moments identify the resulting densities with the
factorial moment measures of the limiting point process. This
identifies each normalized density away from the diagonals, and
continuity identifies it everywhere; hence convergence holds along
the full process-convergent subsequence. Nonnegativity,
plurisubharmonicity and the H\"older estimates persist under local
uniform convergence. The second factorial density assigns zero mass
to the diagonal, so the limiting process is simple. Finally the Janossy
expansion and proof of Theorem~\ref{hier:local-tv} use only local uniform
correlation convergence and exponential compact count bounds, both of
which have just been established. They therefore give the stated local
total variation convergence.
\end{proof}


\begin{lemma}[Deterministic logarithmic derivatives are incompatible with repulsion]
\label{lem:deterministic-score-rigidity}
Let $U\subset\R^2$ be connected and open. Suppose a simple point
process on $U$ has all local count moments, a strictly positive smooth
intensity $\lambda$, and a continuous second correlation density
$\rho_2$ that vanishes on the diagonal. Assume its Ward identity is
\begin{equation}\label{eq:deterministic-score-ward}
 \E\left[\mathcal L_vF(\Xi)
   +F(\Xi)\Xi\bigl(\operatorname{div}v
                           +v\cdot\nabla\log\lambda\bigr)\right]=0
\end{equation}
for every compactly supported $C^1$ real vector field and every bounded
smooth local cylinder function $F$. Then $\rho_2\equiv0$, so the
process has at most one particle almost surely. In particular, no
such process exists if $\int_U\lambda\,dA>1$.
\end{lemma}

\begin{proof}
The local moment assumptions extend the identity to $F(\eta)=\eta(h)$
for every real $h\in C_c^1(U)$. Explicitly, approximate the identity
function on $\R$ by smooth bounded functions $\chi_M$ that agree with
it on $[-M,M]$, satisfy $|\chi_M(s)|\le C|s|$ and
$|\chi'_M(s)|\le C$, and converge with their derivatives pointwise.
Apply the assumed identity to $\chi_M(\Xi(h))$ and use domination by
a fixed compact count for the derivative term and by the square of
a compact count for the other term.

Put $q=\operatorname{div}v+v\cdot\nabla\log\lambda$.
The one- and two-point formulas yield
\[
 0=\int_U\nabla h(a)\cdot v(a)\lambda(a)\,dA(a)
   +\int_Uh(a)q(a)\lambda(a)\,dA(a)
   +\iint_{U^2}h(w)q(a)\rho_2(a,w)\,dA(a)dA(w).
\]
The first two integrals sum to
$\int\operatorname{div}(h v\lambda)dA=0$.
As $h$ and $v$ are arbitrary, the remaining integral states
distributionally that
\[
 \nabla_a\rho_2(a,w)
 =\rho_2(a,w)\nabla\log\lambda(a).
\]
Interchanging the two coordinates gives the analogous identity in $w$.
Dividing by the smooth positive function $\lambda(a)\lambda(w)$ shows
that the continuous function
$Q(a,w)=\rho_2(a,w)/(\lambda(a)\lambda(w))$ has all distributional
first derivatives zero on $U^2$. Such a function is constant on a
connected open set: convolution makes it locally a smooth function
with zero gradient, and overlapping balls identify the constants.
Since $U^2$ is connected, $Q$ is constant there. Its value on the
diagonal is zero, hence $\rho_2\equiv0$.

Every relatively compact set $K$ consequently satisfies
$\E[\Xi(K)(\Xi(K)-1)]=0$, and a countable exhaustion gives
$\Xi(U)\le1$ almost surely. Its mean is therefore at most one, in
contradiction with $\int_U\lambda\,dA>1$ when that condition holds.
\end{proof}


\section{The critical gas also converges to the empty process}
\label{sec:critical-force-empty}

\begin{theorem}[Empty limit at the critical inverse temperature]
\label{thm:critical-empty}
At $\beta=4$, every subsequential limit of $\Xi_N$ is almost surely
empty. Consequently the full sequence converges to the empty
configuration, vaguely in law and locally in total variation.
\end{theorem}

\begin{proof}
Suppose instead that a subsequential limit $\Xi$ is not almost surely
empty, and write $\rho$ for its intensity. Let $b$ be its canonical
complex score, and let $m(z)=\E[b(z,\Xi)\mid z]$ and
$V(z)=\E[|b(z,\Xi)-m(z)|^2\mid z]$ denote the conditional mean and
variance under its Campbell measure. We use the conclusions already
proved at $\beta=4$:
\begin{align}
 &(1-r^2)^2\rho(r)\longrightarrow c\in(0,1/\pi],
 \label{eq:critical-proof-intensity}\\
 &h(r):=\log[(1-r^2)^2\rho(r)]
       \text{ is nondecreasing and }h(e^t)\text{ is convex},
 \label{eq:critical-proof-convexity}\\
 &\int_{r_0<|z|<1}(1-|z|)V(z)\,dA(z)<\infty
       \quad (r_0>0),
 \label{eq:critical-proof-variance-integrability}\\
 &m(z)=\frac2\beta\partial\log\rho(z)
       \quad\text{almost everywhere where }\rho>0.
 \label{eq:critical-proof-mean}
\end{align}
The last identity is the one-point consequence of the canonical Ward
identity. Rotation invariance and uniqueness of the canonical score
give
$b(e^{i\theta}z,(e^{i\theta})_*\eta)=e^{-i\theta}b(z,\eta)$
up to Campbell-null sets: the rotated expression satisfies the same
Ward identity, which determines the score uniquely. Consequently
$V(z)=V(|z|)$ almost everywhere.

\emph{Boundary rescaling.}
Set $H=\{w=x+iy:x>0\}$ and, for $\varepsilon>0$, put
\[
 T_\varepsilon(z)=\frac{1-z}{\varepsilon},\qquad
 \eta_\varepsilon=\sum_{z\in\Xi:\ |z|\le1-\varepsilon^2}
                                  \delta_{T_\varepsilon(z)}.
\]
The trimmed configuration is finite almost surely and hence defines a
point process on all of $H$. This trim is necessary: an untrimmed image
could accumulate on its circular boundary inside $H$. Every fixed
compact subset of $H$ lies inside $T_\varepsilon(B_{1-\varepsilon^2})$
for sufficiently small $\varepsilon$, so trimming changes none of the
local statistics used below. Theorem~\ref{thm:boundary-blowup-compactness},
with the reflected coordinate $w\mapsto-w$, gives tightness of $\eta_\varepsilon$ on $H$, uniform moments
of every compact count, and locally uniform correlation convergence
along subsequences. Its subsequential limits are simple and have
continuous second correlation density vanishing on the diagonal.
Indeed these are exactly the estimates applied after this translation
and dilation, with constants depending only on the chosen compact
subset of $H$.

By \eqref{eq:critical-proof-intensity}, the rescaled intensities obey
\begin{equation}\label{eq:critical-halfplane-intensity}
 \rho_\varepsilon(w)=\varepsilon^2\rho(1-\varepsilon w)
 \longrightarrow \lambda(w):=\frac{c}{4x^2}
 \quad\text{locally uniformly on }H.
\end{equation}
Here
$1-|1-\varepsilon w|^2=2\varepsilon x-\varepsilon^2|w|^2$.
In particular $\rho_\varepsilon$ is uniformly bounded on each compact
subset of $H$, and every rescaled limiting process has intensity
$\lambda$.

\emph{The rescaled score becomes deterministic.}
Attach to each retained point the transformed original score and its
spatial conditional mean:
\[
 b_\varepsilon(w)
       =-\varepsilon b(1-\varepsilon w,\Xi),\qquad
 m_\varepsilon(w)=-\varepsilon m(1-\varepsilon w).
\]
The marks $b_\varepsilon$ are defined on the original probability space;
they are not asserted to be functions of the trimmed configuration.
The minus sign is the derivative of $z=1-\varepsilon w$.
For a compact set $K\Subset H$, take constants $a,B>0$ such that
$a\le x\le B$ and $|y|\le B$ on $K$. For small $\varepsilon$,
$r=|1-\varepsilon w|$ lies between
$1-C_K\varepsilon$ and $1-c_K\varepsilon$, and
\[
 \left|\frac{\partial r}{\partial x}\right|
 =\varepsilon\frac{1-\varepsilon x}{r}
 \in[\varepsilon/2,2\varepsilon].
\]
Using rotation invariance of $V$, Campbell's formula, and the uniform
intensity bound on $K$, gives
\begin{align}
 &\E\sum_{w\in\eta_\varepsilon\cap K}
           |b_\varepsilon(w)-m_\varepsilon(w)|^2
 \nonumber\\
 &\qquad=\int_K\rho_\varepsilon(w)\varepsilon^2
                                     V(|1-\varepsilon w|)\,dA(w)
 \nonumber\\
 &\qquad\le C_K\varepsilon
           \int_{1-C_K\varepsilon}^{1-c_K\varepsilon}V(r)\,dr
 \le C'_K\int_{1-C_K\varepsilon}^{1-c_K\varepsilon}
                                      (1-r)V(r)\,dr
 \longrightarrow0.
 \label{eq:critical-score-residual-vanishes}
\end{align}
The last step follows from
\eqref{eq:critical-proof-variance-integrability}; $r$ is bounded away
from zero on these annuli.

We next show, in essential supremum on every compact subset of $H$,
\begin{equation}\label{eq:critical-score-mean-convergence}
 m_\varepsilon(w)\longrightarrow
       m_\infty(w):=-\frac2{\beta x}.
\end{equation}
Let $\ell=\lim_{r\uparrow1}h(r)=\log c$. Convexity of $h(e^t)$
implies, almost everywhere and with $r_+=(1+r)/2$,
\[
 0\le h'(r)
 \le\frac{h(r_+)-h(r)}{r\log(r_+/r)}
 \le\frac{C[\ell-h(r)]}{1-r}
\]
on an outer annulus. Therefore
$\varepsilon h'(|1-\varepsilon w|)\to0$ in essential supremum on $K$.
The weak derivative chain rule and
\eqref{eq:critical-proof-mean} give
\[
 m_\varepsilon(w)
 =-\frac{4\varepsilon\overline{(1-\varepsilon w)}}
            {\beta(1-|1-\varepsilon w|^2)}
   -\frac{\varepsilon\overline{(1-\varepsilon w)}}
            {\beta|1-\varepsilon w|}
                         h'(|1-\varepsilon w|).
\]
The first term tends uniformly to $-2/(\beta x)$ and the second tends
to zero essentially uniformly, proving
\eqref{eq:critical-score-mean-convergence}. Notice that
$\beta\overline{m_\infty(w)}$, as a real vector, equals
$\nabla\log\lambda(w)=(-2/x,0)$.

\emph{Passage of the Ward identity and contradiction.}
Choose a convergent subsequence of the rescaled processes and call its
limit $\eta$. For each compactly supported vector field $v$ on $H$
and bounded smooth local cylinder $F$, the canonical Ward identity
transforms exactly into the following identity on the original
probability space. For fixed local tests their spatial supports lie
strictly inside the trimming disk for all sufficiently small
$\varepsilon$, so the trim has no effect on their derivatives:
\[
 \E\left[\mathcal L_vF(\eta_\varepsilon)
 +F(\eta_\varepsilon)\eta_\varepsilon(\operatorname{div}v)
 +\beta F(\eta_\varepsilon)\operatorname{Re}
             \sum_{w\in\eta_\varepsilon}v(w)
                         b_\varepsilon(w)\right]=0
\]
for small enough $\varepsilon$. Its force term may be replaced by
the deterministic $m_\infty$. To see this for the centered remainder,
apply Cauchy--Schwarz to the Campbell measure over a compact set
containing $\supp v$:
\[
 \E\sum_{w\in\eta_\varepsilon\cap K}|b_\varepsilon-m_\varepsilon|
 \le\bigl(\E\eta_\varepsilon(K)\bigr)^{1/2}
      \left(\E\sum_{w\in\eta_\varepsilon\cap K}
                      |b_\varepsilon-m_\varepsilon|^2\right)^{1/2}
 \longrightarrow0.
\]
The difference between $m_\varepsilon$ and $m_\infty$ is handled by
\eqref{eq:critical-score-mean-convergence} and the bounded mean count.
The remaining local statistics pass to the limit by uniform compact
count moments. Hence $\eta$ satisfies
\[
 \E\left[\mathcal L_vF(\eta)
   +F(\eta)\eta\bigl(\operatorname{div}v
                          +v\cdot\nabla\log\lambda\bigr)\right]=0.
\]
The deterministic-score rigidity lemma applies: $\eta$ has all local
count moments, positive smooth intensity $\lambda$, and a continuous
repulsive second correlation density. It follows that $\eta(H)\le1$
almost surely. This is impossible because
$\int_H c/(4x^2)\,dA=\infty$; equivalently, a sufficiently long compact
rectangle with $1\le x\le2$ already has expected count greater than one.
The contradiction proves that no nonempty critical subsequential
limit exists.

Finally tightness gives full vague convergence to the empty process.
For every compact $K$, compact-count uniform integrability and
convergence to zero imply $\E\Xi_N(K)\to0$. Thus
$\Pbb\{\Xi_N(K)>0\}\to0$, which is exactly convergence in total
variation of the restriction to the empty law. No quantitative rate
at $\beta=4$ is asserted.
\end{proof}


\begin{corollary}[Qualitative consequences at and above the endpoint]
\label{cor:critical-empty-consequences}
For every fixed $\beta\ge4$ and every integer $p\ge1$,
\[
 \rho_{p,N}\longrightarrow0,\qquad h_{p,N}\longrightarrow0
 \quad\text{locally uniformly on }\D^p.
\]
For every relatively compact Borel set $A\subset\D$, every $s>0$,
and every real $t$,
\[
 \E\Xi_N(A)^s\longrightarrow0,\qquad
 \E e^{t\Xi_N(A)}\longrightarrow1.
\]
For every $q>1$, the boundary-weight statistic
$U_{q,N}=\sum_i(1-|z_i|)^q$ converges to zero in every positive
moment, and $\E e^{tU_{q,N}}\to1$ for every real $t$.
No rate in these statements is asserted at $\beta=4$.
\end{corollary}

\begin{proof}
The full process limit is empty by Theorem~\ref{thm:critical-empty}
and the theorem above $4$. The locally uniform hierarchy convergence
of Theorem~\ref{thm:hierarchy} applies to this full convergent sequence;
the empty process has zero factorial densities. The continuous
normalized densities are zero as well, since they vanish away from
the collision diagonals, a dense subset of $\D^p$.

Choose a compactly supported continuous function at least one on
$\overline A$. Its statistic tends to zero in probability, so
$\Xi_N(A)\to0$ in probability. Uniform compact count moments of
all orders and uniform exponential moments at every real parameter
give uniform integrability of each displayed random variable, proving
the count assertions. For example, the exponential moment at $2t$
controls uniform integrability when $t>0$; for $t\le0$ the integrand
is bounded.

For $U_{q,N}$, the boundary truncation estimates and limiting-statistic
theorem in Section~\ref{sec:weighted-boundary-tilts} give convergence
to the corresponding statistic of the empty process, namely zero.
Their uniform exponential bounds give uniform integrability of every
power and every fixed exponential, proving the remaining claims.
\end{proof}


\section{Critical emptiness through disk automorphisms}
\label{sec:critical-mobius-empty}

This proof uses the exact Bergman bounds from
Theorem~\ref{thm:even-beta-bergman-bound} and
Corollary~\ref{cor:even-beta-determinantal-upper-bound}. The other inputs
are the canonical logarithmic derivative, the critical residual-variance
estimate, and the deterministic-score rigidity lemma proved earlier.
All limiting processes in the argument live on the same open disk.

\begin{theorem}[Uniqueness at the critical value]
\label{thm:critical-mobius-empty}
At $\beta=4$, the unscaled disk gas converges along the full sequence to
the empty point process. On every compact window the convergence is in
total variation.
\end{theorem}

\begin{proof}
Let $\Xi$ be any subsequential limit. We show that it is empty.
Suppose otherwise, and write $\rho$ for its one-point intensity.
The already established critical curvature consequence says that
\[
 c(r)=(1-r^2)^2\rho(r)>0,\qquad
 g(t)=\log c(e^t),\quad -\infty<t<0,
\]
is nondecreasing and convex as a function of $t$. The Bergman bound
gives $c(r)\leq1/\pi$. Thus
\begin{equation}
 c(r)\longrightarrow c_*\in(0,1/\pi]
 \quad\text{as }r\uparrow1.
 \label{eq:critical-mobius-positive-boundary-coefficient}
\end{equation}
Only radii sufficiently close to one are needed in the displays below;
the same argument applies if positivity and the convexity assertion
were initially formulated only on such an outer annulus.

Let $b(z,\eta)$ be the canonical complex logarithmic derivative from
Corollary~\ref{cor:canonical-force-score}, and put
$B(z,\eta)=4\overline{b(z,\eta)}$, regarded as a real two-dimensional
vector. Disintegrate the Campbell measure over $z$ and let
$V(z)=\operatorname{Var}(b(z,\Xi)\mid z)$.
The canonical identity with constant cylinder test implies
\[
 \mathbb E[B(z,\Xi)\mid z]=\nabla\log\rho(z).
\]
The residual-variance result gives, for some $r_0<1$,
\begin{equation}
 \int_{r_0<|z|<1}(1-|z|)V(z)\,dA(z)<\infty.
 \label{eq:critical-mobius-residual-input}
\end{equation}
We may choose $V$ radial up to area-null sets. Indeed rotations preserve
the law, and transforming the canonical score under a rotation gives
another score for the same law. Its uniqueness identifies the two
versions up to Campbell-null sets, so their conditional variances agree.
Since $\rho$ is positive on the annulus, this is also an area-almost-everywhere
assertion there.

For real $a\uparrow1$, define
\[
 \varphi_a(w)=\frac{a+w}{1+aw},\qquad
 \Xi_a=(\varphi_a^{-1})_*\Xi,\qquad \varepsilon=1-a.
\]
The transformed intensity is exactly
\begin{equation}
 \rho_a(w)=|\varphi_a'(w)|^2\rho(\varphi_a(w))
 =\frac{c(|\varphi_a(w)|)}{(1-|w|^2)^2}
 \longrightarrow\lambda(w):=\frac{c_*}{(1-|w|^2)^2},
 \label{eq:critical-mobius-intensity-limit}
\end{equation}
locally uniformly. Here we used
$1-|\varphi_a(w)|^2=(1-a^2)(1-|w|^2)/|1+aw|^2$.
The exact all-correlation Bergman bounds remain valid for $\Xi_a$.
Consequently these processes are tight, all their compact count moments
are uniformly bounded, and along a sequence $a\uparrow1$ they converge
to a process $\Upsilon$ with intensity $\lambda$.
Their two-point bounds also give uniform repulsion of order four on
compact sets.

We record why the limiting second correlation may be taken continuous,
as required by the score-rigidity lemma. At finite $N$, each even-temperature
correlation is the diagonal of a positive holomorphic kernel: replace
the second copy of the marked coordinates in the squared-polynomial
insertion integral by independent complex coordinates. Cauchy--Schwarz
bounds that kernel at $(\mathbf z,\overline{\mathbf w})$ by
$\sqrt{\rho_{p,N}(\mathbf z)\rho_{p,N}(\mathbf w)}$.
The Bergman diagonal bounds make these kernels locally uniformly bounded.
Montel's theorem and Cauchy estimates therefore give locally convergent
holomorphic kernels along subsequences, both in the original
$N\to\infty$ limit and after the subsequent automorphisms. In the latter
case the transformed kernel acquires the holomorphic factor
$\prod_i\varphi_a'(z_i)\overline{\varphi_a'(w_i)}$.
Thus the limiting correlations are real analytic on their diagonals,
in particular continuous. Their integrals agree with the factorial
moments of the limiting process by uniform compact count moments.
The repulsion bound passes to this continuous version, giving
$\rho_{2,\Upsilon}(w,w)=0$. It also ensures simplicity of $\Upsilon$.

The canonical real score for $\Xi_a$ is
\begin{equation}
 B_a(w,\eta)=D\varphi_a(w)^T
       B(\varphi_a(w),(\varphi_a)_*\eta)
       +\nabla\log|\varphi_a'(w)|^2.
 \label{eq:critical-mobius-score-transform}
\end{equation}
This follows directly by changing variables in the Campbell Ward
identity; the Jacobian contributes the second term. Its conditional mean
is $d_a=\nabla\log\rho_a$, and conformality gives the centered variance
\[
 \mathbb E[|B_a-d_a|^2\mid w]
       =16|\varphi_a'(w)|^2V(\varphi_a(w)).
\]

For every compact $K\Subset\mathbb D$, elementary bounds on the displayed
automorphism give positive constants $c_K,C_K$ such that on $K$,
\[
 c_K\varepsilon\leq1-|\varphi_a(w)|\leq C_K\varepsilon,
 \qquad |\varphi_a'(w)|\leq C_K\varepsilon.
\]
Moreover $\varphi_a(K)$ is contained in a sector of angular width
$C_K\varepsilon$ about the point $1$: indeed
$\varphi_a(w)-1=(a-1)(1-w)/(1+aw)$.
Using the Bergman intensity bound and changing variables,
\begin{align*}
 \mathbb E\sum_{w\in\Xi_a\cap K}|B_a(w,\Xi_a)-d_a(w)|^2
 &=16\int_{\varphi_a(K)}
   |\varphi_a'(\varphi_a^{-1}(z))|^2\rho(z)V(z)\,dA(z)\\
 &\leq C_K\int_{\varphi_a(K)}V(z)\,dA(z)\\
 &\leq C_K\varepsilon
          \int_{1-C_K\varepsilon}^{1-c_K\varepsilon}V(r)\,dr
 \longrightarrow0.
\end{align*}
For the last convergence, $r$ is bounded away from zero and
$1-r$ is comparable to $\varepsilon$ on the integration interval;
thus the last expression is bounded by a fixed constant times the
integral of $(1-r)V(r)r\,dr$ over that interval. It tends to zero by
\eqref{eq:critical-mobius-residual-input} and absolute continuity of
the integral.

We also have locally uniform convergence of the mean scores,
\begin{equation}
 d_a(w)\longrightarrow d(w):=\nabla\log\lambda(w)
                      =\frac{4w}{1-|w|^2}.
 \label{eq:critical-mobius-mean-score-limit}
\end{equation}
Here is an estimate that does not require interchanging a limit and a
derivative. For a nondecreasing convex function $g$ with a finite limit
at zero, the convex secant inequality gives
\[
 0\leq(-t)g'_+(t)
 \leq2\big(g(t/2)-g(t)\big)\longrightarrow0,
 \qquad t\uparrow0.
\]
For $w$ in a fixed compact set, $-\log|\varphi_a(w)|$ is comparable
to $\varepsilon$, and $|\varphi_a'|=O_K(\varepsilon)$.
The chain rule in
$\log\rho_a=\log c(|\varphi_a|)-2\log(1-|w|^2)$ therefore shows that
the gradient of the first term tends uniformly to zero. The possible
one-sided formulation is harmless; the intensities are real analytic
where positive by the kernel argument above, so the ordinary derivative
is available on the annuli in use.

Let $F$ be a bounded smooth local cylinder function and $v$ a compactly
supported $C^1$ vector field. The Ward identity for $\Xi_a$ reads
\[
 \mathbb E\left[\mathcal L_vF(\Xi_a)
  +F(\Xi_a)\sum_{w\in\Xi_a}
       \{\operatorname{div}v(w)+v(w)\cdot B_a(w,\Xi_a)\}\right]=0.
\]
Replacing $B_a$ by $d_a$ changes this by a quantity tending to zero:
Campbell Cauchy--Schwarz bounds its absolute value by
\[
 \|F\|_\infty
 \left(\mathbb E\sum_{w\in\Xi_a}|v(w)|^2\right)^{1/2}
 \left(\mathbb E\sum_{w\in\Xi_a\cap\operatorname{supp}v}
                         |B_a-d_a|^2\right)^{1/2}.
\]
The first factor involving the process is uniformly bounded and the
second tends to zero. Uniform compact count moments, vague convergence,
and \eqref{eq:critical-mobius-mean-score-limit} then allow passage to
the limit in every remaining term. We obtain
\[
 \mathbb E\left[\mathcal L_vF(\Upsilon)
  +F(\Upsilon)\Upsilon\bigl(\operatorname{div}v
                                  +v\cdot\nabla\log\lambda\bigr)\right]=0.
\]
All hypotheses of Lemma~\ref{lem:deterministic-score-rigidity} now hold
on the connected disk. Hence $\Upsilon$ has at most one point almost
surely. This contradicts
$\int_{\mathbb D}\lambda\,dA=\infty$, since $c_*>0$.
The assumed nonempty subsequential limit cannot exist.

Finally, tightness of the original sequence and uniqueness of its only
possible subsequential limit imply full convergence to the empty process.
Uniform compact count moments imply convergence of compact count means
to zero, so the probability of a nonempty compact window tends to zero.
This is precisely total variation convergence on that window.
\end{proof}


\section{Empty or infinite limits for \texorpdfstring{$\beta\ge2$}{beta is at least 2}}
\label{sec:finite-count-exclusion}

The local force identity can be restricted to an event prescribing the
\emph{total} number of limiting particles. This restriction is useful
because total particle count is invariant under compactly supported
motions, although it is not a continuous function in the vague topology.

\begin{theorem}[Empty or infinite when $\beta\ge2$]
\label{thm:finite-count-exclusion}
Let $\beta\ge2$, and let $\Xi$ be any subsequential limit of the disk gas.
Then
\[
 \mathbb P\{\Xi(\D)\in\{0,\infty\}\}=1.
\]
Equivalently, every event $\{\Xi(\D)=m\}$ with an integer $m\ge1$ has
probability zero. This does not assert that the process is nonempty:
a mixture of the empty configuration and infinite configurations is
not excluded by the theorem.
\end{theorem}

\begin{proof}
Fix an integer $m\ge1$ and write $E_m=\{\Xi(\D)=m\}$.
We work on a joint limiting probability space carrying $\Xi$, its force
marks, and the Gaussian modes from
Theorem~\ref{thm:gaussian-force-regression}. It is enough to show that
$p_m=\mathbb P(E_m)>0$ is impossible when $\beta\ge2$.

\medskip\noindent
\emph{Local approximation of the total-count event.}
Choose smooth radial functions $\chi_\ell\in C_c^\infty(\D;[0,1])$
increasing pointwise to one, with $\chi_\ell=1$ on disks whose radii
increase to one. Such a choice follows, for example, by applying a
fixed nonincreasing smooth cutoff to $\ell(1-|z|^2)$ and taking its
complement. Choose $f\in C_c^\infty(\mathbb R;[0,1])$ with
$f(m)=1$ and support in $(m-1/2,m+1/2)$, and set
\[
 H_\ell(\eta)=f\left(\int\chi_\ell\,d\eta\right).
\]
Then $H_\ell$ is a bounded continuous cylinder function and
\[
 H_\ell(\eta)\longrightarrow\mathbf1_{\{\eta(\D)=m\}}
\]
for every locally finite configuration $\eta$, by monotone convergence
of $\int\chi_\ell\,d\eta$ to its integer-valued or infinite total mass.
For every fixed compactly supported vector field $v$, one has
$\mathcal L_vH_\ell=0$ for all sufficiently large $\ell$, because
$\chi_\ell$ is identically one on a neighborhood of $\supp v$.

\medskip\noindent
\emph{A restricted force-energy identity.}
Let $0\le\phi\in C_c^\infty(\D)$, and choose $\ell$ large enough that
$\chi_\ell=1$ on a neighborhood of $\supp\phi$. At finite $N$,
\begin{equation}\label{eq:restricted-finite-force-energy}
 \beta^2\E\left[H_\ell(\Xi_N)
                  \sum_i\phi(z_i)|C_{i,N}|^2\right]
 =\E\left[H_\ell(\Xi_N)\sum_i\Delta\phi(z_i)\right].
\end{equation}
Indeed, in the insertion variable $a$ in this neighborhood,
\[
 H_\ell\left(\delta_a+\sum_j\delta_{y_j}\right)
 =f\left(1+\sum_j\chi_\ell(y_j)\right),
\]
which is independent of $a$ and lies between zero and one. Multiplying
the insertion measure by this nonnegative bounded factor preserves
the polynomial distributional identity
\eqref{eq:polynomial-force-laplacian}. Integrating that identity against
$\phi$ proves \eqref{eq:restricted-finite-force-energy}, with exactly
the same collision-integrability justification as in
Proposition~\ref{prop:campbell-force-identities}.

Pass along the joint convergent subsequence with $\ell$ fixed.
Continuous compact mark cutoffs and the uniform compact count moments
justify convergence of each bounded-mark term on the left. Monotone
approximation of $|c|^2$ then gives a lower bound for its full limiting
expectation. The right side converges by bounded continuity of
$H_\ell$ and uniform integrability of compact particle counts. Thus
\[
 \beta^2\E\left[H_\ell(\Xi)
                \sum_{z\in\Xi}\phi(z)|C_z|^2\right]
 \le\E\left[H_\ell(\Xi)\sum_{z\in\Xi}\Delta\phi(z)\right].
\]
Let $\ell\to\infty$. Dominated convergence applies on the left by the
unrestricted local force-energy bound and on the right by the local
count bound. We obtain
\begin{equation}\label{eq:restricted-limiting-force-energy}
 \beta^2\E\left[\mathbf1_{E_m}
                \sum_{z\in\Xi}\phi(z)|C_z|^2\right]
 \le\E\left[\mathbf1_{E_m}\sum_{z\in\Xi}\Delta\phi(z)\right].
\end{equation}

Define the restricted intensity $u_m$ by
\[
 \E\left[\mathbf1_{E_m}\Xi(da)\right]=u_m(a)\,dA(a).
\]
This density exists and satisfies $0\le u_m\le\rho_1$ almost everywhere,
because its intensity measure is dominated by the unrestricted one.
In particular it is locally bounded, and
\begin{equation}\label{eq:restricted-intensity-mass}
 \int_\D u_m\,dA=m p_m<\infty.
\end{equation}
The measure and its density are rotation invariant.

\medskip\noindent
\emph{Restricted Ward identity and Gaussian force variance.}
Apply the unmarked Ward identity for the canonical score $b$ to the
cylinder function $H_\ell$ and a fixed compactly supported $v$.
For all large $\ell$ its derivative term vanishes. The local Campbell
$L^2$ bound for $b$ and the count bound imply absolute integrability,
so dominated convergence gives
\begin{equation}\label{eq:restricted-score-ward}
 \E\left[\mathbf1_{E_m}\sum_{z\in\Xi}
    \left(\operatorname{div}v(z)
                  +\beta\operatorname{Re}(v(z)b(z,\Xi))\right)\right]=0.
\end{equation}
Disintegrate the Campbell measure restricted to $E_m$ over its spatial
marginal $u_m(a)dA(a)$. Denote the conditional mean of $b$ by
$\mathfrak b_m(a)$, and set $J_m=u_m\mathfrak b_m$. Formula
\eqref{eq:restricted-score-ward} says
\[
 \partial_xu_m=\beta\operatorname{Re}J_m,
 \qquad \partial_yu_m=-\beta\operatorname{Im}J_m
\]
in the sense of distributions. Conditional Jensen and local boundedness
of $u_m$ show $J_m\in L^2_{\rm loc}$; hence $u_m\in H^1_{\rm loc}$.
Moreover,
\begin{equation}\label{eq:restricted-fisher-score}
 \int\phi\frac{|\nabla u_m|^2}{u_m}\,dA
 \le\beta^2\E\left[\mathbf1_{E_m}
                    \sum_{z\in\Xi}\phi(z)|b(z,\Xi)|^2\right],
\end{equation}
where the quotient is assigned zero on $\{u_m=0\}$. The gradient is
zero almost everywhere on that zero set, as also follows directly
from the displayed distributional derivative formulas.

The conditional regression
\eqref{eq:gaussian-force-regression} may be multiplied by
$\mathbf1_{E_m}$, since this event is determined by $\Xi$.
The Gaussian modes remain independent of the restricted unmarked
Campbell measure. Conditional Jensen, followed by Gaussian integration
of the square and then monotone convergence in the number of modes,
therefore gives
\begin{align*}
 &\E\left[\mathbf1_{E_m}
                  \sum_{z\in\Xi}\phi(z)|C_z|^2\right]\\
 &\qquad\ge
 \E\left[\mathbf1_{E_m}
                  \sum_{z\in\Xi}\phi(z)|b(z,\Xi)|^2\right]
       +\frac2\beta\int\frac{\phi(a)u_m(a)}{(1-|a|^2)^2}\,dA(a).
\end{align*}
Combining this bound with
\eqref{eq:restricted-limiting-force-energy} and
\eqref{eq:restricted-fisher-score} proves the measure inequality
\begin{equation}\label{eq:restricted-intensity-curvature}
 \Delta u_m\ge\frac{|\nabla u_m|^2}{u_m}\,dA
                    +\frac{2\beta u_m}{(1-|a|^2)^2}\,dA.
\end{equation}
In particular $\Delta u_m$ is a positive Radon measure.

\medskip\noindent
\emph{Boundary growth contradicts finite mass.}
There is a continuous radial representative of $u_m$ on every annulus
compactly contained in the punctured disk. To see this without assuming
regularity of the event $E_m$, write $U_m(t)=u_m(e^t)$, $t<0$.
Radial $H^1$ regularity implies $U_m\in H^1_{\rm loc}(( -\infty,0))$,
so $U_m$ has a continuous representative. The positivity of
$\Delta u_m$, tested against radial functions $\zeta(\log|a|)$,
implies $U_m''\ge0$ distributionally. Thus $U_m$ is convex.
Since $u_m\le\rho_1$ almost everywhere and $\rho_1$ is bounded near
zero, $U_m$ is bounded above as $t\to-\infty$. Convexity then forces
$U_m$ to be nondecreasing: a negative supporting slope would contradict
that upper bound. These representatives agree on intersecting annuli.

If $p_m>0$, equation \eqref{eq:restricted-intensity-mass} implies
$u_m$ is not identically zero. By radial monotonicity it is strictly
positive on some outer annulus $r_*<|a|<1$. On compact subsets of this
annulus it is continuous and bounded below by a positive constant.
The weak logarithmic chain rule from
Theorem~\ref{thm:limiting-intensity-curvature} therefore applies to
\eqref{eq:restricted-intensity-curvature} and gives
\[
 \Delta\log u_m\ge\frac{2\beta}{(1-|a|^2)^2}
 \quad\hbox{on }\{r_*<|a|<1\}.
\]
More explicitly, testing $\Delta u_m$ against $\phi/u_m$ is justified
by uniform and $H^1$ approximation on each such compact subset; the
chain rule subtracts $\phi|\nabla u_m|^2/u_m^2$, which cancels the first
term on the right of \eqref{eq:restricted-intensity-curvature}.
Consequently
\[
 a\longmapsto\log u_m(a)+\frac\beta2\log(1-|a|^2)
\]
is radial subharmonic on that annulus. Its expression as a function of
$t=\log|a|$ is a finite convex function on $(\log r_*,0)$. A supporting
affine function at any interior point supplies a lower bound as
$t\uparrow0$. Exponentiating gives constants $c>0$ and $r_0<1$ such that
\[
 u_m(r)\ge c(1-r^2)^{-\beta/2},\qquad r_0<r<1.
\]
But
\[
 2\pi\int_{r_0}^1 r(1-r^2)^{-\beta/2}\,dr=\infty
 \qquad\hbox{when }\beta\ge2,
\]
contradicting \eqref{eq:restricted-intensity-mass}. Thus $p_m=0$.
Taking the union over all integers $m\ge1$ proves the theorem.
\end{proof}


\section{Schur coordinates and an explicit reference process}
\label{sec:schur-gram-reference}

This section gives a change of variables for the exact disk gas, an
independently solvable reference measure, and the precise likelihood
relating the two. The reference measure is the truncated circular beta ensemble
of Killip and Kozhan~\cite{KillipKozhan}. It is not the disk gas when
$\beta\ne2$. The identities and convergence argument needed here
are proved directly; no claim of bibliographic novelty is made for
the reference limit.

\subsection{The exact coefficient density}
For a monic polynomial $P$ of degree $n$, write
$P^*(z)=z^n\overline{P(1/\overline z)}$. Given
$\alpha_0,\ldots,\alpha_{N-1}\in\mathbb D$, define
\begin{equation}
 P_0=1,\qquad
 P_{j+1}(z)=zP_j(z)-\overline{\alpha_j}P_j^*(z).
 \label{eq:schur-exact-recursion}
\end{equation}

\begin{proposition}[Exact Schur density]
\label{prop:schur-exact-density}
Recursion~\eqref{eq:schur-exact-recursion} is a real analytic
bijection between $\mathbb D^N$ and the coefficient vectors of monic
polynomials of degree $N$ whose roots lie in $\mathbb D$. Its real
Jacobian is
\begin{equation}
 J_N(\alpha)=\prod_{j=0}^{N-1}(1-|\alpha_j|^2)^j.
 \label{eq:schur-coefficient-jacobian}
\end{equation}
Consequently, if $b=\beta/2$, the Schur parameters of the exact gas have
joint density
\begin{equation}
 \frac{1}{Q_{N,\beta}}
 |\operatorname{Disc}P_N|^{b-1}
 \prod_{j=0}^{N-1}(1-|\alpha_j|^2)^j
 \prod_{j=0}^{N-1}dA(\alpha_j).
 \label{eq:schur-actual-density}
\end{equation}
Here $|\operatorname{Disc}P_N|=|\Delta_N|^2$, and the density is understood
almost everywhere. In particular, at $\beta=2$ the parameters are
independent, with respective densities
$(j+1)(1-|\alpha_j|^2)^j/\pi$.
\end{proposition}

\begin{proof}
If the roots of $P_j$ lie in the disk, then on its boundary
$|zP_j(z)|=|P_j^*(z)|$. Rouch\'e's theorem proves that $P_{j+1}$ has all
$j+1$ roots in the disk. Conversely, the constant coefficient of a
stable monic polynomial $P_n$ determines
$\alpha_{n-1}=-\overline{P_n(0)}\in\mathbb D$, and
\[
 P_{n-1}(z)=\frac{P_n(z)+\overline{\alpha_{n-1}}P_n^*(z)}
                   {z(1-|\alpha_{n-1}|^2)}.
\]
The numerator has a zero at zero and, again by Rouch\'e's theorem,
exactly $n$ zeros in the disk. The quotient is monic, of degree $n-1$,
with all roots in the disk. This proves bijectivity and the stated
regularity of the inverse.

Write $P_{n-1}(z)=z^{n-1}+a_1z^{n-2}+\cdots+a_{n-1}$. With
$\alpha=\alpha_{n-1}$ fixed, its first $n-1$ new coefficients are
$a_k-\overline\alpha\,\overline{a_{n-k}}$, whereas the constant
coefficient is $-\overline\alpha$. The real determinant of the first
map is $(1-|\alpha|^2)^{n-1}$: after a scalar unitary change of
coordinates its antilinear involution has $n-1$ real eigenvalues $1$
and $n-1$ real eigenvalues $-1$. The constant-coefficient map has
absolute Jacobian one, and block triangularity proves
\eqref{eq:schur-coefficient-jacobian} by induction.

The map from labeled roots to coefficients has complex Jacobian equal,
up to sign, to the Vandermonde and has multiplicity $N!$ off the
collision set. Thus its real Jacobian is $|\Delta_N|^2$ and its
pushforward density is
$N!|\Delta_N|^{\beta-2}/Z_{N,\beta}$. Combining the two changes of
variables gives~\eqref{eq:schur-actual-density}. The integrability for
$0<\beta<2$, including near the discriminant set, follows directly
from this change of variables; it is not an additional assumption.
The independent normalization at $\beta=2$ follows from
$\int_{\mathbb D}(1-|a|^2)^j\,dA(a)=\pi/(j+1)$.
\end{proof}

\subsection{A Gram determinant is the entire change of measure}
Put
\[
 R_N(z)=\prod_{i,j=1}^N(1-z_i\overline z_j)>0,\qquad
 K(z,w)=\frac{1}{1-z\overline w}.
\]

\begin{proposition}[Resultant and Gram identities]
\label{prop:schur-gram-identity}
For the polynomial in~\eqref{eq:schur-exact-recursion},
\begin{equation}
 R_N=\prod_{j=0}^{N-1}(1-|\alpha_j|^2)^{j+1}.
 \label{eq:schur-resultant-product}
\end{equation}
For distinct points in the disk,
\begin{equation}
 D_N(z):=\det[K(z_i,z_j)]_{i,j=1}^N
 =\frac{|\Delta_N(z)|^2}{R_N(z)}
 =\prod_i\frac{1}{1-|z_i|^2}
  \prod_{i<j}\left|\frac{z_i-z_j}{1-z_i\overline z_j}\right|^2.
 \label{eq:schur-szego-gram}
\end{equation}
Let $\mathsf R_{N,b}$ be the law obtained from independent Schur
parameters of densities
\begin{equation}
 \frac{b(j+1)}{\pi}(1-|a|^2)^{b(j+1)-1}\,dA(a),
 \qquad 0\le j<N.
 \label{eq:schur-reference-parameters}
\end{equation}
Its density on labeled roots is
\begin{equation}
 \frac{b^N}{\pi^N}|\Delta_N(z)|^2R_N(z)^{b-1}\prod_i dA(z_i).
 \label{eq:schur-reference-roots}
\end{equation}
The exact disk gas satisfies
\begin{equation}
 \frac{d\mathsf P_{N,2b}}{d\mathsf R_{N,b}}(z)
 =\frac{D_N(z)^{b-1}}{\mathbb E_{\mathsf R_{N,b}}D_N^{b-1}},
 \qquad
 \mathbb E_{\mathsf R_{N,b}}D_N^{b-1}
 =\frac{b^NZ_{N,2b}}{\pi^N}.
 \label{eq:schur-gram-change-measure}
\end{equation}
Moreover, adjoining one point $a$ changes the determinant by exactly
\begin{equation}
 \frac{D_{N+1}(z_1,\ldots,z_N,a)}{D_N(z_1,\ldots,z_N)}
 =\frac{1}{1-|a|^2}
   \prod_{j=1}^N\left|\frac{a-z_j}{1-a\overline z_j}\right|^2.
 \label{eq:schur-gram-insertion}
\end{equation}
\end{proposition}

\begin{proof}
Write $\operatorname{Res}(A,B)$ for the polynomial resultant. Since
$P_N$ is monic,
$\operatorname{Res}(P_N,P_N^*)=\prod_iP_N^*(z_i)=R_N$.
The reversed recursion gives
$P_n^*=P_{n-1}^*-\alpha_{n-1}zP_{n-1}$, and hence
\[
 P_n^*=(1-|\alpha_{n-1}|^2)P_{n-1}^*-\alpha_{n-1}P_n.
\]
The elementary multiplicative and remainder properties of the
resultant imply
\begin{align*}
 \operatorname{Res}(P_n,P_n^*)
 &=(1-|\alpha_{n-1}|^2)^n
          \operatorname{Res}(P_n,P_{n-1}^*)\\
 &=(1-|\alpha_{n-1}|^2)^n
          \operatorname{Res}(zP_{n-1},P_{n-1}^*)\\
 &=(1-|\alpha_{n-1}|^2)^n
          \operatorname{Res}(P_{n-1},P_{n-1}^*),
\end{align*}
since $P_{n-1}^*(0)=1$. These identities follow either by evaluating
products over roots or by polynomial continuity when degrees drop.
Induction proves~\eqref{eq:schur-resultant-product}.

The Cauchy determinant identity gives
$\det[(1-x_i y_j)^{-1}]=
\prod_{i<j}(x_i-x_j)(y_i-y_j)/\prod_{i,j}(1-x_i y_j)$.
For completeness, after multiplication by the denominator the
numerator is alternating in each family, has degree at most $N-1$
in each individual variable, and is therefore a constant multiple of
the two Vandermondes. Comparing a successive residue at $x_Ny_N=1$
shows inductively that the constant is one. Taking
$y_j=\overline z_j$ proves the first equality in
\eqref{eq:schur-szego-gram}; separating diagonal factors proves the
second. Positivity also follows because $K$ is a Gram kernel:
$K(z,w)=\sum_{k\ge0}z^k\overline w^{\,k}$.

Dividing the product density~\eqref{eq:schur-reference-parameters}
by the Jacobian~\eqref{eq:schur-coefficient-jacobian}, and then
multiplying by $|\Delta_N|^2/N!$, gives
\eqref{eq:schur-reference-roots}. The difference of the exponents in
this division is $(b-1)(j+1)$, so
\eqref{eq:schur-resultant-product} applies. Division of the actual
root density by~\eqref{eq:schur-reference-roots} gives
\eqref{eq:schur-gram-change-measure}. Finally, cancel factors in
\eqref{eq:schur-szego-gram} to obtain
\eqref{eq:schur-gram-insertion}.
\end{proof}

\subsection{The reference process has the Bergman limit for every parameter}
For $P_n$ from the reference construction, the function
$B_n=P_n/P_n^*$ is a finite Blaschke product, and its zeros are exactly
the roots of $P_n$. It satisfies
\begin{equation}
 B_0=1,\qquad
 B_{n+1}(z)=\frac{zB_n(z)-\overline{\alpha_n}}
                 {1-\alpha_nzB_n(z)}.
 \label{eq:schur-blaschke-recursion}
\end{equation}

\begin{theorem}[Solvable reference limit]
\label{thm:schur-reference-bergman}
Fix $b>0$, and use the independent parameters
\eqref{eq:schur-reference-parameters}. In the topology of uniform
convergence on compact subsets of $\mathbb D$,
\begin{equation}
 \sqrt N B_N(z)\ \Longrightarrow\
 -\frac1{\sqrt b}\sum_{\ell=0}^\infty\xi_\ell z^\ell,
 \label{eq:schur-reference-gaf}
\end{equation}
where the $\xi_\ell$ are independent standard circular complex
Gaussians. Consequently the unscaled zero process of
$\mathsf R_{N,b}$ converges to the Bergman process for every $b>0$.
The identification of the limiting zero process follows by comparison
with $b=1$, where the reference is exactly the $\beta=2$ disk gas.
\end{theorem}

\begin{proof}
Direct integration of~\eqref{eq:schur-reference-parameters} gives
\begin{equation}
 \mathbb E|\alpha_n|^2=\frac{1}{b(n+1)+1},\qquad
 \mathbb E|\alpha_n|\le C_b(n+1)^{-1/2}.
 \label{eq:schur-reference-second-moment}
\end{equation}
If $w\in\mathbb D$ and $a$ has uniform argument conditional on
$|a|=s$, a geometric-series integration shows that
\[
 \mathbb E\left[\left|\frac{w-\overline a}{1-aw}\right|^2
                   \,\middle|\,|a|=s\right]
 =|w|^2+\frac{s^2(1-|w|^2)^2}{1-s^2|w|^2}
 \le |w|^2+s^2.
\]
Independence of $\alpha_n$ from the preceding parameters yields, for
$|z|\le r<1$,
\[
 \mathbb E|B_{n+1}(z)|^2
 \le r^2\mathbb E|B_n(z)|^2+\frac{C_b}{n+1}.
\]
Iteration, splitting the resulting geometric sum at $n/2$, gives
\begin{equation}
 \sup_{|z|\le r}\mathbb E|B_n(z)|^2
 \le\frac{C_{b,r}}{n+1}.
 \label{eq:schur-blaschke-second-bound}
\end{equation}
If $r<R<1$, Cauchy's formula and Cauchy--Schwarz imply
\[
 \sup_{|z|\le r}|B_n(z)|^2
 \le C_{r,R}\int_0^{2\pi}|B_n(Re^{i\theta})|^2
                              \frac{d\theta}{2\pi}.
\]
Thus~\eqref{eq:schur-blaschke-second-bound} holds with the supremum
inside the expectation as well, after changing its constant.

Write the nonlinear remainder as
\[
 E_n(z):=B_{n+1}(z)-zB_n(z)+\overline{\alpha_n}
 =\frac{\alpha_nzB_n(z)
               (zB_n(z)-\overline{\alpha_n})}
              {1-\alpha_nzB_n(z)}.
\]
The denominator has modulus at least $1-r$ on $|z|\le r$. Since
$|B_n|\le1$, independence,~\eqref{eq:schur-reference-second-moment},
and the supremum version of~\eqref{eq:schur-blaschke-second-bound}
give
\begin{equation}
 \mathbb E\sup_{|z|\le r}|E_n(z)|
 \le C_{b,r}(n+1)^{-3/2}.
 \label{eq:schur-nonlinear-remainder}
\end{equation}
Iterating the linearized recursion gives the exact identity
\[
 B_N(z)=z^N-\sum_{\ell=0}^{N-1}
       z^\ell\overline{\alpha_{N-1-\ell}}
       +\sum_{n=0}^{N-1}z^{N-1-n}E_n(z).
\]
By~\eqref{eq:schur-nonlinear-remainder}, the expectation of the
supremum of the last sum, multiplied by $\sqrt N$, tends to zero on
every compact disk: for $n\ge N/2$ its bound is $C_{b,r}/N$, and
for $n<N/2$ exponential decay of $r^{N-1-n}$ dominates its bounded
sum. The term $\sqrt Nz^N$ tends to zero as well.

For every fixed $L$, the independent variables
$\sqrt N\alpha_{N-1},\ldots,\sqrt N\alpha_{N-L}$ converge jointly to
independent complex Gaussians of variance $1/b$. This follows
immediately from their rescaled densities
\[
 \frac{b(N-\ell)}{\pi N}
 \left(1-\frac{|u|^2}{N}\right)^{b(N-\ell)-1}
 \mathbf1_{\{|u|<\sqrt N\}},\qquad 0\le\ell<L.
\]
The remaining series tails vanish uniformly in the required order.
Indeed, independence and zero means give for $|z|=R<1$
\[
 \mathbb E\left|\sqrt N\sum_{\ell=L}^{N-1}
           z^\ell\overline{\alpha_{N-1-\ell}}\right|^2
 =N\sum_{\ell=L}^{N-1}
       \frac{R^{2\ell}}{b(N-\ell)+1}.
\]
The part $\ell\le N/2$ is at most $C_{b,R}R^{2L}$, while the other
part is at most $C_bN R^N\log(N+1)$ and tends to zero. A second
application of Cauchy's formula gives the same tail control for the
supremum on $|z|\le r<R$. Finite-dimensional convergence followed by
this tail estimate proves~\eqref{eq:schur-reference-gaf}.

The limiting analytic function is almost surely not identically zero.
For analytic functions converging locally uniformly to a nonzero
analytic function, the zero counting measures converge vaguely: on
any small disk whose boundary avoids the limit zeros, Rouch\'e's
theorem preserves the number of zeros, with multiplicity; finitely
many such disks and uniform approximation of compactly supported
continuous tests prove the assertion. The continuous mapping theorem
therefore proves convergence of the zero processes. The scalar
$-1/\sqrt b$ does not alter the limiting zero set. At $b=1$,
\eqref{eq:schur-reference-roots} is the exact $\beta=2$ gas, whose
Bergman limit was proved earlier. Therefore the zero set of the
limiting analytic function, and hence the reference limit at every
$b>0$, is the Bergman process.
\end{proof}

\subsection{Why direct consistency fails for the exact gas}

\begin{proposition}[Failure of forward Schur consistency]
\label{prop:schur-nonprojectivity}
For every $\beta>0$ with $\beta\ne2$, the marginal law of $\alpha_0$
under the two-particle exact gas is not uniform on $\mathbb D$.
Consequently the exact finite-dimensional Schur laws
\eqref{eq:schur-actual-density} are not the initial-coordinate
marginals of one probability measure on $\mathbb D^{\mathbb N}$.
\end{proposition}

\begin{proof}
At $N=1$, the parameter $\alpha_0$ is uniform. At $N=2$, put $b=\beta/2$.
The polynomial and its discriminant are
\[
 P_2(z)=z^2+(\overline{\alpha_1}\alpha_0-
            \overline{\alpha_0})z-\overline{\alpha_1},\qquad
 \operatorname{Disc}P_2=
 (\overline{\alpha_1}\alpha_0-\overline{\alpha_0})^2
       +4\overline{\alpha_1}.
\]
Up to a common normalizing constant the marginal density at
$\alpha_0=0$ is
\begin{equation}
 I_0=4^{b-1}\int_{\mathbb D}(1-|a|^2)|a|^{b-1}\,dA(a)
 =\frac{\pi4^b}{(b+1)(b+3)}.
 \label{eq:schur-marginal-center}
\end{equation}
Its radial boundary limit is
\begin{equation}
 I_1=\int_{\mathbb D}(1-|a|^2)|1+a|^{2b-2}\,dA(a)
 =\frac{\pi\Gamma(2b+1)}{\Gamma(b+2)^2}.
 \label{eq:schur-marginal-edge}
\end{equation}
Here is a direct evaluation of the last integral. The binomial series
and angular orthogonality give
\[
 \frac{I_1}{\pi}
 =\sum_{k\ge0}\frac{((1-b)_k)^2}{(k!)^2(k+1)(k+2)}
 =\frac12\,{}_2F_1(1-b,1-b;3;1).
\]
For $0<b<1$, Euler's beta integral, first at an argument below one
and then by monotone convergence of this nonnegative series, evaluates
the last expression as
$\Gamma(2b+1)/\Gamma(b+2)^2$. Both the integral and the gamma ratio
are holomorphic in $b$ on $\Re b>0$ (use domination on compact
subsets near the single boundary point $a=-1$), so the identity
theorem proves the formula for all real $b>0$.

We justify the boundary limit used in~\eqref{eq:schur-marginal-edge}.
It suffices by rotation to take $\alpha_0=t\uparrow1$. Write
$s=(1-\sqrt{1-t^2})/(1+\sqrt{1-t^2})$. Its discriminant factors as
\[
 t^2(a-1)^2+4a=\frac4{(1+s)^2}(a+s)(1+sa).
\]
For $b\ge1$ bounded convergence applies. If $0<b<1$, set $p=1-b<1$.
Since $|a+s|\le|1+sa|$ for $|a|<1$,
\[
 |(a+s)(1+sa)|^{-p}\le|a+s|^{-2p}.
\]
For some $q>1$ one has $2pq<2$; hence these integrands have uniformly
bounded $L^q(\mathbb D)$ norms as $s\uparrow1$. Uniform integrability
and pointwise convergence prove the boundary limit. The same
argument with an exponent strictly below two gives continuity of
the marginal near $\alpha_0=0$. Thus differing values $I_0,I_1$
really imply a nonconstant density, rather than only a difference
of versions on null sets.

The gamma duplication identity reduces their ratio to
\[
 \frac{I_0}{I_1}
 =\frac{\sqrt\pi\,(b+1)\Gamma(b+1)}{(b+3)\Gamma(b+1/2)}.
\]
This ratio is strictly increasing for $b>0$: its logarithmic
 derivative is
\[
 \frac1{b+1}-\frac1{b+3}+\psi(b+1)-\psi(b+1/2)>0,
\]
where $\psi$ is the logarithmic derivative of $\Gamma$ and is strictly
increasing because
$\psi'(x)=\sum_{k=0}^\infty(x+k)^{-2}>0$ for $x>0$.
At $b=1$ the ratio is one, so it differs from one for every
$b\ne1$. This proves the proposition.
\end{proof}

\begin{remark}[What this representation does and does not settle]
The Gram likelihood~\eqref{eq:schur-gram-change-measure} is exact.
It gives an explicit reference construction with a proved full-sequence
limit, and~\eqref{eq:schur-gram-insertion} isolates the interaction
created by an additional interior point. Its normalized likelihood
is not shown here to have a limit, or to be uniformly integrable
jointly with the zero process. Those estimates cannot be omitted:
without them, Theorem~\ref{thm:schur-reference-bergman} says nothing
about convergence of the exact disk gas at $\beta\ne2$.
Proposition~\ref{prop:schur-nonprojectivity} rules out direct forward
consistency; it does not rule out convergence of terminal Schur
coordinates, nor a successful reverse-recursion argument.
\end{remark}

\subsection{Only logarithmically many terminal parameters are needed locally}

\begin{proposition}[Uniform reverse contraction]
\label{prop:schur-reverse-contraction}
Consider two recursions of the form~\eqref{eq:schur-blaschke-recursion}
which use the same last $L\ge1$ parameters but may have arbitrary
initial values in the closed unit disk. Their terminal functions
$B$ and $\widetilde B$ satisfy
\begin{equation}
 \sup_{|z|\le r}|B(z)-\widetilde B(z)|
 \le2\operatorname{arctanh}(r)\,r^{L-1},\qquad 0<r<1.
 \label{eq:schur-uniform-contraction}
\end{equation}
\end{proposition}

\begin{proof}
Use the Poincar\'e distance whose infinitesimal metric is
$2|dw|/(1-|w|^2)$. The map
$w\mapsto(w-\overline\alpha)/(1-\alpha w)$ is an isometry.
Multiplication by $z$, with $|z|\le r$, has Lipschitz constant at
most $r$ in this metric, because its infinitesimal norm is
\[
 \frac{|z|(1-|w|^2)}{1-|zw|^2}\le |z|\le r.
\]
After the first shared step the distance is at most
$4\operatorname{arctanh}(r)$, since before its isometry both values
lie in the closed disk of radius $r$. Each remaining step contracts
that distance by at most $r$. Poincar\'e distance is at least twice
Euclidean distance, giving~\eqref{eq:schur-uniform-contraction}.
Boundary initial values cause no difficulty after multiplication
by $z$; alternatively one can take a limit from inside the disk.
\end{proof}

\begin{theorem}[Logarithmic terminal representation of the exact gas]
\label{thm:schur-logarithmic-memory}
Fix $\beta>0$ and $0<r<1$. There is a constant $C_{\beta,r}<\infty$
with the following property. Obtain the Schur parameters of the
exact $N$-particle gas by Proposition~\ref{prop:schur-exact-density},
and put
\[
 L_N=\min\{N,\lceil C_{\beta,r}\log(N+1)\rceil\}.
\]
Restart recursion~\eqref{eq:schur-blaschke-recursion} with value zero,
using only $\alpha_{N-L_N},\ldots,\alpha_{N-1}$, and let
$\widehat\Xi_N$ be the zero counting measure of its terminal rational
function in $\mathbb D$. Then, for every continuous function $f$
with compact support in $B_r$,
\begin{equation}
 \int f\,d\Xi_N-\int f\,d\widehat\Xi_N
 \longrightarrow0\quad\hbox{in probability}.
 \label{eq:schur-terminal-zero-approximation}
\end{equation}
Thus the interior process is approximated, in its natural topology,
by a function of $O(\log N)$ terminal Schur parameters.
\end{theorem}

\begin{proof}
We first record the consequence of the radial comparison already
proved for the exact gas. If
$H_N=\sum_{i=1}^N-\log|z_i|$, it is stochastically bounded above by
$Y_N=\sum_{j=1}^NjE_j$, where the $E_j$ are independent exponentials
of rates $a_j=2j+\beta j(j-1)/2$. Hence
\[
 \mathbb E Y_N=\frac2\beta\log N+O_\beta(1),\qquad
 \operatorname{Var}Y_N=\sum_{j=1}^N\frac{j^2}{a_j^2}
                      \le C_\beta.
\]
For every $A>2/\beta$ it follows that
\begin{equation}
 \mathbb P\{H_N>A\log N\}\longrightarrow0.
 \label{eq:schur-total-depth-logarithmic}
\end{equation}
No coupling between Schur parameters and the dominating exponentials
is required for this probabilistic assertion.

Choose $r<R<R_1<1$. For $|z|\le R$ and $|w|>R_1$, the disk Green
function obeys
\begin{equation}
 -\log\left|\frac{z-w}{1-z\overline w}\right|
 \le C_{R,R_1}(-\log|w|).
 \label{eq:schur-green-depth-bound}
\end{equation}
Indeed, the pseudohyperbolic distance in this expression is bounded
below by $d=(R_1-R)/(1+R)>0$, while
\[
 1-\left|\frac{z-w}{1-z\overline w}\right|^2
 =\frac{(1-|z|^2)(1-|w|^2)}{|1-z\overline w|^2}
 \le\frac{2(1-|w|)}{(1-R)^2}.
\]
The inequality $-\log x\le(1-x^2)/(2d^2)$ for $d\le x\le1$
proves~\eqref{eq:schur-green-depth-bound}.
For an interior root with $|w|\le R_1$ and $|z-w|\ge\delta$, the
same Green function is at most $\log(2/\delta)$. Since $B_N$ is the
finite Blaschke product of the roots, this gives on any such contour
\begin{equation}
 |B_N(z)|\ge
 \exp\{-C_{R,R_1}H_N-M_N(R_1)\log(2/\delta)\}.
 \label{eq:schur-contour-lower-bound}
\end{equation}

Here are explicit contours sufficient to compare point measures.
Fix a mesh $\eta>0$ small enough that every square of side $\eta$
intersecting $B_r$ is contained in $B_R$. Put $M=M_N(R_1)$ and
$\delta=\eta/[8(M+1)]$. Choose the horizontal and vertical shifts
of the square grid so that every grid line has distance at least
$\delta$ from the corresponding coordinate of each of these $M$
roots. Such shifts exist: in either shift interval of length $\eta$,
the forbidden intervals have total length at most
$2M\delta<\eta/4$. Consequently every boundary point of the grid
squares under consideration has distance at least $\delta$ from
every root in $B_{R_1}$.

Compact-count tightness implies
$M_N(R_1)\le\sqrt{\log N}$ with probability tending to one. For this
fixed mesh, therefore,
\[
 M_N(R_1)\log\frac{16(M_N(R_1)+1)}{\eta}=o(\log N)
\]
on events whose probabilities tend to one. Combining with
\eqref{eq:schur-total-depth-logarithmic} and
\eqref{eq:schur-contour-lower-bound}, the modulus on all the chosen
square boundaries is at least
$N^{-C_{R,R_1}A-o(1)}$ with probability tending to one.
On the other hand, Proposition~\ref{prop:schur-reverse-contraction}
bounds the difference between $B_N$ and the restarted function by
$2\operatorname{arctanh}(R)R^{L_N-1}$ on $B_R$.
Choose $C_{\beta,r}$ so large that
\[
 C_{\beta,r}|\log R|>C_{R,R_1}A+2.
\]
This choice is independent of the mesh. Rouch\'e's theorem then shows,
with probability tending to one, that the original and restarted
functions have the same number of zeros in each selected square.
Both functions are analytic in the disk; the restarted one is not
identically zero almost surely, since its value at zero is
$-\overline{\alpha_{N-1}}$ and the original product of root moduli is
almost surely positive.

On this event the zeros can be paired within individual squares,
including multiplicities. If $\omega_f$ is the modulus of continuity
of $f$ extended by zero, the difference of its two integrals is at
most $\omega_f(\sqrt2\eta)M_N(R)$. The latter counts are tight
uniformly in $N$. First take $N\to\infty$ at fixed mesh, then let
$\eta\downarrow0$. This proves
\eqref{eq:schur-terminal-zero-approximation}.
\end{proof}

\subsection{Failure of uniform integrability above the empty-limit threshold}

\begin{corollary}[Asymptotic singularity of the Gram change of measure]
\label{cor:schur-gram-singularity}
For $b\ge2$, write
\[
 W_N=\frac{D_N^{b-1}}{\mathbb E_{\mathsf R_{N,b}}D_N^{b-1}}.
\]
Then
\begin{equation}
 \|\mathsf P_{N,2b}-\mathsf R_{N,b}\|_{\rm TV}\longrightarrow1,
 \qquad
 W_N\longrightarrow0\quad\hbox{in }\mathsf R_{N,b}\hbox{-probability}.
 \label{eq:schur-gram-singularity}
\end{equation}
In particular, the mean-one likelihoods $W_N$ are not uniformly
integrable.
\end{corollary}

\begin{proof}
Let $A_R$ be the event of no points in $B_R$. Since
$W_N\,d\mathsf R_{N,b}=d\mathsf P_{N,2b}$,
\[
 \mathbb E_{\mathsf R_{N,b}}\min(W_N,1)
 \le\mathsf R_{N,b}(A_R)+\mathsf P_{N,2b}(A_R^c).
\]
For fixed $R<1$, the second term tends to zero by the empty-limit
 theorem for $2b\ge4$, including Theorem~\ref{thm:critical-empty}. The first tends to the Bergman hole probability,
by Theorem~\ref{thm:schur-reference-bergman}; the circle $|z|=R$
has zero probability of containing a Bergman point, so the count
convergence indeed implies this hole convergence. These Bergman
hole probabilities tend to zero as $R\uparrow1$, since that process
has infinitely many points almost surely. Thus
$\mathbb E\min(W_N,1)\to0$. For a likelihood of mean one,
\[
 1-\|\mathsf P_{N,2b}-\mathsf R_{N,b}\|_{\rm TV}
 =\mathbb E_{\mathsf R_{N,b}}\min(W_N,1),
\]
where total variation is the supremum over measurable events.
Also, for every $\varepsilon>0$,
\[
 \mathsf R_{N,b}\{W_N>\varepsilon\}
 \le\frac{\mathbb E\min(W_N,1)}{\min(\varepsilon,1)}\longrightarrow0.
\]
Uniform integrability would then force $\mathbb E W_N\to0$, contrary
to its constant value one.
\end{proof}


\section{A lowest-weight decomposition at the critical parameter}
\label{sec:critical-su11-sectors}

We give an exact representation of the one-particle radial distribution
at $\beta=4$. It separates a simple, explicit representation-theoretic
kernel from a probability distribution on primary sectors. The latter
distribution contains the finer asymptotic information that the explicit
kernel alone does not determine.

Use normalized area $d\mu=dA/\pi$ and the polynomial Bergman inner
product on $\mathbb D^n$. On polynomials define
\[
 D_n=\sum_{i=1}^n\partial_{w_i},\qquad
 L_n=\sum_{i=1}^n(w_i^2\partial_{w_i}+2w_i),\qquad
 H_n=\sum_{i=1}^n w_i\partial_{w_i}+n.
\]
All adjoint statements below concern polynomials, so no assertion about
domains of closed unbounded operators is needed. Monomial orthogonality
and $\|w^j\|^2=(j+1)^{-1}$ give
\begin{equation}
 D_n^*=L_n,\qquad [D_n,L_n]=2H_n,\qquad [H_n,L_n]=L_n.
 \label{eq:critical-su11-operators}
\end{equation}
For example, in one variable,
$\langle\partial w^{j+1},w^j\rangle=1
=\langle w^{j+1},(w^2\partial+2w)w^j\rangle$;
all other monomial matrix elements follow from orthogonality.
The commutators follow by differentiation.

Let $\mathcal I_n$ be the space of symmetric polynomials divisible by
$\Delta_n^2$, with $\Delta_1=1$. Write $\mathcal I_{n,d}$ for its
homogeneous part of degree $d$ and
\[
 \mathcal P_{n,d}=\mathcal I_{n,d}\cap\ker D_n.
\]
Translation invariance of $\Delta_n^2$ gives $D_n\Delta_n^2=0$.
The pairwise derivative identity
\[
 \sum_iw_i^2\partial_i\Delta_n^2
 =2(n-1)\left(\sum_iw_i\right)\Delta_n^2
\]
shows that $D_n,L_n,H_n$ all preserve $\mathcal I_n$.
In particular, $\mathcal I_{n,d}=\{0\}$ for $d<n(n-1)$.

\begin{lemma}[Orthogonal decomposition into descendants]
\label{lem:critical-su11-descendants}
For every degree $d$,
\begin{equation}
 \mathcal I_{n,d}=
 \mathop{\bigoplus}_{\substack{\ell\leq d\\\ell\geq n(n-1)}}^{\perp}
       L_n^{d-\ell}\mathcal P_{n,\ell}.
 \label{eq:critical-su11-decomposition}
\end{equation}
If $v\in\mathcal P_{n,\ell}$, $\kappa=n+\ell$, and $K\geq0$, then
\begin{equation}
 \frac{\|D_n^kL_n^K v\|^2}{\|L_n^K v\|^2}
 = (K)_{\underline k}(2\kappa+K-1)_{\underline k},
 \qquad 0\leq k\leq K,
 \label{eq:critical-su11-lowering-norm}
\end{equation}
where $(a)_{\underline k}=a(a-1)\cdots(a-k+1)$; the numerator is zero
for $k>K$.
\end{lemma}

\begin{proof}
The commutators give, by induction,
\[
 D_nL_n^jv=j(2\kappa+j-1)L_n^{j-1}v,
 \qquad
 \|L_n^jv\|^2=j!(2\kappa)_j\|v\|^2,
\]
where the last Pochhammer symbol is rising. Iterating the first identity
and dividing the second proves the norm formula.

In each finite-dimensional homogeneous space, the orthogonal complement
of $L_n\mathcal I_{n,d-1}$ is exactly $\mathcal P_{n,d}$, by the adjoint
identity and preservation of the ideal. This proves that the descendants
span, by induction on degree. They are mutually orthogonal: for primitive
vectors $u,v$ of different degrees, transfer powers of $L_n$ across the
inner product using its adjoint $D_n$. The lowering formula reduces the
inner product to a scalar multiple of an inner product between a primitive
vector and a positive descendant. Transferring one remaining raising
operator kills that primitive vector. Within a fixed primitive degree,
the displayed norm identity preserves orthogonality of primitive vectors.
This proves the orthogonal direct sum.
\end{proof}

Fix $N\geq2$, set $n=N-1$ and $d=N(N-1)$, and expand
\[
 F_N(z,w):=\Delta_N(z,w)^2=\sum_{k=0}^{2N-2}z^k f_{k,N}(w).
\]
Then
\begin{equation}
 f_{0,N}(w)=\left(\prod_{i=1}^n w_i^2\right)\Delta_n(w)^2,
 \qquad
 f_{k,N}=\frac{(-D_n)^k}{k!}f_{0,N}.
 \label{eq:critical-su11-coefficients}
\end{equation}
The second assertion follows by equating coefficients in
$(\partial_z+D_n)F_N=0$. By the lemma there are unique orthogonal vectors
\[
 f_{0,N}=\sum_{K=0}^{2N-2} h_{K,N},\qquad
 h_{K,N}\in L_n^K\mathcal P_{n,d-K}.
\]
Set $b_{K,N}=\|h_{K,N}\|^2$, allowing zero sectors, and put
\begin{equation}
 T_N=2N^2-2,\qquad
 A_N(K)=\frac1{K+1}\binom{T_N}{K},\qquad
 p_N(K)=\frac{b_{K,N}A_N(K)}{\|F_N\|^2}.
 \label{eq:critical-su11-sector-probabilities}
\end{equation}

\begin{theorem}[An explicit hypergeometric kernel for the radial law]
\label{thm:critical-su11-hypergeometric}
The numbers $p_N(K)$ are a probability distribution on
$\{0,\ldots,2N-2\}$. Define
\begin{equation}
 a_{k,N}=\frac{\|f_{k,N}\|^2}{(k+1)\|F_N\|^2}.
 \label{eq:critical-su11-orbital-probability}
\end{equation}
These numbers also sum to one and satisfy
\begin{equation}
 a_{k,N}=\sum_{K=k}^{2N-2}p_N(K)\,
 \frac{\binom{K+1}{k+1}\binom{T_N-K-1}{k}}
      {\binom{T_N}{K}}.
 \label{eq:critical-su11-hypergeometric-mixture}
\end{equation}
Equivalently, first sample $K$ from $p_N$; then sample $J$ from the
hypergeometric distribution obtained by drawing $K$ objects without
replacement from a population of $T_N$, of which $K+1$ are marked;
finally put $k=K-J$.
The one-point intensity with respect to ordinary area is
\begin{equation}
 \rho_{1,N}(z)=\frac N\pi\sum_{k=0}^{2N-2}(k+1)a_{k,N}|z|^{2k}.
 \label{eq:critical-su11-intensity}
\end{equation}
In particular,
\begin{equation}
 0\leq\pi\rho_{1,N}(0)-Np_N(0)
       \leq\frac N{N^2-1}.
 \label{eq:critical-su11-primary-obstruction}
\end{equation}
Thus $\rho_{1,N}(0)\to0$ if and only if $Np_N(0)\to0$.
More generally, for every fixed nonnegative integer $k$,
\begin{equation}
 Na_{k,N}=Np_N(k)+O_k(N^{-1})\qquad(N\to\infty).
 \label{eq:critical-su11-fixed-orbitals}
\end{equation}
\end{theorem}

\begin{proof}
In the sector indexed by $K$, the primitive weight is
$\kappa=n+(d-K)=N^2-1-K$. The lowering norm formula, together with
orthogonality after lowering and \eqref{eq:critical-su11-coefficients},
gives
\[
 \|F_N\|^2
 =\sum_Kb_{K,N}\sum_{k=0}^K
 \frac{(K)_{\underline k}(T_N-K-1)_{\underline k}}
      {k!^2(k+1)}.
\]
The inner summand equals
\[
 \frac1{K+1}\binom{K+1}{k+1}\binom{T_N-K-1}{k}.
\]
Substitute $j=K-k$ and use the elementary binomial convolution: selecting
$K$ objects from disjoint sets of sizes $K+1$ and $T_N-K-1$ gives
\[
 \sum_{k=0}^K\binom{K+1}{k+1}\binom{T_N-K-1}{k}
 =\binom{T_N}{K}.
\]
This proves both the normalization $A_N(K)$ and the probability formula.
The same counting interpretation gives the stated hypergeometric law.

The coefficient functions $f_{k,N}$ are homogeneous of different degrees.
Their inner products vanish under simultaneous rotation of the $w_i$.
Integrating $|F_N(z,w)|^2$ in $w$, and remembering that normalized area
is $dA/\pi$, therefore proves \eqref{eq:critical-su11-intensity}.

At $k=0$, the conditional probability is
$(K+1)/\binom{T_N}{K}$. It equals one for $K=0$. For
$1\leq K\leq2N-2$, it is at most $2/T_N$: the ratio of its values at
$K+1$ and $K$ is $(K+2)/(T_N-K)$, which is at most one in the relevant
range (including $N=2$). Summing against $p_N$ proves
\eqref{eq:critical-su11-primary-obstruction}.

To check the last statement, write the conditional kernel as $w_N(k,K)$.
For fixed $k$,
\[
 w_N(k,k)=\frac{\binom{T_N-k-1}{k}}{\binom{T_N}{k}}
        =1+O_k(T_N^{-1}).
\]
For $K\geq k$, direct cancellation of factorials gives
\[
 \frac{w_N(k,K+1)}{w_N(k,K)}
 =\frac{K+2}{K+1-k}
   \frac{T_N-K-1-k}{T_N-K-1}
   \frac{K+1}{T_N-K}.
\]
Uniformly over $k\leq K\leq2N-3$, this is $O_k(N^{-1})$.
Thus, for all sufficiently large $N$, the kernel is decreasing for
$K\geq k+1$, and its largest value in that range is
\[
 w_N(k,k+1)
 =(k+2)\frac{\binom{T_N-k-2}{k}}{\binom{T_N}{k+1}}
 =O_k(T_N^{-1}).
\]
The mixture formula now gives
$|a_{k,N}-p_N(k)|=O_k(T_N^{-1})$, which proves
\eqref{eq:critical-su11-fixed-orbitals}.
\end{proof}

\begin{corollary}[An exact primary-sector criterion for critical emptiness]
\label{cor:critical-su11-empty-criterion}
For the unweighted disk gas at $\beta=4$, the following are equivalent:
\begin{enumerate}
 \item $\Xi_N$ converges to the empty point process on $\mathbb D$;
 \item for every fixed nonnegative integer $k$, $Np_N(k)\to0$.
\end{enumerate}
When these conditions hold, convergence on each compact window is in
total variation. The criterion uses the previously proved uniform compact
count bounds, but does not presume uniqueness of a nonempty limit.
\end{corollary}

\begin{proof}
Write $n_{k,N}=Na_{k,N}$. Integrating
\eqref{eq:critical-su11-intensity} gives the nonnegative series
\[
 \mathbb E\Xi_N(B_r)=\sum_{k=0}^{2N-2}n_{k,N}r^{2k+2},
 \qquad 0<r<1.
\]
If $\Xi_N$ converges to the empty process, the compact counts converge
to zero in probability. Their uniform second-moment bounds imply
convergence of their means to zero. Thus
$n_{k,N}r^{2k+2}\leq\mathbb E\Xi_N(B_r)\to0$ for each fixed $k$.
Equation~\eqref{eq:critical-su11-fixed-orbitals} implies the second
condition.

Conversely, that equation shows that $n_{k,N}\to0$ for each fixed $k$.
For $0<r<R<1$ and an integer $M\geq0$, positivity gives
\[
 \sum_{k>M}n_{k,N}r^{2k+2}
 \leq(r/R)^{2M+4}\mathbb E\Xi_N(B_R).
\]
The right side tends to zero uniformly in $N$ as $M\to\infty$ by the
compact-count bound. The finite sum over $k\leq M$ tends to zero with
$N$. Hence $\mathbb E\Xi_N(B_r)\to0$, and Markov's inequality gives
$\mathbb P\{\Xi_N(B_r)>0\}\to0$. This is exactly total variation
convergence of the process in $B_r$ to the empty configuration.
A countable exhaustion of the disk proves vague convergence in law.
\end{proof}

Combining this criterion with the independently proved critical
emptiness theorem, Theorem~\ref{thm:critical-mobius-empty}, gives
$Np_N(k)\to0$ for every fixed $k$. This is a qualitative consequence;
the representation does not presently give a decay rate for these
primary-sector probabilities.

\paragraph{Small-size checks.}
For $N=2$ the one-variable primitive space consists only of constants,
so $p_2(2)=1$. Here $T_2=6$, and the formula gives
$(a_{0,2},a_{1,2},a_{2,2})=(1/5,3/5,1/5)$.
This agrees with direct integration of $|z-w|^4$ and gives
$\pi\rho_{1,2}(0)=2/5$.
For an additional check, direct expansion of $\Delta_3^2$ gives
\[
 \|F_3\|^2=\frac{49}{12},\qquad
 (a_{0,3},\ldots,a_{4,3})
 =\left(\frac{23}{245},\frac{44}{245},\frac{108}{245},
                   \frac{10}{49},\frac4{49}\right).
\]
The only sectors are $K=0,2,4$, since a symmetric translation-invariant
polynomial in two variables is a polynomial in $(w_1-w_2)^2$.
Substitution in the triangular formula yields
\[
 (p_3(0),p_3(2),p_3(4))
 =\left(\frac{81}{1001},\frac{128}{273},\frac{104}{231}\right),
\]
whose entries sum to one and reproduce all five orbital probabilities.

\begin{remark}[Primary-sector decay and the remaining quantitative question]
The binomial kernel in \eqref{eq:critical-su11-hypergeometric-mixture}
is completely explicit. The independent critical emptiness argument
proves qualitative decay of the normalized mass of the primary projection
$h_{0,N}\in\ker D_{N-1}$ of
$\prod_iw_i^2\Delta_{N-1}^2$.
It proves $Np_N(0)\to0$, but no direct norm estimate or quantitative
rate for this projection has been obtained from the representation alone.
The condition that a holomorphic state vanish on pairwise collisions
does not alone select an infinite-volume state: every finite-particle
Laughlin polynomial already has that property. A uniqueness argument
based on a positive contact Hamiltonian would therefore need an
additional, justified condition selecting the boundary state.
\end{remark}


\section{Unconditional canonical quasi-Gibbs laws on interior disks}
\label{sec:unconditional-canonical-quasigibbs}

This section gives a conditional-law statement for every subsequential
limit and every $\beta>0$. It requires no estimate on the exterior
inverse-angular field. The extra harmonic field first appears on an
augmented probability space. Averaging this field gives a genuine
conditional density for the unmarked process, but that density need not
factor into one-body harmonic weights.

Fix an open Euclidean disk $B=B(a_0,R)$ with $\overline B\subset\D$.
Write $M=\Xi(B)$ and $\eta=\Xi|_{\D\setminus\overline B}$.
The limiting intensity is locally bounded, so $\Xi(\partial B)=0$
almost surely. Thus $\eta$ records the entire exterior up to a null
event. Whenever $M=m\ge1$, order the $m$ interior points uniformly
at random, independently of all other randomness, and denote the
resulting tuple by $X=(X_1,\ldots,X_m)$. Conditional densities below
refer to this uniform ordering and ordinary area measure.

For a holomorphic function $q$ on $B$ define
\[
 Z_{B,m}(q)=\int_{B^m}|\Delta_m(\mathbf z)|^\beta
                          \prod_{i=1}^m|q(z_i)|^2\,dA^m(\mathbf z),
 \qquad m\ge1.
\]
In particular this definition has no factor $1/m!$.

\begin{lemma}[A normal-family bound for normalized exterior factors]
\label{lem:normalized-exterior-normal-family}
Put
\[
 d_{m,\beta}=2m+\frac\beta2m(m-1),\qquad
 J_{m,\beta}=\int_{\D^m}|\Delta_m(\mathbf u)|^\beta\,dA^m(\mathbf u).
\]
Then $0<J_{m,\beta}<\infty$. If $q$ is holomorphic on $B$ and
$Z_{B,m}(q)=1$, then, whenever $B(a,r)\subset B$,
\begin{equation}\label{eq:normalized-exterior-point-bound}
 |q(a)|^{2m}\le J_{m,\beta}^{-1}r^{-d_{m,\beta}}.
\end{equation}
Consequently the functions with $Z_{B,m}(q)=1$ form a locally uniformly
bounded family in the space $\mathcal H(B)$ of holomorphic functions
with its locally uniform topology. The closure of the subfamily that
is zero-free and satisfies $q(a_0)>0$ is compact. Every member of this
closure is either identically zero or zero-free; a nonzero member
satisfies $q(a_0)>0$. Every member also satisfies $Z_{B,m}(q)\le1$.
\end{lemma}

\begin{proof}
Finiteness and positivity of $J_{m,\beta}$ follow because its integrand
is continuous and bounded on the closed polydisk and strictly positive
off the collision diagonals. For fixed $\mathbf u\in\D_r^m$, the
function
\[
 \lambda\longmapsto\prod_{i=1}^m q(a+\lambda u_i)
\]
is holomorphic on a neighborhood of the closed unit disk. The
mean-value inequality for its squared modulus gives
\[
 \frac1{2\pi}\int_0^{2\pi}
        \prod_i|q(a+e^{i\theta}u_i)|^2\,d\theta
       \ge |q(a)|^{2m}.
\]
The weight $|\Delta_m(\mathbf u)|^\beta$ and the domain $\D_r^m$
are invariant under a common rotation of all coordinates. Integrating
the last inequality with this weight therefore gives
\[
 1\ge\int_{B(a,r)^m}|\Delta_m(\mathbf z)|^\beta
                          \prod_i|q(z_i)|^2\,dA^m(\mathbf z)
       \ge J_{m,\beta}r^{d_{m,\beta}}|q(a)|^{2m}.
\]
This proves the bound. Applying it with a radius uniformly positive on
each compact subset of $B$ proves local boundedness. The Cauchy
integral formula then bounds all derivatives on smaller compact sets;
Arzel\`a--Ascoli and a diagonal exhaustion give compactness of the
closure in $\mathcal H(B)$. The elementary zero-free alternative for
locally uniform limits follows from Hurwitz's theorem. For completeness,
if a nonzero limit had a zero, a small circle enclosing only that zero
and no zero on its boundary would, by uniform convergence and Rouch\'e's
theorem, force a zero of the approximating functions. Finally Fatou's
lemma gives $Z_{B,m}(q)\le1$, and the phase assertion follows by
evaluation at $a_0$ and the zero-free alternative.
\end{proof}

\begin{theorem}[An augmented canonical DLR equation]
\label{thm:augmented-canonical-dlr}
Let $\Xi$ be any subsequential limit of the disk gas at a fixed
$\beta>0$. Fix $B$ as above and $m\ge1$ such that $\mathbb P(M=m)>0$.
On an extension of the probability space under the conditional law
given $M=m$, there is a random zero-free holomorphic function $q$ on
$B$ satisfying
\[
 q(a_0)>0,\qquad Z_{B,m}(q)=1,
\]
such that
\begin{equation}\label{eq:augmented-canonical-dlr}
 \mathbb P\{X\in d\mathbf z\mid\eta,q,M=m\}
   =|\Delta_m(\mathbf z)|^\beta
                       \prod_{i=1}^m|q(z_i)|^2\,dA^m(\mathbf z).
\end{equation}
The joint law can be obtained by a further subsequential limit of the
actual finite-$N$ exterior factors. In particular
$h=(2/\beta)\log|q|$ is a finite harmonic function on $B$, and the
right side is the usual canonical log-gas density with the additional
one-body term $\beta\sum_i h(z_i)$.
On the same augmented limit, the original finite-$N$ forces at all
interior points converge jointly, and are given by
\begin{equation}\label{eq:augmented-harmonic-force}
 C_{X_i}=\sum_{j\ne i}(X_i-X_j)^{-1}
                         +\frac2\beta\frac{q'(X_i)}{q(X_i)}.
\end{equation}

The theorem does not assert that $q$ or $h$ is measurable with respect
to $\eta$, or even with respect to the complete unmarked configuration.
\end{theorem}

\begin{proof}
Write $M_N=\Xi_N(B)$ and $\eta_N=\Xi_N|_{\D\setminus\overline B}$.
Almost surely there is no finite-$N$ particle on $\partial B$.
For a realized finite exterior tuple $\mathbf w$, choose the unique
holomorphic logarithm $L_w$ on $B$ such that
\[
 e^{L_w(z)}=\frac{z-w}{a_0-w},\qquad L_w(a_0)=0.
\]
It exists because $B$ is simply connected and the quotient has no zero
there. Set
\[
 q_N^{\rm raw}(z)=
       \exp\left(\frac\beta2\sum_{w\in\eta_N}L_w(z)\right),\qquad
 q_N(z)=\frac{q_N^{\rm raw}(z)}
              {Z_{B,m}(q_N^{\rm raw})^{1/(2m)}}
 \quad\hbox{on }\{M_N=m\}.
\]
The sum is finite. Its normalization is positive and finite: the raw
factor is bounded on $B$, and is positive in modulus there. Factoring
the original joint density into interior--interior, exterior--exterior
and cross terms proves exactly
\begin{equation}\label{eq:finite-normalized-exterior-kernel}
 \mathbb P\{X_N\in d\mathbf z\mid\eta_N,M_N=m\}
   =|\Delta_m(\mathbf z)|^\beta
                        \prod_i|q_N(z_i)|^2\,dA^m(\mathbf z).
\end{equation}
The factorial arising from uniform ordering cancels in the conditional
normalization.

Let $N_j$ be the original process-convergent subsequence. Since the
limit has no particle on $\partial B$, the restrictions and counts
converge jointly, and
$\mathbb P(M_{N_j}=m)\to\mathbb P(M=m)>0$. Thus the laws conditioned
on this event converge on unmarked configurations. Under these
conditional laws the interior tuples are tight in $B^m$. One direct
verification is to bound the probability that a coordinate lies in a
thin interior collar of $\partial B$ by the expected particle count
in that collar divided by $\mathbb P(M_{N_j}=m)$. The compact intensity
bound makes this uniformly small as the collar shrinks. On each
remaining compact set, ordinary finite-dimensional compactness applies.
Uniform random ordering gives the stated limiting ordering.

Lemma~\ref{lem:normalized-exterior-normal-family} places $q_N$ in a
fixed compact subset of $\mathcal H(B)$ after taking its closure.
Hence, on a further subsequence, $(\eta_N,X_N,q_N)$ converges jointly
under the law conditioned on $M_N=m$ to $(\eta,X,q)$. At this stage
the possible limit $q\equiv0$ has not yet been excluded.

For $\psi\in C_c(B^m)$ define
\[
 T_\psi(q)=\int_{B^m}\psi(\mathbf z)|\Delta_m(\mathbf z)|^\beta
                         \prod_i|q(z_i)|^2\,dA^m(\mathbf z).
\]
On the compact normal family this is bounded and continuous. At finite
$N$, equation~\eqref{eq:finite-normalized-exterior-kernel} implies, for
every bounded continuous function $H$ of exterior configurations and
holomorphic functions,
\[
 \mathbb E[H(\eta_N,q_N)\psi(X_N)\mid M_N=m]
   =\mathbb E[H(\eta_N,q_N)T_\psi(q_N)\mid M_N=m].
\]
Joint convergence passes this identity to the limit. Bounded
continuous functions determine finite measures on these metrizable
spaces, so the identity extends to bounded Borel $H$. Equivalently,
\[
 \mathbb E[\psi(X)\mid\eta,q,M=m]=T_\psi(q).
\]
Choose continuous compactly supported functions
$0\le\psi_j\uparrow1$ on $B^m$. Since $X$ has all $m$ coordinates
in $B$, conditional monotone convergence on the left and ordinary
monotone convergence on the right yield
$1=Z_{B,m}(q)$ almost surely. Thus $q$ is not identically zero, and
the zero-free alternative proves all the assertions about $q$.
The conditional identity now extends to all bounded Borel $\psi$ by
the monotone-class theorem, proving the claimed probability kernel.
At finite $N$, logarithmic differentiation of $q_N$ gives exactly
\[
 C_{X_{i,N},N}=\sum_{j\ne i}(X_{i,N}-X_{j,N})^{-1}
                       +\frac2\beta\frac{q_N'(X_{i,N})}{q_N(X_{i,N})}.
\]
Locally uniform convergence of holomorphic functions implies locally
uniform convergence of their derivatives. Because the limiting $q$
is zero-free, all limiting coordinates lie in $B$, and the limit is
simple, the right side is continuous at almost every limiting tuple.
The continuous-mapping theorem proves the force assertion.
\end{proof}

\begin{theorem}[An unconditional canonical quasi-Gibbs law]
\label{thm:unconditional-canonical-quasigibbs}
For every $\beta>0$, every subsequential limit and every disk
$\overline B\subset\D$, the following holds. For almost every exterior
configuration $\eta$, and every integer $m\ge1$ for which
$p_m(\eta):=\mathbb P(M=m\mid\eta)>0$, the conditional law of the
uniformly ordered interior tuple has a density
\begin{equation}\label{eq:unmarked-canonical-quasigibbs}
 \mathbb P\{X\in d\mathbf z\mid\eta,M=m\}
   =|\Delta_m(\mathbf z)|^\beta
                   g_{\eta,m}(\mathbf z)\,dA^m(\mathbf z).
\end{equation}
Here $g_{\eta,m}$ is symmetric, strictly positive and real analytic on
$B^m$, including the collision diagonals. Its logarithm is
plurisubharmonic. For each compact $K\subset B$ there are a
deterministic finite constant $C_{B,K,m,\beta}$ and a measurable random
constant $c_{\eta,K,m,\beta}>0$ such that
\begin{equation}\label{eq:canonical-quasigibbs-two-sided}
 c_{\eta,K,m,\beta}\le g_{\eta,m}(\mathbf z)
                   \le C_{B,K,m,\beta},\qquad \mathbf z\in K^m.
\end{equation}
Thus the conditional law is equivalent to
$|\Delta_m|^\beta dA^m$ on $B^m$ and has full support on the
space of $m$ distinct interior points. For $m=0$, the conditional law
is simply the point mass at the empty configuration.
\end{theorem}

\begin{proof}
For each $m$ with $\mathbb P(M=m)>0$, disintegrate the augmented law
of Theorem~\ref{thm:augmented-canonical-dlr} with respect to
$(\eta,M=m)$, and let $\nu_{\eta,m}$ be the resulting conditional law
of $q$. For almost every such exterior it is supported on zero-free
functions with $Z_{B,m}(q)=1$ and $q(a_0)>0$. Put
\begin{equation}\label{eq:quasigibbs-mixture-factor}
 g_{\eta,m}(\mathbf z)=
      \int\prod_i|q(z_i)|^2\,\nu_{\eta,m}(dq).
\end{equation}
Taking conditional expectation in the augmented DLR equation proves
\eqref{eq:unmarked-canonical-quasigibbs}. The same construction can be
made for all relevant $m$ simultaneously, since there are countably
many values. Exterior sets null under the law given $M=m$ are null
under $p_m(\eta)\,\mathbb P_\eta(d\eta)$, so this proves the stated
almost-every-exterior quantifier on $\{p_m>0\}$.

Let $r>0$ be small enough that $B(a,r)\subset B$ for every $a\in K$.
The normal-family bound gives, for every admissible $q$,
\[
 \prod_i|q(z_i)|^2\le J_{m,\beta}^{-1}r^{-d_{m,\beta}},
       \qquad\mathbf z\in K^m.
\]
This proves the deterministic upper bound. Zero-freeness and
compactness of $K$ give $\min_K|q|>0$ for every admissible $q$, and
\[
 c_{\eta,K,m,\beta}:=
      \int (\min_K|q|)^{2m}\,\nu_{\eta,m}(dq)>0
\]
gives the lower bound. The minimum is a Borel function of $q$ in the
locally uniform topology; the integral is finite by the same upper
bound. An exhaustion of $B$ by compact disks gives one full-measure
set of exteriors on which these assertions and strict positivity hold
throughout $B^m$.

For regularity, write $F_q(\mathbf z)=\prod_i q(z_i)$ and define
\[
 K_{\eta,m}(\mathbf z,\mathbf w)
       =\int F_q(\mathbf z)\overline{F_q(\mathbf w)}
                                         \,\nu_{\eta,m}(dq).
\]
The deterministic bounds on slightly larger compact sets and the
Cauchy integral formula permit differentiation under the integral to
all orders. More explicitly, the Taylor series of $F_q$ on any
polydisk strictly contained in $B^m$ has coefficients uniformly
bounded by the same geometric majorant; integrating the products of
these series proves that $K_{\eta,m}$ is holomorphic in $\mathbf z$
and antiholomorphic in $\mathbf w$. Its diagonal
$g_{\eta,m}(\mathbf z)=K_{\eta,m}(\mathbf z,\mathbf z)$ is therefore
real analytic.

On any complex line and with a prime denoting its holomorphic
directional derivative, differentiation and Cauchy--Schwarz give
\[
 \partial\bar\partial\log g
   =\frac{\bigl(\int|F_q'|^2d\nu\bigr)
                  \bigl(\int|F_q|^2d\nu\bigr)
            -\bigl|\int F_q'\overline{F_q}\,d\nu\bigr|^2}{g^2}
     \ge0.
\]
This proves log-plurisubharmonicity. Finally strict positivity of $g$
and positivity of the Vandermonde off the area-null collision set
prove equivalence and full support.
\end{proof}

\begin{corollary}[Conditional Poisson absolute continuity and local moves]
\label{cor:canonical-local-quasiinvariance}
Conditioned on the exterior, the law in $B$ is absolutely continuous
with respect to a unit-intensity Poisson process on $B$. It is
equivalent to that Poisson law if and only if
$p_m(\eta)>0$ for every integer $m\ge0$. No assertion that these count
probabilities are all positive is made here.

For any $C^1$ diffeomorphism $T:B\to B$ that equals the identity in a
neighborhood of $\partial B$, moving all interior points by $T$ and
fixing the exterior produces a law equivalent to the original limiting
law. On the event $M=m\ge1$, a version of its Radon--Nikodym derivative
at the interior tuple $\mathbf x$ is
\begin{equation}\label{eq:canonical-local-move-density}
 \left|\frac{\Delta_m(T^{-1}\mathbf x)}{\Delta_m(\mathbf x)}\right|^\beta
 \frac{g_{\eta,m}(T^{-1}\mathbf x)}{g_{\eta,m}(\mathbf x)}
 \prod_{i=1}^m |\det D T^{-1}(x_i)|.
\end{equation}
It equals one when $m=0$.
\end{corollary}

\begin{proof}
The unit-intensity Poisson law has the expansion
\[
 \mathbb E_{\rm Pois}F
       =e^{-|B|}\sum_{m=0}^\infty\frac1{m!}
                         \int_{B^m}F\left(\sum_i\delta_{z_i}\right)dA^m.
\]
Thus the conditional density relative to this law, on configurations
with $m\ge1$ points, is
$e^{|B|}m!p_m(\eta)|\Delta_m|^\beta g_{\eta,m}$; for $m=0$ it is
$e^{|B|}p_0(\eta)$. The equivalence criterion follows from strict
positivity in each allowed sector. A diffeomorphism preserves the
count, the exterior and distinctness. The ordinary change-of-variables
formula in each conditional density gives
\eqref{eq:canonical-local-move-density}, which is finite and strictly
positive almost surely. Integrating over counts and exteriors proves
mutual absolute continuity of the full laws.
\end{proof}

\begin{proposition}[The score of the canonical conditional density]
\label{prop:canonical-quasigibbs-score}
Let $b$ be the canonical logarithmic derivative from
Corollary~\ref{cor:canonical-force-score}. For almost every exterior
and each allowed count $m\ge1$, at almost every distinct interior
tuple under its conditional law,
\begin{equation}\label{eq:canonical-quasigibbs-score}
 b(z_i,\eta+\textstyle\sum_j\delta_{z_j})
   =\sum_{j\ne i}\frac1{z_i-z_j}
          +\frac2\beta\partial_{z_i}\log g_{\eta,m}(\mathbf z).
\end{equation}
The last term is the posterior average of the holomorphic exterior
force in the augmented DLR equation:
\[
 \frac2\beta\partial_{z_i}\log g_{\eta,m}(\mathbf z)
   =\frac2\beta
       \int\frac{q'(z_i)}{q(z_i)}
       \frac{\prod_j|q(z_j)|^2}{g_{\eta,m}(\mathbf z)}
                                                \,\nu_{\eta,m}(dq).
\]
This is an identification with a conditional-density score, not with
an exterior-measurable infinite Coulomb sum.
\end{proposition}

\begin{proof}
Away from collisions, differentiate the positive smooth conditional
density in Theorem~\ref{thm:unconditional-canonical-quasigibbs}:
\[
 \frac2\beta\partial_{z_i}\log
       \bigl(|\Delta_m(\mathbf z)|^\beta g_{\eta,m}(\mathbf z)\bigr)
   =\sum_{j\ne i}(z_i-z_j)^{-1}
                  +\frac2\beta\partial_{z_i}\log g_{\eta,m}(\mathbf z).
\]
Compactly supported integration by parts in the interior variables
identifies this with the conditional logarithmic derivative. The
collision singularities are locally integrable against the density,
since their worst radial factor is $|z_i-z_j|^{\beta-1}$ and $\beta>0$.
The factor $g$ and all its derivatives are bounded on compact sets.
Here are details justifying the disintegration of the global Ward
identity. Fix a vector field $v$ compactly supported in $B$. Choose
$\chi_\ell\in C_c^\infty(\D)$ converging pointwise off $\partial B$
to $\mathbf1_B$, equal to one near $\operatorname{supp}v$, with their
gradients supported in shrinking collars of $\partial B$. Then
$\Xi(\chi_\ell)\to M$ almost surely, and
$\mathcal L_v\Xi(\chi_\ell)=0$ exactly. If a bounded smooth scalar
function $\vartheta_m$ equals one at the integer $m$ and zero at all
other nonnegative integers, then
$\vartheta_m(\Xi(\chi_\ell))\to\mathbf1_{\{M=m\}}$ boundedly, also
with zero derivative in the Ward identity. Exterior cylinders use
test functions supported in $\D\setminus\overline B$ and likewise
have zero derivative. The already proved local Campbell $L^2$ bound
on $b$ permits passage to bounded pointwise limits of these
multipliers. A monotone-class argument therefore allows arbitrary
bounded Borel functions of the exterior and the indicator of the
inside count, without introducing any derivative of these factors.

To identify individual coordinates, cover the set of distinct
interior tuples by countably many coordinate charts made of $m$
pairwise disjoint disks with rational data, compactly contained in
$B$. On a chart the configuration has exactly one point in each disk
and no others in $B$. Local count selectors for these disks are
obtained by the same collar approximation; for a vector field
supported in one chosen disk, all their derivatives vanish. Products
of coordinate tests are represented on the chart by products of
linear statistics with their test functions supported in the
respective disjoint disks. Bounded smooth truncations put these tests
in the original cylinder class. Consequently the global identity,
after disintegration, determines the distributional derivative in
each individual coordinate on every such chart. The conditional
density is strictly positive and smooth there, so its logarithmic
derivative must equal the restriction of $b$. Countable exhaustion
proves the assertion almost everywhere. This argument uses the
integrability of the existing $b$ and does not presuppose a global
integrability bound on the candidate $\partial\log g$.
The posterior formula follows by differentiating
\eqref{eq:quasigibbs-mixture-factor}. Its integrability follows from
the deterministic compact bounds on $q$ and $q'$ after cancellation of
the displayed factor $q(z_i)$ with its product weight.
\end{proof}

\begin{corollary}[Strict positivity of every nonzero repulsion factor]
\label{cor:strict-positive-repulsion-factor}
For each $p\ge1$, let $h_p$ be the continuous function from
Theorem~\ref{thm:hierarchy}, so that $\rho_p=|\Delta_p|^\beta h_p$.
Exactly one of the following alternatives holds:
\begin{enumerate}
\item $\mathbb P\{\Xi(\D)\ge p\}=0$, in which case $h_p\equiv0$;
\item $\mathbb P\{\Xi(\D)\ge p\}>0$, in which case
      $h_p(\mathbf z)>0$ for every $\mathbf z\in\D^p$, including
      collision diagonals.
\end{enumerate}
In the second case, every nonempty open $U\subset\D$ has
$\mathbb P\{\Xi(U)\ge p\}>0$, and the leading constants in
Corollary~\ref{cor:hier-cluster-asymptotics} are strictly positive at
every center. In particular every limit that is not almost surely
empty satisfies $\rho_1(z)>0$ throughout $\D$, including at $z=0$.
\end{corollary}

\begin{proof}
The first alternative follows immediately from the factorial moment
identity, followed by continuity of $h_p$. In the second alternative,
fix $\mathbf a\in\D^p$. By exhaustion there is a disk
$\overline B\subset\D$ containing all its coordinates and satisfying
$\mathbb P\{\Xi(B)\ge p\}>0$. Since $\Xi(B)$ is finite almost surely,
there is $m\ge p$ with $P_m:=\mathbb P\{\Xi(B)=m\}>0$.
Average the canonical mixture factor over the exterior under the law
given this count, and write
\[
 \overline g_{B,m}(\mathbf z)
       =\mathbb E[g_{\eta,m}(\mathbf z)\mid\Xi(B)=m].
\]
The same deterministic compact bounds as in the theorem show that
this function is continuous and strictly positive throughout $B^m$.
The contribution to the $p$-point factorial density from the event
$\{\Xi(B)=m\}$ gives, for almost every $\mathbf z\in B^p$,
\[
 \rho_p(\mathbf z)\ge
 P_m(m)_p\int_{B^{m-p}}|\Delta_m(\mathbf z,\mathbf w)|^\beta
          \overline g_{B,m}(\mathbf z,\mathbf w)\,dA^{m-p}(\mathbf w).
\]
When $m>p$, choose $m-p$ pairwise disjoint closed disks of positive
radius in $B$ avoiding the finite set of coordinates of $\mathbf a$,
and restrict the integral to their product. Divide by
$|\Delta_p(\mathbf z)|^\beta$ off the collision diagonals.
The resulting restricted integral, including its cross-distance
factors, is a continuous strictly positive function in a neighborhood
of $\mathbf a$, even when some coordinates of $\mathbf a$ coincide.
When $m=p$ the same statement holds with no integral. Continuity of
$h_p$ extends the resulting lower bound from almost every distinct
tuple to that entire neighborhood, proving $h_p(\mathbf a)>0$.
Integrating the now positive correlation density over distinct tuples
in $U^p$ proves the local count assertion. The previously established
small-cluster formulas then have strictly positive coefficients.
\end{proof}

\begin{remark}[The remaining DLR problem]
The unmarked equation~\eqref{eq:unmarked-canonical-quasigibbs} is an
unconditional canonical conditional-law theorem. It identifies the
repulsion exponent, proves equivalence of all interior locations at
each allowed count, and gives a mixture representation of its smooth
factor. It does not compute that factor from a prescribed regularized
interaction with the exterior, determine the conditional count
probabilities, or prove uniqueness of the limiting law. In particular,
one may not replace $g_{\eta,m}$ by
$\exp(\beta\sum_i h_\eta(z_i))$ merely because every finite-$N$
conditional density has that form: a nontrivial mixture over limiting
harmonic fields generally fails to have such a product form.
\end{remark}

\section{A Gaussian lower bound on the curvature of conditional laws}
\label{sec:conditional-gaussian-curvature}

The harmonic mixture in the augmented DLR equation is necessarily
nontrivial. Its posterior fluctuations contain the independent Gaussian
harmonic field. This observation strengthens mere
log-plurisubharmonicity to an explicit matrix inequality for every
canonical conditional density.

For $\mathbf x=(x_1,\ldots,x_m)\in B^m$ define
\[
 \mathcal K(\mathbf x)_{ij}=\frac1{(1-x_i\overline{x_j})^2},
 \qquad
 \mathcal R_m(\mathbf x)=\prod_{i,j=1}^m(1-x_i\overline{x_j}).
\]
The second expression is real and strictly positive, since it equals
$\prod_i(1-|x_i|^2)\prod_{i<j}|1-x_i\overline{x_j}|^2$.
For Hermitian matrices, $A\succeq B$ means that $A-B$ is positive
semidefinite. We use the Wirtinger convention
$\partial_z=(\partial_x-i\partial_y)/2$, so
$\Delta_z=4\partial_z\partial_{\bar z}$.

\begin{theorem}[Strict conditional Levi curvature]
\label{thm:conditional-gaussian-levi}
In Theorem~\ref{thm:unconditional-canonical-quasigibbs}, for almost
every exterior $\eta$ and every allowed count $m\ge1$, the positive
real-analytic factor $g_{\eta,m}$ satisfies
\begin{equation}\label{eq:conditional-gaussian-levi}
 \left[\partial_{x_i}\partial_{\overline{x_j}}
                \log g_{\eta,m}(\mathbf x)\right]_{i,j=1}^m
       \succeq \frac\beta2\mathcal K(\mathbf x),
       \qquad \mathbf x\in B^m.
\end{equation}
Consequently
\begin{equation}\label{eq:conditional-relative-psh}
 \log\bigl(g_{\eta,m}\mathcal R_m^{\beta/2}\bigr)
 \quad\hbox{is plurisubharmonic on }B^m.
\end{equation}
For the ordinary-area conditional density
$p_{\eta,m}=|\Delta_m|^\beta g_{\eta,m}$, every coordinate satisfies
\begin{equation}\label{eq:conditional-scalar-curvature}
 \Delta_{x_i}\log p_{\eta,m}(\mathbf x)
       \ge \frac{2\beta}{(1-|x_i|^2)^2}
\end{equation}
off the collision diagonals. The matrix on the right of
\eqref{eq:conditional-gaussian-levi} is positive definite whenever
the coordinates are distinct.
\end{theorem}

\begin{proof}
We first describe the required joint realization. In the construction
of Theorem~\ref{thm:augmented-canonical-dlr}, retain the force-marked
configuration and all positive harmonic modes when extracting the
subsequence. Their joint tightness was established in
Theorems~\ref{thm:force-marked-limit} and
\ref{thm:gaussian-force-regression}. For the fixed event $M=m$, the
normalized exterior factors belong to a fixed compact normal family.
They can be set equal to zero off $\{M_N=m\}$ while taking the joint
subsequence; then condition on that event, whose probability converges
to a positive number. The proof of the augmented DLR theorem applies
without change. Its force identity
\eqref{eq:augmented-harmonic-force} therefore holds for marks that
also obey the Gaussian regression
\eqref{eq:gaussian-force-regression}. Uniform ordering can be supplied
by auxiliary randomness independent of this entire construction.

Fix an exterior and an allowed count for which the conclusions of
the augmented and unmarked DLR theorems hold, and abbreviate the
mixing measure by $\nu=\nu_{\eta,m}$. Put
\[
 A_q(\mathbf x)=\prod_{i=1}^m q(x_i),\qquad
 g(\mathbf x)=\int |A_q(\mathbf x)|^2\,\nu(dq).
\]
The posterior distribution of $q$ given the complete unmarked
configuration, written as $(\eta,m,\mathbf x)$ with its uniform
ordering, is
\begin{equation}\label{eq:conditional-q-posterior}
 \nu_{\mathbf x}(dq)=\frac{|A_q(\mathbf x)|^2}{g(\mathbf x)}\,\nu(dq).
\end{equation}
This is Bayes' formula applied to
\eqref{eq:augmented-canonical-dlr}; the Vandermonde cancels because
it does not depend on $q$. All statements made initially almost
everywhere below can be taken with respect to this posterior.

Let $L_i(q,\mathbf x)=q'(x_i)/q(x_i)$. The zero-free property makes
this well defined. Differentiation under the mixing integral yields
\begin{equation}\label{eq:conditional-levi-as-covariance}
 \partial_{x_i}\partial_{\overline{x_j}}\log g
  =\operatorname{Cov}_{\nu_{\mathbf x}}
                          (L_i,\overline{L_j}).
\end{equation}
For clarity, the covariance on the right means
$\mathbb E[L_i\overline{L_j}]-\mathbb E[L_i]\overline{\mathbb E[L_j]}$.
The differentiation is justified even without a uniform bound on
$q'/q$. Indeed $A_qL_i=\partial_{x_i}A_q$, and the deterministic
normal-family bounds, followed by the Cauchy integral formula on
slightly larger compact polydisks, bound both $A_q$ and all its
first derivatives uniformly in $q$. Thus all integrals of
$|A_q|^2L_i\overline{L_j}$ are absolutely convergent, locally uniformly
in $\mathbf x$, and the asserted derivative formula follows.

Off collisions, the force identity gives
\[
 C_i=\sum_{j\ne i}(x_i-x_j)^{-1}+\frac2\beta L_i.
\]
The first term is fixed by the unmarked configuration. Hence
\eqref{eq:conditional-levi-as-covariance} becomes
\begin{equation}\label{eq:conditional-levi-force-covariance}
 \left[\partial_{x_i}\partial_{\overline{x_j}}\log g\right]_{ij}
       =\frac{\beta^2}{4}
          \operatorname{Cov}\bigl((C_i)_{i=1}^m\mid\Xi,\mathbf x\bigr).
\end{equation}
Here the notation $\mathbf x$ in the conditioning only retains the
independent ordering of points already contained in $\Xi$.

For any integer $L\ge1$, Gaussian regression gives simultaneously
at all the finitely many interior particles
\[
 \mathbb E[C_i\mid\Xi,\mathbf x,G_1,\ldots,G_L]
   =b(x_i,\Xi)-\sum_{k=1}^L\overline{G_k}x_i^{k-1}.
\]
The Campbell almost-everywhere assertion in that theorem holds
simultaneously at these particles: the expected number of exceptions
in any compact subset is zero, and a countable exhaustion covers $B$.
The modes are independent of $\Xi$ and of its independent ordering.
Conditional orthogonal projection, applied to every complex linear
combination of the $C_i$, therefore gives
\[
 \operatorname{Cov}\bigl((C_i)_{i=1}^m\mid\Xi,\mathbf x\bigr)
 \succeq \frac2\beta
       \left[\sum_{k=1}^L kx_i^{k-1}\overline{x_j}^{\,k-1}\right]_{ij}.
\]
The conditional second moments are finite almost surely, either by
local Campbell integrability or directly from the preceding
normal-family derivative bound. Letting $L\to\infty$, the geometric
series converges locally uniformly and gives $\mathcal K$. Combining
this with \eqref{eq:conditional-levi-force-covariance} proves
\eqref{eq:conditional-gaussian-levi} initially almost everywhere
under the conditional density. That density is strictly positive
off collisions, so the same assertion holds Lebesgue-almost
everywhere there. Both matrix sides are continuous on all of $B^m$;
a countable dense set of complex test vectors then extends the
inequality everywhere, including the diagonals.

Direct differentiation gives
\[
 \partial_{x_i}\partial_{\overline{x_j}}
                       \log\mathcal R_m=-\mathcal K(\mathbf x)_{ij}.
\]
This proves \eqref{eq:conditional-relative-psh}. The Vandermonde
logarithm is pluriharmonic off its zero set, so its Levi matrix there
vanishes. Taking the diagonal in
\eqref{eq:conditional-gaussian-levi} and multiplying by four proves
\eqref{eq:conditional-scalar-curvature}.
Finally $\mathcal K$ is the Gram matrix of the vectors
$(\sqrt{k}x_i^{k-1})_{k\ge1}$. For distinct coordinates the first
$m$ entries already form an invertible Vandermonde matrix. Thus
this Gram matrix is positive definite.
\end{proof}

\begin{corollary}[A harmonic exterior alone cannot describe the conditional law]
\label{cor:no-pure-harmonic-unmarked-dlr}
For any allowed count $m\ge1$, the unmarked conditional density cannot
have the form
\[
 p_{\eta,m}(\mathbf x)
   =Z^{-1}|\Delta_m(\mathbf x)|^\beta
             \exp\left\{\beta\sum_{i=1}^m h_\eta(x_i)\right\}
\]
with $h_\eta$ harmonic on $B$. In the augmented DLR construction,
the exterior force $q'/q$ cannot be measurable with respect to the
complete unmarked configuration on any event of positive probability
that is measurable with respect to that configuration and on which
$M\ge1$.
\end{corollary}

\begin{proof}
In the displayed product representation the logarithm of the smooth
factor has zero Levi matrix, contradicting the strictly positive
diagonal in \eqref{eq:conditional-gaussian-levi}. For the second
assertion, conditional on the unmarked configuration the covariance
of $q'/q$ at each interior point is at least
$\beta/(2(1-|x_i|^2)^2)>0$, by
\eqref{eq:conditional-levi-as-covariance} and the theorem. A measurable
force would have zero conditional variance. The contradiction remains
valid after restriction to any positive-probability event measurable
with respect to the unmarked configuration and contained in $M\ge1$.
\end{proof}

\begin{proposition}[A necessary background bound for a compensated Green DLR law]
\label{prop:conditional-green-background-bound}
Let a nonempty subsequential limit satisfy a canonical Green DLR
specification relative to
$d\lambda(x)=(1-|x|^2)^{-2}dA(x)$, with conditional densities
\begin{equation}\label{eq:canonical-green-density-for-bound}
 p_{\eta,m}(\mathbf x)
  =Z_{\eta,m}^{-1}
    \exp\left\{-\beta\sum_{i<j}g_{\D}(x_i,x_j)
                   -\beta\sum_i H_\eta(x_i)\right\}
                   \prod_i(1-|x_i|^2)^{-2},
\end{equation}
where
$g_{\D}(x,y)=-\log|(x-y)/(1-x\overline y)|$.
Assume that the finite exterior field is twice continuously
differentiable in $B$ and satisfies
\[
 \Delta H_\eta(x)=\frac{2\pi c}{(1-|x|^2)^2},
 \qquad x\in B,
\]
for a constant $c\ge0$. Then necessarily
\begin{equation}\label{eq:green-background-density-bound}
 c\le\frac{4-\beta}{\pi\beta}.
\end{equation}
In particular, if $c$ is the intensity with respect to hyperbolic
area in an isometry-invariant compensated DLR description, that
intensity must obey this bound. This is a necessary condition for
such a description; no existence of the assumed field or DLR
specification is inferred from it.
\end{proposition}

\begin{proof}
Since the law is not almost surely empty, some relatively compact
disk has positive probability of an allowed count $m\ge1$. On such
a sector, $g_{\D}(x_i,x_j)$ is harmonic in $x_i$ away from $x_j$.
Moreover
$\Delta[-2\log(1-|x|^2)]=8/(1-|x|^2)^2$.
Taking the coordinate Laplacian of the logarithm of
\eqref{eq:canonical-green-density-for-bound} therefore gives
\[
 \Delta_{x_i}\log p_{\eta,m}(\mathbf x)
   =\frac{8-2\pi\beta c}{(1-|x_i|^2)^2}
\]
off collisions. Comparison with
\eqref{eq:conditional-scalar-curvature} proves
$8-2\pi\beta c\ge2\beta$, which is exactly
\eqref{eq:green-background-density-bound}.
\end{proof}

\begin{remark}[Why the compensated field has this Laplacian]
For an exterior configuration supported in $B^c$, finite sums
$\sum_w g_{\D}(x,w)$ are harmonic in $B$. A compensated field formed
by subtracting $c\int_Eg_{\D}(x,w)\,d\lambda(w)$, with $E$
containing $B$, has distributional Laplacian
$2\pi c(1-|x|^2)^{-2}$ there, because
$-\Delta_xg_{\D}(x,w)=2\pi\delta_w$. If such compensated fields
converge in local $L^1(B)$, their limit retains this identity.
Subtracting the smooth particular solution
$-(\pi c/2)\log(1-|x|^2)$ leaves a distributionally harmonic
function, which has a smooth harmonic representative. Thus the
regularity and Laplacian assumptions of the proposition follow from
this specific local convergence, if it can be proved. The required
convergence has not been established here for general $\beta$.
\end{remark}

\begin{corollary}[Saturation determines the residual force]
\label{cor:conditional-green-force-saturation}
Suppose the canonical Green DLR assumptions of
Proposition~\ref{prop:conditional-green-background-bound} hold on
all relatively compact interior disks, with
$c=(4/\beta-1)/\pi$. In any joint force-and-Gaussian limit as in
Theorem~\ref{thm:gaussian-force-regression}, almost surely every
particle satisfies
\begin{equation}\label{eq:saturated-force-decomposition}
 C_z=b(z,\Xi)-\sum_{k=1}^\infty\overline{G_k}z^{k-1}.
\end{equation}
The series converges locally in Campbell $L^2$. Thus these force marks
are determined by the unmarked configuration and its independent
Gaussian harmonic field, with no additional force randomness.
At $\beta=2$ this conclusion is unconditional: the Bergman DLR law
of Theorem~\ref{thm:bergman-grandcanonical-dlr} and the potential
regularity in Corollary~\ref{cor:bergman-conditional-curvature}
verify the hypotheses with $c=1/\pi$.
\end{corollary}

\begin{proof}
For $i\ne j$, the smooth factor in
\eqref{eq:canonical-green-density-for-bound} has mixed Levi entry
\[
 \partial_{x_i}\partial_{\overline{x_j}}\log g_{\eta,m}
          =\frac\beta{2(1-x_i\overline{x_j})^2}.
\]
Indeed the pair contribution is
$-\beta\log|1-x_i\overline{x_j}|$; its holomorphic--antiholomorphic
part has coefficient $-\beta/2$. No one-body term contributes to
these mixed entries. The diagonal entry is
\[
 \frac{2-\pi\beta c/2}{(1-|x_i|^2)^2}
       =\frac\beta{2(1-|x_i|^2)^2}
\]
at the specified value of $c$. Thus equality holds throughout
\eqref{eq:conditional-gaussian-levi}. Equation
\eqref{eq:conditional-levi-force-covariance} shows that the conditional
force covariance is exactly $(2/\beta)\mathcal K$.

Write $T(z)=\sum_{k\ge1}\overline{G_k}z^{k-1}$. The infinite-mode
regression gives
$\mathbb E[C_z\mid\Xi,(G_k)]=b(z,\Xi)-T(z)$.
Conditional orthogonal projection therefore decomposes each force
variance into the variance $2/\{\beta(1-|z|^2)^2\}$ of $T(z)$ and
the nonnegative mean square of
$C_z-b(z,\Xi)+T(z)$. The preceding covariance equality makes that
mean square zero at every interior particle, in the almost-sure
conditional sense. Restricting first to compact Campbell windows
justifies all $L^2$ computations. A countable exhaustion by interior
disks proves \eqref{eq:saturated-force-decomposition} at every
particle almost surely. The claimed series convergence is exactly
that established in Theorem~\ref{thm:gaussian-force-regression}.
\end{proof}

\section{A consistent augmentation by a meromorphic electric field}
\label{sec:consistent-meromorphic-augmentation}

The one-disk augmented equations can be made simultaneous and
consistent. The field retained here is additional random data. Its
existence does not assert that the exterior of the unmarked point
process determines the conditional interaction.

Choose once and for all $0<r_1<r_2<\cdots\uparrow1$ and put
$B_j=\D_{r_j}$. Almost surely the limiting configuration has no point
on any of the circles $\partial B_j$.

\begin{proposition}[Simultaneous fields and meromorphic gluing]
\label{prop:meromorphic-field-gluing}
For every subsequential limit $\Xi$ there is an augmentation by random
functions $q_j\in\mathcal H(B_j)$ with the following properties.
Writing $m_j=\Xi(B_j)$, set $q_j=0$ if $m_j=0$; if $m_j\ge1$, then
$q_j$ is zero-free, $q_j(0)>0$, and $Z_{B_j,m_j}(q_j)=1$.
The augmented canonical DLR equation of
Theorem~\ref{thm:augmented-canonical-dlr} holds for every $j$ and every
allowed positive value of $m_j$ on this single probability space.

For $j<k$ with $m_j\ge1$, one has on $B_j$
\begin{equation}\label{eq:nested-exterior-field-compatibility}
 \frac2\beta\frac{q_j'(z)}{q_j(z)}
   =\frac2\beta\frac{q_k'(z)}{q_k(z)}
        +\sum_{w\in\Xi\cap(B_k\setminus B_j)}\frac1{z-w}.
\end{equation}
Consequently, on the event $\Xi\ne\varnothing$, the meromorphic
functions
\begin{equation}\label{eq:local-total-meromorphic-field}
 \mathcal C_j(z)=\sum_{w\in\Xi\cap B_j}\frac1{z-w}
                      +\frac2\beta\frac{q_j'(z)}{q_j(z)},
                    \qquad m_j\ge1,
\end{equation}
glue to one meromorphic function $\mathcal C$ on $\D$. It has exactly
the particles of $\Xi$ as its poles, each simple with residue one.
Its regular part at a particle equals the corresponding augmented
force mark. On $\{\Xi=\varnothing\}$ we may set $\mathcal C=0$.

If $m_j=0$ but $\Xi\ne\varnothing$, its missing exterior field is
recovered by restricting this $\mathcal C$ to $B_j$; it is holomorphic
there. More generally
\begin{equation}\label{eq:glued-augmented-exterior-field}
 \mathcal E_j(z)=\mathcal C(z)
              -\sum_{w\in\Xi\cap B_j}\frac1{z-w}
\end{equation}
extends holomorphically to a neighborhood of $\overline B_j$.
For $m_j\ge1$, it is $(2/\beta)q_j'/q_j$ on $B_j$.
\end{proposition}

\begin{proof}
At finite $N$, define $q_{j,N}$ by the normalization in the proof of
Theorem~\ref{thm:augmented-canonical-dlr} on
$\{\Xi_N(B_j)=m\ge1\}$, and set it to zero when the count is zero.
For each fixed $j$ the count laws are tight. For any finite collection
of allowed counts, the corresponding $q_{j,N}$ belong to a finite
union of compact normal-family closures from
Lemma~\ref{lem:normalized-exterior-normal-family}. This proves joint
tightness of the count and function. A countable diagonal extraction
therefore retains $\Xi$ and all the functions $q_j$ jointly. Conditional
on any positive-probability event $\{m_j=m\ge1\}$, the proof of the
one-disk theorem applies without alteration: testing only the exterior
and $q_j$ still gives its exact conditional density. The argument with
compactly supported test functions increasing to one proves
$Z_{B_j,m}(q_j)=1$ on that event. Countably many $j,m$ give one
probability-one set on which all assertions hold simultaneously.

Before taking limits, logarithmic differentiation of the exterior
products gives \eqref{eq:nested-exterior-field-compatibility} exactly.
The annulus between two fixed circles is relatively compact in $\D$.
Its particles converge as a finite configuration and none lies on its
boundary in the limit. On $\{m_j\ge1\}$ both normalized holomorphic
functions have nonzero zero-free limits, so their logarithmic
derivatives converge locally uniformly. The finite annular Cauchy
sum converges locally uniformly on $B_j$. This proves the displayed
compatibility. The countable collection of identities holds on one
probability-one set.

If $\Xi\ne\varnothing$, there is some first disk with positive count,
and all larger disks also have positive count. The compatibility
identity says exactly that their functions $\mathcal C_j$ agree on
overlaps. They therefore define $\mathcal C$ on their union, which is
all of $\D$. Each local expression has simple poles of residue one
at the interior particles and is holomorphic elsewhere. No
cancellation can remove a pole, because the other terms in that local
expression are holomorphic near that particle. Its regular part is
the expression in \eqref{eq:augmented-harmonic-force}.

This also explains the construction for any smaller zero-count disk:
use a larger positive-count disk to define $\mathcal C$ and restrict
it. The result is independent of the chosen larger disk. Finally,
local finiteness and the absence of particles on $\partial B_j$ give
an open neighborhood of $\overline B_j$ containing no particles other
than the finitely many inside $B_j$. Subtracting their simple pole
terms removes every pole there, proving the asserted extension of
$\mathcal E_j$. The construction is measurable: the first positive
count disk is a measurable integer, and values on a countable dense
subset, away from the almost surely absent fixed-point poles,
determine the glued meromorphic function. Equivalently one may use
the usual Borel space of meromorphic functions with the locally
uniform spherical topology.
\end{proof}

\begin{theorem}[A consistent canonical specification on the augmented space]
\label{thm:consistent-augmented-canonical-specification}
On the augmented law of Proposition~\ref{prop:meromorphic-field-gluing},
the conditional distribution in $B_j$ given
\[
 \Xi|_{\D\setminus\overline B_j},\qquad m_j,\qquad
                         \mathcal E_j|_{B_j}
\]
is the following canonical kernel. If $m_j=0$, it fixes the state. If
$m_j=m\ge1$, let $q$ be the unique zero-free holomorphic function on
$B_j$ with
\[
 \frac2\beta\frac{q'}q=\mathcal E_j,
       \qquad q(0)>0,\qquad Z_{B_j,m}(q)=1.
\]
Sample a uniformly ordered tuple $\mathbf x\in B_j^m$ with density
\begin{equation}\label{eq:consistent-augmented-kernel}
 |\Delta_m(\mathbf x)|^\beta\prod_i|q(x_i)|^2\,dA^m(\mathbf x).
\end{equation}
Keep the exterior fixed and replace the meromorphic field by
\begin{equation}\label{eq:augmented-field-resampling-update}
 \mathcal C^{\rm new}(z)=\mathcal C(z)
       +\sum_{i=1}^m\frac1{z-x_i}
       -\sum_{w\in\Xi\cap B_j}\frac1{z-w}.
\end{equation}
Denote this kernel on augmented states by $\gamma_j$. Its result
satisfies the same admissibility conditions as the original state,
and the augmented law $\widehat P$ satisfies
\[
 \widehat P\gamma_j=\widehat P\qquad\hbox{for every }j.
\]
For $j<k$, the canonical kernels are consistent:
\begin{equation}\label{eq:augmented-canonical-consistency}
 \gamma_j\gamma_k=\gamma_k\gamma_j=\gamma_k.
\end{equation}
These are count-preserving equations. They do not determine the
probabilities of the different inside counts, and their conditioning
includes an exterior-field mark in addition to the unmarked exterior.
\end{theorem}

\begin{proof}
A holomorphic function on a disk has a holomorphic primitive.
Exponentiating $(\beta/2)$ times the primitive of $\mathcal E_j$,
chosen zero at the origin, determines $q$ up to a positive constant;
the displayed phase and normalization determine that constant uniquely.
It exists and is finite because $\mathcal E_j$ extends to a
neighborhood of $\overline B_j$. Thus, when $m_j\ge1$, knowing
$(m_j,\mathcal E_j)$ is equivalent to knowing $(m_j,q_j)$. The one-disk
augmented DLR equation proves the asserted conditional distribution
of the interior tuple.

One must also check that this describes the conditional law of the
whole augmented state, rather than just of the point positions. Given
the exterior, count, field $\mathcal E_j$ and an interior tuple, the
formula
\[
 \mathcal C(z)=\mathcal E_j(z)+\sum_{i=1}^{m_j}(z-x_i)^{-1}
                    \quad\hbox{on }B_j
\]
determines the global meromorphic function uniquely by the identity
theorem. Existence of that continuation is part of the original
augmentation. For a resampled tuple it is supplied explicitly by
\eqref{eq:augmented-field-resampling-update}. The meromorphic
continuation, on the set where it exists, is measurable: restriction
of meromorphic functions from $\D$ to $B_j$ is an injective Borel map
between standard Borel spaces, and hence has a Borel inverse on its
image. Thus no further random data in the augmented state need be
specified once these quantities and the tuple are known. This proves
$\widehat P\gamma_j=\widehat P$.

The update cancels precisely the old interior pole terms and inserts
precisely the new ones, preserving the residue-one condition and
leaving $\mathcal E_j$ unchanged. The sampled tuple is distinct and
avoids all the countably many circles almost surely by absolute
continuity. For any smaller disk, subtracting its new interior pole
terms leaves a function holomorphic near its closed disk; for any
larger disk its exterior field is unchanged. All required positive
finite normalizations therefore continue to exist. This verifies
stability of the admissible augmented state space.

For consistency, first observe that a move in $B_j\subset B_k$
preserves the exterior of $B_k$, its total count $m_k$, and its exterior
field $\mathcal E_k$: the two finite Cauchy-sum changes cancel in
\eqref{eq:glued-augmented-exterior-field}. Consequently $\gamma_k$ is the
same kernel before and after applying $\gamma_j$, proving
$\gamma_j\gamma_k=\gamma_k$.

For the reverse order, suppose $m_k\ge1$ and condition within the
$\gamma_k$ density on the particles in $B_k\setminus\overline B_j$
and on the number $m_j$ of remaining particles. The finite
Vandermonde factorization gives a conditional density proportional to
\[
 |\Delta_{m_j}(\mathbf x)|^\beta
       \prod_{i=1}^{m_j}|q_k(x_i)|^2
       \prod_{\substack{i\le m_j\\ w\text{ in the annulus}}}
                                                   |x_i-w|^\beta.
\]
Its holomorphic logarithmic derivative is exactly the nested field
in \eqref{eq:nested-exterior-field-compatibility}; its normalization
is therefore \eqref{eq:consistent-augmented-kernel} for $B_j$.
Conditionally resampling from this density leaves the $\gamma_k$
law unchanged, by ordinary conditional expectation. The field update
is the same finite Cauchy-sum update for both constructions. This
proves $\gamma_k\gamma_j=\gamma_k$. If either relevant inside count
is zero, the corresponding inner resampling is the identity and the
same argument applies. In particular when $m_k=0$ both kernels are
the identity on the interior. This proves all the stated consistency
relations.
\end{proof}

\begin{remark}[What is still extra information]
The field $\mathcal C$ contains the particle positions as its poles,
but may contain an additional holomorphic part coming from the
boundary. Conditioning on $\mathcal E_j$ retains that part. The
consistent augmented specification is therefore not an assertion that
the unmarked exterior alone determines a single harmonic interaction.
Integrating out the field gives precisely the canonical quasi-Gibbs
mixture from Theorem~\ref{thm:unconditional-canonical-quasigibbs}.
\end{remark}

\section{Hyperbolic covariance of the score and of the Gaussian force}
\label{sec:hyperbolic-score-covariance}

Write $\lambda(dz)=(1-|z|^2)^{-2}dA(z)$. This differs by a
constant factor from the area measure of the curvature $-1$ metric;
the group of isometries is unchanged. In this section an automorphism
means an orientation-preserving conformal automorphism of $\D$.
Reflection invariance, which is already known for every limiting law,
then supplies all orientation-reversing isometries as well.

\begin{proposition}[The exact geometric transformation of a score]
\label{prop:hyperbolic-score-transform}
Let $P$ be a point-process law having the locally Campbell-integrable
canonical score $b$ in the normalization of
Corollary~\ref{cor:canonical-force-score}. For an automorphism $\varphi$,
let $P^\varphi$ be the law of $\varphi^{-1}_*\Xi$, where $\Xi$ has law
$P$. Then a canonical score for $P^\varphi$ is
\begin{equation}\label{eq:general-hyperbolic-score-transform}
 b^\varphi(w,\eta)
 =\varphi'(w)b(\varphi(w),\varphi_*\eta)
       +\frac2\beta\frac{\varphi''(w)}{\varphi'(w)}.
\end{equation}
Define the score relative to hyperbolic area by
\[
 b_{\rm H}(z,\eta)
 =b(z,\eta)-\frac4\beta\frac{\overline z}{1-|z|^2}.
\]
It transforms without the inhomogeneous term:
\begin{equation}\label{eq:hyperbolic-score-one-form}
 b_{\rm H}^{\varphi}(w,\eta)
 =\varphi'(w)b_{\rm H}(\varphi(w),\varphi_*\eta).
\end{equation}
Equivalently, for a compactly supported vector field $v$, the Ward
identity is
\begin{equation}\label{eq:hyperbolic-relative-ward}
 \E\left[\mathcal L_vF(\Xi)
  +F(\Xi)\sum_{z\in\Xi}
       \bigl\{\operatorname{div}_{\lambda}v(z)
             +\beta\operatorname{Re}(v(z)b_{\rm H}(z,\Xi))\bigr\}\right]=0,
\end{equation}
where
$\operatorname{div}_{\lambda}v=\operatorname{div}v+
4\operatorname{Re}(v\overline z)/(1-|z|^2)$.
\end{proposition}

\begin{proof}
The real score is $B=\beta\overline b$. Change variables by $z=\varphi(w)$
in its Campbell integration-by-parts identity. The chain rule and the
change of area give the transformed real score
\[
 B^\varphi(w,\eta)
  =D\varphi(w)^T B(\varphi(w),\varphi_*\eta)
       +\nabla_w\log|\varphi'(w)|^2.
\]
This calculation can first be made with the spatial support contained
in a fixed compact set. The map and its inverse have bounded derivatives
there, so all integrations and their changes of variables are justified
by the assumed local integrability. In complex notation,
$D\varphi^TB=\overline{\varphi'}B$ and
$\nabla\log|\varphi'|^2=2\overline{\varphi''/\varphi'}$.
Conjugation and division by $\beta$ prove
\eqref{eq:general-hyperbolic-score-transform}. Uniqueness of the score
identifies this version with the canonical version whenever the latter
has already been constructed.

The identity
\[
 \frac{|\varphi'(w)|^2}{(1-|\varphi(w)|^2)^2}
       =\frac1{(1-|w|^2)^2}
\]
and its logarithmic derivative give
\[
 \varphi'(w)\frac4\beta
       \frac{\overline{\varphi(w)}}{1-|\varphi(w)|^2}
       +\frac2\beta\frac{\varphi''(w)}{\varphi'(w)}
 =\frac4\beta\frac{\overline w}{1-|w|^2}.
\]
Subtracting this identity proves \eqref{eq:hyperbolic-score-one-form}.
Substitution of the definition of $b_{\rm H}$ in the original Ward
identity proves \eqref{eq:hyperbolic-relative-ward}.
\end{proof}

\begin{proposition}[The Gaussian force already has full isometry covariance]
\label{prop:gaussian-force-hyperbolic-covariance}
Let $G_k$ be the independent circular Gaussians of
Theorem~\ref{thm:gaussian-force-regression}, and define
\[
 T(z)=\sum_{k\ge1}\overline{G_k}z^{k-1}.
\]
This series defines a random holomorphic function on $\D$. For every
automorphism $\varphi$,
\begin{equation}\label{eq:gaussian-force-isometry-law}
 \bigl(\varphi'(z)T(\varphi(z))\bigr)_{z\in\D}
       \stackrel{\rm law}=\bigl(T(z)\bigr)_{z\in\D}.
\end{equation}
Thus the independent Gaussian summand in the force regression is
already a covariant holomorphic one-form. This assertion does not
assume, or prove, that the point-process law is invariant.
\end{proposition}

\begin{proof}
For every $r<1$,
$\sum_k \E|G_k|r^{k-1}<\infty$, since
$\E|G_k|\le\sqrt{2k/\beta}$.
Tonelli's theorem gives almost sure absolute and uniform convergence on
$\{|z|\le r\}$, and a countable exhaustion gives holomorphy. The field
is centered circular Gaussian with covariance
\[
 \E[T(z)\overline{T(w)}]
   =\frac2\beta\sum_{k\ge1}k(z\overline w)^{k-1}
   =\frac{2/\beta}{(1-z\overline w)^2},
 \qquad \E[T(z)T(w)]=0.
\]
The elementary automorphism identity
\[
 \frac{\varphi'(z)\overline{\varphi'(w)}}
      {(1-\varphi(z)\overline{\varphi(w)})^2}
      =\frac1{(1-z\overline w)^2}
\]
therefore proves equality of all finite-dimensional Gaussian laws.
Evaluations on a countable dense set generate the Borel sigma-algebra
of $\mathcal H(\D)$, proving the asserted equality of random
holomorphic functions.
\end{proof}

\begin{theorem}[Necessary consequences of hyperbolic invariance]
\label{thm:hyperbolic-invariant-score-constraint}
Suppose a subsequential limiting law $P$ is invariant under every
automorphism and is not concentrated on the empty configuration. Then
there are constants $c>0$ and $v>0$ such that
\begin{align}
 \rho_1(z)&=\frac{c}{(1-|z|^2)^2},
       \label{eq:invariant-intensity-constant}\\
 \E[b(z,\Xi)\mid z]&=\frac4\beta\frac{\overline z}{1-|z|^2},
       \label{eq:invariant-score-mean}\\
 \operatorname{Var}(b(z,\Xi)\mid z)
       &=\frac{v}{(1-|z|^2)^2}\quad\hbox{for area-almost every }z,
       \label{eq:invariant-score-variance}\\
 \intertext{where}
 0<v&\le\frac{8-2\beta}{\beta^2}.
       \label{eq:invariant-score-budget}
\end{align}
The conditional expectations refer to the spatial disintegration of
the Campbell measure. Moreover
$P(0<\Xi(\D)<\infty)=0$.
\end{theorem}

\begin{proof}
Invariance and the continuous intensity imply
$\rho_1(w)=|\varphi'(w)|^2\rho_1(\varphi(w))$.
Choosing an automorphism that sends $w$ to zero proves
\eqref{eq:invariant-intensity-constant}, with $c=\rho_1(0)$; if $c=0$,
the process is empty almost surely by compact exhaustion. The Ward
identity with $F=1$ says
$\beta\overline{\E[b(z,\Xi)\mid z]}=\nabla\log\rho_1(z)$,
which proves \eqref{eq:invariant-score-mean}.

By score uniqueness, invariance identifies the score in
\eqref{eq:general-hyperbolic-score-transform} with $b$.
The deterministic second term does not change conditional variances;
therefore, with $V(z)=\operatorname{Var}(b(z,\Xi)\mid z)$,
\[
 V(w)=|\varphi'(w)|^2 V(\varphi(w))
 \quad\hbox{for almost every }w.
\]
In particular $V(z)dA(z)$ is an invariant locally finite measure. Here
local finiteness follows from the Campbell $L^2$ bound and the positive
lower bound of $\rho_1$ on compact sets.
For completeness, an invariant locally integrable density $V$ must be
$v(1-|z|^2)^{-2}$: differentiate its invariance against a compactly
supported smooth test along the flows of $1-z^2$ and $i(1+z^2)$.
In distributions on $z=x+iy$ this gives
\[
 \begin{pmatrix}1-x^2+y^2&-2xy\\-2xy&1+x^2-y^2\end{pmatrix}
 \binom{\partial_xV}{\partial_yV}
       =4V\binom{x}{y}.
\]
The matrix determinant is $1-|z|^4>0$. Its inverse gives
$\nabla V=4(x,y)V/(1-|z|^2)$, hence
$\nabla[(1-|z|^2)^2V]=0$. A locally integrable function with zero
distributional gradient on the connected disk is constant: convolve
on relatively compact disks, where the convolution is constant, and
let the smoothing parameter tend to zero. This proves
\eqref{eq:invariant-score-variance} with $v\ge0$.

The residual-variance inequality
\eqref{eq:curvature-residual-variance}, together with
$\Delta\log\rho_1=8/(1-|z|^2)^2$, gives
$8\ge2\beta+\beta^2v$. If $v=0$, then
$\beta\overline{b(z,\Xi)}=\nabla\log\rho_1(z)$ Campbell-almost
everywhere. The canonical Ward identity would consequently have a
deterministic score. The continuous two-point correlation and its
vanishing on the diagonal, proved earlier for every subsequential
limit, allow application of Lemma~\ref{lem:deterministic-score-rigidity}.
That lemma implies that the process has at most one point almost
surely. This contradicts
$\E\Xi(\D)=\int_\D c(1-|z|^2)^{-2}dA=\infty$.
Thus $v>0$ and \eqref{eq:invariant-score-budget} follows.

Finally, the event $\{\Xi(\D)=m\}$ is invariant for every finite $m$.
If it had positive probability for $m\ge1$, the conditional law on that
event would have an invariant intensity measure of total mass $m$.
This measure has a locally integrable density, since it is bounded by
the unconditional intensity divided by that positive probability.
The preceding invariant-density argument makes it a constant multiple
of $\lambda$. This is incompatible with finite positive total mass.
Taking the union over $m\ge1$ proves the final assertion.
\end{proof}

\begin{remark}[Why a score identity does not determine the DLR count weights]
Even an explicit canonical score does not in general determine the
activity. For every $a>0$, the Poisson process with intensity
$a\lambda$ has the same score
$b(z,\eta)=4\overline z/[\beta(1-|z|^2)]$.
Indeed, on a compact window its conditional density given its count
$m$ is proportional to $\prod_{j=1}^m(1-|z_j|^2)^{-2}$; ordinary
integration by parts in each coordinate gives this score. The factors
$a^m/m!$ affect the count law and disappear under the spatial
derivative. These Poisson laws are not candidates for the repulsive
limits considered here, but they give an exact demonstration of the
information that a spatial Ward identity omits. A full DLR theorem must
also identify the conditional probabilities of the different counts.
\end{remark}

\begin{remark}[What remains unproved]
The invariance of $T$ and the transformation law for $b_{\rm H}$ are
unconditional statements. Invariance of $P$ is an additional hypothesis
in Theorem~\ref{thm:hyperbolic-invariant-score-constraint}.
For $\beta=2$ the Bergman kernel proves it, and for $\beta\ge4$ it follows
from the empty-limit theorem. For $0<\beta<4$, $\beta\ne2$, neither the
score covariance calculation nor the Gaussian-force regression removes
the missing $o(N^{-1})$ estimate in
Corollary~\ref{cor:mobius-invariance-criterion}. A fixed automorphism
still produces a holomorphic tilt with a coefficient of order $N$.
\end{remark}

\section{Why an invariant DLR equation must be renormalized}
\label{sec:invariant-dlr-obstruction}

The distinction between a Green interaction and a compensated Green
interaction is essential.  The following obstruction applies to every
isometry-invariant point process, without a moment assumption.
Write
\[
 \delta(a,w)=\left|\frac{a-w}{1-a\overline w}\right|,
 \qquad g(a,w)=-\log\delta(a,w),
 \qquad d\lambda(a)=\frac{dA(a)}{(1-|a|^2)^2}.
\]

\begin{theorem}[No nonempty invariant Blaschke configuration]
\label{thm:invariant-green-divergence}
Let $\Xi$ be a locally finite point process on $\mathbb D$ whose law is
invariant under all disk automorphisms. Almost surely, precisely one of
the following alternatives holds:
\[
 \Xi=\varnothing;
 \qquad
 \sum_{w\in\Xi}g(a,w)=+\infty
       \quad\hbox{for every }a\in\mathbb D.
\]
Consequently a nonempty invariant process cannot be supported on
Blaschke configurations, that is, configurations for which
$\sum_w(1-|w|)<\infty$.
\end{theorem}

\begin{proof}
For a deterministic locally finite configuration $\xi$ define
\[
 p_\xi(a)=\prod_{w\in\xi}\delta(a,w)
   =\lim_{R\uparrow1}\prod_{w\in\xi:\,|w|\le R}\delta(a,w)
   =\exp\left\{-\sum_{w\in\xi}g(a,w)\right\}.
\]
All factors belong to $[0,1]$, so this limit exists, is independent of
the order of multiplication, and lies in $[0,1]$. Each finite product is
the modulus of a holomorphic finite Blaschke product, hence is
subharmonic. These subharmonic functions decrease as $R$ increases.
Their bounded nonnegative limit is upper semicontinuous, and passing
to the limit in the submean inequality shows that $p_\xi$ is itself
subharmonic. In particular its distributional Laplacian is nonnegative.

The invariance of pseudohyperbolic distance gives
$p_{\varphi\xi}(\varphi(a))=p_\xi(a)$ for every disk automorphism
$\varphi$. By transitivity and invariance in law,
$\E p_\Xi(a)=c$ is independent of $a$. For every nonnegative
$\chi\in C_c^\infty(\mathbb D)$, Fubini's theorem, justified by
$0\le p_\Xi\le1$, gives
\[
 0\le \E\langle\Delta p_\Xi,\chi\rangle
   =\int_{\mathbb D}c\,\Delta\chi\,dA=0.
\]
Thus $\langle\Delta p_\Xi,\chi\rangle=0$ almost surely.
Choose a countable family of nonnegative smooth compactly supported
functions such that every compact subset of the disk is contained in
the strictly positive set of one of them. Since $\Delta p_\Xi$ is a
positive measure, it follows on one probability-one event that
$\Delta p_\Xi=0$ throughout the disk. A locally integrable function with
zero distributional Laplacian has a harmonic representative. The
subharmonic representative $p_\Xi$ agrees with this harmonic function
everywhere: its averages over shrinking disks converge to its value,
by the submean inequality and upper semicontinuity.

If $w\in\Xi$, then $p_\Xi(w)=0$ because one of the product factors is
zero. The minimum principle for the nonnegative harmonic function
$p_\Xi$ therefore implies $p_\Xi\equiv0$. If $\Xi$ is empty, the
product equals one identically. This proves the dichotomy. Finally,
$g(0,w)=-\log|w|$ is comparable to $1-|w|$ near the boundary; finitely
many interior points do not change convergence of the corresponding
series away from its logarithmic poles. This gives the last assertion.
\end{proof}

\begin{corollary}[The uncompensated invariant Gibbs equation has no solution]
\label{cor:uncompensated-green-dlr-impossible}
Fix $\beta>0$ and an activity $z>0$. For a relatively compact open
$B\subset\mathbb D$ and an exterior configuration $Y\subset B^c$,
consider the unrenormalized grand-canonical kernel
\begin{align}
 \gamma_B(d\eta\mid Y)
 &\propto\sum_{m=0}^\infty\frac{z^m}{m!}
   \exp\left\{-\beta\sum_{i<j}g(a_i,a_j)
       -\beta\sum_{i=1}^m\sum_{w\in Y}g(a_i,w)\right\}
   \prod_{i=1}^m d\lambda(a_i),
 \label{eq:uncompensated-green-specification}
\end{align}
where $(a_1,\ldots,a_m)\in B^m$ and the $m=0$ term is one.
There is no disk-automorphism-invariant locally finite probability law
satisfying the DLR equations for these kernels.
\end{corollary}

\begin{proof}
The kernels are well defined: nonnegativity of $g$ bounds their
normalizing constants between $1$ and $\exp\{z\lambda(B)\}$.
Suppose an invariant law $P$ satisfied the DLR equations. By
Theorem~\ref{thm:invariant-green-divergence}, its nonempty configurations
have infinite Green potential everywhere. Such a configuration is
necessarily infinite. Removing its finitely many points in $B$ leaves
an exterior $Y$ for which
$\sum_{w\in Y}g(a,w)=\infty$ for Lebesgue-almost every $a\in B$:
the removed finite sum is finite outside its finite set of poles.
For such an exterior, every term with $m\ge1$ in
\eqref{eq:uncompensated-green-specification} vanishes. Hence
$\gamma_B(\cdot\mid Y)$ is concentrated on the empty configuration.

If $\Xi\cap B$ is nonempty, then $\Xi$ is nonempty and infinite, so
its exterior is of the preceding type. The DLR equation consequently
gives $P(\Xi\cap B\ne\varnothing)=0$. Applying this to a countable
exhaustion by relatively compact disks yields $P(\Xi=\varnothing)=1$.
But for the empty exterior the $m=1$ term in
\eqref{eq:uncompensated-green-specification} has positive mass
$z\lambda(B)$ whenever $B$ has positive area. Thus the empty law does
not satisfy the same DLR equation, a contradiction.
\end{proof}

\begin{remark}
The corollary rules out only the displayed \emph{uncompensated}
specification. It does not rule out invariant processes, compensated
Gibbs equations, or canonical equations restricted to a fixed number
of particles. In particular, the Bergman process is nonempty and
invariant, and its correct conditional equation uses a regularized
Green field.
\end{remark}

\subsection{The deterministic term in an anchored Green field}

Here is an elementary calculation that locates one normalization issue
in such regularizations. Let
\[
 d\mu(w)=\frac{dA(w)}{\pi(1-|w|^2)^2},
 \qquad I_R(a)=\int_{|w|<R}g(a,w)\,d\mu(w).
\]
For $|a|<R<1$,
\begin{equation}
 I_R(a)=-\frac{R^2\log R}{1-R^2}
         -\frac12\log(1-R^2)
         +\frac12\log(1-|a|^2).
 \label{eq:green-mean-exact}
\end{equation}
Indeed, the circular average of $\log|1-a\overline w|$ is zero,
whereas the circular average of $\log|a-w|$ is
$\log\max\{|a|,|w|\}$. Thus
\[
 I_R(a)=\int_0^R\frac{-2s\log\max\{|a|,s\}}{(1-s^2)^2}\,ds.
\]
Integration at $a=0$ gives the first two terms in
\eqref{eq:green-mean-exact}. Differentiating the difference with respect
to $u=|a|$ gives $-u/(1-u^2)$, and the difference vanishes at $u=0$;
integration yields the last term.

Suppose a point process has intensity $\mu$ and its centered Green
statistics
\[
 H_R(a)=\sum_{w\in\Xi:\,|w|<R}g(a,w)-I_R(a)
\]
converge in $L^2$ to $H(a)$ for each fixed $a$. Then
\begin{equation}
 \lim_{R\uparrow1}
 \sum_{w\in\Xi:\,|w|<R}\bigl(g(a,w)-g(0,w)\bigr)
 =H(a)-H(0)+\frac12\log(1-|a|^2)
 \label{eq:anchored-green-drift}
\end{equation}
in $L^2$. The position-dependent last term must be retained when using
origin-centered cutoffs and anchoring by subtraction at zero.
Full centering by $I_R(a)$ removes this term and, when the centered
limit is independent of the exhaustion, gives an isometry-covariant
field.

For completeness, the required $L^2$ convergence holds for the
Bergman process by an elementary variance estimate. Its determinantal
correlations imply, for real compactly supported $f\ge0$,
\[
 \Var\left(\sum_{w\in\Xi}f(w)\right)
  =\int f(w)^2\,d\mu(w)
    -\iint f(w)f(v)|K(w,v)|^2\,dA(w)dA(v)
  \le\int f(w)^2\,d\mu(w).
\]
For $R<S$, apply this inequality to
$f(w)=g(a,w)\mathbf1_{\{R<|w|<S\}}$. The logarithmic pole is
square integrable on compact sets; near the unit circle,
$g(a,w)=O_a(1-|w|)$. Consequently
$\int g(a,w)^2d\mu(w)<\infty$, and the annular square integral tends
to zero as $R\uparrow1$. This proves the Cauchy property of $H_R(a)$.
The same argument compares arbitrary increasing exhaustions, since
each of the two set differences can be treated separately. Hence the
fully centered Bergman Green field is independent of the exhaustion.

\subsection{A precise compensated specification and its missing input}
\label{subsec:compensated-specification-candidate}

The following conditional construction separates algebraic consistency
from the probabilistic assertion that a limit law satisfies the
resulting equations.

\begin{proposition}[A finite-addition criterion for a DLR specification]
\label{prop:compensated-specification-cocycle}
Let $\Omega$ be a measurable class of locally finite configurations,
stable under finite additions and deletions. Suppose there is a
measurable real-valued function $H(a,Y)$, defined for
$Y\in\Omega$ and $a\notin Y$, with the finite-addition identity
\begin{equation}
 H(a,Y\cup\{b\})=H(a,Y)+g(a,b)
 \qquad(a\ne b,\ a,b\notin Y).
 \label{eq:renormalized-field-addition}
\end{equation}
Assume that for each relatively compact open $B$ and each
$Y\in\Omega$ supported on $B^c$, the function $a\mapsto H(a,Y)$ is
bounded below on $B$. Fix $z>0$ and $\beta>0$. Then
\begin{align}
 \gamma^{H,z}_B(d\eta\mid Y)
 ={}&\frac1{Z_B(Y)}\sum_{m=0}^\infty\frac{z^m}{m!}
  \exp\left\{-\beta\sum_{i<j}g(a_i,a_j)
       -\beta\sum_{i=1}^mH(a_i,Y)\right\}
  \prod_{i=1}^m d\lambda(a_i)
 \label{eq:renormalized-specification-candidate}
\end{align}
defines finite, normalized, consistent kernels on $B$; diagonals are
irrelevant to the Lebesgue integrals. If $\Omega$ is invariant under
disk automorphisms and
\[
 H(\varphi(a),\varphi(Y))=H(a,Y),
\]
the kernels are covariant under these automorphisms as well.
\end{proposition}

\begin{proof}
If $h_B(Y)=\inf_{a\in B}H(a,Y)>-\infty$, nonnegativity of $g$ gives
\[
 1\le Z_B(Y)\le
   \exp\left\{z e^{-\beta h_B(Y)}\lambda(B)\right\}<\infty.
\]
For $B\subset C$, divide a configuration in $C$ into
$\eta\subset B$ and $\zeta\subset C\setminus B$. The energy in
\eqref{eq:renormalized-specification-candidate} splits into the energy
of $\zeta$ in the exterior $Y$, plus the energy of $\eta$ in the
exterior $Y\cup\zeta$, because
\[
 \sum_{a\in\eta}H(a,Y\cup\zeta)
 =\sum_{a\in\eta}H(a,Y)
   +\sum_{a\in\eta}\sum_{b\in\zeta}g(a,b).
\]
The factorial reference measure splits over the two disjoint sets.
Consequently the conditional law of $\eta$, given $\zeta$, under the
$C$-kernel is exactly the $B$-kernel, proving consistency. Covariance
follows by change of variables from invariance of $g$ and $\lambda$
and the assumed covariance of $H$.
\end{proof}

A candidate field associated with a constant hyperbolic background
$c\lambda$ is
\begin{equation}
 H(a,Y)=\lim_{n\to\infty}\left\{
   \sum_{w\in Y\cap E_n}g(a,w)
       -c\int_{E_n}g(a,w)\,d\lambda(w)\right\},
 \label{eq:general-centered-green-candidate}
\end{equation}
where $E_n$ exhausts the disk. Whenever this defines a measurable
field with the stated lower bounds and with a limit independent of
the exhaustion, finite addition proves
\eqref{eq:renormalized-field-addition} and change of variables proves
isometry covariance. Thus the preceding proposition applies. These
are hypotheses in the present discussion for general $\beta$, not
established properties of all gas limits.

There is also a useful necessary identity that retains the particle
count information. If a law $P$ supported on $\Omega$ satisfies these DLR equations, then
for every nonnegative measurable $F$ supported in a relatively compact
set in its first argument,
\begin{equation}
 \E_P\sum_{a\in\Xi}F(a,\Xi\setminus\{a\})
 =z\int_{\mathbb D}\E_P\left[
      e^{-\beta H(a,\Xi)}F(a,\Xi)\right]d\lambda(a).
 \label{eq:compensated-gnz-necessary}
\end{equation}
To prove it, condition on the exterior of a bounded set containing
the support of the first argument, substitute
\eqref{eq:renormalized-specification-candidate}, and distinguish one
particle in the term with $m$ particles. The factor $m/m!$ becomes
$1/(m-1)!$, while \eqref{eq:renormalized-field-addition} gives precisely
the exponential on the right. All integrands are nonnegative, so
Tonelli's theorem justifies every rearrangement.

The activity $z$ cancels from the conditional density when the number
of particles in $B$ is fixed. Consequently, convergence of density
ratios for configurations with the same number of particles cannot,
by itself, identify \eqref{eq:compensated-gnz-necessary} or the full
DLR equation. A proof for the original gas must in addition identify
the relative weights of the different particle counts.

\subsection{A conditional compatibility test for hyperbolic invariance}
\label{subsec:conditional-hyperbolic-intensity}

This subsection states an implication, not a new existence or
invariance theorem for the gas. It records a normalization test that
any proposed canonical DLR derivation should pass.

\begin{proposition}[Covariance forbids an uncompensated one-body factor]
\label{prop:invariance-forces-zero-one-body-exponent}
Suppose $P$ is a nonempty, disk-automorphism-invariant point-process
law. Suppose its conditional law in each relatively compact open set
$B$, given the exterior $Y$ and the inside count $m$, has a strictly
positive density away from collisions of the form
\begin{equation}
 \frac1{Z_{B,Y,m}}
 \prod_{i<j}\delta(a_i,a_j)^\beta
 \prod_{i=1}^m
       e^{-\beta H(a_i,Y)}(1-|a_i|^2)^q\,d\lambda(a_i),
 \label{eq:canonical-covariance-test}
\end{equation}
where $q$ is deterministic and $H$ is isometry covariant. Then $q=0$.
\end{proposition}

\begin{proof}
Some relatively compact open disk $B$ has
$P(\Xi(B)\ge1)>0$. Fix a nonrotation disk automorphism $\varphi$.
Invariance of $P$ identifies the push-forward of its conditional
law on $B$, given $(Y,m)$, with its conditional law on $\varphi(B)$,
given $(\varphi(Y),m)$, almost surely. The Green interaction, $H$,
and $\lambda$ are covariant, so the ratio between the two displayed
densities, pulled back to $B^m$, is proportional to
\[
 \prod_{i=1}^m\left(
    \frac{1-|\varphi(a_i)|^2}{1-|a_i|^2}\right)^q
 =\prod_{i=1}^m|\varphi'(a_i)|^q.
\]
The equality of the normalized conditional laws requires this product
to be constant almost everywhere on $B^m$, for a conditional event
with some $m\ge1$. Their densities are strictly positive there away
from collisions, so continuity makes the product constant everywhere
off the collision diagonals. Fixing all coordinates but one shows
that $|\varphi'|^q$ is constant on an open set. For a nonrotation
automorphism $|\varphi'(a)|=(1-|u|^2)/|1-\overline u a|^2$ with
$u\ne0$, which is nonconstant on every open set. Hence $q=0$.
\end{proof}

To see why this test is relevant to the original gas, suppose, in
addition to assumptions not yet proved for general $\beta$, that its
canonical limiting density were established in the anchored form
\begin{equation}
 \frac1{\widetilde Z_{B,Y,m}}
 \prod_{i<j}\delta(a_i,a_j)^\beta
 \prod_{i=1}^m(1-|a_i|^2)^{2-\beta/2}
                  e^{-\beta C_Y(a_i)}\,d\lambda(a_i).
 \label{eq:anchored-canonical-candidate}
\end{equation}
Here $B$ contains the origin, and $C_Y$ denotes the limit of the
origin-anchored exterior sums
$\sum_{w\in Y,\,|w|<R}(g(a,w)-g(0,w))$. Suppose the invariant law has
intensity $c\lambda$, with $0<c<\infty$, and suppose its fully centered
field $H$ exists independently of the exhaustion. Formula
\eqref{eq:green-mean-exact}, with the intensity multiplied by $c\pi$,
then gives
\[
 C_Y(a)=H(a,Y)-H(0,Y)+\frac{c\pi}{2}\log(1-|a|^2).
\]
The factor involving $H(0,Y)$ is constant when $m$ is fixed. Thus
\eqref{eq:anchored-canonical-candidate} becomes
\eqref{eq:canonical-covariance-test} with
\[
 q=2-\frac\beta2-\frac{\beta c\pi}{2}.
\]
If this same canonical law is established on all relatively compact
sets, covariance consequently forces
\begin{equation}
 c=\frac{4/\beta-1}{\pi}.
 \label{eq:conditional-hyperbolic-intensity-prediction}
\end{equation}
At $\beta=2$ this agrees with the Bergman intensity. At $\beta=4$
it vanishes. Equation~\eqref{eq:conditional-hyperbolic-intensity-prediction}
is only a compatibility consequence of the explicitly stated
canonical-DLR and field hypotheses. Neither those hypotheses nor
hyperbolic invariance at other $0<\beta<4$ have been proved here.

\section{A complete compensated DLR equation at \texorpdfstring{$\beta=2$}{beta=2}}
\label{sec:bergman-complete-dlr}

At the exactly solvable temperature the conditional law can be identified
including its conditional particle number.  We give a derivation, because
it also identifies the interaction that a proposed limiting DLR equation
must reproduce.  Throughout this section $\mathsf B$ denotes the Bergman
process, and
\[
 d\nu(z)=\frac{dA(z)}{\pi(1-|z|^2)^2},\qquad
 \delta(z,w)=\left|\frac{z-w}{1-z\bar w}\right|,
 \qquad g(z,w)=-\log\delta(z,w).
\]
Thus $\nu$ is precisely the first correlation measure of $\mathsf B$.
Our normalization of hyperbolic area includes the factor $1/\pi$.

\subsection{An intrinsic compensated potential}

\begin{lemma}[Centered potentials and exponential integrability]
\label{lem:bergman-centered-potential}
For every $q\in\D$ there is a real random variable $H_q$ such that, for
any increasing exhaustion $E_j\uparrow\D$ by relatively compact Borel sets,
\begin{equation}
 H_q(X)=\lim_{j\to\infty}\left\{
   \sum_{x\in X\cap E_j}g(q,x)-\int_{E_j}g(q,z)\,d\nu(z)
 \right\}
 \quad\hbox{in }L^2(\mathsf B).
 \label{eq:bergman-centered-potential}
\end{equation}
The result is independent of the exhaustion.  There is a jointly
measurable version for Lebesgue-almost every $q$ and $\mathsf B$-almost
every $X$.  Put $Y_q=-2H_q$.  For every $p>0$, the exponentials of
the corresponding truncated $pY_q$ converge in $L^1$.  In particular
$\exp(Y_q)$ is strictly positive and integrable.
\end{lemma}

\begin{proof}
For a real compactly supported test function $f$ the determinantal
variance identity gives
\[
 \operatorname{Var}\sum_{x\in X}f(x)
 =\int f^2\,d\nu-
   \iint f(z)f(w)|K(z,w)|^2\,dA(z)dA(w)
 \leq\int f^2\,d\nu.
\]
The second term is nonnegative: for bounded tests it is the squared
Hilbert--Schmidt norm of the self-adjoint operator $KfK$, since $K$ is
an orthogonal projection.  Truncation extends the bound to the required
$L^2(\nu)$ tests.  Hyperbolic invariance and polar integration give
\[
 \int_\D |\log\delta(q,z)^2|^2\,d\nu(z)
 =\int_0^1\frac{(\log s)^2}{(1-s)^2}\,ds
 =2\sum_{n=1}^{\infty}\frac1{n^2}<\infty.
\]
The variance bound therefore proves the Cauchy property and independence
of the exhaustion.  Applying the same bound after integrating $q$ over
a compact set yields jointly measurable $L^2$ limits, consistently on a
countable exhaustion of the $q$ variable.

For $f\leq0$ of compact support the determinantal generating functional
and $\log\det(I+A)\leq\operatorname{Tr}A$ give
\begin{equation}
 \E\exp\left(p\left\{\sum_{x\in X}f(x)-\int f\,d\nu\right\}\right)
 \leq \exp\left(\int(e^{pf}-1-pf)\,d\nu\right).
 \label{eq:bergman-centered-exp-bound}
\end{equation}
One may first replace $f$ by $\max(f,-M)$ and then pass to the limit;
the sums are finite on its compact support and the deterministic
integrals converge.  Since $0\leq e^{pu}-1-pu\leq p^2u^2/2$ for
$u\leq0$, the right-hand side is bounded uniformly over truncations
of $f(z)=\log\delta(q,z)^2$.  Applying the bound with any larger
exponent proves uniform integrability.  Convergence in probability,
already supplied by $L^2$ convergence of the centered statistics,
then proves the asserted $L^1$ convergence of their exponentials.
\end{proof}

In particular the compensation in \eqref{eq:bergman-centered-potential}
is intrinsic.  A different exhaustion does not change the centered
random variable.  Removing finitely many points is a different operation:
the compensating integral is retained, so the potential changes by the
ordinary finite sum of their pair potentials.

\begin{proposition}[Normalized Palm likelihood]
\label{prop:bergman-palm-likelihood}
Let $\gamma=\lim_{n\to\infty}(\sum_{k=1}^n k^{-1}-\log n)$.
The reduced Palm law $\mathsf B^{!q}$ at $q$ is equivalent to $\mathsf B$,
and its Radon--Nikodym derivative is
\begin{equation}
 \Psi_q(X):=\frac{d\mathsf B^{!q}}{d\mathsf B}(X)
 =e^{-\gamma}\exp(-2H_q(X)).
 \label{eq:bergman-palm-likelihood}
\end{equation}
For finite additions or removals, wherever these expressions are defined,
\begin{equation}
 \Psi_q(X\cup\{p_1,\ldots,p_m\})
 =\Psi_q(X)\prod_{j=1}^m\delta(q,p_j)^2.
 \label{eq:bergman-palm-multiplicativity}
\end{equation}
Moreover, under $\mathsf B$, the single random variable $\Psi_q$ has the
exponential distribution of mean one.
\end{proposition}

\begin{proof}
It suffices to start with $q=0$.  Write $t=r^2$ and let
$f_t(z)=\ind_{\{|z|^2<t\}}\log|z|^2$.  Direct integration gives
\begin{equation}
 I(t):=\int f_t\,d\nu
 =\frac{t\log t}{1-t}+\log(1-t).
 \label{eq:bergman-log-centering}
\end{equation}
Set $h_t=e^{f_t}$.  In the orthonormal Bergman basis
$e_k(z)=\sqrt{(k+1)/\pi}\,z^k$, the operator $K(1-h_t)K$ is diagonal,
with eigenvalues
\[
 a_k(t)=(k+1)\int_0^t(1-s)s^k\,ds,
 \qquad 0\leq a_k(t)\leq\frac1{k+2}.
\]
Its trace is $-\log(1-t)$.  Consequently
\begin{align*}
 \log\E\exp\left(\sum_{x\in X}f_t(x)-I(t)\right)
 &= -\frac{t\log t}{1-t}
    +\sum_{k\geq0}\{\log(1-a_k(t))+a_k(t)\}.
\end{align*}
The summands are dominated in absolute value by a constant times
$(k+2)^{-2}$.  Their limiting sum is
\[
 \lim_{M\to\infty}
 \left\{-\log(M+2)+\sum_{k=0}^{M}\frac1{k+2}\right\}
 =\gamma-1.
\]
Since $-t\log t/(1-t)\to1$, Lemma~\ref{lem:bergman-centered-potential}
proves $\E e^{Y_0}=e^{\gamma}$.

We also identify the tilted law directly.  Weighting the Bergman process
by $\prod_{x\in X}h_t(x)$ and normalizing gives a determinantal process
with kernel
\begin{equation}
 K_t(z,w)=\sqrt{h_t(z)h_t(w)}
   \sum_{k\geq0}\frac{k+1}{\pi(1-a_k(t))}z^k\bar w^k.
 \label{eq:bergman-radial-tilted-kernel}
\end{equation}
For completeness, the finite-rank version follows by multiplying each
row and column of the joint determinant by $\sqrt{h_t}$ and normalizing
the resulting orthogonal functions; radiality makes their Gram matrix
diagonal.  Passing through the increasing Bergman projections gives
\eqref{eq:bergman-radial-tilted-kernel}.  The passage is justified by
convergence of their Fredholm generating functionals on compact sets:
$1-h_t$ has compact support, while $1-a_k(t)\geq1/2$ bounds all inverse
Gram coefficients.  The same bound gives uniform convergence of the
kernel series on compact sets and convergence of its diagonal integrals.

As $t\uparrow1$, the kernel in \eqref{eq:bergman-radial-tilted-kernel}
converges locally, including its diagonal integrals, to
\[
 \frac{|z||w|}{\pi}\sum_{k\geq0}(k+2)z^k\bar w^k.
\]
After the harmless gauge multiplication by
$(z/|z|)(\bar w/|w|)$, this is
\[
 \frac{z\bar w}{\pi}\sum_{k\geq0}(k+2)z^k\bar w^k
 =K(z,w)-\frac1\pi.
\]
The latter is the reduced Palm kernel at zero, as is seen by applying
the Schur-complement identity to each correlation determinant with one
argument fixed at zero.  On the other hand, Lemma~\ref{lem:bergman-centered-potential}
shows that the normalized densities of these tilted laws converge in
$L^1(\mathsf B)$ to $e^{-\gamma}e^{Y_0}$.  This proves
\eqref{eq:bergman-palm-likelihood} at zero.  The Bergman kernel identity
\[
 K(\varphi z,\varphi w)\varphi'(z)\overline{\varphi'(w)}=K(z,w)
\]
for disk automorphisms, together with the exhaustion independence of
the centered statistic, proves the assertion at every $q$.
Positivity of the exponential proves equivalence.  Finite
multiplicativity follows already at the truncated-product level.

Here is also a check of the complete one-point distribution of the
likelihood.  For $p>0$, radial diagonalization gives
\begin{equation}
 \E e^{pY_0}
 =\prod_{n=1}^{\infty}\frac{e^{p/n}}{1+p/n}
 = e^{\gamma p}\Gamma(1+p).
 \label{eq:bergman-likelihood-moments}
\end{equation}
To justify the first limit without a formal infinite-product
interchange, set $b_{n,t}=\int_0^1(e^{pf_t(\sqrt s)}-1)n s^{n-1}\,ds$.
Then $|b_{n,t}|\leq p/(n+p)$ and
\[
 0\leq b_{n,t}-p\int_0^1f_t(\sqrt s)n s^{n-1}\,ds
 \leq\frac{p^2}{2}\int_0^1(\log s)^2n s^{n-1}\,ds
 =\frac{p^2}{n^2}.
\]
The identity $\log(1+b)=b+O(b^2)$ bounds the remaining tail uniformly
by a summable multiple of $n^{-2}$; the finitely many small $n$ cause
no difficulty.  Dominated convergence therefore gives the first
product in \eqref{eq:bergman-likelihood-moments}.  The second equality
is Euler's product for the gamma function, equivalently obtained by
applying the beta integral and letting its integer parameter increase.
It follows that $\E\Psi_q^m=m!$ for every integer $m\geq0$.  These
moments determine a probability law: for $0<s<1$, monotone convergence
gives $\E e^{s\Psi_q}=\sum_{m\geq0}s^m=(1-s)^{-1}$.
They therefore identify the exponential law of mean one.
\end{proof}

\subsection{The full conditional law, including the number of points}

Let $B\Subset\D$ be a nonempty open set and $C=\D\setminus B$.
For an exterior configuration $Y=X\cap C$, define $H_q(Y)$ by the same
centered expression as in \eqref{eq:bergman-centered-potential}, with
the sum over $Y$ and the integral still over the whole truncating set.
Equivalently, if $X=Y\cup\xi$ is a typical full configuration,
\[
 H_q(Y)=H_q(X)-\sum_{p\in\xi}g(q,p),\qquad q\in B.
\]
This definition is measurable and independent of the chosen completion
for exterior-almost every $Y$ and $\nu$-almost every $q$.  Indeed, the
truncated exterior expressions depend only on $(q,Y)$.  Apply the joint
$L^2$ construction in Lemma~\ref{lem:bergman-centered-potential}
along a subsequence and the exact finite-deletion identity.  Their
almost sure limit therefore also depends only on $(q,Y)$; disintegration
gives the stated almost-everywhere assertion.  Set
\[
 W_Y(q)=e^{-\gamma}\exp(-2H_q(Y)).
\]

\begin{theorem}[Compensated grand-canonical Bergman specification]
\label{thm:bergman-grandcanonical-dlr}
For $\mathsf B|_C$-almost every $Y$, the conditional law of $X\cap B$
is
\begin{equation}
 \Gamma_B(d\xi\mid Y)
 =\frac1{Z_B(Y)}\sum_{m=0}^{\infty}\frac1{m!}
   \left[\prod_{i=1}^m W_Y(q_i)
   \prod_{1\leq i<j\leq m}\delta(q_i,q_j)^2\right]
   \prod_{i=1}^m d\nu(q_i),
 \label{eq:bergman-grandcanonical-dlr}
\end{equation}
where the $m$-th term is pushed forward by
$(q_1,\ldots,q_m)\mapsto\{q_1,\ldots,q_m\}$, the $m=0$ term is a
unit mass at the empty configuration, and $1\leq Z_B(Y)<\infty$.
Every integer $m\geq0$ has strictly positive conditional probability.
The conditional law is therefore equivalent to a Poisson process with
intensity $\nu|_B$.  In particular the Bergman process is not number rigid.

Conditionally also on $\#(X\cap B)=m$, its symmetric ordered density
with respect to $\nu^{\otimes m}$ is proportional to
\begin{equation}
 \exp\left[-2\sum_{i<j}g(q_i,q_j)
            -2\sum_{i=1}^m H_{q_i}(Y)\right].
 \label{eq:bergman-canonical-dlr}
\end{equation}
Thus the pair potential in this conditional law is the hyperbolic Green
potential.  Both its divergent exterior interaction and its conditional
particle-number normalization are explicitly accounted for.
\end{theorem}

\begin{proof}
The reduced Campbell identity and Proposition~\ref{prop:bergman-palm-likelihood}
give the following Georgii--Nguyen--Zessin identity, first for
nonnegative measurable tests $F$ and then for integrable ones:
\begin{equation}
 \E\sum_{q\in X}F(q,X\setminus\{q\})
 =\int_\D\E[F(q,X)\Psi_q(X)]\,d\nu(q).
 \label{eq:bergman-gnz}
\end{equation}
Indeed both sides equal the integral of $F$ against the reduced Campbell
measure, after disintegration at its distinguished point.
Restrict $q$ to $B$, multiply the test by any bounded measurable function
of $Y=X\cap C$, and disintegrate over $Y$.  For almost every $Y$ the
conditional finite process $Q_Y$ on $B$ then has insertion rate, relative
to $\nu$, given by
\[
 c_Y(q,\xi)=W_Y(q)\prod_{p\in\xi}\delta(q,p)^2.
\]
This follows from \eqref{eq:bergman-palm-multiplicativity}.  The passage
can be performed first for a countable class of cylinder tests that
generates the relevant sigma fields, and then for all nonnegative tests
by a monotone-class argument.

We spell out why this identity determines also the probabilities of
the different counts.  Let $J_{Y,m}$ be the symmetric Janossy measure:
its total mass is $m!Q_Y(\#B=m)$.  The conditional Campbell identity
implies the recursion, in the sense of measures,
\begin{equation}
 J_{Y,m+1}(dq,dq_1\cdots dq_m)
 =c_Y(q,\{q_1,\ldots,q_m\})\,d\nu(q)
   J_{Y,m}(dq_1\cdots dq_m).
 \label{eq:bergman-janossy-recursion}
\end{equation}
For example, this follows by using a test that vanishes unless the
remaining configuration has exactly $m$ points and symmetrizing its
arguments; the factors $m!$ cancel.  It also shows inductively that
all the Janossy measures have densities once $J_{Y,0}$ is specified.
If $J_{Y,0}=0$, the recursion makes every $J_{Y,m}$ zero, contradicting
the fact that $Q_Y$ is supported on finite configurations.  Hence
$J_{Y,0}>0$.  Iteration gives
\[
 J_{Y,m}(dq_1\cdots dq_m)
 =J_{Y,0}\prod_i W_Y(q_i)\prod_{i<j}\delta(q_i,q_j)^2
   \prod_i d\nu(q_i).
\]
Summing its mass divided by $m!$ gives
$Z_B(Y)=J_{Y,0}^{-1}<\infty$ and proves
\eqref{eq:bergman-grandcanonical-dlr}.  Every integrand is positive
off the diagonals, for $\nu^{\otimes m}$-almost every tuple.  As
$0<\nu(B)<\infty$, every term has strictly positive integral.
The same formula proves equivalence to Poisson measure and the
canonical conditional law.  Nonrigidity follows because an
exterior-measurable count would have a conditionally deterministic
distribution, whereas here every possible finite count has positive
conditional probability.
\end{proof}

\begin{corollary}[Correctly normalized $L$-ensemble]
\label{cor:bergman-conditional-lensemble}
The positive operator on $L^2(B,\nu)$ with kernel
\begin{equation}
 L_Y(z,w)=
 \frac{\sqrt{(1-|z|^2)W_Y(z)}\,
       \sqrt{(1-|w|^2)W_Y(w)}}{1-z\bar w}
 \label{eq:bergman-correct-l-kernel}
\end{equation}
is trace class and satisfies
\[
 Z_B(Y)=\det(I+L_Y),\qquad
 Q_Y(\#B=0)=\det(I+L_Y)^{-1}.
\]
The conditional process is determinantal, with kernel
$L_Y(I+L_Y)^{-1}$.  In particular $L_Y(z,z)=W_Y(z)$; this diagonal
identity fixes the normalization unambiguously.
\end{corollary}

\begin{proof}
The Cauchy determinant identity gives
\[
 \det(L_Y(q_i,q_j))_{i,j=1}^m
 =\prod_iW_Y(q_i)\prod_{i<j}\delta(q_i,q_j)^2.
\]
Positivity follows also from the power-series expansion of
$(1-z\bar w)^{-1}$.  The $m=1$ coefficient of the already finite
partition function shows $\int_B W_Y\,d\nu<\infty$.
Writing $L_Y$ as the sum of the rank-one kernels generated by
$\sqrt{(1-|z|^2)W_Y(z)}z^k$, the sum of their traces is exactly
$\int_BW_Y\,d\nu$.  Thus the sum converges in trace norm and
defines a positive trace-class operator.  Its Fredholm expansion is
\eqref{eq:bergman-grandcanonical-dlr}; the generating functional then
identifies the correlation kernel $L_Y(I+L_Y)^{-1}$.
\end{proof}

\begin{corollary}[A curvature obstruction to a naive conditional limit]
\label{cor:bergman-conditional-curvature}
For exterior-almost every $Y$, the compensated potential has a smooth
representative on $B$, and
\[
 h_Y(q):=H_q(Y)+\tfrac12\log(1-|q|^2)
\]
is harmonic there.  The one-particle conditional density $p_Y$ with
respect to ordinary area, after additionally conditioning on $\#B=1$,
therefore satisfies
\begin{equation}
 p_Y(q)=\frac{1}{\mathcal Z_Y}\,
       \frac{e^{-2h_Y(q)}}{1-|q|^2},
 \qquad
 \Delta\log p_Y(q)=\frac4{(1-|q|^2)^2}.
 \label{eq:bergman-onepoint-conditional-curvature}
\end{equation}
If $B$ is connected, this conditional probability cannot be obtained
as a weak limit on $B$ of the finite-$N$, $\beta=2$ conditional
one-particle laws with any sequence of deterministic exterior
configurations in $\D\setminus B$.
\end{corollary}

\begin{proof}
Use centered disks for the truncation.  The distributional identity
$\Delta_qg(q,z)=-2\pi\delta_z(q)$ gives, for compactly supported tests
in $B$ and all truncating disks containing $\overline B$,
\[
 \Delta_q\left\{\sum_{y\in Y\cap E_j}g(q,y)
                   -\int_{E_j}g(q,z)\,d\nu(z)\right\}
 =\frac2{(1-|q|^2)^2}.
\]
Use the joint $L^2$ construction and finite deletion to pass to the
limit against the test functions.  Explicitly, on each compact
$q$-window the deleted sum has finite joint second moment: expand its
square, use the compact bounds on the first two correlation functions,
and the local integrability of $\log^2|q-p|$.  Thus finite deletion
preserves the required joint $L^2$ control.  Since
$\Delta\log(1-|q|^2)=-4/(1-|q|^2)^2$, the displayed $h_Y$ is
distributionally harmonic.  Weyl's lemma gives its smooth harmonic
representative.  Substitution into the $m=1$ term of
\eqref{eq:bergman-grandcanonical-dlr}, including the factor
$d\nu/dA$, proves \eqref{eq:bergman-onepoint-conditional-curvature}.

For a finite exterior configuration $\{y_1,\ldots,y_n\}$, the
conditional one-particle density is $|F_n(q)|^2$, where
\[
 F_n(q)=\frac{\prod_{j=1}^{n}(q-y_j)}
 {\left(\int_B\prod_{j=1}^{n}|u-y_j|^2\,dA(u)\right)^{1/2}}.
\]
Each $F_n$ is holomorphic and has no zero in $B$, and
$\int_B|F_n|^2\,dA=1$.  The submean inequality makes this family
locally bounded.  Montel's theorem supplies a locally uniformly
convergent subsequence with holomorphic limit $F$.  If the probability
measures $|F_n|^2dA$ converge weakly to a probability measure on $B$,
testing against compactly supported continuous functions identifies
that measure as $|F|^2dA$, so $F$ is not identically zero.  Hurwitz's
theorem implies that $F$ has no zero in connected $B$.  Its limiting
log density is consequently harmonic, in contradiction with
\eqref{eq:bergman-onepoint-conditional-curvature}.
\end{proof}

This obstruction does not contradict convergence of the point
processes.  It demonstrates that conditioning on the full exterior
and taking the infinite-particle limit do not commute in the required
way; the correct limiting conditional law retains an effect of
boundary randomness.

\subsection{Isometry covariance and the scope of this specification}

\begin{proposition}[Covariance under every hyperbolic isometry]
\label{prop:bergman-dlr-isometries}
The Bergman process is invariant under every isometry of the hyperbolic
disk, including orientation-reversing ones.  For each such isometry
$\varphi$, compatible almost-everywhere versions satisfy
\[
 H_{\varphi q}(\varphi X)=H_q(X),\qquad
 \Psi_{\varphi q}(\varphi X)=\Psi_q(X),
\]
and the conditional kernels satisfy
\[
 \Gamma_{\varphi B}(\varphi A\mid\varphi Y)
 =\Gamma_B(A\mid Y)
\]
for exterior-almost every $Y$ and every measurable event $A$ of
configurations in $B$.
\end{proposition}

\begin{proof}
For orientation-preserving isometries, invariance follows from the
Bergman kernel transformation identity used above, after changing the
area variables in every correlation determinant.  Reflection follows
from $K(\bar z,\bar w)=\overline{K(z,w)}$: the determinants are real.
Together these generate all hyperbolic isometries.  Both $g$ and $\nu$
are invariant.  Transform the exhaustion in
\eqref{eq:bergman-centered-potential}; exhaustion independence gives
the equality of centered potentials in $L^2$ and hence of their
compatible almost-everywhere versions.  Finite deletion gives the
same relation for exterior configurations.  Changing variables in
each term of \eqref{eq:bergman-grandcanonical-dlr} proves the claimed
covariance, including equality of its partition functions.
\end{proof}

The formulas above identify the conditional laws of one specific,
already constructed process.  They do not prove uniqueness among all
infinite-volume solutions of a compensated Green specification, and
they do not identify a specification for the nonsolvable disk gases.
In particular, finite-volume Euclidean pair factors and their
conditional exterior fields cannot simply be passed term by term to
the limit: already at $\beta=2$, the factor
\[
 \prod_{i<j}\delta(q_i,q_j)^2
 =|\Delta_m(q)|^2\prod_{i<j}|1-q_i\bar q_j|^{-2}
\]
contains a two-body image term absent from the finite canonical
conditional density with a fixed exterior configuration.

\begin{remark}[Relation to the literature]
The conditional Bergman law was explicitly studied by A.~I.~Bufetov
\cite{BufetovBergman}. The preceding self-contained derivation fixes
the regularization, reference measure, factorial convention, and
$L$-kernel normalization used here. The Palm likelihood has expectation
one, and the $L$-kernel has diagonal $W_Y$.
\end{remark}

\section{Interpolation lower bounds and nonempty limits for \texorpdfstring{$0<\beta\le1$}{0<beta <= 1}}
\label{sec:subadditive-interpolation}

This section gives a lower bound, rather than a compactness or
conditional-law argument. Its only temperature restriction comes from
the elementary inequality $|u+v|^\beta\le |u|^\beta+|v|^\beta$ for
$0<\beta\le1$. In particular, it proves that an empty subsequential
limit is impossible throughout this interval.

\begin{theorem}[An interpolation lower bound]
\label{thm:subadditive-interpolation}
Fix $0<\beta\le1$. If $f$ is a nonzero polynomial of degree strictly
less than $N$, then, for every $a\in\D$,
\begin{equation}
 \rho_{1,N}(a)\ge
 \frac{|f(a)|^\beta}{\displaystyle\int_\D |f(z)|^\beta\,dA(z)}.
 \label{eq:subadditive-interpolation-one}
\end{equation}
More generally, let $1\le p\le N$, and let $F$ be a nonzero
alternating polynomial in $p$ variables, of degree strictly less than
$N$ in each individual variable. Then
\begin{equation}
 \rho_{p,N}(\mathbf a)\ge
 \frac{p!\,|F(\mathbf a)|^\beta}
 {\displaystyle\int_{\D^p}|F(\mathbf z)|^\beta\,dA^p(\mathbf z)}
 \qquad(\mathbf a\in\D^p).
 \label{eq:subadditive-interpolation-many}
\end{equation}
In particular,
\begin{equation}
 \rho_{p,N}(\mathbf a)\ge
 \frac{p!}{Z_{p,\beta}}|\Delta_p(\mathbf a)|^\beta
 \qquad(N\ge p).
 \label{eq:subadditive-interpolation-basic}
\end{equation}
\end{theorem}

\begin{proof}
For a configuration $\mathbf z=(z_1,\ldots,z_N)$ of distinct points,
write
\[
 \ell_i(a;\mathbf z)=
 \prod_{j\ne i}\frac{a-z_j}{z_i-z_j}.
\]
Polynomial interpolation gives
$f(a)=\sum_{i=1}^N\ell_i(a;\mathbf z)f(z_i)$ whenever
$\deg f<N$. Subadditivity therefore gives the pointwise inequality
\[
 |f(a)|^\beta\le
 \sum_{i=1}^N |\ell_i(a;\mathbf z)|^\beta|f(z_i)|^\beta.
\]
All summands are nonnegative. Multiplication by the gas density and
Tonelli's theorem are therefore legitimate without a preliminary
integrability assumption. For each $i$, cancellation of the factors
involving $z_i$ in the Vandermonde shows that
\begin{align*}
 &\mathbb E_{N,\beta}\sum_{i=1}^N
       |\ell_i(a;\mathbf z)|^\beta|f(z_i)|^\beta\\
 &\quad=\frac{N}{Z_{N,\beta}}
 \left(\int_\D|f(z)|^\beta\,dA(z)\right)
 \int_{\D^{N-1}}|\Delta_{N-1}(\mathbf y)|^\beta
                  \prod_{j=1}^{N-1}|a-y_j|^\beta\,dA^{N-1}(\mathbf y)\\
 &\quad=\rho_{1,N}(a)\int_\D|f(z)|^\beta\,dA(z).
\end{align*}
This proves~\eqref{eq:subadditive-interpolation-one}. The collision
set discarded at the start has probability zero.

For completeness, the exact multilinear identity behind the
$p$-point version is
\begin{equation}
 F(\mathbf a)=
 \sum_{\substack{I\subset\{1,\ldots,N\}\\|I|=p}}
 \det[\ell_{i_k}(a_j;\mathbf z)]_{j,k=1}^p\,
 F(z_{i_1},\ldots,z_{i_p}),
 \label{eq:subadditive-alternating-interpolation}
\end{equation}
where $i_1<\cdots<i_p$ orders $I$. To prove it, apply the
one-variable interpolation formula successively in all $p$ variables.
Terms with repeated indices vanish because $F$ is alternating.
Grouping the remaining ordered indices by their underlying set gives
the determinant in~\eqref{eq:subadditive-alternating-interpolation}.

The Cauchy determinant identity, or direct cancellation in the
Lagrange formulas, gives in absolute value
\begin{equation}
 \left|\det[\ell_{i_k}(a_j;\mathbf z)]_{j,k=1}^p\right|
 =\frac{|\Delta_p(\mathbf a)|}{|\Delta_p(\mathbf z_I)|}
 \frac{\displaystyle\prod_{j=1}^p\prod_{k\notin I}|a_j-z_k|}
 {\displaystyle\prod_{i\in I}\prod_{k\notin I}|z_i-z_k|}.
 \label{eq:subadditive-cauchy-cancellation}
\end{equation}
Initially the formula can be read away from all zero denominators;
the determinant formula extends it everywhere else by continuity.
Subadditivity applied to~\eqref{eq:subadditive-alternating-interpolation},
followed by integration against $|\Delta_N|^\beta/Z_{N,\beta}$,
cancels both the cross factors and the Vandermonde on $\mathbf z_I$.
The variables indexed by $I$ consequently contribute exactly
$\int_{\D^p}|F|^\beta\,dA^p$. The remaining variables contribute
the integral in the definition of $h_{p,N}$. As there are
$\binom Np$ sets $I$, the resulting identity is
\begin{align*}
 &\mathbb E_{N,\beta}
 \sum_{|I|=p}
 \left|\det[\ell_{i_k}(a_j;\mathbf z)]\right|^\beta
 |F(\mathbf z_I)|^\beta\\
 &\qquad=\frac{\rho_{p,N}(\mathbf a)}{p!}
                         \int_{\D^p}|F(\mathbf z)|^\beta\,dA^p(\mathbf z).
\end{align*}
This proves~\eqref{eq:subadditive-interpolation-many} for distinct
$a_j$, and continuity proves it on the diagonals as well.
Choosing $F=\Delta_p$, whose degree in each variable is $p-1$, proves
\eqref{eq:subadditive-interpolation-basic}.
\end{proof}

\begin{corollary}[Every limiting intensity dominates the Bergman intensity]
\label{cor:subadditive-bergman-lower}
Let $0<\beta\le1$, and let $\Xi$ be any subsequential limit of the
disk gas. Its one-point density satisfies
\begin{equation}
 \rho_1(a)\ge\frac1{\pi(1-|a|^2)^2},\qquad a\in\D.
 \label{eq:subadditive-bergman-lower}
\end{equation}
In particular, $\Xi$ is not the almost surely empty process, and
$\mathbb E\Xi(\D)=+\infty$. Its factorial densities satisfy
\begin{equation}
 \rho_p(\mathbf a)\ge
       \frac{p!}{Z_{p,\beta}}|\Delta_p(\mathbf a)|^\beta
       \qquad(p\ge1).
 \label{eq:subadditive-limit-many}
\end{equation}
Every nonempty open $U\subset\D$ consequently has
$\mathbb P\{\Xi(U)\ge p\}>0$ for every positive integer $p$.
\end{corollary}

\begin{proof}
Fix $a\in\D$. The zero-free function
\[
 f_a(z)=(1-\overline a z)^{-4/\beta}
\]
uses the holomorphic logarithm equal to zero at $z=0$ and is
holomorphic on a neighborhood of $\overline\D$. Its Taylor
polynomials converge uniformly on $\overline\D$. For each fixed
polynomial, first pass to the subsequential limit in
\eqref{eq:subadditive-interpolation-one}, using the local uniform
convergence of correlation densities. Then pass to $f_a$ by uniform
convergence. The change of variables by a disk automorphism, or
angular integration of the geometric series, gives
\[
 \int_\D |1-\overline a z|^{-4}\,dA(z)
 =\frac{\pi}{(1-|a|^2)^2}.
\]
Also $|f_a(a)|^\beta=(1-|a|^2)^{-4}$. Their quotient proves
\eqref{eq:subadditive-bergman-lower}. Its integral over $\D$ is
infinite.

Pass to the limit directly in
\eqref{eq:subadditive-interpolation-basic} to obtain
\eqref{eq:subadditive-limit-many}. Any nonempty open $U$ contains
$p$ disjoint disks with compact closures in $U$. The Vandermonde is
nonzero on their product, which has positive area. Thus the integral
of $\rho_p$ over that product is strictly positive. This is an
expectation of a nonnegative factorial count and vanishes if
$\Xi(U)<p$ almost surely. The last assertion follows.
\end{proof}

\begin{remark}[What is and is not concluded]
The conclusion that the limiting law is not concentrated on the
empty configuration does not by itself say that its empty-configuration
probability is zero. Similarly, an infinite expected total count does
not by itself prove an almost surely infinite total count. No
monotonicity in $N$, no uniqueness, and no hyperbolic invariance is
used or asserted here. The proof applies to $0<\beta\le1$; extending
its subadditivity step to $1<\beta<2$ is not legitimate.
\end{remark}

\section{Nonempty limits throughout \texorpdfstring{$0<\beta<2$}{0<beta<2}}
\label{sec:weighted-interpolation-nonemptiness}

The interpolation argument extends beyond its pointwise subadditivity
range by using a boundary weight. The important input is a uniform
radial bound under holomorphic product weights, including weights
with zeros. The argument also allows the boundary weight to be
centered at any point in the hyperbolic metric.

\begin{lemma}[Radial domination with an arbitrary holomorphic weight]
\label{lem:holomorphic-weight-radial}
Let $n\ge1$, $\beta>0$, and let $h$ be holomorphic on a neighborhood
of $\overline\D$, not identically zero. Give $\D^n$ density
proportional to
\begin{equation}
 |\Delta_n(\mathbf z)|^\beta\prod_{j=1}^n|h(z_j)|^\beta.
 \label{eq:holomorphic-weight-gas}
\end{equation}
The ordered radial depths are jointly stochastically dominated by
\[
 T_j^{(n)}=\sum_{k=j}^n E_k,\qquad
 E_k\text{ independent exponential of rate }
 a_k=2k+\frac\beta2 k(k-1).
\]
The bound is independent of $h$. Consequently, for $p>1$ and $q>0$,
there is a finite constant $A_{\beta,p,q}$, independent of $n,h$,
such that
\begin{equation}
 \mathbb E\left[\sum_{j=1}^n(1-|z_j|^2)^p\right]^q
 \le A_{\beta,p,q}.
 \label{eq:holomorphic-weight-depth-moment}
\end{equation}
The same bound holds with $|z_j|$ replaced by $|\phi(z_j)|$, where
$\phi$ is any automorphism of $\D$.
\end{lemma}

\begin{proof}
Use the ordered ratios $t_k=r_k/r_{k+1}$, $r_{n+1}=1$, as in the
radial comparison proof. The density is an angular factor times
$\prod_k t_k^{a_k-1}$. To check that this factor is nondecreasing
in $t_k$, rotate the first $k$ angles together. The factors crossing
from these $k$ particles to the others, together with
$\prod_{i\le k}h(z_i)$, are the modulus to the power $\beta$ of a
holomorphic function of the common variable $t_ke^{i\theta}$.
Every remaining factor is nonnegative and independent of this
variable. The circular submean inequality therefore gives the same
monotonicity as before. For $k=n$ there are no crossing factors,
but the product of the $h$ factors is still holomorphic under the
common dilation and rotation. Zeros of $h$ cause no difficulty:
$|h|^\beta$ is subharmonic for every $\beta>0$. Product association
now proves the asserted domination, with exactly the same beta
reference law.

For precision, the moment bound can be written with the finite
constant
\[
 A_{\beta,p,q}=
 \mathbb E\left[\sum_{j=1}^\infty(1-e^{-2T_j})^p\right]^q,
 \qquad T_j=\sum_{k=j}^\infty E_k.
\]
Its finiteness follows from the already established weighted-depth
moment estimate for this reference process. Alternatively, for
$j\ge2$ the exponential-moment estimate
$\log\mathbb E e^{\lambda T_j}\le4\lambda/[\beta(j-1)]$
when $0\le\lambda\le\beta j(j-1)/4$ gives
$\mathbb E T_j^s\le C_{\beta,s}j^{-s}$ for every $s>0$.
For $q\ge1$, Minkowski's inequality and $1-e^{-2t}\le2t$ then give
\[
 \left\|\sum_j(1-e^{-2T_j})^p\right\|_{L^q}
 \le C_{\beta,p,q}\sum_j j^{-p}<\infty;
\]
for $0<q<1$, use Jensen's inequality after the first-moment bound.
The term $j=1$ is bounded by one and does not require a moment
bound on $T_1$ of a different type.

Finally, after a disk automorphism $z=\phi^{-1}(u)$, the new density
has the form~\eqref{eq:holomorphic-weight-gas}, with a new
holomorphic product weight. Indeed,
\[
 |\phi^{-1}(u)-\phi^{-1}(v)|
 =|u-v|\,| (\phi^{-1})'(u)|^{1/2}
              |(\phi^{-1})'(v)|^{1/2}.
\]
The derivative factor in each variable is consequently
$| (\phi^{-1})'|^{\beta(n-1)/2+2}$. Since the derivative is
zero-free on a neighborhood of $\overline\D$, its required real
power is the modulus of a holomorphic function, using a
holomorphic logarithm on a slightly larger disk. Multiplying this
function by $h\circ\phi^{-1}$ gives the new weight. The bound just
proved does not depend on that weight or on $\phi$.
\end{proof}

\begin{theorem}[A positive hyperbolic intensity bound below $2$]
\label{thm:nonemptiness-below-two}
For every $0<\beta<2$ there exists $c_\beta>0$ such that every
subsequential limiting process of the unweighted disk gas satisfies
\begin{equation}
 \rho_1(a)\ge
 \frac{c_\beta}{\pi(1-|a|^2)^2},\qquad a\in\D.
 \label{eq:nonemptiness-below-two}
\end{equation}
For $0<\beta\le1$ one can take $c_\beta=1$. For $1<\beta<2$, put
$q=\beta-1$ and choose any $p>1$ such that $pq<1$. One may take
\begin{equation}
 c_\beta=
 \left(1+\frac{A_{\beta,p,q}}{1-pq}\right)^{-1}>0,
 \label{eq:nonemptiness-below-two-constant}
\end{equation}
with the constant from Lemma~\ref{lem:holomorphic-weight-radial}.
In particular, no subsequence can converge to the almost surely
empty process, and every limiting law has infinite expected total
count.
\end{theorem}

\begin{proof}
The case $\beta\le1$ is
Corollary~\ref{cor:subadditive-bergman-lower}. Suppose
$1<\beta<2$, and fix $a\in\D$. Let $\phi_a$ be a disk
automorphism taking $a$ to zero and set
\[
 w_a(z)=(1-|\phi_a(z)|^2)^p,\qquad
 W_a(\mathbf z)=\sum_{i=1}^N w_a(z_i).
\]
For every polynomial $f$ of degree less than $N$, Lagrange
interpolation and weighted H\"older give the pointwise inequality
\begin{equation}
 |f(a)|^\beta
 \le W_a(\mathbf z)^q
       \sum_{i=1}^N
       \frac{|\ell_i(a;\mathbf z)f(z_i)|^\beta}{w_a(z_i)^q}.
 \label{eq:weighted-lagrange-holder}
\end{equation}
All weights are strictly positive in the open disk. To check the
H\"older step, write each summand $u_i$ as
$(u_i/w_i^{q/\beta})w_i^{q/\beta}$ and use exponents
$\beta$ and $\beta/q$.

The exact cancellation in the one-point interpolation proof now
gives
\begin{align}
 |f(a)|^\beta
 &\le\rho_{1,N}(a)
 \mathbb E_{N}^{!a}\int_\D
 |f(z)|^\beta
 \frac{(W_a(\mathbf y)+w_a(z))^q}{w_a(z)^q}\,dA(z)
 \notag\\
 &\le\rho_{1,N}(a)
 \left(\int_\D |f(z)|^\beta\,dA(z)
 +A_{\beta,p,q}\int_\D
             |f(z)|^\beta w_a(z)^{-q}\,dA(z)\right).
 \label{eq:weighted-interpolation-bound}
\end{align}
Here $\mathbb P_N^{!a}$ is the reduced Palm distribution at $a$,
which is the $(N-1)$-particle gas with holomorphic weight
$h(z)=z-a$. Its normalizer is positive. The first line follows
directly by Tonelli and cancellation; no Palm regularity theorem
is needed. The second line uses $0<q<1$, hence
$(x+y)^q\le x^q+y^q$, and the last assertion of
Lemma~\ref{lem:holomorphic-weight-radial} under this Palm law.
For $N=1$ the exterior sum is zero and the same calculation applies.
The second integral is finite because $pq<1$ and $\phi_a$ is a
smooth diffeomorphism of the closed disk.

Choose the holomorphic branch
\[
 f_a(z)=\phi_a'(z)^{2/\beta}.
\]
It is holomorphic on a neighborhood of $\overline\D$ and satisfies
$|f_a(z)|^\beta=|\phi_a'(z)|^2$. Its Taylor polynomials converge
uniformly on $\overline\D$. For each fixed Taylor polynomial,
first pass to the chosen particle subsequence in
\eqref{eq:weighted-interpolation-bound}. Then pass to $f_a$.
Dominated convergence applies to the second integral because
$w_a^{-q}$ is integrable. Changing variables $u=\phi_a(z)$ gives
\[
 \int_\D|f_a(z)|^\beta\,dA(z)=\pi,
 \qquad
 \int_\D|f_a(z)|^\beta w_a(z)^{-q}\,dA(z)
 =\int_\D(1-|u|^2)^{-pq}\,dA(u)
 =\frac\pi{1-pq}.
\]
Since $|f_a(a)|^\beta=(1-|a|^2)^{-2}$, this proves
\eqref{eq:nonemptiness-below-two} with
\eqref{eq:nonemptiness-below-two-constant}.
\end{proof}

\begin{remark}[Scope of this lower bound]
The intensity estimate does not by itself exclude an empty atom.
The exterior-only interpolation argument below will prove almost sure
infiniteness throughout the same interval. Neither argument proves
uniqueness or hyperbolic invariance. The present weighted H\"older
estimate requires both $p>1$ and $p(\beta-1)<1$, which are
incompatible for $\beta\ge2$.
\end{remark}

\subsection{Strict positivity of every factorial correlation below \texorpdfstring{$2$}{2}}

\begin{theorem}[Uniform lower bounds for every fixed order]
\label{thm:nonemptiness-all-orders-below-two}
Fix $0<\beta<2$ and a positive integer $k$. There is a constant
$b_{\beta,k}>0$ such that
\begin{equation}
 \rho_{k,N}(\mathbf a)\ge
 b_{\beta,k}|\Delta_k(\mathbf a)|^\beta
 \qquad(N\ge k,\ \mathbf a\in\D^k).
 \label{eq:nonemptiness-all-orders-below-two}
\end{equation}
The same bound holds for every subsequential limiting factorial
density. Consequently, for every nonempty open $U\subset\D$ and
every $k\ge1$, every limiting law satisfies
$\mathbb P\{\Xi(U)\ge k\}>0$.
\end{theorem}

\begin{proof}
For $\beta\le1$, use
\eqref{eq:subadditive-interpolation-basic}. Suppose $1<\beta<2$,
put $q=\beta-1$, and fix $s>1$ with $sq<1$. Set
\[
 v(z)=(1-|z|^2)^s,\qquad W(\mathbf z)=\sum_i v(z_i),
 \qquad V_I=\prod_{i\in I}v(z_i).
\]
The elementary symmetric polynomial satisfies
\[
 S_k(\mathbf z):=\sum_{|I|=k}V_I\le W(\mathbf z)^k/k!.
\]
Weighted H\"older applied to the alternating interpolation
identity~\eqref{eq:subadditive-alternating-interpolation} gives
\[
 |F(\mathbf a)|^\beta\le S_k(\mathbf z)^q
 \sum_{|I|=k}
 \frac{\left|\det[\ell_{i_t}(a_j;\mathbf z)]
                         F(\mathbf z_I)\right|^\beta}{V_I^q}
\]
for every alternating polynomial $F$ of degree less than $N$ in
each variable. As in the preceding cancellation calculation,
expectation of the right side equals
\begin{equation}
 \frac{\rho_{k,N}(\mathbf a)}{k!}
 \int_{\D^k}|F(\mathbf z)|^\beta
       \prod_{i=1}^k v(z_i)^{-q}
       \mathbb E_N^{!\mathbf a}S_k(\mathbf y,\mathbf z)^q
                                      \,dA^k(\mathbf z).
 \label{eq:nonemptiness-multipoint-cancellation}
\end{equation}
This identity is first read at distinct anchors $a_j$.
The reduced Palm law in this formula is the $(N-k)$-particle gas
with holomorphic product weight
$h(z)=\prod_{j=1}^k(z-a_j)$. Lemma~\ref{lem:holomorphic-weight-radial}
therefore bounds every moment of $W(\mathbf y)$, uniformly in
$N$ and in the anchors. Since $v\le1$,
\[
 S_k(\mathbf y,\mathbf z)^q
 \le (k!)^{-q}(W(\mathbf y)+k)^{kq}.
\]
Consequently the Palm expectation in
\eqref{eq:nonemptiness-multipoint-cancellation} is at most a finite
constant $B_{\beta,k,s}$, independent of all anchors, the removed
variables, and $N$. The same assertion is immediate for $N=k$,
when $\mathbf y$ is empty.

Take $F=\Delta_k$. The integral
\[
 J_{\beta,k,s}:=\int_{\D^k}|\Delta_k(\mathbf z)|^\beta
           \prod_{i=1}^k(1-|z_i|^2)^{-sq}\,dA^k(\mathbf z)
\]
is finite and positive: the Vandermonde is bounded on the closed
polydisk and each one-variable boundary weight is integrable
because $sq<1$. Thus
$b_{\beta,k}=k!/(B_{\beta,k,s}J_{\beta,k,s})$ is a valid strictly
positive constant. Continuity extends the inequality to colliding
anchors. Local uniform convergence of factorial densities proves
the limiting statement. Integrating over a product of $k$ disjoint
disks inside $U$, as before, proves the final assertion.
\end{proof}

\subsection{A sharp pathwise upper bound on boundary counts}

\begin{proposition}[Sharp leading upper bound for the comparison count]
\label{prop:radial-comparison-sharp-count}
For every subsequential limit at any $\beta>0$, its centered-disk
counts $M(r)=\Xi(\{z:|z|<r\})$ satisfy
\begin{equation}
 \limsup_{r\uparrow1}(1-r)M(r)\le\frac2\beta
 \quad\text{almost surely}.
 \label{eq:radial-comparison-sharp-count}
\end{equation}
Moreover, the family $\{(1-r)M(r):1/2<r<1\}$ is uniformly
integrable.
\end{proposition}

\begin{proof}
For the reference depths $T_m=\sum_{k=m}^\infty E_k$ one has
\[
 \mathbb E T_m=\frac2{\beta m}+O_\beta(m^{-2}),\qquad
 \operatorname{Var}(T_m)=O_\beta(m^{-3}).
\]
These follow by summing the expansions of $a_k^{-1}$ and
$a_k^{-2}$, with $a_k=2k+\beta k(k-1)/2$.
Choose any $\delta>0$ and indices $m_j=\lceil(1+\delta)^j\rceil$.
Chebyshev's inequality shows that the probabilities of
$|m_jT_{m_j}-m_j\mathbb ET_{m_j}|>\varepsilon$ are summable.
Borel--Cantelli therefore gives $m_jT_{m_j}\to2/\beta$ almost
surely. Monotonicity of $T_m$, followed by a countable choice of
$\delta\downarrow0$, proves
\begin{equation}
 mT_m\longrightarrow2/\beta\quad\text{almost surely}.
 \label{eq:reference-depth-sharp-asymptotic}
\end{equation}

The joint ordered-depth domination passes to limiting processes.
Here is one precise way to make this passage. For every finite
$N$, product association gives stochastic domination of the entire
ordered depth vector by $(T_1,\ldots,T_N)$, after padding the
actual vector with zeros. The usual monotone coupling theorem on
countable products gives a coupling with this coordinatewise
inequality. Take a subsequential weak limit of these couplings,
using tightness of the point processes and the fixed reference
law. At every rational radius, the count inequality passes almost
surely through a Skorokhod representation: both the actual limiting
process and the reference process almost surely avoid that circle.
The latter assertion follows since each $T_m$ has a continuous
distribution. The resulting count inequality, first at countably
many rational radii and then at all radii by monotonicity, is
\[
 M(r)\le L(r):=\#\{m:T_m>\log(1/r)\}.
\]
Alternatively, it follows directly by passing the finite-dimensional
stochastic order of the rational-radius count vectors to the limit.
Equation~\eqref{eq:reference-depth-sharp-asymptotic} implies
$(1-r)L(r)\to2/\beta$, proving
\eqref{eq:radial-comparison-sharp-count}.

For uniform integrability one does not need the coupling. The
reference exponential bound already established in the radial
comparison gives
\[
 \mathbb P\{M(r)\ge m\}
 \le\exp\left(m-\frac\beta4m(m-1)\log(1/r)\right)
 \qquad(m\ge2).
\]
Put $s=1-r\in(0,1/2)$ and use $\log(1/r)\ge s$. For all
$t$ larger than a constant depending only on $\beta$, taking
$m=\lfloor t/s\rfloor$ shows that
\[
 \sup_{1/2<r<1}\mathbb P\{(1-r)M(r)>t\}
 \le C_\beta e^{-c_\beta t^2}.
\]
For example, once $t\ge\max(4,32/\beta)$, one has
$m\ge t/(2s)$, $m-1\ge m/2$, and the negative quadratic term
absorbs the positive $m$ term. This uniform Gaussian tail proves
uniform integrability, and indeed uniform moments of every order.
\end{proof}



\subsection{Almost sure infiniteness below \texorpdfstring{$2$}{2}}

The preceding interpolation inequalities can also be applied to a
polynomial which depends on the particles in a fixed interior disk.
The dependence is legitimate because those particles make zero
contributions to the interpolation sum. Every surviving term removes
and reinserts an exterior particle, leaving the polynomial unchanged.

\begin{theorem}[Every limiting process is almost surely infinite below $2$]
\label{thm:almost-sure-infinite-below-two}
For every $0<\beta<2$ and every subsequential limiting process,
\begin{equation}
 \mathbb P\{\Xi(\D)=\infty\}=1.
 \label{eq:almost-sure-infinite-below-two}
\end{equation}
More quantitatively, set $\tau=1$ when $0<\beta\le1$. When
$1<\beta<2$, let $\tau$ be any number satisfying
$0<\tau<2-\beta$. There is a finite constant
$C_{\beta,\tau}$ such that, for every integer $m\ge0$, every
$R\in(0,1)$, and every $N>m$,
\begin{equation}
 \mathbb E_{N,\beta}
 \left[\mathbf1_{\{M_N(R)\le m\}}
                  \prod_{|z_j|<R}|z_j|^\beta\right]
 \le C_{\beta,\tau}2^{m\beta}(1-R^2)^\tau.
 \label{eq:finite-count-product-bound}
\end{equation}
The same inequality holds for every limiting process, with its
finite product over $\Xi\cap B_R$. Empty products equal one.
\end{theorem}

\begin{proof}
Fix $m$, $R$, and $N>m$. On the event $M_N(R)\le m$, form the
monic polynomial
\[
 P_{R,\mathbf z}(t)=\prod_{j:|z_j|<R}(t-z_j).
\]
Its degree is at most $m<N$, so Lagrange interpolation at zero
gives
\begin{equation}
 P_{R,\mathbf z}(0)
 =\sum_{i:|z_i|\ge R}\ell_i(0;\mathbf z)P_{R,\mathbf z}(z_i).
 \label{eq:exterior-only-interpolation}
\end{equation}
Indeed every omitted term contains the zero
$P_{R,\mathbf z}(z_i)=0$. Fixed circles and collision sets have
zero probability, so the distinction between strict and non-strict
inequalities in the exterior index set is immaterial.

Suppose first $0<\beta\le1$. Apply subadditivity to
\eqref{eq:exterior-only-interpolation}, multiply by
$\mathbf1_{\{M_N(R)\le m\}}$, and take expectations. In the
summand indexed by $i$, the variable $z_i$ ranges only over
$\D\setminus B_R$. The polynomial and the count event therefore
depend only on the other variables $\mathbf y$, and are unchanged
when $z_i$ is integrated over this exterior region. The cancellation
already used for deterministic polynomials is thus still exact.
It gives
\begin{align}
 &\mathbb E_{N,\beta}
 \left[\mathbf1_{\{M_N(R)\le m\}}
                       |P_{R,\mathbf z}(0)|^\beta\right]
 \notag\\
 &\quad\le \rho_{1,N}(0)\,
 \mathbb E_N^{!0}\left[
  \mathbf1_{\{M_{\mathbf y}(R)\le m\}}
  \int_{\D\setminus B_R}|P_{R,\mathbf y}(z)|^\beta\,dA(z)
                         \right].
 \label{eq:exterior-interpolation-small-beta}
\end{align}
On the event inside this expectation,
$|P_{R,\mathbf y}(z)|\le2^m$ for all $z\in\D$. The uniform
compact one-point bound gives $\sup_N\rho_{1,N}(0)<\infty$.
The right side of~\eqref{eq:exterior-interpolation-small-beta}
is consequently at most $C_\beta2^{m\beta}(1-R^2)$, proving
\eqref{eq:finite-count-product-bound} in this range.

Now suppose $1<\beta<2$. Set $q=\beta-1$ and choose $s>1$ so
that $sq=1-\tau<1$. Put
\[
 v(z)=(1-|z|^2)^s,\qquad W(\mathbf z)=\sum_jv(z_j).
\]
Weighted H\"older on the surviving exterior summands gives
\[
 |P_{R,\mathbf z}(0)|^\beta
 \le W(\mathbf z)^q
 \sum_{i:|z_i|\ge R}
 \frac{|\ell_i(0;\mathbf z)P_{R,\mathbf z}(z_i)|^\beta}
      {v(z_i)^q}.
\]
Using the larger sum $W$ in place of the exterior weight sum only
increases the right side. The same exterior-only cancellation now
gives the upper bound
\begin{align*}
 \rho_{1,N}(0)\,
 \mathbb E_N^{!0}\left[
 \mathbf1_{\{M_{\mathbf y}(R)\le m\}}
 \int_{\D\setminus B_R}|P_{R,\mathbf y}(z)|^\beta
    \frac{(W(\mathbf y)+v(z))^q}{v(z)^q}\,dA(z)\right].
\end{align*}
The reduced Palm law has holomorphic product weight $h(z)=z$.
Lemma~\ref{lem:holomorphic-weight-radial} bounds
$\mathbb E_N^{!0}W^q$ uniformly. Using
$(W+v)^q\le W^q+v^q$ and $|P_{R,\mathbf y}(z)|\le2^m$ on the
count event therefore gives
\begin{align*}
 &\mathbb E_{N,\beta}
 \left[\mathbf1_{\{M_N(R)\le m\}}
                         |P_{R,\mathbf z}(0)|^\beta\right]\\
 &\qquad\le C_{\beta,s}2^{m\beta}
 \left(\pi(1-R^2)+
               \int_{\D\setminus B_R}(1-|z|^2)^{-sq}\,dA(z)\right)\\
 &\qquad=C_{\beta,s}2^{m\beta}
 \left(\pi(1-R^2)+\frac\pi\tau(1-R^2)^\tau\right).
\end{align*}
Absorb the constants to prove
\eqref{eq:finite-count-product-bound}.

For each fixed $R$, the functional on the left of
\eqref{eq:finite-count-product-bound} is bounded by one and is
continuous at every locally finite configuration avoiding
$\partial B_R$. Every limiting process almost surely avoids this
circle. Thus convergence in law passes the inequality to the
limiting process.

Finally let $R\uparrow1$. On $\{\Xi(\D)>m\}$ the count
indicator is eventually zero. On $\{\Xi(\D)\le m\}$ all
particles are eventually in $B_R$, and the functional converges
to $\prod_{z\in\Xi}|z|^\beta$. Dominated convergence gives
\[
 \mathbb E\left[
 \mathbf1_{\{\Xi(\D)\le m\}}\prod_{z\in\Xi}|z|^\beta\right]=0.
\]
On the indicated finite-count event this finite product is
strictly positive almost surely: the process has no point at the
fixed location zero. Hence $\mathbb P\{\Xi(\D)\le m\}=0$.
A countable union over $m$ proves
\eqref{eq:almost-sure-infinite-below-two}.
\end{proof}

\begin{corollary}[An almost sure logarithmic lower growth bound]
\label{cor:pathwise-logarithmic-count-growth}
Let $0<\beta<2$, and put
\[
 \kappa_\beta=\begin{cases}
 1,&0<\beta\le1,\\
 2-\beta,&1<\beta<2.
 \end{cases}
\]
For every subsequential limiting process, the innermost radius
$r_*:=\inf\{|z|:z\in\Xi\}$ belongs to $(0,1)$ almost surely,
and
\begin{equation}
 \liminf_{r\uparrow1}
 \frac{M(r)}{\log(1/(1-r^2))}
 \ge\frac{\kappa_\beta}{\beta\log(2/r_*)}>0
 \quad\text{almost surely}.
 \label{eq:pathwise-logarithmic-count-growth}
\end{equation}
\end{corollary}

\begin{proof}
Almost sure infiniteness gives $r_*<1$. If $r_*=0$, there are
either infinitely many points in a fixed compact disk or a point
at zero; both alternatives have probability zero. Thus $r_*>0$.
Fix $0<\varepsilon<1$, and let $\tau$ be allowed in
Theorem~\ref{thm:almost-sure-infinite-below-two}. On the event
$\{r_*\ge\varepsilon,\ M(R)\le m\}$ the finite product in
\eqref{eq:finite-count-product-bound} is at least
$\varepsilon^{m\beta}$. Consequently
\[
 \mathbb P\{r_*\ge\varepsilon,\ M(R)\le m\}
 \le C_{\beta,\tau}(2/\varepsilon)^{m\beta}(1-R^2)^\tau.
\]
Take $1-R_n^2=e^{-n}$ and $m_n=\lfloor an\rfloor$, with
$0<a<\tau/[\beta\log(2/\varepsilon)]$. The upper bounds are
summable in $n$, since their exponential rate is at most
$-\tau+a\beta\log(2/\varepsilon)<0$. Borel--Cantelli shows
that, on $\{r_*\ge\varepsilon\}$, eventually $M(R_n)>an$.
Monotonicity of $M$ between $R_n$ and $R_{n+1}$ then gives a
lower bound $a$ for the liminf in
\eqref{eq:pathwise-logarithmic-count-growth}. Make these
conclusions simultaneously for countably many rational
$\varepsilon,a$ and for a countable sequence of allowed $\tau$
increasing to $\kappa_\beta$. On the resulting probability-one
event, let $a$ tend to its upper endpoint and then let rational
$\varepsilon$ increase to $r_*$. This proves the claimed bound.
\end{proof}

\begin{corollary}[A lower tail for the number of interior particles]
\label{cor:lower-count-tail-below-two}
Under the hypotheses and notation of
Theorem~\ref{thm:almost-sure-infinite-below-two}, for every fixed
$m\ge0$ there is a finite constant $C_{\beta,\tau,m}$ such that
\begin{equation}
 \mathbb P\{M(R)\le m\}
 \le C_{\beta,\tau,m}
                 (1-R^2)^{2\tau/(m\beta+2)}.
 \label{eq:lower-count-tail-below-two}
\end{equation}
The same bound holds at finite $N>m$. In particular the hole
probability is $O_{\beta,\tau}((1-R^2)^\tau)$.
\end{corollary}

\begin{proof}
For $0<\varepsilon<1/2$, the uniform compact intensity bound
gives
\[
 \mathbb P\{\Xi(B_\varepsilon)>0\}\le C_\beta\varepsilon^2,
\]
and likewise for every finite gas. On its complement, if
$M(R)\le m$, the product in
\eqref{eq:finite-count-product-bound} is at least
$\varepsilon^{m\beta}$. Consequently
\[
 \mathbb P\{M(R)\le m\}
 \le C_\beta\varepsilon^2+
 C_{\beta,\tau}2^{m\beta}\varepsilon^{-m\beta}(1-R^2)^\tau.
\]
For $R$ sufficiently close to one, choose
$\varepsilon=(1-R^2)^{\tau/(m\beta+2)}$. This proves the
bound near one; increasing its constant covers the remaining
$R$. For $m=0$ the product is identically one on the hole event,
so~\eqref{eq:finite-count-product-bound} already gives the stated
hole bound directly.
\end{proof}

\section{A non-Blaschke theorem for every limit below two}
\label{sec:non-blaschke-limits}

The random interpolation polynomial can be improved by approximating
a finite Blaschke product. This removes the factor exponential in the
interior count and proves a boundary-density statement substantially
stronger than almost sure infiniteness.

\begin{lemma}[Uniform polynomial approximation with prescribed interior zeros]
\label{lem:finite-blaschke-polynomial-approximation}
Fix $0<R<1$, an integer $m\ge0$, and $\varepsilon>0$. There is
an integer $D\ge m$ such that the following holds for every
$0\le n\le m$ and every tuple $\mathbf a=(a_1,\ldots,a_n)$ in
$B_R$. Set
\[
 P_{\mathbf a}(t)=\prod_{j=1}^n(t-a_j),\qquad
 Q_{\mathbf a}(t)=\prod_{j=1}^n(1-\overline{a_j}t),
\]
and let $T_{D-n}(Q_{\mathbf a}^{-1})$ be the Taylor polynomial at
zero through degree $D-n$. The polynomial
\[
 f_{\mathbf a,D}(t)=P_{\mathbf a}(t)
                         T_{D-n}(Q_{\mathbf a}^{-1})(t)
\]
has degree at most $D$, vanishes at each $a_j$, satisfies
$f_{\mathbf a,D}(0)=P_{\mathbf a}(0)$, and obeys
\[
 \sup_{|t|\le1}|f_{\mathbf a,D}(t)|\le1+\varepsilon.
\]
The choice of $D$ is uniform over the indicated tuples.
\end{lemma}
\begin{proof}
For $n\ge1$, the coefficient of $t^k$ in $Q_{\mathbf a}^{-1}$
has modulus at most
$\binom{k+n-1}{n-1}R^k\le\binom{k+m-1}{m-1}R^k$.
This follows by expanding each factor as a geometric series and
counting the weak compositions of $k$ into $n$ parts. Since $R<1$,
\[
 \sup_{|t|\le1}\left|
 Q_{\mathbf a}^{-1}(t)-T_{D-n}(Q_{\mathbf a}^{-1})(t)\right|
 \le\sum_{k>D-m}\binom{k+m-1}{m-1}R^k\longrightarrow0
\]
uniformly in the tuples as $D\to\infty$. On the closed unit disk,
$|P_{\mathbf a}|\le2^m$. Each factor
$(t-a_j)/(1-\overline{a_j}t)$ has modulus at most one there,
so $|P_{\mathbf a}/Q_{\mathbf a}|\le1$. Choose $D$ to make
the Taylor error at most $2^{-m}\varepsilon$ and multiply by
$P_{\mathbf a}$. The remaining assertions follow immediately from
the definition and $Q_{\mathbf a}^{-1}(0)=1$. When $n=0$ take
$f_{\varnothing,D}=1$; the case $m=0$ is included.
\end{proof}

\begin{theorem}[Uniform Laplace decay for accumulated radial depth]
\label{thm:non-blaschke-limits}
Let $0<\beta<2$ and let $\Xi$ be any subsequential limiting
process of the original disk gas. Write
\[
 H(R)=\sum_{z\in\Xi\cap B_R}-\log|z|.
\]
There is almost surely no point at zero, so this is finite for every
$R<1$. Put $\kappa_\beta=1$ if $0<\beta\le1$, and
$\kappa_\beta=2-\beta$ if $1<\beta<2$. For $\beta\le1$ take
$\tau=1$; for $1<\beta<2$ take any $0<\tau<2-\beta$.
There is a finite constant $C_{\beta,\tau}$, uniform over all
subsequential limits, such that
\begin{equation}
 \E e^{-\beta H(R)}
 \le C_{\beta,\tau}(1-R^2)^\tau,
 \qquad 0<R<1.
 \label{eq:non-blaschke-depth-laplace}
\end{equation}
Consequently,
\begin{equation}
 \sum_{z\in\Xi}(1-|z|)=\infty
 \quad\text{almost surely},
 \label{eq:non-blaschke-divergence}
\end{equation}
and the more quantitative pathwise bound is
\begin{equation}
 \liminf_{R\uparrow1}
 \frac{H(R)}{\log(1/(1-R^2))}
 \ge\frac{\kappa_\beta}{\beta}
 \quad\text{almost surely}.
 \label{eq:non-blaschke-depth-growth}
\end{equation}
\end{theorem}

\begin{proof}
Fix $R,m,\varepsilon$, choose $D$ from
Lemma~\ref{lem:finite-blaschke-polynomial-approximation}, and take
$N>D$. On $\{M_N(R)\le m\}$ use $f_{\mathbf a,D}$ with
$\mathbf a$ the tuple of particles in $B_R$. Lagrange interpolation
at zero is valid because the polynomial has degree less than $N$.
Its interior-particle terms vanish, leaving
\[
 P_{\mathbf a}(0)
 =\sum_{i:|z_i|\ge R}\ell_i(0;\mathbf z)
                                      f_{\mathbf a,D}(z_i).
\]
For each surviving term, the particle indexed by $i$ lies outside
$B_R$. Removing it or integrating it over $\D\setminus B_R$
therefore leaves the tuple $\mathbf a$, its count, its polynomial,
and its Taylor truncation degree unchanged. Thus the exact Palm
cancellation used in the proof of
Theorem~\ref{thm:almost-sure-infinite-below-two} applies without
alteration. The polynomial depends measurably and symmetrically on
the interior tuple, and all integrands after the absolute-value
inequality are nonnegative, so Tonelli justifies the integration.

For $0<\beta\le1$, subadditivity and that cancellation give
\begin{align*}
 &\E_{N,\beta}\left[
  \mathbf1_{\{M_N(R)\le m\}}
                      \prod_{|z_j|<R}|z_j|^\beta\right]\\
 &\quad\le\rho_{1,N}(0)\,
 \E_N^{!0}\left[\mathbf1_{\{M_{\mathbf y}(R)\le m\}}
   \int_{\D\setminus B_R}|f_{\mathbf a(\mathbf y),D}(z)|^\beta
                                            \,dA(z)\right]\\
 &\quad\le C_\beta(1+\varepsilon)^\beta(1-R^2).
\end{align*}
Here $\sup_N\rho_{1,N}(0)<\infty$ is the proved compact intensity
bound.

For $1<\beta<2$, put $q=\beta-1$ and
$s=(1-\tau)/q>1$, and use the weights
$v(z)=(1-|z|^2)^s$, $W(\mathbf y)=\sum_jv(y_j)$.
Weighted H\"older followed by the same cancellation bounds the
preceding left side by
\[
 \rho_{1,N}(0)(1+\varepsilon)^\beta
 \E_N^{!0}\int_{\D\setminus B_R}
                 \frac{(W(\mathbf y)+v(z))^q}{v(z)^q}\,dA(z).
\]
The reduced Palm law has holomorphic product weight $h(z)=z$.
Lemma~\ref{lem:holomorphic-weight-radial} bounds its $q$-th
moment of $W$ uniformly in $N$. Since $0<q<1$,
$(W+v)^q\le W^q+v^q$. The last display is therefore at most
\[
 C_{\beta,\tau}(1+\varepsilon)^\beta
 \left((1-R^2)+\int_{\D\setminus B_R}
                           (1-|z|^2)^{-sq}\,dA(z)\right)
 \le C'_{\beta,\tau}(1+\varepsilon)^\beta(1-R^2)^\tau,
\]
because $sq=1-\tau$ and the integral equals
$\pi(1-R^2)^\tau/\tau$. These constants do not depend on $m$,
$D$, or $\varepsilon$.

For fixed $R,m,D$, pass along the chosen particle-number
subsequence. The bounded functional on the left is continuous at
locally finite configurations with no particle on $\partial B_R$,
an event of full limiting probability. Next let
$\varepsilon\downarrow0$, choosing $D$ separately before each
particle-number limit. Finally let $m\to\infty$. Local finiteness
and monotone convergence give
\eqref{eq:non-blaschke-depth-laplace}. This order of limits is
essential: the polynomial degree need not be bounded uniformly in
$m$ or $R$.

As $R\uparrow1$, $H(R)$ increases to the total radial depth.
Dominated convergence in \eqref{eq:non-blaschke-depth-laplace}
shows that the exponential of minus $\beta$ times that total depth
is zero almost surely. Thus $\sum_z-\log|z|=\infty$.
There are only finitely many particles with $|z|\le1/2$, none at
zero, and on $1/2\le r<1$ one has
$1-r\le-\log r\le2(1-r)$. This proves
\eqref{eq:non-blaschke-divergence}.

For the quantitative assertion, take $1-R_n^2=e^{-n}$ and
$0<a<\tau/\beta$. Markov's inequality gives
\[
 \mathbb P\{H(R_n)\le an\}
 \le e^{\beta an}\E e^{-\beta H(R_n)}
 \le C_{\beta,\tau}e^{-(\tau-\beta a)n}.
\]
These bounds are summable. Borel--Cantelli and monotonicity of $H$
between consecutive radii imply a lower bound $a$ for the liminf
in \eqref{eq:non-blaschke-depth-growth}. Take countably many
$a\uparrow\tau/\beta$, and, when $\beta>1$, countably many
$\tau\uparrow2-\beta$. The asserted bound follows.
\end{proof}

\begin{corollary}[Necessary boundary density and divergence of the Green field]
\label{cor:non-blaschke-boundary-consequences}
Every limit in Theorem~\ref{thm:non-blaschke-limits} satisfies
\begin{equation}
 \frac{\kappa_\beta}{\beta}
 \le\limsup_{R\uparrow1}(1-R)\Xi(B_R)
 \le\frac2\beta
 \quad\text{almost surely}.
 \label{eq:non-blaschke-count-limsup}
\end{equation}
Moreover, on one probability-one event,
\begin{equation}
 \sum_{z\in\Xi}g(a,z)=\infty
 \quad\text{for every }a\in\D,
 \qquad g(a,z)=-\log\left|\frac{a-z}{1-\overline a z}\right|.
 \label{eq:non-blaschke-green-divergence}
\end{equation}
No invariance assumption is needed in either assertion.
\end{corollary}
\begin{proof}
The upper bound in \eqref{eq:non-blaschke-count-limsup} is
Proposition~\ref{prop:radial-comparison-sharp-count}. For the lower
bound, write $M(R)=\Xi(B_R)$. Since the innermost point has a
strictly positive radius, Stieltjes integration by parts gives
\[
 H(R)=(-\log R)M(R)+\int_0^R\frac{M(t)}t\,dt.
\]
If $\limsup_{R\uparrow1}(1-R)M(R)<c$, then eventually
$M(R)\le c/(1-R)$. The displayed formula yields
$\limsup_{R\uparrow1}H(R)/\log(1/(1-R^2))\le c$.
Any $c<\kappa_\beta/\beta$ contradicts
\eqref{eq:non-blaschke-depth-growth}, proving the lower bound.

Finally, for $a,z\in\D$, the automorphism identity gives
\[
 1-\left|\frac{a-z}{1-\overline a z}\right|^2
 =\frac{(1-|a|^2)(1-|z|^2)}{|1-\overline a z|^2}
 \ge\frac{1-|a|}{1+|a|}(1-|z|^2).
\]
The inequality $-\log r\ge(1-r^2)/2$ for $0<r\le1$ and
\eqref{eq:non-blaschke-divergence} show that every such Green
sum diverges. The reasoning is deterministic for each configuration
satisfying that divergence, and therefore holds for all $a$
simultaneously. A pole only makes the corresponding sum infinite
immediately.
\end{proof}

\begin{corollary}[Uncompensated grand-canonical equations cannot describe these limits]
\label{cor:actual-limit-uncompensated-green-impossible}
For $0<\beta<2$, no actual subsequential disk-gas limit satisfies
the uncompensated grand-canonical Green specification
\eqref{eq:uncompensated-green-specification} at a positive activity.
This conclusion does not assume hyperbolic invariance.
\end{corollary}
\begin{proof}
For every compact disk $B$, the process has only finitely many
particles in $B$. Removing them leaves an exterior whose Green
potential is infinite at every point of $B$ other than possibly
those finitely many removed particles, by
\eqref{eq:non-blaschke-green-divergence}. Thus every positive-count
term in the uncompensated kernel vanishes almost everywhere on
its tuple space, and that kernel assigns count zero with probability
one. The DLR identity would give $\Xi(B)=0$ almost surely for
every disk in an exhaustion. This contradicts
Theorem~\ref{thm:non-blaschke-limits}.
\end{proof}

\subsection{Local equivalence of all nontrivial candidate limits}
\label{sec:local-poisson-equivalence}

\begin{proposition}[Every local particle count occurs]
\label{prop:local-poisson-equivalence}
Let $P$ be a simple, locally finite point-process law on $\D$
satisfying the strictly positive canonical densities of
Theorem~\ref{thm:unconditional-canonical-quasigibbs}. Suppose that
its total particle count has unbounded support, in the precise sense
that $P\{\Xi(\D)\ge m\}>0$ for every integer $m\ge1$.
Then, for every nonempty disk $B$ with $\overline B\subset\D$,
\begin{equation}
 P\{\Xi(B)=m\}>0\qquad\text{for every integer }m\ge0.
 \label{eq:all-unconditional-local-counts-positive}
\end{equation}
The restriction of $P$ to $B$ is mutually absolutely continuous
with the unit-area Poisson process on $B$.

Consequently every subsequential limit of the disk gas that is not
concentrated on the empty configuration has this local equivalence
property. In particular it holds for every limit when $0<\beta<2$,
and for the Bergman limit at $\beta=2$. Any two nontrivial candidate
limits, even at different inverse temperatures, have mutually
absolutely continuous restrictions to every relatively compact disk.
\end{proposition}

\begin{proof}
Fix $m\ge0$. If $m\ge1$, choose a disk $D_0$ compactly contained
in $\D$, containing $\overline B$ in its interior, such that
$P\{\Xi(D_0)\ge m\}>0$. Such a disk exists by monotone convergence
along a centered exhaustion containing $\overline B$ and the
unbounded-support assumption. If $m=0$, choose any such $D_0$.
Local finiteness yields an integer $n\ge m$ with
$P\{\Xi(D_0)=n\}>0$; for $m=0$ this follows simply because
the probabilities over all finite $n$ sum to one.

Conditional on the exterior of $D_0$ and this count $n$, the
interior law has a strictly positive density away from collisions
on $D_0^n$. The set of tuples with exactly $m$ coordinates in $B$
and $n-m$ coordinates in $D_0\setminus\overline B$ has positive
area when $n>0$, because both regions have positive area; if $n=0$
the desired event already holds. Hence the conditional probability
of $\Xi(B)=m$ is strictly positive for almost every exterior
occurring with count $n$. Integrating proves
\eqref{eq:all-unconditional-local-counts-positive}.

For each $m\ge1$, condition now on the exterior of $B$ and
$\Xi(B)=m$. The canonical density is equivalent to area measure
on $B^m$, since the collision set has zero area and the remaining
factor is strictly positive. Averaging this density over the
exterior shows that the unconditional $m$th Janossy measure is
equivalent to area measure on $B^m$: a set of positive area has
positive conditional probability for every admissible exterior,
and the unconditional mass of this count sector is positive.
The empty sector also has positive mass by the first part. The
Poisson law has exactly these count sectors, each with a strictly
positive constant Janossy density. Its countable decomposition
therefore proves mutual absolute continuity of the complete local
laws.

For an actual subsequential limit, the canonical hypothesis is
Theorem~\ref{thm:unconditional-canonical-quasigibbs}. For
$0<\beta<2$, its total count is infinite almost surely by
Theorem~\ref{thm:almost-sure-infinite-below-two}. For $\beta\ge2$,
Theorem~\ref{thm:finite-count-exclusion} says that its total count
belongs to $\{0,\infty\}$ almost surely. Unless the law is the
empty point mass, it thus has positive probability of infinite
total count and satisfies the same unbounded-support hypothesis.
The final assertion follows by transitivity of local equivalence.
\end{proof}

\begin{remark}[Local overlap does not exclude different phases]
The proposition concerns unconditional local laws. It does not prove
that every count has positive probability after conditioning on a
fixed exterior. Nor does local mutual absolute continuity imply
global mutual absolute continuity or equality of point-process laws.
If distinct candidate states exist, the proposition says they cannot
be separated by an event of probability zero in a fixed compact
window. Their probabilities of positive-probability local events
may still differ, and they may be globally singular because of
tail information.
\end{remark}

\section{A Blaschke interpolation criterion for the interval below \texorpdfstring{$4$}{4}}
\label{sec:blaschke-boundary-moment-criterion}

The next result is a conditional reduction. It isolates a single
uniform boundary estimate which would prove nonemptiness in the
remaining interval $2<\beta<4$. The estimate itself is not proved
here. Ordinary Blaschke interpolation suffices; a Hermite formula,
which would introduce derivatives of the exterior field, is unnecessary.

\begin{lemma}[The total radial depth cannot vanish]
\label{lem:radial-depth-not-vanishing}
For every fixed $\beta>0$ there is a finite $C_\beta$ such that
\begin{equation}
 \mathbb P_{N,\beta}\left\{\sum_{i=1}^N-\log|z_i|\le u\right\}
 \le (C_\beta u)^N,\qquad N\ge1,\quad u>0.
 \label{eq:radial-depth-small-simplex}
\end{equation}
Consequently, for every $b>0$ there is $d_{\beta,b}>0$ such that
\begin{equation}
 \inf_{N\ge1}\mathbb E_{N,\beta}
 \left(1-\prod_{i=1}^N|z_i|^2\right)^b
 \ge d_{\beta,b}.
 \label{eq:blaschke-numerator-lower}
\end{equation}
\end{lemma}

\begin{proof}
The established partition asymptotic implies, after decreasing its
constant to cover the finitely many small values of $N$, that
\[
 Z_{N,\beta}\ge c_\beta^N N^{(\beta/2-1)N}
\]
for some $c_\beta>0$ and every $N\ge1$. Hadamard's inequality gives
$|\Delta_N|^\beta\le N^{\beta N/2}$ on the disk. Put
$t_i=-\log|z_i|$. After angular integration the area element is
$(2\pi)^N e^{-2\sum_i t_i}\,dt_1\cdots dt_N$. The simplex
$\{t_i\ge0:\sum_i t_i\le u\}$ has volume $u^N/N!$. Therefore
\[
 \mathbb P\left\{\sum_i t_i\le u\right\}
 \le \frac{(2\pi)^N N^{\beta N/2}u^N}
 {N!c_\beta^N N^{(\beta/2-1)N}}
 \le \left(\frac{2\pi e}{c_\beta}u\right)^N,
\]
where $N!\ge(N/e)^N$. Choose $u_\beta>0$ with
$C_\beta u_\beta\le1/2$. On the complementary event,
$1-\prod_i|z_i|^2\ge1-e^{-2u_\beta}$. Thus
$d_{\beta,b}=\tfrac12(1-e^{-2u_\beta})^b$ is valid.
\end{proof}

For $n\ge0$, let $\mathbb P_{n,\beta}^{[0]}$ denote the probability
law on $\D^n$ with density proportional to
\[
 |\Delta_n(\mathbf y)|^\beta\prod_{j=1}^n|y_j|^\beta.
\]
It is exactly the reduced Palm distribution at zero of the
$(n+1)$-particle unweighted gas. For $n=0$ it is the empty
configuration. Define its reflected characteristic polynomial by
\[
 Q_{\mathbf y}(z)=\prod_{j=1}^n(1-z\overline{y_j}),
 \qquad z\in\D.
\]

\begin{theorem}[A sufficient uniform boundary moment estimate]
\label{thm:blaschke-boundary-moment-criterion}
Fix $1<\beta<4$. Suppose that there is $K_\beta<\infty$ such that
\begin{equation}
 \sup_{n\ge0}\mathbb E_{n,\beta}^{[0]}
       |Q_{\mathbf y}(z)|^\beta
 \le K_\beta(1-|z|^2)^{-\beta/2},\qquad z\in\D.
 \label{eq:blaschke-boundary-moment-hypothesis}
\end{equation}
Then there is $c>0$ such that $\rho_{1,N}(0)\ge c$ for every
$N\ge1$. In particular every subsequential limiting law is
different from the almost surely empty law.

The hypothesis is a bound uniform both in particle number and as
$|z|\uparrow1$. Pointwise convergence of the left side for fixed
$z$, or uniform bounds on fixed compact subdisks, do not imply it.
\end{theorem}

\begin{proof}
For a configuration of distinct points put
\[
 P(\zeta)=\prod_{i=1}^N(\zeta-z_i),\quad
 Q(\zeta)=\prod_{i=1}^N(1-\zeta\overline{z_i}),\quad
 B(\zeta)=P(\zeta)/Q(\zeta).
\]
The polynomial part of the partial-fraction expansion of $Q/P$
is the constant $\overline{B(0)}$. Its residue at $z_i$ is
$1/B'(z_i)$. Hence, initially away from the zeros of $P$,
\[
 1-\overline{B(0)}B(a)
 =\sum_{i=1}^N\frac{B(a)}{(a-z_i)B'(z_i)}.
\]
Equivalently, writing the Lagrange polynomials as before,
\begin{equation}
 1-\overline{B(0)}B(a)
 =\sum_{i=1}^N\ell_i(a;\mathbf z)\frac{Q(z_i)}{Q(a)}.
 \label{eq:blaschke-lagrange-identity}
\end{equation}
Both sides extend continuously to all $a\in\D$. At $a=0$,
$Q(0)=1$ and the left side is $1-\prod_i|z_i|^2$.

Set $q=\beta-1>0$. Choose
\begin{equation}
 1<p<\frac{\beta+2}{2(\beta-1)}.
 \label{eq:blaschke-weight-range}
\end{equation}
Such a choice is possible precisely because $\beta<4$. Put
$v(z)=(1-|z|^2)^p$ and $W(\mathbf z)=\sum_i v(z_i)$.
Weighted H\"older applied to~\eqref{eq:blaschke-lagrange-identity}
gives
\[
 \left(1-\prod_i|z_i|^2\right)^\beta
 \le W(\mathbf z)^q
 \sum_i\frac{|\ell_i(0;\mathbf z)Q(z_i)|^\beta}{v(z_i)^q}.
\]
No subadditivity assumption on $q$ is used here.

In the summand indexed by $i$, let $z=z_i$ and let $\mathbf y$
be the other particles. Then
$Q(z_i)=(1-|z|^2)Q_{\mathbf y}(z)$. Cancelling the Lagrange
denominators against the gas density and summing over $i$ gives,
by Tonelli,
\begin{align}
 &\mathbb E_{N,\beta}
       \left(1-\prod_i|z_i|^2\right)^\beta\notag\\
 &\quad\le\rho_{1,N}(0)
 \int_\D (1-|z|^2)^{\beta-pq}
 \mathbb E_{N-1,\beta}^{[0]}
 \left[(W(\mathbf y)+v(z))^q|Q_{\mathbf y}(z)|^\beta\right]dA(z).
 \label{eq:blaschke-exact-palm-bound}
\end{align}
The auxiliary moment of $W$ is harmless. Indeed, tilting
$\mathbb P_{n,\beta}^{[0]}$ by $|Q_{\mathbf y}(z)|^\beta$
produces a gas with holomorphic product weight
\[
 h_z(w)=w(1-\overline z w).
\]
Lemma~\ref{lem:holomorphic-weight-radial} therefore implies
\begin{equation}
 \mathbb E_{n,\beta}^{[0]}
 \left[(W+v(z))^q|Q_{\mathbf y}(z)|^\beta\right]
 \le A_{\beta,p,q}'\,
       \mathbb E_{n,\beta}^{[0]}|Q_{\mathbf y}(z)|^\beta,
 \label{eq:blaschke-weighted-moment-removal}
\end{equation}
with a finite constant independent of $n,z$. For example, take
$A_{\beta,p,q}'$ to be a uniform bound for
$\mathbb E(1+W)^q$ under arbitrary holomorphic product weights.
The tilted normalizers are positive, and the case $n=0$ is immediate.

By the hypothesis, the right side of
\eqref{eq:blaschke-exact-palm-bound} is at most
\[
 \rho_{1,N}(0)A_{\beta,p,q}'K_\beta
 \int_\D(1-|z|^2)^{\beta/2-pq}\,dA(z)
 =\rho_{1,N}(0)A_{\beta,p,q}'K_\beta
       \frac{\pi}{1+\beta/2-pq}.
\]
The denominator is strictly positive by
\eqref{eq:blaschke-weight-range}. Lemma~\ref{lem:radial-depth-not-vanishing},
with $b=\beta$, bounds the left side away from zero. This proves
the asserted uniform lower bound on $\rho_{1,N}(0)$. Local uniform
convergence of the correlation densities then proves the limiting
conclusion.
\end{proof}

\begin{corollary}[Any uniform boundary exponent strictly below $2$ suffices]
\label{cor:blaschke-weaker-boundary-exponent}
In Theorem~\ref{thm:blaschke-boundary-moment-criterion}, the
hypothesis can be replaced by
\begin{equation}
 \sup_{n\ge0}\mathbb E_{n,\beta}^{[0]}|Q_{\mathbf y}(z)|^\beta
 \le K(1-|z|^2)^{-\kappa},\qquad z\in\D,
 \label{eq:blaschke-weaker-boundary-exponent}
\end{equation}
for any finite $K$ and any $\kappa<2$. The same uniform positive
lower bound on $\rho_{1,N}(0)$ follows.
\end{corollary}

\begin{proof}
Choose
\[
 1<p<\frac{\beta+1-\kappa}{\beta-1},
\]
which is possible because $\kappa<2$ and $\beta>1$. The proof of
the theorem, through~\eqref{eq:blaschke-weighted-moment-removal},
is unchanged. The final area integral is now
\[
 \int_\D(1-|z|^2)^{\beta-p(\beta-1)-\kappa}\,dA(z)
 =\frac{\pi}{1+\beta-p(\beta-1)-\kappa}<\infty.
\]
The same numerator lower bound completes the proof.
\end{proof}

\begin{remark}[A Gaussian coefficient contraction already fails for one particle]
Write $b=\beta/2$. For the one-particle charged law
$\mathbb P_{1,\beta}^{[0]}$, direct radial integration gives
\[
 \mathbb E_{1,\beta}^{[0]}|y|^2
 =\frac{\beta+2}{\beta+4}=\frac{b+1}{b+2}.
\]
The coefficient of $z$ in $(1-z\overline y)^b$ is
$-b\overline y$. Rotational symmetry therefore gives
\[
 \mathbb E_{1,\beta}^{[0]}|1-z\overline y|^\beta
 =1+\frac{b^2(b+1)}{b+2}|z|^2+O(|z|^4).
\]
By comparison, $(1-|z|^2)^{-b}=1+b|z|^2+O(|z|^4)$.
The first nonconstant coefficient of the former is larger when
$b(b+1)/(b+2)>1$, equivalently $b>\sqrt2$. Thus the proposed
Gaussian domination with multiplicative constant exactly one is
false already for this charged one-particle law when
$2\sqrt2<\beta<4$. This does not disprove either boundary-moment
hypothesis with an unspecified finite multiplicative constant.
\end{remark}

\begin{remark}[The unresolved step]
The criterion is a proved implication, not a proof of
\eqref{eq:blaschke-boundary-moment-hypothesis} for $2<\beta<4$.
The reduction avoids any exterior electric-field estimate and
identifies an explicit charged characteristic-polynomial moment
as a possible next target. It does not prove uniqueness, hyperbolic
invariance, or almost sure infiniteness in that remaining interval.
\end{remark}

\section{The corrected meromorphic field is independent of the Gaussian modes}
\label{sec:gaussian-separation-meromorphic-field}

The independence of the harmonic modes from the point configuration can
be strengthened to include a corrected electric field. The correction
is intrinsic at finite particle number: it is the derivative of the
uncompensated hyperbolic Green potential. The required small-transport
estimate uses its vanishing at the boundary, rather than an estimate
on inverse angular statistics.

For a finite configuration $\xi$ put
\[
 J_w(z)=\frac1{z-w}+\frac{\overline w}{1-z\overline w}
       =\frac{1-|w|^2}{(z-w)(1-z\overline w)},\qquad
 \mathcal A_\xi(z)=\sum_{w\in\xi}J_w(z).
\]
Write $\mathcal A_N=\mathcal A_{\Xi_N}$ and
\[
 \mathcal C_N(z)=\sum_{w\in\Xi_N}\frac1{z-w},\qquad
 T_N(z)=\sum_{w\in\Xi_N}\frac{\overline w}{1-z\overline w}
       =\sum_{k\ge1}\overline{S_{k,N}}z^{k-1}.
\]
Thus $\mathcal A_N=\mathcal C_N+T_N$ exactly. These are meromorphic
functions on the disk, with their poles interpreted in the usual sense.

\begin{lemma}[Stability of the corrected field under small boundary transports]
\label{lem:green-force-small-transport-stability}
Let $T_{N,t}$ be any one of the mode transports used in the proof of
Theorem~\ref{thm:all-harmonic-modes}, including the small M\"obius
transport for the first mode. For every finite list of deterministic
points $u_1,\ldots,u_l\in\D$,
\[
 \bigl(\mathcal A_{(T_{N,t})_*\Xi_N}(u_i)
             -\mathcal A_{\Xi_N}(u_i)\bigr)_{i=1}^l
       \longrightarrow0
\]
in probability under both the original law $P_N$ and its exact
pullback law $Q_{N,t}$. Each evaluation is finite almost surely.
\end{lemma}

\begin{proof}
Fix $\max_i|u_i|<R<1$. On a slightly larger outer annulus the maps
and their inverses have $C^1$ distance $O(N^{-1})$ from the identity.
This follows directly from their displayed formula outside the
shrinking patch near zero; for the first mode it follows from the
M\"obius formula. They map the unit circle onto itself.
Globally, $\epsilon_N:=\sup|T_{N,t}-\mathrm{id}|\to0$.

For $|w|\ge R$, all derivatives in $w$ up to order two of $J_w(u_i)$
are bounded uniformly. For $N$ large the same is true along the
image of the annulus. The real $C^1$ norm of
\[
 F_{N,i}(w)=J_{T_{N,t}(w)}(u_i)-J_w(u_i)
\]
is consequently $O(N^{-1})$. Moreover $F_{N,i}=0$ on $|w|=1$,
because $J_w(u_i)=0$ there and the transport preserves that circle.
Integration of its radial derivative gives
\begin{equation}\label{eq:green-force-boundary-transport-bound}
 |F_{N,i}(w)|\le \frac{C}{N}(1-|w|),\qquad |w|\ge R.
\end{equation}
The outer-particle contribution is at most $CD_N/N$, where
$D_N=\sum_w(1-|w|)$. It tends to zero in probability under both laws:
$D_N=O_{P_N}(\log(eN))$, and the same assertion under $Q_{N,t}$
was proved in the mode-transport argument using its exact pullback
and uniform preservation of boundary distance.

For the inner particles, fix $\epsilon>0$ and exclude the event that
some particle lies within $2\epsilon$ of an evaluation point. Under
$P_N$ its probability is $O(\epsilon^2)$ by the compact intensity
bound. Under $Q_{N,t}$ the same bound, with an error tending to zero,
follows by transporting to $P_N$ and enlarging the finitely many balls
by $\epsilon_N$. Away from these balls, the derivatives of $J_w(u_i)$
are bounded by a constant depending on $R,\epsilon$, so the remaining
inner contribution is at most
$C_{R,\epsilon}\epsilon_N\Xi_N(\D_R)$. The compact count is tight
under both laws, hence this tends to zero in probability. First let
$N\to\infty$, then $\epsilon\downarrow0$. The finite list of
evaluations causes no change to the argument.
\end{proof}

\begin{theorem}[Gaussian separation of the corrected meromorphic field]
\label{thm:corrected-field-gaussian-independence}
Every subsequential point-process limit admits a simultaneous
augmentation by its meromorphic field $\mathcal C$ on
$\{\Xi\ne\varnothing\}$ and the Gaussian sequence $G=(G_k)_{k\ge1}$
such that
\[
 \mathcal A=\mathcal C+T_G
       \quad\hbox{on }\{\Xi\ne\varnothing\},\qquad
 T_G(z)=\sum_{k\ge1}\overline{G_k}z^{k-1},
\]
has the following properties. Set $\mathcal A=0$ on the empty event.
Then
\begin{equation}\label{eq:corrected-field-gaussian-independence}
             (\Xi,\mathcal A)\ \hbox{is independent of }G.
\end{equation}
On the nonempty event, $\mathcal A$ is a meromorphic function with
simple residue-one poles exactly at $\Xi$. It is obtained as the
limit of $\mathcal A_N$ on every positive-count localization.
The assertion does not say that $\mathcal A$ is measurable from $\Xi$.
\end{theorem}

\begin{proof}
Retain all normalized exterior factors on the fixed countable disk
exhaustion in Proposition~\ref{prop:meromorphic-field-gluing}, along
with all the harmonic modes. Theorem~\ref{thm:analytic-harmonic-laplace}
also gives $T_N\Rightarrow T_G$ in $\mathcal H(\D)$: conjugate its
coefficients and divide the analytic resolvent by its argument,
with the removable value at zero. For each positive-count disk the
zero-free limits of the normalized factors identify the locally
meromorphic limit of $\mathcal C_N$. Gluing therefore gives the
asserted limit of $\mathcal A_N$ on the nonempty event.

We give the localization needed to avoid assuming tightness of
$\mathcal C_N$ on the empty branch. Fix a radius $r<1$ from the
exhaustion, a finite deterministic evaluation list $\mathbf u$, and
write $E_{r,N}=\{\Xi_N(\D_r)\ge1\}$.
Introduce a cemetery symbol $*$, give $\C^l\sqcup\{*\}$ its Polish
disjoint-union topology, and define the censored vector
\[
 A_{r,N}=\begin{cases}
 (\mathcal A_N(u_1),\ldots,\mathcal A_N(u_l)),&E_{r,N},\\
 *,&E_{r,N}^{\mathrm c}.
 \end{cases}
\]
Its law jointly with $\Xi_N$ is tight along the chosen subsequence.
Indeed, on every fixed positive count the preceding simultaneous
normal-family construction gives a nonzero normalized limit on a
larger disk containing the evaluation list; the remaining high counts
have arbitrarily small probability by the uniform count bounds.
No deterministic evaluation point is a limiting particle, by absolute
continuity of the limiting intensity. Thus the limiting censored
vector is $(\mathcal A(u_i))_i$ on $E_r=\{\Xi(\D_r)\ge1\}$ and
$*$ otherwise. These conclusions follow for all countably many
radii and all finite lists from a fixed countable dense set on a
single diagonal extraction.

Under the mode transport, $E_{r,N}$ changes with probability tending
to zero under both $P_N$ and $Q_{N,t}$. A change requires a particle
in a shrinking collar of $\partial\D_r$; the compact intensity bound
and exact pullback control this probability. Lemma
\ref{lem:green-force-small-transport-stability} shows that the
censored field vector is also unchanged in probability. The point
configuration is unchanged in the vague metric in probability.
For every fixed mode $j$ other than the transported mode $k$, its
increment tends to zero in probability as well. For $j<k$ this is
\eqref{eq:allmodes-lower-invariant}; for $j>k$ the same expansion
\eqref{eq:allmodes-mode-increment} has only positive indices and
the previously proved macroscopic mode bounds apply. For the first
mode the identical calculation follows by expanding the M\"obius
map. All these assertions hold under $Q_{N,t}$.

Consequently, one may add $A_{r,N}$ and any finite list of modes
other than $k$ to the unaltered observables $Y_N$ in the proof of
\eqref{eq:allmodes-test-identity}. The tightness just established and
the exact pullback imply the same tightness under $Q_{N,t}$, because
the transported observables have their original $P_N$ laws and the
differences tend to zero in probability. Bounded continuous tests
therefore give, in the limit,
\begin{equation}\label{eq:corrected-field-mode-laplace-identity}
 \mathbb E\left[F(\Xi,A_r,(G_j)_{j\ne k})
                     e^{-2\operatorname{Re}(\overline tG_k)}\right]
 =e^{2k|t|^2/\beta}\,\mathbb EF(\Xi,A_r,(G_j)_{j\ne k}),
\end{equation}
where initially only finitely many of the other modes are tested.
Precisely the same uniform integrability of the exact likelihood
ratios as in \eqref{eq:allmodes-ui} justifies this passage; no bound
on exponential moments of the meromorphic mark is required.

Uniqueness of the Gaussian Laplace transform identifies the
conditional distribution of $G_k$ given these observables as its
unconditional Gaussian distribution. Bounded continuous tests extend
to bounded Borel tests by the monotone-class theorem. Applying the
identity successively to each coordinate of a finite Gaussian vector
shows that this entire vector is independent of $(\Xi,A_r)$: at
each step the remaining mode tests and the bounded test of
$(\Xi,A_r)$ are admissible. Countable-product uniqueness then gives
independence from the full sequence $G$.

Finally, on $E_r$ the censored mark is the actual field evaluation.
Let $r\uparrow1$ through the exhaustion and use bounded convergence
to obtain independence of all finite field evaluations together with
$\Xi$, restricted to $\{\Xi\ne\varnothing\}$. On the empty event
the field is assigned zero, and its contribution factors because
$\Xi$ is already known to be independent of $G$.
Evaluations on a countable dense set determine a meromorphic function
and generate its Borel sigma-algebra. This proves
\eqref{eq:corrected-field-gaussian-independence}.
\end{proof}

\begin{remark}[The convention on the empty branch]
The earlier gluing construction set $\mathcal C=0$ on the empty
configuration. The present theorem instead specifies $\mathcal A=0$
there, so it is convenient to set $\mathcal C=-T_G$ on that branch
when the relation $\mathcal A=\mathcal C+T_G$ is desired everywhere.
There are no particle force marks on that event, and every canonical
resampling has count zero. Thus this convention changes none of the
previous force identities or positive-count DLR equations. No
convergence of the uncensored $\mathcal C_N$ on that branch is claimed.
\end{remark}

\section{Eliminating the independent Gaussian field from the canonical equations}
\label{sec:gaussian-elimination-green-specification}

The independence statement of
Theorem~\ref{thm:corrected-field-gaussian-independence} is stronger
than regression of the force at the particles. It permits exact Gaussian
elimination in the augmented canonical equations. The resulting kernels
retain a meromorphic field, but no longer retain the independent Gaussian
modes. Their pair interaction is the hyperbolic Green interaction.

Write
\[
 s_k(\mathbf x)=\sum_{i=1}^m x_i^k,\qquad
 T_g(z)=\sum_{k\ge1}\overline{g_k}z^{k-1},\qquad
 J_p(z)=\frac1{z-p}+\frac{\overline p}{1-z\overline p}.
\]
Thus $J_p=\partial_z\log\delta(z,p)^2$, where
$\delta(z,p)=|(z-p)/(1-z\overline p)|$.
Let $\mathcal A=\mathcal C+T_G$ on the nonempty event, with the empty
branch defined as in the independence theorem for $\mathcal A$.
For $B=B_j$ set
\begin{equation}\label{eq:gaussian-free-exterior-field}
 \mathcal R_B(z)=\mathcal A(z)-\sum_{p\in\Xi\cap B}J_p(z).
\end{equation}
This function is holomorphic on a neighborhood of $\overline B$.
On the nonempty event, if
$H_{B,k}=G_k-s_k(\Xi\cap B)$, then
\begin{equation}\label{eq:R-E-H-identity}
 \mathcal R_B=\mathcal E_B+T_{H_B}.
\end{equation}
Here $T_{H_B}$ converges locally uniformly: the Gaussian coefficients
have at most polynomial growth almost surely, and the subtracted finite
power sums decay geometrically.

\begin{lemma}[Retaining the exterior positive modes]
\label{lem:canonical-retain-exterior-positive-modes}
In the simultaneous construction of the meromorphic augmentation, all
the exterior sequences $H_B$ can also be retained. For every positive
allowed count $m$, the conditional interior density given
$(\eta,m,\mathcal E_B,H_B)$ is the original canonical density with
harmonic field $\mathcal E_B$. Consequently, given
$(\eta,m,\mathcal R_B,H_B=h)$ it is proportional to
\begin{equation}\label{eq:canonical-density-before-gaussian-elimination}
 |\Delta_m(\mathbf x)|^\beta
 \exp\left\{\beta\operatorname{Re}\sum_i V_{\mathcal R_B}(x_i)
       -\beta\operatorname{Re}
            \sum_{k\ge1}\frac{\overline{h_k}s_k(\mathbf x)}k\right\}
 \,dA^m(\mathbf x),
\end{equation}
where $V_{\mathcal R_B}'=\mathcal R_B$ and $V_{\mathcal R_B}(0)=0$.
\end{lemma}

\begin{proof}
At finite $N$ the quantities
$H_{B,k,N}=S_{k,N}-s_k(\Xi_N\cap B)$ depend only on the exterior.
The finite conditional density therefore remains exactly the same
when any finite number of them is included in its conditioning.
In the normal-family proof of
Theorem~\ref{thm:augmented-canonical-dlr}, multiply the exterior test
function by a bounded continuous function of any finite list of these
coordinates. Their joint tightness follows from the Gaussian-mode
tightness and the local count bounds. Along the simultaneous extraction,
the finite interior power sums converge because $\partial B$ contains
no limiting particle. Hence their limits are precisely the displayed
$H_{B,k}$.

The compact-test identity passes to the limit just as in that theorem:
the integral against an interior compactly supported test is continuous
in the holomorphic factor, and the additional bounded continuous
exterior-mode test causes no change. Exhausting the interior tuple
space proves that the limiting normalizing integral is one. The
monotone-class theorem, first in each finite list of mode coordinates
and then in their countable union, gives the asserted conditional
identity with the whole sequence. Countably many disks and counts
can be treated on the same extraction.

On the nonempty event, knowing $(\mathcal R_B,H_B)$ is equivalent to
knowing $(\mathcal E_B,H_B)$ by \eqref{eq:R-E-H-identity}. A primitive
of $T_h$ vanishing at zero is
$\sum_{k\ge1}\overline{h_k}z^k/k$. Substitution into the harmonic
canonical density proves
\eqref{eq:canonical-density-before-gaussian-elimination}.
\end{proof}

\begin{lemma}[Gaussian Bayes formula on a compact tuple space]
\label{lem:infinite-gaussian-bayes-tuples}
Let $X\in B^m$, where $B=\D_r$, $r<1$, and let $Z$ be any random
variable taking values in a standard Borel space. Suppose that
$G=(G_k)_{k\ge1}$ is independent of $(X,Z)$, with independent circular
complex coordinates of variance $2k/\beta$. Put $H=G-s(X)$.
Fix a deterministic tuple $\mathbf a\in B^m$ and write
\[
 \|s(\mathbf x)\|_*^2=\sum_{k\ge1}\frac{|s_k(\mathbf x)|^2}k.
\]
The Gaussian law of $H$ given $X=\mathbf x$ has density, relative to
the law $\gamma_{-s(\mathbf a)}$ of $G-s(\mathbf a)$,
\begin{equation}\label{eq:gaussian-tuple-shift-density}
 r(\mathbf x,h)=\exp\left\{
 -\beta\operatorname{Re}\sum_{k\ge1}
       \frac{\overline{h_k}[s_k(\mathbf x)-s_k(\mathbf a)]}k
 -\frac\beta2\bigl(\|s(\mathbf x)\|_*^2
                      -\|s(\mathbf a)\|_*^2\bigr)\right\}.
\end{equation}
For almost every $h$ this is continuous in $\mathbf x$ on
$\overline B^m$, strictly positive, and bounded above and below by
positive finite constants depending on $h$. If $\Pi_Z$ denotes the
conditional law of $X$ given $Z$, then
\begin{equation}\label{eq:gaussian-bayes-tuples}
 \mathbb P(X\in d\mathbf x\mid Z,H=h)
 =\frac{r(\mathbf x,h)\Pi_Z(d\mathbf x)}
        {\int r(\mathbf y,h)\Pi_Z(d\mathbf y)}.
\end{equation}
\end{lemma}

\begin{proof}
For finitely many coordinates, divide their explicit Gaussian densities.
This gives the finite truncation of
\eqref{eq:gaussian-tuple-shift-density}. All shifts belong to the
Gaussian Cameron--Martin space because
$|s_k(\mathbf x)|\le m r^k$ and
$\sum_k r^{2k}/k<\infty$.
For completeness, under the reference law write $h=G-s(\mathbf a)$.
The logarithm of the truncated density is a centered real Gaussian
linear statistic minus half its variance. The variances are uniformly
bounded, since they equal
$\beta\sum_{k\le L}|s_k(\mathbf x)-s_k(\mathbf a)|^2/k$.
Its exponentials therefore converge in $L^1$, for example by their
uniform $L^2$ bound, and the limit has expectation one. Agreement on
all finite-coordinate cylinder events proves that the limit is the
Radon--Nikodym density of the full translated product measure.

Almost surely $|G_k|\le k^2$ eventually, by the Gaussian tail estimate
and the summability of those tails. Consequently the series involving
$h$ in \eqref{eq:gaussian-tuple-shift-density} converges absolutely and
uniformly for $\mathbf x\in\overline B^m$. The norm series does too.
The stated continuity and positive finite bounds follow. Independence
of $G$ and $(X,Z)$ now expresses their joint conditional law, given
$Z$, as
$\Pi_Z(d\mathbf x)r(\mathbf x,h)\gamma_{-s(\mathbf a)}(dh)$.
Disintegration gives \eqref{eq:gaussian-bayes-tuples}; its denominator
is finite and strictly positive by the bounds just proved.
\end{proof}

\begin{theorem}[A Gaussian-free augmented Green specification]
\label{thm:gaussian-free-green-specification}
For every subsequential limit, use the augmentation from
Theorem~\ref{thm:corrected-field-gaussian-independence}, for which
$(\Xi,\mathcal A)$ is independent of all the Gaussian modes $G_k$,
together with the retained-mode canonical identity of
Lemma~\ref{lem:canonical-retain-exterior-positive-modes}.
For every $B=B_j$ and every allowed count $m\ge1$, the interior
conditional density given
$(\eta,m,\mathcal R_B)$ is
\begin{equation}\label{eq:gaussian-free-green-kernel}
 \frac1{\mathcal Z_{B,m}(\mathcal R_B)}
 \prod_{i<j}\delta(x_i,x_j)^\beta
 \prod_{i=1}^m
    \frac{\exp\{\beta\operatorname{Re}V_{\mathcal R_B}(x_i)\}}
         {(1-|x_i|^2)^{\beta/2}}\,dA^m(\mathbf x).
\end{equation}
The normalizing constant is strictly positive and finite. Resample
this tuple and update the meromorphic field by
\begin{equation}\label{eq:gaussian-free-A-resampling}
 \mathcal A^{\rm new}(z)=\mathcal A(z)
   +\sum_{i=1}^m J_{x_i}(z)-\sum_{p\in\Xi\cap B}J_p(z).
\end{equation}
For $m=0$ use the identity kernel. These kernels preserve the
augmented law of $(\Xi,\mathcal A)$ and form a consistent canonical
specification on the chosen nested disks.

Equivalently, their ordinary-area densities are proportional to
\begin{equation}\label{eq:gaussian-elimination-moment-norm}
 |\Delta_m(\mathbf x)|^\beta
 \exp\left\{\beta\operatorname{Re}\sum_i V_{\mathcal R_B}(x_i)
                  +\frac\beta2\|s(\mathbf x)\|_*^2\right\}.
\end{equation}
No hypothesis about convergence of an inverse exterior field is used
in this elimination. The hypothesis of independence of the corrected
field is essential; independence of $\Xi$ and $G$ alone would not
justify conditioning additionally on $\mathcal R_B$.
\end{theorem}

\begin{proof}
On the event $\Xi(B)=m$, use the auxiliary ordering independent of
all other variables and put
$Z=(\eta,m,\mathcal R_B)$. This is a measurable function of
$(\Xi,\mathcal A)$ and the fixed count. The asserted Gaussian
independence therefore remains valid under this count conditioning,
and Lemma~\ref{lem:infinite-gaussian-bayes-tuples} applies.

For almost every $Z$ and almost every $h$ in its conditional law,
\eqref{eq:gaussian-bayes-tuples} agrees with
\eqref{eq:canonical-density-before-gaussian-elimination}. The latter
is a strictly positive density off the diagonals. Since
$r(\mathbf x,h)$ is positive and bounded above and below, it follows
first that $\Pi_Z$ is absolutely continuous and then that its density
is proportional to that density divided by
$r(\mathbf x,h)$. Substitution cancels every term containing $h$.
The terms containing only the reference tuple $\mathbf a$ are part
of the normalizing constant. What remains is exactly
\eqref{eq:gaussian-elimination-moment-norm}. This also shows directly
that the resulting density does not depend on the selected good value
of $h$.

The elementary series $\sum_{k\ge1}u^k/k=-\log(1-u)$ gives
\[
 \|s(\mathbf x)\|_*^2
 =-\sum_i\log(1-|x_i|^2)
                   -2\sum_{i<j}\log|1-x_i\overline{x_j}|.
\]
This proves \eqref{eq:gaussian-free-green-kernel}. The primitive
$V_{\mathcal R_B}$ is bounded on $\overline B$, as is
$(1-|x|^2)^{-\beta/2}$, while $0\le\delta(x,y)\le1$ on the disk.
Thus the normalizing integral is finite. It is positive because the
integrand is positive on the open set of distinct tuples.

Given the exterior, its count, $\mathcal R_B$ and a resampled tuple,
the global meromorphic field is uniquely determined by its restriction
$\mathcal R_B+\sum_iJ_{x_i}$ to $B$. Its continuation exists by
\eqref{eq:gaussian-free-A-resampling}, and is measurable by the same
Borel restriction argument used for the original meromorphic
augmentation. This proves that the displayed kernel is the
conditional law of the entire augmented state and hence preserves it.

For nested disks $B_j\subset B_k$ one has
\[
 \mathcal R_{B_j}(z)=\mathcal R_{B_k}(z)
                     +\sum_{p\in\Xi\cap(B_k\setminus B_j)}J_p(z).
\]
A move in $B_j$ leaves the exterior, total count, and field
$\mathcal R_{B_k}$ unchanged, by
\eqref{eq:gaussian-free-A-resampling}. Conversely, conditioning the
$B_k$ kernel on its annular points factors its Green pair products
into the inner pair product and the cross factors
$\prod_{i,p}\delta(x_i,p)^\beta$. Each cross factor is the exponential
of the real part of a primitive of $\beta J_p$ on $B_j$.
The fixed one-body weight $(1-|x|^2)^{-\beta/2}$ is unchanged.
Therefore the inner conditional density is precisely the $B_j$
kernel. These two facts prove both orders of canonical consistency,
including zero-count cases as in
Theorem~\ref{thm:consistent-augmented-canonical-specification}.
\end{proof}

\begin{proposition}[Isometry covariance of the Gaussian-free kernels]
\label{prop:gaussian-free-green-covariance}
For this covariance statement, use the ambient admissible state space
in which an empty configuration may carry any holomorphic field
$\mathcal A$. The law constructed above is supported on the particular
choice $\mathcal A=0$ on its empty event. That convention need not be
preserved by the following action. For a holomorphic automorphism
$\varphi$ of $\D$, transform an augmented state by
\begin{equation}\label{eq:gaussian-free-affine-connection}
 \Xi^\varphi=\varphi^{-1}\Xi,\qquad
 \mathcal A^\varphi(z)
  =\varphi'(z)\mathcal A(\varphi(z))
       +\left(\frac2\beta-\frac12\right)
                              \frac{\varphi''(z)}{\varphi'(z)}.
\end{equation}
The kernels \eqref{eq:gaussian-free-green-kernel} transform covariantly
between a relatively compact disk and its image or preimage under
$\varphi$. The statement applies to any such disk for which the
kernel is defined by its exterior field, regardless of the particular
nested exhaustion used above. It is covariance of the specification;
it does not assert invariance of its Gibbs laws. This affine action is
chosen to compensate the one-body reference measure in the kernel. It
is distinct from the pure one-form identity obeyed by the finite
physical sum $\mathcal A_N=\sum_pJ_p$ under direct transport of all
its particles. No identification of those two actions is asserted.
\end{proposition}

\begin{proof}
Hyperbolic distance invariance gives
\[
 \varphi'(z)J_{\varphi(p)}(\varphi(z))=J_p(z).
\]
Consequently the exterior field transforms by the same affine formula
as \eqref{eq:gaussian-free-affine-connection}. A primitive of the
additional term is
$(2/\beta-1/2)\log\varphi'$, up to an irrelevant constant.
The identity
$1-|\varphi(z)|^2=|\varphi'(z)|(1-|z|^2)$ shows that the one-body
area measure in the kernel transforms as
\[
 \frac{dA(\varphi(z))}{(1-|\varphi(z)|^2)^{\beta/2}}
 =|\varphi'(z)|^{2-\beta/2}
                      \frac{dA(z)}{(1-|z|^2)^{\beta/2}}.
\]
The logarithm of this extra factor is exactly $\beta$ times the
real part of the additional primitive. The pair factors are invariant,
and normalization then proves the claim. Reflection is handled by
complex conjugation, so the same conclusion holds for the full
isometry group with the corresponding conjugated field.
\end{proof}

\begin{remark}[The remaining information in the field]
Relative to hyperbolic area $d\lambda=(1-|z|^2)^{-2}dA$, the one-body
factor in \eqref{eq:gaussian-free-green-kernel} is
$(1-|z|^2)^{2-\beta/2}\exp\{\beta\operatorname{Re}V_{\mathcal R_B}(z)\}$.
Its deterministic curvature is the one associated with a local
neutralizing density $(4/\beta-1)/\pi$ relative to $\lambda$.
This identifies the background coefficient in the augmented equations.
It neither proves that this coefficient is the particle intensity nor
identifies $\mathcal R_B$ as a compensated sum determined by the
unmarked exterior. Such an identification, or a uniqueness theorem for
these augmented equations with specified boundary behavior and count
weights, would be an additional result.
\end{remark}

\begin{corollary}[Positive-kernel factorization of the unmarked conditional law]
\label{cor:gaussian-divisible-conditional-kernel}
For almost every exterior $\eta$ and every allowed count $m\ge1$, let
$g_{\eta,m}$ be the analytic factor in
Theorem~\ref{thm:unconditional-canonical-quasigibbs}. There is a positive
Hermitian kernel $\mathcal H_{\eta,m}(\mathbf x,\mathbf y)$, holomorphic
in $\mathbf x\in B^m$ and antiholomorphic in $\mathbf y\in B^m$, such that
\begin{equation}\label{eq:gaussian-divisible-diagonal-factor}
 g_{\eta,m}(\mathbf x)
 =\exp\left\{\frac\beta2\|s(\mathbf x)\|_*^2\right\}
                  \mathcal H_{\eta,m}(\mathbf x,\mathbf x).
\end{equation}
Equivalently, the polarized positive kernel with diagonal
$g_{\eta,m}$ factors as
\begin{equation}\label{eq:gaussian-divisible-kernel-factor}
 \mathcal G_{\eta,m}(\mathbf x,\mathbf y)
 =\exp\left\{\frac\beta2\sum_{k\ge1}
       \frac{s_k(\mathbf x)\overline{s_k(\mathbf y)}}k\right\}
                         \mathcal H_{\eta,m}(\mathbf x,\mathbf y).
\end{equation}
In particular the quotient by the displayed Gaussian kernel remains
positive as a kernel, a stronger assertion than positivity of its
logarithmic Levi form.
\end{corollary}

\begin{proof}
Disintegrate the conditional Green kernels over the law
$\mu_{\eta,m}$ of $\mathcal R_B$ given $(\eta,m)$. Put
\[
 F_R(\mathbf x)=\mathcal Z_{B,m}(R)^{-1/2}
             \exp\left\{\frac\beta2\sum_iV_R(x_i)\right\},\qquad
 \mathcal H_{\eta,m}(\mathbf x,\mathbf y)
      =\int F_R(\mathbf x)\overline{F_R(\mathbf y)}
                                                \,d\mu_{\eta,m}(R).
\]
The diagonal identity follows from
\eqref{eq:gaussian-elimination-moment-norm}, initially almost everywhere.
The deterministic compact upper bounds for $g_{\eta,m}$ from the
canonical quasi-Gibbs theorem and the positivity of the Gaussian factor
bound these diagonal integrals locally almost everywhere. The
mean-value inequality for each $|F_R|^2$, integrated in $R$ and then
over a slightly larger polydisk, upgrades this to a bound at every
point of a smaller polydisk. Cauchy--Schwarz bounds the
corresponding off-diagonal integrals. Local integration of holomorphic
functions, or the mean-value formula followed by Fubini on smaller
polydisks, proves joint holomorphic--antiholomorphic regularity. It also
extends the diagonal identity to every tuple, including the diagonals.
Positivity of $\mathcal H$ follows directly by integrating the square
of any finite linear combination of the $F_R$. The exponential kernel
in \eqref{eq:gaussian-divisible-kernel-factor} has diagonal equal to the
Gaussian factor in \eqref{eq:gaussian-divisible-diagonal-factor}; its
positivity follows from its power-series expansion as an exponential
of an inner-product kernel. The product therefore has diagonal
$g_{\eta,m}$ and is positive. Uniqueness of polarization of a real
analytic diagonal gives \eqref{eq:gaussian-divisible-kernel-factor}.
\end{proof}

\begin{corollary}[A conditionally independent Gaussian component of every force]
\label{cor:conditional-force-gaussian-convolution}
Given the unmarked configuration, the joint force law at any finite
list of distinct particles is the convolution of a probability measure
and a nondegenerate circular complex Gaussian law. More precisely, if
$\mathbf z=(z_1,\ldots,z_m)$ are distinct particles and
$A_i$ is the regular part of $\mathcal A$ at $z_i$, then
\[
 (C_{z_i})_{i=1}^m=(A_i)_{i=1}^m-(T_G(z_i))_{i=1}^m,
\]
and the two vectors on the right are conditionally independent given
$\Xi$. The Gaussian covariance matrix is
\[
 \Sigma_{ij}=\frac2\beta\frac1{(1-z_i\overline{z_j})^2}.
\]
Thus this conditional force law has a strictly positive smooth density
on $\C^m$, bounded above by $(\pi^m\det\Sigma)^{-1}$. Its
characteristic function satisfies
\[
 \left|\mathbb E\left[
  e^{\,i\operatorname{Re}\sum_i\overline{t_i}C_{z_i}}\mid\Xi\right]\right|
 \le \exp\left\{-\frac1{2\beta}
           \sum_{i,j}\overline{t_i}
                     \frac{t_j}{(1-z_i\overline{z_j})^2}\right\}.
\]
Statements at selected particles may be interpreted under finite
Campbell windows, or using any measurable ordering of the particles.
\end{corollary}

\begin{proof}
Adding the holomorphic function $T_G$ changes the regular part at a
pole by its value at that pole, so the displayed decomposition follows
from the definition of $\mathcal A$. The independence of
$(\Xi,\mathcal A)$ and $G$ proves the conditional independence.
The covariance follows by summing
$\sum_{k\ge1}k(z_i\overline{z_j})^{k-1}$.
For distinct points this matrix is positive definite: a vanishing
quadratic form would give
$\sum_i\overline{t_i}z_i^{k-1}=0$ for every $k$, and the first $m$
equations form an invertible Vandermonde system.

The Gaussian density is everywhere positive and bounded by its value
at zero. Integrating its translates gives the asserted positive bounded
density. Each derivative of a nondegenerate Gaussian density is bounded,
so differentiation under the probability integral proves smoothness of
all orders. Finally a circular complex Gaussian with covariance
$\Sigma$ has real-linear characteristic function
$\exp(-t^*\Sigma t/4)$. The other independent factor has modulus at most
one, yielding the stated bound.
\end{proof}

\section{Why the corrected canonical equations still need a state-selection condition}
\label{sec:finite-green-gibbs-obstruction}

The Gaussian elimination gives substantially more information than a
spatial Ward equation. Nevertheless, the resulting canonical equations
alone cannot imply either uniqueness or hyperbolic invariance. The
following explicit family demonstrates this point for $0<\beta<2$.
Each fixed-$m$ law below is excluded as an original-gas limit by
Theorem~\ref{thm:nonemptiness-below-two}: that theorem gives infinite
expected total count, whereas the example has total count $m$.

\begin{proposition}[Distinct finite-particle solutions of the corrected equations]
\label{prop:finite-green-canonical-solutions}
Suppose $0<\beta<2$, and let $m\ge1$ be an integer. There is a
rotationally invariant probability law $P^{[m]}_\beta$ with exactly $m$
particles in $\D$, whose uniformly ordered tuple has density
\begin{equation}\label{eq:finite-green-canonical-density}
 \frac1{Z_{\beta,m}}
 \prod_{1\le i<j\le m}\delta(z_i,z_j)^\beta
 \prod_{i=1}^m(1-|z_i|^2)^{-\beta/2}\,dA^m(\mathbf z).
\end{equation}
Augment it by
\[
 \mathcal A(z)=\sum_{p\in\Xi}J_p(z).
\]
This augmented law satisfies all the Gaussian-free canonical kernels
of Theorem~\ref{thm:gaussian-free-green-specification}, on every disk
$B\Subset\D$. These laws are distinct for distinct $m$. No one of
them is invariant under all disk automorphisms.

More generally, one may replace the density in
\eqref{eq:finite-green-canonical-density} by its product with
\[
 \exp\Bigl\{\beta\operatorname{Re}\sum_i h(z_i)\Bigr\},
\]
where $h$ is any holomorphic function on a neighborhood of
$\overline\D$, and set
\[
 \mathcal A=h'+\sum_{p\in\Xi}J_p.
\]
The same canonical assertion holds.
Arbitrary probability mixtures over the integer $m\ge0$ of the laws
with $h=0$, with the empty state at $m=0$, satisfy these equations as
well.
\end{proposition}

\begin{proof}
Put $b=\beta/2<1$. Since $0\le\delta(z,w)\le1$,
\[
 Z_{\beta,m}\le
 \left(\int_\D(1-|z|^2)^{-b}\,dA(z)\right)^m
 =\left(\frac\pi{1-b}\right)^m<\infty.
\]
The integral is strictly positive, because its integrand is positive
on the nonempty open set of distinct tuples. The density is symmetric,
absolutely continuous, and invariant under simultaneous rotations.
It therefore defines the stated simple point process.

Fix a disk $B\Subset\D$. Its boundary has no particles almost surely.
Given its exterior $\eta$ and its inside count $n$, factor the finite
pair product in \eqref{eq:finite-green-canonical-density}. The
conditional density of its inside tuple $\mathbf x$ is proportional to
\[
 \prod_{i<j}\delta(x_i,x_j)^\beta
 \prod_i(1-|x_i|^2)^{-b}
 \prod_{\substack{i\le n\\p\in\eta}}\delta(x_i,p)^\beta.
\]
The retained exterior field is
$\mathcal R_B=\sum_{p\in\eta}J_p$.
Each function $\log\delta(z,p)$ is harmonic on $B$, and
$\partial_z\log\delta(z,p)^2=J_p(z)$. Hence the last product is,
up to a constant independent of $\mathbf x$, the exponential of
$\beta\operatorname{Re}\sum_i V_{\mathcal R_B}(x_i)$.
This is exactly \eqref{eq:gaussian-free-green-kernel}.
Replacing the inside tuple changes $\mathcal A$ by precisely
\eqref{eq:gaussian-free-A-resampling}. Thus the assertion holds for
the entire augmented state, not only for its position marginal.
The argument applies without alteration when the additional field
$h'$ is retained. That density remains normalizable because
$\operatorname{Re}h$ is bounded on $\overline\D$.

The laws for different total counts have disjoint supports. A finite
nonzero invariant point process cannot exist here: its intensity
would be a nonzero finite invariant measure on $\D$. In the present
case this measure has a continuous density, and invariance would
force that density to be $c(1-|z|^2)^{-2}$ for some $c>0$ by the
change of variables under an automorphism sending zero to a specified
point. Such a measure has infinite mass. This contradicts its total
mass $m$. Finally, the specification preserves the total count and is
the same kernel on each of these invariant sectors, so integration of
the stationarity identities against arbitrary mixture weights proves
the mixture assertion.
\end{proof}

\begin{example}[The one-particle solution saturates the proved curvature bound]
For $m=1$, the density in the proposition is explicitly
\[
 \rho^{[1]}_\beta(z)=\frac{1-\beta/2}{\pi}
                         (1-|z|^2)^{-\beta/2}.
\]
It satisfies
\[
 \Delta\log\rho^{[1]}_\beta(z)
           =\frac{2\beta}{(1-|z|^2)^2},
\]
so the unconditional Gaussian curvature inequality alone does not
exclude it. Its intensity is not proportional to hyperbolic area.
The example shows concretely that the remaining state-selection issue
is not merely a missing formulation of the spatial equations.
\end{example}

\begin{remark}[Adding back the independent Gaussian does not remove the examples]
Attach to any of the preceding laws an independent Gaussian sequence
with $\mathbb E|G_k|^2=2k/\beta$, and set
$\mathcal C=\mathcal A-T_G$. This gives exactly the asserted independence
of $(\Xi,\mathcal A)$ from $G$.
For completeness, these augmented laws also satisfy the earlier
Euclidean harmonic canonical equations. On a disk $B$, its exterior
fixes $\mathcal R_B$, and write $H=G-s(\Xi\cap B)$. Conditional on the
inside tuple $\mathbf x$, the translated Gaussian density in $H$
contributes the factor
\[
 \exp\left\{-\beta\operatorname{Re}\sum_{k\ge1}
                 \frac{\overline{H_k}s_k(\mathbf x)}k
                  -\frac\beta2\|s(\mathbf x)\|_*^2\right\},
\]
up to a factor independent of $\mathbf x$; this is the Gaussian Bayes
formula of Lemma~\ref{lem:infinite-gaussian-bayes-tuples}.
The norm term cancels the positive norm term in
\eqref{eq:gaussian-elimination-moment-norm}. The resulting density is
$|\Delta_n(\mathbf x)|^\beta$ times the exponential of a harmonic
one-body potential with derivative
$\mathcal E_B=\mathcal R_B-T_H$.
The field $H$ is determined by $\mathcal E_B$ and the unmarked exterior
in these examples: $\mathcal R_B$ is already determined by that
exterior, and the restriction of $T_H$ to $B$ uniquely determines its
holomorphic continuation to $\D$, whose Taylor coefficients determine
$H$. This continuation is measurable by the Borel restriction argument
used for the meromorphic augmentation.
Thus conditioning on the original exterior-field mark gives exactly
the original canonical kernel. The count probabilities and the
admissible boundary behavior of the corrected field remain additional
information to be selected by the finite-$N$ model.
\end{remark}

\begin{corollary}[All these finite-count mixtures are excluded as gas limits]
\label{cor:finite-green-mixtures-excluded}
For $0<\beta<2$, no probability mixture
$\sum_{m\ge0}p_mP^{[m]}_\beta$ of the preceding finite-particle laws
can be a subsequential limit of the original disk gas, even if the
mixture has infinite expected total count.
\end{corollary}

\begin{proof}
Every such mixture has finite total count almost surely, whereas
Theorem~\ref{thm:almost-sure-infinite-below-two} proves that every
actual subsequential limit has infinite total count almost surely.
\end{proof}

\section{Infinitely many infinite solutions of the canonical Green equations}
\label{sec:infinite-green-gibbs-nonuniqueness}

The finite-count examples do not address the almost sure infinitude of
actual limits for $0<\beta<2$.  The following construction
does.  For every $0<\beta<4$ it produces distinct, almost surely infinite,
rotationally invariant position laws satisfying exactly the same augmented
canonical Green equations as the disk-gas limits.  These are counterexamples
to uniqueness of that specification, not counterexamples to uniqueness of
the original finite-$N$ limit.  The almost sure non-Blaschke property of actual limits will in fact make
every new example mutually singular with every original-gas limit when
$0<\beta\le2$.

Throughout this section put
\[
 d(z)=1-|z|^2,\qquad
 g(z,w)=-\log\delta(z,w),\qquad b=\beta/2.
\]
A grand-canonical kernel permits both insertion and deletion of particles;
its additional parameter $\lambda>0$ is the activity.  Conditioning such a
kernel on its inside count removes this parameter.

\subsection{A construction with summable Green tails}

\begin{lemma}[Existence with a summable boundary weight]
\label{lem:green-gibbs-summable-activity}
Let $\beta>0$, and let $w$ be a strictly positive continuous function on
$\D$ such that
\begin{equation}\label{eq:green-gibbs-weight-hypotheses}
 \int_\D d(z)w(z)\,dA(z)<\infty.
\end{equation}
For every $\lambda>0$ there exists a simple locally finite point process
$P_{\lambda,w}$ satisfying the following assertions.
\begin{enumerate}
\item Its factorial moment densities satisfy, almost everywhere,
\begin{equation}\label{eq:green-gibbs-ruelle-bound}
 \rho_p(z_1,\ldots,z_p)\le
       \lambda^p\prod_{i=1}^p w(z_i),\qquad p\ge1.
\end{equation}
In particular,
\begin{equation}\label{eq:green-gibbs-finite-depth}
 \mathbb E\sum_{p\in\Xi}d(p)
       \le\lambda\int_\D d(z)w(z)\,dA(z)<\infty.
\end{equation}
\item For every disk $B\Subset\D$, given $\eta=\Xi\cap B^c$, its
interior grand-canonical density, with the usual unordered reference
$\sum_{m\ge0}dA^m/m!$, is
\begin{equation}\label{eq:green-grand-canonical-density}
 \frac1{Z_B(\eta)}\lambda^m
 \prod_{i=1}^m w(x_i)
 \exp\left\{-\beta\sum_{i<j}g(x_i,x_j)
             -\beta\sum_{i=1}^m\sum_{p\in\eta}g(x_i,p)\right\}.
\end{equation}
The exterior sums are finite, uniformly on $\overline B$, almost surely,
and $1\le Z_B(\eta)\le\exp\{\lambda\int_Bw\,dA\}$.
\item Its insertion identity is
\begin{equation}\label{eq:green-grand-gnz}
 \mathbb E\sum_{x\in\Xi} F(x,\Xi\setminus\{x\})
 =\lambda\mathbb E\int_\D F(x,\Xi)
             e^{-\beta\sum_{p\in\Xi}g(x,p)}w(x)\,dA(x)
\end{equation}
for every nonnegative measurable $F$, with equality in $[0,\infty]$.
\item If $\int_\D w\,dA=\infty$, then $\#\Xi=\infty$ almost surely.
For a fixed $w$, states obtained with different activities have distinct
position laws, regardless of the choices made in taking limits.
\end{enumerate}
The process may be obtained as a subsequential limit of the free
 grand-canonical laws on concentric disks.  Consequently it may be chosen
rotationally invariant whenever $w$ is rotationally invariant.
\end{lemma}

\begin{proof}
Write $\sigma(dz)=w(z)dA(z)$ and $D_R=\{|z|<R\}$.  On $D_R$ use the
probability density proportional to
\[
 \lambda^n e^{-\beta\sum_{i<j}g(z_i,z_j)}\prod_iw(z_i)
\]
relative to $\sum_{n\ge0}dA^n/n!$.  The partition function lies between
one and $e^{\lambda\sigma(D_R)}$, since $g\ge0$.  Its ordered factorial
density is
\[
 \lambda^p\prod_iw(z_i)e^{-\beta\sum_{i<j}g(z_i,z_j)}
 \frac{Z_R[z_1,\ldots,z_p]}{Z_R},
\]
where the numerator inserts into every remaining-particle integrand the
factor
\[
 \exp\{-\beta\sum_{i,j}g(z_i,y_j)\}\le1.
\]
This proves
\eqref{eq:green-gibbs-ruelle-bound} uniformly in $R$.  These bounds imply
local tightness: for every compact $K$ and every $t\ge0$,
\[
 \mathbb E(1+t)^{\Xi_R(K)}
 =\sum_{p\ge0}\frac{t^p}{p!}\int_{K^p}\rho_{p,R}\,dA^p
 \le e^{\lambda t\sigma(K)}.
\]
The second factorial bound excludes collisions in every limiting process:
on any compact set its integral over ordered pairs at distance less than
$\varepsilon$ tends to zero as $\varepsilon\downarrow0$.  The first bound implies that deterministic area-null sets contain no
limiting particle.  A diagonal extraction therefore gives a simple locally
finite limit.  Testing factorial measures on compact sets and using the
exponential count bound passes \eqref{eq:green-gibbs-ruelle-bound} to it.
This also proves \eqref{eq:green-gibbs-finite-depth} and the bounds
\begin{equation}\label{eq:green-gibbs-uniform-tail}
 \sup_R\mathbb E\sum_{p\in\Xi_R\setminus D_s}d(p)
 \le\lambda\int_{\D\setminus D_s}d\,d\sigma
       \longrightarrow0\quad(s\uparrow1),
\end{equation}
with the same inequality for the limit.

We give the details needed to pass the conditional kernels, because the
infinite mass of $\sigma$ prevents replacing this step by an assertion of
global Poisson finiteness.  For each compact $K\Subset\D$ there are
$s_K<1$ and $C_K<\infty$ for which
\begin{equation}\label{eq:green-gibbs-green-tail-bound}
 \sup_{z\in K}g(z,p)\le C_Kd(p),\qquad |p|>s_K.
\end{equation}
Indeed,
\[
 1-\delta(z,p)^2
   =\frac{d(z)d(p)}{|1-z\overline p|^2}\le C'_Kd(p),
\]
and $-\log(1-u)\le2u$ for $0\le u\le1/2$.  The identity
$g=-\frac12\log\delta^2$ proves the claim.  Thus the total contribution
of points outside $D_s$ to each exterior potential on $K$ tends to zero
in expectation, uniformly in the finite-volume index, by
\eqref{eq:green-gibbs-uniform-tail}.

One can consequently carry out the extraction in a topology that records
both vague convergence of configurations and their finite weighted measures
$\sum_pd(p)\delta_p$ on the closed disk.  To justify this without any
additional compactness assumption, these finite measures have uniformly
bounded total mass; their restrictions outside $D_s$ have uniformly
vanishing expected mass by \eqref{eq:green-gibbs-uniform-tail}.  Tightness
of finite measures on the compact closed disk and local configuration
tightness give a joint extraction.  The first moment tail bound forces
the limiting weighted measure to put no mass on the boundary and to equal
$\sum_pd(p)\delta_p$ in the interior.  Along a Skorokhod representation
of this joint extraction, convergence of weighted total masses and of
weighted measures on each $\overline D_s$ whose boundary is empty implies
\[
 \lim_{s\uparrow1}\limsup_{R\uparrow1}
       \sum_{p\in\Xi_R\setminus D_s}d(p)=0
 \quad\hbox{almost surely}.
\]
Here and below $R\uparrow1$ refers to the extracted sequence.

Fix $B\Subset\D$.  Its boundary is empty almost surely.  There are only
finitely many exterior points in any larger compact disk, and their
positive distance from $\overline B$ is positive.  Matching these points
under vague convergence and using
\eqref{eq:green-gibbs-green-tail-bound} shows that their exterior potentials
converge uniformly on $\overline B$.  Each fixed-$m$ density in
\eqref{eq:green-grand-canonical-density} therefore converges almost
everywhere.  It is bounded by $\lambda^m\prod_iw(x_i)$, whose summed
integrals divided by $m!$ equal $e^{\lambda\sigma(B)}$.  Dominated
convergence gives convergence of the normalizing constants and of all
bounded interior test integrals.  The original finite-volume conditional
identity can now be multiplied by a bounded continuous exterior cylinder
function and passed to the limit.  A monotone-class argument proves the
conditional identity for the full exterior sigma field.  This works for
every specified disk $B$; a countable determining family can be handled
simultaneously.  It proves the second assertion, including uniform
finiteness of its exterior sums.

For completeness, conditional on the exterior of a disk, shift the
summation index $m$ in the left side of
\eqref{eq:green-grand-gnz} and single out one of its $m$ coordinates.
The factor $m/m!$ becomes $1/(m-1)!$, and the singled-out particle
contributes precisely
$\lambda w(x)\exp\{-\beta\sum_pg(x,p)\}$.  This proves the insertion
identity for nonnegative $F$ supported in that disk in its first
coordinate.  Monotone convergence over disks proves the full identity.

Suppose now $\sigma(\D)=\infty$.  For each deterministic finite
configuration $\eta$, the function
\[
 a_\eta(x)=\prod_{p\in\eta}\delta(x,p)^\beta
\]
tends uniformly to one as $|x|\uparrow1$.  Since $\sigma$ is finite on
compact subsets, $\int a_\eta\,d\sigma=\infty$.  Apply
\eqref{eq:green-grand-gnz} with
$F(x,\eta)=\mathbf1_{\{\#\eta=m\}}$.  Its left side is
$(m+1)\mathbb P(\#\Xi=m+1)$ and is finite; its right side is infinite
if $\mathbb P(\#\Xi=m)>0$.  Therefore
$\mathbb P(\#\Xi=m)=0$ for every integer $m\ge0$.

Finally take a nonempty relatively compact disk $B$.  Its insertion
integral
\[
 I_w(\Xi)=\int_B e^{-\beta\sum_pg(x,p)}w(x)\,dA(x)
\]
is bounded by $\sigma(B)$ and is strictly positive almost surely.
The strict positivity follows from finite weighted depth, the Green tail
bound, and local finiteness: the potential is finite at every $x$ except
possibly the countable set of particles.  The identity gives
$\mathbb E\Xi(B)=\lambda\mathbb EI_w(\Xi)$.  Both expectations are
functionals of the position law, and the second is strictly positive.
The same position law therefore cannot satisfy the identity with two
different activities.  Rotation invariance, when available in finite
volume, passes to the extracted law.
\end{proof}

\begin{theorem}[Distinct infinite canonical states at every $0<\beta<4$]
\label{thm:infinite-green-canonical-nonuniqueness}
For each $0<\beta<4$ there are infinitely many distinct, almost surely
infinite, rotationally invariant position laws which admit an augmentation
satisfying the Gaussian-free canonical specification
\eqref{eq:gaussian-free-green-kernel}--
\eqref{eq:gaussian-free-A-resampling}.  Their corrected meromorphic fields
can be taken in the particularly restrictive form
\begin{equation}\label{eq:infinite-green-absolutely-summed-field}
 \mathcal A(z)=h'(z)+\sum_{p\in\Xi}J_p(z),
 \qquad \sum_{p\in\Xi}d(p)<\infty\quad\hbox{almost surely},
\end{equation}
where the sum converges absolutely and locally uniformly off the particles.
For $2\le\beta<4$ one may take $h=0$.  For $0<\beta<2$ one may take a
single random boundary point $\xi$, uniformly distributed on the unit
circle, and
\begin{equation}\label{eq:infinite-green-boundary-pole}
 h_\xi(z)=-\frac{2-\beta/2}{\beta}\log(1-\overline\xi z),
 \qquad
 h_\xi'(z)=\frac{2-\beta/2}{\beta}
                   \frac{\overline\xi}{1-\overline\xi z}.
\end{equation}
The branch of the logarithm vanishes at zero.  In the latter construction,
$\xi$ is almost surely the unique accumulation point of $\Xi$ on the
boundary; in particular it is determined by the positions.

None of these rotationally invariant states is invariant under the full
group of disk automorphisms.  For $0<\beta\le2$, each of them is mutually singular with every
subsequential limit of the original unweighted disk gas.
\end{theorem}

\begin{proof}
For a holomorphic function $h$ on $\D$ set
\[
 w_h(z)=d(z)^{-b}e^{\beta\operatorname{Re}h(z)}.
\]
For $1\le b<2$, $h=0$ gives
\[
 \int_\D w_h\,dA=\infty,
 \qquad
 \int_\D d\,w_h\,dA=\frac\pi{2-b}<\infty.
\]
The preceding lemma therefore supplies infinite rotationally invariant
states, distinct for different $\lambda$.

Suppose $0<b<1$, put $s=2-b\in(1,2)$, and first fix $\xi=1$.  The
weight is $w_h(z)=d(z)^{-b}|1-z|^{-s}$.  It is integrable outside every
neighborhood of $1$.  Near $1$ use $z=(1-t)e^{i\theta}$, with $t>0$
and $|\theta|$ small.  Uniform comparability of the corresponding
positive quantities gives
\[
 d(z)\asymp t,\qquad |1-z|^2\asymp t^2+\theta^2,
 \qquad dA(z)\asymp dt\,d\theta.
\]
Since $s>1$, for fixed small $\varepsilon>0$ and $0<t<\varepsilon/2$,
\[
 \int_{-\varepsilon}^{\varepsilon}
          (t^2+\theta^2)^{-s/2}\,d\theta\asymp t^{1-s};
\]
substitution $\theta=tu$ proves both bounds by comparison with the
positive finite integral over the real line.  Consequently the local
integral of $d^qw_h$ has the same finiteness as
\[
 \int_0^{\varepsilon/2} t^{q-b+1-s}\,dt
       =\int_0^{\varepsilon/2}t^{q-1}\,dt.
\]
It diverges for $q=0$ and converges for $q=1$.  The weight thus meets
all hypotheses of the lemma.  Its point process has infinitely many
points, but only finitely many outside every neighborhood of $1$, by
the first moment bound.  It follows that $1$ is its unique boundary
accumulation point.  This point is a measurable function of the process.
Indeed, for a fixed sequence $r_n\uparrow1$, on this event it equals
\[
 \lim_{n\to\infty}
   \frac{\sum_{p\in\Xi\cap D_{r_n}}p}{\Xi(D_{r_n})};
\]
the denominators tend to infinity, and all but finitely many summands lie
in any prescribed neighborhood of the unique accumulation point.  Terms
with zero denominator can be assigned the value zero.
Rotate a chosen such law by $\xi$ and average over uniform $\xi$.  This produces a rotationally invariant law, with the
same finite expected weighted depth and the asserted measurable boundary
parameter.  The law at $\xi=1$ may also be chosen invariant under complex
conjugation, since all its finite-volume laws have that symmetry; the
rotational mixture is then invariant under reflection as well.

To verify the canonical equations in all cases, note the algebraic bound
\[
 |J_p(z)|
 =\frac{d(p)}{|z-p|\,|1-z\overline p|}
 \le C_Kd(p),\qquad z\in K,\quad |p|>s_K,
\]
for every compact $K\Subset\D$ and suitable $s_K<1$.
The finite expected depth proves absolute local uniform convergence of
\eqref{eq:infinite-green-absolutely-summed-field}.  The series is
meromorphic with a simple pole of residue one at each particle.  Its
retained exterior field is
\[
 \mathcal R_B=h'+\sum_{p\in\eta}J_p.
\]
The grand-canonical density of the preceding lemma, conditioned on its
count $m$, is
\[
 \prod_{i<j}\delta(x_i,x_j)^\beta
 \prod_i d(x_i)^{-b}
 \exp\left\{\beta\operatorname{Re}\sum_i h(x_i)
                   -\beta\sum_{i,p\in\eta}g(x_i,p)\right\}.
\]
As $J_p$ is the holomorphic derivative corresponding to
$\log\delta(\,\cdot\,,p)$, this is precisely the required kernel with
field $\mathcal R_B$, up to a constant independent of the inside tuple.
Replacing that tuple updates $\mathcal A$ by the required finite sum.
For fixed $h$ the exterior field is already a function of the exterior.
In the rotational mixture, $h_\xi$ is also a function of the exterior,
because deleting finitely many particles leaves the unique boundary
accumulation point unchanged.  Thus conditioning additionally on the
field does not change the displayed conditional density.

For the rotational mixture with $b<1$, distinctness still holds at the
level of position laws.  Its insertion identity is
\[
 \mathbb E\sum_{x\in\Xi}F(x,\Xi\setminus\{x\})
 =\lambda\mathbb E\int F(x,\Xi)e^{-\beta\sum_pg(x,p)}
                     w_{h_{\xi(\Xi)}}(x)\,dA(x).
\]
This follows by averaging the identities at fixed $\xi$; finite insertion
or deletion does not change $\xi(\Xi)$.  On a fixed compact disk the
integral on the right, without $\lambda$, is positive and bounded,
uniformly in $\xi$.  The same argument as in the lemma identifies
$\lambda$ from the position law.  This proves distinctness for the
rotational mixtures too.

All constructed laws have finite expected weighted depth.  If one were
invariant under every disk automorphism, its locally finite intensity
would be $c\,dA/d^2$, by transitivity and change of variables.  The
insertion identity makes this measure nonzero, so $c>0$.  But its
expected weighted depth would then be
$c\int_\D dA/d=\infty$, a contradiction.  Finally, Theorem~\ref{thm:non-blaschke-limits} proves that every actual
limit for $0<\beta<2$ has $\sum_pd(p)=\infty$ almost surely.  At
$\beta=2$ the unique original limit is Bergman; its full hyperbolic
invariance, almost sure nonemptiness, and
Theorem~\ref{thm:invariant-green-divergence} imply
$\sum_pg(0,p)=\infty$ almost surely.  Since $g(0,p)=-\log|p|$
is comparable to $d(p)$ outside any fixed compact disk, and local
finiteness accounts for the remaining finitely many points, Bergman
also has $\sum_pd(p)=\infty$ almost surely.  (The origin is almost
surely not a particle.)  All the constructed laws instead satisfy
$\sum_pd(p)<\infty$ almost surely.  These disjoint measurable events
prove the asserted mutual singularity.
\end{proof}

\subsection{An explicit determinantal family at \texorpdfstring{$\beta=2$}{beta=2}}

\begin{proposition}[Activity interpolates from dilute infinite states to Bergman]
\label{prop:explicit-szego-activity-states}
At $\beta=2$ and $h=0$ the preceding grand-canonical disk exhaustion has
a unique limit for each fixed activity $\lambda>0$.  It is the
determinantal process, relative to ordinary area, with kernel
\begin{equation}\label{eq:explicit-szego-activity-kernel}
 K_\lambda(z,w)=\sum_{n=0}^{\infty}
       \frac{\lambda(n+1)}{n+1+\pi\lambda}(z\overline w)^n.
\end{equation}
It satisfies the same augmented canonical Green equations as the Bergman
process, has infinitely many points almost surely, and satisfies
\[
 \rho_\lambda(0)=\frac{\lambda}{1+\pi\lambda},\qquad
 \rho_\lambda(z)\le\frac{\lambda}{1-|z|^2},\qquad
 \mathbb E\sum_{p\in\Xi}d(p)\le\pi\lambda.
\]
The parameter $\lambda$ therefore gives explicit distinct infinite
solutions of those equations even at the exactly solvable temperature.
As $\lambda\uparrow\infty$ these position laws converge locally to the
Bergman process.
\end{proposition}

\begin{proof}
The Cauchy determinant identity gives
\begin{equation}\label{eq:szego-cauchy-grand-identity}
 \det\left[\frac1{1-z_i\overline z_j}\right]_{i,j=1}^m
 =\prod_i d(z_i)^{-1}\prod_{i<j}\delta(z_i,z_j)^2.
\end{equation}
For example, multiplying the determinant by
$\prod_{i,j}(1-z_i\overline z_j)$ yields a polynomial alternating in
the $z_i$ and separately in their conjugates.  Its degrees force it to
be the product of the two Vandermondes; comparing a coefficient fixes
the sign and constant, and rearranging gives the displayed positive
identity.

On $D_R$ the grand-canonical density is consequently an $L$-ensemble
with $L(z,w)=\lambda/(1-z\overline w)$.  The functions
\[
 e_{n,R}(z)=\sqrt{\frac{n+1}{\pi R^{2n+2}}}\,z^n
\]
are orthonormal, and the nonzero eigenvalues of $L$ are
$\ell_{n,R}=\lambda\pi R^{2n+2}/(n+1)$.  They are summable when
$R<1$.  Expanding the determinant in this orthonormal basis and
integrating each term proves
\[
 Z_R=\prod_{n\ge0}(1+\ell_{n,R}).
\]
The same expansion, with a multiplicative test function inserted, or
with finitely many coordinates left unintegrated, shows that its
correlation kernel is $L(I+L)^{-1}$.  More explicitly,
\[
 K_{\lambda,R}(z,w)
 =\sum_{n\ge0}
       \frac{\lambda(n+1)}{n+1+\lambda\pi R^{2n+2}}
                       (z\overline w)^n.
\]
These calculations can first be performed with a finite orthonormal
sum; passage to the displayed formulas is justified by positivity of
the Cauchy--Binet expansions and the summability of the eigenvalues.

The coefficients are at most $\lambda$, so the kernels converge
uniformly on compact subsets of $\D\times\D$ to
\eqref{eq:explicit-szego-activity-kernel}.  All fixed correlation
functions converge locally uniformly.  Hadamard's inequality bounds
their absolute values on a compact set $K$ by
$[\lambda/(1-\sup_K|z|^2)]^p$.
The factorial expansion of each compactly supported Laplace functional
therefore converges absolutely, and proves uniqueness and convergence
of the position law.  It agrees with the state constructed above and
inherits its infinite total count and canonical equations.  Evaluation
at zero and summation of the coefficient bound give the asserted
intensity estimates and, by integration, the depth estimate.

Lastly, as $\lambda\uparrow\infty$ each coefficient tends to
$(n+1)/\pi$ and is bounded by that quantity.  Hence $K_\lambda$
converges locally uniformly to
\[
 K_{\rm Berg}(z,w)=\frac1{\pi(1-z\overline w)^2}.
\]
The same locally dominated Laplace-functional expansion proves local
convergence to the Bergman point process.
\end{proof}

\begin{remark}[What this establishes about nonuniqueness]
The canonical equations, almost sure infinitude, rotational invariance,
and even an absolutely convergent corrected exterior field do not select
a unique state.  This failure is proved here for every $0<\beta<4$.
At $\beta=2$ the explicit states also show why this failure is compatible
with convergence of the original model to a unique limit: the original
model selects Bergman, whereas every finite-activity state in
\eqref{eq:explicit-szego-activity-kernel} is mutually singular with Bergman.  The same almost sure boundary-mass
separation holds between every constructed state and every original-gas
limit for $0<\beta<2$.
No two different subsequential limits of the original unweighted disk gas
are constructed by this argument.  Moreover none of these counterstates
is fully hyperbolically invariant.  The question of uniqueness within a
correctly selected invariant Gibbs class remains separate.
\end{remark}

\section{Canonical tails, extremal states, and the precise nonuniqueness question}
\label{sec:canonical-tail-nonuniqueness}

There are three logically different questions: uniqueness of the
augmented canonical Gibbs state, uniqueness of its position marginal,
and uniqueness of the subsequential limit selected by the original
finite-particle sequence. The following results keep these questions
separate. They also give an exact test for whether a decomposition into
Gibbs states changes the position law.

Fix $\beta>0$ and nested disks $B_j\uparrow\D$. On the admissible
augmented state space of pairs $\omega=(\Xi,\mathcal A)$, write
$K_j$ for the Gaussian-free canonical kernel in
Theorem~\ref{thm:gaussian-free-green-specification}, and set
\[
 I_j(\omega)=\bigl(\Xi\vert_{B_j^c},\Xi(B_j),\mathcal R_{B_j}\bigr),
 \qquad \mathcal F_j=\sigma(I_j).
\]
We work on a common measurable state space on which these kernels and
invariants are defined. A law $P$ is called a Gibbs law for these
kernels if $K_jF$ is a version of $\E_P[F\mid\mathcal F_j]$ for
every bounded measurable $F$ and every $j$. The conditional-kernel
and consistency assertions already proved give this property to the
constructed augmented subsequential limits. Boundary-null assumptions
are understood for the chosen countable exhaustion.

\begin{lemma}[The canonical information decreases along an exhaustion]
\label{lem:canonical-decreasing-information}
For $j<k$, one has $\mathcal F_k\subseteq\mathcal F_j$.
\end{lemma}
\begin{proof}
The exterior of $B_j$ determines both the exterior of $B_k$ and
all the particles in $B_k\setminus B_j$. Moreover,
\[
 \Xi(B_k)=\Xi(B_j)+\Xi(B_k\setminus B_j),\qquad
 \mathcal R_{B_k}=\mathcal R_{B_j}
                  -\sum_{p\in\Xi\cap(B_k\setminus B_j)}J_p
\]
on $B_j$. The latter restriction determines the exterior field on
$B_k$ by its unique holomorphic continuation; the restriction and
continuation maps on the admissible function spaces are measurable,
as established for the augmentation. Thus $I_k$ is a measurable
function of $I_j$.
\end{proof}

For a Gibbs law $P$, let
\[
 \mathcal T_P=\bigcap_{j\ge1}\mathcal F_j^{P},
\]
where the superscript denotes completion under $P$. This is a
\emph{canonical} tail: unlike the ordinary point-configuration tail,
it retains the count information and corrected field information in
each $I_j$.

\begin{theorem}[An exact canonical-tail criterion]
\label{thm:canonical-tail-extremality}
Let $P$ be a Gibbs law for the kernels $K_j$.
\begin{enumerate}
\item For every bounded measurable real $F$,
\[
 K_jF\longrightarrow H_F:=\E_P[F\mid\mathcal T_P]
       \quad\text{almost surely and in }L^1(P)\text{ and }L^2(P).
\]
\item The law $P$ is an extreme point of the convex set of Gibbs
laws for these same kernels if and only if $\mathcal T_P$ is
$P$-trivial. Distinct extreme Gibbs laws are mutually singular on
the augmented state space.
\item If $F=F(\Xi)$ is a bounded real position observable and
$\operatorname{Var}_P(H_F)>0$, there are two Gibbs laws obtained by
conditioning $P$ on complementary canonical-tail events whose
position marginals differ. Conversely, if $H_F=\E_PF$ almost
surely for every bounded position observable $F$, then every
probability law of the form $fP$, with $f$ nonnegative,
$\mathcal T_P$-measurable and $\E_Pf=1$, has the same position
marginal as $P$.
\end{enumerate}
\end{theorem}

\begin{proof}
The first assertion is the reverse martingale convergence theorem for
the decreasing completed sigma-fields in
Lemma~\ref{lem:canonical-decreasing-information}. Boundedness gives
both claimed norm convergences.

If $A\in\mathcal T_P$ has probability strictly between zero and one,
then $P(\,\cdot\mid A)$ and $P(\,\cdot\mid A^c)$ are Gibbs laws:
for bounded $F$, proper conditioning gives
\[
 \E_P[\mathbf1_AK_jF]=\E_P[\mathbf1_AF]
\]
for every $j$. The same identity with an additional bounded
$\mathcal F_j$-measurable multiplier proves the conditional-kernel
assertion, including for completed representatives of $A$.
They decompose $P$ nontrivially, so $P$ is not extreme.

Conversely, suppose $P=tQ+(1-t)R$, where $0<t<1$ and $Q,R$ are
Gibbs laws for the same kernels. The density $f=dQ/dP$ is bounded
by $1/t$. For every bounded measurable $F$, stationarity of $Q$
and the conditional-kernel identity under $P$ give
\begin{align*}
 \E_P[fF]
 &=\E_QF=\E_QK_jF
   =\E_P[fK_jF]\\
 &=\E_P[\E_P(f\mid\mathcal F_j)K_jF]
   =\E_P[\E_P(f\mid\mathcal F_j)F].
\end{align*}
Consequently $f=\E_P(f\mid\mathcal F_j)$ almost surely for every
$j$, so $f$ is $\mathcal T_P$-measurable. If that tail is trivial,
$f=1$ almost surely, and $Q=P=R$.

For mutual singularity, take two extreme laws $Q,R$ and put
$P=(Q+R)/2$. The preceding argument shows that $f=dQ/dP$ is
measurable with respect to every $\mathcal F_j^P$, hence also to
the corresponding completed canonical tails under $Q$ and $R$.
Extremality makes $f$ constant $Q$-almost surely and constant
$R$-almost surely. If these two constants agree, then $f=1$
$P$-almost surely and $Q=R$. If they differ, their respective
level sets separate $Q$ and $R$, proving singularity.

For the third assertion, choose a real $c$ such that
$A=\{H_F>c\}$ has probability strictly between zero and one; this
is possible because $H_F$ is not constant. Then
\[
 \E_P[F\mid A]=\E_P[H_F\mid A]>c
 \ge\E_P[H_F\mid A^c]=\E_P[F\mid A^c].
\]
The conditional laws are Gibbs by the preceding argument, and their
position expectations differ. In the converse situation, conditioning
on $\mathcal T_P$ gives
\[
 \E_P[fF(\Xi)]=\E_P[f\E_P(F(\Xi)\mid\mathcal T_P)]
             =\E_PF(\Xi).
\]
This holds for every bounded position observable and proves equality
of the position marginals. The same conditioning argument shows that
$fP$ remains a Gibbs law; truncation justifies it for an unbounded
integrable density $f$.
\end{proof}

\begin{proposition}[An explicit family of distinct extreme states]
\label{prop:finite-green-extreme-states}
Let $0<\beta<2$. Every law of
Proposition~\ref{prop:finite-green-canonical-solutions} with fixed
finite total count $m$ and fixed function $h$ holomorphic on a
neighborhood of $\overline\D$ is extreme for the augmented canonical
kernels. In particular, for $m=1$ the
choices $h(z)=tz$, $t\in\mathbb R$, give distinct extreme laws
at the same temperature and the same total count. Their position
laws are mutually absolutely continuous, whereas their augmented
laws are mutually singular.
\end{proposition}

\begin{proof}
Write $P=P_{\beta,m,h}$ and let
$E_j=\{\Xi(B_j)=m\}$. The events $E_j$ increase to the entire
space almost surely. On $E_j$, the invariant $I_j$ has the single
value $(\varnothing,m,h')$. Therefore, for bounded $F$, the
conditional expectation $K_jF$ equals the deterministic number
$\E_P[F\mid E_j]$ on $E_j$, whenever $P(E_j)>0$. That number
tends to $\E_PF$, since $P(E_j)\to1$. Every configuration lies
in $E_j$ eventually, so $K_jF\to\E_PF$ almost surely. The first
part of Theorem~\ref{thm:canonical-tail-extremality} now implies
that the canonical tail is trivial: apply the conclusion to the
indicator of any completed tail event, using a measurable version.
The extremality conclusion follows.

For $m=1$ and $h(z)=tz$, the position density is proportional to
\[
 (1-|z|^2)^{-\beta/2}e^{\beta t\operatorname{Re}z}.
\]
Distinct real $t$ give distinct densities; their ratio is strictly
positive and bounded above and below. The augmented state, however,
satisfies the exact identity
$\mathcal A-J_z=t$ everywhere away from its pole. This identity
provides disjoint measurable supports for distinct $t$.
\end{proof}

\begin{remark}[Symmetry of a mixture does not imply extremality]
For example, average the one-particle states with
$h_\theta(z)=te^{-i\theta}z$, $t>0$, over a uniform angle.
The result is a rotation-invariant Gibbs law. It is not extreme:
the nonconstant global field $\mathcal A-J_z=te^{-i\theta}$ is
retained by every canonical resampling. Every component is extreme
by the same proof as the proposition. All these finite-count states
are excluded as actual gas limits by
Theorem~\ref{thm:almost-sure-infinite-below-two}.
\end{remark}

\begin{proposition}[A sufficient criterion for uniqueness of position laws]
\label{prop:canonical-kernel-oscillation-uniqueness}
Let $\mathcal G$ be a class of Gibbs laws for the same kernels and
suppose a measurable set $\Omega_0$ has full measure under every
member. Let $(F_\ell)$ be a countable family of bounded position
observables that determines point-process laws. If
\[
 \lim_{j\to\infty}
 \sup_{\omega,\omega'\in\Omega_0}
       |K_jF_\ell(\omega)-K_jF_\ell(\omega')|=0
 \quad\text{for every }\ell,
\]
then all laws in $\mathcal G$ have the same position marginal.
\end{proposition}
\begin{proof}
For $P,Q\in\mathcal G$, the Gibbs identity gives
$\E_PF_\ell=\E_PK_jF_\ell$ and
$\E_QF_\ell=\E_QK_jF_\ell$. Their difference is bounded in
absolute value by the displayed oscillation. It therefore vanishes.
The determining property proves equality of the position laws.
\end{proof}

\begin{remark}[What would establish nonuniqueness for the original sequence]
The finite-particle gas is invariant under rotations and reflections,
and these actions are continuous for vague convergence. Every actual
subsequential position law therefore has these symmetries. Asymmetric
extreme states can be constituents of an invariant law without making
the original sequence nonconvergent. Likewise, different isometry
transforms of a Gibbs state, when distinct, solve covariantly related
boundary problems; their existence alone does not show that the
specified finite-particle sequence selects both of them.

A proof of multiple subsequential limits must exhibit a bounded
continuous local position observable and two particle-number
subsequences on which its expectations have different limits. A
canonical-tail decomposition supplies candidate states, but does not
supply those subsequences. Conversely, proving tail triviality for
each candidate separately does not prove they coincide, as the explicit
extreme states above demonstrate. The oscillation condition in
Proposition~\ref{prop:canonical-kernel-oscillation-uniqueness} is a
condition that would compare different candidates; it has not been
proved for the actual infinite-particle states.
\end{remark}

\section{An independent perturbative route: exact first response at two}
\label{sec:beta2-linear-response}

This section studies perturbation in the inverse temperature itself.
The order of operations is essential: we first differentiate the exact
finite-particle gas at $\beta=2$, and then let $N$ increase.  We obtain
a unique limit of every such correlation response, with full hyperbolic
covariance.  No interchange with a fixed nonzero temperature perturbation
is asserted.

Write
\[
 U_N(z)=\sum_{i<j}\log|z_i-z_j|,
 \qquad K_N(u,v)=\frac1\pi\sum_{k=0}^{N-1}(k+1)(u\bar v)^k,
 \qquad K(u,v)=\frac1{\pi(1-u\bar v)^2}.
\]
The $p$-point correlation density with respect to $dA^{\otimes p}$ is
denoted by $\rho_{N,\beta}^{(p)}$.  All tuples below have distinct
coordinates.

\begin{theorem}[The intensity response is explicitly hyperbolic]
\label{thm:beta2-intensity-response}
For every $0<r<1$ and $0<\varepsilon<1$,
\begin{equation}
 \sup_{|a|\le r}\left|
 \left.\partial_\beta\rho_{N,\beta}^{(1)}(a)\right|_{\beta=2}
 +\frac1{\pi(1-|a|^2)^2}\right|
 \le C_{r,\varepsilon}N^{-1+\varepsilon}.
 \label{eq:beta2-intensity-response}
\end{equation}
At the center the finite-$N$ answer is exactly
\begin{equation}
 \left.\partial_\beta\rho_{N,\beta}^{(1)}(0)\right|_{\beta=2}
 =\frac1\pi\left(-1+\frac1{2N}+\frac1{N+1}\right).
 \label{eq:beta2-center-response}
\end{equation}
Consequently, for every compactly supported bounded measurable $f$,
\[
 \lim_{N\to\infty}
 \left.\partial_\beta\E_{N,\beta}\sum_i f(z_i)\right|_{\beta=2}
 =-\int_\D\frac{f(a)}{\pi(1-|a|^2)^2}\,dA(a).
\]
\end{theorem}

We first prove the integral estimate needed to justify passage through
the noncompact pair-energy integral.

\begin{lemma}[A weighted kernel-tail estimate]
\label{lem:beta2-weighted-kernel-tail}
For $0<\varepsilon<1$,
\begin{align}
 &\iint_{\D^2}(1+|\log|u-v||)|K(u,v)|\,dA(u)dA(v)<\infty,
 \label{eq:beta2-kernel-integrable}\\
 &\iint_{\D^2}(1+|\log|u-v||)|K_N(u,v)-K(u,v)|\,dA(u)dA(v)
 \le C_\varepsilon N^{-1+\varepsilon}.
 \label{eq:beta2-kernel-tail}
\end{align}
\end{lemma}

\begin{proof}
Put $t=u\bar v$.  Summing the geometric derivative gives
\[
 \pi(K-K_N)=\frac{t^N}{(1-t)^2}+
                  \frac{N t^N}{1-t}.
\]
The region where $\min(|u|,|v|)\le1/2$ contributes at most
$C(N+1)2^{-N}$ to the second assertion: the denominators are bounded
below, and $\sup_u\int_\D(1+|\log|u-v||)\,dA(v)<\infty$.
It contributes finitely to the first assertion as well.

On the remaining region write $u=re^{i\theta_1}$,
$v=se^{i\theta_2}$, let $\theta\in[-\pi,\pi]$ be their angle
difference, and put $d=1-rs$.  There are absolute positive constants
such that
\[
 |1-rse^{i\theta}|\ge c(d+|\theta|),\qquad
 |re^{i\theta}-s|\ge c|\theta|.
\]
The second inequality holds because
$|re^{i\theta}-s|^2=(r-s)^2+2rs(1-\cos\theta)$ and $r,s\ge1/2$.
The elementary bound $1+|\log x|\le C_\varepsilon(1+x^{-\varepsilon})$
for $0<x\le2$ therefore bounds the logarithmic weight by
$C_\varepsilon(1+|\theta|^{-\varepsilon})$.  Splitting the angular
integrals at $|\theta|=d$ gives
\[
 \int_{-\pi}^{\pi}
 \frac{1+|\theta|^{-\varepsilon}}{(d+|\theta|)^2}\,d\theta
 \le C_\varepsilon d^{-1-\varepsilon},\qquad
 \int_{-\pi}^{\pi}
 \frac{1+|\theta|^{-\varepsilon}}{d+|\theta|}\,d\theta
 \le C_\varepsilon d^{-\varepsilon}.
\]
If $x=1-r$ and $y=1-s$, then $d$ is comparable with $x+y$ and
$(rs)^N\le e^{-cN(x+y)}$.  Integrating any nonnegative function of
$x+y$ over $[0,1/2]^2$ is bounded above by its integral against
$(x+y)\,d(x+y)$ on $[0,1]$.  Thus the two tail terms are bounded by
\[
 C_\varepsilon\int_0^1 e^{-cNh}h^{-\varepsilon}\,dh
 +C_\varepsilon N\int_0^1e^{-cNh}h^{1-\varepsilon}\,dh
 \le C_\varepsilon N^{-1+\varepsilon}.
\]
Without the factors $(rs)^N$, the first angular estimate reduces
the integrability assertion to $\int_0^1h^{-\varepsilon}\,dh<\infty$.
\end{proof}

\begin{lemma}[Finite Palm response]
\label{lem:beta2-finite-palm-response}
For a distinct tuple $\mathbf a=(a_1,\ldots,a_p)$ and $N\ge p$, let
$P_{N,\mathbf a}$ be the orthogonal projection in $L^2(\D,dA)$ onto
the span of $K_N(\cdot,a_1),\ldots,K_N(\cdot,a_p)$.  Its rank is $p$.
Set $P_N=P_{N,\mathbf a}$ and
\[
 V_N(u)=\int_\D\log|u-v|K_N(v,v)\,dA(v)
       =\frac12\sum_{n=1}^N\frac{|u|^{2n}-1}{n}.
\]
Then
\begin{align}
 R_{N,p}(\mathbf a)
 &:=\frac{\left.\partial_\beta\rho_{N,\beta}^{(p)}(\mathbf a)
                       \right|_{\beta=2}}
                         {\rho_{N,2}^{(p)}(\mathbf a)}\notag\\
 &=\sum_{i<j}\log|a_i-a_j|+
      \sum_i V_N(a_i)-\int_\D V_N(u)P_N(u,u)\,dA(u)\notag\\
 &\quad-\sum_i\int_\D\log|a_i-u|P_N(u,u)\,dA(u)\notag\\
 &\quad+\frac12\iint_{\D^2}\log|u-v|
       \{P_N(u,u)P_N(v,v)-|P_N(u,v)|^2\}\,dA(u)dA(v)\notag\\
 &\quad+\operatorname{Re}\iint_{\D^2}\log|u-v|
                  K_N(u,v)\overline{P_N(u,v)}\,dA(u)dA(v).
 \label{eq:beta2-finite-palm-response}
\end{align}
\end{lemma}

\begin{proof}
For fixed $N$ differentiation under the finite-dimensional integrals
is valid on every compact positive-$\beta$ interval.  Indeed
$|\Delta|^\beta(1+|\log|\Delta||)$ is bounded on $\D^N$ uniformly on
such an interval, after using $|\Delta|\le2^{N(N-1)/2}$ and
$\sup_{0<x\le1}x^c|\log x|<\infty$.  The same observation applies
when a distinct tuple is fixed and only the remaining particles are
integrated.  Differentiation of that quotient gives
\[
 R_{N,p}(\mathbf a)=\sum_{i<j}\log|a_i-a_j|
  +\E_{N,2}^{!\mathbf a}\left[
       \sum_i\sum_{u\in X}\log|a_i-u|+U_{N-p}(X)\right]
  -\E_{N,2}U_N.
\]
The reduced Palm kernel is $K_N-P_N$: expanding correlation
determinants with the distinguished tuple by the Schur-complement
identity proves this directly.  Its one-point function is
$K_N(u,u)-P_N(u,u)$.  Expanding the difference of its two-point
function and that of $K_N$ gives
\begin{align*}
 &-K_N(u,u)P_N(v,v)-P_N(u,u)K_N(v,v)\\
 &\qquad+P_N(u,u)P_N(v,v)
 +2\operatorname{Re}\{K_N(u,v)\overline{P_N(u,v)}\}
 -|P_N(u,v)|^2.
\end{align*}
Integration against $\tfrac12\log|u-v|$ proves the stated formula.
Finally angular averaging gives
$\int\log|u-v|(n/\pi)|v|^{2n-2}\,dA(v)
=(|u|^{2n}-1)/(2n)$, proving the expression for $V_N$.
\end{proof}

\begin{theorem}[All first correlation responses have an invariant limit]
\label{thm:beta2-all-response}
Let $P_{\mathbf a}$ be the rank-$p$ projection onto the span of
$K(\cdot,a_1),\ldots,K(\cdot,a_p)$, and put
\[
 v(u)=-\tfrac12\log(1-|u|^2).
\]
Define $R_p(\mathbf a)$ by the right side of
\eqref{eq:beta2-finite-palm-response}, replacing $P_N,K_N,V_N$ by
$P_{\mathbf a},K,v$, respectively.  Its integrals are absolutely
convergent.  Uniformly on compact subsets of the distinct-tuple space,
\begin{equation}
 R_{N,p}(\mathbf a)\longrightarrow R_p(\mathbf a),\qquad
 \left.\partial_\beta\rho_{N,\beta}^{(p)}(\mathbf a)\right|_{\beta=2}
 \longrightarrow\det[K(a_i,a_j)]_{i,j=1}^pR_p(\mathbf a).
 \label{eq:beta2-all-response-limit}
\end{equation}
For every hyperbolic isometry $\varphi$,
\begin{equation}
 R_p(\varphi a_1,\ldots,\varphi a_p)=R_p(a_1,\ldots,a_p).
 \label{eq:beta2-response-invariant}
\end{equation}
In particular the signed limiting correlation response in
\eqref{eq:beta2-all-response-limit} transforms as an invariant
$p$-point correlation measure.
\end{theorem}

\begin{proof}
On a compact set of distinct tuples, the Gram matrices
$[K_N(a_i,a_j)]$ converge uniformly to strictly positive definite
matrices.  Strict positivity follows, for example, from polynomial
interpolation at the distinct points.  Their inverses are uniformly
bounded.  Moreover $K_N(u,a_i)$ converges to $K(u,a_i)$ uniformly for
all $u\in\D$, with a geometrically small error, because the $a_i$
remain in a compact subdisk.  The finite-rank projection formula
therefore gives bounded kernels $P_N$, uniformly on $\D^2$, and their
uniform convergence to $P_{\mathbf a}$.

Since $\int P_N(u,u)dA(u)=p$, the additive constant $-H_N/2$ in
$V_N$ cancels exactly in
$\sum_i V_N(a_i)-\int V_NP_N(u,u)dA$.  Its remaining part
$v_N(u)=\tfrac12\sum_{n\le N}|u|^{2n}/n$ increases to $v(u)$ and
\[
 \int_\D(v-v_N)\,dA
 =\frac\pi2\sum_{n>N}\frac1{n(n+1)}=\frac\pi{2(N+1)}.
\]
This proves convergence of the $V_N$ terms.  Logarithmic kernels
are uniformly integrable when one argument ranges over a compact
set, and integrable on $\D^2$, so all terms containing only $P_N$
converge.  Lemma~\ref{lem:beta2-weighted-kernel-tail} treats the last
term and also proves its absolute convergence.  These estimates
prove both limits and, for $p=1$, the rate $O(N^{-1+\varepsilon})$.

The limiting signed correlation response is locally integrable also
across partial collisions.  Indeed, for $|a_i|\le r<1$, every function
in the range of $P_{\mathbf a}$ has the form $f=P/Q_{\mathbf a}$,
where
\[
 Q_{\mathbf a}(u)=\prod_{i=1}^p(1-u\bar a_i)^2,
 \qquad \deg P\le2p-2.
\]
On the entire disk this denominator is bounded above and below by
positive constants depending only on $p,r$.  Equivalence of norms
on the finite-dimensional polynomial space therefore gives
$\|f\|_{L^\infty(\D)}\le C_{p,r}\|f\|_{L^2(\D)}$.  An orthonormal
basis of the rank-$p$ range then proves
$|P_{\mathbf a}(u,v)|\le C_{p,r}$, uniformly even as some nodes
coalesce.  The absolutely convergent integral formula consequently
shows
\[
 \left|R_p(\mathbf a)-\sum_{i<j}\log|a_i-a_j|\right|\le C_{p,r}.
\]
The logarithmic terms are locally integrable, and the Bergman
correlation determinant is locally bounded, as required.

We verify invariance rather than infer it from formal cancellation
of divergent energies.  First take an orientation-preserving
isometry and put $\ell(u)=\tfrac12\log|\varphi'(u)|$.  The function
$\ell$ is bounded and harmonic on $\D$, and
\[
 \log|\varphi u-\varphi v|=\log|u-v|+\ell(u)+\ell(v),
 \qquad v(\varphi u)=v(u)-\ell(u).
\]
The Bergman unitary $f\mapsto\varphi'(f\circ\varphi)$ carries the
projection associated with $\varphi\mathbf a$ to that associated
with $\mathbf a$.  It follows that all projection and kernel factors
in the defining integrals have the required change-of-variables
covariance.  Set $T=\int_\D\ell(u)P_{\mathbf a}(u,u)dA(u)$ and
$A=\sum_i\ell(a_i)$.  The five lines of the response formula change,
respectively, by
\[
 (p-1)A,\qquad -A+T,\qquad -pA-pT,
 \qquad (p-1)T,\qquad 2T.
\]
For the fourth change use $P^2=P$ and $\operatorname{Tr}P=p$;
for the last use $KP=P$.  The sum is $2(T-A)$.

For completeness $T=A$.  If $h$ is bounded holomorphic, multiplication
by $h$ has adjoint mapping $K(\cdot,a_i)$ to
$\overline{h(a_i)}K(\cdot,a_i)$.  On their finite-dimensional span
its adjoint compression therefore has trace
$\sum_i\overline{h(a_i)}$, and taking conjugates gives
\[
 \int_\D h(u)P_{\mathbf a}(u,u)dA(u)=\sum_i h(a_i).
\]
Here $\ell$ is the real part of a bounded holomorphic function
(use the explicit logarithm of the derivative of a disk
automorphism).  Taking real parts proves $T=A$.  Thus the response
is unchanged.  Complex conjugation leaves every term unchanged and
accounts for the orientation-reversing isometries.
\end{proof}

\begin{proof}[Proof of Theorem~\ref{thm:beta2-intensity-response}]
For $p=1$ the fourth line of
\eqref{eq:beta2-finite-palm-response} vanishes, because $P_N$ has
rank one.  The preceding theorem proves that its limiting response
$R_1$ is invariant, hence constant.  It remains to compute that
constant, and the same computation gives the exact finite-$N$
central response.

At $a=0$ one has $P_N(u,v)=1/\pi$.  Direct radial integration gives
\[
 V_N(0)-\frac1\pi\int_\D V_N(u)dA(u)
 =-\frac12+\frac1{2(N+1)},\qquad
 \frac1\pi\int_\D\log|u|dA(u)=-\frac12.
\]
For the last term, angular averaging gives
\[
 \frac1{\pi^2}\iint_{\D^2}\log|u-v|\,dA(u)dA(v)=-\frac14,
\]
and for every integer $k\ge1$,
\begin{align*}
 \frac1{\pi^2}\iint_{\D^2}
       \log|u-v|u^k\bar v^k\,dA(u)dA(v)
 &=-\frac4k\int_0^1r\,dr\int_0^r s^{2k+1}\,ds\\
 &=-\frac1{k(k+1)(k+2)}.
\end{align*}
Indeed the $k$-th angular Fourier coefficient of
$\log|re^{i\theta}-s|$ is
$-(2k)^{-1}(\min(r,s)/\max(r,s))^k$; this follows by expanding
$\log(1-z)$ for $|z|<1$, and the equality of radii is a null set.
Thus the last term is
\[
 -\frac14-\sum_{k=1}^{N-1}\frac1{k(k+2)}
 =-1+\frac1{2N}+\frac1{2(N+1)}.
\]
Adding the three terms gives
$R_{N,1}(0)=-1+(2N)^{-1}+(N+1)^{-1}$ and hence $R_1\equiv-1$.
Since $\rho_{N,2}^{(1)}(0)=1/\pi$, this proves the exact central
formula.  The rate already established in the preceding proof,
together with locally geometric convergence of $K_N(a,a)$ to
$K(a,a)$, proves~\eqref{eq:beta2-intensity-response}.  Integrating
against $f$ proves the final assertion.
\end{proof}

\begin{corollary}[Explicit response at particle fusion]
\label{cor:beta2-fusion-response}
Let $H_n=\sum_{j=1}^n j^{-1}$.  If the distinct coordinates of
$\mathbf a\in\D^p$ all tend to zero, then
\begin{equation}
 R_p(\mathbf a)-\sum_{i<j}\log|a_i-a_j|\longrightarrow C_p,
 \qquad
 C_p=\frac{2p-1}{4}H_p-\frac{2p+1}{2}H_{2p}+p.
 \label{eq:beta2-fusion-response}
\end{equation}
In particular
\[
 R_2(0,a)=\log|a|-\frac{25}{12}+o(1),\qquad a\to0.
\]
At an arbitrary common limiting point $q\in\D$, the corresponding
constant, after subtracting the Euclidean logarithmic Vandermonde,
is $C_p-\tfrac12p(p-1)\log(1-|q|^2)$.
\end{corollary}

\begin{proof}
As all nodes tend to zero, the projection $P_{\mathbf a}$ converges
uniformly on $\D^2$ to $K_p$.  One way to verify this without an
assumption on the relative speeds of the nodes is to replace the
spanning functions $K(\cdot,a_i)$ by their Newton divided differences
in the variables $\bar a_i$.  For a holomorphic function $F$ on a
circle containing all those variables, the divided difference has
the contour formula
\[
 F[t_0,\ldots,t_j]
 =\frac1{2\pi i}\int_{|\zeta|=r_0}
           \frac{F(\zeta)}{\prod_{i=0}^j(\zeta-t_i)}\,d\zeta.
\]
For $F(\zeta)=\pi^{-1}(1-u\zeta)^{-2}$, $|u|<1$ and fixed
$r_0<1$, this formula proves convergence uniformly in $u$ to
$(j+1)u^j/\pi$.  The limiting Gram matrix is nonsingular, so the
projection convergence follows.  The integral estimates used in
Theorem~\ref{thm:beta2-all-response} now permit passage to the limit
after removing the explicit logarithmic Vandermonde.

We evaluate the remaining terms.  Write $n=k+1$ for the normalized
monomial of degree $k$.  Its expectations of $v(u)$ and $\log|u|$
are respectively $H_n/2$ and $-1/(2n)$.  Thus the second and third
terms of the response formula contribute
\[
 -\frac12\sum_{n=1}^p H_n+\frac p2H_p.
\]
The fourth term is $\E_{p,2}U_p$, and the last is
$-\sum_{n=1}^pH_n$.  Here are direct computations of both, including
the infinite sum in the last term.  If
\[
 J_{kl}=\frac1{\pi^2}\iint_{\D^2}\log|u-v|
             u^k\bar u^l v^l\bar v^k\,dA(u)dA(v),
\]
angular averaging followed by radial integration gives
\[
 (k+1)(l+1)J_{kl}=
 \begin{cases}
 -1/[4(k+1)],&k=l,\\[2pt]
 -\min(k+1,l+1)/[|k-l|(k+l+2)],&k\ne l.
 \end{cases}
\]
For fixed $n=l+1$ the sum over $k\ge0$ is $-H_n$: separate $k=l$
and use
\[
 \sum_{j>n}\frac{n}{j^2-n^2}=\frac12H_{2n},\qquad
 \sum_{j<n}\frac{j}{n^2-j^2}
 =\frac12(H_{n-1}+H_n-H_{2n-1}).
\]
The sum is absolutely convergent.  Equivalently it is the limit of
the kernel integral, justified by
Lemma~\ref{lem:beta2-weighted-kernel-tail}.  This proves the value
of the last term.

The radial product contribution to $\E_{p,2}U_p$ is
$-\tfrac14\sum_{n,m=1}^p(n+m)^{-1}$, because the larger of two
independent squared radii with respective beta parameters $n,m$
has beta parameter $n+m$.  Subtracting the finite kernel term gives
\begin{align*}
 \E_{p,2}U_p
 &=-\frac14\sum_{n,m=1}^p\frac1{n+m}
   +\frac18H_p+\sum_{1\le n<j\le p}\frac{n}{j^2-n^2}\\
 &=\frac{6p+5}{4}H_p-\frac{2p+1}{2}H_{2p}-\frac p2.
\end{align*}
The last equality follows by substituting
\[
 \sum_{n,m=1}^p\frac1{n+m}
 =(2p+1)H_{2p}-2(p+1)H_p
\]
and the preceding partial-fraction sum; for the latter simplification
one can use
$\sum_{j=1}^pH_{2j-1}=(p+\tfrac12)H_{2p}-\tfrac14H_p-p$.
Combining all four contributions and using
$\sum_{n=1}^pH_n=(p+1)H_p-p$ proves the value of $C_p$.
For fusion at $q$, apply a disk automorphism taking $q$ to zero.
Its derivative there has modulus $(1-|q|^2)^{-1}$; the assertion
then follows from~\eqref{eq:beta2-response-invariant}.
\end{proof}

\begin{remark}[What the response settles and what it leaves open]
The one-point response has exactly the invariant shape, with density
slope $-1/\pi$.  Thus failure of hyperbolic invariance cannot be
detected at the first perturbative order about $\beta=2$; the same
statement holds for every fixed correlation order.  If full
convergence on a temperature interval and an exchange of the
temperature derivative with that limit were separately established,
the derivative of the limiting intensity at two would necessarily
be the value proved above.  Neither assumption follows from the
first-response theorem.  In particular no fixed $\beta\ne2$ is
settled merely by a Taylor formula whose second derivative has not
been bounded uniformly in $N$.
\end{remark}

\begin{remark}[A separate monotonicity route]
For this unscaled problem, proving that all correlation functions
increase with $N$ for a fixed $0<\beta<2$ would establish uniqueness
using the already proved local bounds and correlation-moment
determination.  Monotonicity at $\beta=2$ follows from the increasing
projection kernels.  The broader $\beta<2$ statement was explicitly
conjectured by Hedenmalm in his 2015 lectures, and is not assumed here;
see \url{https://irma.math.unistra.fr/~klevtsov/past/Hedenmalm.pdf},
slide~6.  The exact central first response above is compatible with
this direction, but does not prove monotonicity at another temperature.
\end{remark}

\section{A uniform second temperature response at the determinantal point}
\label{sec:beta2-second-response}

The next estimate controls a second connected energy diagram.  It
still concerns derivatives at the single temperature $\beta=2$;
it does not give a uniform interval of analyticity.

\begin{theorem}[Uniform second response of local correlation logarithms]
\label{thm:beta2-second-response}
Let $p\ge1$ and let $A$ be a compact subset of the space of distinct
tuples in $\D^p$.  There is a finite constant $C_A$ such that
\begin{equation}
 \sup_{N\ge p}\ \sup_{\mathbf a\in A}
 \left|\left.\partial_\beta^2
       \log\rho_{N,\beta}^{(p)}(\mathbf a)\right|_{\beta=2}\right|
 \le C_A.
 \label{eq:beta2-second-response}
\end{equation}
Consequently the second derivatives of the correlation densities
themselves are uniformly bounded on $A$ at $\beta=2$.
\end{theorem}

We use ordinary area and the notation of
Section~\ref{sec:beta2-linear-response}.  Write $\mathcal H=L^2(\D,dA)$,
$\phi_n(z)=\sqrt{(n+1)/\pi}\,z^n$, $K_N$ also for its orthogonal
projection, and $Q_N=I-K_N$.  Put $f(u,v)=\log|u-v|$.  For a finite-rank
projection $J$ define
\begin{align*}
 V_J(u)&=\int_\D f(u,v)J(v,v)dA(v),\\
 (A_Jh)(u)&=\int_\D f(u,v)J(u,v)h(v)dA(v),
 \qquad F_J=M_{V_J}-A_J.
\end{align*}
Here $M_s$ denotes multiplication by $s$.  These expressions below
are applied to the finite occupied subspaces, where all their
matrix elements and norms are finite.  All logarithmic singularities
are square integrable with respect to the bounded finite-volume
correlation densities.

\begin{lemma}[The Slater variance identity]
\label{lem:beta2-slater-variance}
Let $J$ be a finite-rank orthogonal projection onto bounded functions,
and let $Q=I-J$.  For its determinantal process, the variance of
\[
 \sum_{i<j}f(z_i,z_j)+\sum_i s(z_i)
\]
is
\begin{equation}
 \|Q(F_J+M_s)J\|_{\mathrm{HS}}^2+
 \|Q^{\wedge2}M_fJ^{\wedge2}\|_{\mathrm{HS}}^2.
 \label{eq:beta2-slater-variance}
\end{equation}
The second term is an operator between the antisymmetric tensor
squares, with their usual Hilbert norms.
\end{lemma}

\begin{proof}
Choose orthonormal bases of $J\mathcal H$ and $Q\mathcal H$.
The normalized determinant of the occupied basis is a unit vector
in the antisymmetric tensor power of $\mathcal H$, and its squared
modulus is the determinantal joint density.  Multiplication by the
displayed symmetric statistic changes at most two occupied basis
vectors.  Expand the resulting vector in the orthonormal determinant
basis.  Its component on the original determinant is its mean.
The components with exactly one occupied vector replaced are the
matrix elements of $Q(F_J+M_s)J$: direct integration over each
unchanged orbital gives the direct term $V_J$ and the exchange term
$-A_J$.  Components with exactly two replacements are the matrix
elements of $Q^{\wedge2}M_fJ^{\wedge2}$.  There are no other
components.  Parseval's identity after removing the mean gives
\eqref{eq:beta2-slater-variance}.  To justify the expansion one can
first use bounded $f,s$ and a finite ambient orthonormal basis, then
pass through its increasing spans and bounded truncations.  The
square-integrability just stated gives convergence of the original
random statistic in $L^2$, hence of its orthogonal components.
\end{proof}

\begin{lemma}[A uniformly bounded particle--hole operator]
\label{lem:beta2-hartree-fock-bound}
The operators
\[
 T_N=Q_NF_{K_N}K_N
\]
satisfy
\begin{equation}
 \sup_N\|T_N\|_{\mathrm{op}}\le\frac\pi{\sqrt6}+\frac14.
 \label{eq:beta2-hartree-fock-bound}
\end{equation}
After adding the irrelevant scalar $H_N/2$ to $F_{K_N}$,
one also has
\begin{equation}
 \|(F_{K_N}+H_N/2)\phi_n\|_2\le C\log(n+2),
 \qquad 0\le n<N.
 \label{eq:beta2-hartree-fock-column}
\end{equation}
\end{lemma}

\begin{proof}
Rotation covariance makes $F_{K_N}\phi_n$ a radial multiplier of
$\phi_n$.  Put $t=|u|^2$.  Angular integration as in
Section~\ref{sec:beta2-linear-response} gives, for $0\le n<N$,
\begin{align*}
 \frac{A_{K_N}\phi_n(u)}{\phi_n(u)}
 &=-\frac12H_{n+1}
   -\frac12\sum_{j=1}^{N-n-1}\frac{t^j}{j}
   +\frac1{2(n+1)}\sum_{j=N-n}^N t^j.
\end{align*}
The quotient is interpreted by its radial continuous formula at zero.
For clarity, the summand of $A_{K_N}$ indexed by $k$ contributes
\[
 \begin{cases}
 -1/[2(n-k)]+t^{k+1}/[2(n+1)],&k<n,\\
 (t^{n+1}-1)/[2(n+1)],&k=n,\\
 -(k+1)t^{k-n}/[2(k-n)(n+1)]+t^{k+1}/[2(n+1)],&k>n.
 \end{cases}
\]
Their sum is the preceding identity.  Since
$V_{K_N}=-H_N/2+\tfrac12\sum_{j=1}^Nt^j/j$, the multiplier after
discarding $-H_N/2$ is
\begin{equation}
 b_{n,N}(t)=\frac12H_{n+1}
 +\frac12\sum_{j=1}^N\frac{t^j}{j}
 +\frac12\sum_{j=1}^{N-n-1}\frac{t^j}{j}
 -\frac1{2(n+1)}\sum_{j=N-n}^Nt^j.
 \label{eq:beta2-hartree-fock-multiplier}
\end{equation}
The last term has range in $[-1/2,0]$ and hence standard deviation
at most $1/4$.  Under $|\phi_n(u)|^2dA(u)$, the squared radius has
density $(n+1)t^n dt$.  For this law
\[
 \E[-\log(1-t)]=H_{n+1},\qquad
 \Var[-\log(1-t)]=\sum_{j=1}^{n+1}\frac1{j^2}\le\frac{\pi^2}{6}.
\]
These identities follow by differentiating twice the elementary
beta integral $\int_0^1t^n(1-t)^sdt$ at $s=0$.
For each partial sum $a_M(t)=\sum_{j\le M}t^j/j$,
\[
 |a_M(t)-a_M(s)|\le
 |\log(1-t)-\log(1-s)|,
\]
because its derivative is at most $(1-t)^{-1}$.
The independent-copy identity for variance therefore bounds its
standard deviation by $\pi/\sqrt6$.  Taking standard deviations
in~\eqref{eq:beta2-hartree-fock-multiplier} proves the claimed bound
on each centered radial multiplier.  The resulting vectors for
different $n$ have different angular frequencies and are orthogonal;
$Q_N$ centers precisely that radial multiplier.  This proves the
operator norm bound.  The same formula, its means, and the variance
bound give~\eqref{eq:beta2-hartree-fock-column}.
\end{proof}

\begin{lemma}[Two-body vertices with one controlled analytic orbital]
\label{lem:beta2-two-body-vertex}
Let $h=\sum_{l=0}^{N-1}c_l\phi_l$ and put
$\|h\|_* =\sum_l|c_l|\sqrt{l+1}$.  There is an absolute constant
$C$ such that
\begin{align}
 \|(Q_N\otimes Q_N)M_f(h\otimes K_N)\|_{\mathrm{HS}}
 &\le C\|h\|_*,\label{eq:beta2-loss-vertex}\\
 \|(\langle h|\otimes Q_N)M_f(K_N\otimes K_N)\|_{\mathrm{HS}}
 &\le C\|h\|_*.
 \label{eq:beta2-gain-vertex}
\end{align}
Here $h\otimes K_N$ maps $v$ to $h\otimes K_Nv$, and the bra in
the second formula integrates the first output coordinate against
$\bar h$.
\end{lemma}

\begin{proof}
Decompose $f(u,v)=\log|1-u\bar v|-g(u,v)$, where $g$ is the
positive hyperbolic Green function.  Since $K_N(v,v)\le K(v,v)$,
\[
 \int_\D g(u,v)^2K_N(v,v)dA(v)
 \le\int_\D g(u,v)^2\,d\nu(v)=\frac{\pi^2}{12}.
\]
The last equality follows by a disk automorphism and the radial
integral $\tfrac14\int_0^1(\log t)^2(1-t)^{-2}dt$.
For the loss estimate one can omit the output projections, and
the Hilbert--Schmidt norm squared is bounded by the last constant
times $\|h\|_2^2$.  For the gain estimate, fix the second input
point $v$.  The norm squared of the first-coordinate functional
on the occupied subspace is
$\|K_N(g(\cdot,v)h)\|_2^2\le\int g(u,v)^2|h(u)|^2dA(u)$.
Integration against $K_N(v,v)dA(v)$ proves the same bound.

For the image kernel use its angular series
\[
 \log|1-u\bar v|=-\frac12\sum_{k=1}^\infty
             \frac{u^k\bar v^k+\bar u^kv^k}{k}.
\]
First take $h=\phi_l$, $l<N$.  Terms with different angular
frequencies are orthogonal in the squared norms below.  In the
loss estimate the first series survives $Q_N$ in the first
coordinate only for $k\ge N-l$.  Its squared norm is at most
\[
 \frac14\sum_{k\ge N-l}\frac{l+1}{k^2(l+k+1)}
       \sum_{n=0}^{N-1}\frac{n+1}{n+k+1}
 \le C(l+1).
\]
Indeed the inner sum is at most $N$, $l+k+1\ge N+1$, and
$\sum_{k\ge N-l}k^{-2}\le C/(N-l)\le C$.
For the second series the second-coordinate condition is
$n+k\ge N$.  Its squared norm is at most
\[
 \frac14\sum_{k\ge1}\frac{l+1}{k^2(l+k+1)}
        \sum_{n=\max(0,N-k)}^{N-1}\frac{n+1}{n+k+1}
 \le\frac{l+1}{4}\sum_{k\ge1}\frac1{k(l+k+1)}
 \le C(l+1).
\]
Projecting the other coordinate can only decrease these bounds.

For the gain estimate the first series has nonzero first-coordinate
matrix element only when its occupied index is $l-k$, with
$1\le k\le l$.  The resulting squared norm is
\[
 \frac14\sum_{k=1}^l\frac{l-k+1}{k^2(l+1)}
             \|Q_NM_{\bar z^k}K_N\|_{\mathrm{HS}}^2
 \le\frac14\sum_{k=1}^l\frac1k\le C(l+1).
\]
Here $\|Q_NM_{\bar z^k}K_N\|_{\mathrm{HS}}^2\le k$.
To see this, the squared column norm for $\phi_n$ equals
$(n+1)/(n+k+1)$ for $n<k$, and
$k^2/[(n+1)(n+k+1)]$ for $n\ge k$.  Summing all $n\ge0$
telescopes to $k$; restricting to $n<N$ decreases the sum.
The second image series has first occupied index $l+k<N$ and
second index $n$ with $n+k\ge N$.  Its squared norm is bounded by
\[
 \frac14\sum_{k=1}^{N-l-1}\frac{l+1}{k^2(l+k+1)}
       \sum_{n=\max(0,N-k)}^{N-1}\frac{n+1}{n+k+1}
 \le C(l+1).
\]
The triangle inequality between the two series gives both desired
bounds for $\phi_l$.  Their sums converge in the indicated
Hilbert--Schmidt spaces; alternatively start with an angular
Abel factor and pass to the limit.  Finally apply the triangle
inequality to the expansion of $h$, proving the asserted norm
$\|h\|_*$.
\end{proof}

\begin{lemma}[The screened Palm perturbation is summable]
\label{lem:beta2-screened-palm-trace}
For $\mathbf a\in A$, let $P_N=P_{N,\mathbf a}$ and
$P=P_{\mathbf a}$ be the projections in
Lemma~\ref{lem:beta2-finite-palm-response}.  Set
\[
 s_{\mathbf a}(u)=\sum_i\log|u-a_i|,
 \qquad B_N=M_{s_{\mathbf a}-V_{P_N}}+A_{P_N}.
\]
Then, uniformly over $\mathbf a\in A$ and $N\ge p$,
\begin{equation}
 \sum_{n=0}^{N-1}\|B_N\phi_n\|_2\le C_A.
 \label{eq:beta2-screened-column-sum}
\end{equation}
There are orthonormal bases $h_{j,N}$ of the ranges of $P_N$ such
that $\|h_{j,N}\|_*\le C_A$ and
\begin{equation}
 \|(F_{K_N}+H_N/2)h_{j,N}\|_2\le C_A.
 \label{eq:beta2-screened-F-column}
\end{equation}
\end{lemma}

\begin{proof}
The Gram-matrix argument in
Theorem~\ref{thm:beta2-all-response} gives orthonormal bases whose
coefficients in the $\phi_l$ basis satisfy
$|c_{j,l,N}|\le C_A\sqrt{l+1}\,r^l$ for some $r<1$ independent
of $N$ and $\mathbf a\in A$.  After enlarging the constants,
the same bounds hold for all $N\ge p$.  They prove the last two
assertions using~\eqref{eq:beta2-hartree-fock-column}.
They also give $\|P_N-P\|_{L^\infty(\D^2)}\le C_A r_1^N$ for
some $r_1<1$.

Put $S=s_{\mathbf a}-V_P$.  The harmonic trace identity proved in
Theorem~\ref{thm:beta2-all-response} implies
\[
 V_P(u)=\sum_i\log|u-a_i|,\qquad |u|>1.
\]
For example expand the logarithm in inverse powers of $u$, use
$\int P(v,v)dA(v)=p$ and
$\int v^kP(v,v)dA(v)=\sum_i a_i^k$, and then use continuity at
the boundary.  Thus $S$ vanishes identically outside the disk.
The potential of the smooth bounded density $P(v,v)$ on the disk
is continuously differentiable across its boundary, because the
gradient of the logarithmic kernel is locally integrable.  The
first normal derivative of $S$ is consequently zero there as well.
Inside a fixed boundary annulus, $S$ solves
$\Delta S=-2\pi P(u,u)$ with zero boundary data.  The smooth
Dirichlet regularity for the disk (equivalently its Poisson and
Green integral formulas differentiated in boundary charts) gives
a bounded second derivative up to the boundary.  All density
derivatives and source distances are uniform over $A$, so
\begin{equation}
 |S(u)|\le C_A(1-|u|^2)^2
 \quad\text{on a fixed boundary annulus}.
 \label{eq:beta2-screened-quadratic-boundary}
\end{equation}
On its complementary compact disk, $S$ has only logarithmic
singularities, with uniformly bounded $L^2$ norm.  Consequently
\[
 \sum_{n\ge0}\|S\phi_n\|_2<\infty
\]
uniformly over $A$: on the annulus use
$(n+1)\int_0^1t^n(1-t)^4dt=O((n+1)^{-4})$ and then take its
square root; on the compact disk use the geometric factor
$\sqrt{n+1}\,r^n$.

It remains to treat $A_P$.  Write its kernel as
$f(u,v)\sum_j h_j(u)\overline{h_j(v)}$, where the finitely many
$h_j$ extend holomorphically to a fixed disk of radius greater
than one and have uniformly bounded analytic norms there.  If
$L$ is the integral operator with kernel $f$, then
$A_P=\sum_jM_{h_j}LM_{\bar h_j}$.  We claim
\begin{equation}
 \sum_{n\ge0}\|L(\bar z^k\phi_n)\|_2\le C(k+1).
 \label{eq:beta2-log-analytic-column}
\end{equation}
For $n>k$, direct angular integration gives
\[
 L(\bar z^k\phi_n)(u)=
 -\frac{\sqrt{\pi(n+1)}\,u^{n-k}}{2(n-k)}
 \left\{\frac{1-|u|^{2k+2}}{k+1}
                    +\frac{|u|^{2k+2}}{n+1}\right\}.
\]
If $n\ge2k+1$, its norm is at most $C(n+1)^{-2}$, since the
bracket is bounded by $(1-|u|^2)+(n+1)^{-1}$ and the relevant
radial beta integrals have that order.  For $n\le2k$, the
Hilbert--Schmidt norm of $L$ bounds each column by an absolute
constant.  This proves~\eqref{eq:beta2-log-analytic-column}.
Expanding the analytic functions $h_j$ in their geometrically
decaying Taylor series and summing proves
$\sum_n\|A_P\phi_n\|_2\le C_A$.

The replacement of $P$ by $P_N$ has error at most $C_A r_1^N$
in the sup norm of its kernel.  The associated potential error
is bounded in sup norm by the same order, using the uniform
integrability of $|\log|u-v||$; the exchange operator error is
bounded in operator norm by its Hilbert--Schmidt norm of that
order.  Summing its first $N$ column bounds adds at most
$C_A N r_1^N$.  Combining these bounds proves
\eqref{eq:beta2-screened-column-sum}.
\end{proof}

\begin{proof}[Proof of Theorem~\ref{thm:beta2-second-response}]
Differentiation of the fixed-tuple partition-function quotient
twice gives
\[
 \left.\partial_\beta^2\log\rho_{N,\beta}^{(p)}(\mathbf a)
 \right|_{\beta=2}
 =\Var_{K_N-P_N}\left(U_{N-p}+\sum_{u\in X}s_{\mathbf a}(u)\right)
       -\Var_{K_N}(U_N).
\]
The deterministic pair energy of the distinguished tuple drops
out.  The differentiations are legitimate by the same logarithmic
domination as before, now with its square.

Abbreviate $K=K_N$, $P=P_N$, $J=K-P$, $Q=I-K$ and $B=B_N$.
Then $I-J=Q+P$ and
$F_J+M_{s_{\mathbf a}}=F_K+B$.  In all one-body expressions below
replace $F_K$ by $F=F_K+H_N/2$, which does not change a matrix
element between an occupied and an unoccupied subspace.
The difference of the single-replacement variances is exactly
\begin{align*}
 &-\|QFP\|_{\mathrm{HS}}^2
 +2\operatorname{Re}\langle QFJ,QBJ\rangle_{\mathrm{HS}}
 +\|QBJ\|_{\mathrm{HS}}^2
 +\|P(F+B)J\|_{\mathrm{HS}}^2.
\end{align*}
The first term is bounded by $p\sup_N\|T_N\|_{\mathrm{op}}^2$.
The column sum in~\eqref{eq:beta2-screened-column-sum} bounds the
trace norm of $BK$ and hence of $BJ$, so the cross term is bounded
by $2\|T_N\|_{\mathrm{op}}\|BJ\|_1$.  It bounds the third term
as well.  Finally $F$ and $B$ are self-adjoint on the vectors
in question, so
\[
 \|P(F+B)J\|_{\mathrm{HS}}
 =\|J(F+B)P\|_{\mathrm{HS}}
 \le\|FP\|_{\mathrm{HS}}+\|BP\|_{\mathrm{HS}}\le C_A
\]
by Lemma~\ref{lem:beta2-screened-palm-trace}.  Thus the entire
single-replacement difference is uniformly bounded.

For the double-replacement term use the orthogonal decompositions
\[
 \wedge^2K\mathcal H
 =\wedge^2J\mathcal H\oplus(J\mathcal H\wedge P\mathcal H)
                         \oplus\wedge^2P\mathcal H,
\]
and the corresponding decomposition of $\wedge^2(Q+P)\mathcal H$.
Its variance difference is the sum of four signed squared norms:
the negative blocks have input $J\wedge P$ or $\wedge^2P$ and
output $\wedge^2Q$, while the positive blocks have input
$\wedge^2J$ and output $Q\wedge P$ or $\wedge^2P$.
Each negative block is bounded by the loss estimate in
Lemma~\ref{lem:beta2-two-body-vertex}, summed over an orthonormal
basis of $P\mathcal H$.  The positive block with output $Q\wedge P$
is bounded by its gain estimate, enlarging the input projection
from $J$ to $K$.  The remaining block has output $\wedge^2P$;
by taking the adjoint and omitting its input projection, it is
bounded by $\|M_fP^{\wedge2}\|_{\mathrm{HS}}$.  This is uniformly
finite since the kernel of $P$ is uniformly bounded on $\D^2$
and $f^2$ is integrable there.  Passing between tensor and antisymmetric tensor
norms costs only an absolute factor.  The uniform analytic norm
bounds on that basis give $C_A$ for the sum of all four squared
norms.  This proves~\eqref{eq:beta2-second-response}.

The identity
$\partial_\beta^2\rho=\rho\{\partial_\beta^2\log\rho+
(\partial_\beta\log\rho)^2\}$, the first-response estimates,
and boundedness of the Bergman determinants on $A$ prove the
last assertion.
\end{proof}

\begin{remark}[The remaining analytic step]
The preceding theorem bounds the first two temperature derivatives
at two after the bulk contributions have canceled.  It does not
bound derivatives on a neighborhood of two, and it does not give
factorial bounds at all orders.  Either a uniform remainder estimate
on such a neighborhood or a convergent connected expansion would
be needed to deduce convergence at a fixed nearby temperature.
In particular, this estimate alone supplies no additional
fixed-temperature uniqueness range.
\end{remark}

\section{Global temperature singularity and the failure of direct interpolation}
\label{sec:global-temperature-singularity}

The following conclusion concerns the full finite-particle laws.  It
does not assert singularity of their restrictions to a compact disk,
or of their limiting interior processes.

\begin{theorem}[Different temperatures are globally distinguishable]
\label{thm:global-temperature-singularity}
Let $\beta_0,\beta_1>0$ be distinct.  Then
\begin{equation}
 \|\mathsf P_{N,\beta_1}-\mathsf P_{N,\beta_0}\|_{\mathrm{TV}}
 \longrightarrow1,
 \label{eq:global-temperature-tv}
\end{equation}
where total variation is the supremum of differences of event
probabilities.  The same assertion holds for the origin-charged laws
\[
 d\widehat{\mathsf P}_{n,\beta}(y)
 =\widehat Z_{n,\beta}^{-1}|\Delta_n(y)|^\beta
                  \prod_{j=1}^n|y_j|^\beta\,dA(y_j).
\]
For either family let $L_N$ be the likelihood of temperature
$\beta_1$ with respect to temperature $\beta_0$.  Then
\begin{equation}
 L_N\longrightarrow0\quad\text{in probability under temperature }
 \beta_0,
 \qquad
 \E_{N,\beta_0}L_N^p\longrightarrow\infty
 \quad\text{for every }p>1.
 \label{eq:global-temperature-likelihood}
\end{equation}
The likelihood nevertheless has expectation one for every $N$.
\end{theorem}

\begin{proof}
For each fixed $m$, the already proved harmonic-mode theorem gives
\[
 W_{N,k}:=\sqrt{\frac{\beta_0}{2k}}\,S_{k,N},\qquad 1\le k\le m,
\]
converging under temperature $\beta_0$ to independent standard
circular complex Gaussians, and under temperature $\beta_1$ to
independent circular complex Gaussians of variance
$r=\beta_0/\beta_1\ne1$.  For the charged family the same statement
follows from Theorem~\ref{thm:inserted-charge-screening}: the charge
is at zero, so every limiting harmonic mean is zero.

Put $T_{N,m}=m^{-1}\sum_{k=1}^m|W_{N,k}|^2$ and
$d=|1-r|/3$.  Under the two limiting Gaussian laws this statistic
has respective means $1,r$ and variances $1/m,r^2/m$.
Chebyshev's inequality shows that the event
$|T_{N,m}-r|<d$, after $N\to\infty$ with $m$ fixed, has probability
at least $1-r^2/(md^2)$ under the second temperature and at most
$1/(4md^2)$ under the first.  The interval endpoints have zero
probability under both limiting laws, so convergence in distribution
indeed gives those event probabilities.  Taking first $N\to\infty$
and then $m\to\infty$ proves~\eqref{eq:global-temperature-tv}.

For equivalent probability measures with likelihood $L$,
$\E\min(1,L)=1-\|L\mathsf P-\mathsf P\|_{\mathrm{TV}}$.
Consequently
\[
 \mathsf P_{N,\beta_0}(L_N>\varepsilon)
 \le\frac{\E_{N,\beta_0}\min(1,L_N)}{\min(1,\varepsilon)}
 \longrightarrow0.
\]
If the $p$-th moments were bounded along an infinite subsequence,
they would make the likelihoods uniformly integrable on that
subsequence.  Convergence to zero in probability would then force
their expectations to tend to zero, contradicting expectation one.
Thus every $p$-th moment diverges along the full sequence.
\end{proof}

\begin{proposition}[An explicit finite-mode lower bound]
\label{prop:global-temperature-renyi}
Suppose $b=\beta_1/\beta_0>1$ and $p>1$.  For either family in
Theorem~\ref{thm:global-temperature-singularity}, and every integer
$m\ge1$,
\begin{equation}
 \liminf_{N\to\infty}\E_{N,\beta_0}L_N^p
 \ge\left(\frac{b^p}{1+p(b-1)}\right)^m.
 \label{eq:global-temperature-renyi}
\end{equation}
The quantity in parentheses is strictly greater than one.
\end{proposition}

\begin{proof}
Let $q=p/(p-1)$.  For any bounded continuous nonnegative function
$h$ of the first $m$ normalized harmonic modes, H\"older's
inequality gives
\[
 \E_{N,\beta_0}L_N^p\ge
 \frac{(\E_{N,\beta_1}h)^p}
      {(\E_{N,\beta_0}h^q)^{p-1}}.
\]
The limiting likelihood of the second Gaussian vector with respect
to the first is
\[
 \ell_m(w)=b^m\exp\left(-(b-1)\sum_{k=1}^m|w_k|^2\right).
\]
Choose $h=\ell_m^{p-1}$, which is bounded and continuous because
$b>1$, and pass to the Gaussian limits in the displayed inequality.
The resulting ratio equals
\[
 \E_{\mathrm{CN}(0,I_m)}\ell_m^p
 =\left(\frac{b^p}{1+p(b-1)}\right)^m.
\]
Strict convexity of $x\mapsto x^p$, applied to the nonconstant
mean-one likelihood $\ell_1$, proves strict inequality over one.
\end{proof}

\begin{corollary}[The exact obstruction in a global H\"older bound]
\label{cor:global-temperature-holder-obstruction}
Fix $\beta>2$, $p>1$, $q=p/(p-1)$, and set
$\widehat A_n(t)=\log\widehat Z_{n,t}$.  The likelihood factor in
the direct interpolation from the charged temperature-two law is
\begin{equation}
 \mathcal C_{n,\beta,p}
 =\exp\left\{\frac{
 \widehat A_n(2+p(\beta-2))+(p-1)\widehat A_n(2)
                     -p\widehat A_n(\beta)}p\right\}
 \longrightarrow\infty.
 \label{eq:global-holder-free-energy-cost}
\end{equation}
In particular the otherwise valid estimate
\[
 \widehat\E_{n,\beta}|Q_y(z)|^\beta
 \le\mathcal C_{n,\beta,p}
       \left(\widehat\E_{n,2}|Q_y(z)|^{\beta q}\right)^{1/q},
 \qquad Q_y(z)=\prod_j(1-z\bar y_j),
\]
does not supply an $n$-uniform moment bound with finite amplitude.
More strongly, for each fixed $0<|z|<1$,
\begin{equation}
 \inf_{p>1}\left\{
 \|\widehat L_{n,\beta}\|_{L^p(\widehat{\mathsf P}_{n,2})}
 \big\||Q_y(z)|^\beta\big\|_{L^{p/(p-1)}(\widehat{\mathsf P}_{n,2})}
 \right\}\longrightarrow\infty.
 \label{eq:global-holder-adaptive-failure}
\end{equation}
Thus even optimizing the H\"older exponent separately for every
$n$ fails to give a uniform bound at a single nonzero interior point.
\end{corollary}

\begin{proof}
The exact likelihood under the charged law is
\[
 \widehat L_{n,\beta}
 =\frac{\widehat Z_{n,2}}{\widehat Z_{n,\beta}}
   \exp\left((\beta-2)
       \left\{\log|\Delta_n(y)|+\sum_j\log|y_j|\right\}\right).
\]
Thus $\mathcal C_{n,\beta,p}
=(\widehat\E_{n,2}\widehat L_{n,\beta}^p)^{1/p}$ exactly.
The theorem proves divergence.  H\"older gives the stated moment
inequality.  To prove the stronger adaptive assertion, note that
for fixed $s>0$ and $|z|<1$, the charged harmonic exponential-moment
theorem gives
\[
 \lim_{n\to\infty}\widehat\E_{n,2}|Q_y(z)|^s
 =(1-|z|^2)^{-s^2/4}.
\]
Indeed apply~\eqref{eq:inserted-charge-laplace} with the analytic
function $h(u)=(s/2)\log(1-\bar zu)$; its value at zero is zero.
It follows, for each fixed $q_0>1$, that
\[
 \lim_{n\to\infty}\big\||Q_y(z)|^\beta\big\|_{L^{q_0}}
 =(1-|z|^2)^{-\beta^2q_0/4}.
\]
For $0<|z|<1$ these limits grow without bound as $q_0$ increases.
Given a target bound, choose $q_0$ so that this limit exceeds it,
and put $p_0=q_0/(q_0-1)$.  If $1<p\le p_0$, monotonicity of
norms makes the second H\"older factor at least its $q_0$-norm,
while the likelihood norm is at least one.  If $p\ge p_0$, its
likelihood factor is at least the divergent fixed $p_0$-norm,
while the second factor is at least its $L^1$ norm, whose limit
is the positive number $(1-|z|^2)^{-\beta^2/4}$.  This proves
\eqref{eq:global-holder-adaptive-failure}, without applying a
fixed-exponent limit theorem to moving exponents.
\end{proof}

\begin{remark}[Scope of the obstruction]
Different limiting variances in infinitely many boundary harmonic
modes cause the global singularity.  Those modes are independent
of every interior limit already constructed.  The theorem is
therefore compatible with local absolute continuity, with local
moment bounds, and with full convergence on an interval of
temperatures.  It rules out the particular use of a globally
uniformly integrable change of temperature.  A successful
perturbation would have to integrate out or renormalize this
boundary likelihood before estimating its local effect.
\end{remark}

\section{The hyperbolic spectral gap and the boundary modes it does not remove}
\label{sec:hyperbolic-operator-uniqueness}

Failure of absolute integrability is not the whole story for the Green
interaction.  Its integral operator is bounded on a natural Hilbert
space.  The following results establish that fact, identify an exact
obstruction to using it as a uniqueness proof, and give a separate
uniqueness statement for the choice of renormalized potential.
Throughout this section,
\[
 g(z,w)=-\log\left|\frac{z-w}{1-z\overline w}\right|,
 \qquad d\lambda(z)=\frac{dA(z)}{(1-|z|^2)^2},
 \qquad (Gf)(z)=\int_{\mathbb D}g(z,w)f(w)\,d\lambda(w).
\]

\subsection{A bounded Green operator despite the divergent row integral}

\begin{proposition}[The sharp $L^2$ Green bound]
\label{prop:hyperbolic-green-operator-bound}
The operator $G$, initially defined on bounded compactly supported
functions, extends to a bounded positive self-adjoint operator on
$L^2(\lambda)$, with
\begin{equation}
 \|G\|_{L^2(\lambda)\to L^2(\lambda)}=2\pi.
 \label{eq:hyperbolic-green-sharp-norm}
\end{equation}
For every $\beta>0$, the operator with kernel
$k_\beta(z,w)=1-e^{-\beta g(z,w)}$ is also bounded on this space and
satisfies
\begin{equation}
 \|K_\beta\|_{L^2(\lambda)\to L^2(\lambda)}\le2\pi\beta.
 \label{eq:hyperbolic-mayer-operator-bound}
\end{equation}
The second assertion is an operator norm bound; it does not assert
positive definiteness of $k_\beta$.
\end{proposition}

\begin{proof}
We first prove the sharp inequality
\begin{equation}
 \int_{\mathbb D}|\nabla u|^2\,dA
 \ge \int_{\mathbb D}|u|^2\,d\lambda,
 \qquad u\in C_c^\infty(\mathbb D).
 \label{eq:hyperbolic-hardy-gap}
\end{equation}
Under a conformal map to the upper half-plane
$\mathbb H=\{x+iy:y>0\}$, Dirichlet energy is unchanged and
$d\lambda=dx\,dy/(4y^2)$.  The identity
\[
 0\le\int_0^\infty
       \left|\partial_y v-\frac{v}{2y}\right|^2dy
   =\int_0^\infty|\partial_y v|^2dy
      -\frac14\int_0^\infty\frac{|v|^2}{y^2}dy
\]
holds for compactly supported smooth $v$; integration by parts proves
the equality, for real or complex $v$.  Integrating in $x$ and adding
the nonnegative horizontal derivative term proves
\eqref{eq:hyperbolic-hardy-gap}.

Here the constant one is optimal.  Choose nonzero real functions
$\chi,\psi\in C_c^\infty((-1,1))$, put $A_L=e^{2L}$, and use the
upper-half-plane test functions
\[
 v_L(x,y)=\chi(x/A_L)y^{1/2}\psi((\log y)/L).
\]
Writing $C_\chi=\int\chi^2$, $C_\psi=\int\psi^2$, and using the same
notation with primes for the corresponding derivative integrals gives
\[
 \|v_L\|_{L^2(\lambda)}^2
       =\frac{A_LL}{4}C_\chi C_\psi,
 \qquad
 \int|\partial_yv_L|^2dx\,dy
       =A_LLC_\chi\left(\frac{C_\psi}{4}
                                      +\frac{C_{\psi'}}{L^2}\right).
\]
The mixed term in the latter equality is zero because
$\int\psi\psi'=0$.  Moreover,
\[
 \int|\partial_xv_L|^2dx\,dy
   =\frac{C_{\chi'}}{A_L}
       \int_0^\infty y\psi((\log y)/L)^2dy
   \le\frac{C_{\chi'}}{A_L}Le^{2L}C_\psi.
\]
Thus their Rayleigh quotients tend to one.

Let $A$ be the positive self-adjoint operator associated with the
closure of the Dirichlet form
$\mathcal E(u,u)=\int|\nabla u|^2dA$ on $L^2(\lambda)$.  The
inequality just proved says $A\ge I$, and the test functions say that
the lower endpoint of its spectrum is one.  Consequently its bounded
inverse has norm one.  We now identify $2\pi A^{-1}$ with the
displayed Green kernel, so no separate boundary condition is being
assumed in this argument.

For a nonnegative smooth compactly supported $f$, let $u_R$ be the
zero-boundary solution on $\{|z|<R\}$ of
\[
 -\Delta u_R=2\pi f(1-|z|^2)^{-2},
 \qquad
 u_R(z)=\int_{|w|<R}g_R(z,w)f(w)\,d\lambda(w),
\]
where $R$ contains the support of $f$ and $g_R$ is the disk Green
function.  Its extension by zero belongs to the form domain.
Testing the equation with $u_R$ and applying
\eqref{eq:hyperbolic-hardy-gap} gives
\[
 \|u_R\|_2^2\le\mathcal E(u_R,u_R)
       =2\pi\langle f,u_R\rangle
       \le2\pi\|f\|_2\|u_R\|_2.
\]
The Green kernels increase to $g$ as $R\uparrow1$.  Monotone
convergence and the preceding bounds show that $u_R$ tends in
$L^2(\lambda)$ to $Gf$.  Weak compactness in the form domain and the
weak equation identify $Gf=2\pi A^{-1}f$.  Taking differences,
then using density, proves this identity for the extended operator.
The spectral conclusion yields \eqref{eq:hyperbolic-green-sharp-norm}.

Finally, $0\le k_\beta\le\beta g$ pointwise.  For a compactly
supported $f$,
\[
 |K_\beta f|\le\beta G|f|,
 \qquad \|K_\beta f\|_2\le2\pi\beta\|f\|_2.
\]
Density proves \eqref{eq:hyperbolic-mayer-operator-bound}.
\end{proof}

\subsection{Exact positive solutions outside the contraction space}

\begin{proposition}[Boundary modes survive strict Hilbert-space contraction]
\label{prop:hyperbolic-boundary-eigenmodes}
For $\xi\in\partial\mathbb D$ and $0<s<1$, define
\[
 u_{s,\xi}(z)
       =\left(\frac{1-|z|^2}{|\xi-z|^2}\right)^s.
\]
For every $z\in\mathbb D$ the following integral is finite, and
\begin{equation}
 Gu_{s,\xi}(z)=\frac{\pi}{2s(1-s)}u_{s,\xi}(z).
 \label{eq:hyperbolic-positive-boundary-eigenfunction}
\end{equation}
None of these functions belongs to $L^2(\lambda)$.  In particular,
for every $0<t<1/(2\pi)$ there are nonzero positive smooth solutions
of $u=tGu$, although $\|tG\|_{2\to2}<1$.
\end{proposition}

\begin{proof}
The conformal map from the disk to $\mathbb H$ which sends $\xi$ to
infinity identifies the Poisson kernel
$(1-|z|^2)/|\xi-z|^2$ with the height coordinate $y$, up to a
positive constant.  The Laplace--Beltrami operator for $\lambda$ is
$(1-|z|^2)^2\Delta$ in the disk and $4y^2\Delta$ on $\mathbb H$.
Since $4y^2\Delta y^s=-4s(1-s)y^s$, it follows that
\[
 -\Delta u_{s,\xi}
    =4s(1-s)u_{s,\xi}(1-|z|^2)^{-2}.
\]
Also $y^s\notin L^2(dx\,dy/(4y^2))$, already because its integral
over $x\in\mathbb R$ and $1<y<2$ is infinite.  This proves the
nonintegrability assertion.

We justify the Green representation, including the absence of an
additional harmonic boundary term.  Rotating, take $\xi=1$, and write
$\varepsilon=1-R$.  For $R\ge1/2$ and $-\pi\le\theta\le\pi$,
\[
 \frac{1-R^2}{|1-Re^{i\theta}|^2}
       \le C\frac{\varepsilon}{\varepsilon^2+\theta^2}.
\]
Therefore the circular average of $u_{s,1}(Re^{i\theta})$ is at
most a constant times
\[
 \varepsilon^s\int_0^\pi
             (\varepsilon^2+\theta^2)^{-s}d\theta.
\]
This tends to zero: it is $O(\varepsilon^s)$ if $s<1/2$,
$O(\varepsilon^{1/2}\log(1/\varepsilon))$ if $s=1/2$, and
$O(\varepsilon^{1-s})$ if $s>1/2$.

The Green representation on the disk of radius $R>|z|$ gives
\[
 u_{s,1}(z)
  =h_R(z)+\frac{4s(1-s)}{2\pi}
        \int_{|w|<R}g_R(z,w)u_{s,1}(w)\,d\lambda(w),
\]
where $h_R$ is the harmonic extension of the boundary values of
$u_{s,1}$.  For fixed $z$, the normalized Poisson kernel of that
disk is bounded uniformly as $R\uparrow1$.  The circular-average
estimate implies $h_R(z)\to0$.  Monotone convergence of $g_R$
now proves both finiteness of the integral and
\eqref{eq:hyperbolic-positive-boundary-eigenfunction}.
For the final assertion choose either root in $(0,1)$ of
$4s(1-s)=2\pi t$.
\end{proof}

This is a concrete obstruction to a proposed proof, not a
nonuniqueness theorem for the gas.  An $L^2$ operator contraction
does rule out $L^2$ differences.  It cannot rule out the displayed
positive boundary modes.  In particular, a proof based on such a
contraction must first place the actual differences of Gibbs states
in the contraction space, or establish a boundary condition which
excludes these modes.  A homogeneous nonzero intensity difference
is itself outside $L^2(\lambda)$ because $\lambda(\mathbb D)=\infty$.

There is a stronger issue: finite variance at each point, even together
with full isometry invariance, does not supply the spatial $L^2$
hypothesis needed for the contraction.

\begin{proposition}[Invariant random modes below the spatial spectral gap]
\label{prop:invariant-subgap-gaussian-fields}
For every $0<\kappa<1$, there is a real smooth centered Gaussian
field $U$ on the disk such that
\begin{enumerate}
\item its law is invariant under every hyperbolic isometry and
$\E U(z)^2=1$;
\item $-(1-|z|^2)^2\Delta U=\kappa U$ almost surely;
\item for every fixed $z$,
\begin{equation}
 \lim_{R\uparrow1}\int_{|w|<R}g(z,w)U(w)\,d\lambda(w)
       =\frac{2\pi}{\kappa}U(z)
       \quad\hbox{in }L^2\hbox{ of probability}.
 \label{eq:invariant-subgap-green-limit}
\end{equation}
\end{enumerate}
Thus these fields are nonzero solutions of $U=tGU$ in the displayed
truncated, mean-square sense, where $t=\kappa/(2\pi)$ and
$\|tG\|_{L^2(\lambda)\to L^2(\lambda)}=\kappa<1$.
\end{proposition}

\begin{proof}
Choose $1/2<s<1$ with $\kappa=4s(1-s)$, put
$\alpha=2(1-s)\in(0,1)$, and let $\sigma$ be normalized arc length
on the unit circle.  The kernel
\[
 k(\xi,\eta)=|\xi-\eta|^{-\alpha}
\]
is integrable on the product circle and is positive definite.
For the latter assertion use, off the diagonal,
\[
 |\xi-\eta|^{-\alpha}
  =\frac1{\Gamma(\alpha/2)}\int_0^\infty
        t^{\alpha/2-1}e^{-t|\xi-\eta|^2}\,dt.
\]
Each Gaussian kernel in this integral is positive definite: writing
circle points as vectors in $\mathbb R^2$, expand
$e^{-t|\xi-\eta|^2}=e^{-2t}\sum_{n\ge0}(2t)^n
(\xi\cdot\eta)^n/n!$, and each summand is the Gram kernel of
a tensor power.  For bounded functions, absolute integrability of
$k$ justifies integration of the quadratic forms.  This proves the
claim, with the diagonal irrelevant to the arc-length integrals.

Let
\[
 J=\iint k(\xi,\eta)\,d\sigma(\xi)d\sigma(\eta)\in(0,\infty),
 \qquad
 [f,h]_s=J^{-1}\iint f(\xi)h(\eta)k(\xi,\eta)
                          \,d\sigma(\xi)d\sigma(\eta).
\]
Complete the continuous real functions modulo their zero-norm
subspace in this positive semidefinite inner product.  There is a
centered real Gaussian linear map $W$ on the resulting separable
Hilbert space with $\E[W(f)W(h)]=[f,h]_s$.  Explicitly, choose an
orthonormal basis $(e_n)$ and independent standard real Gaussians
$(Z_n)$ and set $W(f)=\sum_n[f,e_n]_sZ_n$, with convergence in
mean square.  Define
\[
 P_z(\xi)=\frac{1-|z|^2}{|\xi-z|^2},
 \qquad U(z)=W(P_z^s).
\]
In particular $\E U(0)^2=[1,1]_s=1$.

We verify invariance directly.  If $\varphi$ is a disk automorphism
and $j(\xi)=|\varphi'(\xi)|$ is its boundary Jacobian, then
\[
 P_{\varphi(z)}(\varphi(\xi))=\frac{P_z(\xi)}{j(\xi)},
 \qquad
 |\varphi(\xi)-\varphi(\eta)|
       =\sqrt{j(\xi)j(\eta)}|\xi-\eta|.
\]
Changing variables in $[P_{\varphi(z)}^s,P_{\varphi(w)}^s]_s$
therefore gives a factor
$j(\xi)^{1-s-\alpha/2}j(\eta)^{1-s-\alpha/2}=1$.
The covariance is invariant, so the centered Gaussian law is
invariant.  Complex conjugation preserves the same covariance and
provides the orientation-reversing isometries.  In particular all
pointwise variances equal one.

For each compact disk, the functions $P_z(\xi)^s$ and their
$z$-derivatives of every order are bounded uniformly in $\xi$.
Because $k$ is integrable, their Hilbert-space derivatives are
bounded as well.  Gaussian moment estimates and the continuity
criterion, applied to the corresponding derivative fields on a
countable exhaustion, give a smooth version of $U$.  Equivalently,
one can first obtain a continuous version from
$\E|U(z)-U(w)|^{2m}\le C_{K,m}|z-w|^{2m}$ on compact sets and
then use the weak equation below and interior elliptic regularity.
The calculation in
Proposition~\ref{prop:hyperbolic-boundary-eigenmodes} shows that
$-(1-|z|^2)^2\Delta P_z(\xi)^s=\kappa P_z(\xi)^s$.
It may be passed through $W$ against smooth compactly supported
test functions, since $W$ is continuous into $L^2$ of probability.
On a countable dense set of test functions it holds on one
probability-one event, and then on all test functions by continuity.
This proves the asserted almost sure differential equation.

Finally, for fixed $z$ define the boundary functions
\[
 h_R(\xi)=\int_{|w|<R}g(z,w)P_w(\xi)^s\,d\lambda(w).
\]
By Proposition~\ref{prop:hyperbolic-boundary-eigenmodes},
\[
 0\le h_R(\xi)\uparrow\frac{2\pi}{\kappa}P_z(\xi)^s.
\]
The limit is bounded uniformly on the boundary because $z$ is
fixed.  Dominated convergence in the double integral defining
$[\cdot,\cdot]_s$ shows convergence to this limit in the boundary
Hilbert space.  The logarithmic pole of $g(z,w)$ is integrable on
compact sets, so stochastic Fubini, or approximation by finite
sums in the Hilbert space, gives
\[
 W(h_R)=\int_{|w|<R}g(z,w)U(w)\,d\lambda(w).
\]
The continuity of $W$ now proves
\eqref{eq:invariant-subgap-green-limit}.
\end{proof}

The proposition does not construct two Gibbs states: its charge
source is a smooth Gaussian field rather than a point measure.
It proves a narrower but useful negative assertion.  Spatial
$L^2(\lambda)$ and the space of invariant random fields with finite
one-point variance are different spaces, and the spectral-gap
contraction cannot be transferred from the former to the latter.

\subsection{Uniqueness of an invariant renormalized field up to a constant}

\begin{lemma}[An integrable invariant harmonic field is spatially constant]
\label{lem:integrable-invariant-harmonic-field}
Let $U$ be a random real harmonic function on $\mathbb D$.  Suppose
its law is invariant under precomposition by every disk automorphism,
and $\E|U(0)|<\infty$.  Then almost surely $U$ is a constant function.
The constant may be random.
\end{lemma}

\begin{proof}
Let $\Phi(t)=\sqrt{1+t^2}$.  Transitivity and invariance imply that
$\E\Phi(U(z))$ is a finite constant independent of $z$.  On each
sample, harmonicity gives
\[
 \Delta\Phi(U)=\Phi''(U)|\nabla U|^2\ge0,
 \qquad \Phi''(t)=(1+t^2)^{-3/2}>0.
\]
For nonnegative $\chi\in C_c^\infty(\mathbb D)$, Fubini is justified
by the constant finite expectation and gives
\[
 0\le\E\langle\Delta\Phi(U),\chi\rangle
     =\int\E\Phi(U(z))\Delta\chi(z)\,dA(z)=0.
\]
Using a countable collection of such functions positive on a compact
exhaustion proves $\Delta\Phi(U)=0$ everywhere almost surely.
The displayed identity and strict positivity of $\Phi''$ imply
$\nabla U=0$.  Connectedness completes the proof.
\end{proof}

\begin{proposition}[Ambiguity of an integrable covariant Green field]
\label{prop:covariant-green-field-ambiguity}
Let $P$ be an isometry-invariant point-process law.  On a measurable
configuration class $\Omega$ of full $P$-measure, stable under finite
additions and deletions, suppose $H_1(a,X),H_2(a,X)$ are measurable
real functions off the points of $X$, define locally integrable
distributions in $a$, and are taken in their smooth off-pole
representatives.  Suppose they satisfy all the following conditions:
\begin{enumerate}
\item They obey the same distributional equation
\[
 \Delta H_j=2\pi(c\lambda-\Xi),\qquad j=1,2,
\]
for every $X\in\Omega$, where measures are identified with
distributions relative to $dA$.  Their difference has a measurable
harmonic representative $U_X$ on the whole disk for every
$X\in\Omega$.
\item Their joint law is covariant under isometries, and their
difference is represented by a random harmonic function $U$ with
$\E|U(0)|<\infty$.
\item On $\Omega$, their finite-addition identities hold pointwise
off the poles:
\[
 H_j(a,X\cup\{b\})=H_j(a,X)+g(a,b),\qquad a\ne b.
\]
\end{enumerate}
Then $H_1(a,X)-H_2(a,X)=C(X)$ almost surely, independent of $a$.
The random variable $C$ is measurable with respect to the tail
$\sigma$-field of $X$.  In particular, if that tail is trivial,
$C$ is deterministic.  The two resulting canonical specifications
are the same, wherever their finite-volume normalizing constants
are finite.  Their grand-canonical specifications agree after
replacing the activity $z$ for $H_1$ by $z e^{-\beta C}$ for $H_2$.
\end{proposition}

\begin{proof}
The first two hypotheses and
Lemma~\ref{lem:integrable-invariant-harmonic-field} make $U$ spatially
constant.  The third makes its harmonic representative unchanged by
any finite modification: subtract the identities off their finitely
many poles, then extend the difference harmonically across those
points.  For every relatively compact $B$, deleting the finitely
many points in $B$ therefore gives
$C(X)=C(X\vert_{B^c})$.  The right side is measurable in the
exterior configuration.  A countable exhaustion proves tail
measurability.  For a fixed inside count $m$, the change of potential
multiplies every unnormalized weight by the same factor
$e^{-\beta mC}$, which cancels on normalization.  Across all
counts it changes $z^m$ to $(ze^{-\beta C})^m$, as claimed.
\end{proof}

\begin{remark}[What this uniqueness statement does and does not identify]
The last proposition concerns two fields defined over the same
invariant point-process law.  It does not compare two different
point-process laws, prove tail triviality, or establish existence of
the compensated field for the original gas.  Its value is to isolate
the remaining ambiguity: under the stated integrability and
covariance assumptions it is a tail-dependent activity, rather than
an arbitrary harmonic function of the insertion point.
\end{remark}

\section{A direct test for different subsequential limits}
\label{sec:boundary-insertion-nonuniqueness}

The existence of several solutions of a limiting DLR specification is
different from the existence of several subsequential limits of the
particular finite-volume sequence in \eqref{eq:model}.  This section
gives exact tests for the latter question.  It also identifies the
additional estimate needed to exclude a simple alternation between two
different limits.  No such estimate is asserted without proof.

Write $P_n$ for the law of $\Xi_n$.  We use the boundary insertion
weight $G_n$ of \eqref{transfer:boundary-weight} and put
$\widehat G_n=G_n/\mathbb E_nG_n$.  The exact outermost-particle
identity is already proved in Proposition~\ref{prop:outermost-transfer}.
In particular, for $f\in C_c^1(\D)$ with $f\geq0$, Proposition
\ref{prop:transfer-error} gives
\begin{equation}
 L_{n+1}(f)-L_n(f)=c_n(f)+e_n(f),\qquad
 \sum_{n\geq1}|e_n(f)|<\infty,
 \label{eq:boundary-covariance-summable-transfer}
\end{equation}
where
\[
 L_n(f)=\mathbb E_n e^{-\Xi_n(f)},\qquad
 c_n(f)=\mathbb E_n[e^{-\Xi_n(f)}(\widehat G_n-1)].
\]
These are explicit finite-dimensional integral ratios.  Namely, if
\begin{align*}
 \mathcal Z_n(f)&=\int_{\D^n}|\Delta_n(z)|^\beta
                    e^{-\sum_jf(z_j)}\,dA^{\otimes n}(z),\\
 \mathcal B_n(f)&=\frac1{2\pi}\int_0^{2\pi}\int_{\D^n}
       |\Delta_n(z)|^\beta e^{-\sum_jf(z_j)}
       \prod_j|e^{it}-z_j|^\beta\,dA^{\otimes n}(z)\,dt,
\end{align*}
then
\begin{equation}
 c_n(f)=\frac{\mathcal B_n(f)}{\mathcal B_n(0)}
              -\frac{\mathcal Z_n(f)}{\mathcal Z_n(0)}.
 \label{eq:explicit-boundary-covariance-ratios}
\end{equation}

\begin{proposition}[An exact oscillation criterion for nonuniqueness]
\label{prop:boundary-covariance-nonuniqueness}
Choose a countable collection $\mathcal F\subset C_c^1(\D)$ of
nonnegative functions whose Laplace functionals determine point-process
laws.  The following statements are equivalent:
\begin{enumerate}
\item the sequence $P_n$ has at least two distinct subsequential limits;
\item for some $f\in\mathcal F$, the bounded partial sums
      $\sum_{n=1}^{m}c_n(f)$ do not converge as $m\to\infty$;
\item for some $f\in\mathcal F$,
\begin{equation}
 \lim_{K\to\infty}\sup_{q\geq p\geq K}
           \left|\sum_{n=p}^{q}c_n(f)\right|>0.
 \label{eq:exact-nonuniqueness-covariance-tail}
\end{equation}
\end{enumerate}
For any test witnessing these equivalent statements, both
$\sum_n[c_n(f)]_+$ and $\sum_n[c_n(f)]_-$ diverge.
In particular, $\limsup_n|c_n(f)|>0$ for even one such $f$ is a
sufficient condition for two distinct subsequential limits.
\end{proposition}

\begin{proof}
Summing \eqref{eq:boundary-covariance-summable-transfer} shows that
$\sum_{n=1}^{m}c_n(f)$ differs from $L_{m+1}(f)$ by a convergent
sequence.  It is consequently bounded, and it converges if and only
if $L_m(f)$ converges.  Cauchy's criterion gives equivalence with
\eqref{eq:exact-nonuniqueness-covariance-tail}.  If every $L_m(f)$,
$f\in\mathcal F$, converges, tightness implies that all subsequential
limits have the same determining Laplace functionals, hence are equal.
Conversely, distinct subsequential limits differ on one test in
$\mathcal F$, so the corresponding $L_m(f)$ does not converge.

If a bounded sequence of partial sums has a finite sum of positive
parts or a finite sum of negative parts, the other sum is finite too:
otherwise the partial sums tend to minus or plus infinity.  Both sums
being finite would imply convergence.  This proves the divergence
assertion.  Finally, if $c_n(f)$ does not tend to zero, its partial
sums cannot converge.  Equivalently, select indices along which it
is bounded away from zero in absolute value; joint extraction of
$P_n$ and $P_{n+1}$ and the vanishing transfer error gives two
limits with different $f$-Laplace functionals.
\end{proof}

\begin{proposition}[All fixed regularizations are screened]
\label{prop:regularized-boundary-screening}
For $0\leq s<1$, set
\[
 G_{n,s}(z)=\frac1{2\pi}\int_0^{2\pi}
           \prod_{j=1}^n|1-se^{-it}z_j|^\beta\,dt,
 \qquad \widehat G_{n,s}=\frac{G_{n,s}}{\mathbb E_n G_{n,s}}.
\]
For each fixed $s<1$,
\begin{align}
 \mathbb E_n G_{n,s}&\longrightarrow(1-s^2)^{-\beta/2},
 \label{eq:regularized-boundary-mass-limit}\\
 \mathbb E_n[F(\Xi_n)(\widehat G_{n,s}-1)]&\longrightarrow0
 \label{eq:regularized-boundary-screening}
\end{align}
for every bounded continuous $F$.  If the additional estimate
\begin{equation}
 \lim_{s\uparrow1}\limsup_{n\to\infty}
    \mathbb E_n|\widehat G_n-\widehat G_{n,s}|=0
 \label{eq:unproved-boundary-mass-approximation}
\end{equation}
holds, then $P_n$ and $P_{n+1}$ have asymptotically zero distance
in every metric for weak convergence on their relatively compact
closure.  Estimate~\eqref{eq:unproved-boundary-mass-approximation}
is a sufficient hypothesis, not a proved statement about the gas.
\end{proposition}

\begin{proof}
For fixed $s<1$ and $t$, the function
$h_{s,t}(z)=(\beta/2)\log(1-se^{-it}z)$ is holomorphic on a
neighborhood of $\overline\D$ and vanishes at zero.  Its coefficients
are $-(\beta/2)s^k e^{-ikt}/k$.  Theorem
\ref{thm:analytic-harmonic-laplace} gives
\[
 \mathbb E_n e^{Y_n(h_{s,t})}
 \longrightarrow
 \exp\left(\frac\beta2\sum_{k\geq1}\frac{s^{2k}}k\right)
 =(1-s^2)^{-\beta/2}.
\]
Rotation invariance makes the expectation on the left independent of
$t$, proving \eqref{eq:regularized-boundary-mass-limit}.  Along an
arbitrary subsequence on which $\Xi_n\Rightarrow\Xi$, the mixed
Laplace identity in the same theorem gives the additional factor
$\mathbb EF(\Xi)$.  Integration in $t$ is legitimate: for fixed
$s$, the absolute value of each mixed expectation is bounded by
$\|F\|_\infty\mathbb E_n e^{Y_n(h_{s,0})}$, uniformly in $n,t$.
Dominated convergence proves
\eqref{eq:regularized-boundary-screening} along that subsequence.
Tightness and a subsequence contradiction prove it along the whole
sequence.

Finally, the difference between the unregularized and regularized
mixed expectations is bounded by
$\|F\|_\infty\mathbb E_n|\widehat G_n-\widehat G_{n,s}|$.
The additional hypothesis, followed by the fixed-$s$ result, gives
vanishing normalized covariance for every $F$.  Applying this to
compact Laplace functionals and using
\eqref{eq:boundary-covariance-summable-transfer} gives the asserted
adjacent-size conclusion by tightness.
\end{proof}

\begin{proposition}[A connectedness consequence of adjacent-size screening]
\label{prop:adjacent-screening-connected-limits}
Suppose that $P_n$ and $P_{n+1}$ have asymptotically zero weak distance.
Then the set of subsequential limits of $P_n$ is compact and connected.
If this set is not a singleton, it contains continuum many different
probability laws.  Thus under this hypothesis a two-phase parity
alternation is impossible, although full convergence does not follow.
\end{proposition}

\begin{proof}
The closure of the sequence is compact by tightness, and its limit
set $C$ is a nonempty compact set.  If $C=C_0\cup C_1$ were a
disconnection into nonempty disjoint compact sets, their distance
would be positive.  Choose disjoint neighborhoods of $C_0,C_1$
at positive distance.  All sufficiently late terms of the sequence
lie in their union, since otherwise a subsequence outside the union
would have a limit point outside $C$.  Both neighborhoods are visited
infinitely often.  Thus infinitely many consecutive terms cross from
one neighborhood to the other, contradicting the vanishing of their
distances.  This proves connectedness.

If $x,y\in C$ are distinct, the continuous function
$z\mapsto d(x,z)$ maps $C$ onto a connected subset of the real line
containing $0$ and $d(x,y)>0$, and hence onto a set containing an
interval.  Therefore $C$ has at least continuum cardinality.  A
separable metric space has at most continuum cardinality, which
proves the final assertion.
\end{proof}

\paragraph{What would constitute an actual disproof of uniqueness.}
A positive oscillation in
\eqref{eq:exact-nonuniqueness-covariance-tail} would be a disproof for the original disk gas, rather than merely
a multiplicity statement for an auxiliary DLR equation.  Convergence to zero of each individual covariance rules out adjacent-size phase
switches but still allows slow variation through a connected family
of limits.  The fixed-$s$ calculation does not permit interchange of
$n\to\infty$ and $s\uparrow1$.  The missing estimate concerns
fluctuations at boundary scales tending to zero with $n$; it is not
supplied by convergence of every fixed harmonic mode or by the
leading partition-function asymptotic.


\subsection{An independent sufficient monotonicity criterion}
\label{subsec:correlation-monotonicity-criterion}

For the fixed-potential planar gas, monotonicity of all correlation
measures in particle number is known at $\beta=2$; Hedenmalm
\cite[slide 6]{Hedenmalm2015} proposed the analogous statement for
$\beta<2$. We do not assume that proposal holds for the disk gas.
The following implication explains precisely why a disk-specific
proof would settle uniqueness below $2$.

\begin{proposition}
\label{prop:all-correlation-monotonicity-uniqueness}
Fix $\beta>0$. If, for every $p\ge1$, the correlation measures
$\rho_{p,N}dA^p$ are nondecreasing in $N\ge p$, then the disk gas
has a unique subsequential limit and hence converges along the full
sequence.
\end{proposition}
\begin{proof}
Uniform local bounds make the increasing correlation measures converge
to locally finite measures, uniquely for every $p$. Every subsequential
point-process limit has these measures by the already proved local
correlation convergence. They determine its law: for a compactly
supported nonnegative bounded function $f$, expand
$\prod_{z\in\Xi}(1+(e^{-f(z)}-1))$ into factorial sums. The absolute
sum is bounded by $2^{\Xi(\operatorname{supp}f)}$, which is integrable
by the established compact exponential moments. The resulting Laplace
functional is therefore uniquely specified by the factorial measures.
Laplace functionals determine point-process laws. Tightness then
upgrades uniqueness of every subsequential limit to full convergence.
\end{proof}

This is a sufficient criterion, not a monotonicity theorem for
$\beta\ne2$. The interpolation inequalities above bound correlations
from below; they do not compare successive particle numbers.


\section{An exact second temperature response at the center}
\label{sec:beta2-exact-second-center}

The intensity conjecture has a second, nontrivial perturbative check.
The order of limits in the following theorem is part of its statement.

\begin{theorem}[Exact second response at the center]
\label{thm:beta2-exact-second-center}
For the unweighted disk gas, with intensity measured against ordinary
area,
\begin{align}
 \lim_{N\to\infty}\left.\partial_\beta^2
       \log\rho_{N,\beta}^{(1)}(0)\right|_{\beta=2}&=0,
       \label{eq:beta2-exact-second-log-center}\\
 \lim_{N\to\infty}\left.\partial_\beta^2
       \rho_{N,\beta}^{(1)}(0)\right|_{\beta=2}&=\frac1\pi.
       \label{eq:beta2-exact-second-center}
\end{align}
Thus the first two finite-volume temperature responses at the center
agree, after passage to the limit, with the conjectured function
$(4/\beta-1)/\pi$.  No interchange of a temperature derivative with
an infinite-volume limit is asserted.
\end{theorem}

We give all cancellations, including a boundary contribution that
would be lost by simply replacing finite Bergman projections by the
infinite one.  For this proof use normalized area $d\mu=dA/\pi$ and
orthonormal functions $h_n(z)=\sqrt{n+1}\,z^n$.  Let $K_N$ project onto
their first $N$ terms, let $K$ be the full Bergman projection, and put
$Q_N=I-K_N$, $Q=I-K$, and $P=|h_0\rangle\langle h_0|$.
All operator norms and integrals in this section use $\mu$.
Changing the reference measure multiplies the intensity by $\pi$ and
does not change its logarithmic temperature derivatives.
Write $\zeta_2=\sum_{j\ge1}j^{-2}$ and
\[
 f(u,v)=\log|u-v|=I(u,v)-g(u,v),\qquad
 I(u,v)=\log|1-u\bar v|,\qquad
 g(u,v)=\log\left|\frac{1-u\bar v}{u-v}\right|.
\]
For $m<n$ put $\psi_{mn}=h_m\wedge h_n$, with the convention
$h_m\wedge h_n=(h_m\otimes h_n-h_n\otimes h_m)/\sqrt2$.

\begin{lemma}[The single-replacement contribution]
\label{lem:beta2-exact-center-single}
The difference between the single-replacement variances in the
Slater identity, for the Palm process at zero with statistic
$U_{N-1}+\sum\log|z_i|$ and for the original $N$-particle process,
converges to
\[
 S=\frac{47}{16}-\frac74\zeta_2.
\]
\end{lemma}

\begin{proof}
The Palm occupied space is $J_N=K_N-P$.  The screened perturbation
$B=M_{\log|z|-V_P}+A_P$ is independent of $N$.  For $n\ge1$, with
$t=|z|^2$, its action is
\[
 Bh_n=\left\{\frac12\log t+\frac12(1-1/n)
                    -\frac{n}{2(n+1)}t\right\}h_n.
\]
Indeed $V_P=(t-1)/2$ and
$A_Ph_n/h_n=-1/(2n)+t/[2(n+1)]$, by direct angular integration.
The Hartree--Fock multiplier of Lemma~\ref{lem:beta2-hartree-fock-bound},
after removing its scalar, converges for each fixed $n$ in
$L^2((n+1)t^n dt)$ to
\[
 b_n(t)=\tfrac12H_{n+1}-\log(1-t).
\]
The uniform bound for the centered multipliers in that lemma gives
strong convergence of the single-replacement operators on the
Bergman space.  The trace-norm column bound for $BK$ in
Lemma~\ref{lem:beta2-screened-palm-trace} then justifies passage to
the limit in their cross term.  More explicitly, approximate $BK$
by its first $M$ columns, pass to the limit for those columns, and
bound the remainder by the uniform operator norm times
$\sum_{n>M}\|Bh_n\|_2$.  The latter tends to zero.  The quadratic
$B$ term passes by the same argument in Hilbert--Schmidt norm.
The lost $h_0$ column passes by its $L^2$ convergence.  The
$P(F+B)J_N$ term is zero, since its input and output have different
angular frequencies.  Consequently the limit is
\[
 -1+\sum_{a=2}^{\infty}
 \left\{\operatorname{Var}_{a t^{a-1}dt}
  \left[-\log(1-t)+\tfrac12\log t-\tfrac{a-1}{2a}t\right]
       -\operatorname{Var}_{a t^{a-1}dt}[-\log(1-t)]\right\}.
\]
The series is absolutely convergent by the same trace-norm argument.
The elementary beta integral and its first two derivatives give
\[
 \operatorname{Var}(\log t)=a^{-2},\quad
 \operatorname{Var}(t)=\frac{a}{(a+1)^2(a+2)},\quad
 \operatorname{Cov}(\log t,t)=\frac1{(a+1)^2},
\]
\[
 \operatorname{Cov}(-\log(1-t),\log t)=\sum_{j>a}\frac1{j^2},
 \qquad
 \operatorname{Cov}(-\log(1-t),t)=\frac{a}{(a+1)^2}.
\]
The summand therefore equals
\[
 \sum_{j>a}\frac1{j^2}
 +\frac5{8a}+\frac1{4a^2}-\frac1{2(a+1)}-\frac9{8(a+2)}.
\]
For completeness,
\[
 \sum_{a=2}^M\sum_{j>a}j^{-2}
 =H_M-H_M^{(2)}+M\sum_{j>M}j^{-2}-\zeta_2+1,
\]
where $H_M^{(2)}=\sum_{j\le M}j^{-2}$ and
$M\sum_{j>M}j^{-2}\to1$ by integral comparison.
Summing the rational terms and then letting $M$ increase gives
$S=-1+(2-2\zeta_2)+(31/16+\zeta_2/4)$, as claimed.
\end{proof}

\begin{lemma}[The infinite Green replacement terms]
\label{lem:beta2-exact-center-green}
Define
\[
 L_g=\sum_{n\ge1}\|Q^{\wedge2}M_g\psi_{0n}\|_2^2,
 \qquad
 G_g=\|\operatorname{Proj}_{h_0\wedge Q\mathcal H}
                 M_g(K-P)^{\wedge2}\|_{\mathrm{HS}}^2.
\]
Both are finite and
\[
 L_g=\frac12-\frac14\zeta_2,\qquad
 G_g=\frac32\zeta_2-\frac{39}{16}.
\]
\end{lemma}

\begin{proof}
Set
\[
 A=\sum_{n\ge1}\|M_g\psi_{0n}\|_2^2,\quad
 B=\sum_{n\ge1}\|(K\otimes I)M_g\psi_{0n}\|_2^2,\quad
 C=\sum_{n\ge1}\|(K\otimes K)M_g\psi_{0n}\|_2^2.
\]
Since the vectors are antisymmetric, $L_g=A-2B+C$.
The Bergman kernel against $\mu$ is $(1-u\bar v)^{-2}$.  Summing
the nonnegative squared differences defining $A$ gives
\[
 A=\iint g(u,v)^2\{K(v,v)-\operatorname{Re}K(u,v)\}
                         d\mu(u)d\mu(v)=\frac{\zeta_2}{2}-\frac12.
\]
Here $\int g(u,v)^2 K(v,v)d\mu(v)=\zeta_2/2$ by disk
automorphisms and
$\tfrac14\int_0^1(\log t)^2(1-t)^{-2}dt=\zeta_2/2$.
For the other integral use $v=(u-w)/(1-\bar u w)$.
The product $K(u,v)d\mu(v)$ transforms to
$(1-\bar u w)^{-2}d\mu(w)$, whose angular average is one;
thus its integral against $\log^2|w|$ is $1/2$.
All integrations and limits are justified, for example, by
$|K(u,v)|\le\sqrt{K(u,u)K(v,v)}$ and Cauchy--Schwarz with
weight $g^2$, since each corresponding diagonal integral is finite.

To evaluate $B$, let $T_v=KM_{g(\cdot,v)}K$ on the Bergman space.
Conformal covariance makes $T_v$ unitarily equivalent to $T_0$,
and direct radial integration gives
\[
 T_0h_n=\frac{h_n}{2(n+1)},\qquad
 \|T_v\|_{\mathrm{HS}}^2=\frac14\zeta_2.
\]
The reproducing vector at $v$ is an eigenvector of $T_v$ with
eigenvalue $1/2$: a disk automorphism taking zero to $v$ transports
the constant eigenvector to a nonzero multiple of that vector.
Hence $(T_v^2h_0)(v)=1/4$.  Finally,
\[
 (T_vh_0)(u)=\frac{1-|v|^2}{2(1-u\bar v)},\qquad
 K(v,v)\|T_vh_0\|_2^2
       =\frac{-\log(1-|v|^2)}{4|v|^2},
\]
with the continuous interpretation at zero.  Expanding the squared
contractions in the first coordinate now gives
\begin{align*}
 B&=\frac12\int\{K(v,v)\|T_vh_0\|_2^2
       +\|T_v\|_{\mathrm{HS}}^2
       -2\operatorname{Re}(T_v^2h_0)(v)\}\,d\mu(v)\\
  &=\frac12\left(\frac{\zeta_2}{4}
                      +\frac{\zeta_2}{4}-\frac12\right)
    =\frac{\zeta_2-1}{4}.
\end{align*}
These are finite quantities, so the expansion is legitimate by
Parseval and Cauchy--Schwarz.

For $C$ and $G_g$ it is useful to record the elementary Green integral
\begin{equation}
 \int g(u,v)u^l\bar u^k\,d\mu(u)
 =\frac{H_{l-k}(v)-v^{l+1}\bar v^{k+1}}{2(l+1)(k+1)},
 \label{eq:beta2-monomial-green-integral}
\end{equation}
where $H_j(v)=v^j$ for $j\ge0$ and
$H_j(v)=\bar v^{-j}$ for $j<0$.
One proof applies $-\Delta_v$ to both sides and uses zero boundary
values; uniqueness follows from the maximum principle.  Alternatively
the angular expansion of $g$ gives the same expression directly.
For $k\ge0$, contraction of $g\psi_{0n}$ against $h_k$ in its
first coordinate is
\[
 D_{kn}(v)=\frac{(n+1-n|v|^2)v^n\bar v^k-H_{n-k}(v)}
                         {2\sqrt{2(k+1)(n+1)}}.
\]
Only $0\le k\le n$ survive the second Bergman projection.
Their squared norms after that projection sum to
\[
 C_n=\sum_{k=0}^n
 \frac{(n-2k)^2}{8(k+1)(n+1)(n-k+1)(n+2)^2}
 =\frac{H_{n+1}}{4(n+1)(n+2)}-\frac1{2(n+2)^2}.
\]
Using $\sum_{a\ge1}H_a/[a(a+1)]=\zeta_2$, which follows by
interchanging positive sums, gives $C=1/2-\zeta_2/4$.
Combining $A,B,C$ proves the formula for $L_g$.

For $1\le m<n$, contraction of $g\psi_{mn}$ against $h_0$ is
\[
 \frac{n-m}{2\sqrt{2(m+1)(n+1)}}v^{m+n}(1-|v|^2).
\]
Projecting off the analytic monomial of degree $m+n$, and accounting
for both antisymmetric coordinates, gives the squared norm
\[
 \frac{(n-m)^2}
 {4(m+1)(n+1)(m+n+2)^2(m+n+3)}.
\]
Summation with $s=m+n+2$ therefore yields
\begin{align*}
 G_g&=\sum_{s=4}^{\infty}
 \left\{\frac{H_{s-2}-1}{4s(s+1)}
                    -\frac{s-3}{2s^2(s+1)}\right\}\\
 &=\frac32\zeta_2-\frac{39}{16}.
\end{align*}
To check the last equality directly, the two required sums are
$\sum_{s\ge4}H_{s-2}/[s(s+1)]=2/3$ and
$\sum_{s\ge4}(s-3)/[s^2(s+1)]=61/12-3\zeta_2$;
both follow by telescoping and interchanging positive harmonic sums.
\end{proof}

\begin{lemma}[The finite-projection boundary contributions cancel]
\label{lem:beta2-exact-center-image}
The difference of the finite-$N$ double-replacement variances tends
to $G_g-L_g$.
\end{lemma}

\begin{proof}
The difference in question is the gain with input $J_N^{\wedge2}$
and output $h_0\wedge Q_N\mathcal H$, minus the loss with input
$h_0\wedge J_N\mathcal H$ and output $Q_N^{\wedge2}$.
The Green parts of these two operators converge in Hilbert--Schmidt
norm to the corresponding infinite-projection operators.  For the
loss, the unprojected input operator is Hilbert--Schmidt by the finite
value of $A$ above.  For the gain, even the unprojected contracted
operator
$(\langle h_0|\otimes I)M_g(K\otimes K)$ is Hilbert--Schmidt.
Indeed, Bessel's inequality gives the bound
\[
 \sum_{m,n\ge0}\left\|h_n(v)
       \int g(u,v)h_m(u)d\mu(u)\right\|_{L^2(d\mu(v))}^2
 \le\iint g(u,v)^2 K(v,v)d\mu(u)d\mu(v)<\infty.
\]
Strong convergence of the input and output projections applied to
these fixed Hilbert--Schmidt operators proves the assertion.

The image kernel contributes finite norms that must be retained.
Its angular series shows that the squared loss norm is exactly
\begin{align}
 L_{I,N}
 &=\frac14\sum_{n=1}^{N-1}\sum_{k\ge N-n}
 \frac{n+1}{n+k+1}\int_0^1t^k
                   \left(\frac1k-\frac{t^n}{n+k}\right)^2dt.
 \label{eq:beta2-exact-center-image-loss}
\end{align}
For example the output with angular frequencies $(-k,n+k)$
and its antisymmetric partner have radial coefficient
proportional to $\bar u^k v^{n+k}(1/k-|u|^{2n}/(n+k))$.
These are exactly the surviving terms: the analytic output index
must be at least $N$.  Different such angular pairs are orthogonal.
For each fixed $k$, the $k$ indices $n=N-k,\ldots,N-1$ give a
total tending to $1/[4k(k+1)]$.  Their total is always at most
$1/[4k(k+1)]$, since the integrand is bounded by $t^k/k^2$
and there are at most $k$ indices.  Dominated convergence gives
$L_{I,N}\to1/4$.

For the gain, contraction of $I\psi_{mn}$ against $h_0$, with
$1\le m<n<N$ and $s=m+n$, is the analytic function
\[
 -\frac{(n-m)(s+1)}{2\sqrt2\,mn\sqrt{(m+1)(n+1)}}v^s.
\]
It survives $Q_N$ exactly when $s\ge N$.  Thus
\begin{equation}
 G_{I,N}=
 \sum_{\substack{1\le m<n<N\\m+n\ge N}}
 \frac{(n-m)^2(m+n+1)}{4m^2n^2(m+1)(n+1)}.
 \label{eq:beta2-exact-center-image-gain}
\end{equation}
For each fixed $m$, the $m$ indices $n=N-m,\ldots,N-1$
give a total tending to $1/[4m(m+1)]$.  For all $N$, the
corresponding total is at most $1/[2m(m+1)]$: there are at
most $m$ indices, $(n-m)^2/n^2\le1$, and
$(m+n+1)/(n+1)\le2$.  Hence $G_{I,N}\to1/4$ as well.

Finally the Green--image cross terms tend to zero.  Regard the
gain and loss operators as operators between fixed tensor spaces
by including their projections.  Their image parts have bounded
Hilbert--Schmidt norm by the preceding two formulas and converge
weakly to zero in that space.  To see weak convergence, first test
against operators with finitely many analytic input columns and
finitely many output vectors.  In the loss, each fixed input column
has image norm tending to zero directly from
\eqref{eq:beta2-exact-center-image-loss}; in the gain each fixed
input pair is annihilated once $N>m+n$.  Finite-rank tests are
dense, and the uniform Hilbert--Schmidt bound completes the argument.
Their inner products with the strongly convergent Green parts
therefore vanish.  Since $f=I-g$, the limiting double-replacement
difference is $(G_g+1/4)-(L_g+1/4)=G_g-L_g$.
\end{proof}

\begin{proof}[Proof of Theorem~\ref{thm:beta2-exact-second-center}]
Differentiating the finite-volume Palm partition quotient twice,
as in Theorem~\ref{thm:beta2-second-response}, gives a difference
of the two energy variances.  The Slater identity splits it into
the single- and double-replacement terms evaluated above.  Their
sum tends to
\[
 \left(\frac{47}{16}-\frac74\zeta_2\right)
 +\left(\frac32\zeta_2-\frac{39}{16}
                      -\frac12+\frac14\zeta_2\right)=0.
\]
The original ordinary-area intensity satisfies $\rho_{N,2}(0)=1/\pi$.
The exact first-response formula gives
$\partial_\beta\log\rho_{N,\beta}(0)|_2\to-1$.
Using $\rho''=\rho\{(\log\rho)''+((\log\rho)')^2\}$ proves
\eqref{eq:beta2-exact-second-center}.
\end{proof}


\section{The exact second response at every interior anchor}
\label{sec:beta2-exact-second-all-anchors}

\begin{theorem}[Locally uniform second intensity response]
\label{thm:beta2-exact-second-all-anchors}
For every $0<r<1$,
\begin{align}
 \lim_{N\to\infty}\sup_{|a|\le r}
 \left|\left.\partial_\beta^2\log\rho_{N,\beta}^{(1)}(a)
                                      \right|_{\beta=2}\right|&=0,
 \label{eq:beta2-exact-second-all-log}\\
 \lim_{N\to\infty}\sup_{|a|\le r}
 \left|\left.\partial_\beta^2\rho_{N,\beta}^{(1)}(a)
 \right|_{\beta=2}-\frac1{\pi(1-|a|^2)^2}\right|&=0.
 \label{eq:beta2-exact-second-all-intensity}
\end{align}
These statements take the finite-volume derivatives before the
particle-number limit.
\end{theorem}

Continue with normalized area, projections, and kernels from
Section~\ref{sec:beta2-exact-second-center}.  Put
$L(z)=-\log(1-|z|^2)$, $U_k=M_{z^k}$,
$H_k=QM_{\bar z^k}K$, and $T_L=KM_LK$ on its natural domain.
The normalized coherent vector and its finite version are
\[
 h_a(z)=\frac{1-|a|^2}{(1-\bar a z)^2},\qquad
 h_{a,N}(z)=\frac{K_N(z,a)}{\sqrt{K_N(a,a)}}.
\]
Their projections are denoted by $P_a$ and $P_{a,N}$.

\begin{lemma}[An image-kernel boundary formula]
\label{lem:beta2-general-image-boundary}
Let $h_N\in K_N\mathcal H$ be unit vectors converging to a unit
analytic vector $h$ in the norm
$\|\sum_l c_lh_l\|_* =\sum_l|c_l|\sqrt{l+1}$.
Suppose $h$ extends holomorphically to a neighborhood of the closed
disk.  Write $P_N=|h_N\rangle\langle h_N|$, $J_N=K_N-P_N$,
$P=|h\rangle\langle h|$, and $J=K-P$.  Then the image-kernel loss
and gain satisfy
\begin{align}
 \|Q_N^{\wedge2}M_I(h_N\wedge K_N)\|_{\mathrm{HS}}^2
   &\longrightarrow\frac14,\label{eq:beta2-general-image-loss}\\
 \|\operatorname{Proj}_{h_N\wedge Q_N\mathcal H}
                        M_IJ_N^{\wedge2}\|_{\mathrm{HS}}^2
   &\longrightarrow
     \|\operatorname{Proj}_{h\wedge Q\mathcal H}
                        M_IJ^{\wedge2}\|_{\mathrm{HS}}^2
       +\frac14\left\{\int L|h|^2d\mu
              -\sum_{k\ge1}\frac{|\mu_k(h)|^2}{k}\right\},
       \label{eq:beta2-general-image-gain}
\end{align}
where $\mu_k(h)=\int z^k|h(z)|^2d\mu(z)$.  The first operator is
viewed as a map from $K_N\mathcal H$ by $v\mapsto h_N\wedge v$.
Its kernel contains $h_N$, so this has precisely the loss norm on
$h_N\wedge J_N\mathcal H$.

The image loss converges weakly to zero in Hilbert--Schmidt space.
The image gain is the sum of a strongly convergent operator, whose
limit is its displayed infinite-projection counterpart, and a
bounded operator converging weakly to zero.  For the coherent
vectors $h=h_a$, the infinite-projection gain is zero and the
additional term in braces equals one.
\end{lemma}

\begin{proof}
We first prove the claims for a fixed unit polynomial $h$ of degree
at most $d$ and $N>d$.  Define the analytic tail operator
$R_{k,N}=Q_NU_kK_N$.  Since multiplication by $\bar z^k$ lowers
the angular index,
\[
 Q_NM_{\bar z^k}K_N=H_kK_N,\qquad
 \|R_{k,N}\|_{\mathrm{HS}}^2
 =\sum_{n=\max(0,N-k)}^{N-1}\frac{n+1}{n+k+1}
 \le\min(k,N).
\]
For fixed $k$ this last quantity tends to $k$.
Moreover $H_k$ is Hilbert--Schmidt and $\|H_k\|_{\mathrm{HS}}^2=k$,
as follows by summing the squared column norms
\[
 \|H_kh_l\|_2^2=
 \begin{cases}
 k^2/[(l+1)(l+k+1)],&k\le l,\\
 (l+1)/(l+k+1),&k>l.
 \end{cases}
\]
The same formulas, summed first over $k$ for each fixed $l$, give
the bounded-operator identity
\begin{equation}
 \sum_{k\ge1}\frac{H_k^*H_k}{k}=I
 \quad\hbox{on the Bergman space}.
 \label{eq:beta2-hankel-resolution}
\end{equation}
Indeed the summands are diagonal in $h_l$ and their diagonal sum is
\[
 \sum_{k=1}^l\frac{k}{(l+1)(l+k+1)}
       +\sum_{k>l}\frac{l+1}{k(l+k+1)}=1.
\]

The image angular expansion is
$M_I=-\tfrac12\sum_{k\ge1}(U_k\otimes U_k^*
                         +U_k^*\otimes U_k)/k$.
The loss, apart from a term tending to zero, is therefore
\[
 -\frac12\sum_{k\ge1}\frac{(H_kh)\wedge(R_{k,N}K_N)}{k}.
\]
For fixed $k$ its squared Hilbert--Schmidt norm tends to
$\|H_kh\|_2^2/(4k)$.  Terms with different $k$ are orthogonal
on each input monomial: their high analytic output indices differ,
and their other output indices are at most $d<N$.
The sum of their squared norms is bounded by
$\tfrac14\sum_k\|H_kh\|_2^2/k$.
The omitted term contains $R_{k,N}h$ and can occur only when
$k\ge N-d$.  Its squared norm is $O_d(N^{-1})$.
For clarity, $\|U_kh\|_2^2\le C_h/(k+1)$ and
$\|H_kK_N\|_{\mathrm{HS}}^2\le\min(k,N)$; summing these
bounds against $k^{-2}$ over $k\ge N-d$ gives $O_h(N^{-1})$.
Cross terms in this omitted sum cost at most $2d+1$, because equal
high output angular indices require $|k-k'|\le d$.
Equation~\eqref{eq:beta2-hankel-resolution} proves the loss limit.
Each fixed input column tends to zero, so boundedness gives its
weak Hilbert--Schmidt convergence to zero.

For the gain first allow input $K_N^{\wedge2}$.  Identify the output
$h\wedge Q_N\mathcal H$ with $Q_N\mathcal H$ by contraction.
For an input $h_m\wedge h_n$, $m<n$, its output is
\begin{align*}
 -\frac12\sum_{k\ge1}\frac1k\{&
  a_k(m)H_kh_n-a_k(n)H_kh_m\\
 &+b_k(m)R_{k,N}h_n-b_k(n)R_{k,N}h_m\},
\end{align*}
where
$a_k(m)=\langle h,U_kh_m\rangle$ and
$b_k(m)=\langle h,U_k^*h_m\rangle$; use the inner product linear
in its second entry.  The $a$ part has $k\le d$ and converges
strongly in Hilbert--Schmidt norm, since each $H_k$ is
Hilbert--Schmidt.  It is the infinite-projection gain operator.

The $b$ part converges weakly to zero.  For finitely many $k$,
its nonzero low input indices satisfy $m\le d+k$, whereas the
tail output comes from $n\ge N-k$; these two ranges are disjoint
for sufficiently large $N$.  Orthogonality of the high output
angular indices then gives its limiting squared norm
\[
 \frac14\sum_{k\ge1}\frac{\|U_kh\|_2^2}{k}
       =\frac14\int L|h|^2d\mu.
\]
To justify the infinite sum, write $h$ in its $d+1$ monomials and
use Cauchy--Schwarz for that finite sum.  For a fixed monomial,
different output angular indices are orthogonal; the two terms
that interchange $m,n$ cost at most a factor two.  Hence the
squared norm of the part $k>K_0$ is at most
\[
 C_h\sum_{k>K_0}\frac{1}{k^2}
                 \frac{\min(k,N)}{k+1}
 \le C_h\sum_{k>K_0}\frac1{k(k+1)}.
\]
This bound is uniform in $N$ and tends to zero with $K_0$.
It proves both passage to the infinite sum and the asserted weak
convergence.  Its cross term with the strongly convergent $a$ part
therefore vanishes.

It remains to remove the input space $h\wedge K_N\mathcal H$.
Its restriction is obtained by putting $h$ in one input slot.
The $H_k$ terms again converge strongly; the sole surviving tail
term is
\[
 -\frac12\sum_{k\ge1}\frac{\overline{\mu_k(h)}}{k}R_{k,N}K_N.
\]
For polynomial $h$, the moments vanish when $k>d$.  Orthogonality
on each input column shows that the squared norm of this term tends
to $\tfrac14\sum_k|\mu_k(h)|^2/k$; it converges weakly to zero.
The term containing $R_{k,N}h$ has squared norm $O_h(N^{-1})$
by the same large-$k$ estimate used for the loss.  Subtract the
squared norms of this orthogonal input block from those with full
input $K_N^{\wedge2}$.  This proves the gain formula and the
strong-plus-weak description for polynomial $h$.

We extend the result to the stated analytic vectors by polynomial
approximation.  The image versions of the two vertex bounds in
Lemma~\ref{lem:beta2-two-body-vertex} bound the full loss and full
gain operators by $C\|h\|_*$, uniformly in $N$; the same estimates
bound their differences when $h$ is replaced by another vector.
Restriction to $h\wedge K_N$ changes continuously as well, by
the triangle inequality and these bounds.  Approximate $h$ by
normalized Taylor polynomials, first let $N$ increase, then let
the polynomial degree increase.  The assumption on $h_N$ permits
its replacement by the corresponding normalized truncations of
$h$.  The scalar term $\int L|h|^2d\mu$ is continuous in this
approximation, since $\|h\|_\infty\le\|h\|_*$ and $L$ is
integrable.  Also
\[
 \sum_{k\ge1}\frac{|\mu_k(h)|^2}{k}
 =\iint -\log|1-u\bar v|\,|h(u)|^2|h(v)|^2
                                    d\mu(u)d\mu(v)
\]
by the angular series, or its Abel regularizations followed by
dominated convergence.  The kernel is absolutely integrable,
so this expression is continuous under uniform convergence of
the analytic vectors.  This completes the approximation argument.

For $h=h_a$, reproducing gives $\mu_k(h_a)=a^k$.
Furthermore the squared coefficients are
$(1-|a|^2)^2(l+1)|a|^{2l}$, whence
\[
 \int L|h_a|^2d\mu
 =(1-|a|^2)^2\sum_{l\ge0}(l+1)H_{l+1}|a|^{2l}
 =1-\log(1-|a|^2).
\]
Thus the expression in braces is one.  Finally, if $u,v$ belong
to $J_a\mathcal H$, then $u(a)=v(a)=0$.  In the angular series
for $I$, contraction against $h_a$ of an analytic factor
$z^ku(z)$ vanishes; contraction of an antianalytic factor leaves
an analytic function in the other coordinate, which is killed
by $Q$.  This proves that the infinite-projection image gain is
zero, first for finite angular sums and then by the vertex bounds.
\end{proof}

\begin{lemma}[Conformal covariance of the limiting single contribution]
\label{lem:beta2-single-conformal}
The limiting difference of the single-replacement variances at
an anchor $a$ exists and equals the constant $S$ in
Lemma~\ref{lem:beta2-exact-center-single}.
\end{lemma}

\begin{proof}
The limiting renormalized Hartree--Fock operator on analytic vectors
is
\[
 F K=(M_L+\tfrac12T_L)K.
\]
This follows from the radial multiplier formula, since
$T_Lh_n=H_{n+1}h_n$.
Let $P=P_a$, $J=K-P$, and
\[
 B_a=M_{V_g(P)-g(\cdot,a)}+A_I(P)-A_g(P),
\]
where $V_g(P)(u)=\int g(u,v)P(v,v)d\mu(v)$ and
$A_j(P)$ has kernel $j(u,v)P(u,v)$.
This is precisely the previous screened Palm perturbation: the
harmonic reproducing identity gives
$V_I(P)(u)=I(u,a)$, so
$\log|u-a|-V_f(P)=V_g(P)-g(u,a)$.
The finite-Palm trace-norm bounds, local uniform convergence of
the coherent vectors, and the uniform bound on $Q_NF_{K_N}K_N$
show that the single contribution converges to
\begin{align*}
 S_a={}&-\|QFP\|_{\mathrm{HS}}^2
 +2\operatorname{Re}\langle QFJ,QB_aJ\rangle_{\mathrm{HS}}
 +\|QB_aJ\|_{\mathrm{HS}}^2
 +\|P(F+B_a)J\|_{\mathrm{HS}}^2.
\end{align*}
The cross pairing is well defined because $QFK$ is bounded and
$B_aK$ is trace class.  Convergence follows by truncating the
trace-class factor and the finitely many coherent columns, as
in the center proof.  The bounds used there hold uniformly on
compact sets of anchors.

Take a disk automorphism $\phi$ with $\phi(0)=a$, and put
$d=\log|\phi'|$.  Pull back by the unitary
$u\mapsto\phi'(u)f(\phi(u))$ on $L^2(\mu)$.
It preserves the Bergman space and takes $P_a$ to $P_0$.
Since $L\circ\phi=L-d$, the pulled-back $F$ on the Bergman space
is $F-M_d-T_d/2$, where $T_d=KM_dK$.
The Green part of $B_a$ is invariant, whereas
$I(\phi(u),\phi(v))=I(u,v)+(d(u)+d(v))/2$.
Thus its pulled-back screened perturbation is
\[
 B_0+\tfrac12(M_dP_0+P_0M_d).
\]
Put $P=P_0$, $J=K-P$ for the rest of the proof.
Combining these transformations shows that the full
single-replacement operator changes from $QFK$ to
$QFK-QM_dK$, and the Palm operator changes from
$(Q+P)(F+B_0)J$ to that operator minus $(Q+P)M_dJ$.
For the $P$ output this uses $PT_dJ=PM_dJ$; its coefficient
cancels the extra $P M_dJ/2$ from $B_0$.

The cross terms in the difference of these two squared norms
vanish.  This cancellation must be made before pairing: $QFK$
itself need not be Hilbert--Schmidt.  Let $T=QFK$ and $H=QM_dK$.
At the center $P(F+B_0)J=0$.  The cross difference reduces to
\[
 \langle QB_0J,HJ\rangle_{\mathrm{HS}}
                         -\langle TP,HP\rangle_{\mathrm{HS}}.
\]
Both pairings are absolutely defined by the trace-class and
finite-rank bounds.  Both are zero because $T$ and $B_0$ preserve
angular frequency, whereas multiplication by the nonconstant
harmonic part of $d$ changes it.  Its constant part is annihilated
by the off-diagonal projections.

The remaining squared-norm difference is
\[
 -\|QM_dP\|_{\mathrm{HS}}^2+\|PM_dJ\|_{\mathrm{HS}}^2=0.
\]
Indeed write $d=d_0+\sum_{k\ge1}(d_kz^k+\bar d_k\bar z^k)$.
The function $d$ extends harmonically past the closed disk, so
all relevant series converge absolutely, and the two displayed
norms both equal $\sum_{k\ge1}|d_k|^2/(k+1)$.
The off-diagonal operator $QM_dK$ is Hilbert--Schmidt, with
squared norm $\sum_{k\ge1}k|d_k|^2$.
Consequently $S_a=S_0=S$.
\end{proof}

\begin{proof}[Proof of Theorem~\ref{thm:beta2-exact-second-all-anchors}]
Use the single- and double-replacement variance decomposition at
the anchor $a$.  The preceding lemma handles the single term.
The Green gain and loss converge strongly in Hilbert--Schmidt
norm, by the unprojected vertex bounds from the center proof and
convergence of the coherent vectors.  Their limiting norms equal
$G_g,L_g$, since Green multiplication, $K$, and $P_a$ transform
unitarily under disk automorphisms.  Lemma~\ref{lem:beta2-general-image-boundary}
shows that the image parts each have limiting squared norm $1/4$
and converge weakly to zero for coherent vectors.  Their cross
terms with the strongly convergent Green parts vanish.  Thus the
double difference tends to $G_g-L_g=-S$, proving the first limit.

All passages are uniform for $|a|\le r<1$.  To see this directly,
the coherent vectors and their finite versions have geometrically
decaying coefficients in $\|\cdot\|_*$, uniformly on that compact
set.  They therefore admit simultaneous polynomial approximations.
For each fixed polynomial degree and angular cutoff the preceding
limits concern finitely many coefficients, uniformly bounded
radial integrals, and convergent finite-dimensional expressions;
hence they are uniform on this compact parameter set.  The
discarded tails are controlled uniformly by the vertex bounds and
trace-norm bounds used above.  The same compactness argument applies
to the Green operators, whose unprojected Hilbert--Schmidt kernels
depend continuously on $a$.  This proves local uniform convergence.
Finally use the locally uniform first-response theorem and
$\rho_{N,2}(a)\to1/[\pi(1-|a|^2)^2]$ in
$\rho''=\rho\{(\log\rho)''+((\log\rho)')^2\}$.
\end{proof}


\section{The conjectured density coefficient as a charge exponent}
\label{sec:charged-radial-intensity-exponent}

Let $0<\beta<2$ and use ordinary planar area.  Set
\[
 H_N=\sum_{i=1}^N-\log|z_i|,\qquad
 M_N(t)=\mathbb E_{N,\beta}e^{-tH_N},\qquad t>-2.
\]
These moments are finite: $|\Delta_N|^\beta$ is bounded on
$\overline\D^N$ and $|z|^t$ is integrable for $t>-2$.
The $t$-charged law is
\begin{equation}
 d\mathbb P_{N,\beta}^{[t]}
   =M_N(t)^{-1}e^{-tH_N}\,d\mathbb P_{N,\beta}.
 \label{eq:general-central-charge-law}
\end{equation}
The charge may be negative, positive, or zero.

\begin{lemma}[Uniform relative intensities for central charges]
\label{lem:general-charge-relative-intensity}
Let $\widehat\rho_{N}^{[t]}$ denote the one-point density of
$\mathbb P_{N,\beta}^{[t]}$ relative to $|a|^t dA(a)$, continued
continuously to $a=0$ by its defining integral.  For every compact
interval $I\subset(-2,\infty)$ there are constants
$0<c_{\beta,I}\le C_{\beta,I}<\infty$ such that
\begin{equation}
 c_{\beta,I}\le\widehat\rho_N^{[t]}(0)
               \le C_{\beta,I},\qquad N\ge1,\quad t\in I.
 \label{eq:general-charge-relative-bounds}
\end{equation}
Moreover, with $R_{n,\beta}=Q_{n+1,\beta}/Q_{n,\beta}$,
\begin{equation}
 M_N(t)=\frac{M_{N-1}(t+\beta)}
 {R_{N-1,\beta}\widehat\rho_N^{[t]}(0)},
 \qquad N\ge1,
 \label{eq:general-charge-recursion}
\end{equation}
where $M_0(u)=1$.
\end{lemma}

\begin{proof}
The ordered radial-ratio representation of the $t$-charged gas has
exactly the same nondecreasing angular factor as the unweighted gas.
Its reference rates are
\[
 a_k^{[t]}=2k+\frac\beta2k(k-1)+tk
         =k\left(2+t+\frac\beta2(k-1)\right)>0.
\]
Thus the radial domination proof applies without alteration.
Uniformly for $t$ in a compact subinterval of $(-2,\infty)$,
$a_k^{[t]}\ge c a_k$ for every $k$, with a fixed $c>0$.
The reference independent exponential gaps are consequently bounded
by $c^{-1}$ times the original reference gaps in a coordinatewise
coupling.  The compact-count and superlinear weighted-count moment
arguments therefore give uniform bounds on
\[
 \mathbb E_{N,\beta}^{[t]}\Xi_N(B_r)
 \quad\hbox{and}\quad
 \mathbb E_{N,\beta}^{[t]}
       \left[\sum_{i=1}^N(1-|z_i|^2)^s\right]^q
\]
for every fixed $r<1$, $s>1$, and $q>0$.  These bounds also hold
for the charges in $I+\beta$.

For clarity, the relative intensity is the continuous subharmonic
function
\[
 \widehat\rho_N^{[t]}(a)=
 \frac{N}{Z_{N,\beta}M_N(t)}
 \int_{\D^{N-1}} |\Delta_{N-1}(y)|^\beta
        \prod_{j=1}^{N-1}|y_j|^t|a-y_j|^\beta\,dA^{N-1}(y).
\]
Continuity and subharmonicity follow by the subharmonicity of powers
of polynomial moduli and domination on each compact set for fixed
$N,t$.  Its circular mean is nondecreasing in radius.  Hence, for
$A=\{1/4<|a|<1/2\}$,
\[
 \widehat\rho_N^{[t]}(0)
 \le |A|^{-1}\int_A\widehat\rho_N^{[t]}(a)\,dA(a)
 \le C_I\mathbb E_{N,\beta}^{[t]}\Xi_N(A)\le C_{\beta,I}.
\]
The middle inequality uses a uniform positive lower bound on $|a|^t$
for $a\in A$ and $t\in I$.  Using an annulus avoids any division
by a vanishing weight at the origin.

For the lower bound, interpolate the constant polynomial at zero.
If $\beta\le1$, subadditivity and Lagrange cancellation give
\[
 1\le\widehat\rho_N^{[t]}(0)\int_\D |z|^t\,dA(z).
\]
If $1<\beta<2$, put $q=\beta-1$ and choose $s>1$ with $sq<1$.
Apply weighted H\"older with $w(z)=(1-|z|^2)^s$ to the same
Lagrange identity.  The reduced Palm integral at zero has precisely
the $(N-1)$-particle law with central charge $t+\beta$, because
insertion at zero multiplies its weight by $|y|^\beta$.
The uniform weighted-count bound just proved and
$(x+y)^q\le x^q+y^q$ yield
\[
 1\le\widehat\rho_N^{[t]}(0)
 \left\{\int_\D |z|^t\,dA(z)
  +C_{\beta,I,s}\int_\D |z|^t(1-|z|^2)^{-sq}\,dA(z)\right\}.
\]
Both integrals are uniformly finite for $t\in I$: the condition
$t>-2$ handles zero and $sq<1$ handles the boundary.  This gives
the required positive lower bound, including $N=1$ when the exterior
sum is empty.

Finally, evaluating the defining integral at $a=0$ gives
\[
 \widehat\rho_N^{[t]}(0)
 =\frac{N Z_{N-1,\beta}M_{N-1}(t+\beta)}
        {Z_{N,\beta}M_N(t)}
 =\frac{M_{N-1}(t+\beta)}{R_{N-1,\beta}M_N(t)}.
\]
This is the claimed recursion, with no asymptotic normalization.
\end{proof}

\begin{theorem}[Signed charge exponents and the radial-depth law]
\label{thm:signed-charge-radial-exponent}
Fix $0<\beta<2$ and set $d_\beta=2/\beta-1/2$.
There is a set $\mathcal E_\beta\subset\mathbb N$ of natural
 density one such that
\begin{equation}
 \frac{\log M_N(t)}{\log N}\longrightarrow-d_\beta t
 \label{eq:signed-charge-free-energy}
\end{equation}
locally uniformly for $t\in(-2,\infty)$, as $N\to\infty$ through
$\mathcal E_\beta$.  For every fixed $t>-2$ and every finite $q>0$,
\begin{equation}
 \mathbb E_{N,\beta}^{[t]}
 \left|\frac{H_N}{\log N}-d_\beta\right|^q\longrightarrow0.
 \label{eq:signed-charge-depth-convergence}
\end{equation}
For every $\varepsilon>0$ and $0<s<t+2$,
\begin{equation}
 \mathbb P_{N,\beta}^{[t]}
 \left\{\left|\frac{H_N}{\log N}-d_\beta\right|>\varepsilon\right\}
 \le2N^{-s\varepsilon+o(1)}.
 \label{eq:signed-charge-depth-tail}
\end{equation}
In particular these statements apply to the original gas, $t=0$,
with any tail exponent $s\varepsilon$ for $0<s<2$.
\end{theorem}

\begin{proof}
For every positive integer $p$, exact confluence of the $p$ anchors
in the factorial correlation identity gives
\begin{equation}
 h_{p,N+p}(0,\ldots,0)
    =\frac{Q_{N,\beta}}{Q_{N+p,\beta}}M_N(p\beta).
 \label{eq:charge-confluent-identity}
\end{equation}
Indeed, after division by the Vandermonde of the anchors at distinct
locations, its integral contains $\prod_{i,j}|a_i-y_j|^\beta$.
Letting all $a_i$ tend to zero gives the charge $p\beta$; the
factorial prefactor is $(N+p)!/(N!Z_{N+p,\beta})$.
Continuity of $h_{p,N+p}$ justifies this evaluation on the collision
diagonal.  The hierarchy upper bound and
Theorem~\ref{thm:nonemptiness-all-orders-below-two} give
\begin{equation}
 0<b_{\beta,p}\le h_{p,N+p}(0,\ldots,0)
                    \le B_{\beta,p}<\infty
 \label{eq:charge-confluent-bounds}
\end{equation}
uniformly in $N$.  For the lower bound, divide the stated correlation
inequality at distinct anchors and then use continuity.

Let $\mathcal D_\beta$ be the density-one set of
Theorem~\ref{thm:partition-ratio-density-one}.  There is a density-one
set $\mathcal E_\beta$ with the following property: for every fixed
integer $j$, all sufficiently large $N\in\mathcal E_\beta$ have
$N+j\in\mathcal D_\beta$.  To verify this diagonal selection, put
\[
 A_m=\{N\ge m+1:N+j\in\mathcal D_\beta\text{ for }|j|\le m\}.
\]
These are nested density-one sets.  Choose strictly increasing
$L_m\to\infty$ so that
$\#([1,x]\setminus A_m)\le x/m$ for every integer $x\ge L_m$.
For $L_m\le N<L_{m+1}$ retain $N$ exactly when $N\in A_m$.
If $L_m\le x<L_{m+1}$, all excluded indices between $L_1$ and $x$
lie outside $A_m$ by nesting.  Their number is at most $x/m$;
adding the fixed initial segment proves density one.  The eventual
shift property is immediate.

It follows, for each fixed integer $p\ge1$, that
\[
 \frac{\log(Q_{N+p,\beta}/Q_{N,\beta})}{\log N}
 \longrightarrow p\left(\frac\beta2-2\right)=-p\beta d_\beta.
\]
Equations~\eqref{eq:charge-confluent-identity} and
\eqref{eq:charge-confluent-bounds} therefore prove
$\log M_N(p\beta)/\log N\to-p\beta d_\beta$.
At $p=0$ this is exact.  The functions
$f_N(u)=\log M_N(u)/\log N$ are convex by H\"older's inequality.
For $p\beta\le u\le(p+1)\beta$, convexity gives an upper bound by
the chord at those endpoints and a lower bound by the line through
$((p+1)\beta,f_N((p+1)\beta))$ with slope
$[f_N((p+2)\beta)-f_N((p+1)\beta)]/\beta$.
Both lines converge uniformly on the interval to $-d_\beta u$.
Thus convergence is uniform on every bounded subinterval of
$[0,\infty)$.  The same proof applies with $N$ replaced by $N-\ell$
for any fixed integer $\ell$, since the chosen set has the eventual
property for every integer shift of $\mathcal D_\beta$.

Now fix a compact interval $I\subset(-2,\infty)$.  Choose a fixed
integer $\ell\ge0$ such that $t+\ell\beta\ge0$ for all $t\in I$.
Iterating \eqref{eq:general-charge-recursion} gives
\begin{align*}
 \log M_N(t)
 &=\log M_{N-\ell}(t+\ell\beta)
   -\sum_{j=0}^{\ell-1}\log R_{N-j-1,\beta}\\
 &\hspace{12mm}
   -\sum_{j=0}^{\ell-1}
                 \log\widehat\rho_{N-j}^{[t+j\beta]}(0).
\end{align*}
The final sum is uniformly $O_{\beta,I}(1)$ by
Lemma~\ref{lem:general-charge-relative-intensity}.
After division by $\log N$, the other terms converge uniformly in
$t\in I$ to $-d_\beta(t+\ell\beta)+\ell\beta d_\beta$.
This proves \eqref{eq:signed-charge-free-energy}.

For fixed $t>-2$ and $0<s<t+2$, exponential Markov gives
\begin{align*}
 \mathbb P_{N,\beta}^{[t]}
 \{H_N>(d_\beta+\varepsilon)\log N\}
 &\le N^{-s(d_\beta+\varepsilon)}\frac{M_N(t-s)}{M_N(t)}
       =N^{-s\varepsilon+o(1)},\\
 \mathbb P_{N,\beta}^{[t]}
 \{H_N<(d_\beta-\varepsilon)\log N\}
 &\le N^{s(d_\beta-\varepsilon)}\frac{M_N(t+s)}{M_N(t)}
       =N^{-s\varepsilon+o(1)}.
\end{align*}
The lower event is empty when its threshold is nonpositive.
This proves the probability estimate.  To justify all moments,
for sufficiently large retained $N$ the first inequality with an
arbitrary threshold $u$ gives
\[
 \mathbb P_{N,\beta}^{[t]}\{H_N/\log N>u\}
 \le N^{-s(u-d_\beta)+1}
 \le 2^{-s(u-d_\beta)+1},\qquad u\ge d_\beta+2/s.
\]
This common exponential upper tail proves uniform integrability of
every finite power.  The finitely many earlier $N\ge2$ also have
all moments, because $M_N(t-s)<\infty$ and $x^q\le C_{q,s}e^{sx}$.
Convergence in probability now proves the asserted absolute-moment
convergence.
\end{proof}

\begin{remark}[A coefficient is identified, but the central prefactor is not]
The conjectured intensity coefficient is
$(4/\beta-1)/\pi=2d_\beta/\pi$.  The theorem identifies $d_\beta$
in the actual uncharged finite-particle gas and in all fixed charged
laws with $t>-2$, without assuming hyperbolic invariance or uniqueness.
It still does not identify the local limiting intensity.  The exact
central insertion formula is
\[
 \rho_{N+1,\beta}(0)=
       \frac{M_N(\beta)}{Q_{N+1,\beta}/Q_{N,\beta}}.
\]
The logarithmic asymptotics of numerator and denominator cancel.
Their finite prefactor ratio is not determined by
\eqref{eq:signed-charge-free-energy}.  The theorem also retains an
exceptional set of particle numbers of natural density zero.
\end{remark}


\subsection{The determinantal charge benchmark}
\label{subsec:charged-depth-beta2-benchmark}

At $\beta=2$ the signed-charge depth statement has a full-sequence
version and an exact finite prefactor.

\begin{proposition}[All central charges at the determinantal point]
\label{prop:charged-depth-beta2-benchmark}
For $t>-2$ and $\beta=2$,
\[
 M_N(t)=\prod_{k=1}^N\frac{2k}{2k+t}
       =\frac{\Gamma(N+1)\Gamma(1+t/2)}{\Gamma(N+1+t/2)}.
\]
Consequently $\log M_N(t)/\log N\to-t/2$ locally uniformly
on $(-2,\infty)$. Under each fixed $t$-charged law,
$H_N/\log N\to1/2$ in every positive absolute moment along the
full sequence. The deviation estimate of
Theorem~\ref{thm:signed-charge-radial-exponent} holds with
$d_2=1/2$ and any $0<s<t+2$.
\end{proposition}

\begin{proof}
Expand the two Vandermonde determinants and integrate all angles.
Orthogonality eliminates the terms with different monomial indices,
and the remaining radial norms are
\[
 \int_{\D}|z|^{2(k-1)+t}\,dA(z)=\frac{2\pi}{2k+t}.
\]
Thus the charged partition function is $N!$ times the product of
these norms. Division by its value at $t=0$ proves the product
formula; repeated use of $\Gamma(x+1)=x\Gamma(x)$ proves the
gamma-function form. Uniformly for $t$ in a fixed compact subset
of $(-2,\infty)$, Taylor expansion for all sufficiently large $k$
gives $\log(1+t/(2k))=t/(2k)+O(k^{-2})$.
The finitely many earlier terms are uniformly bounded, so
$\log M_N(t)=-(t/2)\log N+O(1)$ uniformly on that compact set.
The exponential Markov inequalities and the uniform exponential
tail argument in the preceding theorem now apply without an
exceptional set of particle numbers and prove the remaining claims.
\end{proof}

In particular, $M_N(2)=1/(N+1)$ and
$R_{N,2}=\pi/(N+1)$, so the exact insertion quotient is
$\rho_{N+1,2}(0)=1/\pi$. Away from $\beta=2$ the logarithmic
charge asymptotic alone does not identify this quotient's finite
prefactor.


\subsection{The same radial-depth coefficient at every hyperbolic center}
\label{subsec:all-anchor-radial-depth}

For a deterministic $a\in\D$ put
\[
 H_N(a)=\sum_{w\in\Xi_N}g(a,w),\qquad
 g(a,w)=-\log\left|\frac{a-w}{1-a\bar w}\right|.
\]
Thus $H_N(0)=\sum_w-\log|w|$.  Each $H_N(a)$ is finite almost surely,
although as a function of $a$ it has logarithmic singularities at the
particle positions.

\begin{lemma}[Tightness of anchored potential differences below $2$]
\label{lem:anchored-potential-difference-tight}
For $0<\beta<2$ and each fixed $a\in\D$, the random variables
$H_N(a)-H_N(0)$ form a tight family.  The same is true jointly for any
finite list of deterministic anchors.
\end{lemma}

\begin{proof}
Take an arbitrary sequence of particle numbers.  By tightness of the
point processes, the almost-sure infinitude of every limiting process
below $2$, and Theorem~\ref{thm:corrected-field-gaussian-independence},
it has a further subsequence with a simultaneous augmentation by the
nonempty limiting configuration $\Xi$ and the corrected meromorphic
field $\mathcal A$.  The normal-family construction underlying that
theorem and the meromorphic gluing proposition give locally uniform
convergence of $\mathcal A_N$ away from the limiting poles, on each
positive-count localization.  These localizations exhaust an event of
probability one because $\Xi$ is nonempty.  Equivalently one may retain
the normalized holomorphic exterior factors in the countable disk
exhaustion used in that construction.

Use a Skorohod representation of these jointly convergent Polish-space
objects.  Almost surely none of the finitely many deterministic anchors
or the origin belongs to $\Xi$, because its intensity is absolutely
continuous.  On such a realization, choose for each anchor a piecewise
smooth path from $0$ to that anchor whose compact image avoids $\Xi$.
Such a path exists since the configuration is locally finite.  The path
can be covered by finitely many compact subsets of the complement of
the poles.  The convergences of the configurations and fields imply that
for all sufficiently large retained $N$ it also avoids the finite-$N$
poles, and $\mathcal A_N$ converges uniformly on that path.  The identity
$2\partial_z g(z,w)=-J_w(z)$ gives
\[
 H_N(a)-H_N(0)
       =-\operatorname{Re}\int_0^a\mathcal A_N(z)\,dz.
\]
The real part is independent of the avoiding path: the residue of each
pole is one, so the real part of every closed contour integral is zero.
Uniform convergence on the chosen paths therefore gives finite limits
for all the anchored differences.  This proves subsequential tightness
along every original sequence and hence tightness of the original
families.  Random selection of the paths in this argument creates no
measurability issue: the left-hand sides are measurable random variables
and the almost-sure limits are their pointwise limits; the paths only
verify finiteness and convergence on each realization.
\end{proof}

\begin{corollary}[A center-independent depth coefficient]
\label{cor:all-anchor-depth-coefficient}
Fix $0<\beta<2$ and let $\mathcal E_\beta$ be a natural-density-one set
on which the uncharged depth conclusion of
Theorem~\ref{thm:signed-charge-radial-exponent} holds.  Put
$d_\beta=2/\beta-1/2$.  For every fixed $a\in\D$ and every $p>0$,
\begin{equation}
 \mathbb E\left|\frac{H_N(a)}{\log N}-d_\beta\right|^p
       \longrightarrow0,
 \qquad N\to\infty,\ N\in\mathcal E_\beta.
 \label{eq:all-anchor-depth-moments}
\end{equation}
The convergence holds jointly in probability for finitely many anchors.
Moreover, for every compact $K\Subset\D$ and every $p>0$,
\begin{equation}
 \mathbb E\int_K
 \left|\frac{H_N(a)}{\log N}-d_\beta\right|^p dA(a)
       \longrightarrow0
 \qquad(N\in\mathcal E_\beta).
 \label{eq:all-anchor-depth-spatial-moments}
\end{equation}
This statement concerns deterministic anchors and integrated spatial
norms; it asserts no supremum bound over a compact set containing random
particle positions.
\end{corollary}

\begin{proof}
The tightness lemma shows
$(H_N(a)-H_N(0))/\log N\to0$ in probability.  Combined with the
uncharged depth law at zero, it proves the asserted pointwise and finite
joint convergence in probability.

We verify uniform integrability and its uniformity over compact sets of
anchors.  Fix $K\Subset\D$ and choose
$\sup_{a\in K}|a|<R<1$.  There is a finite constant $C_K$ such that
\begin{equation}
 g(a,w)\le C_K(-\log|w|),\qquad a\in K,\ |w|\ge R.
 \label{eq:green-depth-outer-comparison}
\end{equation}
Indeed, the ratio of the two positive functions is continuous on the
closed outer annulus after extension to $|w|=1$.  Its boundary value is
$(1-|a|^2)/|a-w|^2$, uniformly bounded for $a\in K$.  The extension
follows either by differentiating radially at the unit circle or from
\[
 1-\left|\frac{a-w}{1-a\bar w}\right|^2
       =\frac{(1-|a|^2)(1-|w|^2)}{|1-a\bar w|^2}
\]
and $-\log x\sim1-x$ as $x\uparrow1$.
On the inner disk,
\[
 g(a,w)\le\log2+\log^+\frac1{|a-w|}.
\]
Consequently
\begin{equation}
 H_N(a)\le C_K H_N(0)+L_N(a),\qquad
 L_N(a)=\sum_{w\in\Xi_N\cap B_R}
       \left(\log2+\log^+\frac1{|a-w|}\right).
 \label{eq:green-depth-inner-outer-bound}
\end{equation}
For every $q>0$,
\begin{equation}
 \sup_N\sup_{a\in K}\mathbb E L_N(a)^q<\infty.
 \label{eq:local-log-singularity-moments}
\end{equation}
To prove this, first take an integer $k\ge\max(1,q)$.  Expand the
$k$th power of the nonnegative sum according to which of its $k$ indices
coincide.  Each resulting term is an integral against a factorial
density $\rho_{j,N}$ on $B_R^j$, where $1\le j\le k$, of a product
of powers of $\log2+\log^+(1/|a-w|)$ of total degree $k$.
The densities have uniform bounds on this compact polydisk.  Each
logarithmic power has an area integral bounded uniformly in $a\in K$,
because all singularities are translates of the integrable functions
$(\log^+(1/|w|))^l$ on a bounded planar set.  There are finitely many
index partitions, which proves the claim for $k$.  H\"older's
inequality then proves it for $q$.

The uncharged depth theorem gives bounded moments of every order for
$H_N(0)/\log N$ on $\mathcal E_\beta$, with $N\ge2$.
Equations~\eqref{eq:green-depth-inner-outer-bound} and
\eqref{eq:local-log-singularity-moments} therefore imply, for every
$q>0$,
\[
 \sup_{\substack{N\in\mathcal E_\beta\\N\ge2}}
 \sup_{a\in K}\mathbb E\left(\frac{H_N(a)}{\log N}\right)^q
 <\infty.
\]
Choosing $q>p$ proves uniform integrability to the $p$th power.
Pointwise convergence in probability now proves
\eqref{eq:all-anchor-depth-moments}.  Finally its integrand's expectation
is bounded uniformly in $a\in K$ and retained $N$.  Dominated convergence
in $a$, followed by Tonelli's theorem, proves
\eqref{eq:all-anchor-depth-spatial-moments}.
\end{proof}

\begin{remark}[The order-one term remains undetermined]
The preceding result fixes the leading coefficient $d_\beta\log N$ of
the positive Green potential at every center.  The intensity is encoded
by the order-one spatial variation of its expectation, namely the exact
mean defect of Proposition~\ref{prop:intensity-mean-potential-defect}.
Thus the center-independent leading coefficient, by itself, does not
prove that the mean defect vanishes or determine the local intensity.
\end{remark}


\section{Two exact reductions of the intensity conjecture}
\label{sec:intensity-palm-dipole-affine}

The conjectural formula separates into a spatial-shape problem and a
coefficient problem. The first reduction below uses only a one-point
Palm statistic. The second determines the coefficient from symmetry
of a meromorphic field, without assuming that this field is measurable
from the particle configuration or constructing a centered Green
potential. Neither symmetry hypothesis is proved for the original gas.

\subsection{The exact Palm dipole controlling the radial profile}

Write $\rho_N(a)=\rho_{1,N}(a)$ and
$c_N=2+\beta(N-1)/2$. Let $\E_N^{!a}$ denote expectation under
the reduced Palm law of the $N$-particle gas at $a\in\D$, and put
\[
 d_N(a)=a+\E_N^{!a}\sum_{w\in\Xi_N}w.
\]
Thus the sum in this display contains the other $N-1$ particles.
Write $\rho_N(a)=f_N(|a|^2)$.

\begin{proposition}[A finite-particle Palm dipole identity]
\label{prop:exact-intensity-palm-dipole}
For every $N\ge1$ and $\beta>0$, $f_N$ is positive and continuously
differentiable on $[0,1)$. For $a\ne0$,
\begin{equation}
 d_N(a)=\frac{a}{c_N}
 \left(2-(1-|a|^2)\frac{f_N'(|a|^2)}{f_N(|a|^2)}\right).
 \label{eq:exact-intensity-palm-dipole}
\end{equation}
In particular, the quotient $d_N(a)/a$ is real and radial and has a
continuous extension at zero. For $0\le r<1$,
\begin{equation}
 \log\frac{(1-r^2)^2\rho_N(r)}{\rho_N(0)}
 =-c_N\int_0^{r^2}
       \frac{d_N(\sqrt{s})}{\sqrt{s}(1-s)}\,ds,
 \label{eq:exact-intensity-profile-dipole}
\end{equation}
with that extension understood at $s=0$.
\end{proposition}

\begin{proof}
The finite-$N$ insertion formula expresses $\rho_N(a)$ as a positive
constant times the integral of
$\prod_{j=1}^{N-1}|a-w_j|^\beta|\Delta_{N-1}(\mathbf w)|^\beta$.
First and second derivatives in $a$ are locally integrable, because
the worst single-factor singularity has modulus a constant times
$|a-w_j|^{\beta-2}$, integrable in two dimensions for $\beta>0$.
Terms with two differentiated factors have exponents $\beta-1$ in
two separate variables and are likewise integrable. The other factors
are bounded on the finite product of disks. Splitting small balls
around $a$ from their complements justifies differentiation and the
continuity of the resulting derivatives. Thus $\rho_N$ is $C^2$.
Rotation invariance gives the asserted regularity of $f_N$, including
at zero by the second-order radial Taylor formula. Strict positivity
follows from the insertion formula.

Take a real $h\in C_c^1(\D)$ and apply the exact M\"obius Ward
identity to $F(\xi)=\sum_{a\in\xi}h(a)$. This functional is bounded
for fixed $N$; alternatively it can be truncated outside the possible
range of its argument before applying that identity. Since
$\E S_{1,N}=0$, the Campbell formula gives, for every $t\in\C$,
\[
 \int_\D dh(a)[t-\overline t a^2]\rho_N(a)\,dA(a)
 =2c_N\int_\D h(a)\rho_N(a)
             \operatorname{Re}(\overline t d_N(a))\,dA(a).
\]
The real divergence of $a\mapsto t-\overline t a^2$ equals
$-4\operatorname{Re}(\overline t a)$. Since
$\rho_N(a)=f_N(|a|^2)$, integration by parts therefore makes the
left side equal to
\[
 \int_\D h(a)\operatorname{Re}(\overline t a)
       \left[4f_N(|a|^2)-2(1-|a|^2)f_N'(|a|^2)\right]dA(a).
\]
The Palm expectation in $d_N$ is continuous in $a$ by its bounded
insertion numerator and positive denominator. The arbitrariness of
$h$ and $t$ proves~\eqref{eq:exact-intensity-palm-dipole} everywhere.
Rearranging that equality gives
\[
 \frac{d}{ds}\log\bigl[(1-s)^2 f_N(s)\bigr]
 =-\frac{c_N}{1-s}\frac{d_N(\sqrt{s})}{\sqrt{s}}.
\]
Integration from zero to $r^2$ proves the second formula.
\end{proof}

\begin{corollary}[An integrated dipole criterion for the spatial shape]
\label{cor:intensity-integrated-dipole-criterion}
Suppose $N_j\to\infty$, the intensities converge locally uniformly
to $\rho$, and $\rho(0)>0$. Then the following two assertions are
equivalent:
\begin{enumerate}
\item $\rho(r)=\rho(0)(1-r^2)^{-2}$ for every $0\le r<1$;
\item for every $0\le r<1$,
\[
 c_{N_j}\int_0^{r^2}
 \frac{d_{N_j}(\sqrt{s})}{\sqrt{s}(1-s)}\,ds\longrightarrow0.
\]
\end{enumerate}
A sufficient condition is, for every $r<1$,
\begin{equation}
 N_j\int_0^{r^2}
     \frac{|d_{N_j}(\sqrt{s})|}{\sqrt{s}(1-s)}\,ds
     \longrightarrow0.
 \label{eq:intensity-absolute-dipole-criterion}
\end{equation}
These criteria concern only the first correlation function and do
not require M\"obius invariance of all local observables.
\end{corollary}

\begin{proof}
The left side of~\eqref{eq:exact-intensity-profile-dipole} converges
to $\log[(1-r^2)^2\rho(r)/\rho(0)]$ whenever $\rho(r)>0$.
Conversely, if the integral tends to zero, that identity gives
$(1-r^2)^2\rho_{N_j}(r)/\rho_{N_j}(0)\to1$, hence positivity of
$\rho(r)$ and the asserted shape. Finally $c_{N_j}/N_j\to\beta/2$,
so~\eqref{eq:intensity-absolute-dipole-criterion} suffices.
\end{proof}

\begin{remark}[The missing order of Palm screening]
The harmonic Palm theorem proves $d_N(a)\to0$ for each fixed $a$.
The exact identity above explains why this leading screening is
insufficient: the profile involves a dipole multiplied by $c_N$,
which is of order $N$. In particular, even an $O(N^{-1})$ dipole
can produce a nonconstant limiting factor multiplying
$(1-|a|^2)^{-2}$. At $\beta=2$ the criterion can be checked
directly: putting $s=|a|^2$,
\[
 \pi(1-s)^2 f_N(s)=1-(N+1)s^N+Ns^{N+1},
\]
and consequently
\[
 \frac{d_N(a)}a
 =\frac{N(1-s)^2s^{N-1}}
       {1-(N+1)s^N+Ns^{N+1}},\qquad c_N=N+1.
\]
The right side is exponentially small on every fixed compact
subdisk as $N\to\infty$. No corresponding estimate at other
temperatures is claimed.
\end{remark}

\begin{proposition}[An exact susceptibility deficit at the center]
\label{prop:intensity-center-susceptibility-deficit}
Let
\[
 V_N=\E_N|S_{1,N}|^2,\qquad
 V_N^{!0}=\E_N^{!0}\left|\sum_w w\right|^2.
\]
For every $N\ge1$ and $\beta>0$,
\begin{equation}
 V_N=\frac N{c_N},\qquad
 c_N^2(V_N^{!0}-V_N)=\frac{f_N'(0)}{f_N(0)}-2.
 \label{eq:intensity-center-susceptibility-deficit}
\end{equation}
Consequently $f_N'(0)/f_N(0)\to2$ is equivalent to
$V_N^{!0}-V_N=o(N^{-2})$. The equality $V_N^{!0}=V_N$ holds
at $\beta=2$ for every $N\ge2$.
\end{proposition}

\begin{proof}
Applying the unweighted M\"obius Ward identity to $S_{1,N}$ and
using rotation invariance gives $N=c_NV_N$.
For the other identity put $n=N-1$ and, under the Palm law at zero,
write $S=\sum_{j=1}^n w_j$ and $T=\sum_{j=1}^n w_j^{-1}$.
For $n=0$ both sums are zero. For $n>0$ these variables and $ST$
are integrable, because the Palm density contains
$\prod_j|w_j|^\beta$, and a single inverse power is locally
integrable against each of these weighted area factors.
Differentiating the Palm insertion formula at $a=0$ gives
\[
 \lim_{a\to0}\frac{d_N(a)}a=1-\frac\beta2\E_N^{!0}(ST).
\]
Indeed the derivative of the insertion factor with respect to $a$
is $-(\beta/2)T$ there; rotation invariance makes the derivative
of its normalizer vanish and removes the conjugate-linear term.

Under the disk flow $w\mapsto (w+t)/(1+tw)$ with real $t$,
the logarithmic derivative at zero of the Palm density, including
its Jacobian, equals
$-2c_N\operatorname{Re}S+\beta\operatorname{Re}T$.
The differentiated statistic $S$ has mean derivative $n$.
The change-of-variables identity therefore yields
\[
 0=n-c_NV_N^{!0}+\frac\beta2\E_N^{!0}(ST).
\]
The derivative is integrable also at the inserted zero, by the
integrability noted above; it may equivalently be justified first
off a shrinking disk and then by dominated integral estimates.
It follows that
\[
 \lim_{a\to0}\frac{d_N(a)}a=N-c_NV_N^{!0}.
\]
Combine this with~\eqref{eq:exact-intensity-palm-dipole} at zero
and $V_N=N/c_N$ to obtain
\eqref{eq:intensity-center-susceptibility-deficit}.
The asymptotic equivalence uses $c_N/N\to\beta/2$.
At $\beta=2$, the explicit density
$f_N(s)=\pi^{-1}\sum_{k=0}^{N-1}(k+1)s^k$ gives
$f_N'(0)/f_N(0)=2$ when $N\ge2$, proving the last assertion.
\end{proof}

\begin{remark}[A finer cancellation than the harmonic central limit theorem]
The common limiting variance $2/\beta$ does not determine the
susceptibility difference in
\eqref{eq:intensity-center-susceptibility-deficit}. Even identifying
the $N^{-1}$ correction to both variances would leave the
$N^{-2}$ scale unresolved. The proposition is an exact statement
about finite-$N$ spatial derivatives. Turning it into a necessary
criterion for a proposed limiting profile additionally requires
convergence of those derivatives; local uniform intensity
convergence alone does not justify that step.
\end{remark}

\subsection{Affine meromorphic-field symmetry fixes the coefficient}

\begin{theorem}[Intensity fixed by an affine field symmetry]
\label{thm:affine-meromorphic-field-intensity}
Let $\kappa\in\R$. Let $\Xi$ be a simple locally finite random
configuration in $\D$, with no particle at any fixed deterministic
point. Let $\mathcal A$ be a random meromorphic function whose poles
are exactly the points of $\Xi$, each pole being simple with residue
one. Assume that, for every disk automorphism $\varphi$, the pair
\begin{equation}
 \left(\varphi^{-1}\Xi,
       \varphi'\,\mathcal A\circ\varphi
             +\kappa\frac{\varphi''}{\varphi'}\right)
 \quad\hbox{has the same law as }(\Xi,\mathcal A).
 \label{eq:affine-meromorphic-field-stationarity}
\end{equation}
Assume also $\E|\mathcal A(0)|<\infty$. Then
\begin{equation}
 \E\mathcal A(z)=\frac{2\kappa\overline z}{1-|z|^2},
 \qquad
 \E\Xi(dz)=\frac{2\kappa}{\pi(1-|z|^2)^2}\,dA(z).
 \label{eq:affine-meromorphic-exact-intensity}
\end{equation}
In particular $\kappa\ge0$. If $\kappa=0$, the configuration is
empty almost surely; if $\kappa>0$, its total count is infinite
almost surely. For the affine action of the Gaussian-eliminated disk-gas
specification, $\kappa=2/\beta-1/2$, so the asserted intensity is
exactly
\[
 \rho(z)=\frac{4/\beta-1}{\pi(1-|z|^2)^2}.
\]
\end{theorem}

\begin{proof}
Take $\varphi(z)=e^{i\theta}z$ in the symmetry assumption.
The random variables $\mathcal A(0)$ and
$e^{i\theta}\mathcal A(0)$ have the same law. Integrability gives
$\E\mathcal A(0)=0$. For a fixed $a\in\D$ use
$\varphi_a(z)=(z+a)/(1+\overline a z)$. Since
\[
 \varphi_a'(0)=1-|a|^2,\qquad
 \frac{\varphi_a''(0)}{\varphi_a'(0)}=-2\overline a,
\]
the symmetry assumption implies the distributional identity
\[
 \mathcal A(0)\ \stackrel{\rm law}=\
 (1-|a|^2)\mathcal A(a)-2\kappa\overline a.
\]
It proves integrability at every deterministic $a$, the first
identity in~\eqref{eq:affine-meromorphic-exact-intensity}, and the
locally uniform bound
\[
 \E|\mathcal A(a)|
 \le\frac{\E|\mathcal A(0)|+2|\kappa||a|}{1-|a|^2}.
\]
Tonelli therefore gives
$\E\int_K|\mathcal A(a)|\,dA(a)<\infty$ for every compact
$K\Subset\D$.

Every sample meromorphic function is locally integrable; subtracting
its finitely many residue-one principal parts on a neighborhood of
a test function's support leaves a holomorphic function there. The
distributional identity $\partial_{\overline z}(1/(z-a))=\pi\delta_a$
thus gives, sample by sample,
\[
 \partial_{\overline z}\mathcal A=\pi\Xi.
\]
The local expected-integrability bound permits expectation to pass
through its distributional pairing. It also ensures local finiteness
of $\E\Xi$: pair the displayed identity with a nonnegative smooth
function equal to one on a given compact set. Consequently
\[
 \pi\E\Xi
 =\partial_{\overline z}\E\mathcal A
 =\frac{2\kappa}{(1-|z|^2)^2}\,dA(z),
\]
as asserted. Positivity of the left side forces $\kappa\ge0$.
Zero intensity implies the almost surely empty configuration by a
countable compact exhaustion. For $\kappa>0$, integrating over
$B_r$ yields
$\E\Xi(B_r)=2\kappa r^2/(1-r^2)$, which diverges as $r\uparrow1$.
If some finite total count $m$ had positive probability, condition
on $\{\Xi(\D)=m\}$. This event is invariant under every action
in~\eqref{eq:affine-meromorphic-field-stationarity}, and integrability
of $\mathcal A(0)$ persists under this conditioning. Applying the
already proved intensity identity to that conditional law would
give infinite expected total count, contradicting the deterministic
count $m$. All finite counts therefore have probability zero.
\end{proof}

\begin{remark}[Scope of the field criterion]
The proof needs only the corresponding symmetry identities for first
moments and the local expected-integrability that they supply. Full
joint stationarity is a convenient sufficient hypothesis. Neither a
centered Green potential, a count-resolved DLR equation, ergodicity,
nor measurability of $\mathcal A$ from $\Xi$ is required.
Conversely, covariance of the canonical resampling kernels does not
prove~\eqref{eq:affine-meromorphic-field-stationarity} for a chosen
Gibbs law. Nor does the proved leading independence of the Gaussian
harmonic modes imply this affine symmetry. These remain substantive
hypotheses for an actual disk-gas limit.
\end{remark}


\subsection{An exact normalized mean-potential defect}
\label{subsec:intensity-mean-potential-defect}

The coefficient in the proposed intensity can be tested without assuming
that the corrected field is measurable from the point configuration, or
that its limit has pointwise integrable evaluations.  The following
identities are finite-particle identities followed by limits of their
expectations.  They do not interchange an expectation with convergence
of an unbounded random field.

Put $q_\beta=4/\beta-1$ and $\kappa_\beta=q_\beta/2$.  For a finite
configuration $\xi\subset\D$ having no particle at $0$ or $z$, define
\begin{align*}
 V_\xi(z)&=\sum_{w\in\xi}\{g(z,w)-g(0,w)\},\\
 U_\xi(z)&=V_\xi(z)-\kappa_\beta\log(1-|z|^2),\qquad
 g(z,w)=-\log\left|\frac{z-w}{1-z\bar w}\right|.
\end{align*}
Write $m_N(r)=\mathbb E\Xi_N(B_r)$ and
$u_N(z)=\mathbb E U_{\Xi_N}(z)$, where $B_r=\{|z|<r\}$.

\begin{proposition}[The mean defect is determined exactly by the intensity]
\label{prop:intensity-mean-potential-defect}
For every $N$, $\beta>0$, and $z\in\D$, the expectation defining
$u_N(z)$ is finite and
\begin{equation}
 u_N(z)=-\int_0^{|z|}\frac{m_N(s)}s\,ds
            -\kappa_\beta\log(1-|z|^2).
 \label{eq:finite-intensity-mean-defect}
\end{equation}
Let $\Xi_{N_j}$ converge to a subsequential point-process limit with
intensity $\rho(z)dA(z)$, and set
$m(r)=2\pi\int_0^r s\rho(s)\,ds$.  Then $u_{N_j}$ converges locally
in $C^2(\D)$ to the deterministic radial function
\begin{equation}
 u(z)=-\int_0^{|z|}\frac{m(s)}s\,ds
            -\kappa_\beta\log(1-|z|^2).
 \label{eq:limiting-intensity-mean-defect}
\end{equation}
It satisfies $u(0)=0$ and
\begin{equation}
 \Delta u(z)=2\pi\left\{
       \frac{q_\beta}{\pi(1-|z|^2)^2}-\rho(z)\right\}.
 \label{eq:intensity-exact-poisson-defect}
\end{equation}
Consequently, for this subsequential limit the proposed formula
$\rho(z)=q_\beta/[\pi(1-|z|^2)^2]$ holds everywhere if and only if
$u$ vanishes identically.
\end{proposition}

\begin{proof}
At fixed $N$ the density is bounded on $\D^N$: the Vandermonde is
bounded there and its partition function is strictly positive.  The
logarithmic singularities at $w=0,z$ are integrable with respect to
area, while $\log|1-z\bar w|$ is bounded for fixed $z\in\D$.
Thus all expectations and the angular integrations below are absolutely
convergent at fixed $N$.

Rotation invariance makes the intensity radial.  Jensen's elementary
angular formula gives, for $r,s>0$ with $r,s<1$,
\[
 \frac1{2\pi}\int_0^{2\pi}\log|r-se^{i\theta}|\,d\theta
       =\log\max(r,s),\qquad
 \frac1{2\pi}\int_0^{2\pi}\log|1-rse^{i\theta}|\,d\theta=0.
\]
For completeness, away from $r=s$ both formulas follow by factoring out
the larger modulus and integrating the convergent power series for
$\operatorname{Re}\log(1-ae^{i\theta})$, $|a|<1$.  At $r=s$, the
identity follows by direct integration of
$\log|1-e^{i\theta}|=\log(2|\sin(\theta/2)|)$, whose integral is zero;
alternatively, dominated convergence after separating a small interval
about $\theta=0$ gives the same conclusion.  It follows that the
angular mean of $g(r,se^{i\theta})-g(0,se^{i\theta})$ is
$-\log(r/s)$ for $s<r$ and zero for $s\ge r$.  Therefore
\[
 \mathbb E V_{\Xi_N}(z)
 =-2\pi\int_0^r s\rho_N(s)\log(r/s)\,ds
 =-\int_0^r\frac{m_N(s)}s\,ds,
 \qquad r=|z|,
\]
where Tonelli's theorem applied to the nonnegative integrand establishes
the last equality.  This proves
\eqref{eq:finite-intensity-mean-defect}.

The proved local uniform convergence of the intensities implies local
uniform convergence of $m_N(r)/r^2$ including its continuous value at
zero.  In particular the integral in
\eqref{eq:finite-intensity-mean-defect} converges uniformly on compact
radial intervals.  To check two derivatives without assuming derivatives
of $\rho_N$ converge, set
$h_N(r)=-\int_0^r m_N(s)\,ds/s$.  For $r>0$,
\[
 h_N'(r)=-\frac{m_N(r)}r,\qquad
 h_N''(r)=-2\pi\rho_N(r)+\frac{m_N(r)}{r^2}.
\]
The corresponding Hessian of $h_N(|z|)$ is a combination of
$h_N''(r)$ and $h_N'(r)/r$, both controlled by the locally uniform
convergence of $\rho_N$.  At zero both tend to $-\pi\rho_N(0)$;
continuity follows from continuity of $\rho_N$.  These formulas prove
local $C^2$ convergence.  They also give
$\Delta h=-2\pi\rho$.  Since
$\Delta\log(1-|z|^2)=-4/(1-|z|^2)^2$, this proves
\eqref{eq:intensity-exact-poisson-defect}.

If the conjectured intensity holds, the right-hand side of
\eqref{eq:intensity-exact-poisson-defect} vanishes.  The function $u$
is radial, harmonic on the entire disk, and zero at the origin, so it
is zero everywhere.  This also follows directly by integrating
$m(r)=q_\beta r^2/(1-r^2)$.  The converse follows immediately from
\eqref{eq:intensity-exact-poisson-defect}.
\end{proof}

\begin{remark}[What this does and does not prove]
The quantity $u$ is the limit of the deterministic expectations
$\mathbb E U_{\Xi_N}$.  Identifying it with the expectation of a limiting
random potential requires a separate uniform-integrability argument.
Moreover, tightness or existence of that random potential would not
itself show that its mean is zero.  Thus the exact background coefficient
in the augmented canonical kernels is compatible with the proposed
intensity precisely when the mean defect above vanishes; kernel
covariance alone does not establish that vanishing.
\end{remark}

One can express the same missing information in terms of the corrected
meromorphic field
$\mathcal A_N(z)=\sum_{w\in\Xi_N}J_w(z)$.  At finite $N$ its first
moment exists at every fixed interior point and angular integration gives
\begin{equation}
 \mathbb E\mathcal A_N(z)=\frac{m_N(|z|)}z\quad(z\ne0),
 \qquad \mathbb E\mathcal A_N(0)=0.
 \label{eq:mean-corrected-field-radial-count}
\end{equation}
Indeed the angular mean of $(z-w)^{-1}$ is $z^{-1}$ when $|w|<|z|$
and zero when $|w|>|z|$, while the angular mean of
$\bar w/(1-z\bar w)$ is zero.  Its singularity is locally integrable,
so this calculation uses ordinary Fubini at each finite $N$.
For $a\in\D$ let $\varphi_a(z)=(a+z)/(1+\bar a z)$.  The affine
transformation prescribed by the Gaussian-free kernels has mean at zero
\begin{align}
 &\mathbb E\left[
 \varphi_a'(0)\mathcal A_N(\varphi_a(0))
       +\kappa_\beta\frac{\varphi_a''(0)}{\varphi_a'(0)}\right]
 \nonumber\\
 &\hspace{2em}=
 \bar a\left\{\frac{(1-|a|^2)m_N(|a|)}{|a|^2}-q_\beta\right\},
 \qquad a\ne0.
 \label{eq:affine-mean-covariance-defect}
\end{align}
Its limit exists along the point-process subsequence by local count
convergence.  Vanishing of that limit for every $a$ is another exact
equivalent of the proposed intensity.  This assertion concerns limits
of finite-particle means and requires no unproved passage of an
expectation through a limiting meromorphic mark.


\subsection{Exact changes of geometry and a finite-size prediction at the critical value}
\label{subsec:intensity-literature-geometry}

The two-dimensional counterion literature contains the same finite
disk ensemble, but several limits and changes of geometry must be
distinguished.  In particular, a finite-size extrapolation at the
critical value should be confronted with the qualitative emptiness
theorems, rather than regarded as an established limiting formula.

\begin{proposition}[Inversion and excess charge]
\label{prop:intensity-inversion-excess-charge}
Under the map $w_i=1/z_i$, the unweighted disk gas has density proportional to
\begin{equation}\label{eq:intensity-inverted-density}
 |\Delta_N(w)|^\beta
 \prod_{i=1}^N |w_i|^{-\beta(N-1)-4}
 \mathbf 1_{\{|w_i|>1\}}\,dA(w_i).
\end{equation}
Writing the external-disk counterion density as
$|\Delta_N(w)|^{2\Xi}\prod_i|w_i|^{-2\xi}$, its parameters are
\[
 \Xi=\frac\beta2,\qquad
 \xi=\frac\beta2(N-1)+2.
\]
Consequently the central charge in units of a counterion charge is
\[
 Q=\frac\xi\Xi=N-1+\frac4\beta,
 \qquad Q-N=\frac4\beta-1.
\]
If $\rho^{\mathrm{ext}}_{N,\beta}$ denotes the inverted intensity, then
\[
 \rho^{\mathrm{ext}}_{N,\beta}(w)
   =|w|^{-4}\rho_{1,N}(1/w),\qquad
 \lim_{|w|\to\infty}|w|^4\rho^{\mathrm{ext}}_{N,\beta}(w)
   =\rho_{1,N}(0).
\]
\end{proposition}

\begin{proof}
The origin has probability zero.  The identities
$|z_i-z_j|=|w_i-w_j|/(|w_i||w_j|)$ and
$dA(z_i)=|w_i|^{-4}dA(w_i)$ give
\eqref{eq:intensity-inverted-density}.  Matching the exponents proves
the parameter statements.  Changing variables in the one-point
intensity gives the last identity, and its limit follows from
continuity of $\rho_{1,N}$ at zero.
\end{proof}

The external-disk notation above agrees with the Boltzmann factor
in equation (4.1) of
\href{https://arxiv.org/pdf/1502.05291}{Uma\~na Mallarino--T\'ellez (2015)}.
Thus the conjectured coefficient $\pi\rho_\beta(0)=4/\beta-1$
has an exact interpretation as the excess charge in this finite-model
reformulation.  This matching does not prove equality between the excess
charge and the limiting density amplitude.  The condensation thresholds
and limiting regimes treated there do not supply that equality along
$\xi=\Xi(N-1)+2$ with $N\to\infty$ and fixed $\Xi$.

\begin{proposition}[The single-wall cylinder at coupling four]
\label{prop:intensity-cylinder-critical-identification}
Consider the cylinder $\{u=x+iy:x>0,\ y\in\mathbb R/W\mathbb Z\}$,
with $N$ counterions, a boundary line charge density $\sigma=N/W$,
and coupling $\Gamma>0$.  The density
\[
 \exp\left(-\Gamma\pi\sigma\sum_i x_i\right)
 \prod_{i<j}\left|2\sinh\frac{\pi(u_i-u_j)}W\right|^\Gamma
 \prod_i dx_i\,dy_i
\]
is transformed by $\zeta_i=e^{-2\pi u_i/W}$ into a density proportional to
\[
 |\Delta_N(\zeta)|^\Gamma
 \prod_i|\zeta_i|^{\Gamma/2-2}
 \mathbf 1_{\{0<|\zeta_i|<1\}}\,dA(\zeta_i).
\]
In particular, at $\Gamma=4$ this is exactly the original unweighted
disk gas at $\beta=4$.
\end{proposition}

\begin{proof}
Put $r_i=|\zeta_i|$.  The wall weight is
$r_i^{\Gamma N/2}$.  Each pair contributes
\[
 \left|2\sinh\frac{\pi(u_i-u_j)}W\right|^\Gamma
 =|\zeta_i-\zeta_j|^\Gamma(r_i r_j)^{-\Gamma/2}.
\]
Finally $dx_i\,dy_i=(W/(2\pi))^2r_i^{-2}dA(\zeta_i)$.
The exponent for each $r_i$ is therefore
$\Gamma N/2-\Gamma(N-1)/2-2=\Gamma/2-2$.
\end{proof}

In \href{https://arxiv.org/pdf/1012.4983}{\v{S}amaj--Trizac (2011)},
Section IV.B, equations (4.21)--(4.23), the dimensionless occupation
coefficient at $\Gamma=4$ satisfies
\[
 f^{(2)}_{0,N}=2\pi\rho_{1,N}(0).
\]
Indeed their moments are $w_j=1/(j+1)$ and their cylinder density,
after the change of area in the preceding proof, gives
$\rho_{1,N}(r)=\pi^{-1}\sum_j\langle\Xi_j\Psi_j\rangle r^{2j}$;
their definition is $f^{(2)}_{0,N}=2w_0\langle\Xi_0\Psi_0\rangle$.
Here the symbols $\Xi_j,\Psi_j$ are the Grassmann composites of that
paper, not point processes.

The paper extrapolates its exact finite-$N$ data with $N\le10$ to
$f^{(2)}_{0,N}\to0.195$; see Figure 5 and its accompanying discussion.
This would predict a positive limiting disk intensity
$0.195/(2\pi)$ at the origin.  It is a numerical extrapolation, not a
proof of that limit.  Theorems~\ref{thm:critical-empty} and
\ref{thm:critical-mobius-empty}, together with locally uniform intensity
convergence, instead imply
\[
 f^{(2)}_{0,N}\longrightarrow0.
\]
Neither critical proof supplies a convergence rate, so this qualitative
conclusion does not by itself quantify the finite-size error in the
published fit.

\begin{remark}[What the Poisson--Boltzmann comparison establishes]
The disk profile in \cite[Section 4.2]{Samaj} is
$4a^2/[\pi\beta(a^2-r^2)^2]$, with
$a^2=R^2+4R^2/(\beta N)$, within the Poisson--Boltzmann approximation.
At $R=1$ its formal $N\to\infty$ limit has coefficient
$4/(\pi\beta)$.  This supports the leading weak-coupling term but
does not determine the proposed constant correction $-1/\pi$.
Moreover, the planar limit with fixed physical distance from a wall
and the limit at fixed relative disk coordinate are distinct limiting
operations.  Neither interchange is assumed here.
\end{remark}


\section{Finite-particle numerical checks of the proposed intensity}
\label{sec:intensity-finite-n-numerics}

This section reports finite-particle computations, not additional limit theorems.
The proposed intensity would imply
\begin{equation}
 \lim_{N\to\infty}\mathbb E M_N(r)
 =\left(\frac4\beta-1\right)\frac{r^2}{1-r^2}.
 \label{eq:numerical-count-prediction}
\end{equation}
The goal of the computation was to test this prediction across several
particle numbers and temperatures, while checking the sampler against the
exact finite-particle answer at $\beta=2$.

\subsection{The target law and the transition kernel}

The target was exactly the unweighted disk law
$Z_{N,\beta}^{-1}|\Delta_N(z)|^\beta\prod_i dA(z_i)$.
Write $z_i=e^{-t_i}e^{\mathrm i\theta_i}$, with $t_i>0$ and
$\theta_i\in\mathbb R/(2\pi\mathbb Z)$.  In these coordinates its density,
up to a constant, is
\[
 \exp\!\left(\beta\sum_{i<j}\log|z_i-z_j|-2\sum_i t_i\right).
\]
A randomly selected particle was updated by one of three proposals:
with probability $0.4$, change its angle by an independent normal variable
of standard deviation $2\pi/N$; with probability $0.4$, multiply its depth
$t_i$ by $e^{1.4G}$, where $G$ is standard normal; and with probability $0.2$,
redraw its angle uniformly and its depth from
\[
 q_N(t)=0.85N e^{-Nt}+0.3e^{-2t},\qquad t>0.
\]
The logarithm of the Metropolis acceptance ratio contains the change in
$\beta\sum_{j\ne i}\log|z_i-z_j|$.  Its additional term is respectively
\[
 0,\qquad
 \log(t_i'/t_i)-2(t_i'-t_i),\qquad
 -2(t_i'-t_i)+\log q_N(t_i)-\log q_N(t_i').
\]
These terms include the polar-area Jacobian and the proposal-density ratio.
For the ideal kernels in exact arithmetic, detailed balance follows directly:
if $p$ denotes the target density and $q$ a proposal density, the off-diagonal
accepted probability flux is
$\min\{p(x)q(x,y),p(y)q(y,x)\}$, which is symmetric in $x,y$.
Thus this computation samples the original finite-particle model, subject to
ordinary floating-point and finite-run errors; it introduces no extra radial
potential and no spatial scaling.

\subsection{Run length, diagnostics, and uncertainty}

Four chains were run independently for each reported pair $(N,\beta)$,
using seeds $93017$, $46023$, $75263$, and $18349$.
A sweep consisted of $N$ random particle-update attempts.
For $N\le128$, each chain discarded $10000$ sweeps and then recorded every
tenth sweep during $100000$ further sweeps.  For $N=512$, the corresponding
numbers were $15000$ discarded and $60000$ retained sweeps, again recording
every tenth sweep.  Initial depths were independent exponential variables
of rates $2(i+1)$, $i=0,\ldots,N-1$, and initial angles were independent
uniform points in the $N$ consecutive equal angular cells.  This
initialization was not assumed to have the target distribution.

The errors displayed below are one estimated Monte Carlo standard error.
They were computed from the sample variance of nonoverlapping batch means,
using batches of $500$ recorded observations, or $5000$ sweeps.  There were
$80$ such batches in total when $N\le128$ and $48$ when $N=512$.
For the counts at radii $0.25,0.5,0.7,0.8,0.9$, classical split-chain
$\widehat R$ was below $1.001$ in every reported run.  Estimated integrated
autocorrelation times, measured in recorded observations and using the initial
monotone positive sequence of adjacent autocorrelation sums, were between
$1$ and $2.6$.
These are empirical diagnostics.  No finite-time mixing bound or rigorous
coverage statement for the displayed error bars is asserted.
In particular, the error bars do not include finite-$N$ bias relative to an
infinite-particle limit.

\begin{center}
\small
\begin{tabular}{rrcc}
\hline
$\beta$ & $N$ & Estimated $\mathbb E M_N(0.5)$ & Estimated $\mathbb E M_N(0.8)$\\
\hline
1 & 32 & $0.91173\pm 0.00452$ & $4.44265\pm 0.00750$ \\
1 & 128 & $0.97515\pm 0.00419$ & $5.07308\pm 0.00683$ \\
1 & 512 & $0.99667\pm 0.00480$ & $5.26092\pm 0.01133$ \\
1 & $\infty$ (conj.) & $1.00000$ & $5.33333$ \\\hline
2 & 128 & $0.33278\pm 0.00268$ & $1.76978\pm 0.00611$ \\
2 & $\infty$ (exact) & $0.33333$ & $1.77778$ \\\hline
3 & 32 & $0.12258\pm 0.00200$ & $0.70400\pm 0.00475$ \\
3 & 128 & $0.11253\pm 0.00197$ & $0.62623\pm 0.00410$ \\
3 & 512 & $0.11096\pm 0.00280$ & $0.61133\pm 0.00545$ \\
3 & $\infty$ (conj.) & $0.11111$ & $0.59259$ \\\hline
3.5 & 32 & $0.06630\pm 0.00137$ & $0.41600\pm 0.00346$ \\
3.5 & 128 & $0.05683\pm 0.00144$ & $0.32870\pm 0.00340$ \\
3.5 & 512 & $0.05671\pm 0.00156$ & $0.30562\pm 0.00515$ \\
3.5 & $\infty$ (conj.) & $0.04762$ & $0.25397$ \\\hline
\end{tabular}
\end{center}

The rows marked ``conj.'' are evaluations of
\eqref{eq:numerical-count-prediction}; they are not extrapolations from the data.

\subsection{Exact calibration and what the experiment does not establish}

At $\beta=2$, the finite-particle kernel is
\[
 K_N(z,w)=\frac1\pi\sum_{k=1}^N k(z\overline w)^{k-1}.
\]
Integrating its diagonal over $B_r$ gives the exact finite-$N$ formula
\[
 \mathbb E M_N(r)=\sum_{k=1}^N r^{2k}
 =\frac{r^2(1-r^{2N})}{1-r^2}.
\]
Similarly, radial integration of $-\log|z|$ gives
\[
 \mathbb E\sum_{i=1}^N-\log|z_i|=\frac12\sum_{k=1}^N\frac1k.
\]
At $N=128$ the latter exact value is $2.7165735463$; the simulation gave
$2.717462\pm0.003440$.
The exact count means at radii $0.5$ and $0.8$ differ negligibly, at the
precision displayed, from $1/3$ and $16/9$; both are compatible with the
Monte Carlo estimates in the table.

The other temperatures exhibit visible dependence on $N$.
For example, at $\beta=3.5$ and $N=512$, the estimated count at radius
$0.5$ exceeds the conjectured limiting value by about $5.8$ estimated
Monte Carlo standard errors.  This is a discrepancy between a finite-$N$
mean and a proposed limit, so it is not a counterexample to the conjecture.
Conversely, the proximity of some of the other entries to their conjectured
limits is not evidence sufficient to identify an exact coefficient.
The experiment neither supplies a justified bias expansion nor rules out
nearby alternative limiting intensities.  In particular, fitting a universal
$1/N$ correction would be an additional unproved assumption.

The reproducibility bundle includes the C++ sampler, the Python run and
analysis scripts, the summary statistics, and the recorded observable traces.


\section{Euclidean coordinates for the intensity and its screening defect}
\label{sec:screening-euclidean-coordinates}

The following change of coordinates turns the intensity conjecture into
an assertion of constant density relative to ordinary Euclidean area.
The interaction and the elliptic operator must be transformed at the
same time.  We make both transformations explicit before testing the
proposed energy argument.

\begin{proposition}[Area coordinates and the transformed interaction]
\label{prop:screening-euclidean-area-map}
The maps
\[
 T(z)=\frac{z}{\sqrt{1-|z|^2}},\qquad
 F(x)=\frac{x}{\sqrt{1+|x|^2}}
\]
are mutually inverse smooth diffeomorphisms between $\D$ and
$\mathbb R^2$.  If $d\lambda(z)=dA(z)/(1-|z|^2)^2$, then
$F^*d\lambda=dx$.  The image of the finite disk gas has density
$\widehat Z_{N,\beta}^{-1}e^{-\beta\widehat H_{N,\beta}}$ relative
to $dx_1\cdots dx_N$, where
\begin{equation}
 \widehat H_{N,\beta}(x_1,\ldots,x_N)
 =-\sum_{i<j}\log|F(x_i)-F(x_j)|
       +\frac2\beta\sum_i\log(1+|x_i|^2).
 \label{eq:screening-euclidean-finite-hamiltonian}
\end{equation}
For any point process with intensity $\rho(z)dA(z)$, its image has
intensity
\begin{equation}
 \widetilde\rho(x)=\frac{\rho(F(x))}{(1+|x|^2)^2}.
 \label{eq:screening-euclidean-intensity-transform}
\end{equation}
Consequently the proposed disk intensity is equivalent to
$\widetilde\rho\equiv c_\beta=(4/\beta-1)/\pi$.
\end{proposition}

\begin{proof}
The radial map $R\mapsto R/\sqrt{1+R^2}$ is strictly increasing
from $[0,\infty)$ onto $[0,1)$ and has the stated inverse.  At
$x\ne0$, the derivative $DF$ has radial eigenvalue
$(1+|x|^2)^{-3/2}$ and tangential eigenvalue
$(1+|x|^2)^{-1/2}$.  Both extend smoothly and positively at zero.
Thus $\det DF=(1+|x|^2)^{-2}$, whereas
$1-|F(x)|^2=(1+|x|^2)^{-1}$.  These identities prove the area
assertion, the transformed Hamiltonian by changing variables in each
particle coordinate, and the intensity formula by the Campbell
identity.  The last conclusion follows by substitution.
\end{proof}

\begin{proposition}[The transformed Green operator and mean defect]
\label{prop:screening-euclidean-elliptic-operator}
Define
\[
 A(x)=\frac{I+(2+|x|^2)xx^{\mathsf T}}{1+|x|^2},\qquad
 \mathcal L=\operatorname{div}(A\nabla).
\]
The matrix $A$ has radial and tangential eigenvalues $1+|x|^2$
and $(1+|x|^2)^{-1}$, respectively, and determinant one.  For
smooth functions $v,w$ of compact support in $\D$, writing
$V=v\circ F$ and $W=w\circ F$, one has
\begin{equation}
 \int_\D\nabla v\cdot\nabla w\,dA
       =\int_{\mathbb R^2}\nabla V\cdot A\nabla W\,dx.
 \label{eq:screening-dirichlet-change-of-variables}
\end{equation}
The transformed disk Green function
\[
 \Gamma(x,y)=-\log\left|
       \frac{F(x)-F(y)}{1-F(x)\overline{F(y)}}\right|
\]
satisfies $-\mathcal L_x\Gamma(x,y)=2\pi\delta_y$ in
distributions and tends to zero as $|x|\to\infty$ for fixed $y$.
For the normalized mean defect $u$ of
Proposition~\ref{prop:intensity-mean-potential-defect},
$U=u\circ F$ satisfies exactly
\begin{equation}
 \mathcal LU=2\pi(c_\beta-\widetilde\rho),\qquad U(0)=0.
 \label{eq:screening-euclidean-defect-equation}
\end{equation}
The same identities hold for every finite-particle mean defect.
\end{proposition}

\begin{proof}
The chain rule and change of area show that the matrix in the
Dirichlet form is
$(\det DF)DF^{-1}DF^{-\mathsf T}$.  The derivative eigenvalues
computed above give the stated matrix and its eigenvalues.  This
proves~\eqref{eq:screening-dirichlet-change-of-variables}.
Apply this identity locally to the weak disk Green equation
$-\Delta_z g(z,w)=2\pi\delta_w$.  The logarithmic singularity
and its first derivative are locally integrable, so the change of
variables is valid after excising a small circle and then letting its
radius decrease to zero.  Equivalently it follows directly from the
weak equation tested against a smooth compactly supported function
composed with $T$.  This proves the asserted point source, including
its coefficient.  As $|x|\to\infty$, $|F(x)|\to1$ and
$g(F(x),F(y))\to0$ uniformly in the angular coordinate, proving
the boundary assertion.

In general, for smooth $u$ the same weak change of variables yields
$\mathcal L(u\circ F)=(\det DF)(\Delta u)\circ F$.
Substitute the proved equation
$\Delta u=2\pi(c_\beta(1-|z|^2)^{-2}-\rho(z))$ and the
Jacobian identity.  This gives
\eqref{eq:screening-euclidean-defect-equation}; the normalization
is unchanged because $F(0)=0$.
\end{proof}

The coefficients are uniformly elliptic on every bounded set, but
their eigenvalue ratio grows like $(1+|x|^2)^2$.  Thus ordinary
whole-plane Coulomb screening estimates do not transfer without
checking the operator and the boundary terms.  The next results
identify these boundary terms exactly.


\subsection{Radial screening energy and the boundary at infinity}
\label{subsec:screening-boundary-energy}

The area-flattening coordinates separate an exact local identity from
the boundary assumptions needed for a global energy argument.  Write
\[
 F(x)=\frac{x}{\sqrt{1+|x|^2}},\qquad
 q=\frac4\beta-1,\qquad c=\frac q\pi,
 \qquad \widetilde\rho(x)=\frac{\rho(F(x))}{(1+|x|^2)^2}.
\]
Let $U(x)=u(F(x))$, with $u$ as in
Proposition~\ref{prop:intensity-mean-potential-defect}, and set
\[
 A(x)=\frac{I+(2+|x|^2)xx^{\mathsf T}}{1+|x|^2},\qquad
 \mathcal L=\operatorname{div}(A\nabla),\qquad
 M(R)=2\pi\int_0^R s\widetilde\rho(s)\,ds.
\]
Here and below $U(R)$ denotes the value of a radial function on the
circle of Euclidean radius $R$.  The radial eigenvalue of $A$ is
$1+R^2$, so $\mathcal LU=R^{-1}(R(1+R^2)U')'$.

\begin{proposition}[Exact local energy identity]
\label{prop:screening-local-energy}
For an actual subsequential limit,
\begin{align}
 D(R)&:=qR^2-M(R)=R(1+R^2)U'(R),
 \label{eq:screening-discrepancy-flux}\\
 \mathcal E(R)&:=\int_{B_R}\nabla U\cdot A\nabla U\,dx
       =2\pi\int_0^R\frac{D(s)^2}{s(1+s^2)}\,ds.
 \label{eq:screening-radial-energy}
\end{align}
For every real constant $a$ and every finite $R>0$,
\begin{equation}
 \mathcal E(R)=2\pi\{U(R)-a\}D(R)
       +2\pi\int_{B_R}(U-a)(\widetilde\rho-c)\,dx.
 \label{eq:screening-local-energy-boundary}
\end{equation}
In particular, the proposed intensity is equivalent to
$\mathcal E(R)=0$ for every finite $R$.
\end{proposition}

\begin{proof}
The mean-defect equation and the change of variables give
$\mathcal LU=2\pi(c-\widetilde\rho)$.  Integrating this equation
over $B_R$ and using the divergence theorem yields the first identity;
the origin contributes no flux since $U$ is $C^2$ and radial.
Polar integration of the energy gives the second identity.
Multiply the equation by $U-a$ and integrate by parts on $B_R$.
The boundary integral is
$2\pi R(1+R^2)(U(R)-a)U'(R)$, which is the first term in
\eqref{eq:screening-local-energy-boundary}.  All integrals are finite
on a bounded disk, so no assumption at infinity has entered.
Finally, vanishing of all the nonnegative energies forces $U'=0$;
the normalization $U(0)=0$ then forces $U=0$.  The mean-defect
proposition identifies this with the proposed intensity.
\end{proof}

A radial finite-energy function has a trace at infinity.  Subtracting
that trace removes the need to postulate a separate boundary-flux
condition.  It also changes the sign estimate one would need to prove.

\begin{proposition}[A finite-energy criterion with the correct constant]
\label{prop:screening-energy-trace-criterion}
Suppose that $U$ is radial and $C^2$ on $\mathbb R^2$, that
$\mathcal LU=2\pi(c-\widetilde\rho)$ with continuous
$\widetilde\rho$, and that
\[
 \mathcal E(\infty):=\int_{\mathbb R^2}\nabla U\cdot A\nabla U\,dx
 <\infty.
\]
Then a finite limit $\ell=\lim_{R\to\infty}U(R)$ exists.  If
$\mathcal E_{\rm tail}(R)=\mathcal E(\infty)-\mathcal E(R)$, then
\begin{equation}
 |U(R)-\ell|^2\le
 \frac{\mathcal E_{\rm tail}(R)}{4\pi}
          \log(1+R^{-2}),\qquad R>0.
 \label{eq:screening-finite-energy-trace}
\end{equation}
Fix a smooth function $\chi:[0,\infty)\to[0,1]$ that equals one
on $[0,1]$ and zero on $[2,\infty)$, and put
$\chi_R(x)=\chi(|x|/R)$.  Then
\begin{equation}
 \mathcal E(\infty)=2\pi\lim_{R\to\infty}
       \int_{\mathbb R^2}\chi_R(U-\ell)
                         (\widetilde\rho-c)\,dx.
 \label{eq:screening-renormalized-energy-identity}
\end{equation}
The limit in this formula exists without an assumption of absolute
integrability of its integrand over the entire plane.  Consequently,
the additional estimate
\begin{equation}
 \limsup_{R\to\infty}\int\chi_R(U-\ell)
                         (\widetilde\rho-c)\,dx\le0
 \label{eq:screening-renormalized-sign-criterion}
\end{equation}
forces $\widetilde\rho=c$ everywhere.  If $U(0)=0$, it also forces
$U=0$ and $\ell=0$.
\end{proposition}

\begin{proof}
For $S>R>0$, Cauchy--Schwarz gives
\[
 |U(S)-U(R)|^2
 \le\left(\int_R^S s(1+s^2)U'(s)^2\,ds\right)
     \left(\int_R^S\frac{ds}{s(1+s^2)}\right).
\]
The first factor is bounded by
$\mathcal E_{\rm tail}(R)/(2\pi)$, and the second is bounded by
$\tfrac12\log(1+R^{-2})$.  Thus $U(R)$ is Cauchy as $R$ tends
to infinity.  Taking $S\to\infty$ proves
\eqref{eq:screening-finite-energy-trace}.

Set $v=U-\ell$.  Testing the equation against the compactly
supported function $\chi_Rv$ gives
\[
 \int\chi_R\nabla U\cdot A\nabla U
 =2\pi\int\chi_Rv(\widetilde\rho-c)
       -\int v\,\nabla\chi_R\cdot A\nabla U.
\]
The final integral is supported on $R\le |x|\le2R$.  For $R\ge1$,
the trace bound, $\log(1+s^{-2})\le s^{-2}$, and
$|\chi_R'(s)|\le\|\chi'\|_\infty/R$ show that
\begin{align*}
 \int v^2\nabla\chi_R\cdot A\nabla\chi_R\,dx
 &=2\pi\int_R^{2R}v(s)^2s(1+s^2)\chi_R'(s)^2\,ds\\
 &\le C_\chi\mathcal E_{\rm tail}(R).
\end{align*}
Another Cauchy--Schwarz inequality therefore bounds the absolute
value of the final integral by
$C_\chi^{1/2}\mathcal E_{\rm tail}(R)$, which tends to zero.
Dominated convergence sends the energy term on the left to
$\mathcal E(\infty)$ and proves
\eqref{eq:screening-renormalized-energy-identity}.  Under
\eqref{eq:screening-renormalized-sign-criterion}, this nonnegative
energy is zero.  Positive definiteness of $A$ gives $\nabla U=0$;
the differential equation then gives $\widetilde\rho=c$.
\end{proof}

\begin{remark}[The constant in the test function matters]
Replacing $U-\ell$ by $U$ adds the term
$\ell\int\chi_R(\widetilde\rho-c)$ to the work integral.
Finite energy does not, by itself, control this net-charge term.
Thus a sign estimate using $U$ and one using $U-\ell$ are not
interchangeable without further information.  No finite-energy
estimate or either sign estimate is asserted here for an actual
limit away from $\beta=2$.
\end{remark}

\subsection{Energy escapes to infinity even at the determinantal point}
\label{subsec:screening-beta2-energy-escape}

Uniform finite-particle control of the energy on the whole plane
cannot be the missing estimate: that assertion already fails at
the exactly solved temperature.

\begin{proposition}[Finite-particle obstruction and an exact benchmark]
\label{prop:screening-finite-particle-energy-escape}
Use the finite-particle mean defect $U_N=u_N\circ F$ and the
corresponding expected count $M_N$ and energy $\mathcal E_N$.
For every fixed $N$ and every $\beta$ with $q\ne0$,
\begin{equation}
 \lim_{R\to\infty}\frac{\mathcal E_N(R)}{R^2}=\pi q^2,
 \qquad \mathcal E_N(\infty)=\infty.
 \label{eq:screening-finite-n-infinite-energy}
\end{equation}
Writing $H_N=\sum_{z\in\Xi_N}-\log|z|$, one also has
\begin{equation}
 0\le U_N(R)-\frac q2\log(1+R^2)+\mathbb EH_N
       \le\frac N2\log(1+R^{-2}).
 \label{eq:screening-finite-n-mean-boundary}
\end{equation}
If $0<\beta<4$ and $a>q^{-1/2}$, then, without using the
determinantal structure,
\begin{equation}
 \liminf_{N\to\infty}\frac{\mathcal E_N(a\sqrt N)}N
 \ge\pi\int_{1/q}^{a^2}\left(q-\frac1v\right)^2\,dv>0.
 \label{eq:screening-all-beta-mesoscopic-lower}
\end{equation}
At $\beta=2$, put $t_R=R^2/(1+R^2)$.  Then exactly
\begin{align}
 M_N(R)&=R^2(1-t_R^N),\qquad D_N(R)=R^2t_R^N,
 \label{eq:screening-beta2-count-discrepancy}\\
 U_N(R)&=\frac12\sum_{k=N+1}^\infty\frac{t_R^k}{k},
 \label{eq:screening-beta2-potential}\\
 \mathcal E_N(R)&=\pi\int_0^{R^2}
             \left(\frac v{1+v}\right)^{2N+1}\,dv.
 \label{eq:screening-beta2-energy}
\end{align}
Hence $\mathcal E_N(R)\to0$ for each fixed $R$, whereas, for each
$a>0$,
\begin{equation}
 \lim_{N\to\infty}\frac{\mathcal E_N(a\sqrt N)}N
       =\pi\int_0^{a^2}e^{-2/v}\,dv>0.
 \label{eq:screening-beta2-mesoscopic-energy}
\end{equation}
For arbitrary $R>0$ one has the explicit bound
\begin{equation}
 \mathcal E_N(R)\le\pi R^2
       \exp\left(-\frac{2N+1}{1+R^2}\right).
 \label{eq:screening-beta2-growing-cutoff}
\end{equation}
In particular, along any growing radii $R_N$ for which the
right-hand side tends to zero, the energy on $B_{R_N}$ vanishes.
\end{proposition}

\begin{proof}
The local energy identity applies unchanged at finite $N$.
Since $0\le M_N(R)\le N$ and $M_N(R)\to N$, its integrand
satisfies
\[
 \frac{D_N(R)^2}{R(1+R^2)}=q^2R+O_{N,q}(R^{-1})
       \qquad(R\to\infty).
\]
Integration proves \eqref{eq:screening-finite-n-infinite-energy}.
For the potential bound, put $r=R/\sqrt{1+R^2}$.  Tonelli's
theorem gives
$\mathbb EH_N=\int_0^1m_N(s)\,ds/s<\infty$.
The finite mean-defect formula consequently yields
\[
 U_N(R)=\frac q2\log(1+R^2)-\mathbb EH_N
                  +\int_r^1\frac{m_N(s)}s\,ds.
\]
 The last integral lies between zero and $N(-\log r)$, which is
the upper bound in \eqref{eq:screening-finite-n-mean-boundary}.

For $q>0$ and $s\ge\sqrt{N/q}$ the elementary bound
$M_N(s)\le N$ implies $D_N(s)\ge qs^2-N\ge0$.  Restrict
the energy integral to this range and then set $v=s^2/N$.
For $a>q^{-1/2}$ this gives
\[
 \frac{\mathcal E_N(a\sqrt N)}N
 \ge\pi\int_{1/q}^{a^2}
        \frac{(qv-1)^2}{v(v+1/N)}\,dv.
\]
Dominated convergence on this fixed interval proves
\eqref{eq:screening-all-beta-mesoscopic-lower}.

At $\beta=2$ the finite Bergman kernel gives
$m_N(r)=\sum_{k=1}^Nr^{2k}=r^2(1-r^{2N})/(1-r^2)$.
This proves \eqref{eq:screening-beta2-count-discrepancy}.
Integrating $m_N(s)/s$ in the mean-defect formula and using
$-\log(1-t)=\sum_{k\ge1}t^k/k$ proves
\eqref{eq:screening-beta2-potential}.  Substitute the discrepancy
into \eqref{eq:screening-radial-energy}, then set $v=s^2$.
This gives \eqref{eq:screening-beta2-energy}.

For $R=a\sqrt N$, set $v=Nw$ in that identity.  For each $w>0$,
\[
 \left(\frac{Nw}{1+Nw}\right)^{2N+1}\longrightarrow e^{-2/w}.
\]
The integrands lie between zero and one, so dominated convergence
on $[0,a^2]$ proves \eqref{eq:screening-beta2-mesoscopic-energy}.
Finally, the integrand in \eqref{eq:screening-beta2-energy} is
increasing in $v$, and
$\log(v/(1+v))\le-1/(1+v)$.  These two facts prove
\eqref{eq:screening-beta2-growing-cutoff}, which also proves
fixed-radius energy convergence.
\end{proof}

\subsection{Why the established one-point bounds do not supply the sign}
\label{subsec:screening-one-point-obstruction}

The following examples concern candidate intensity functions only.
They are not asserted to be intensities of Gibbs laws or of actual
disk-gas limits.

\begin{proposition}[Localized defects pass the one-point tests]
\label{prop:screening-localized-profile-defects}
Fix $0<\beta<4$.  There exist positive smooth radial intensity
functions $\rho_\varepsilon$ on $\D$, distinct from
$q/[\pi(1-|z|^2)^2]$, such that:
\begin{enumerate}
\item they agree with that proposed intensity outside a fixed compact
annulus in $\D$;
\item their total signed difference from the proposed intensity is zero;
\item their normalized mean defects have finite energy, vanish outside
a compact annulus, and have zero boundary term in
\eqref{eq:screening-local-energy-boundary} for all sufficiently large $R$;
\item they satisfy
$\Delta\log\rho_\varepsilon(z)\ge
2\beta/(1-|z|^2)^2$ everywhere.
\end{enumerate}
For each such nonzero defect, the work integral appearing in either
global sign criterion is strictly positive.
\end{proposition}

\begin{proof}
Choose a nonzero smooth radial function $\psi$ on $\mathbb R^2$
supported in the annulus $1<|x|<2$.  Set
\[
 U_\varepsilon(x)=\varepsilon\psi(x),\qquad
 \widetilde\rho_\varepsilon(x)
       =c-\frac\varepsilon{2\pi}\mathcal L\psi(x),\qquad
 \rho_\varepsilon(z)
       =\frac{\widetilde\rho_\varepsilon(F^{-1}(z))}
                    {(1-|z|^2)^2}.
\]
Since $c>0$ and $\mathcal L\psi$ is bounded and compactly
supported, $\widetilde\rho_\varepsilon$ is strictly positive for
all sufficiently small $|\varepsilon|>0$.  The equations give
$\mathcal LU_\varepsilon=2\pi(c-\widetilde\rho_\varepsilon)$
and $U_\varepsilon(0)=0$.  Uniqueness of a smooth radial solution
with this normalization, obtained by integrating its radial flux,
shows that $U_\varepsilon$ is exactly the mean defect defined from
this candidate intensity.  Integration of $\mathcal L\psi$ over
the plane gives zero, because $\psi$ is compactly supported.
This proves the first three assertions, including the zero signed
mass after the change of variables.

For the last assertion, the unperturbed profile $\rho_0$ satisfies
\[
 \Delta\log\rho_0(z)=\frac8{(1-|z|^2)^2}.
\]
The difference from the required lower bound is
$(8-2\beta)/(1-|z|^2)^2$, strictly positive.  On the fixed compact
support of the perturbation, $\rho_\varepsilon\to\rho_0$ in
$C^2$ and all these functions have a common positive lower bound
for sufficiently small $|\varepsilon|$.  Thus the strict inequality
persists there; outside that support the profiles agree exactly.
Finally $\ell=0$ and integration by parts gives
\[
 \int U_\varepsilon(\widetilde\rho_\varepsilon-c)\,dx
       =\frac1{2\pi}\int\nabla U_\varepsilon\cdot
                         A\nabla U_\varepsilon\,dx>0.
\]
The strict positivity follows because $\psi$ is nonzero and
compactly supported, hence is not constant.
\end{proof}

These examples isolate the logical issue: neither finite energy,
vanishing boundary terms, the correct exterior profile, zero total
charge discrepancy, nor the proved one-point curvature bound forces
the required sign.  Any successful sign argument must use additional
information about the interacting gas.  This proposition does not
rule out that additional information or the intensity conjecture.


\section{An exact neutral-plasma representation and its energy obstruction}
\label{sec:neutral-ocp-screening}

This section gives an exact finite-volume interpretation of the proposed
intensity coefficient.  The interpretation does not assume screening.
It also identifies two obstacles to importing a screening argument: the
change from the hyperbolic Green interaction is a critical quadratic
tilt, and the finite-volume mean electrostatic energy need not vanish.

Put $\delta(z)=1-|z|^2$ and
\[
 d\lambda(z)=\delta(z)^{-2}dA(z),\qquad
 c_\beta=\frac{4/\beta-1}{\pi}.
\]
The metric with area measure $\lambda$ is
$ds_\lambda^2=|dz|^2/\delta(z)^2$, of curvature $-4$.
Its distance is denoted by $d_\lambda$.
The curvature $-1$ convention is $ds_{\mathbb H}^2=4ds_\lambda^2$;
thus $d_{\mathbb H}=2d_\lambda$ and $dV_{\mathbb H}=4d\lambda$.
In the Euclidean coordinates $x=z/\sqrt{\delta(z)}$,
$d\lambda=d^2x$ and $dV_{\mathbb H}=4d^2x$.
Define the two invariant interactions
\begin{align}
 v_{\rm s}(z,w)
 &=-\log\sinh d_\lambda(z,w)
 =-\log|z-w|+\tfrac12\log\delta(z)
                    +\tfrac12\log\delta(w),
 \label{eq:neutral-sinh-potential}\\
 g(z,w)
 &=-\log\tanh d_\lambda(z,w)
 =-\log|z-w|+\log|1-z\bar w|.
 \label{eq:neutral-green-potential}
\end{align}
Equivalently, the first interaction is $-\log\sinh(d/(2a))$
when the curvature is $-1/a^2$, and $a=1/2$ in our $\lambda$
convention.

\begin{proposition}[Exact neutral one-component plasma]
\label{prop:exact-neutral-ocp}
Let $0<\beta<4$, $c=c_\beta$, and
\[
 R_N^2=\frac{N}{N+\pi c},\qquad D_N=\{z:|z|<R_N\}.
\]
Consider $N$ unit positive charges in $D_N$, a fixed background charge
$-c\,1_{D_N}\lambda$, pair interaction $v_{\rm s}$, and reference
measure $\lambda^{\otimes N}$.  Including the background self-energy,
define
\begin{align}
 H_{N}^{\rm s}(z_1,\ldots,z_N)
 ={}&\sum_{i<j}v_{\rm s}(z_i,z_j)
       -c\sum_i\int_{D_N}v_{\rm s}(z_i,w)d\lambda(w)
 \nonumber\\
 &+\frac{c^2}{2}\iint_{D_N^2}v_{\rm s}(w,w')
                         d\lambda(w)d\lambda(w').
 \label{eq:neutral-ocp-energy}
\end{align}
This background has total charge $-N$.
There is a finite constant $C_{N,c}^{\rm s}$, independent of all particle
positions, such that
\[
 H_N^{\rm s}
 =-\sum_{i<j}\log|z_i-z_j|
   -\frac{1+\pi c}{2}\sum_i\log\delta(z_i)+C_{N,c}^{\rm s}.
\]
Consequently the canonical probability density
$e^{-\beta H_N^{\rm s}}/\mathcal Z_N^{\rm s}$ relative to
$\lambda^{\otimes N}$ is exactly the unweighted disk gas
\[
 \frac{1}{Z_{N,\beta}(R_N)}
 |\Delta_N(z)|^\beta\prod_i1_{D_N}(z_i)dA(z_i).
 \label{eq:neutral-ocp-unweighted}
\]
If $X_{N,\beta}$ denotes the original unit-disk gas, this law is the
law of $R_NX_{N,\beta}$.  The two sequences have the same subsequential
limits in the local topology.
\end{proposition}

\begin{proof}
Direct radial integration gives
\[
 \lambda(B_R)=\frac{\pi R^2}{1-R^2},
\]
so $c\lambda(D_N)=N$.  For $R<1$ define the finite constants
\[
 I_R^0=\int_{B_R}\log|w|\,d\lambda(w),\qquad
 J_R=\int_{B_R}\log\delta(w)\,d\lambda(w).
\]
The logarithmic singularity is integrable and $\delta$ is bounded below
on $B_R$, so both constants are finite.  The circular mean identity
$\frac1{2\pi}\int_0^{2\pi}\log|r-se^{it}|dt=\log\max(r,s)$
implies, for $0<r<R$,
\[
 \frac{d}{dr}\int_{B_R}\log|r-w|d\lambda(w)
 =\frac{\lambda(B_r)}{r}
 =\frac{\pi r}{1-r^2}.
\]
Integration from zero proves
\begin{equation}
 \int_{B_R}\log|z-w|d\lambda(w)
 =I_R^0-\frac\pi2\log\delta(z),\qquad |z|<R.
 \label{eq:neutral-log-background}
\end{equation}
The pair part of \eqref{eq:neutral-ocp-energy} contributes
\[
 -\sum_{i<j}\log|z_i-z_j|
 +\frac{N-1}{2}\sum_i\log\delta(z_i).
\]
The particle-background part, by
\eqref{eq:neutral-sinh-potential} and
\eqref{eq:neutral-log-background}, contributes
\[
 -\frac{\pi c+N}{2}\sum_i\log\delta(z_i)
 +cNI_{R_N}^0-\frac{cN}{2}J_{R_N}.
\]
Thus one may take
\[
 C_{N,c}^{\rm s}
 =cNI_{R_N}^0-\frac{cN}{2}J_{R_N}
 +\frac{c^2}{2}\iint_{D_N^2}v_{\rm s}(w,w')
                                      d\lambda(w)d\lambda(w').
\]
This is finite by the same local logarithmic integrability.
Multiplying the Boltzmann factor by
$\prod_i\delta(z_i)^{-2}dA(z_i)$ leaves the exponent
$\beta(1+\pi c)/2-2=0$ on every $\delta(z_i)$.
In particular,
$\mathcal Z_N^{\rm s}=e^{-\beta C_{N,c}^{\rm s}}Z_{N,\beta}(R_N)$.

The substitution $z_i=R_Nu_i$ gives
\[
 Z_{N,\beta}(R_N)
 =R_N^{2N+\beta N(N-1)/2}Z_{N,\beta}(1),
\]
and proves the asserted equality in law.  Finally $R_N\to1$.
For $f\in C_c^1(\mathbb D)$ choose a compact set $K\Subset\mathbb D$
containing both its support and all its inverse images under the
scalings $z\mapsto R_Nz$, for sufficiently large $N$.
Then
\[
 |(R_NX_{N,\beta})(f)-X_{N,\beta}(f)|
 \le (1-R_N)\|\nabla f\|_\infty X_{N,\beta}(K).
\]
The compact count tightness already proved makes this tend to zero in
probability.  A countable convergence-determining collection of such
$f$ proves the local-limit assertion.
\end{proof}

\begin{remark}[What neutrality does and does not establish]
The proposition identifies the conjectured coefficient with a genuine
finite-volume background, rather than with a formal mean-field
parameter.  It does not establish that the local particle density
converges to that background density.  Global neutrality says only that
the integral of the finite-volume charge density vanishes.  The boundary
of these hyperbolic disks has size comparable to their volume, so a
boundary layer cannot be discarded merely by dividing by volume.
\end{remark}

\begin{proposition}[The Green comparison is a critical quadratic tilt]
\label{prop:neutral-critical-mode-tilt}
At the same $c,N,D_N$, define $H_N^{\rm g}$ by replacing
$v_{\rm s}$ by $g$ throughout \eqref{eq:neutral-ocp-energy}.
There are finite position-independent constants $C_{N,c}^{\rm g}$ and
$C_{N,c}^{\rm s}$ such that
\begin{align*}
 H_N^{\rm g}-C_{N,c}^{\rm g}
 &=\sum_{i<j}g(z_i,z_j)-\frac{\pi c}{2}\sum_i\log\delta(z_i),\\
 (H_N^{\rm s}-C_{N,c}^{\rm s})
 -(H_N^{\rm g}-C_{N,c}^{\rm g})
 &=\frac12\sum_{k=1}^{\infty}\frac{|S_k|^2}{k},
 \qquad S_k=\sum_i z_i^k.
\end{align*}
If $P_N^{\rm s}$ and $P_N^{\rm g}$ are the corresponding canonical
probability laws, then
\begin{equation}
 \frac{dP_N^{\rm g}}{dP_N^{\rm s}}
 =\frac{\exp\{\frac\beta2\sum_{k\ge1}|S_k|^2/k\}}
 {\mathbb E_{P_N^{\rm s}}
       \exp\{\frac\beta2\sum_{k\ge1}|S_k|^2/k\}}.
 \label{eq:neutral-critical-tilt}
\end{equation}
The denominator is finite for every $N$ but tends to infinity as
$N\to\infty$.  In fact, for every fixed integer $K\ge1$,
\begin{equation}
 \mathbb E_{P_N^{\rm s}}
   \exp\left\{\frac\beta2\sum_{k=1}^{K}\frac{|S_k|^2}{k}\right\}
 \longrightarrow\infty.
 \label{eq:neutral-critical-moment-divergence}
\end{equation}
\end{proposition}

\begin{proof}
For every $|z|,|w|<1$,
$\log|1-z\bar w|=-\operatorname{Re}\sum_{k\ge1}z^k\bar w^k/k$.
The convergence is uniform on $\overline D_N^2$.
The integral of each positive angular mode of the radial background
vanishes.  Thus the background potential of $g$ is the negative of
\eqref{eq:neutral-log-background}, and one may take
\[
 C_{N,c}^{\rm g}
 =cNI_{R_N}^0
 +\frac{c^2}{2}\iint_{D_N^2}g(w,w')d\lambda(w)d\lambda(w').
\]
Subtracting the two configuration-dependent energies gives
\[
 -\sum_{i<j}\log|1-z_i\bar z_j|
 -\frac12\sum_i\log\delta(z_i)
 =-\frac12\sum_{i,j}\log|1-z_i\bar z_j|
 =\frac12\sum_{k\ge1}\frac{|S_k|^2}{k}.
\]
This proves \eqref{eq:neutral-critical-tilt}.
For fixed $N$ the sum is bounded above by
$N^2\sum_{k\ge1}R_N^{2k}/k<\infty$, so its exponential is integrable.

Under $P_N^{\rm s}$, the already proved harmonic-mode central limit
theorem and the scaling $R_N\to1$ give
$(S_1,\ldots,S_K)\Rightarrow(G_1,\ldots,G_K)$, where the $G_k$
are independent circular complex Gaussians with
$\mathbb E|G_k|^2=2k/\beta$.
For a single such Gaussian,
\[
 \mathbb E\exp\{\beta|G_k|^2/(2k)\}=\infty:
\]
the positive quadratic exponent cancels its Gaussian density exactly.
Apply weak convergence to the bounded continuous truncations of the
exponential in \eqref{eq:neutral-critical-moment-divergence}, and then
let the truncation level increase.  Their limiting expectations are
unbounded, proving that the liminf is infinite.  The full denominator
dominates each truncated-mode moment, and hence also diverges.
\end{proof}

\begin{remark}
Consequently, convergence and asymptotic independence of finitely many
Gaussian modes cannot justify replacing the sinh plasma by the Green
plasma through \eqref{eq:neutral-critical-tilt}.  The required positive
quadratic exponential lies exactly at the Gaussian integrability
threshold.  No assertion about a local difference of the two limiting
point processes follows from this failure of uniform integrability.
\end{remark}

\begin{proposition}[Mean boundary energy persists even at the solved temperature]
\label{prop:neutral-mean-boundary-energy}
Take $\beta=2$, $c=1/\pi$, and $R_N^2=N/(N+1)$.
Let $\sigma_N$ be the signed mean charge measure of the neutral sinh
plasma:
\[
 \sigma_N=\mathbb E_{P_N^{\rm s}}\sum_i\delta_{z_i}
             -\pi^{-1}1_{D_N}\lambda.
\]
Its Green energy
\[
 \mathcal E_N=\frac12\iint g(z,w)d\sigma_N(z)d\sigma_N(w)
\]
is finite and nonnegative, and
\begin{equation}
 \liminf_{N\to\infty}\frac{\mathcal E_N}{N}
 \ge\frac14\int_2^3
 \left(\frac{1-e^{-(t-1)}}{t-1}-\frac1t\right)^2dt>0.
 \label{eq:neutral-energy-lower-bound}
\end{equation}
Nevertheless its particle intensity converges locally to
$\pi^{-1}\delta(z)^{-2}$.
\end{proposition}

\begin{proof}
The last assertion is the Bergman limit and the deterministic scaling
in Proposition~\ref{prop:exact-neutral-ocp}.
The mean measure is radial, supported in $\overline D_N$, has a bounded
density for each fixed $N$, and has total charge zero.
Let $Q_N(r)=\sigma_N(B_r)$ and
$V_N(z)=\int g(z,w)d\sigma_N(w)$.
The angular mean of the image term in $g$ is zero, so direct angular
integration gives
$rV_N'(r)=-Q_N(r)$.
Because $Q_N(r)=0$ for $r\ge R_N$, the potential vanishes outside
$D_N$: it is constant there and has zero boundary value at $|z|=1$.
Integration by parts, or first integration on compact subdisks using
the bounded density, therefore gives
\begin{equation}
 \mathcal E_N
 =\frac1{4\pi}\int_{\mathbb D}|\nabla V_N|^2dA
 =\frac12\int_0^{R_N}\frac{Q_N(r)^2}{r}\,dr.
 \label{eq:neutral-radial-energy}
\end{equation}
Near zero $Q_N(r)=O(r^2)$, so the final integral is finite.

The truncated Bergman kernel on $B_{R_N}$ gives
\[
 \mathbb E_{P_N^{\rm s}}\#B_r
 =\sum_{k=1}^N\left(\frac{r^2}{R_N^2}\right)^k,
 \qquad
 Q_N(r)=\sum_{k=1}^N\left(\frac{r^2}{R_N^2}\right)^k
                  -\frac{r^2}{1-r^2},\qquad r<R_N.
\]
For $t\in[2,3]$ set $r_N(t)^2=1-t/N$.  For all sufficiently large
$N$, these radii lie in $(0,R_N)$, and
\[
 \frac{r_N(t)^2}{R_N^2}
 =1-\frac{t-1}{N}-\frac{t}{N^2}.
\]
The logarithm of this ratio is $-(t-1)/N+O(N^{-2})$, uniformly in
$t\in[2,3]$.  A uniform Riemann-sum limit therefore gives
\[
 \frac{Q_N(r_N(t))}{N}\longrightarrow
 h(t)=\int_0^1e^{-(t-1)s}ds-\frac1t
     =\frac{1-e^{-(t-1)}}{t-1}-\frac1t.
\]
For $t>1$, the inequality $\log t<t-1$ implies $h(t)>0$.
Restrict \eqref{eq:neutral-radial-energy} to the annulus parametrized
by $2\le t\le3$, and use
$|dr|/r=dt/[2N(1-t/N)]$.  This yields
\[
 \frac{\mathcal E_N}{N}
 \ge\frac14\int_2^3
 \frac{(Q_N(r_N(t))/N)^2}{1-t/N}\,dt,
\]
whose limit is the strictly positive right-hand side of
\eqref{eq:neutral-energy-lower-bound}.
\end{proof}

\begin{remark}[Consequence for the proposed screening strategy]
An argument that tries to prove the vanishing of the entire
finite-volume mean-charge energy, or even its vanishing after division
by $N$, cannot work: Proposition~\ref{prop:neutral-mean-boundary-energy}
rules out both conclusions in the explicitly solved case.
This does not rule out an energy argument for the limiting defect.
Such an argument must first isolate the local limit from the moving
boundary layer and justify the limiting energy and boundary terms.
\end{remark}

\paragraph{Scope of the primary literature.}
The algebraic plasma correspondence above is consistent with the
Hamiltonian in Section~7.1 of Jancovici and T\'ellez,
\emph{Two-Dimensional Coulomb Systems on a Surface of Constant Negative
Curvature}, \href{https://arxiv.org/pdf/cond-mat/9801322}{arXiv:cond-mat/9801322}.
Their Section~5 assumes screening in deriving the sum rules; the
general-temperature argument in Section~4.3 for equality of the sinh
and Green bulk correlations likewise uses screening.  Their exact
calculations establish the special-temperature case.
Fantoni, Jancovici and T\'ellez,
\emph{Pressures for a One-Component Plasma on a Pseudosphere},
\href{https://arxiv.org/pdf/cond-mat/0207525}{arXiv:cond-mat/0207525},
derive an explicit grand-canonical density at coupling two;
equation~(4.52) yields the background density deep in the bulk while
allowing boundary excess charge.  These results do not supply the
missing screening estimate at general $\beta$.


\section{The connected-force correction to the screening equation}
\label{sec:screening-connected-force}

The mean Poisson equation does not contain the first Ward equation.
Combining the two equations isolates a connected two-point force whose
work would control the proposed screening. The formulation below keeps
the summation at finite particle number until the limiting force has
been defined; no absolutely convergent infinite pair integral is assumed.

Identify a complex number $a$ with the real vector
$\mathfrak v(a)=(\operatorname{Re}a,-\operatorname{Im}a)$. Put
$\delta(z)=1-|z|^2$, $q=4/\beta-1$, and
\[
 g(z,w)=-\log\left|\frac{z-w}{1-z\bar w}\right|,
 \qquad
 j(z,w)=-\nabla_zg(z,w)
 =\mathfrak v\left(\frac1{z-w}+\frac{\bar w}{1-z\bar w}\right).
\]
Let $\rho_N$ and $\rho_{2,N}$ denote the first and ordered second
correlation densities of the unweighted $N$-particle disk gas. Define
\begin{align}
 h_N(z,w)&=\rho_{2,N}(z,w)-\rho_N(z)\rho_N(w),\notag\\
 I_N(z)&=\int_\D j(z,w)h_N(z,w)\,dA(w),
 &F_N(z)&=I_N(z)/\rho_N(z),
 \label{eq:screening-finite-connected-force}\\
 T_N(z)&=\sum_{w\in\Xi_N}\frac{\bar w}{1-z\bar w},
 &D_N(z)&=\mathbb E_N^{z}\mathfrak v(T_N(z)).\notag
\end{align}
The superscript $z$ in the last expression denotes the full Palm law,
including the particle at $z$. Thus $D_N$ includes the self contribution
$z/\delta(z)$. Each integral in
\eqref{eq:screening-finite-connected-force} is an ordinary absolutely
convergent integral at fixed $N$: the first-order singularity is locally
integrable, the correlation densities are bounded at fixed $N$, and
$j(z,w)$ is bounded near the boundary for fixed interior $z$.

\begin{proposition}[A finite-particle Ward equation with a connected force]
\label{prop:screening-finite-connected-ward}
Let $u_N$ be the normalized mean defect of
Proposition~\ref{prop:intensity-mean-potential-defect}, and put
$f_N(z)=\delta(z)^2\rho_N(z)$. Then, for every $N$ and $\beta>0$,
\begin{equation}
 \nabla\log f_N=-\beta\nabla u_N+\beta F_N-\beta D_N
 \quad\hbox{on }\D.
 \label{eq:screening-finite-connected-ward}
\end{equation}
\end{proposition}

\begin{proof}
Differentiating the finite insertion integral, justified by the locally
integrable powers $|z-w|^{\beta-1}$, gives the first Ward equation
\[
 \nabla\log\rho_N(z)
 =\beta\,\mathbb E_N^{!z}
       \mathfrak v\left(\sum_{w\in\Xi_N}\frac1{z-w}\right).
\]
The reduced Palm expectation of the complete Green force equals
\[
 \mathbb E_N^{!z}\sum_wj(z,w)
 =\int_\D j(z,w)\rho_N(w)\,dA(w)+F_N(z).
\]
Rotation invariance and the angular mean computation in
\eqref{eq:mean-corrected-field-radial-count} give
\[
 a_N(z):=\int_\D j(z,w)\rho_N(w)\,dA(w)
 =\frac{m_N(|z|)}{|z|^2}z,
\]
with value zero at $z=0$. Removing the image force from the complete
Green force therefore gives
\[
 \nabla\log\rho_N
 =\beta\left(a_N+F_N+\frac z\delta-D_N\right).
\]
The exact mean-defect formula gives
$\nabla u_N=-a_N+qz/\delta$. Finally,
$\nabla\log f_N=\nabla\log\rho_N-4z/\delta$ and
$\beta(q+1)=4$, which prove the assertion.
\end{proof}

\begin{proposition}[A limiting connected force without an infinite sum]
\label{prop:screening-limiting-connected-force}
Suppose $\Xi_{N_j}$ converges to a subsequential limit with intensity
$\rho$, and let $\Omega=\{z\in\D:\rho(z)>0\}$. On every compact
subset of $\Omega$, $F_{N_j}$ is bounded in $L^2$, converges weakly
in $L^2$ to a uniquely determined vector field $F$, and $D_{N_j}$
converges weakly in $L^2$ to zero. With $f=\delta^2\rho$ and $u$
the limiting mean defect,
\begin{equation}
 \nabla\log f=-\beta\nabla u+\beta F
 \quad\hbox{weakly on }\Omega.
 \label{eq:screening-limiting-connected-ward}
\end{equation}
This definition does not assert that $F$ equals the result of first
passing to limiting correlation functions in
\eqref{eq:screening-finite-connected-force} and then integrating up
to $\partial\D$.
\end{proposition}

\begin{proof}
Fix $K\Subset\Omega$ and a slightly larger compact set contained in
$\Omega$. Local uniform convergence of the intensities bounds
$\rho_{N_j}$ above and below there. Corollary~\ref{cor:uniform-campbell-force} and conditional
Jensen's inequality give a uniform local bound on
$\int |\nabla\log\rho_{N_j}|^2dA$. The logarithms converge uniformly
on compact sets, and consequently their gradients converge weakly in
local $L^2$ to $\nabla\log\rho$. This last statement follows by weak
compactness and distributional differentiation; every weak cluster
point is the same distributional derivative.

Here is an explicit justification of the analytic supremum moments.
Choose $r<s<1$ with $K\subset\{|z|\le r\}$. Cauchy's integral
formula and Jensen's inequality give, for $p\ge1$,
\[
 \sup_{|z|\le r}|T_N(z)|^p
 \le\left(\frac{s}{s-r}\right)^p
          \frac1{2\pi}\int_0^{2\pi}|T_N(se^{i\theta})|^p\,d\theta.
\]
Rotation invariance makes the expectation of the integrand independent
of $\theta$. At the real point $s$,
$T_N(s)=\overline{\sum_{w\in\Xi_N} w/(1-sw)}$.
Theorem~\ref{thm:analytic-harmonic-laplace}, applied to arbitrary real
and imaginary scalar multiples of the function $w\mapsto w/(1-sw)$,
bounds the exponential moments of both real components uniformly in
$N$. Hence $\sup_N\mathbb E|T_N(s)|^p<\infty$ for every $p\ge1$.
The preceding Cauchy estimate gives the required uniform supremum
moments. The compact-count bounds then imply
\[
 \sup_N\mathbb E\sum_{z\in\Xi_N\cap K}|T_N(z)|^2<\infty
\]
by Cauchy--Schwarz. Palm conditional Jensen gives a uniform bound on
$\int_K\rho_N|D_N|^2dA$, hence on $\|D_N\|_{L^2(K)}$.
For a continuous compactly supported vector test $v$,
\[
 \int\rho_N(z)v(z)\cdot D_N(z)\,dA(z)
 =\mathbb E\sum_{z\in\Xi_N}v(z)\cdot\mathfrak v(T_N(z)).
\]
Joint convergence of the analytic harmonic field and the point process,
and their proved independence in the limit, make the last expression
converge to zero. Indeed, the limiting analytic Gaussian field is
centered. The product is uniformly integrable by the same higher-moment
and count bounds, so this conclusion concerns expectations, not just
convergence in law. To pass from $\rho_ND_N$ to $D_N$, use the fixed
continuous test $v/\rho$, where $v$ has compact support in $K$. The
resulting expression differs from $\int v\cdot D_N$ by at most
\[
 \left\|\frac{\rho_N}{\rho}-1\right\|_{L^\infty(K)}
      \|v\|_{L^2(K)}\|D_N\|_{L^2(K)},
\]
which tends to zero. Density of continuous tests in $L^2$, together
with the uniform local $L^2$ bound, proves that $D_N$ converges weakly
in local $L^2$ to zero.

Now use \eqref{eq:screening-finite-connected-ward}. The local $C^2$
convergence of $u_N$, the weak convergence of the logarithmic gradients,
and the preceding bound on $D_N$ give both the claimed local $L^2$
bound on $F_N$ and the unique weak limit
$F=\nabla u+\beta^{-1}\nabla\log f$.
\end{proof}

\subsection{A compact entropy--energy identity in the plane}

Let $\Psi(x)=x/\sqrt{1+|x|^2}$ map $\mathbb R^2$ onto $\D$, and set
\[
 p(x)=\frac{\rho(\Psi(x))}{(1+|x|^2)^2},\qquad
 U(x)=u(\Psi(x)),\qquad
 \mathcal F(x)=D\Psi(x)^{\mathsf T}F(\Psi(x)).
\]
The letter $p$ here denotes an intensity, not a probability density
of total mass one. Define the positive definite conductivity
\[
 A(x)=\frac{I+(2+|x|^2)xx^{\mathsf T}}{1+|x|^2}.
\]
The Jacobian of $\Psi$ is $(1+|x|^2)^{-2}$ and direct change of
variables in the Dirichlet form gives
$A=(\det D\Psi)D\Psi^{-1}D\Psi^{-\mathsf T}$. Thus the two equations
are exactly
\begin{equation}
 \operatorname{div}(A\nabla U)=2\pi(c-p),\qquad
 \nabla\log p=-\beta\nabla U+\beta\mathcal F,
 \qquad c=\frac{4/\beta-1}{\pi}.
 \label{eq:screening-plane-coupled-equations}
\end{equation}

\begin{theorem}[Entropy, mean-field energy, and connected-force work]
\label{thm:screening-entropy-energy}
Assume $0<\beta<4$. Let $0\le\chi\in C_c^1(\mathbb R^2)$ have
support in the positive-intensity set $\{p>0\}$. Then all the following
integrals are finite, and
\begin{align}
 &2\pi\int\chi(p-c)\log(p/c)\,dx
       +\beta\int\chi\nabla U\cdot A\nabla U\,dx\notag\\
 &\qquad=\beta\int\chi\mathcal F\cdot A\nabla U\,dx
       +\int\log(p/c)\nabla\chi\cdot A\nabla U\,dx.
 \label{eq:screening-compact-entropy-energy}
\end{align}
Suppose in addition that $p>0$ on the plane and there are
$0\le\chi_k\le1$ of compact support with $\chi_k\to1$ pointwise,
such that
\begin{align}
 &\int\log(p/c)\nabla\chi_k\cdot A\nabla U\,dx\longrightarrow0,
 \label{eq:screening-cutoff-boundary-condition}\\
 &\limsup_{k\to\infty}
       \int\chi_k\mathcal F\cdot A\nabla U\,dx\le0.
 \label{eq:screening-connected-work-sign}
\end{align}
Then $p\equiv c$, and hence the proposed disk intensity holds.
No prior assumption of finite total energy or finite total entropy is
needed for this sufficient criterion.
\end{theorem}

\begin{proof}
On the support of $\chi$, $p$ is continuous and bounded away from zero,
$\log(p/c)$ is in $H^1$, $U$ is $C^2$, and $\mathcal F$ is in $L^2$.
The matrices $A$ and $A^{-1}$ are bounded there. Thus all products are
integrable. Taking the scalar product of the second equation in
\eqref{eq:screening-plane-coupled-equations} with $\chi A\nabla U$
and integrating gives
\[
 \int\chi\nabla\log(p/c)\cdot A\nabla U
       +\beta\int\chi\nabla U\cdot A\nabla U
 =\beta\int\chi\mathcal F\cdot A\nabla U.
\]
Weak integration by parts in the first term, using the first equation
of \eqref{eq:screening-plane-coupled-equations}, gives
\[
 \int\chi\nabla\log(p/c)\cdot A\nabla U
 =2\pi\int\chi(p-c)\log(p/c)
       -\int\log(p/c)\nabla\chi\cdot A\nabla U.
\]
This proves \eqref{eq:screening-compact-entropy-energy}.
Both integrands on its left are nonnegative, since $c>0$, $A$ is
positive definite, and $(t-c)\log(t/c)\ge0$ for $t>0$.
Under the two cutoff hypotheses their sum has upper limit at most zero.
Fatou's lemma then gives
$\int(p-c)\log(p/c)=0$. Its integrand vanishes only where $p=c$.
Continuity proves $p\equiv c$. Equivalently, the energy vanishes,
$U$ is constant, and $U(0)=0$ fixes that constant.
\end{proof}

\begin{remark}[The correlation estimate which remains unproved]
The right side of \eqref{eq:screening-compact-entropy-energy} is an
exact correction to an entropy--energy screening argument. Neither
\eqref{eq:screening-cutoff-boundary-condition} nor
\eqref{eq:screening-connected-work-sign} has been proved for a general
disk-gas limit. In particular, $\mathcal F$ is not set to zero by the
independence of the harmonic Gaussian modes: that independence removed
$D_N$ and left the connected Green force $F_N$.
\end{remark}

\subsection{Two checks on proposed sign shortcuts}

Proposition~\ref{prop:screening-localized-profile-defects} already
shows that the normalized mean Poisson equation, the known curvature
inequality, and even compact support of the mean defect do not imply
the screening sign. The following finite-particle example separately
checks that a sign condition on connected correlations alone is also
insufficient.

\begin{example}[Negative connected correlations alone do not determine force work]
\label{ex:screening-negative-correlation-insufficient}
For the actual one-particle disk gas, $\rho_1=1/\pi$ and
$h_1(z,w)=-1/\pi^2<0$ for all distinct $z,w$. Nevertheless,
$F_1(z)=-z$ and
\[
 u_1(z)=-\frac{|z|^2}2-\frac q2\log\delta(z),\qquad
 F_1(z)\cdot\nabla u_1(z)
 =|z|^2\left(1-\frac q{\delta(z)}\right).
\]
For $2<\beta<4$, this work is strictly positive whenever
$0<|z|^2<1-q$. Thus even pointwise negative connected correlations do
not, by themselves, give a nonpositive connected-force work for this
potential. This finite-particle example is not a counterexample to a
limiting screening assertion: its full-Palm term $D_1$ is present.
The example only rules out inferring the desired sign directly from
$h_N\le0$ without an additional argument.
\end{example}


\section{Gibbs variation and the correlation contribution to screening}
\label{sec:screening-gibbs-variation}
\subsection{What Gibbs convexity controls, and the missing fluctuation term}
\label{subsec:screening-gibbs-variational-audit}

An exponential change of an external potential has a definite response
sign.  Its comparison density, however, is the density of another Gibbs
law.  It cannot be replaced by a proposed constant background without
an additional theorem.

\begin{proposition}[The exact response and symmetric entropy identities]
\label{prop:screening-exact-tilt-response}
Let $P$ be a point-process law, let $v$ be a bounded compactly supported
real function, and put $X=\sum_{x\in\Xi}v(x)$.  Assume that
$\mathbb E_P\exp\{a\sum_{x\in\Xi}|v(x)|\}<\infty$ for every $a>0$.  For $0\le t\le1$ set
\[
 \frac{dP_t}{dP}=\frac{e^{\beta tX}}{\mathbb E_Pe^{\beta tX}}.
\]
Write $\rho_t(x)dx$ for the intensity, when it exists.  Then
\begin{align}
 \int v(x)\{\rho_1(x)-\rho_0(x)\}\,dx
 &=\beta\int_0^1\operatorname{Var}_{P_t}(X)\,dt,\label{eq:screening-response-variance}\\
 \operatorname{Ent}(P_1\mid P)+\operatorname{Ent}(P\mid P_1)
 &=\beta\int v(x)\{\rho_1(x)-\rho_0(x)\}\,dx.
 \label{eq:screening-response-symmetric-entropy}
\end{align}
The same formulas hold for finite canonical Gibbs laws without any
extra moment hypothesis.
\end{proposition}

\begin{proof}
The exponential-moment assumption permits two differentiations of
$\psi(t)=\log\mathbb E_Pe^{\beta tX}$ under the expectation, uniformly
on a neighborhood of $[0,1]$.  Thus
$\psi'(t)=\beta\mathbb E_{P_t}X$ and
$\psi''(t)=\beta^2\operatorname{Var}_{P_t}(X)$.  Integrating the
second identity proves~\eqref{eq:screening-response-variance}, using
the Campbell formula.  Its absolute integrability follows from the
assumption, which also implies $\mathbb E_{P_t}\sum_x|v(x)|<\infty$
uniformly for $0\le t\le1$.  The two logarithmic likelihoods are
$\beta X-\psi(1)$ and its negative; integrating them under their
respective laws proves~\eqref{eq:screening-response-symmetric-entropy}.
For a fixed finite number of particles, $X$ is bounded.
\end{proof}

\begin{remark}[The comparison law is part of the missing theorem]
For a compactly supported truncation $v$ of the mean defect, the
proposition gives $\int v(\rho_0-\rho_1)\le0$.  The desired comparison
with a constant $c$ would follow from this particular argument if one
could identify $\rho_1=c$ on the relevant region and control removal
of the truncation.  Neither equality follows from convexity.  In
particular, removing a mean electrostatic potential from the energy
does not remove the fluctuating electrostatic potential.
\end{remark}

The next finite-volume example isolates that distinction using the
same Green pair interaction.  It is a comparison model, not the
unweighted fixed-$N$ disk gas, and its conclusion does not disprove the
intensity conjecture for that gas.

Let $\Pi_c$ be the Poisson process of constant intensity $c>0$ on
$\D$.  Its total count is finite almost surely.  Write
\[
 g(x,y)=-\log\left|\frac{x-y}{1-x\overline y}\right|,
 \qquad v_c(x)=c\int_\D g(x,y)\,dA(y)
             =\frac{\pi c}{2}(1-|x|^2).
\]
The last equality follows either by angular integration or by solving
$-\Delta v_c=2\pi c$ with zero boundary values.  Define
\begin{equation}
 H_c(\eta)=\sum_{\{x,y\}\subset\eta}g(x,y)
                  -\sum_{x\in\eta}v_c(x),\qquad
 \frac{dP_{\beta,c}}{d\Pi_c}
             =Z_{\beta,c}^{-1}e^{-\beta H_c},\quad\beta\ge0.
 \label{eq:screening-neutral-comparison-gas}
\end{equation}
This interaction contains a constant compensating background in its
one-particle term, and its reference activity is also $c$.

\begin{proposition}[An exact free-energy decomposition with correlations]
\label{prop:screening-correlation-energy-obstruction}
The law in~\eqref{eq:screening-neutral-comparison-gas} is well-defined
for every $\beta\ge0$.  For $\beta>0$, let $\rho$ and $\rho^{(2)}$
be its first and second factorial correlation densities, and put
\[
 C^{(2)}(x,y)=\rho^{(2)}(x,y)-\rho(x)\rho(y),\qquad
 f(x)=\rho(x)-c.
\]
All terms below are finite, and
\begin{align}
 &\operatorname{Ent}(P_{\beta,c}\mid\Pi_c)
      +\frac\beta2\iint_{\D^2}g(x,y)f(x)f(y)\,dA(x)dA(y)
 \nonumber\\
 &\hspace{3em}\le
       -\frac\beta2\iint_{\D^2}g(x,y)C^{(2)}(x,y)\,dA(x)dA(y).
 \label{eq:screening-free-energy-correlation-cost}
\end{align}
Moreover,
\begin{align}
 \operatorname{Ent}(P_{\beta,c}\mid\Pi_c)
 &\ge\int_\D\left\{\rho\log\frac\rho c-\rho+c\right\}\,dA,
 \label{eq:screening-ideal-entropy-bound}\\
 \iint_{\D^2}g(x,y)f(x)f(y)\,dA(x)dA(y)&\ge0.
 \label{eq:screening-positive-mean-energy}
\end{align}
Thus the variational principle bounds a mean-density defect by a
correlation-energy contribution.  It does not make the mean-density
energy nonpositive.
\end{proposition}

\begin{proof}
Since $g\ge0$ and $0\le v_c\le\pi c/2$, the bound
$H_c(\eta)\ge-(\pi c/2)\eta(\D)$ implies
$Z_{\beta,c}<\infty$ by the exponential moments of a Poisson random
variable.  The empty configuration has positive Poisson probability,
so $Z_{\beta,c}>0$.  The resulting law has all exponential count
moments.  For every fixed integer $k$, all $k$th moments of $H_c$
under a Poisson law whose intensity has been multiplied by a finite
constant are finite.  Indeed, after conditioning on its finite count,
the inequality $(\sum_{j=1}^m a_j)^k\le m^{k-1}\sum_j a_j^k$ for
$a_j\ge0$ reduces this assertion to integrability of powers of the
logarithmic singularity of $g$ and to Poisson count moments.  On $\D$,
$g(x,y)\le\log(2/|x-y|)$, and every positive power of the latter
function is integrable on $\D^2$.  These observations also prove the
energy and entropy integrability used below.

The Poisson insertion formula gives, for every nonnegative measurable
$f_0$,
\[
 \mathbb E\sum_{x\in\Xi}f_0(x,\Xi\setminus\{x\})
 =c\int_\D\mathbb E\left[
 f_0(x,\Xi)e^{\beta v_c(x)-\beta\sum_{y\in\Xi}g(x,y)}
 \right]dA(x).
\]
To verify the formula, expand the Poisson expectation by total count,
choose the distinguished point among the $n$ coordinates, and cancel
$n/n!$ against $1/(n-1)!$; the quotient of Gibbs weights is the
exponential displayed above.  In particular
$0<\rho(x)\le c e^{\beta\pi c/2}$.  Two insertions similarly give
$0\le\rho^{(2)}(x,y)\le c^2e^{\beta\pi c}$.  The event of no other
points also gives a positive constant lower bound on $\rho$.
Consequently all the displayed Green integrals are absolutely finite.

Jensen's inequality yields
$\log Z_{\beta,c}\ge-\beta\mathbb E_{\Pi_c}H_c$.  The exact entropy
identity is
\[
 \operatorname{Ent}(P_{\beta,c}\mid\Pi_c)
       =-\beta\mathbb E_{P_{\beta,c}}H_c-\log Z_{\beta,c}.
\]
The Poisson factorial moment formulas and the definition of $v_c$
give
\[
 \mathbb E_{\Pi_c}H_c=-\frac{c^2}{2}\iint g,
\]
and a direct expansion gives
\[
 \mathbb E_{P_{\beta,c}}H_c-\mathbb E_{\Pi_c}H_c
   =\frac12\iint g f f+\frac12\iint g C^{(2)}.
\]
Combining these equalities proves
\eqref{eq:screening-free-energy-correlation-cost}.

For any bounded real $h$, nonnegativity of relative entropy with
respect to the exponential tilt of $\Pi_c$ by $\sum h(x)$ gives
\[
 \operatorname{Ent}(P_{\beta,c}\mid\Pi_c)
 \ge\int h\rho\,dA-c\int(e^h-1)\,dA.
\]
Take $h=\log(\rho/c)$, which is bounded by the intensity bounds just
proved, to obtain~\eqref{eq:screening-ideal-entropy-bound}.
Finally, for $u(x)=\int_\D g(x,y)f(y)\,dA(y)$ one has
$u\in H_0^1(\D)$ and $-\Delta u=2\pi f$.  Testing this weak equation
with $u$ gives
\[
 \iint g f f=\int uf\,dA
              =\frac1{2\pi}\int_\D|\nabla u|^2\,dA\ge0.
\]
These assertions follow first for smooth $f$, by the disk Green
formula, and then for bounded $f$ by approximation in $L^2$ and the
energy estimate for the zero-boundary Poisson equation.
\end{proof}

\begin{proposition}[A constant background does not remove fluctuations]
\label{prop:screening-neutral-comparison-response}
For each fixed $x\in\D$, the intensity $\rho_{\beta,c}(x)$ of
\eqref{eq:screening-neutral-comparison-gas} has two right derivatives
at $\beta=0$, with
\begin{equation}
 \rho_{0,c}(x)=c,\qquad
 \left.\partial_\beta\rho_{\beta,c}(x)\right|_{0+}=0,\qquad
 \left.\partial_\beta^2\rho_{\beta,c}(x)\right|_{0+}
       =c^2\int_\D g(x,y)^2\,dA(y)>0.
 \label{eq:screening-neutral-comparison-second-response}
\end{equation}
In particular, $\rho_{\beta,c}(x)>c$ for every sufficiently small
positive $\beta$, with the allowed neighborhood depending on $x$.
For any $\beta>0$, the identity $\rho_{\beta,c}\equiv c$ on $\D$ is
impossible.
\end{proposition}

\begin{proof}
Write $H=H_c$ and
\[
 S_x(\eta)=\sum_{y\in\eta}g(x,y),\qquad T_x=S_x-v_c(x).
\]
The insertion formula gives the exact identity
\begin{equation}
 \frac{\rho_{\beta,c}(x)}c
 =\mathbb E_{P_{\beta,c}}e^{-\beta T_x}
 =\frac{\mathbb E_{\Pi_c}e^{-\beta(H+T_x)}}
        {\mathbb E_{\Pi_c}e^{-\beta H}}.
 \label{eq:screening-neutral-comparison-insertion}
\end{equation}
Here $H\ge-C\eta(\D)$ and
$H+T_x\ge-C\eta(\D)-v_c(x)$.  Polynomial moments of $H$ and $S_x$
under any finite-intensity Poisson law are finite, by the logarithmic
integrability argument in the preceding proof.  Dominated
convergence therefore permits two right differentiations of numerator
and denominator, including continuity of those second derivatives at
zero.

Under $\Pi_c$, the mean of $T_x$ is zero and its variance is
$c\int g(x,y)^2dA(y)$.  We claim also that
$\operatorname{Cov}_{\Pi_c}(T_x,H)=0$.  Indeed, the same insertion
calculation for a linear statistic $L_a=\sum_y a(y)$ gives
\[
 \operatorname{Cov}_{\Pi_c}(L_a,H)
 =c\int_\D a(y)\mathbb E_{\Pi_c}\{H(\eta+\delta_y)-H(\eta)\}
       \,dA(y).
\]
The insertion difference is $S_y-v_c(y)$, whose Poisson mean is zero.
Taking $a(y)=g(x,y)$ is justified by the integrability already proved;
it establishes the claim.  Thus $\mathbb E_{\Pi_c}HT_x=0$.
The first derivative of the quotient
in~\eqref{eq:screening-neutral-comparison-insertion} is zero.  Its
second derivative is
\[
 \mathbb E_{\Pi_c}(H+T_x)^2-\mathbb E_{\Pi_c}H^2
 =\mathbb E_{\Pi_c}T_x^2
 =c\int_\D g(x,y)^2\,dA(y),
\]
which proves~\eqref{eq:screening-neutral-comparison-second-response}
and the asserted strict inequality for small positive $\beta$.

Finally suppose, for a fixed $\beta>0$, that $\rho_{\beta,c}=c$
almost everywhere.  Then for each fixed $x$,
$\mathbb E_{P_{\beta,c}}T_x=\int g(x,y)c\,dA(y)-v_c(x)=0$.
The random variable $T_x$ is not almost surely constant: the empty
configuration has positive probability and gives $S_x=0$, whereas
exactly one particle in any nonempty open set away from $x$ has
positive probability and gives $S_x>0$.  The integrable variable
$e^{-\beta T_x}$ is bounded above by $e^{\beta v_c(x)}$.
Strict Jensen inequality in
\eqref{eq:screening-neutral-comparison-insertion} consequently gives
$\rho_{\beta,c}(x)>c$, a contradiction.
\end{proof}

\begin{remark}[The implication for the intensity conjecture]
The comparison gas has a finite number of particles, a specified
activity, and a zero-boundary Green interaction.  It is not an actual
subsequential limit of the original model.  Its purpose is to establish
precisely why a constant compensating background and convexity of Gibbs
free energy do not by themselves identify the intensity: the potential
variance and the connected pair correlation survive after subtracting
the mean field.  A successful screening proof for the disk-gas limit
must control their contribution and the selected conditional count
weights; replacing them by zero is a mean-field approximation, not an
exact inference.
\end{remark}


\section{A coercive screening response at the determinantal point}
\label{sec:screening-berezin-response}

This section proves the proposed screening sign in a precise class of
one-body deformations of the Bergman process.  The result is nonlinear,
but its application to an actual change of inverse temperature would
require an additional identification that is not established here.
Put $\delta(z)=1-|z|^2$ and $d\lambda=\delta^{-2}dA$, and let
$\mathbb P_2$ be the Bergman process, with kernel
$K(z,w)=\pi^{-1}(1-z\bar w)^{-2}$ relative to $dA$.

\subsection{An elementary strict contraction for the Berezin operator}

Define
\[
 (Bf)(z)=\int_\D b(z,w)f(w)\,d\lambda(w),\qquad
 b(z,w)=\frac{\delta(z)^2\delta(w)^2}{\pi|1-z\bar w|^4}.
\]
The kernel is symmetric, nonnegative, invariant under simultaneous disk
automorphisms, and has row integral one.  The last assertion follows
from the reproducing identity for $K$, or by integrating at $z=0$ and
then applying invariance.

\begin{proposition}[Coercive Bergman susceptibility]
\label{prop:screening-berezin-coercivity}
The operator $B$ extends to a positive self-adjoint operator on
$L^2(\lambda)$ and
\begin{equation}
 0\le B\le \frac\pi4 I.
 \label{eq:screening-berezin-contraction}
\end{equation}
For every real bounded compactly supported $f$,
\begin{equation}
 \operatorname{Var}_{\mathbb P_2}\Xi(f)
  =\frac1\pi\langle f,(I-B)f\rangle_{L^2(\lambda)},\qquad
 \frac{1-\pi/4}{\pi}\|f\|_{L^2(\lambda)}^2
 \le\operatorname{Var}_{\mathbb P_2}\Xi(f)
 \le\frac1\pi\|f\|_{L^2(\lambda)}^2.
 \label{eq:screening-berezin-variance}
\end{equation}
\end{proposition}

\begin{proof}
For $\xi\in\partial\D$, write
$P_\xi(z)=\delta(z)/|\xi-z|^2$.  For $0<s<1$, invariance and the
Poisson-kernel transformation identity give
\[
 B(P_\xi^s)(z)=c_sP_\xi(z)^s,\qquad
 c_s=\frac1\pi\int_\D P_1(w)^s\,dA(w).
\]
Indeed, under a disk automorphism taking $0$ to $z$, $P_\xi$ becomes
$P_\xi(z)P_\eta$ for some $\eta\in\partial\D$, while
$b(0,w)d\lambda(w)=dA(w)/\pi$.  Expanding
$(1-w)^{-s}$ and integrating in angle and radius, with nonnegative
terms after angular integration, gives
\[
 c_s=\frac1{s+1}\sum_{k=0}^\infty
       \frac{(s)_k^2}{k!(s+2)_k}
     =\Gamma(1+s)\Gamma(2-s).
\]
Here $(a)_k=a(a+1)\cdots(a+k-1)$, with $(a)_0=1$.  To verify the
sum without a summation formula, insert the beta integral for
$(s)_k/(s+2)_k$.  Tonelli's theorem and
$\sum_k(s)_kt^k/k!=(1-t)^{-s}$ reduce the sum to
$s\int_0^1t^{s-1}(1-t)^{1-s}\,dt$.
In particular, the strictly positive function $h=P_1^{1/2}$ satisfies
$Bh=(\pi/4)h$.

The weighted Schur estimate is elementary.  Cauchy--Schwarz gives
\[
 |Bf(z)|^2\le (Bh)(z)
       \int b(z,w)\frac{|f(w)|^2}{h(w)}\,d\lambda(w).
\]
Integrating, using symmetry and $Bh=(\pi/4)h$ a second time, proves
$\|Bf\|_2\le(\pi/4)\|f\|_2$.  All exchanges for nonnegative
integrands use Tonelli.  Symmetry therefore gives a bounded
self-adjoint extension.  Positivity first follows for real or complex
compactly supported $f$ by expanding
$|1-z\bar w|^{-4}$: the quadratic form is a sum of nonnegative
weighted squared moments.  Uniform absolute convergence on the
product of the compact supports justifies integration of the series.
Positivity then extends by density.

The determinantal two-point identity gives
\[
 \operatorname{Var}\Xi(f)=\int f(z)^2K(z,z)dA(z)
       -\iint f(z)f(w)|K(z,w)|^2dA(z)dA(w).
\]
Since $K(z,z)=\pi^{-1}\delta(z)^{-2}$, this is precisely
\eqref{eq:screening-berezin-variance}; the operator inequalities give
the two bounds.
\end{proof}

The constant function is not in $L^2(\lambda)$.  Thus the strict
$L^2$ contraction does not contradict $B1=1$, and it does not control
a change in the constant hyperbolic intensity.

\subsection{Regularized one-body tilts and their exact screening sign}

For real $V\in L^\infty(\lambda)\cap L^2(\lambda)$, define the centered
linear statistic $Y(V)$ as the $L^2(\mathbb P_2)$ limit of
\[
 \Xi(V\mathbf1_{B_r})-\frac1\pi
          \int_{B_r}V\,d\lambda,\qquad r\uparrow1.
\]
Proposition~\ref{prop:screening-berezin-coercivity} proves existence,
independence of the compact approximation in $L^2(\lambda)$, and
\begin{equation}
 \frac{1-\pi/4}{\pi}\|V\|_2^2
 \le\mathbb E_2Y(V)^2\le\frac1\pi\|V\|_2^2.
 \label{eq:screening-centered-variance}
\end{equation}
The map extends linearly and continuously from real $L^2(\lambda)$
to centered $L^2(\mathbb P_2)$ random variables.

\begin{lemma}[Exponential integrability and density defects]
\label{lem:screening-tilt-integrability}
For real $V\in L^\infty\cap L^2(\lambda)$ and every real $t$,
\begin{equation}
 \mathbb E_2e^{tY(V)}
 \le \exp\left\{\frac{t^2e^{|t|\|V\|_\infty}}{2\pi}
                          \|V\|_2^2\right\}.
 \label{eq:screening-centered-exponential}
\end{equation}
Consequently
\[
 \frac{d\mathbb P^V}{d\mathbb P_2}
   =\frac{e^{-2Y(V)}}{\mathbb E_2 e^{-2Y(V)}}
\]
defines a probability measure equivalent to $\mathbb P_2$.  It has
locally finite intensity $\rho_VdA$, and
\[
 r_V:=\delta^2\rho_V-\frac1\pi\in L^2(\lambda),\qquad
 \langle f,r_V\rangle=\mathbb E_{\mathbb P^V}Y(f)
 \quad(f\in L^2(\lambda)).
\]
\end{lemma}

\begin{proof}
For bounded compactly supported $f$, the determinantal generating
functional and $\log(1+x)\le x$ for $x>-1$ give
\[
 \log\mathbb E_2e^{t(\Xi(f)-\mathbb E_2\Xi(f))}
 \le\int_\D (e^{tf}-1-tf)K(z,z)\,dA(z)
 \le \frac{t^2 e^{|t|\|f\|_\infty}}{2\pi}\|f\|_2^2.
\]
For precision, restrict the kernel to a compact set containing the
support of $f$.  Its compression $K_C$ is a positive trace-class
contraction.  The determinant can be written
$\det(I+K_C^{1/2}M_{e^{tf}-1}K_C^{1/2})$.
Its operator inside the determinant is bounded below by
$\min(1,e^{-|t|\|f\|_\infty})I$, so the displayed logarithmic
inequality applies eigenvalue by eigenvalue, also when $f$ changes
sign.

Apply this bound to $f=V\mathbf1_{B_r}$.  An almost surely convergent
subsequence of the $L^2$-convergent centered statistics and Fatou's
lemma prove~\eqref{eq:screening-centered-exponential}.  Bounds with
an exponent larger than $|t|$ also give uniform integrability, so the
exponential expectations converge along the entire approximation.
The likelihood $D_V=d\mathbb P^V/d\mathbb P_2$ is strictly positive
and belongs to every finite $L^p(\mathbb P_2)$.

For compactly supported bounded $f$, Cauchy--Schwarz and the variance
upper bound give
\[
 |\mathbb E_2[(D_V-1)Y(f)]|
 \le\frac{\|D_V-1\|_{L^2(\mathbb P_2)}}{\sqrt\pi}
                                      \|f\|_{L^2(\lambda)}.
\]
The Riesz representation theorem supplies $r_V\in L^2(\lambda)$.
Applying the identity to indicators of relatively compact Borel sets
identifies the intensity as
$(\pi^{-1}+r_V)d\lambda$, which is nonnegative because it is an
intensity measure.  Approximation in $L^2(\lambda)$ proves the
identity for every $f$ in that space.
\end{proof}

\begin{theorem}[Strict nonlinear screening monotonicity]
\label{thm:screening-tilt-monotonicity}
For real $V,W\in L^\infty\cap L^2(\lambda)$,
\begin{equation}
 \langle V-W,r_V-r_W\rangle
 =-2\int_0^1\operatorname{Var}_{\mathbb P^{W+t(V-W)}}
                              Y(V-W)\,dt\le0.
 \label{eq:screening-nonlinear-monotonicity}
\end{equation}
Equality holds if and only if $V=W$ in $L^2(\lambda)$.  In particular,
$\langle V,r_V\rangle<0$ for every nonzero $V$ in this class.
Equivalently, the sum of the two finite relative entropies is
\begin{equation}
 D(\mathbb P^V\Vert\mathbb P^W)
 +D(\mathbb P^W\Vert\mathbb P^V)
       =-2\langle V-W,r_V-r_W\rangle.
 \label{eq:screening-symmetric-entropy}
\end{equation}
At $V=0$ the weak directional derivative is
\begin{equation}
 \left.\frac{d}{dt}\right|_{t=0}
      \langle f,r_{tV}\rangle
 =-\frac2\pi\langle f,(I-B)V\rangle,
 \qquad f\in L^2(\lambda).
 \label{eq:screening-linear-response}
\end{equation}
The operator $\frac2\pi(I-B)$ is bounded and invertible on
$L^2(\lambda)$, with inverse norm at most
$\pi/[2(1-\pi/4)]$.
\end{theorem}

\begin{proof}
The exponential bounds justify differentiation under expectation on
every compact $t$-interval, with all polynomial factors in the
statistics absorbed by exponential moments.  Differentiate
$\mathbb E_{\mathbb P^{W+t(V-W)}}Y(V-W)$ to obtain minus twice its
variance.  Integrating from zero to one gives
\eqref{eq:screening-nonlinear-monotonicity}.  If $V-W\ne0$, its
centered statistic has strictly positive variance under $\mathbb P_2$
by~\eqref{eq:screening-centered-variance}.  Equivalence of all the
tilted measures implies that the statistic remains nonconstant under
every tilt, so the variance integrand is strictly positive.
The logarithm of the likelihood ratio between the two tilts is
$-2Y(V-W)$ plus a deterministic normalization constant.  Integrating
it under both laws with opposite signs cancels that constant and
proves~\eqref{eq:screening-symmetric-entropy}.  The entropy integrals
are finite by the proved exponential bounds.

Differentiation at zero and polarization of
\eqref{eq:screening-berezin-variance} prove
\eqref{eq:screening-linear-response}.  For arbitrary $f\in L^2$,
no exponential moment of $Y(f)$ is needed: the exponential bounds for
$Y(V)$ imply that the tilted likelihood is differentiable in
$L^2(\mathbb P_2)$ at zero with derivative $-2Y(V)$, and
$Y(f)\in L^2(\mathbb P_2)$ can be paired with that derivative by
Cauchy--Schwarz.  The claimed
inverse bound follows from
\eqref{eq:screening-berezin-contraction} and the spectral theorem.
\end{proof}

\begin{corollary}[Uniqueness for self-consistent one-body screening]
\label{cor:screening-one-body-fixedpoint}
Let $G$ be the positive $L^2(\lambda)$ Green operator of
Proposition~\ref{prop:hyperbolic-green-operator-bound} and let
$j\in L^2(\lambda)$ be fixed.  The equation
\begin{equation}
 V=G(r_V+j),\qquad V\in L^\infty\cap L^2(\lambda),
 \label{eq:screening-one-body-fixedpoint}
\end{equation}
has at most one solution.  For $j=0$ its unique solution is $V=0$.
\end{corollary}

\begin{proof}
If $V,W$ are solutions, then
\[
 \langle V-W,r_V-r_W\rangle
 =\langle G(r_V-r_W),r_V-r_W\rangle\ge0.
\]
Strict monotonicity forces $V=W$.  For $j=0$, the zero function is
a solution because $\mathbb P^0=\mathbb P_2$ and $r_0=0$.
\end{proof}

\subsection{The precise limitation for the temperature conjecture}

These results establish the desired screening sign for an actual,
explicit family of interacting point processes, and a strict linear
response gap on the whole hyperbolic disk.  They do not assert that a
disk-gas limit at $\beta\ne2$ belongs to that family.  A temperature
change alters the pair interaction and cannot in general be represented
by a one-body tilt.  Moreover, a nonzero constant change of hyperbolic
intensity is outside $L^2(\lambda)$.  Applying
Corollary~\ref{cor:screening-one-body-fixedpoint} to the intensity
conjecture would require both a justified reference process at the new
temperature and a proved self-consistent one-body representation of its
defect, with the stated boundary selection.  Neither assertion follows
from the DLR kernels.

There is one rigorous connection to the previously computed temperature
responses.  If $u_{N,\beta}$ denotes the normalized finite mean defect
of~\eqref{eq:finite-intensity-mean-defect}, then, for $j=0,1,2$,
\begin{equation}
 \left.\partial_\beta^j u_{N,\beta}\right|_{\beta=2}
          \longrightarrow0\quad\hbox{locally in }C^2(\D).
 \label{eq:screening-mean-defect-second-order-flatness}
\end{equation}
Indeed $q_2=1$, $q'_2=-1$, $q''_2=1$, and the proved local uniform
zeroth, first, and second intensity responses are respectively
$1,-1,1$ times $[\pi\delta^2]^{-1}$.  Substitute these in the radial
integral formula for $u_{N,\beta}$.  The same first- and second-radial
derivative formulas used in
Proposition~\ref{prop:intensity-mean-potential-defect} prove local
$C^2$ convergence including the origin.  This differentiates before
taking $N\to\infty$ and supplies no remainder estimate on a nonzero
temperature interval.  In particular it does not imply the missing
fixed-temperature screening sign.


\subsection{A uniform response gap under bounded one-body tilts}
\label{subsec:screening-uniform-berezin-tilts}

The strict sign just obtained has a quantitative version on every set
of uniformly bounded potentials.  It also has an operator-theoretic
proof that identifies the point processes in the reference family.
For a real essentially bounded function $V$, put
\[
 \operatorname{osc}V=\operatorname*{ess\,sup}V-
                         \operatorname*{ess\,inf}V.
\]

\begin{proposition}[Weighted projections and a uniform variance gap]
\label{prop:screening-weighted-projection-gap}
Let $V\in L^\infty(\lambda)\cap L^2(\lambda)$ be real and let
$\mathbb P^V$ be the regularized tilt of
Lemma~\ref{lem:screening-tilt-integrability}.  On $L^2(\D,dA)$ let
$P$ be the Bergman projection and put $m=e^{-V}$.  The process
$\mathbb P^V$ is determinantal with orthogonal projection onto the
closed subspace $m\operatorname{Ran}P$.  More precisely, if
\[
 S=\left.Pm^2P\right|_{\operatorname{Ran}P},
 \qquad
 P_V=mP S^{-1}Pm,
\]
where the inverse acts on $\operatorname{Ran}P$, then $P_V$ is its
correlation operator.  For every real $f\in L^2(\lambda)$,
\begin{equation}
 e^{-2\operatorname{osc}V}\operatorname{Var}_{\mathbb P_2}Y(f)
 \le \operatorname{Var}_{\mathbb P^V}Y(f)
 \le e^{2\operatorname{osc}V}\operatorname{Var}_{\mathbb P_2}Y(f).
 \label{eq:screening-weighted-variance-comparison}
\end{equation}
Consequently,
\begin{equation}
 \operatorname{Var}_{\mathbb P^V}Y(f)
 \ge e^{-2\operatorname{osc}V}
          \frac{1-\pi/4}{\pi}\|f\|_{L^2(\lambda)}^2.
 \label{eq:screening-uniform-tilted-gap}
\end{equation}
\end{proposition}

\begin{proof}
Write $a=\operatorname*{ess\,inf}m>0$ and
$b=\operatorname*{ess\,sup}m<\infty$.  Then
$a^2 I\le S\le b^2 I$ on $\operatorname{Ran}P$, so $S$ is
invertible.  The operator $U=mP S^{-1/2}$, viewed as a map from
$\operatorname{Ran}P$ to $L^2(dA)$, is an isometry onto
$m\operatorname{Ran}P$.  Hence $P_V=UU^*$ is its orthogonal
projection.  The range is closed because multiplication by $m$ is
boundedly invertible.

We verify that this projection describes the probabilistic tilt.
First suppose that $V$ is compactly supported.  On
$\operatorname{Ran}P$ the operator $S-I=P(m^2-1)P$ is trace class:
write it as the difference of the two positive compressions obtained
from the positive and negative parts of $m^2-1$, each having trace
$\int (m^2-1)_\pm K(z,z)dA(z)<\infty$.  The determinantal
multiplicative-functional identity gives
\[
 \mathbb E_2 e^{-2\Xi(V)}=\det_{\operatorname{Ran}P} S.
\]
For a nonnegative bounded compactly supported function $h$, the same
identity applied with the additional factor $e^{-h}$ yields
\[
 \frac{\mathbb E_2 e^{-2\Xi(V)-\Xi(h)}}
      {\mathbb E_2 e^{-2\Xi(V)}}
 =\det_{\operatorname{Ran}P}
    \bigl(I+S^{-1/2}Pm(e^{-h}-1)mP S^{-1/2}\bigr).
\]
The equality follows by conjugating the trace-class perturbation of
$S$ by $S^{-1/2}$ and using multiplicativity of Fredholm
determinants for invertible trace-class perturbations of the identity.
The cyclic determinant identity, with the localized Hilbert--Schmidt
factors, turns this into
\[
 \det_{L^2(dA)}
  \bigl(I-\sqrt{1-e^{-h}}\,P_V\sqrt{1-e^{-h}}\bigr).
\]
The deterministic centering in $Y(V)$ cancels between the numerator
and denominator.  These are exactly the Laplace functionals of the
projection determinantal process with kernel $P_V$.

For general bounded $V\in L^2(\lambda)$, set
$V_R=V\mathbf1_{B_R}$ and $m_R=e^{-V_R}$.  All $m_R$ and their
inverses are bounded by common constants.  Multiplication by $m_R$
converges strongly to multiplication by $m$, by dominated
convergence.  Thus $S_R\to S$ strongly, and their inverses converge
strongly as well: this follows from
$S_R^{-1}-S^{-1}=S_R^{-1}(S-S_R)S^{-1}$ and their uniform norm bound.
If $k_z=K(\cdot,z)\in\operatorname{Ran}P$, the kernels of the
weighted projections are given almost everywhere by
\[
 K_{V_R}(z,w)=m_R(z)m_R(w)
       \langle S_R^{-1}k_w,k_z\rangle,
\]
with the inner-product convention chosen to agree with the displayed
Bergman kernel.  They converge almost everywhere to the analogous
kernel for $P_V$.  Moreover,
\[
 |K_{V_R}(z,w)|\le C
                \sqrt{K(z,z)K(w,w)}
\]
with $C$ independent of $R$.  On every relatively compact set the
correlation determinants therefore converge by dominated convergence;
Hadamard's inequality bounds their order-$p$ integrals by the $p$th
power of a fixed finite local trace.  The Fredholm expansion for each
compactly supported Laplace functional is consequently dominated by
an exponential series and converges to that of $P_V$.
On the other hand,
$e^{-2Y(V_R)}/\mathbb E_2 e^{-2Y(V_R)}$ converges in
$L^1(\mathbb P_2)$ to the likelihood of $\mathbb P^V$, by the
exponential bounds in Lemma~\ref{lem:screening-tilt-integrability}.
The two limits of the Laplace functionals coincide.  This proves the
claimed determinantal representation.  The kernel diagonal bound
also proves local trace class for the limiting projection.

Put $Q=I-P$ and $Q_V=I-P_V$.  For any $u\in L^2(dA)$, distance
to the two closed ranges gives
\begin{equation}
 a\|Qu\|_2\le \|Q_Vmu\|_2\le b\|Qu\|_2.
 \label{eq:screening-distance-comparison}
\end{equation}
Indeed the middle quantity is
$\inf_{v\in\operatorname{Ran}P}\|m(u-v)\|_2$, and multiplication
by $m$ changes every norm by a factor between $a$ and $b$.
For bounded compactly supported real $f$, the projection
variance identity and $UU^*=P_V$ imply
\[
 \operatorname{Var}_{\mathbb P^V}\Xi(f)
 =\|Q_V fP_V\|_{\mathrm{HS}}^2
 =\|Q_V fU\|_{\mathrm{HS}}^2.
\]
Since $f$ commutes with $m$, apply
\eqref{eq:screening-distance-comparison} to each column of
$fP S^{-1/2}$.  This bounds the last squared norm between
$a^2\|QfP S^{-1/2}\|_{\mathrm{HS}}^2$ and
$b^2\|QfP S^{-1/2}\|_{\mathrm{HS}}^2$.
The inequalities $b^{-2}I\le S^{-1}\le a^{-2}I$ then give
\[
 \frac{a^2}{b^2}\|QfP\|_{\mathrm{HS}}^2
 \le\operatorname{Var}_{\mathbb P^V}\Xi(f)
 \le\frac{b^2}{a^2}\|QfP\|_{\mathrm{HS}}^2.
\]
Here $a^2/b^2=e^{-2\operatorname{osc}V}$, and
$\|QfP\|_{\mathrm{HS}}^2$ is the Bergman variance.  Subtracting
any deterministic center leaves the variances unchanged.

Finally, approximate an arbitrary real $f\in L^2(\lambda)$ by
bounded compactly supported $f_j$ in $L^2(\lambda)$.  The variance
upper comparison, together with
$\mathbb E_{\mathbb P^V}Y(f_j)=\langle f_j,r_V\rangle$, makes
$Y(f_j)$ Cauchy in $L^2(\mathbb P^V)$.  Its limit agrees with
$Y(f)$: the convergence in $L^2(\mathbb P_2)$ implies convergence
in probability under $\mathbb P_2$, and absolute continuity implies
convergence in probability under $\mathbb P^V$ to the same random
variable.  Passing to the limit proves
\eqref{eq:screening-weighted-variance-comparison} for every such $f$.
The Bergman lower variance bound then proves
\eqref{eq:screening-uniform-tilted-gap}.
\end{proof}

\begin{corollary}[Strong monotonicity on bounded potential sets]
\label{cor:screening-uniform-monotonicity}
If $V,W\in L^\infty\cap L^2(\lambda)$ are real and
$\|V\|_\infty,\|W\|_\infty\le M$, then
\begin{equation}
 \langle V-W,r_V-r_W\rangle
 \le -\frac{2e^{-4M}(1-\pi/4)}{\pi}
                  \|V-W\|_{L^2(\lambda)}^2.
 \label{eq:screening-uniform-monotonicity}
\end{equation}
\end{corollary}

\begin{proof}
Every potential $W+t(V-W)$ on the joining segment has essential
oscillation at most $2M$.  Insert
\eqref{eq:screening-uniform-tilted-gap} into the exact integrated
variance identity~\eqref{eq:screening-nonlinear-monotonicity}.
\end{proof}

\begin{example}[A boundary ambiguity outside the square-integrable class]
\label{ex:screening-bounded-harmonic-gauge}
The weighted-projection construction itself makes sense for every
bounded real $V$, even when the centered exponential tilt is not
defined by the preceding $L^2$ construction.  In this larger class the
intensity response need not determine $V$.  Indeed, let $V=\Re h$ be
bounded and harmonic on the disk, with $h$ holomorphic.  Then
$a=e^{-h}$ and $a^{-1}$ are bounded analytic functions, so
$a\operatorname{Ran}P=\operatorname{Ran}P$.  Writing $m=e^{-V}$,
\[
 m\operatorname{Ran}P
 =e^{i\Im h}a\operatorname{Ran}P
 =e^{i\Im h}\operatorname{Ran}P.
\]
Consequently the weighted projection has kernel
\[
 K_V(z,w)=e^{i\Im h(z)}K(z,w)e^{-i\Im h(w)}.
\]
All its correlation determinants equal the Bergman determinants.
Thus the resulting point-process law is exactly $\mathbb P_2$, and
its intensity defect is zero, even for nonconstant $V$.  For example,
$V(z)=\Re z$ is bounded, has $V(0)=0$, and is nonzero, while
$-\delta^2\Delta V=0$.  Hence a local screening PDE alone allows
this nonzero potential with zero density defect.  This $V$ is outside
$L^2(\lambda)$, since its angular squared mean is $r^2/2$ and
$\int_0^1r^3(1-r^2)^{-2}dr=\infty$.
The example does not contradict the preceding uniqueness theorem,
which uses the $L^2$ Green operator and its boundary selection.
Nor does it provide a radial counterexample to the intensity
conjecture: a radial harmonic function regular at the origin is
constant, and the normalization there fixes that constant.
\end{example}


\section{Status of the uniqueness question}\label{sec:limits}
\begin{enumerate}
\item \emph{Uniqueness proved for $\beta\ge4$.}
Theorem~\ref{thm:empty-limit-beta-above-four} gives full-sequence
convergence to the empty configuration for $\beta>4$, locally in total
variation, with an explicit vanishing rate. The endpoint $4$ follows
from Theorem~\ref{thm:critical-empty}, without a rate. Every fixed
correlation and compact count moment vanishes throughout $\beta\ge4$
by Corollary~\ref{cor:critical-empty-consequences}; the explicit
bounds of Corollary~\ref{cor:empty-correlation-rates} apply above $4$. This is a statement on the
open disk: all $N$ particles remain in the finite system and leave
each fixed compact set as $N$ increases. A second, independent
proof, without a rate, follows from the Gaussian force regression
and curvature inequality in Section~\ref{sec:gaussian-force-curvature}.
\item \emph{Almost sure infiniteness proved below $2$.}
Theorem~\ref{thm:almost-sure-infinite-below-two} proves that every
subsequential limit is almost surely infinite for $0<\beta<2$.
The proof uses an exterior-only interpolation identity with a random
polynomial, and yields quantitative hole and lower-count bounds.
Theorems~\ref{thm:nonemptiness-below-two}
and~\ref{thm:nonemptiness-all-orders-below-two} give positive intensity
and factorial lower bounds. The limiting intensity dominates
$c_\beta/[\pi(1-|z|^2)^2]$, with $c_\beta=1$ for $\beta\le1$.
There is also a positive almost sure logarithmic lower growth rate
for boundary counts, and the sharp comparison upper bound
$\limsup_{r\uparrow1}(1-r)M(r)\le2/\beta$.
The stronger Theorem~\ref{thm:non-blaschke-limits} proves infinite
Blaschke mass almost surely, with quantitative Laplace decay and
logarithmic growth of accumulated radial depth. Consequently the
uncompensated Green potential diverges at every anchor, without an
invariance assumption, and uncompensated grand-canonical equations
cannot describe these actual limits.
At $\beta=2$ the full limit is the almost surely infinite Bergman process.
\item \emph{The range still unresolved here.}
Full-sequence uniqueness and hyperbolic invariance remain unproved
for $0<\beta<4$, $\beta\ne2$. Nonemptiness remains unresolved for
$2<\beta<4$ only. Theorem~\ref{thm:blaschke-boundary-moment-criterion}
reduces that nonemptiness question to one uniform boundary moment
estimate for a charged characteristic polynomial. By
Corollary~\ref{cor:blaschke-weaker-boundary-exponent}, any boundary
exponent strictly below $2$ suffices; such an estimate is still
an unproved hypothesis. No claim is made that $4$ is the exact threshold
for emptiness. Gaussian mode convergence and intensity lower bounds
do not by themselves identify the interior law.
\item \emph{What the additional identities settle.}
Fixed inserted charges shift the harmonic Gaussian means exactly as in
Theorem~\ref{thm:inserted-charge-screening}; this does not identify the
local conditional law. The full-density ratio theorem controls most
particle numbers, and its radial-depth corollary concerns a noncompact
observable. Neither assertion proves nonemptiness or uniqueness in the
remaining range. The boundary-statistic estimates require a power
strictly greater than one, whereas the centered exterior Green field
contains the first power of the boundary distance.
The proved local M\"obius covariance is $O(N^{-1})$; the conditional
symmetry criterion requires $o(N^{-1})$. These estimates cannot be
substituted for one another.
\item \emph{New constraints on the limiting law.}
The canonical logarithmic derivative in
Corollary~\ref{cor:canonical-force-score} is a measurable function of
the configuration and is uniquely determined by its integration-by-parts
identity. This uniqueness is a statement about the score of a given law,
not about the law itself. The force regression gives
\eqref{eq:intensity-log-curvature-lower} and monotonicity of
$(1-r^2)^{\beta/2}\rho_1(r)$. Assuming a nonempty critical limit gives
$(1-r^2)^2\rho_1(r)\to c$ with $0<c\le1/\pi$. The residual score
variance then vanishes under boundary rescaling; the resulting
deterministic score is incompatible with pair repulsion. This closes
the critical case by Theorems~\ref{thm:critical-empty}
and~\ref{thm:critical-mobius-empty}. Moreover,
Theorem~\ref{thm:finite-count-exclusion} excludes every finite positive
total count for every limiting law with $\beta\ge2$. Together with
the new infiniteness theorem below $2$, this means that every limit
at any $\beta>0$ has total count either zero or infinity.
\item \emph{What the polynomial representations reduce.}
Theorem~\ref{thm:schur-logarithmic-memory} shows that a compact portion
of the actual process is determined to vanishing error by a terminal
block of $O(\log N)$ Schur parameters. Their dependent law still needs
to be analyzed. The independent reference has a Bergman limit for
every parameter, but its Gram likelihood is asymptotically singular
for $\beta\ge4$; direct transfer of that limit is therefore impossible
in that range. At $\beta=4$, the explicit hypergeometric mixture in
Theorem~\ref{thm:critical-su11-hypergeometric} reduces fixed orbital
occupations to primary-sector masses. Critical emptiness implies decay
of every fixed primary-sector occupation, although no direct norm
asymptotic or quantitative critical rate is obtained.
\item \emph{A remaining sufficient criterion.}
Proposition~\ref{prop:transfer-uniqueness-criterion} gives an explicit
sufficient covariance estimate, \eqref{transfer:missing-covariance-bound}.
That estimate remains unproved. It is unnecessary for the proof above
$4$, and may be relevant to the remaining range. The deterministic
example in Proposition~\ref{prop:boundary-memory} explains why exterior
fields do not follow from vague convergence alone.
\item \emph{What is now proved about DLR equations.}
Theorem~\ref{thm:unconditional-canonical-quasigibbs} proves canonical
quasi-Gibbs densities for every subsequential limit at every positive
$\beta$, without the inverse-field variance hypothesis. The densities
are positive in every allowed count sector, real analytic after
division by the Vandermonde factor, and give local diffeomorphism
quasi-invariance. Every nonzero repulsion factor $h_p$ is strictly
positive even on the collision diagonals. Theorem~\ref{thm:conditional-gaussian-levi}
gives the stronger conditional Levi bound, which rules out a single
harmonic exterior potential and forces genuine posterior randomness
in the physical force marks.
Theorem~\ref{thm:consistent-augmented-canonical-specification} retains
this extra information as a meromorphic electric field and proves
consistent canonical DLR kernels on the augmented space.
Theorem~\ref{thm:corrected-field-gaussian-independence} further proves
that $(\Xi,\mathcal A)$ is independent of the Gaussian modes, where
$\mathcal A=\mathcal C+T_G$ on the nonempty branch. Exact Gaussian
Bayes then gives Theorem~\ref{thm:gaussian-free-green-specification}:
an explicit Green pair interaction and the one-body factor
$(1-|z|^2)^{-\beta/2}$, conditional on a retained exterior-field mark.
The positive-kernel factorization and conditional Gaussian force
convolution are stronger than the earlier covariance bounds.
At $\beta=2$,
Theorem~\ref{thm:bergman-grandcanonical-dlr} gives the complete
compensated DLR law; every finite conditional count has positive
probability. This theorem concerns the Bergman law and does not prove
uniqueness among all solutions of a specified infinite-volume equation.
For other $0<\beta<4$, a Green interaction is now identified on the
augmented space. Recovering its exterior field from the unmarked
exterior and determining the relative conditional count weights remain
unproved. Explicit finite-particle solutions show why the canonical
equations alone do not give uniqueness; the new almost sure
infiniteness theorem excludes those examples as actual gas limits.
The stronger Theorem~\ref{thm:infinite-green-canonical-nonuniqueness}
constructs distinct infinite, rotationally invariant position laws
satisfying the same canonical Green equations for every $0<\beta<4$.
For $0<\beta\le2$, their finite Blaschke mass makes them mutually
singular with every actual gas limit. Thus uniqueness of the canonical
specification fails, while uniqueness of the original limiting
sequence remains a separate question. Canonical-tail extremality and
position-changing decompositions are characterized in
Theorem~\ref{thm:canonical-tail-extremality}.
Proposition~\ref{prop:local-poisson-equivalence} shows that all
nontrivial actual candidates are nevertheless locally equivalent to
area Poisson; this does not establish equality of their laws.
\item \emph{What is now proved about hyperbolic symmetry.}
The score relative to hyperbolic area transforms as a complex one-form,
and the independent Gaussian force has the corresponding invariant law.
These statements do not establish invariance of the point process.
Such invariance is known here at $\beta=2$ and for the empty limits
at $\beta\ge4$. Theorem~\ref{thm:invariant-green-divergence} proves
that the uncompensated Green sum of every nonempty invariant locally
finite process diverges everywhere, without a moment hypothesis.
Consequently its positive-activity uncompensated DLR equation has no
invariant solution. Propositions~\ref{prop:conditional-green-background-bound}
and~\ref{prop:invariance-forces-zero-one-body-exponent} impose
necessary tests on a compensated description. The formula
$c=(4/\beta-1)/\pi$ remains conditional on the explicitly unproved
canonical-DLR and field hypotheses in
Section~\ref{subsec:conditional-hyperbolic-intensity}; it is not an
established intensity formula for the general gas. The new affine-field
criterion, Theorem~\ref{thm:affine-meromorphic-field-intensity}, removes
the need to construct an intrinsic centered potential from that conditional
argument, but still assumes symmetry of the field law. The exact finite-particle
mean defect and Palm dipole identities isolate the remaining expectations
without assuming this symmetry. The Gaussian-free
canonical kernels have an affine isometry covariance, which is distinct
from invariance of the selected augmented law.
\item \emph{Independent perturbative and spectral routes.}
Theorem~\ref{thm:beta2-all-response} identifies the limiting first
inverse-temperature response of every fixed correlation order at
$\beta=2$ and proves its full hyperbolic covariance. In particular,
$\partial_\beta\rho_{1,N}|_{\beta=2}$ converges locally uniformly to
$-1/[\pi(1-|z|^2)^2]$.
Theorem~\ref{thm:beta2-second-response} also bounds second derivatives
of the correlation logarithms uniformly in $N$ on compact sets of
distinct tuples, at $\beta=2$.
Theorems~\ref{thm:beta2-exact-second-center}
and~\ref{thm:beta2-exact-second-all-anchors} evaluate the second
logarithmic response as zero and the second intensity response as
$1/[\pi(1-|z|^2)^2]$, locally uniformly throughout the disk. The proofs
explicitly retain two finite-projection boundary terms, each converging
to $1/4$, which cancel at every coherent Palm anchor. A fixed nonzero perturbation
in $\beta$ still requires control on a temperature interval.
Theorem~\ref{thm:signed-charge-radial-exponent} identifies
$2/\beta-1/2$ as the radial-depth coefficient in the actual gas for
$0<\beta<2$, in all positive moments on a density-one set of particle
numbers, and proves signed-charge exponents throughout $t>-2$.
The same leading depth coefficient holds at every fixed hyperbolic center.
These logarithmic statements do not fix the order-one insertion prefactor
or the mean-potential defect that determines the limiting intensity. The global gas laws at any two
distinct temperatures become asymptotically singular, by
Theorem~\ref{thm:global-temperature-singularity}. Thus a direct global
change of temperature cannot have a uniformly integrable likelihood;
this does not preclude local absolute continuity or local convergence.
The sharp hyperbolic Green
operator bound is $\|G\|_{L^2(\lambda)}=2\pi$, but explicit invariant
random boundary modes show why this spatial contraction does not
prove uniqueness among invariant random fields. Under stated
integrability and covariance hypotheses, two renormalized fields over
the same configuration differ only by a tail-dependent constant; this
does not identify the configuration law or the activity.
Proposition~\ref{prop:all-correlation-monotonicity-uniqueness} gives a
different sufficient uniqueness criterion, whose monotonicity input
has not been proved here.
Proposition~\ref{prop:boundary-covariance-nonuniqueness} also gives an
exact necessary and sufficient oscillation test for multiple original
gas subsequences, expressed in finite partition-integral ratios.
Its required nonzero oscillation has not been proved. Fixed boundary
regularizations are screened, but the boundary-uniform estimate needed
to control adjacent particle numbers remains unproved.
\item \emph{What the screening and energy method establishes.}
The exact Euclidean-coordinate equation is
$\operatorname{div}(A\nabla U)=2\pi(c-\widetilde\rho)$, with
the anisotropic matrix $A$ specified in
Proposition~\ref{prop:screening-euclidean-elliptic-operator}.
For a radial finite-energy defect, its finite trace $\ell=U(\infty)$
exists and testing with $U-\ell$ removes the cutoff boundary term;
this yields the sufficient sign criterion in
Proposition~\ref{prop:screening-energy-trace-criterion}.
The finite-energy hypothesis and the required sign have not been
proved for the selected limits. Whole-plane finite-particle energy
is infinite, and even a neutral finite-volume mean-charge energy is
of order $N$ at $\beta=2$. Such global estimates cannot replace a
localized argument.
The exact Ward identity gives a different sufficient criterion,
Theorem~\ref{thm:screening-entropy-energy}, that requires no prior
finite total energy but retains connected-force work and a cutoff
boundary condition. The harmonic Gaussian modes do not remove this
connected force. The variational decomposition in
Proposition~\ref{prop:screening-correlation-energy-obstruction}
shows explicitly where the correlation energy enters.
For bounded square-integrable one-body Bergman tilts, strict
screening monotonicity and positivity of the Green operator prove
uniqueness of $V=G(r_V+j)$ for each prescribed source $j$.
The weighted-projection formula makes this sign quantitatively
coercive on uniformly bounded potential sets. Changing inverse
temperature changes the pair interaction, so this theorem supplies
no identification of the original general-temperature process.
Neither the full intensity conjecture nor its uniqueness question
is resolved by this revision.
\item \emph{The temperature assumption.}
At $\beta=0$ the positions are independent uniform points of the disk
and $M_N(r)$ is $\operatorname{Binomial}(N,r^2)$. Its mean is $Nr^2$
and variance at most $N$, so Chebyshev's inequality gives
$M_N(r)\to\infty$ in probability for every $r>0$. The fixed positive
$\beta$ hypothesis is therefore essential for unscaled local compactness.
\end{enumerate}

\begin{thebibliography}{9}
\raggedright
\bibitem{Samaj}
L. \v{S}amaj,
\emph{Counter-Ions Near a Charged Wall: Exact Results for Disc and Planar
Geometries}, Journal of Statistical Physics \textbf{161} (2015), 227--249.
\href{https://doi.org/10.1007/s10955-015-1308-8}{doi:10.1007/s10955-015-1308-8}.

\bibitem{EPW}
J. D. Esary, F. Proschan, and D. W. Walkup,
\emph{Association of Random Variables, with Applications},
Annals of Mathematical Statistics \textbf{38} (1967), 1466--1474.
\href{https://doi.org/10.1214/aoms/1177698701}{doi:10.1214/aoms/1177698701}.

\bibitem{Kallenberg}
O. Kallenberg,
\emph{Random Measures, Theory and Applications},
Probability Theory and Stochastic Modelling, vol.~77, Springer, 2017.
\href{https://doi.org/10.1007/978-3-319-41598-7}{doi:10.1007/978-3-319-41598-7}.

\bibitem{Strassen}
V. Strassen,
\emph{The Existence of Probability Measures with Given Marginals},
Annals of Mathematical Statistics \textbf{36} (1965), 423--439.
\href{https://doi.org/10.1214/aoms/1177700153}{doi:10.1214/aoms/1177700153}.

\bibitem{PV}
Y. Peres and B. Vir\'ag,
\emph{Zeros of the i.i.d. Gaussian power series: a conformally invariant
determinantal process}, Acta Mathematica \textbf{194} (2005), 1--35.
\href{https://doi.org/10.1007/BF02392515}{doi:10.1007/BF02392515}.

\bibitem{BGZNW}
R. Butez, D. Garc\'ia-Zelada, A. Nishry, and A. Wennman,
\emph{Universality of Outliers in Weakly Confined Coulomb-Type Systems},
Communications in Mathematical Physics \textbf{406} (2025), article 127.
\href{https://doi.org/10.1007/s00220-025-05293-7}{doi:10.1007/s00220-025-05293-7}.
\bibitem{KillipKozhan}
R. Killip and R. Kozhan,
\emph{Matrix models and eigenvalue statistics for truncations of classical
ensembles of random unitary matrices},
Communications in Mathematical Physics \textbf{349} (2017), 991--1027.
\href{https://arxiv.org/abs/1501.05160}{arXiv:1501.05160}.


\bibitem{BufetovBergman}
A. I. Bufetov,
\emph{The conditional measures for the determinantal point process with
the Bergman kernel}, Rendiconti Lincei -- Matematica e Applicazioni
\textbf{34} (2023), 159--173.
\href{https://doi.org/10.4171/RLM/1002}{doi:10.4171/RLM/1002}.


\bibitem{Hedenmalm2015}
H. Hedenmalm,
\emph{Coulomb gas ensembles in 2D}, lecture slides, December 11, 2015.
\url{https://irma.math.unistra.fr/~klevtsov/past/Hedenmalm.pdf}.

\end{thebibliography}
\end{document}
